\documentclass[11pt,openany]{book}
\setcounter{secnumdepth}{3}
\usepackage[margin=1.1in]{geometry}
\usepackage[T1]{fontenc}
\usepackage[utf8]{inputenc}
\usepackage{lmodern}
\usepackage[brazilian]{babel}
\AtBeginDocument{\shorthandoff{"}} % keep " literal (TikZ labels, code)
\usepackage{amsmath,amssymb,amsthm}
\usepackage{graphicx}
\usepackage{tikz}
\usetikzlibrary{arrows.meta,decorations.markings,decorations.pathmorphing,decorations.pathreplacing,calc}
\input{figures/icons}
\usepackage[unicode,colorlinks=true,linkcolor=blue!60!black,citecolor=blue!60!black]{hyperref}
\usepackage{microtype}
\setlength{\emergencystretch}{3em} % long names (licences, references) in narrow lines
\usepackage{empheq}
\usepackage{array}
\usepackage{booktabs}
\usepackage{multirow}
\usepackage{longtable}
\usepackage{circuitikz}
\definecolor{keyyellow}{rgb}{1.0,0.97,0.78}
% key equations: \begin{empheq}[box=\keybox]{equation} ... \end{empheq} -- light-yellow box, number stays outside
\newcommand{\keybox}[1]{{\setlength{\fboxsep}{7pt}\colorbox{keyyellow}{#1}}}
\usepackage[most]{tcolorbox}
% key statements (propositions etc.): the same light-yellow background, breakable across pages
\newtcolorbox{keyblock}{enhanced,breakable,colback=keyyellow,colframe=keyyellow,boxrule=0pt,arc=0pt,left=6pt,right=6pt,top=2pt,bottom=2pt,boxsep=0pt}
\renewcommand{\chaptermark}[1]{\markboth{\small\chaptername~\thechapter.\ #1}{}}
\renewcommand{\sectionmark}[1]{\markright{\small\thesection.\ #1}}
\renewcommand{\topfraction}{0.95}
\renewcommand{\textfraction}{0.05}
\renewcommand{\floatpagefraction}{0.75}

% part pages: title at the top third, and text may follow on the same page (used for the part introductions)
\makeatletter
\renewcommand\part{\if@openright\cleardoublepage\else\clearpage\fi\thispagestyle{plain}\if@twocolumn\onecolumn\@tempswatrue\else\@tempswafalse\fi\null\vspace{0.18\textheight}\secdef\@part\@spart}
\renewcommand\@endpart{\par\vspace{3em}\if@tempswa\twocolumn\fi}
\makeatother
\newtheorem{theorem}{Teorema}[chapter]
\newtheorem{proposition}{Proposição}[chapter]
\newtheorem{lemma}{Lema}[chapter]
\theoremstyle{remark}
\newtheorem{remark}{Observação}[chapter]
% exercises: \begin{exercise}[\lvlB] ... \end{exercise}, numbered "Exercise 4.3"; difficulty marks
\theoremstyle{definition}
\newtheorem{exercise}{Exercício}[chapter]
\newcommand{\lvlA}{$\star$}
\newcommand{\lvlB}{$\star\star$}
\newcommand{\lvlC}{$\star\star\star$}
% answer to an exercise in the Answers appendix: \answer{ex:NN:k}{text}
\newcommand{\answer}[2]{\par\noindent\textbf{\ref{#1}.}~#2\par\smallskip}
\newcommand{\dd}{\mathrm{d}}
% neuron type code: first four letters of the primary mediator
\newcommand{\nt}[1]{\textsf{\textbf{#1}}}
\newcommand{\mV}{\,\text{mV}}
\newcommand{\ms}{\,\text{ms}}
\newcommand{\mM}{\,\text{mM}}
\newcommand{\um}{\,\text{\textmu m}}
\newcommand{\ion}[2]{\mathrm{#1}^{#2}}
\newcommand{\Na}{\ion{Na}{+}}
\newcommand{\K}{\ion{K}{+}}
\newcommand{\Cl}{\ion{Cl}{-}}
\newcommand{\Ca}{\ion{Ca}{2+}}
\newcommand{\conc}[1]{[\mathrm{#1}]}

% textbook look: chapter openers, headings, running heads, figure labels (figures/booklook.tex)
\input{figures/booklook}

\title{\Huge O computador de 20 watts\\[16pt] \Large Como o cérebro calcula: uma abordagem quantitativa para engenheiros}
\hypersetup{pdftitle={O computador de 20 watts. Como o cérebro calcula: uma abordagem quantitativa para engenheiros},pdfauthor={Konstantin Kladko}}
\author{\Large Konstantin~Kladko}
% cover: a salad of the book's figures around the title card (figures/bookcover.tex)
\newcommand{\covertitle}{O computador de 20 watts}
\newcommand{\coversubtitle}{Como o cérebro calcula: uma abordagem quantitativa para engenheiros}
\newcommand{\coverauthor}{Konstantin~Kladko}
\date{}

\begin{document}
\frontmatter
\input{figures/bookcover}
\thispagestyle{empty}
\vspace*{\fill}
\noindent{\small \copyright~2026 Konstantin~Kladko.\\[4pt]
O texto, as figuras, os exercícios e as respostas são distribuídos sob a licença Creative Commons Atribuição-NãoComercial-CompartilhaIgual 4.0 Internacional (CC BY-NC-SA 4.0): \url{https://creativecommons.org/licenses/by-nc-sa/4.0/}. O livro pode ser copiado, usado no ensino, traduzido e adaptado livremente para fins não comerciais, com crédito ao autor e distribuição das obras derivadas sob a mesma licença.\\[4pt]
Os programas dos modelos e de compilação são distribuídos sob a licença MIT.\\[4pt]
Fontes, modelos e novas versões: \url{https://github.com/kladkogex/20-watt-computer}.}
\clearpage

\tableofcontents
\chapter*{Como ler este livro}
\addcontentsline{toc}{chapter}{Como ler este livro}

O livro não precisa ser lido em ordem. Os capítulos~1--4 são a base comum: os tipos de neurônios, o que calculam os \nt{GLUT} e os \nt{GABA} e em que linguagem eles se comunicam. Depois disso o livro se ramifica. O mapa da fig.~\ref{fig:roadmap} mostra quais capítulos se apoiam em quais: uma seta do capítulo $A$ ao capítulo $B$ significa que $B$ usa conceitos e fórmulas de $A$. O mapa mostra apenas as dependências principais; pequenas referências a capítulos anteriores não atrapalham a leitura.

\begin{figure}[h]
\centering
\begin{tikzpicture}[>=Stealth,scale=0.95,
  ch/.style={draw,thick,rounded corners=3pt,align=center,font=\scriptsize,text width=1.85cm,minimum height=0.62cm,inner sep=2pt},
  base/.style={ch,fill=blue!8}, bus/.style={ch,fill=orange!14}, sum/.style={ch,fill=gray!18},
  io/.style={ch,fill=green!10}, sort/.style={ch,fill=cyan!8}, lab/.style={ch,fill=violet!10},
  dep/.style={->,thick,gray!60!black}]
% foundation
\node[base] (c1) at (0,0) {1 Tipos de neurônios};
\node[base] (c2) at (2.3,0) {2 \nt{GLUT}};
\node[base] (c3) at (4.6,0) {3 \nt{GLUT} e \nt{GABA}};
\node[base] (c4) at (6.9,0) {4 Código Morse};
% broadcast channels and the synthesis
\node[bus] (c5) at (4.6,-1.5) {5 \nt{SERO}};
\node[bus] (c6) at (7.3,-1.5) {6 \nt{DOPA}};
\node[bus] (c7) at (9.2,-0.9) {7 Teorias da dopamina};
\node[bus] (c8) at (9.2,-2.1) {8 \nt{NORA}};
\node[bus] (c9) at (11.5,-2.1) {9 \nt{ACET}};
\node[sum] (c10) at (13.8,-0.9) {10 A máquina inteira};
% inputs and outputs
\node[io] (c12) at (2.3,-4.1) {12 Reconhe-\\cimento};
\node[io] (c11) at (4.6,-4.1) {11 Entradas};
\node[io] (c13) at (6.9,-4.1) {13 Músculos};
\node[io] (c14) at (9.2,-3.05) {14 Dor};
\node[io] (c15) at (9.2,-4.1) {15 Aprender\\movimentos};
\node[io] (c16) at (9.2,-5.15) {16 Química\\do corpo};
% lab, sorting, sleep
\node[lab] (c17) at (11.5,-4.1) {17 Depuração};
\node[lab] (c21) at (13.8,-4.1) {21 Máquinas vivas};
\node[lab] (c22) at (13.8,-5.15) {22 DishBrain};
\node[sort] (c18) at (11.5,-5.15) {18 Neurônios-\\instrumentos};
\node[sort] (c19) at (13.8,-6.2) {19 Número de\\dendritos};
\node[bus] (c20) at (11.5,-6.2) {20 Sono};
% dependencies
\draw[dep] (c1) -- (c2); \draw[dep] (c2) -- (c3); \draw[dep] (c3) -- (c4);
\draw[dep] (c1) -- (c5); \draw[dep] (c4) -- (c6); \draw[dep] (c5) -- (c6);
\draw[dep] (c6) -- (c7); \draw[dep] (c6) -- (c8); \draw[dep] (c8) -- (c9);
\draw[dep] (c7) -- (c10); \draw[dep] (c9) -- (c10);
\draw[dep] (c4) -- (c11); \draw[dep] (c11) -- (c12); \draw[dep] (c11) -- (c13); \draw[dep] (c5) -- (c13);
\draw[dep] (c13) -- (c14); \draw[dep] (c13) -- (c15); \draw[dep] (c13) -- (c16);
\draw[dep] (c15) -- (c17); \draw[dep] (c9) -- (c17); \draw[dep] (c17) -- (c21); \draw[dep] (c21) -- (c22);
\draw[dep] (c16) -- (c18); \draw[dep] (c18) -- (c19); \draw[dep] (c16) -- (c20);
% legend
\begin{scope}[yshift=-7.25cm]
\foreach \s/\t [count=\i] in {base/base, bus/canais e modos, sum/síntese, io/entradas e saídas, sort/classificações de neurônios, lab/laboratório}{
  \pgfmathtruncatemacro{\col}{mod(\i-1,3)}\pgfmathtruncatemacro{\row}{(\i-1)/3}
  \node[\s,text width=0.3cm,minimum height=0.32cm] at ({\col*4.8+0.2},{-\row*0.55}) {};
  \node[anchor=west,font=\scriptsize] at ({\col*4.8+0.55},{-\row*0.55}) {\t};}
\end{scope}
\end{tikzpicture}
\caption{Mapa dos capítulos. A seta vai de um capítulo para aquele que se apoia nele. As demais referências entre capítulos são indicações, não dependências.}
\label{fig:roadmap}
\end{figure}

\section*{Caminhos pelo livro}

\begin{center}
\footnotesize
\setlength{\tabcolsep}{4pt}
\renewcommand{\arraystretch}{1.2}
\begin{tabular}{>{\raggedright}p{3.0cm}>{\raggedright}p{3.9cm}>{\raggedright\arraybackslash}p{6.4cm}}
\toprule
Caminho & Para quem & Capítulos em ordem \\
\midrule
curso de um trimestre & professor, 10 semanas & semana 1: cap.~1--2; 2: cap.~3; 3: cap.~4; 4: cap.~5--6; 5: cap.~7--8; 6: cap.~9; 7: cap.~10; 8: cap.~11--12; 9: cap.~13 e~15; 10: cap.~17 \\
IA e aprendizado de máquina & quem conhece redes neurais e quer entender como o cérebro aprende & 1--4, 5 (leitura rápida), 6, 7, 8, 9, 11 (leitura rápida), 12, 13 (leitura rápida), 15 \\
eletrônica e hardware & projetista de circuitos, desenvolvedor de chips neuromórficos & 1--4, 5, 11, 13, 16, 18, 19, 17 (capítulos 9 e 15 em leitura rápida) \\
neurointerfaces e computadores vivos & quem constrói interfaces cérebro-computador e chips com neurônios vivos & 1, 2, 4, 11, 13, 17 (capítulos 9 e 15 em leitura rápida), 21, 22 (capítulo 12 em leitura rápida) \\
visão geral rápida & engenheiro com um fim de semana livre & 1, 2, 3, 6, 10; as referências do capítulo~6 aos capítulos~4--5 se entendem pelo contexto \\
\bottomrule
\end{tabular}
\end{center}

``Leitura rápida'' significa: basta ler o início do capítulo, as proposições nos quadros amarelos e a tabela final.

No fim de cada capítulo há exercícios: estimativas, deduções, modelagem e projeto; as respostas curtas estão reunidas no apêndice~\ref{sec:answers}. Toda a notação está no apêndice~\ref{sec:notation}, e os experimentos que sustentam as principais afirmações de cada capítulo estão no apêndice~\ref{sec:supplement}. Os números do livro vêm de pequenos programas na pasta \texttt{models/}; eles podem ser executados e modificados.

\mainmatter

\chapter{Tipos de neurônios}
\label{sec:neurontypes}

\section{O neurônio como caixa-preta}

Suponha que lhe entreguem um saco com oitenta e seis bilhões de neurônios~\cite{Azevedo2009} e peçam para separá-los em pilhas. Como você faria isso?

Um engenheiro que recebe um saco de microchips desconhecidos não os separaria pelo formato do encapsulamento. Ele perguntaria: quantas entradas o componente tem, quantas saídas e o que ele entrega na saída. Vamos desenhar o neurônio da mesma forma, como um bloco num diagrama (fig.~\ref{fig:blackbox}).

\begin{figure}[t]
\centering
\begin{tikzpicture}[>=Stealth,x=1cm,y=1cm]
% the box
\node[draw,very thick,minimum width=2.6cm,minimum height=2.6cm,fill=blue!6] (N) at (0,0) {};
\node at (0,0.55) {\small neurônio};
\node at (0,-0.15) {$\displaystyle u=\sum_i w_i x_i$};
\node at (0,-0.8) {\small $y=\Theta(u-\theta)$};
% inputs
\foreach \k/\y/\lab in {1/1.05/{x_1},2/0.55/{x_2},3/0.05/{x_3}}{
  \draw[->,thick,blue!60!black] (-4.2,\y) -- (-1.3,\y);
  \node[left] at (-4.2,\y) {$\lab$};
  \node[above,blue!60!black] at (-2.1,\y-0.06) {\scriptsize $w_{\k}$};}
\node at (-4.45,-0.4) {$\vdots$};
\node at (-2.1,-0.4) {$\vdots$};
\draw[->,thick,blue!60!black] (-4.2,-1.05) -- (-1.3,-1.05);
\node[left] at (-4.2,-1.05) {$x_n$};
\node[above,blue!60!black] at (-2.1,-1.11) {\scriptsize $w_{n}$};
\draw[decorate,decoration={brace,amplitude=5pt},gray!60!black] (-4.9,-1.2) -- (-4.9,1.2);
\node[left,align=right,gray!40!black] at (-5.1,0) {\small $n\sim10^3$--$10^5$\\[-2pt]\small entradas};
% single output wire, then fan-out
\draw[very thick,red!70!black] (1.3,0) -- (3.2,0);
\foreach \x in {1.6,2.2,2.34,2.9}{\draw[thick,red!70!black] (\x,0) -- (\x,0.4);}
\node[above right,font=\tiny\itshape,gray!30!black,inner sep=1pt] at (1.42,0.45) {clique\dots\ clique-clique\dots};
\node[below,red!70!black,align=center] at (2.25,-0.05) {\scriptsize uma saída $y$};
\fill[red!70!black] (3.2,0) circle (0.06);
\foreach \y in {1.3,0.65,0,-0.65,-1.3}{
  \draw[->,thick,red!70!black] (3.2,0) -- (5.0,\y);}
\draw[decorate,decoration={brace,amplitude=5pt},gray!60!black] (5.3,1.45) -- (5.3,-1.45);
\node[right,align=left,gray!40!black] at (5.5,0) {\small cópias de $y$\\[-2pt]\small para $10^3$--$10^4$\\[-2pt]\small outros neurônios};
\end{tikzpicture}
\caption{O neurônio visto por um engenheiro. Muitas entradas $x_i$, cada uma com seu peso $w_i$; dentro, a soma ponderada $u$ e a comparação com o limiar $\theta$ ($\Theta$ é o degrau: $1$ se o argumento é positivo e $0$ caso contrário); uma única saída $y$, uma sequência de disparos multiplicada por ramificação para milhares de destinatários.}
\label{fig:blackbox}
\end{figure}

Na saída do bloco não há um número, e sim uma sequência de pulsos curtos de tensão: ``clique'', pausa, ``clique-clique'', pausa longa. Todos os pulsos têm a mesma altura, então toda a informação está em \emph{quando} eles ocorrem. Dentro do bloco, em primeira aproximação grosseira, há um somador e um comparador: as entradas são somadas com pesos e, se a soma passa do limiar, o bloco emite um disparo. No próximo capítulo vamos substituí-lo por um modelo que funciona de fato em tempo contínuo.

A saída é uma só, mas o fio se ramifica, e cada ramo leva \emph{a mesma} cópia do sinal. Em eletrônica isso se chama fan-out, o leque de saída. Num neurônio do córtex ele é da ordem de alguns milhares; uma porta lógica num chip costuma alimentar poucas outras portas.

\emph{O que}, então, o fio de saída transmite ao destinatário? O disparo corre até a ponta de cada ramo e ali vira um sinal químico: o ramo libera, na fenda estreita entre as células, uma pitada de uma substância específica, o \emph{neurotransmissor}, que altera a corrente de entrada do destinatário. Eis o critério de classificação: \emph{qual neurotransmissor a saída do neurônio libera}.

\section{Uma saída, um tipo de sinal}

A classificação só faz sentido se cada neurônio tiver um único tipo de saída. Será assim? A experiência diz que quase.

\begin{keyblock}
\begin{proposition}[Um neurônio, um tipo de saída]
\label{prop:dale}
Quase todo neurônio tem um neurotransmissor principal e o libera em todos os ramos de sua saída. Por isso o tipo do neurônio é definido pelo seu neurotransmissor principal~\cite{Strata1999}.
\end{proposition}
\end{keyblock}

Vamos combinar que o tipo de um neurônio é indicado \emph{pelas quatro primeiras letras do nome em inglês de seu neurotransmissor principal}. Assim, um neurônio cujo neurotransmissor principal é o glutamato é um neurônio \nt{GLUT}. Todos os tipos estão reunidos na tabela~\ref{tab:types}.

\begin{table}[t]
\centering
\footnotesize
\setlength{\tabcolsep}{3pt}
\renewcommand{\arraystretch}{1.15}
\begin{tabular}{>{\raggedright}p{1.0cm}>{\raggedright}p{2.05cm}>{\raggedright}p{1.75cm}>{\raggedright}p{1.6cm}>{\centering}p{0.7cm}>{\raggedright}p{1.55cm}>{\raggedright}p{1.8cm}>{\raggedright\arraybackslash}p{3.2cm}}
\toprule
Tipo & Neurotrans\-missor principal & Barramento & Velocidade & Sinal & Neurônios no cérebro & Saídas por neurônio & Papel no cálculo \\
\midrule
\nt{GLUT} & glutamato & de endereços & rápido e lento & $+$ & $\sim8\cdot10^{10}$ & $\sim10^4$ & principal peso positivo \\
\nt{GABA} & GABA & de endereços & rápido e lento & $-$ & $\sim3\cdot10^{9}$ & $10^3$--$10^4$ & principal peso negativo \\
\nt{GLYC} & glicina & de endereços & rápido & $-$ & $10^6$--$10^7$ & $10^2$--$10^3$ & peso negativo na medula espinal e no tronco encefálico; segunda chave de um dos receptores de glutamato \\
\nt{ACET} & acetilcolina & ambos & rápido e lento & $\pm$ & $\sim10^6$ & músculo: $10$--$10^3$; cérebro: $\sim10^5$ & saída para o músculo; no cérebro, taxa de aprendizado (hipótese) \\
\nt{DOPA} & dopamina & de difusão & lento & $\pm$ & $\sim5\cdot10^{5}$ & $10^5$--$10^6$ & erro de predição de recompensa $\delta$ \\
\nt{SERO} & serotonina & de difusão & lento (um tipo de receptor é rápido) & $\pm$ & $\sim3\cdot10^{5}$ & $\sim10^5$ & horizonte de planejamento (hipótese) \\
\nt{HIST} & histamina & de difusão & lento & $+$ & $\sim6\cdot10^{4}$ & sem estimativa & sinal global ``fique acordado'' \\
\nt{NORA} & noradrenalina & de difusão & lento & $\pm$ & $\sim5\cdot10^{4}$ & $\sim10^5$ & ganho (hipótese) \\
\nt{ADRE} & adrenalina & de difusão & lento & $\pm$ & $\sim10^4$ & sem estimativa & poucos neurônios; papel no cérebro incerto \\
\nt{ASPA} & aspartato & de endereços & rápido & $+$ & sem estimativa & sem estimativa & peso positivo fraco; papel controverso \\
\bottomrule
\end{tabular}
\caption{Tipos de neurônios, em ordem decrescente do número de neurônios no cérebro. \emph{Barramento}: de endereços, o neurotransmissor é liberado numa fenda estreita para um destinatário definido; de difusão, vai de um pequeno grupo de neurônios-fonte para grandes regiões (seção~\ref{sec:buses}). \emph{Velocidade}: rápido, por receptores ionotrópicos, milissegundos; lento, por metabotrópicos, centenas de milissegundos ou mais. \emph{Sinal}: o habitual; quem o fixa em última instância é o destinatário (seção~\ref{sec:receiver}), e $\pm$ significa que ocorrem ambos os sinais. \emph{Neurônios no cérebro}: quantos neurônios desse tipo há em todo o cérebro humano, em ordem de grandeza; no total são cerca de $8.6\cdot10^{10}$ neurônios~\cite{Azevedo2009}. Quase todos os neurônios \nt{GLUT}, cerca de $7\cdot10^{10}$, são os pequenos neurônios de entrada do cerebelo; os números de \nt{GLUT} e \nt{GABA} vêm dessa contagem e da fração de neurônios \nt{GABA} no córtex~\cite{Hendry1987}. Dos tipos pouco numerosos, só foram contados diretamente os \nt{HIST}~\cite{Airaksinen1991} e o principal dos dois grupos de neurônios \nt{DOPA}~\cite{Pakkenberg1991}, de modo que para \nt{DOPA} o valor é um limite inferior; para \nt{SERO}, é a soma de duas contagens parciais~\cite{Baker1990,Baker1991}. Os números de \nt{NORA}, \nt{GLYC}, \nt{ACET} e \nt{ADRE} são estimativas grosseiras de ordem de grandeza, sem contagem direta. \emph{Saídas por neurônio}: quantos terminais sinápticos um neurônio tem, isto é, pontos de liberação do neurotransmissor nos ramos do fio de saída (fig.~\ref{fig:blackbox}), em ordem de grandeza. Nos tipos de difusão os terminais não foram contados diretamente: para \nt{DOPA} o número foi deduzido da fração do alvo coberta pela ramificação de um único fio de saída~\cite{Matsuda2009}; para os demais, as estimativas são ainda mais grosseiras. Para \nt{ACET} há dois casos: o neurônio que comanda um músculo e o neurônio de difusão no cérebro.}
\label{tab:types}
\end{table}

% the pie is a table-type float with a figure caption, so it can never float ahead of Table~\ref{tab:types}
\begin{table}[tb]
\makeatletter\def\@captype{figure}\makeatother
\centering
\definecolor{vizglut}{HTML}{2A78D6}
\definecolor{vizgaba}{HTML}{E34948}
\definecolor{vizrest}{HTML}{8A8986}
\begin{tikzpicture}[>=Stealth]
% pie: GLUT 96.4 %, GABA 3.6 % (13.0 degrees), the rest 0.006 % (a hairline at 0 degrees)
\fill[vizglut] (0,0) -- (13:1.9) arc (13:360:1.9) -- cycle;
\fill[vizgaba] (0,0) -- (0:1.9) arc (0:13:1.9) -- cycle;
\draw[white,line width=1pt] (0,0) -- (13:1.9);
\draw[vizrest,line width=1.2pt] (0,0) -- (0:1.9);
\node[white,align=center,font=\small] at (190:0.95) {\nt{GLUT}\\$96.4\,\%$};
\draw[gray!60!black] (6.5:1.75) -- (2.05,2.1);
\node[anchor=west,font=\small] at (2.05,2.1) {\nt{GABA} $3.6\,\%$};
% zoom box for the remaining types
\draw[densely dashed,vizrest] (1.9,0) -- (3.1,1.65);
\draw[densely dashed,vizrest] (1.9,0) -- (3.1,-2.2);
\draw[vizrest,rounded corners=3pt] (3.1,-2.2) rectangle (10.1,1.65);
\node[anchor=west,font=\scriptsize] at (3.2,1.4) {todos os demais juntos: $\approx0.006\,\%$; escala logarítmica};
\foreach \code/\lg/\y/\val/\lx in {GLYC/6.477/1.0/{$10^6$--$10^7$}/4.05,ACET/6.0/0.6/{$\sim10^6$}/3.0,DOPA/5.699/0.2/{$\sim5\cdot10^5$}/2.699,SERO/5.477/-0.2/{$\sim3\cdot10^5$}/2.477,HIST/4.778/-0.6/{$\sim6\cdot10^4$}/1.778,NORA/4.699/-1.0/{$\sim5\cdot10^4$}/1.699,ADRE/4.0/-1.4/{$\sim10^4$}/1.0}{
  \node[anchor=east,font=\scriptsize] at (4.35,\y) {\nt{\code}};
  \fill[vizrest] (4.45,\y-0.13) rectangle ({4.45+\lg-3},\y+0.13);
  \node[anchor=west,font=\scriptsize] at ({4.5+\lx},\y) {\val};}
\draw[vizrest,thick] (7.45,1.0) -- (8.45,1.0);
\draw[vizrest,thick] (7.45,0.92) -- (7.45,1.08);
\draw[vizrest,thick] (8.45,0.92) -- (8.45,1.08);
\draw[gray!60!black] (4.45,-1.72) -- (8.45,-1.72);
\foreach \e in {3,4,5,6,7}{
  \draw[gray!60!black] ({4.45+\e-3},-1.72) -- ({4.45+\e-3},-1.66);
  \node[below,font=\scriptsize] at ({4.45+\e-3},-1.72) {$10^{\e}$};}
\end{tikzpicture}
\caption{Proporções dos tipos de neurônios no cérebro segundo as estimativas da tabela~\ref{tab:types}. O círculo é quase todo \nt{GLUT}; \nt{GABA} é um setor estreito, e todos os outros tipos juntos são um fio de cabelo na borda. À direita, esse fio de cabelo ampliado, em escala logarítmica do número de neurônios; para \nt{GLYC} a barra vai até o meio da faixa $10^6$--$10^7$, e o segmento mostra a faixa inteira. \nt{ASPA} não tem estimativa e não aparece na figura.}
\label{fig:typepie}
\end{table}

Quão desiguais são essas pilhas se vê na fig.~\ref{fig:typepie}: de cada mil neurônios, $964$ são \nt{GLUT}, $36$ são \nt{GABA}, e todos os outros tipos juntos somam cerca de seis neurônios em cem mil.

O GABA já tem um nome de quatro letras. As substâncias que quase nunca são o neurotransmissor principal (neuropeptídeos, gases, lipídios, purinas) não recebem código próprio; elas estão no apêndice~\ref{sec:othermediators}. A palavra ``quase'' na proposição~\ref{prop:dale} é importante. Muitos neurônios liberam, junto com o neurotransmissor principal, substâncias auxiliares~\cite{Yao2023}. Alguns liberam também um segundo neurotransmissor rápido, até de sinal oposto: parte dos neurônios \nt{DOPA} libera também GABA~\cite{Tritsch2012}. E em alguns neurônios neurotransmissores diferentes saem de ramos diferentes da mesma saída~\cite{Zhang2015}. Tudo isso foi medido, então ``um neurônio, um neurotransmissor'' é uma regra com exceções, não uma lei. Mas como primeira divisão ela resiste ao teste: no atlas celular do cérebro do camundongo, em que os neurônios são separados pelos genes ativos em mais de cinco mil grupos, eles continuam se dividindo em grandes classes antes de tudo pelo neurotransmissor principal~\cite{Yao2023}.

Essa regra tem uma consequência que distingue de imediato o cérebro de uma rede neural artificial. Numa rede artificial, o mesmo neurônio pode ter peso positivo para um destinatário e negativo para outro: os pesos $w_{ij}$ são apenas números numa matriz. No cérebro, o sinal da ação é determinado principalmente pelo neurotransmissor, e o neurônio tem um só. Logo, se o neurônio $j$ excita um destinatário, ele excita todos os outros. Toda a \emph{coluna} da matriz de pesos que sai do neurônio $j$ tem o mesmo sinal:
\begin{empheq}[box=\keybox]{equation}
\operatorname{sign} w_{ij} \;=\; s_j \quad\text{para todos os destinatários } i, \qquad s_j\in\{+1,-1\}.
\label{eq:dalesign}
\end{empheq}
Os neurônios se dividem em excitatórios ($s_j=+1$) e inibitórios ($s_j=-1$), e isso é propriedade do próprio neurônio, não da conexão. Se a rede precisa de uma conexão negativa a partir de um neurônio excitatório, ela tem de passar por um intermediário: o neurônio excitatório excita um inibitório, e este inibe o alvo: um peso negativo com atraso de um elo.

Conhecem-se mais de cem neurotransmissores, mas quase todo o trabalho do cérebro se apoia em meia dúzia deles, e eles se dividem em duas classes completamente diferentes.

\section{Dois barramentos: de endereços e de difusão}
\label{sec:buses}

Nas redes de computadores há duas maneiras de entregar uma mensagem. A primeira é a \emph{endereçada}, unicast: o pacote vai de um remetente para um destinatário definido, rápido, e ninguém mais o vê. A segunda é a \emph{de difusão}, broadcast: o remetente envia uma mensagem de uma vez a todos os nós da rede. Os neurônios usam as duas, e os neurotransmissores se dividem exatamente por esse critério (fig.~\ref{fig:phone_radio}).

\begin{figure}[t]
\centering
\begin{tikzpicture}[>=Stealth,scale=0.95]
% ---- unicast: point-to-point wiring
\node[above] at (2.0,3.3) {\small \textbf{De endereços}: glutamato e GABA};
\foreach \i in {0,...,4}{
  \fill[blue!60!black] (0.2+0.9*\i,2.5) circle (0.1);
  \fill[blue!60!black] (0.2+0.9*\i,0.6) circle (0.1);}
\foreach \a/\b in {0/0,0/1,1/1,1/3,2/2,3/2,3/4,4/4,2/0}{
  \draw[->,blue!60!black] (0.2+0.9*\a,2.38) -- (0.2+0.9*\b,0.72);}
\fill[red!70!black] (1.1,1.55) circle (0.09);
\pic[scale=1.2] at (1.75,1.9) {letter};
\draw[->,red!70!black,thick] (1.1,1.55) -- (0.28,0.7);
\draw[->,red!70!black,thick] (1.1,1.55) -- (1.95,0.72);
\node[align=center,below] at (2.0,0.2) {\small bilhões de neurônios,\\[-2pt]\small cada um com seus destinatários,\\[-2pt]\small tempo: milissegundos};
% ---- broadcast: one small source, global targets
\begin{scope}[xshift=7.2cm]
\node[above] at (1.8,3.3) {\small \textbf{De difusão}: dopamina, \dots};
\foreach \i in {0,...,8}{\fill[blue!60!black] (-0.2+0.5*\i,2.75) circle (0.08);}
\fill[orange!85!black] (1.8,0.35) ellipse (0.35 and 0.18);
\pic[scale=1.5] at (2.7,0.75) {mast};
\foreach \x in {-0.2,0.3,0.8,1.3,1.8,2.3,2.8,3.3,3.8}{
  \draw[->,orange!85!black] (1.8,0.5) .. controls (1.8,1.3) and (\x,1.6) .. (\x,2.62);}
\node[align=center,below] at (1.8,0.1) {\small de milhares a centenas de milhares de neurônios,\\[-2pt]\small uma mensagem para todos,\\[-2pt]\small tempo: centenas de milissegundos ou mais};
\end{scope}
\end{tikzpicture}
\caption{Duas maneiras de transmitir. À esquerda, o barramento de endereços, parecido com o correio com o endereço no envelope; em vermelho, um neurônio e dois de seus destinatários. À direita, o de difusão, como uma estação de rádio: um pequeno grupo de neurônios-fonte, e todos recebem a mesma mensagem.}
\label{fig:phone_radio}
\end{figure}

O \emph{barramento de endereços} é o glutamato e o GABA. Eles são liberados pela esmagadora maioria dos neurônios. Cada saída leva a células definidas, o neurotransmissor age só na fenda estreita do contato, e age rápido: em um milissegundo abrem-se os canais iônicos do destinatário, em poucos milissegundos tudo termina. É o barramento de \emph{dados}: por ele passa aquilo que o cérebro calcula.

O \emph{barramento de difusão} é a dopamina, a serotonina, a noradrenalina e em parte a acetilcolina. Eles são liberados por grupos minúsculos de neurônios, mas a saída de cada um se ramifica de forma incrivelmente ampla. Muitas vezes o neurotransmissor não é liberado numa fenda estreita, e sim no espaço extracelular, onde se espalha e age sobre todos os que estão por perto. Age devagar: não abre canais diretamente, mas dispara dentro do destinatário uma cadeia de reações químicas. Por ele não passam dados, e sim \emph{parâmetros de controle} comuns a grandes regiões da rede, que mudam seu modo de funcionamento: o ganho, a taxa de aprendizado, o sinal ``isso foi bom''. Por isso esses neurotransmissores são chamados de \emph{moduladores}. Por muito tempo se pensou que cada canal desses fosse um único número para o cérebro inteiro. As medições modernas refinaram o quadro: o canal não é um fio só, mas um pequeno feixe de linhas paralelas. Os neurônios-fonte de um mesmo neurotransmissor se dividem em subsistemas, cada um com seus destinatários e sua mensagem, e a quantidade liberada no local é ajustada também por sinais locais~\cite{Ren2018,Poe2020}. O número dessas linhas é de unidades ou dezenas, não de bilhões, então o canal continua sendo de difusão.

Se você precisa de uma única imagem deste capítulo, é esta. O cérebro é uma enorme rede de transmissão endereçada de dados, sobre a qual pairam alguns canais de difusão com sinais de controle. Um programador reconhecerá a arquitetura familiar: o cálculo mais um pequeno conjunto de variáveis de controle, cada uma comum a uma grande região da rede.

\section{\nt{GLUT}: peso positivo}

O glutamato é o principal neurotransmissor excitatório do cérebro. Quando a saída de um neurônio libera glutamato, abrem-se no destinatário canais pelos quais o sódio entra, a tensão na membrana do destinatário sobe e aumenta a chance de ele emitir seu próprio disparo. Na linguagem da fig.~\ref{fig:blackbox}, é uma entrada com peso \emph{positivo}, $w_i>0$.

Quem libera glutamato? Quase todos os que transmitem dados a longa distância: de três quartos a quatro quintos dos neurônios do córtex e a maioria dos neurônios que ligam diferentes regiões do cérebro. A rede do cérebro é antes de tudo uma rede de conexões glutamatérgicas.

Um detalhe de engenharia. Após a liberação, a concentração de glutamato na fenda salta milhares de vezes, até cerca de um milimolar, e cai com constante de tempo de cerca de um milissegundo~\cite{Clements1992}: o glutamato se difunde para fora da fenda, e proteínas-bomba especiais o recolhem de volta para as células. Para quê? Pelo mesmo motivo pelo qual, numa linha de comunicação, o sinal volta a zero depois de cada símbolo. Se o neurotransmissor não for removido, o próximo disparo chegará a uma linha ``ocupada'', e disparos separados por poucos milissegundos se fundirão.

\section{\nt{GABA}: peso negativo}

Imagine uma rede em que todos os pesos são positivos. Se, em média, cada disparo gera mais de um disparo seguinte, a atividade cresce exponencialmente e em poucos passos toma a rede inteira. Um eletrônico reconhecerá um laço de realimentação positiva com ganho maior que um: o circuito satura. Para evitar isso, são necessários pesos negativos. Sua principal fonte é o ácido gama-aminobutírico, abreviado GABA.

O GABA abre no destinatário canais que deixam passar íons cloreto. O potencial de equilíbrio do cloreto fica perto do potencial de repouso ou um pouco abaixo, por isso os canais de cloreto abertos mantêm a membrana perto do repouso e impedem que a tensão suba até o limiar. É uma entrada com peso \emph{negativo}, $w_i<0$. Os neurônios GABAérgicos são de um quinto a um quarto dos neurônios do córtex~\cite{Hendry1987}. Suas saídas são curtas, e eles inibem seus vizinhos mais próximos.

A inibição no cérebro faz muito mais do que uma simples subtração. Ela funciona como uma realimentação negativa que mantém toda a rede no ponto de operação. Considera-se que, num córtex funcionando normalmente, as correntes excitatória e inibitória que chegam a um neurônio se equilibram quase exatamente: esse regime foi previsto pela teoria~\cite{vanVreeswijk1996} e aparece nas medições de correntes no córtex vivo~\cite{Haider2006}, e o neurônio fica o tempo todo perto do limiar. Um circuito estabilizado por forte realimentação negativa em torno de um ponto instável reage rápido e com sensibilidade a sinais pequenos: um velho truque de engenharia.

\section{\nt{DOPA}: sinal de erro}

O primeiro canal de difusão é a dopamina. Os neurônios dopaminérgicos no ser humano são da ordem de meio milhão, cerca de um em cento e setenta mil; esse é o número contado no principal de seus dois grupos~\cite{Pakkenberg1991}, então é um limite inferior. Suas saídas se espalham pelas regiões que escolhem as ações.

O que esse canal transmite? A resposta vem de experimentos em que se registram os disparos de neurônios dopaminérgicos num animal que recebe uma recompensa~\cite{Schultz1997}. De vez em quando, o animal recebe inesperadamente uma gota de suco, e os neurônios dopaminérgicos respondem com uma rajada curta de disparos. Depois, antes do suco, passa-se a acender uma lâmpada, sempre um segundo antes do suco. Quando o animal aprende essa associação, os neurônios passam a responder com uma rajada à \emph{lâmpada}, e ao próprio suco, que chega na hora certa, não respondem de modo algum. E se depois da lâmpada o suco não vem, no momento em que deveria chegar os neurônios \emph{silenciam}, e sua frequência cai abaixo da habitual.

A regra que explica os três casos (fig.~\ref{fig:rpe}): os neurônios não informam quão bom algo é, mas quanto \emph{melhor ou pior do que o esperado}.

\begin{figure}[t]
\centering
\begin{tikzpicture}[>=Stealth,x=3.4cm,y=1cm]
\foreach \y/\lab in {2.8/{suco inesperado},1.4/{suco após a lâmpada},0/{lâmpada, mas sem suco}}{
  \draw[gray!70!black] (0,\y) -- (3,\y);
  \draw[densely dotted,gray!60!black] (0,\y+0.3) -- (3,\y+0.3);
  \node[left,align=right,font=\small] at (-0.03,\y+0.35) {\lab};}
% row 1: juice without a cue -> burst at the juice
\draw[very thick,orange!85!black] plot[domain=0:3,samples=120] (\x,{2.8+0.3+0.75*exp(-((\x-2.08)/0.07)^2)});
\pic[scale=0.8] at (2,2.58) {juice};
% row 2: cue then juice -> burst at the cue, nothing at the juice
\draw[very thick,orange!85!black] plot[domain=0:3,samples=120] (\x,{1.4+0.3+0.75*exp(-((\x-1.08)/0.07)^2)});
\pic[scale=0.85] at (1,1.15) {lamp}; \pic[scale=0.85] at (2,1.15) {juice};
% row 3: cue, juice omitted -> burst at the cue, dip when the juice was due
\draw[very thick,orange!85!black] plot[domain=0:3,samples=120] (\x,{0.3+0.75*exp(-((\x-1.08)/0.07)^2)-0.28*exp(-((\x-2.1)/0.12)^2)});
\pic[scale=0.85] at (1,-0.25) {lamp}; \pic[scale=0.85] at (2,-0.25) {nojuice};
\node[right,font=\scriptsize] at (2.36,0.1) {queda};
\node[right,font=\scriptsize] at (2.13,3.75) {surpresa!};
\node[left,font=\scriptsize] at (0.97,2.25) {pico na lâmpada};
\node[above,font=\scriptsize] at (2,1.75) {nada};
\draw[->,gray!70!black] (0,-0.75) -- (3.05,-0.75) node[right,black,font=\small] {$t$};
\draw[gray!60!black] (1,-0.8) -- (1,-0.7) node[below=2pt,font=\scriptsize] {lâmpada, $1\,\text{s}$};
\draw[gray!60!black] (2,-0.8) -- (2,-0.7) node[below=2pt,font=\scriptsize] {suco, $2\,\text{s}$};
\end{tikzpicture}
\caption{Frequência de disparos dos neurônios \nt{DOPA} em três experimentos (esquemático, segundo~\cite{Schultz1997}). O pontilhado é a frequência de fundo constante. A lâmpada é o sinal, a gota é o suco, a gota riscada é o suco prometido mas não entregue. A resposta é o erro de predição~\eqref{eq:rpe}.}
\label{fig:rpe}
\end{figure}

Um programador que conhece o aprendizado por reforço~\cite{SuttonBarto2018} reconhece essa grandeza de imediato. Um agente que aprende a escolher ações guarda uma estimativa $V(s)$: quanta recompensa ele espera estando no estado $s$. Depois de cada passo, ele compara a expectativa com o que recebeu:
\begin{empheq}[box=\keybox]{equation}
\delta_t \;=\; r_t \;+\; \gamma\,V(s_{t+1}) \;-\; V(s_t) ,
\label{eq:rpe}
\end{empheq}
e corrige a estimativa: $V(s_t)\leftarrow V(s_t)+\alpha\,\delta_t$. Aqui $r_t$ é a recompensa recebida, $\gamma<1$ é o fator de desconto (quanto uma recompensa futura vale menos que uma imediata), $\alpha$ é a taxa de aprendizado. A grandeza $\delta_t$ se chama \emph{erro de diferença temporal}.

Ela se comporta exatamente como os neurônios dopaminérgicos. Suco inesperado: $V$ ainda é zero, $r>0$, $\delta>0$. Suco aprendido: a lâmpada previu a recompensa, $V(s_t)$ antes do suco já vale $r$, e $\delta=0$. O suco não veio: $V(s_t)=r$, mas $r_t=0$, $\delta<0$. E o pico se desloca para a lâmpada porque é justamente no momento da lâmpada que a expectativa sobe de repente: $\gamma V(s_{t+1})-V(s_t)>0$. Os passos $t,t+1,\dots$ são aqui uma convenção do algoritmo: no organismo não há um relógio que marque os passos, e a mesma grandeza em tempo contínuo será obtida no capítulo~\ref{sec:dofa}.

Portanto, a dopamina é a difusão de um sinal de erro escalar $\delta$ para todos os que devem aprender com ele. Em primeira aproximação, um fio para o cérebro inteiro e um número a cada instante; o quanto esse fio é na verdade um feixe é discutido no capítulo~\ref{sec:dofaalt}.

\section{Os demais canais de difusão}

Para os outros moduladores há apenas hipóteses de trabalho, que traduzem seus papéis para a mesma linguagem de parâmetros globais~\cite{Doya2002}. Trate-as exatamente como hipóteses.

\emph{\nt{NORA}, noradrenalina: o ganho.} Os neurônios noradrenérgicos são ainda menos numerosos que os dopaminérgicos, por estimativa grosseira algumas dezenas de milhares, e suas saídas se espalham por praticamente todo o cérebro. Eles se ativam bruscamente quando acontece algo inesperado ou importante. A noradrenalina muda a inclinação da característica de transferência dos neurônios destinatários: entradas fracas ficam mais fracas, e fortes, mais fortes. Isso se mede assim: os disparos de fundo do destinatário são suprimidos mais do que a resposta à entrada, e a relação sinal-ruído cresce~\cite{Foote1975NE}. Num modelo simples, isso se descreve por um único número: se o neurônio responde com frequência $f=F\bigl(g\cdot u\bigr)$ à corrente de entrada $u$, a noradrenalina aumenta o coeficiente $g$~\cite{AstonJones2005} (mais detalhes no capítulo~\ref{sec:nora}). Uma rede com $g$ grande trabalha com mais ``contraste'' e decide mais rápido; com $g$ pequeno, de modo ``borrado'', e testa mais alternativas, como um agente de aprendizado por reforço em modo de exploração.

\emph{\nt{ACET}, acetilcolina: a taxa de aprendizado.} A acetilcolina controla quanto a rede escuta a entrada atual e quanto escuta seus próprios dados armazenados. Com nível alto de acetilcolina, as entradas dos órgãos dos sentidos são reforçadas, as conexões internas, de ``lembrança'', são enfraquecidas e, ao mesmo tempo, fica mais fácil alterar os pesos~\cite{Hasselmo2006}. Grosso modo, é uma chave ``gravação/leitura'' e, junto com ela, a taxa de aprendizado $\alpha$.

\emph{\nt{SERO}, serotonina: o horizonte de planejamento.} O que a serotonina transmite é o que menos se entende. Uma das hipóteses a liga ao fator de desconto $\gamma$ de~\eqref{eq:rpe}: um nível alto de serotonina corresponde à disposição de esperar por uma recompensa grande em vez de agarrar uma pequena agora. Os neurônios serotoninérgicos trabalham de forma regular e lenta, alguns disparos por segundo na vigília, e quase silenciam numa das fases do sono~\cite{JacobsAzmitia1992}.

Os seis neurotransmissores estão reunidos na tabela~\ref{tab:transmitters}.

\begin{table}[t]
\centering
\small
\begin{tabular}{>{\raggedright}p{2.4cm}>{\raggedright}p{3.0cm}>{\raggedright}p{2.0cm}>{\raggedright\arraybackslash}p{5.2cm}}
\toprule
Tipo de neurônio & Fração dos neurônios & Barramento & Papel no cálculo \\
\midrule
\nt{GLUT} (glutamato) & a maioria & de endereços & peso positivo, $w>0$ \\
\nt{GABA} (GABA) & 20--25\% no córtex & de endereços & peso negativo, $w<0$; estabilização \\
\nt{DOPA} (dopamina) & $\sim 6\cdot10^{-6}$ & de difusão & erro de predição $\delta$ \\
\nt{NORA} (noradrenalina) & $\sim 6\cdot10^{-7}$ & de difusão & ganho $g$ (hipótese) \\
\nt{ACET} (acetilcolina) & pequena & ambos & taxa de aprendizado $\alpha$; ``gravação/leitura'' (hipótese) \\
\nt{SERO} (serotonina) & $\sim 3\cdot10^{-6}$ & de difusão & desconto $\gamma$ (hipótese) \\
\bottomrule
\end{tabular}
\caption{Os principais neurotransmissores do cérebro na linguagem do cálculo. A coluna da direita é uma forte simplificação: confiável para glutamato, GABA e dopamina; para os demais são hipóteses.}
\label{tab:transmitters}
\end{table}

\section{O sinal é definido pelo destinatário}
\label{sec:receiver}

Eu os enganei um pouco ao dizer: glutamato é peso positivo, GABA é negativo. Nenhuma molécula carrega um sinal em si. A molécula é apenas uma chave. O que acontece quando ela é capturada é decidido pela \emph{fechadura}: o receptor na célula destinatária.

O melhor exemplo é a dopamina. Ela tem receptores de duas famílias~\cite{Beaulieu2011}. Em uns, facilita a excitação da célula; em outros, a dificulta. Na região que escolhe as ações, alguns neurônios têm receptores do primeiro tipo, outros do segundo, e o mesmo sinal de dopamina reforça um grupo e enfraquece o outro ao mesmo tempo. É como um único pacote de difusão que diferentes nós da rede interpretam de modos diferentes.

\begin{keyblock}
\begin{proposition}[O sentido da mensagem é definido pelo destinatário]
\label{prop:receptor}
O sinal e a velocidade de ação de um neurotransmissor não são determinados pelo próprio neurotransmissor, mas pelo receptor ao qual ele se liga. Há dois tipos de receptores. Os \emph{ionotrópicos} são eles próprios canais iônicos e se abrem em frações de milissegundo: é o controle direto da corrente, como a porta de um transistor. Os \emph{metabotrópicos} disparam dentro da célula uma cadeia de reações químicas e agem em centenas de milissegundos ou mais: é antes uma escrita num registrador de configuração, que muda o funcionamento posterior da célula. O barramento de endereços fala principalmente pelos primeiros; o de difusão, principalmente pelos segundos.
\end{proposition}
\end{keyblock}

Por que, então, se pode dizer ``o glutamato excita''? Porque quase todos os receptores de glutamato encontrados na prática são excitatórios, e quase todos os receptores de GABA são inibitórios. Não é uma lei da natureza, mas uma convenção muito confiável, e a regra~\eqref{eq:dalesign} se apoia justamente nela.

\section{O que levar deste capítulo}

A esmagadora maioria dos neurônios forma o barramento de endereços de dados, com pesos de um só sinal por neurônio; uma minoria ínfima forma os canais de difusão, que transmitem a grandes regiões do cérebro alguns sinais de controle comuns. Cerca de dois neurônios em cada cem mil, e deles depende como aprende e em que modo funciona todo o resto da rede. Para um engenheiro isso não surpreende: em qualquer computador há incomparavelmente menos registradores de controle do que células de memória, e mesmo assim a máquina não funciona sem eles.

No próximo capítulo vamos verificar o que pode calcular uma rede formada só por neurônios \nt{GLUT}, o tipo mais numeroso, funcionando em tempo contínuo.

\section{Exercícios}

\begin{exercise}[\lvlA]
\label{ex:01:1}
\emph{Pilhas e fios.} Pela tabela~\ref{tab:types} (meio da faixa em escala logarítmica; \nt{ACET}: $10^5$ saídas): (a)~se cada neurônio fosse uma pessoa, quantas populações da Terra formariam os neurônios \nt{GLUT} e de que tamanho seria a cidade de todos os de difusão? (b)~Quantas vezes a fração de terminais de \nt{DOPA}, \nt{SERO}, \nt{NORA}, \nt{ACET} supera a fração de seus neurônios?
\end{exercise}

\begin{exercise}[\lvlA]
\label{ex:01:2}
\emph{Retorno a zero.} O glutamato na fenda cai como $c_0e^{-t/\tau_c}$, $c_0=1\mM$, $\tau_c=1\ms$; o destinatário percebe concentrações acima de $10$~\textmu M. (a)~Por quanto tempo a linha fica ``ocupada'' após uma liberação? (b)~As bombas estão parcialmente bloqueadas, $\tau_c=10\ms$. Até que frequência de liberações a linha ainda consegue se desocupar?
\end{exercise}

\begin{exercise}[\lvlB]
\label{ex:01:3}
\emph{Para onde vai o pico.} No experimento da fig.~\ref{fig:rpe}, depois da lâmpada vêm os estados $s_0\to\dots\to s_{K-1}$ com passo $\Delta$, $T=K\Delta$; a recompensa $r$ chega na última transição, e o momento da lâmpada é imprevisível. Encontre os $V(s_k)$ aprendidos e a resposta $\delta$ à lâmpada. Tomando $\gamma=e^{-\Delta/\tau_\gamma}$, encontre $\tau_\gamma$ para $T=1$~s, $\Delta=0.1$~s, $\gamma=0.95$.
\end{exercise}

\begin{exercise}[\lvlB]
\label{ex:01:4}
\emph{Modelagem: aprendizado pelo erro.} Simule a cadeia do exercício~\ref{ex:01:3} com $K=10$, $\gamma=0.95$, $r=1$ e a regra $V(s_t)\leftarrow V(s_t)+\alpha\delta_t$. Em que tentativa $n^*$ a resposta à lâmpada supera pela primeira vez a resposta ao suco, para $\alpha=0.05$, $0.1$, $0.2$, $0.5$? Como $n^*$ depende de $\alpha$?
\end{exercise}

\begin{exercise}[\lvlB]
\label{ex:01:5}
\emph{Projeto: difundir um número.} O número $\delta$ precisa chegar a todos os $10^{10}$ neurônios; cada destinatário, incluindo os retransmissores, deve receber $10$ terminais da camada anterior, e cada elo acrescenta $1\ms$. Quantos neurônios e elos tem a árvore de retransmissores mais econômica com $10^4$ saídas por neurônio (\nt{GLUT}) e com $10^{5.5}$ (\nt{DOPA})?
\end{exercise}

\begin{exercise}[\lvlB]
\label{ex:01:6}
\emph{Por que o equilíbrio é sensível.} Numa rede com $80\,\%$ \nt{GLUT} e $20\,\%$ \nt{GABA}, todos ligados a todos com pesos $J_E/N$ e $-J_I/N$; $\tau\dot r=-r+Wr+h$. Com $J_E=10$, o único autovalor não nulo de $W$ é $0.5$. Encontre a razão entre a corrente inibitória e a excitatória e em quantos por cento basta enfraquecer a inibição para a rede entrar em avalanche.
\end{exercise}

\begin{exercise}[\lvlC]
\label{ex:01:7}
\emph{Pesquisa: \nt{SERO} é $\gamma$?} Pelo exercício~\ref{ex:01:3}, a resposta dos neurônios \nt{DOPA} à lâmpada é $re^{-T/\tau_\gamma}$. Atrasos $T=1,\dots,5$~s, com $n$ apresentações cada; $\ln\delta$ é medido com dispersão $0.5$. Quanto deve ser $n$ para distinguir $\tau_\gamma=2$~s de $2.4$~s (nível $0.05$, poder $0.8$)? Que mecanismo, além de $\gamma$, daria a mesma queda com $T$?
\end{exercise}

\chapter{O que os neurônios \nt{GLUT} calculam}
\label{sec:glutcomputer}

\section{As regras do jogo}

No capítulo~\ref{sec:neurontypes} vimos que a esmagadora maioria dos neurônios do cérebro é \nt{GLUT}: eles só excitam, seus pesos são só positivos. Tomemos quantos neurônios desses quisermos e perguntemos: o que essa rede pode calcular se nos permitirmos apenas o que existe num organismo vivo?

Num organismo vivo não há gerador de clock que comute todos os elementos ao mesmo tempo. Cada neurônio trabalha por conta própria, em tempo contínuo: os disparos chegam e saem quando calha, com atrasos diferentes. Há \emph{entradas}, neurônios sensoriais que disparam com luz, som ou pressão, e \emph{saídas}, neurônios que comandam os músculos. Entre elas está a rede, e ninguém lhe diz quando começar e quando terminar o cálculo. Para um eletrônico, é um circuito \emph{assíncrono} com elementos cujos atrasos não são conhecidos com precisão de antemão.

Tomemos o modelo de neurônio mais simples entre os que funcionam de fato em tempo contínuo. A tensão $V$ na membrana do neurônio em repouso é zero. Cada disparo que chega do neurônio $j$ a eleva de um salto $w_j\ge0$, e depois a tensão escoa sozinha de volta a zero com constante de tempo $\tau$:
\begin{empheq}[box=\keybox]{equation}
\tau\,\frac{\dd V}{\dd t} \;=\; -V \;+\; \tau\sum_j w_j \sum_k \delta\bigl(t-t_j^{(k)}-d_j\bigr), \qquad w_j\ge 0 .
\label{eq:glutneuron}
\end{empheq}
Aqui $t_j^{(k)}$ são os instantes dos disparos do neurônio $j$, $d_j$ é o atraso no caminho a partir dele, $\delta$ é a função delta, que representa justamente o salto instantâneo. Quando $V$ atinge o limiar $\theta$, o neurônio emite um disparo, $V$ volta a zero e, por cerca de um milissegundo, o neurônio não pode disparar de novo. Números típicos: $\tau$ vai de frações de milissegundo a vinte milissegundos em diferentes neurônios; os atrasos, de frações de milissegundo a dezenas de milissegundos. É o \emph{neurônio integrador com vazamento}: um circuito RC com limiar. É uma simplificação consciente: nos neurônios do córtex, os ramos de entrada emitem eles mesmos disparos locais~\cite{Smith2013}, e um neurônio se comporta mais como uma pequena rede de várias camadas (capítulo~\ref{sec:synthesis}). Para as questões deste capítulo, um único circuito RC basta.

A principal diferença em relação a uma porta lógica é o vazamento: o neurônio não lembra um disparo para sempre, mas por cerca de $\tau$, e só se soma o que chegou dentro dessa janela.

\section{OU e E}

Comecemos pelas portas. Se um único disparo eleva a tensão acima do limiar, $w\ge\theta$, o neurônio dispara a cada disparo de qualquer entrada. É o \emph{OU}.

O E é mais interessante. Seja $w<\theta<2w$: um disparo é pouco, são necessários dois. Mas agora a pergunta ``chegaram os dois?'' não tem sentido enquanto não se diz \emph{dentro de que intervalo de tempo}. Suponha que o primeiro disparo chegou no instante $0$ e o segundo, depois de um tempo $\Delta$. Quando o segundo chega, resta do primeiro $we^{-\Delta/\tau}$, e o neurônio dispara se $we^{-\Delta/\tau}+w\ge\theta$. Daí a janela de coincidência:
\begin{empheq}[box=\keybox]{equation}
\Delta \;\le\; \Delta_{\max} \;=\; \tau\,\ln\frac{w}{\theta-w} .
\label{eq:window}
\end{empheq}
Isto não é um E lógico, mas um \emph{detector de coincidência}: um E no tempo (fig.~\ref{fig:coincidence}). Com $\theta=1.5w$, a janela vale $\tau\ln2\approx0.7\tau$. Num neurônio do córtex com $\tau=10\ms$, são cerca de sete milissegundos. Nos neurônios do sistema auditivo, em que $\tau$ é menor que um milissegundo, a janela encolhe para frações de milissegundo.

\begin{figure}[t]
\centering
\begin{tikzpicture}[>=Stealth,x=0.9cm,y=1.6cm]
\foreach \dx/\lab/\c [count=\i] in {0.5/{coincidiram}/red!70!black, 2.6/{não coincidiram}/blue!60!black}{
 \begin{scope}[xshift={(\i-1)*7.2cm}]
  \draw[->,gray!70!black] (0,0) -- (6.3,0) node[below left,black] {\small $t$};
  \draw[->,gray!70!black] (0,0) -- (0,1.75) node[left,black] {\small $V$};
  \draw[densely dashed,gray!60!black] (0,1.5) -- (6.0,1.5) node[right,black] {\small $\theta$};
  \draw[very thick,\c] (0,0) -- (1,0) -- (1,1)
      plot[domain=1:1+\dx,samples=30] (\x,{exp(-(\x-1)/1.5)});
  \pgfmathsetmacro{\v}{exp(-\dx/1.5)+1}
  \draw[very thick,\c] (1+\dx,{exp(-\dx/1.5)}) -- (1+\dx,{min(\v,1.5)});
  \ifdim\v pt<1.5pt
    \draw[very thick,\c] plot[domain=1+\dx:6,samples=40] (\x,{\v*exp(-(\x-1-\dx)/1.5)});
  \else
    \draw[very thick,\c] (1+\dx,1.5) -- (1+\dx,0) -- (6,0);
    \draw[->,thick,\c] (1+\dx,1.55) -- (1+\dx,1.72) node[above,font=\scriptsize] {disparo};
  \fi
  \draw[->,thick] (1,-0.35) -- (1,-0.08); \draw[->,thick] (1+\dx,-0.35) -- (1+\dx,-0.08);
  \draw[decorate,decoration={brace,amplitude=3pt,mirror}] (1,-0.42) -- (1+\dx,-0.42) node[midway,below=3pt] {\scriptsize $\Delta$};
  \node[\c] at (3.5,-0.95) {\small \lab};
 \end{scope}}
\end{tikzpicture}
\caption{Detector de coincidência, limiar $\theta=1.5w$. À esquerda, o segundo disparo chegou após $\Delta<\Delta_{\max}$, a soma atingiu o limiar e o neurônio disparou. À direita, $\Delta>\Delta_{\max}$: do primeiro disparo quase nada restou, e o neurônio ficou calado.}
\label{fig:coincidence}
\end{figure}

Outra maneira de obter um E é pela duração, e não pelo instante exato. Enquanto a luz está acesa, o neurônio sensorial emite disparos frequentes e regulares; se uma entrada dessas não basta para o limiar, mas duas bastam, o neurônio trabalha enquanto ambas estão ativas. Assim a rede calcula por \emph{níveis}, e o instante exato de chegada dos disparos não importa. A maior parte do que vem a seguir funciona justamente assim.

\section{O que não se pode fazer}

Tentemos agora fazer um NÃO: um neurônio que trabalha quando a entrada está calada. Não dá: sem disparos de entrada, em~\eqref{eq:glutneuron} não há nenhum salto, e a tensão fica em zero, abaixo do limiar. E uma rede de muitos neurônios? Também não, e isso se prova de uma vez para todas as redes.

É mais cômodo falar de frequências. Seja $r_i$ a frequência de disparos do neurônio $i$, $I_i$ o aporte das entradas e $f_i$ sua característica de transferência: como a frequência depende do aporte total. No neurônio integrador, essa característica é não decrescente: mais aporte, disparos mais frequentes. As frequências numa rede que chegou ao equilíbrio satisfazem as equações
\begin{equation}
r_i \;=\; f_i\Bigl(\sum_j w_{ij}\,r_j + I_i\Bigr), \qquad w_{ij}\ge 0 .
\label{eq:ratenet}
\end{equation}

\begin{keyblock}
\begin{proposition}[Monotonicidade]
\label{prop:monotone}
Suponha que a rede~\eqref{eq:ratenet} seja formada só por neurônios \nt{GLUT} e parta do silêncio total. Se as entradas forem reforçadas, a frequência de regime de nenhum neurônio diminuirá. Tudo o que essa rede calcula são funções \emph{monótonas} das entradas.
\end{proposition}
\end{keyblock}

Demonstração. Partimos do silêncio, $r^{(0)}=0$, e substituímos as frequências no lado direito de~\eqref{eq:ratenet}: $r^{(k+1)}=F\bigl(r^{(k)}\bigr)$. O lado direito $F$ não decresce em nenhum argumento, porque os pesos são não negativos e as $f_i$ não decrescem. Como $r^{(1)}\ge0=r^{(0)}$, por indução $r^{(k+1)}\ge r^{(k)}$: as frequências só crescem e convergem para a menor solução de~\eqref{eq:ratenet}. Se as entradas $I$ forem trocadas por entradas maiores $I'$, então a cada passo $r^{(k)}(I)\le r^{(k)}(I')$, e no limite também. A rede real chega ao equilíbrio não em passos, mas continuamente, e para ela vale o mesmo: com conexões positivas, uma desigualdade entre duas execuções que vale no início se mantém o tempo todo.

Para disparos individuais, a proposição não vale ao pé da letra: um disparo de entrada a mais pode fazer o neurônio disparar um pouco antes e, por causa do milissegundo de refratariedade, o disparo seguinte se perde. Mas o efeito é fraco e pouco confiável, não se pode calcular com ele. Para frequências, a monotonicidade é estrita.

As consequências são as mesmas de qualquer circuito sem inversores. Não há NÃO. Não há OU exclusivo: acrescentar a segunda entrada não pode apagar a saída. E o mais desagradável: \emph{atividade não pode ser detida por atividade}, só se pode retirar o que a sustenta.

\section{Dois fios: ligar e desligar}

A saída do impasse é sugerida pelo próprio organismo. Em muitos órgãos dos sentidos, o estímulo é transmitido por dois conjuntos de neurônios: na visão, umas células respondem ao aumento da luz, outras à sua diminuição~\cite{Schiller1992}. Vamos chamá-las de \emph{liga} e \emph{desliga}. Em vez de um fio $x$, a rede recebe um par $(x,\bar x)$, e a ausência, que ela não sabe calcular, é informada pelo próprio sensor.

Com um par de fios, o NÃO sai de graça: basta trocar os fios. E e OU se fazem com dois neurônios cada:
\begin{empheq}[box=\keybox]{equation}
\begin{aligned}
x\wedge y \;&\to\; \bigl(\;x\wedge y,\;\; \bar x\vee\bar y\;\bigr),\\
x\vee y \;&\to\; \bigl(\;x\vee y,\;\; \bar x\wedge\bar y\;\bigr).
\end{aligned}
\label{eq:dualrail}
\end{empheq}
À direita todas as operações são monótonas, de modo que a proposição~\ref{prop:monotone} não é violada.

Para um circuito assíncrono, o código de dois fios tem uma segunda vantagem, ainda mais importante. Um par de fios tem três estados: ``sim'', ``não'' e, quando os dois estão calados, \emph{``ainda sem resposta''}. O destinatário sabe que o cálculo terminou não por um relógio, mas pelo próprio sinal: em cada par, exatamente um fio foi ativado. Entre estímulos, os sensores ficam calados, e a rede volta ao silêncio, pronta para a próxima vez. Na eletrônica assíncrona, esse truque é a base de circuitos que funcionam corretamente com quaisquer atrasos dos elementos~\cite{Sparso2001}.

\begin{keyblock}
\begin{proposition}[Cálculo assíncrono]
\label{prop:dualrail}
Se cada entrada chega por um par de fios ``liga--desliga'', uma rede de neurônios \nt{GLUT} pode calcular qualquer função lógica das entradas. A resposta também chega em pares de fios, e sua prontidão se vê pelo fato de que, em cada par, um fio foi ativado. Não é preciso relógio para isso. O preço é cerca do dobro de neurônios.
\end{proposition}
\end{keyblock}

Exemplo: o OU exclusivo, ``mudou exatamente um dos dois estímulos''. O fio direto da resposta é $(x\wedge\bar y)\vee(\bar x\wedge y)$; o inverso, $(x\wedge y)\vee(\bar x\wedge\bar y)$. São seis neurônios, e nenhum deles precisa de inibição (fig.~\ref{fig:dualxor}).

\begin{figure}[t]
\centering
\begin{tikzpicture}[>=Stealth,x=1cm,y=1cm,
  nrn/.style={circle,draw,very thick,fill=blue!6,minimum size=0.75cm,inner sep=0pt,font=\small},
  inp/.style={font=\small}]
\node[inp] (X)  at (0,3.3) {$x$};
\node[inp] (Xb) at (0,2.2) {$\bar x$};
\node[inp] (Y)  at (0,1.1) {$y$};
\node[inp] (Yb) at (0,0)   {$\bar y$};
\node[nrn] (A1) at (3.2,3.3) {$2$};
\node[nrn] (A2) at (3.2,2.2) {$2$};
\node[nrn] (A3) at (3.2,1.1) {$2$};
\node[nrn] (A4) at (3.2,0)   {$2$};
\node[nrn] (O)  at (7.2,2.75) {$1$};
\node[nrn] (Ob) at (7.2,0.55) {$1$};
\foreach \s/\t in {X/A1,Yb/A1,Xb/A2,Y/A2,X/A3,Y/A3,Xb/A4,Yb/A4}{\draw[->,thick,blue!60!black] (\s) -- (\t);}
\draw[->,thick,blue!60!black] (A1) -- node[pos=0.45,above,sloped,font=\scriptsize,black] {$x\wedge\bar y$} (O);
\draw[->,thick,blue!60!black] (A2) -- node[pos=0.45,below,sloped,font=\scriptsize,black] {$\bar x\wedge y$} (O);
\draw[->,thick,blue!60!black] (A3) -- node[pos=0.45,above,sloped,font=\scriptsize,black] {$x\wedge y$} (Ob);
\draw[->,thick,blue!60!black] (A4) -- node[pos=0.45,below,sloped,font=\scriptsize,black] {$\bar x\wedge\bar y$} (Ob);
\draw[->,very thick,red!70!black] (O) -- (8.6,2.75) node[right,black] {$x\oplus y$: ``sim''};
\draw[->,very thick,red!70!black] (Ob) -- (8.6,0.55) node[right,black] {$\overline{x\oplus y}$: ``não''};
\node[below,font=\small,gray!40!black] at (3.2,-0.55) {E: limiar $2$};
\node[below,font=\small,gray!40!black] at (7.2,-0.55) {OU: limiar $1$};
\node[below,font=\small,gray!40!black] at (0,-0.55) {pares de fios};
\end{tikzpicture}
\caption{OU exclusivo em código de dois fios, só com neurônios \nt{GLUT}. O número no círculo é o limiar em unidades do peso de uma entrada. Quatro neurônios E reconhecem as quatro combinações de entradas, e dois neurônios OU as reúnem no par de fios da resposta.}
\label{fig:dualxor}
\end{figure}

\section{Atrasos em vez de relógio}

Para cálculos com o tempo bastam atrasos nos fios e detectores de coincidência. O melhor exemplo é como o cérebro determina de onde veio um som.

O som de uma fonte situada de lado chega ao ouvido mais próximo antes que ao mais distante. A diferença depende da direção: de zero, quando a fonte está bem à frente, a $0.22\,\text{m}/343\,\text{m/s}\approx0.6\ms$ no ser humano, quando ela está exatamente de lado. O cérebro distingue diferenças de cerca de dez \emph{microssegundos}~\cite{KlumppEady1956,Grothe2010}, embora os próprios neurônios disparem com precisão não melhor que frações de milissegundo. Como?

Passemos um fio vindo do ouvido esquerdo ao longo de uma fileira de detectores de coincidência, da esquerda para a direita, e outro vindo do ouvido direito, da direita para a esquerda (fig.~\ref{fig:itd}). O sinal corre pelo fio com velocidade $v$. Até o detector situado à distância $x$ da borda esquerda de uma fileira de comprimento $L$, o sinal esquerdo leva o tempo $x/v$, o direito, $(L-x)/v$. Se o som chegou ao ouvido direito $\delta t$ mais tarde, os dois sinais se encontram ao mesmo tempo no ponto em que
\begin{empheq}[box=\keybox]{equation}
\frac{x}{v} \;=\; \frac{L-x}{v} + \delta t
\quad\Longrightarrow\quad
x \;=\; \frac{L}{2} + \frac{v\,\delta t}{2} .
\label{eq:itd}
\end{empheq}
Só dispara o detector ao qual os dois sinais chegaram dentro da janela~\eqref{eq:window}. O número do detector que disparou é a direção da fonte. A diferença de tempo virou \emph{lugar}.

\begin{figure}[t]
\centering
\begin{tikzpicture}[>=Stealth,
  nrn/.style={circle,draw,very thick,fill=blue!6,minimum size=0.7cm,inner sep=0pt}]
\foreach \k in {0,...,6}{\node[nrn] (D\k) at (1.4*\k,0) {\scriptsize $2$};}
\node[nrn,fill=red!15] at (D4) {\scriptsize $2$};
\draw[->,thick,blue!60!black] (-1.4,0.9) node[left,black,align=right] {\small ouvido\\\small esquerdo} -- (9.2,0.9);
\draw[->,thick,orange!85!black] (9.8,-0.9) node[right,black,align=left] {\small ouvido\\\small direito} -- (-0.8,-0.9);
\foreach \k in {0,...,6}{
  \draw[->,blue!60!black] (1.4*\k,0.9) -- (D\k.north);
  \draw[->,orange!85!black] (1.4*\k,-0.9) -- (D\k.south);}
\draw[->,thick,red!70!black] (D4.north east) ++(0.1,0.05) -- ++(0.5,0.45) node[above right,font=\small,black] {disparou};
\node[font=\small,gray!40!black] at (4.2,-1.5) {o atraso cresce $\longrightarrow$ \hspace{2.2cm} $\longleftarrow$ o atraso cresce};
\pic[scale=1.4] at (-2.3,1.75) {speaker};
\node[font=\scriptsize,align=center] at (-2.3,2.35) {fonte à esquerda};
\end{tikzpicture}
\caption{Determinação da direção de um som. Os sinais dos dois ouvidos correm um ao encontro do outro; cada círculo é um detector de coincidência com limiar $2$. Se o som chegou mais tarde ao ouvido direito, os sinais se encontram à direita do meio, pela fórmula~\eqref{eq:itd}. A saída da rede é o número do neurônio que disparou. Aqui a fonte está à esquerda: o som chegou antes ao ouvido esquerdo.}
\label{fig:itd}
\end{figure}

Aqui não há relógio, nem inibição, nem mesmo código de dois fios: a rede não responde ``sim'' ou ``não'', mas aponta para um neurônio entre muitos. Esse circuito de linhas de atraso e detectores de coincidência parece de fato funcionar no tronco auditivo das aves~\cite{Carr1990}. A precisão de microssegundos não vem da precisão de cada neurônio, mas do fato de haver muitos detectores, cada um olhando para seu próprio atraso. Nos mamíferos, essa fileira de atrasos não foi encontrada: lá a mesma tarefa parece ser resolvida por detectores de coincidência ajustados por uma inibição com tempo calculado com precisão, vinda de neurônios \nt{GLYC}~\cite{Brand2002,Grothe2010}. Como a inibição estreita a janela de coincidência veremos no capítulo~\ref{sec:gaba}, no exemplo do \nt{GABA}; a glicina inibe do mesmo modo, abrindo canais de cloreto no destinatário.

\section{Memória}

Um computador precisa de memória. Numa rede assim, a memória é uma atividade que se sustenta sozinha. Tomemos um grupo de neurônios \nt{GLUT} fortemente ligados entre si e vamos descrevê-lo por uma única frequência $r$ (fig.~\ref{fig:recurrentgroup}a). No equilíbrio, $r=f(wr+I)$, em que $w$ é a força da realimentação do grupo sobre si mesmo e $I$ é o aporte externo. Resolvemos essa equação graficamente, cruzando a curva $f(wr+I)$ com a reta $r$ (fig.~\ref{fig:bistable}).

\begin{figure}[t]
\centering
% (b): tau dr/dt = -r + f(w r + I), f as in fig:bistable, w=1, I=0.6; Euler dt=0.001 tau.
\begin{tikzpicture}[>=Stealth,
  nrn/.style={circle,draw,thick,fill=blue!6,minimum size=0.5cm,inner sep=0pt}]
\begin{scope}
\node[font=\small] at (-0.5,1.55) {(a)};
\foreach \k in {1,...,6}{\node[nrn] (n\k) at ({1.4+1.0*cos(60*\k)},{1.0*sin(60*\k)}) {};}
\foreach \a/\b in {1/2,2/3,3/4,4/5,5/6,6/1,1/4,2/5,3/6}{\draw[<->,blue!60!black] (n\a) -- (n\b);}
\foreach \k in {2,3,4}{\draw[->,thick,green!45!black] ($(n\k)+(-0.95,0)$) -- (n\k);}
\node[font=\small,green!45!black] at (-0.35,-0.45) {$I$};
\node[font=\Large] at (3.1,0) {$\approx$};
\node[nrn,minimum size=0.85cm,font=\small] (R) at (4.35,-0.1) {$r$};
\draw[->,thick,blue!60!black] (R) edge[loop above,looseness=6] node[above,font=\small] {$w$} (R);
\draw[->,thick,green!45!black] (3.55,-0.95) node[left,font=\small] {$I$} -- (R);
\end{scope}
\begin{scope}[xshift=6.3cm,y=2.6cm,x=1.05cm]
\node[font=\small] at (-0.6,1.25) {(b)};
\draw[->,gray!70!black] (0,0) -- (7.3,0) node[below left,font=\small,black] {$t/\tau$};
\draw[->,gray!70!black] (0,0) -- (0,1.12) node[left,font=\small,black] {$r$};
\foreach \t in {0,2,4,6}{\draw[gray!60!black] (\t,-0.02) -- (\t,0.02); \node[below,font=\scriptsize] at (\t,0) {\t};}
\draw[densely dotted,gray!60!black] (0,0.5) -- (7,0.5) node[right,font=\scriptsize,black,align=left] {limiar de\\gravação};
\fill[green!45!black,opacity=0.25] (1,-0.2) rectangle (2,-0.13);
\fill[gray!60!black,opacity=0.35] (1,-0.29) rectangle (1.5,-0.22);
\draw[very thick,blue!60!black] plot coordinates {(0.0,0.002) (0.1,0.002) (0.2,0.002) (0.3,0.002) (0.4,0.002) (0.5,0.002) (0.6,0.002) (0.7,0.002) (0.8,0.002) (0.9,0.002) (1.0,0.002) (1.1,0.083) (1.2,0.165) (1.3,0.242) (1.4,0.313) (1.5,0.378) (1.6,0.437) (1.7,0.491) (1.8,0.539) (1.9,0.583) (2.0,0.623) (2.1,0.643) (2.2,0.665) (2.3,0.688) (2.4,0.71) (2.5,0.732) (2.6,0.753) (2.7,0.773) (2.8,0.792) (2.9,0.81) (3.0,0.826) (3.2,0.855) (3.4,0.88) (3.6,0.9) (3.8,0.917) (4.0,0.931) (4.4,0.953) (4.8,0.967) (5.2,0.977) (5.6,0.984) (6.0,0.988) (6.5,0.992) (7.0,0.994)};
\draw[very thick,gray!60!black,densely dashed] plot coordinates {(0.0,0.002) (0.1,0.002) (0.2,0.002) (0.3,0.002) (0.4,0.002) (0.5,0.002) (0.6,0.002) (0.7,0.002) (0.8,0.002) (0.9,0.002) (1.0,0.002) (1.1,0.083) (1.2,0.165) (1.3,0.242) (1.4,0.313) (1.5,0.378) (1.6,0.358) (1.7,0.336) (1.8,0.313) (1.9,0.291) (2.0,0.269) (2.2,0.228) (2.4,0.191) (2.6,0.159) (2.8,0.133) (3.0,0.11) (3.2,0.091) (3.4,0.076) (3.6,0.063) (3.8,0.052) (4.0,0.043) (4.4,0.03) (4.8,0.021) (5.2,0.015) (5.6,0.011) (6.0,0.008) (6.5,0.006) (7.0,0.004)};
\node[font=\scriptsize,blue!60!black,anchor=west] at (3.3,0.8) {aporte $1\tau$: gravado};
\node[font=\scriptsize,gray!50!black,anchor=west] at (2.3,0.3) {aporte $0.5\tau$: apagou};
\end{scope}
\end{tikzpicture}
\caption{Um grupo fortemente ligado a si mesmo. (a)~Muitos neurônios \nt{GLUT} que se excitam mutuamente se comportam como um único neurônio de frequência $r$ que excita a si mesmo com força $w$; de fora chega o aporte $I$. (b)~Frequência do grupo no tempo, com $w=1$ e a mesma curva $f$ da fig.~\ref{fig:bistable}. O aporte $I=0.6$ é aplicado durante o tempo marcado pela faixa sob o eixo. Se ele dura $1\tau$, o grupo consegue passar o ponto instável $r=0.5$, acende e continua aceso quando o aporte termina. Se dura só $0.5\tau$, não chega a esse ponto e se apaga de volta: para gravar um bit, o aporte precisa durar mais de cerca de $0.7\tau$.}
\label{fig:recurrentgroup}
\end{figure}

\begin{figure}[t]
\centering
\begin{tikzpicture}[>=Stealth,x=5cm,y=3.2cm]
\draw[->,gray!70!black] (0,0) -- (1.1,0) node[below left,black] {\small $r$};
\draw[->,gray!70!black] (0,0) -- (0,1.1);
\draw[thick,gray!60!black] (0,0) -- (1.05,1.05) node[right,black] {\small $r$};
\draw[very thick,blue!60!black] plot[domain=0:1.05,samples=80] (\x,{1/(1+exp(-(\x-0.5)/0.08))});
\draw[very thick,red!70!black,densely dashed] plot[domain=0:1.05,samples=80] (\x,{1/(1+exp(-(\x+0.3-0.5)/0.08))});
\fill[blue!60!black] (0.002,0.002) circle (0.018);
\draw[blue!60!black,fill=white,thick] (0.5,0.5) circle (0.018);
\fill[blue!60!black] (0.998,0.998) circle (0.018);
\node[below right,font=\small] at (0.02,0.0) {deslig.};
\node[right,font=\small] at (0.52,0.46) {instável};
\node[below,font=\small] at (0.99,0.96) {lig.};
\node[anchor=east,font=\small,red!70!black] at (0.19,0.62) {$f(wr+I)$};
\node[anchor=west,font=\small,blue!60!black] at (0.45,0.1) {$f(wr)$};
\end{tikzpicture}
\caption{Memória feita de um grupo de neurônios \nt{GLUT} com forte realimentação. Curva contínua: sem aporte externo, duas interseções estáveis com a reta $r$, ``desligado'' e ``ligado'', e uma instável entre elas. Um aporte breve $I$ (curva tracejada) deixa só a de cima.}
\label{fig:bistable}
\end{figure}

Se a realimentação é forte, há três interseções: a de baixo é o silêncio, a de cima é o grupo aceso a todo vapor, a do meio é instável. Temos um elemento com dois estados estáveis, um flip-flop. Gravar nele um um é fácil: um aporte breve levanta a curva, a interseção de baixo some e o grupo acende, e continua aceso quando o aporte termina (fig.~\ref{fig:recurrentgroup}b). Um aporte curto demais não consegue levar o grupo até o ponto instável, e ele se apaga de volta. Essa atividade autossustentada, que dura segundos depois que o estímulo sumiu, é de fato observada no cérebro e é considerada a base da memória de curto prazo~\cite{Wang2001}. Mas essa não é a única maneira de lembrar por segundos. O registro também pode ser guardado numa variável lenta das próprias sinapses: depois de uma rajada de disparos, a facilitação, como na sinapse ``traço'' do capítulo~\ref{sec:code}, dura cerca de um segundo, e entre breves surtos a rede pode ficar quase calada: não um flip-flop, mas um capacitor carregado~\cite{Mongillo2008}. Esses intervalos ``silenciosos'' durante a retenção na memória são observados~\cite{Stokes2015}, e guardar sem disparos contínuos custa menos energia. Qual das duas maneiras o córtex usa em cada caso ainda está em discussão.

E como apagar? Pela proposição~\ref{prop:monotone}, por atividade não há como; resta \emph{retirar} algo. A primeira maneira é retirar o suporte: se o grupo só fica aceso junto com o aporte vindo de outro grupo, a memória se apaga com ele. A segunda é a fadiga. Nas sinapses reais, com trabalho prolongado, o estoque de neurotransmissor se esgota~\cite{Tsodyks1997}, e nos neurônios crescem correntes lentas que reduzem a excitabilidade. A realimentação enfraquece, a interseção de cima some, e o grupo se apaga sozinho em segundos. O resultado é uma memória que se liga por um sinal e só se desliga com o tempo.

\section{Sequências}

Em vez de um relógio, pode-se ter um \emph{temporizador disparado por evento}. Liguemos grupos de neurônios \nt{GLUT} em cadeia: o primeiro excita o segundo, o segundo o terceiro, e assim por diante. Um sinal externo acende o primeiro grupo, e uma onda de atividade percorre a cadeia, cada elo um ou dois atrasos depois do anterior, e uma única vez: logo após o disparo os neurônios estão refratários, e a onda já foi adiante.

Essa cadeia mede o tempo a partir de um evento: se disparou o elo número $k$, desde o início se passaram $k$ atrasos. Suas saídas podem ir para os músculos e produzir uma sequência de movimentos que se desenrola sozinha. Nas aves canoras, durante o canto, neurônios \nt{GLUT} individuais de uma das regiões do cérebro disparam estritamente em sequência, cada um no seu momento do canto, o que se parece muito com uma cadeia assim~\cite{Hahnloser2002}.

Ao contrário de um relógio, a cadeia fica calada até ser disparada, executa uma vez e mede o tempo só para quem está ligado a ela. A rede continua sem um ciclo de clock comum.

\section{O que falta}

Só com neurônios \nt{GLUT}, sem inibição, obtém-se muita coisa. Mas faltam três coisas, e as três por causa da proposição~\ref{prop:monotone}.

\emph{Escolha.} A rede precisa escolher uma direção, uma ação. Se dois estímulos são quase iguais, na rede da fig.~\ref{fig:itd} disparam as duas saídas, e não há com que suprimir a mais fraca em favor da mais forte.

\emph{Reset.} Não se pode apagar a memória com um sinal.

\emph{Estabilidade.} Numa rede em que as conexões surgem sem o plano de um engenheiro, laços de realimentação positiva são inevitáveis, e a atividade neles pode crescer em avalanche (capítulo~\ref{sec:neurontypes}). O único freio é a mesma fadiga, lenta e grosseira.

Os três males têm um só remédio: um neurônio que apaga um disparo com um disparo. É o \nt{GABA}, e ele não é necessário para calcular, mas para escolher, zerar e manter o cálculo sob controle.

\section{Exercícios}

\begin{exercise}[\lvlA]
\label{ex:02:1}
\emph{Fileira de detectores para o som.} Os sinais dos ouvidos correm pelos fios da fig.~\ref{fig:itd} a $5$~m/s; a maior diferença de chegada é $0.64\ms$. Os detectores, neurônios~\eqref{eq:glutneuron} com $\tau=0.5\ms$, $\theta=1.5w$, estão dispostos de modo que vizinhos distinguem $10$~\textmu s. Quantos detectores há na fileira e quantos disparam para um som?
\end{exercise}

\begin{exercise}[\lvlB]
\label{ex:02:2}
\emph{Entradas desiguais deslocam a janela.} Um detector de coincidência~\eqref{eq:glutneuron} com $\tau=0.5\ms$, $\theta=1.5$ recebe um disparo de cada uma de duas entradas com pesos $1$ e $0.8$. Encontre a janela de coincidência e seu centro. Em quantos microssegundos erra a fileira da fig.~\ref{fig:itd} se a entrada de um ouvido for $20\,\%$ mais fraca que a do outro?
\end{exercise}

\begin{exercise}[\lvlB]
\label{ex:02:3}
\emph{Que redes não oscilam.} A rede $\tau\dot r_i=-r_i+f_i\bigl(\textstyle\sum_jw_{ij}r_j+I_i\bigr)$ com $f_i$ não decrescentes e $w_{ij}\ge0$ para $i\ne j$ é monótona e, estavelmente, não oscila. Prove: pela troca $r_i\to\pm r_i$, a ela se reduz uma rede com número par de conexões inibitórias em cada ciclo. A inibição mútua de dois grupos pode oscilar? E o laço \nt{GLUT}--\nt{GABA}?
\end{exercise}

\begin{exercise}[\lvlB]
\label{ex:02:4}
\emph{Projeto: somador de dois fios.} Cada bit é transmitido por um par de fios $(x,\bar x)$. Construa um somador completo de três bits, que entregue os pares $(S,\bar S)$ da soma e $(C,\bar C)$ do vai-um, com neurônios \nt{GLUT}, indicando pesos e limiares. Use no máximo oito neurônios. Quantos exigiria um circuito de elementos E e OU, como na fig.~\ref{fig:dualxor}?
\end{exercise}

\begin{exercise}[\lvlB]
\label{ex:02:5}
\emph{Modelagem: quanto dura a gravação.} O grupo da fig.~\ref{fig:recurrentgroup}: $\tau\dot r=-r+f(r+I)$, $f(u)=1/\bigl(1+e^{-(u-0.5)/0.08}\bigr)$, antes da gravação no estado de baixo; o aporte $I$ é aplicado durante o tempo $T$. Encontre numericamente o menor $T$ que grava o bit com $I=0.6$ e o limiar $I_c$ abaixo do qual o bit não é gravado para nenhum $T$.
\end{exercise}

\begin{exercise}[\lvlA]
\label{ex:02:6}
\emph{A cadeia como temporizador.} Um elo da cadeia tem $100$ neurônios \nt{GLUT}, atraso de $2\ms$ com dispersão independente de $0.2\ms$. (a)~Quantos neurônios são necessários para medir $1$~s, e qual é o erro relativo? Como ele depende do intervalo? (b)~Num temporizador real, o erro é de $5\,\%$ para qualquer intervalo. Que fonte de dispersão produz isso?
\end{exercise}

\begin{exercise}[\lvlC]
\label{ex:02:7}
\emph{Pesquisa: flip-flop ou capacitor.} $100$ neurônios guardam um bit por $10$~s. Flip-flop: cada um dá $20$ disparos por segundo. Capacitor: a facilitação $u$ decai com $\tau_F=1$~s, o bit é lido se $u\ge u_{\max}/2$; renovação: rajada de três disparos por neurônio. Quantas vezes o capacitor é mais econômico, e o que acontece com o bit se o grupo for silenciado por $200\ms$?
\end{exercise}

\chapter{\nt{GLUT} e \nt{GABA} juntos}
\label{sec:gaba}

\section{As regras do jogo}

Uma rede só de neurônios \nt{GLUT} não sabe escolher, zerar a memória nem impedir que a atividade vire avalanche, e tudo por causa da monotonicidade (proposição~\ref{prop:monotone}). Vamos acrescentar a ela o segundo tipo de neurônio mais numeroso, o \nt{GABA}. As regras são as mesmas: só o que existe num organismo vivo, e nenhum sinal de clock.

O modelo de neurônio é o mesmo de~\eqref{eq:glutneuron}, mas o disparo de um neurônio \nt{GABA} não eleva a tensão do destinatário, e sim a abaixa. Agora os pesos têm dois sinais, e o sinal é definido pelo remetente, regra~\eqref{eq:dalesign}.

Alguns parâmetros do circuito. Os neurônios \nt{GABA} são de um quinto a um quarto dos neurônios do córtex~\cite{Hendry1987}. Suas saídas são curtas: eles inibem os vizinhos, não regiões distantes. A inibição chega em um milissegundo e dura alguns milissegundos. Sem entrada, o neurônio \nt{GABA} fica calado: para inibir alguém, ele mesmo precisa ser excitado, em geral por neurônios \nt{GLUT} da mesma rede.

Por isso, uma conexão negativa de um neurônio \nt{GLUT} $A$ para um neurônio $B$ tem de passar por um intermediário: $A$ excita um neurônio \nt{GABA}, e este inibe $B$. A troca de sinal custa um neurônio e um atraso, cerca de um milissegundo.

Na inibição importa não só \emph{de quem} vem a conexão, mas também \emph{onde} ela se instala: o local de contato determina em grande parte o que a conexão faz~\cite{Markram2004}. As saídas \nt{GLUT} se instalam principalmente nos ramos de entrada do destinatário, seus \emph{dendritos}, e levam dados: cada entrada dessas é uma parcela da soma. Uma saída no corpo do neurônio age sobre a soma inteira: assim muitos neurônios \nt{GABA} rápidos dividem a resposta do destinatário e definem quando ele vai disparar. Uma saída no local onde nasce o disparo do destinatário é a última palavra, um veto à saída inteira de uma vez. E uma saída no terminal do fio de outro neurônio muda a probabilidade de liberação $p$ de uma única linha de entrada, sem tocar nas demais~\cite{Rudomin1999}; é assim, em particular, que funciona a porta da dor no capítulo~\ref{sec:pain}. Fora do cérebro, as saídas se instalam em fibras musculares (capítulo~\ref{sec:muscles}) e em órgãos (capítulo~\ref{sec:chemoutput}), e algumas não se instalam em lugar nenhum e liberam o neurotransmissor no volume do tecido ou no sangue (capítulos~\ref{sec:serocomputer} e~\ref{sec:chemoutput}).

\section{NÃO e veto}

Dá para fazer um NÃO agora? Quase. Um neurônio que trabalha com a entrada calada precisa tirar excitação de algum lugar. No cérebro vivo ela existe: a cada neurônio chega o tempo todo uma atividade de fundo de milhares de outros. Suponha que o fundo mantenha o neurônio $Y$ ativo, e que a entrada $x$ o iniba por meio de um intermediário \nt{GABA}. Isso é um NÃO.

Mais útil, porém, é o \emph{veto}: ``$a$, mas não com $b$'':
\begin{empheq}[box=\keybox]{equation}
y \;=\; a \wedge \bar b .
\label{eq:veto}
\end{empheq}
O neurônio $Y$ recebe excitação de $a$ e inibição de um neurônio \nt{GABA} excitado por $b$. Aqui não é preciso fundo: a excitação vem do próprio $a$. Com vetos e OU monta-se tudo. Por exemplo, o OU exclusivo: $x\oplus y=(x\wedge\bar y)\vee(\bar x\wedge y)$, dois neurônios de veto, dois intermediários \nt{GABA} e um neurônio OU.

\begin{keyblock}
\begin{proposition}[Completude de \nt{GLUT}+\nt{GABA}]
\label{prop:gabacomplete}
Uma rede de neurônios \nt{GLUT} e \nt{GABA} pode calcular qualquer função lógica das entradas, e cada bit é transmitido por um único fio. Os sensores ``liga--desliga'' do capítulo~\ref{sec:glutcomputer} não são mais necessários para isso. Se a função deve dar um com todas as entradas caladas, é preciso uma fonte de atividade de fundo.
\end{proposition}
\end{keyblock}

Demonstração: o veto~\eqref{eq:veto} junto com o OU é funcionalmente completo se houver uma constante um, porque NÃO $x$ é o veto ``um, mas não com $x$''. E as funções que dão zero com todas as entradas caladas se montam com vetos e OU mesmo sem a constante um.

\section{Direção do movimento}

O veto no tempo, assim como o E no tempo no capítulo~\ref{sec:glutcomputer}, dá um detector de direção do movimento.

Sejam dois sensores vizinhos, $A$ e $B$, a uma distância $\Delta x$ um do outro. O neurônio de saída $Y$ recebe excitação de $B$ e inibição de $A$ por meio de um intermediário \nt{GABA}, com atraso $d$ (fig.~\ref{fig:direction}). Uma mancha de luz que corre com velocidade $v$ de $A$ para $B$ passa por $A$ no instante $0$, e a inibição chega a $Y$ no instante $d$; passa por $B$ no instante $\Delta x/v$, e nesse mesmo momento chega a excitação. Se esses instantes coincidem dentro da janela~\eqref{eq:window}, a inibição anula a excitação, e $Y$ fica calado. Isso acontece a uma velocidade de cerca de
\begin{empheq}[box=\keybox]{equation}
v^* \;=\; \frac{\Delta x}{d} .
\label{eq:vstar}
\end{empheq}
No movimento de $B$ para $A$, a excitação de $B$ chega no instante $0$, e a inibição de $A$ só no instante $\Delta x/v+d$, quando $Y$ já disparou.

\begin{figure}[t]
\centering
\begin{tikzpicture}[>=Stealth,
  nrn/.style={circle,draw,very thick,fill=blue!6,minimum size=0.9cm,inner sep=0pt},
  gab/.style={nrn,fill=red!10},
  inh/.style={-{Bar[width=7pt]},thick,red!70!black}]
\node[nrn] (A) at (0,2) {$A$};
\node[nrn] (B) at (3,2) {$B$};
\node[gab] (G) at (0.9,0.6) {\small \nt{GABA}};
\node[nrn] (Y) at (3,-0.6) {$Y$};
\draw[->,thick] (A) -- (G);
\draw[inh] (G) -- node[below left=-2pt] {\scriptsize atraso $d$} (Y);
\draw[->,thick] (B) -- (Y);
\draw[->,very thick,red!70!black] (Y) -- (4.6,-0.6) node[right,black] {\small saída};
\draw[->,thick,orange!85!black] (-0.6,2.9) -- (3.6,2.9) node[right,black] {\small calado};
\draw[<-,thick,orange!85!black] (-0.6,3.4) -- (3.6,3.4) node[right,black] {\small responde};
\foreach \yy/\dd in {2.9/-1,3.4/1}{
  \foreach \k in {1,2,3}{\draw[orange!70!yellow,line width=0.8pt] (1.5+\dd*0.12+\dd*0.13*\k,\yy+0.1-0.1*\k+0.1) -- ++(\dd*0.25,0);}
  \fill[yellow!70!orange,draw=orange!70!black] (1.5,\yy) circle (0.14);}
\node[font=\scriptsize,align=center] at (1.5,3.95) {mancha de luz};
\end{tikzpicture}
\caption{Detector de direção do movimento. $B$ excita $Y$, $A$ o inibe por meio de um intermediário \nt{GABA} com atraso $d$. No movimento de $A$ para $B$ com velocidade de cerca de $\Delta x/d$, a inibição anula a excitação; no movimento contrário, ela chega atrasada. A mancha amarela com tracinhos é a luz que corre sobre os sensores.}
\label{fig:direction}
\end{figure}

Na retina, a seletividade à direção do movimento parece ser criada exatamente assim: por uma inibição assimétrica vinda do lado de onde o movimento não deve provocar resposta~\cite{Barlow1965}.

\section{Subtração ou divisão}

Até agora, nossa inibição subtraía. Na verdade, ela é mais sutil.

Uma sinapse abre um canal, isto é, liga uma condutância. Numa sinapse excitatória, o potencial de equilíbrio fica muito acima do limiar, cerca de $70\mV$ acima do repouso: o canal aberto puxa a tensão para cima. Numa sinapse \nt{GABA} inibitória, o canal deixa passar cloreto, cujo potencial de equilíbrio fica perto do repouso, e o canal aberto puxa a tensão não para baixo, mas \emph{para o repouso}. Contando a tensão a partir do repouso, a tensão de regime de uma membrana com condutância de vazamento $g_L$, excitatória $g_E$ e inibitória $g_I$ é
\begin{empheq}[box=\keybox]{equation}
V_\infty \;=\; \frac{g_E\,E_E}{g_L+g_E+g_I} ,
\label{eq:shunt}
\end{empheq}
em que $E_E\approx70\mV$ é o potencial de equilíbrio da sinapse excitatória. É a lei de Ohm para um divisor de tensão: $g_E$ puxa para $E_E$, e $g_L$ e $g_I$, para zero.

$g_I$ não aparece no numerador com sinal de menos, e sim no denominador: a inibição não \emph{subtrai} da excitação, ela a \emph{divide} (fig.~\ref{fig:shunt}). Essa inibição se chama inibição por derivação: os canais de cloreto abertos curto-circuitam a membrana. Para a rede, é um controle de ganho. Se $g_I$ é proporcional à atividade total dos vizinhos, cada neurônio responde não à força absoluta de sua entrada, mas à sua fração na atividade total. Essa divisão pela atividade total aparece em muitas regiões do cérebro e é considerada uma de suas operações padrão~\cite{Carandini2012}: a mesma rede funciona tanto com entrada fraca quanto com entrada forte.

Uma ressalva: a fórmula~\eqref{eq:shunt} trata de um neurônio que não emite disparos. Num neurônio que dispara, a tensão fica o tempo todo perto do limiar, a corrente pela condutância inibitória é quase constante, e sobre a \emph{frequência} de disparos a derivação age principalmente por subtração~\cite{HoltKoch1997}. A frequência é dividida se ao neurônio forem acrescentadas ao mesmo tempo condutâncias de fundo excitatória e inibitória equilibradas, como no regime trêmulo e equilibrado do córtex vivo~\cite{Chance2002}. Assim, o cérebro obtém a divisão não de uma derivação isolada, mas da inibição junto com o ruído de fundo~\cite{Carandini2012}.

\begin{figure}[t]
\centering
\begin{tikzpicture}[>=Stealth,x=1.25cm,y=0.05cm]
\draw[->,gray!70!black] (0,0) -- (5.4,0) node[right,black] {\small $g_E/g_L$};
\draw[->,gray!70!black] (0,0) -- (0,68) node[above,black] {\small $V_\infty$, mV};
\foreach \v in {20,40,60}{\draw[gray!60!black] (-0.05,\v) -- (0.05,\v) node[left=3pt,font=\scriptsize] {\v};}
\foreach \x in {1,2,3,4,5}{\draw[gray!60!black] (\x,-1.5) -- (\x,1.5) node[below=5pt,font=\scriptsize] {\x};}
\draw[very thick,blue!60!black] plot[domain=0:5,samples=60] (\x,{70*\x/(1+\x)});
\draw[very thick,red!70!black] plot[domain=0:5,samples=60] (\x,{70*\x/(3+\x)});
\draw[thick,densely dashed,gray!50!black] plot[domain=0.56:5,samples=60] (\x,{70*\x/(1+\x)-25});
\node[right,font=\small,blue!60!black] at (5.02,58.3) {sem inibição};
\node[right,font=\small,red!70!black] at (5.02,45.5) {divide: $g_I=2g_L$};
\node[right,font=\small,gray!40!black] at (5.02,32.5) {subtrai $25\mV$};
\end{tikzpicture}
\caption{A inibição por derivação divide a tensão em vez de subtrair dela. A curva azul é a tensão de regime~\eqref{eq:shunt} sem inibição; a vermelha, com condutância inibitória $g_I=2g_L$. A vermelha é a mesma azul esticada três vezes na horizontal: como se a entrada fosse dividida por $1+g_I/g_L=3$. A tracejada é uma inibição que subtrai $25\mV$ constantes: a curva desce inteira, e a inclinação não muda.}
\label{fig:shunt}
\end{figure}

Há também uma segunda consequência. A constante de tempo da membrana é $\tau_{\text{ef}}=C/(g_L+g_E+g_I)$, e a inibição a encurta. Como a janela de coincidência~\eqref{eq:window} é proporcional a $\tau$, uma inibição que chega junto com a excitação \emph{estreita a janela}. O circuito em que a entrada excita o neurônio diretamente e, ao mesmo tempo, com atraso de um milissegundo, o inibe por meio de um intermediário \nt{GABA} é um dos mais comuns no cérebro: o neurônio só tem tempo de responder ao que chegou no primeiro ou segundo milissegundo~\cite{Pouille2001}.

\section{Escolha}

Agora o principal que faltava à rede \nt{GLUT}: escolher uma entre várias opções. Tomemos dois grupos de neurônios \nt{GLUT} com entradas $I_1$ e $I_2$. Cada um, por meio de seus intermediários \nt{GABA}, inibe o outro com força $\beta$ (fig.~\ref{fig:wta}). Para as frequências $r_1,r_2$ obtemos
\begin{equation}
\tau\,\frac{\dd r_1}{\dd t} = -r_1 + \bigl[I_1-\beta r_2\bigr]_+ ,\qquad
\tau\,\frac{\dd r_2}{\dd t} = -r_2 + \bigl[I_2-\beta r_1\bigr]_+ ,
\label{eq:wta}
\end{equation}
em que $[u]_+=\max(u,0)$: a frequência não pode ser negativa.

\begin{figure}[t]
\centering
\begin{tikzpicture}[>=Stealth,
  nrn/.style={circle,draw,very thick,fill=blue!6,minimum size=0.9cm,inner sep=0pt},
  gab/.style={nrn,fill=red!10,minimum size=0.75cm},
  inh/.style={-{Bar[width=7pt]},thick,red!70!black}]
\node[nrn] (E1) at (0,0) {$r_1$};
\node[nrn] (E2) at (4,0) {$r_2$};
\node[gab] (G1) at (1.4,1.1) {};
\node[gab] (G2) at (2.6,-1.1) {};
\draw[->,thick] (E1) -- (G1); \draw[inh] (G1) -- (E2);
\draw[->,thick] (E2) -- (G2); \draw[inh] (G2) -- (E1);
\draw[->,thick,blue!60!black] (-1.6,0) node[left,black] {$I_1$} -- (E1);
\draw[->,thick,blue!60!black] (5.6,0) node[right,black] {$I_2$} -- (E2);
\end{tikzpicture}
\caption{Escolha entre dois. Dois grupos de neurônios \nt{GLUT} (azuis) inibem um ao outro por meio de intermediários \nt{GABA} (vermelhos).}
\label{fig:wta}
\end{figure}

Enquanto os dois neurônios estão ativos, o sistema~\eqref{eq:wta} é linear, e seu comportamento é dado por dois autovalores: $-(1+\beta)/\tau$ para desvios iguais das duas frequências e $-(1-\beta)/\tau$ para desvios opostos. Com inibição fraca, $\beta<1$, os dois autovalores são negativos: os dois neurônios trabalham, e a inibição apenas abafa o mais fraco. Com inibição forte, $\beta>1$, o segundo autovalor fica positivo: qualquer diferença entre $r_1$ e $r_2$ cresce até que um neurônio se cale por completo.

\begin{keyblock}
\begin{proposition}[O vencedor leva tudo]
\label{prop:wta}
Se a inibição mútua é maior que um, $\beta>1$, o estado em que os dois grupos trabalham é instável. A rede chega a um estado em que um grupo trabalha e o outro fica calado. O vencedor, com frequência $r_1=I_1$, mantém a vitória enquanto $I_2<\beta I_1$.
\end{proposition}
\end{keyblock}

Verificamos isso num modelo. Com $\beta=0.5$ e entradas $1.0$ e $0.9$, os dois grupos trabalham, com frequências $0.73$ e $0.53$. Com $\beta=2$ e as mesmas entradas, o segundo grupo fica totalmente calado. Essa escolha tem memória: se depois as entradas forem trocadas, o vencedor anterior continua vencedor, porque $1.0<2\cdot0.9$.

Com cem grupos é igual: cada um inibe os demais, e sobrevive um.

\section{Reset da memória}

A memória do capítulo~\ref{sec:glutcomputer} só se desligava pela fadiga. Vamos acrescentar a ela um neurônio \nt{GABA} que é excitado por um sinal de ``reset'' e inibe o grupo. A inibição abaixa a curva $f(wr+I)$ da fig.~\ref{fig:bistable}, a interseção de cima some e o grupo se apaga, e fica no estado de baixo quando o reset termina. Agora a memória é gravada por um sinal e apagada por outro, como num flip-flop com entradas de set e reset.

Mais elegante ainda é ligar dois desses grupos por inibição mútua, como na fig.~\ref{fig:wta}, e dar a cada um realimentação sobre si mesmo. Um sinal na entrada de um dos grupos leva a célula ao estado dele, e a célula lembra a última decisão. É o análogo direto de um latch de dois elementos NOU acoplados em cruz.

\section{Estabilidade e ritmo}

Resta o terceiro mal da rede \nt{GLUT}: a avalanche. Tomemos um grupo de neurônios \nt{GLUT} com realimentação sobre si mesmo $w_{EE}$ e um grupo de neurônios \nt{GABA} que o primeiro excita com força $w_{IE}$ e que inibem o primeiro com força $w_{EI}$. Para as frequências $E$ e $I$ obtemos
\begin{equation}
\tau_E\frac{\dd E}{\dd t} = -E + w_{EE}E - w_{EI}I + h, \qquad
\tau_I\frac{\dd I}{\dd t} = -I + w_{IE}E ,
\label{eq:ei}
\end{equation}
em que $h$ é a entrada externa~\cite{WilsonCowan1972}. Sem inibição, com $w_{EE}>1$, a atividade cresce sem limite: cada disparo gera mais de um disparo seguinte. Com inibição, a frequência de regime
\begin{equation}
E^* \;=\; \frac{h}{1-w_{EE}+w_{EI}w_{IE}}
\label{eq:eistar}
\end{equation}
é finita se o denominador for positivo. Mas isso não basta: o equilíbrio ainda precisa atrair. Da análise linear de~\eqref{eq:ei} saem duas condições.

\begin{keyblock}
\begin{proposition}[Condições de estabilidade]
\label{prop:ei}
O equilíbrio~\eqref{eq:eistar} é estável se a inibição for
\begin{enumerate}
\item[(i)] \emph{forte o bastante}: $w_{EI}\,w_{IE} > w_{EE}-1$;
\item[(ii)] \emph{rápida o bastante}: $\tau_I < \tau_E/(w_{EE}-1)$.
\end{enumerate}
Se (i) é violada, a atividade cresce em avalanche. Se (i) vale mas (ii) é violada, a inibição chega atrasada, e a atividade começa a oscilar.
\end{proposition}
\end{keyblock}

A condição~(i) é o determinante da matriz do sistema~\eqref{eq:ei}; a condição~(ii), o seu traço. O sentido da segunda é claro para quem já ajustou um controlador: uma realimentação atrasada não amortece o desvio, mas o amplifica.

\begin{figure}[t]
\centering
\begin{tikzpicture}[>=Stealth]
\draw[->,gray!70!black] (0,0) -- (10.9,0) node[below left,black] {\small $t$, ms};
\draw[->,gray!70!black] (0,0) -- (0,3.3) node[left,black] {\small $E$};
\foreach \t in {0,50,100,150}{\draw[gray!70!black] (\t*0.07,0.06) -- (\t*0.07,-0.06) node[below,font=\scriptsize] {\t};}
\input{figures/ei_traces}
\node[font=\small,gray!40!black,anchor=west] at (0.9,3.15) {sem inibição: avalanche};
\node[font=\small,blue!60!black,anchor=west] at (7.4,1.25) {inibição rápida};
\node[font=\small,red!70!black,anchor=west] at (7.4,2.55) {inibição lenta};
\end{tikzpicture}
\caption{Frequência de um grupo \nt{GLUT} com realimentação $w_{EE}=1.5$, $\tau_E=10\ms$, após a entrada ser ligada. Sem inibição (tracejado), a atividade cresce sem limite. Com inibição de força $w_{EI}w_{IE}=2$ e $\tau_I=2\ms$ (curva azul), ela se estabiliza no nível $E^*=h/(1-1.5+2)=0.67h$. Com a mesma inibição, mas lenta, $\tau_I=30\ms$ (curva vermelha), a atividade oscila.}
\label{fig:ei}
\end{figure}

Considera-se que o córtex trabalha exatamente nesse regime: sua própria realimentação é forte o bastante para cair em avalanche sem inibição, e só a inibição rápida o mantém no ponto de operação~\cite{Tsodyks1997isn}.

Esse regime tem uma assinatura paradoxal pela qual pode ser reconhecido de fora~\cite{Tsodyks1997isn}. Acrescentemos em~\eqref{eq:ei} uma entrada externa $h_I$ diretamente ao grupo \nt{GABA}. As frequências de regime passam a ser
\begin{equation}
E^* \;=\; \frac{h-w_{EI}\,h_I}{D},\qquad
I^* \;=\; \frac{w_{IE}\,h+(1-w_{EE})\,h_I}{D},\qquad
D \;=\; 1-w_{EE}+w_{EI}w_{IE} .
\label{eq:isn}
\end{equation}
Se $w_{EE}>1$, o coeficiente de $h_I$ na segunda fórmula é negativo: empurre os neurônios \nt{GABA}, e sua frequência \emph{cai}. A causa está no laço. A inibição extra abafa o grupo \nt{GLUT}, este excita menos os neurônios \nt{GABA}, e eles perdem mais do que ganharam. Com os números da fig.~\ref{fig:ei} ($w_{EE}=1.5$, $w_{EI}w_{IE}=2$, $D=1.5$), cada unidade de entrada acrescentada reduz $I^*$ em um terço de unidade. Essa assinatura foi encontrada no córtex: quando a resposta de neurônios do córtex visual é suprimida por um estímulo ao redor, a corrente inibitória que eles recebem não cresce, mas cai junto com a excitatória~\cite{Ozeki2009}. Mais tarde a mesma assinatura foi encontrada em diferentes áreas e camadas do córtex~\cite{Sanzeni2020}.

Oscilações quando a condição~(ii) é violada também ocorrem no cérebro. O laço ``\nt{GLUT} excita \nt{GABA}, \nt{GABA} inibe \nt{GLUT}'' com atraso de alguns milissegundos gera um ritmo com período da ordem de $12$--$50\ms$ (frequência de $20$--$80$\,Hz), e esses ritmos rápidos no córtex são bem conhecidos~\cite{Cardin2009,Buzsaki2012}. Não é o clock de um computador: o ritmo é local e só surge quando a rede está ativa. Mas, ao longo de alguns ciclos, ele divide o tempo em janelas nas quais os neurônios podem disparar.

\section{Resumo}

\begin{table}[t]
\centering
\small
\begin{tabular}{>{\raggedright}p{3.3cm}>{\raggedright}p{4.4cm}>{\raggedright\arraybackslash}p{4.4cm}}
\toprule
& só \nt{GLUT} & \nt{GLUT} + \nt{GABA} \\
\midrule
NÃO & só por sensores ``liga--desliga'' & veto $a\wedge\bar b$; NÃO com atividade de fundo \\
fios por bit & dois & um \\
direção do movimento & não & veto com atraso \\
controle de ganho & não & divisão: derivação junto com ruído de fundo \\
janela de coincidência & $\sim\tau$ & estreitada pela inibição \\
escolha de um entre muitos & não & inibição mútua, $\beta>1$ \\
reset da memória & só por fadiga & por sinal via \nt{GABA} \\
estabilidade & só por fadiga & realimentação negativa rápida \\
ritmo & desnecessário & surge com inibição lenta \\
\bottomrule
\end{tabular}
\caption{O que o \nt{GABA} acrescenta a uma rede de neurônios \nt{GLUT}.}
\label{tab:gaba}
\end{table}

O \nt{GABA} quase não acrescenta à rede novos \emph{cálculos} (tudo isso também podia ser calculado com sensores de dois fios), mas acrescenta \emph{controle} (tabela~\ref{tab:gaba}): escolha, reset, controle de ganho e estabilidade. No cérebro, a excitação leva os dados, e a inibição decide quais dados passam, quando e com que volume. Para um em cada quatro ou cinco neurônios do córtex, é uma divisão de trabalho muito sensata.

\section{Exercícios}

\begin{exercise}[\lvlA]
\label{ex:03:1}
\emph{A derivação em números.} $E_E=70\mV$. Por~\eqref{eq:shunt}, encontre $V_\infty$ para $g_E=g_L$ e $4g_L$, sem inibição e com derivação $g_I=2g_L$. Que $V_\infty$ em $g_E=4g_L$ daria uma subtração que, em $g_E=g_L$, retira tanto quanto a derivação? Quantas vezes a derivação estreita a janela de coincidência?
\end{exercise}

\begin{exercise}[\lvlB]
\label{ex:03:2}
\emph{Dedução das condições de estabilidade.} Encontre o traço e o determinante da matriz do sistema linear~\eqref{eq:ei} e deduza deles as duas condições da proposição~\ref{prop:ei}. Na fronteira da condição~(ii), o equilíbrio perde estabilidade de forma oscilatória. Encontre o período dessas oscilações com os números da fig.~\ref{fig:ei}.
\end{exercise}

\begin{exercise}[\lvlB]
\label{ex:03:3}
\emph{Ritmo a partir do atraso.} Troquemos a inibição lenta do modelo~\eqref{eq:ei} por uma rápida com atraso $d$: $\tau_E\dot E=-E(t)-GE(t-d)+h$. Com $\tau_E=10\ms$, encontre o menor atraso $d_c$ que produz oscilações e seu período para $G=2$ e $G=5$. Eles caem nos ritmos rápidos do córtex, $12$--$50\ms$, ao contrário do exercício~\ref{ex:03:2}?
\end{exercise}

\begin{exercise}[\lvlB]
\label{ex:03:4}
\emph{Modelagem: escolha com memória.} Integre~\eqref{eq:wta} com $\beta=2$, $\tau=1$. (a)~Com $I_1=1$, suba $I_2$ lentamente de $0$ a $3$ e desça de volta: em que valores de $I_2$ a rede comuta? (b)~Quanto cresce o tempo de decisão a partir do silêncio quando a diferença das entradas $\varepsilon$ diminui dez vezes?
\end{exercise}

\begin{exercise}[\lvlB]
\label{ex:03:5}
\emph{Projeto: detector para uma faixa de velocidades.} O detector da fig.~\ref{fig:direction} deve ficar calado no movimento de $A$ para $B$ com tempo de percurso de $5$--$20\ms$. Um intermediário \nt{GABA} com atraso $d_k$ veta $Y$ no intervalo $[d_k,\,d_k+5\ms]$ após o disparo de $A$; a excitação de $B$ chega de imediato. Quantos intermediários são necessários, e com que atrasos?
\end{exercise}

\begin{exercise}[\lvlB]
\label{ex:03:6}
\emph{Vetos e fundo.} Quantos neurônios exige a paridade $x\oplus y\oplus z$ no esquema geral ``intermediário \nt{GABA} por entrada, veto para cada combinação com $f=1$, OU comum''? Construa-a com dois neurônios, \nt{GABA} e \nt{GLUT}, que recebem as entradas diretamente, indicando pesos e limiares.
\end{exercise}

\begin{exercise}[\lvlC]
\label{ex:03:7}
\emph{Projeto: escolha entre cem.} Cada um de $100$ grupos \nt{GLUT} deve inibir igualmente todos os outros, mas não a si mesmo. Intermediários \nt{GABA} lineares são excitados pelos grupos e os inibem; os grupos podem excitar uns aos outros, mas não a si mesmos. Qual é o menor número de intermediários? E se não houver nenhuma conexão excitatória direta?
\end{exercise}

\begin{exercise}[\lvlC]
\label{ex:03:8}
\emph{Pesquisa: quando a inibição é paradoxal.} No modelo da fig.~\ref{fig:ei} ($w_{EI}=2$, $w_{IE}=1$), o grupo \nt{GLUT} recebeu a função de transferência $f(u)=[u]_+^2$, e o grupo \nt{GABA} continuou linear. Acima de que entrada $h_c$ um acréscimo ao grupo \nt{GABA} reduz a própria frequência dele? Verifique por simulação.
\end{exercise}

\chapter{Código Morse: o que sabemos sobre a linguagem dos neurônios \nt{GLUT} e \nt{GABA}}
\chaptermark{Código Morse}
\label{sec:code}

\section{Um alfabeto de um só sinal}

O código Morse tem dois sinais, ponto e traço, e pausas de três durações entre eles. O neurônio, como vimos no capítulo~\ref{sec:neurontypes}, tem um alfabeto ainda mais pobre: um único sinal. O disparo dura cerca de um milissegundo, e todos os disparos de um neurônio têm a mesma altura e a mesma forma. Por isso a mensagem inteira está nas pausas, isto é, na sequência de instantes $t_1,t_2,t_3,\dots$ É um código sem traços, em que só o ritmo dos pontos carrega sentido (fig.~\ref{fig:morse}).

\begin{figure}[t]
\centering
\begin{tikzpicture}[>=Stealth,x=0.08cm,y=1cm]
% Morse letter: dot, dash, dot
\node[left,font=\small] at (-6,2.55) {Morse:};
\fill[black] (0,2.45) rectangle (6,2.65);
\fill[black] (14,2.45) rectangle (32,2.65);
\fill[black] (40,2.45) rectangle (46,2.65);
\node[right,font=\small,gray!40!black] at (52,2.55) {ponto, traço, ponto: o sentido está nas durações e pausas};
% stimulus onset
\node[left,font=\small] at (-6,0.4) {neurônio:};
\draw[->,thick,green!45!black] (0,-0.55) -- (0,-0.05);
\node[below,font=\scriptsize,green!45!black] at (0,-0.55) {estímulo};
\draw[gray!70!black] (0,0) -- (152,0) node[right,black] {\small $t$};
% spikes: single ones blue, the burst red
\foreach \s in {14,26,41,88,112,131}{\draw[very thick,blue!60!black] (\s,0) -- (\s,0.8);}
\foreach \s in {60,63,66,69}{\draw[very thick,red!70!black] (\s,0) -- (\s,0.8);}
% latency
\draw[<->,thick] (0,1.1) -- node[above,font=\scriptsize] {$t_1$: tempo do primeiro disparo} (14,1.1);
% burst
\draw[decorate,decoration={brace,amplitude=4pt},red!70!black] (58,0.95) -- (71,0.95) node[midway,above=4pt,font=\scriptsize] {rajada: ``traço''};
% counting windows
\foreach \a/\b/\n in {0/50/3,50/100/5,100/150/2}{
  \draw[decorate,decoration={brace,amplitude=4pt,mirror},gray!50!black] (\a+1,-0.2) -- (\b-1,-0.2) node[midway,below=4pt,font=\small] {$n=\n$};}
\node[font=\scriptsize,gray!40!black] at (75,-1.05) {código de taxa: número de disparos $n$ em cada janela};
\foreach \x in {0,50,100,150}{\draw[gray!60!black] (\x,-0.06) -- (\x,0.06);}
\end{tikzpicture}
\caption{O código Morse e o código do neurônio. A mesma sequência de disparos pode ser lida de maneiras diferentes: contando os disparos em janelas de $50\ms$ (seção~\ref{sec:rate}), medindo a latência $t_1$~\eqref{eq:latency} ou notando a rajada, o ``traço''~\eqref{eq:burst}.}
\label{fig:morse}
\end{figure}

Que pausas são possíveis? Por baixo, elas são limitadas pela refratariedade: depois de um disparo, o neurônio não pode disparar de novo por cerca de um milissegundo, então não passa de algumas centenas de disparos por segundo. Por cima, nada as limita: um neurônio pode ficar calado por minutos. Nos neurônios \nt{GLUT} do córtex, as frequências vão de frações de disparo a algumas dezenas de disparos por segundo, e a maioria dos neurônios passa a maior parte do tempo perto do limite inferior.

Nos capítulos~\ref{sec:glutcomputer} e~\ref{sec:gaba}, de passagem, adotamos ora uma, ora outra maneira de escrever números com disparos. Agora perguntemos diretamente: em que linguagem os neurônios \nt{GLUT} e \nt{GABA} falam entre si? Não há resposta completa; essa linguagem não foi decifrada. Mas algumas coisas sobre ela são sabidas com segurança, e quase tudo o que se sabe pode ser obtido raciocinando como um engenheiro de comunicações: capacidade do canal, ruído, receptor, custo da transmissão.

\section{Quanto se pode transmitir}

Comecemos pelo limite superior. Suponha que o destinatário distinga os instantes dos disparos com precisão $\Delta t$: deslocamentos menores que $\Delta t$ são indistinguíveis para ele. Cortemos o tempo em células de comprimento $\Delta t$, menores que o tempo de refratariedade, de modo que em cada célula haja um disparo ou nenhum (fig.~\ref{fig:cells}). Com frequência média $r$, um disparo cai numa célula com probabilidade $p=r\Delta t$. O fluxo carrega o máximo de informação quando as células são independentes, e então cada célula dá a entropia $-p\log_2p-(1-p)\log_2(1-p)\approx p\log_2(e/p)$ bits para $p\ll1$. Dividindo por $\Delta t$, obtemos a taxa de transmissão máxima e sua parcela por disparo~\cite{MacKay1952}:
\begin{empheq}[box=\keybox]{equation}
R_{\max} \;\approx\; r\,\log_2\frac{e}{r\,\Delta t}\ \ \text{bits/s},
\qquad
\frac{R_{\max}}{r} \;\approx\; \log_2\frac{e}{r\,\Delta t}\ \ \text{bits por disparo}.
\label{eq:spikeentropy}
\end{empheq}
Com $r=10$ disparos por segundo e $\Delta t=1\ms$, são cerca de oito bits por disparo e oitenta bits por segundo.

\begin{figure}[t]
\centering
\begin{tikzpicture}[>=Stealth,x=0.5cm,y=1cm]
% time cut into cells of width Delta t; a cell holds one impulse or none
\foreach \k in {0,...,23}{\draw[gray!55] (\k,0) rectangle (\k+1,0.7);}
\foreach \k in {2,7,8,15,21}{
  \fill[red!10] (\k,0) rectangle (\k+1,0.7);
  \draw[very thick,red!70!black] (\k+0.5,0.05) -- (\k+0.5,0.65);}
\foreach \k in {0,...,23}{
  \pgfmathtruncatemacro{\bit}{(\k==2)||(\k==7)||(\k==8)||(\k==15)||(\k==21)}
  \ifnum\bit=1 \def\bitcol{red!70!black}\else \def\bitcol{gray!60!black}\fi
  \node[font=\scriptsize,text=\bitcol] at (\k+0.5,-0.25) {\bit};}
\draw[<->,gray!60!black] (0,0.95) -- (1,0.95) node[midway,above,font=\scriptsize,black] {$\Delta t$};
\node[font=\small,anchor=east] at (-0.3,0.35) {disparos};
\node[font=\small,anchor=east] at (-0.3,-0.25) {bits};
\draw[decorate,decoration={brace,amplitude=5pt,mirror},gray!60!black] (0,-0.55) -- (24,-0.55)
  node[midway,below=6pt,font=\scriptsize,black,align=center] {destinatário que só conta: ``$n=5$'';\\ destinatário que vê as células: quais $5$ de $24$, isto é, $\log_2\binom{24}{5}\approx15$ bits};
\end{tikzpicture}
\caption{De onde vem a capacidade~\eqref{eq:spikeentropy}. O tempo é cortado em células de comprimento $\Delta t$; em cada célula há um disparo ($1$) ou não ($0$): o fluxo de disparos é uma cadeia binária. Um disparo cai numa célula com probabilidade $p=r\Delta t$. Quanto menores as células, mais longa a cadeia e mais bits; um contador de disparos perde quase tudo o que está escrito em \emph{onde} ficam os uns.}
\label{fig:cells}
\end{figure}

Os bits por disparo dependem da precisão temporal só logaritmicamente: com precisão de $10\ms$ restam cerca de cinco bits por disparo. Mas se o destinatário só conta os disparos num segundo, com $r=10$ ele recebe menos de dois bits no segundo inteiro (seção~\ref{sec:rate}).

\begin{keyblock}
\begin{proposition}[Capacidade do canal de disparos]
\label{prop:capacity}
Um fluxo de disparos com frequência média $r$, lido com precisão temporal $\Delta t$, não carrega mais que~\eqref{eq:spikeentropy}. Quase toda essa capacidade está no \emph{tempo} dos disparos: um destinatário que só conta disparos extrai do mesmo fluxo dezenas de vezes menos do que um que distingue seus instantes com precisão de milissegundo.
\end{proposition}
\end{keyblock}

Isto é um limite superior, não uma medição. As medições em neurônios sensoriais, para os quais o estímulo é conhecido, dão de um a alguns bits por disparo; nos melhores casos, até metade do limite calculado para sua precisão~\cite{Rieke1997}. Logo, pelo menos na entrada do sistema nervoso, o tempo dos disparos é usado, e não se perde.

\section{Um canal ruidoso}

A capacidade~\eqref{eq:spikeentropy} foi calculada sem ruído. Sobre onde ele surge, três fatos são sabidos com segurança.

\emph{O gerador de disparos em si é preciso.} Se num neurônio do córtex, retirado do cérebro com um pedaço de tecido, se injeta muitas vezes a mesma corrente, que varia rápido no tempo, os disparos se repetem com precisão de cerca de um milissegundo~\cite{Mainen1995}. Se a corrente é constante, os instantes dos disparos se espalham de uma vez para outra: o tempo preciso é gerado pelas variações rápidas da entrada, não pelo próprio neurônio.

\emph{A sinapse não é confiável.} Um disparo que chega a uma sinapse \nt{GLUT} libera uma porção de neurotransmissor apenas com certa probabilidade $p$. Nas sinapses do hipocampo, mais da metade dos disparos simplesmente não chega ao destinatário, e a causa é justamente o acaso da liberação, não falhas de condução~\cite{AllenStevens1994}. Para um engenheiro de comunicações, é um canal com apagamentos; várias sinapses entre dois neurônios salvam em parte a situação.

\emph{No córtex vivo, o fluxo de disparos parece aleatório.} Os intervalos entre disparos se espalham mais ou menos como num fluxo aleatório de Poisson, e, quando o mesmo estímulo é repetido, o número de disparos numa janela varia de modo que sua variância fica próxima da média~\cite{Softky1993}.

Como um elemento preciso gera um fluxo aleatório? O neurônio recebe milhares de entradas, cada uma pouco confiável, e a excitação e a inibição estão equilibradas, como no capítulo~\ref{sec:gaba}. A tensão da membrana treme perto do limiar, e os instantes dos disparos são definidos por esses tremores. Se isso é ruído ou uma mensagem que não sabemos ler, é em grande parte uma questão em aberto. Parte do ``ruído'' já se revelou mensagem: no córtex visual, uma parcela considerável da variação acompanha os movimentos do próprio animal~\cite{Stringer2019}. A dificuldade aqui é de princípio: uma mensagem bem comprimida é indistinguível de um fluxo aleatório para quem não conhece o código.

\section{Taxa: lento, mas confiável}
\label{sec:rate}

O código mais simples é o \emph{código de taxa}: o número está escrito em quantos disparos o neurônio emite numa janela $T$. Nem apagamentos nem tremores dos instantes afetam esse código. Mas ele tem um preço. Se o fluxo é de Poisson, o número de disparos $n$ numa janela tem média $rT$ e dispersão $\sqrt{rT}$:
\begin{empheq}[box=\keybox]{equation}
\frac{\sigma_n}{\bar n} \;=\; \frac{1}{\sqrt{r\,T}} .
\label{eq:rateerror}
\end{empheq}
Para conhecer a frequência com precisão de dez por cento, são necessários cem disparos; a vinte disparos por segundo, cinco segundos. Níveis vizinhos são distinguíveis se estão a um par de dispersões um do outro, por isso, até a frequência $r$, no tempo $T$ distinguem-se cerca de $\sqrt{rT}$ níveis: com $r=10$ e $T=1\,\text{s}$, são três ou quatro níveis, menos de dois bits.
O cérebro não trabalha tão devagar. Uma pessoa reconhece um animal numa imagem mostrada por um instante, e a resposta no cérebro aparece após cerca de $150\ms$~\cite{Thorpe1996}. Por estimativa grosseira, nesse tempo o sinal passa por cerca de uma dezena de estágios sucessivos de processamento, $10$--$20\ms$ por estágio. Com as frequências habituais do córtex, cada neurônio consegue emitir zero ou um disparo por estágio, e sua frequência numa janela dessas não está definida. Há duas saídas: usar muitos neurônios ou olhar para o tempo.

\emph{Muitos neurônios.} Suponha que $N$ neurônios levem o mesmo número e que o destinatário some seus disparos. Se seus ruídos são independentes, em~\eqref{eq:rateerror} entra $NrT$ no lugar de $rT$: cem neurônios com $r=20$ em $15\ms$ dão trinta disparos e precisão de cerca de vinte por cento. O espaço substitui o tempo. Porém, os ruídos de neurônios vizinhos do córtex não são independentes: o coeficiente de correlação dos números de disparos de um par de vizinhos costuma ser cerca de $0.1$~\cite{Zohary1994}. Se ele vale $c$, a variância da média de $N$ neurônios é
\begin{empheq}[box=\keybox]{equation}
\sigma^2_N \;=\; \sigma^2_1\,\frac{1+(N-1)\,c}{N}
\;\xrightarrow[N\to\infty]{}\; c\,\sigma^2_1 .
\label{eq:correlated}
\end{empheq}

\begin{keyblock}
\begin{proposition}[Limite da média]
\label{prop:pooling}
Com ruído correlacionado, a média sobre um grupo reduz o erro no máximo $1/\sqrt{c}$ vezes. Com $c=0.1$, são três vezes, e esse ganho já é alcançado com algumas dezenas de neurônios; acrescentar mais é inútil.
\end{proposition}
\end{keyblock}

Nem toda correlação é igualmente nociva. Só limita a precisão a parte do ruído comum que \emph{se parece com o próprio sinal}, isto é, que desloca o grupo inteiro do mesmo modo que o deslocaria uma mudança da grandeza medida~\cite{MorenoBote2014}. Quanto desse ruído existe de fato no cérebro ainda está sendo investigado.

Há também outra maneira de usar muitos neurônios: escrever o número em \emph{qual} neurônio disparou, como no detector da direção do som (fig.~\ref{fig:itd}). É o código de \emph{lugar}, o mais confiável que se conhece: nos neurônios sensoriais, o número do fio diz qual ponto da pele foi tocado ou qual a altura do som, e isso foi estabelecido muitas vezes. Ele é caro em número de neurônios, mas rápido: uma mensagem leva um atraso.

\section{Tempo: rápido, mas contado a partir de quê}

A segunda saída é escrever o número no instante do disparo. Tomemos o neurônio~\eqref{eq:glutneuron} e apliquemos a ele, a partir de $t=0$, uma entrada constante que elevaria a tensão até o nível $U$. A tensão cresce como $V=U\bigl(1-e^{-t/\tau}\bigr)$ e atinge o limiar $\theta$ no instante
\begin{empheq}[box=\keybox]{equation}
t_1 \;=\; \tau\,\ln\frac{U}{U-\theta}, \qquad U>\theta .
\label{eq:latency}
\end{empheq}
Entrada forte, disparo cedo; fraca, disparo tarde; sublimiar, nenhum. Com $\tau=10\ms$, a entrada $U=4\theta$ dá o primeiro disparo após $2.9\ms$; $U=2\theta$, após $6.9\ms$; $U=1.2\theta$, após $17.9\ms$. Um único disparo carrega um número.

Mas a partir de que instante contar $t_1$? Não há relógio, e o destinatário não sabe quando o estímulo começou. Há três respostas.

\emph{Latência relativa.} Dois neurônios veem o mesmo estímulo, e a diferença de suas latências $t_1^A-t_1^B$ depende só de suas entradas, não do instante de início. Os instantes são comparados por detectores de coincidência com atrasos (fig.~\ref{fig:itd}). Se $N$ neurônios emitiram um disparo cada, a própria \emph{ordem} de chegada carrega até $\log_2N!$ bits: para dez neurônios, cerca de vinte e dois bits, enquanto a resposta ``disparou ou não'' daria no máximo dez. Na retina, mostrou-se que a imagem pode ser reconstruída de forma rápida e confiável a partir das latências relativas dos primeiros disparos de seus neurônios de saída~\cite{Gollisch2008}.

\emph{Uma salva como início de palavra.} Uma mudança brusca do estímulo faz disparar quase ao mesmo tempo uma multidão de neurônios. Essa salva serve ela mesma de marca de início, como a pausa longa entre palavras no código Morse, e o destinatário conta as latências a partir dela.

\emph{Ritmo local.} No capítulo~\ref{sec:gaba}, o laço \nt{GLUT}--\nt{GABA} gerava um ritmo local com período de $12$--$50\ms$. Não é um relógio, mas dentro de seu ciclo um neurônio com entrada forte consegue disparar antes de um com entrada fraca, e~\eqref{eq:latency} vira um código de \emph{fase}: o número está escrito na parte do ciclo em que o disparo chegou. Esse código foi observado: neurônios responsáveis por um lugar definido no espaço disparam cada vez mais cedo no ciclo de um ritmo local lento à medida que o animal atravessa o ``seu'' lugar~\cite{OKeefe1993}. Quanto o cérebro usa a fase de modo geral ainda não se sabe.

\section{Quem escolhe o código é o destinatário}

Numa sequência de disparos há frequência, instantes e ordem. O que será lido é decidido pelo destinatário, e ele decide com um único parâmetro: sua constante de tempo $\tau$.

Tomemos a média da equação~\eqref{eq:glutneuron} no tempo. Se ao neurônio chegam fluxos independentes com frequências $r_j$, a tensão média e seu tremor relativo são
\begin{empheq}[box=\keybox]{equation}
\bar V \;=\; \tau\sum_j w_j\,r_j ,
\qquad
\frac{\sigma_V}{\bar V} \;\approx\; \frac{1}{\sqrt{2\,r_{\text{ent}}\,\tau}} ,
\label{eq:reader}
\end{empheq}
em que a segunda igualdade vale para pesos iguais, e $r_{\text{ent}}$ é a frequência total de todas as entradas. Quando $\tau$ é grande, a tensão quase não treme e acompanha a soma ponderada das frequências: o destinatário é um \emph{integrador}, lê frequências e esquece instantes. Quando $\tau$ é pequeno, a tensão tem tempo de cair entre disparos, e o limiar só é atingido quando vários disparos chegam dentro da janela~\eqref{eq:window}: o destinatário é um \emph{detector de coincidência}, lê instantes (fig.~\ref{fig:reader}). Mil entradas a cinco disparos por segundo com $\tau=10\ms$ dão um tremor de cerca de dez por cento, uma integração quase pura. Se a inibição, como no capítulo~\ref{sec:gaba}, encurtou $\tau$ para $2\ms$, o tremor passa de vinte por cento, e o neurônio começa a notar coincidências.

\begin{figure}[t]
\centering
\begin{tikzpicture}[>=Stealth,x=0.115cm,y=1cm]
% common input: the same spike train for both receivers
\foreach \s in {6,21,22,38,52,55.5,59,62.5,66,69.5,73,76.5,80,83.5,87,90.5,94,97.5}{\draw[thick,gray!40!black] (\s,5.25) -- (\s,5.65);}
\draw[gray!70!black] (0,5.25) -- (105,5.25);
\node[left,font=\small] at (-1,5.45) {entrada};
\node[above,font=\scriptsize,gray!40!black] at (13.5,5.65) {raro};
\node[above,font=\scriptsize,gray!40!black] at (21.5,6.0) {par};
\node[above,font=\scriptsize,gray!40!black] at (75,5.65) {frequente};
% axes and thresholds
\foreach \y/\th in {2.7/2.5,0/1.5}{
  \draw[gray!70!black] (0,\y) -- (106,\y) node[right,black] {\small $t$};
  \draw[densely dashed,gray!60!black] (0,\y+0.6*\th) -- (104,\y+0.6*\th) node[right,black,font=\small] {$\theta$};
}
\node[anchor=south west,font=\small] at (0,4.3) {$\tau=10\ms$: integrador};
\node[anchor=south east,font=\small] at (104,1.0) {$\tau=2\ms$: detector de coincidência};
% integrator: tau=10 ms, theta=2.5w
\begin{scope}[yshift=2.7cm]
\foreach \a/\b/\v in {0/6/0.000,6/21/1.000,21/22/1.223,22/38/2.107,38/52/1.425,52/55.5/1.351,55.5/59/1.952,59/62.5/2.376,62.5/66/0.000,66/69.5/1.000,69.5/73/1.705,73/76.5/2.201,76.5/80/0.000,80/83.5/1.000,83.5/87/1.705,87/90.5/2.201,90.5/94/0.000,94/97.5/1.000,97.5/104/1.705}{\draw[very thick,blue!60!black] plot[domain=\a:\b,samples=14] (\x,{0.6*\v*exp(-(\x-\a)/10)});}
\foreach \s/\p/\q in {6/0.000/1.000,21/0.223/1.223,22/1.107/2.107,38/0.425/1.425,52/0.351/1.351,55.5/0.952/1.952,59/1.376/2.376,62.5/1.674/2.500,62.5/2.500/0.000,66/0.000/1.000,69.5/0.705/1.705,73/1.201/2.201,76.5/1.551/2.500,76.5/2.500/0.000,80/0.000/1.000,83.5/0.705/1.705,87/1.201/2.201,90.5/1.551/2.500,90.5/2.500/0.000,94/0.000/1.000,97.5/0.705/1.705}{\draw[very thick,blue!60!black] (\s,0.6*\p) -- (\s,0.6*\q);}
\foreach \s in {62.5,76.5,90.5}{\draw[->,thick,red!70!black] (\s,1.58) -- (\s,1.95);}
\end{scope}
% coincidence: tau=2 ms, theta=1.5w
\begin{scope}[yshift=0cm]
\foreach \a/\b/\v in {0/6/0.000,6/21/1.000,21/22/1.001,22/38/0.000,38/52/1.000,52/55.5/1.001,55.5/59/1.174,59/62.5/1.204,62.5/66/1.209,66/69.5/1.210,69.5/73/1.210,73/76.5/1.210,76.5/80/1.210,80/83.5/1.210,83.5/87/1.210,87/90.5/1.210,90.5/94/1.210,94/97.5/1.210,97.5/104/1.210}{\draw[very thick,blue!60!black] plot[domain=\a:\b,samples=14] (\x,{0.6*\v*exp(-(\x-\a)/2)});}
\foreach \s/\p/\q in {6/0.000/1.000,21/0.001/1.001,22/0.607/1.500,22/1.500/0.000,38/0.000/1.000,52/0.001/1.001,55.5/0.174/1.174,59/0.204/1.204,62.5/0.209/1.209,66/0.210/1.210,69.5/0.210/1.210,73/0.210/1.210,76.5/0.210/1.210,80/0.210/1.210,83.5/0.210/1.210,87/0.210/1.210,90.5/0.210/1.210,94/0.210/1.210,97.5/0.210/1.210}{\draw[very thick,blue!60!black] (\s,0.6*\p) -- (\s,0.6*\q);}
\foreach \s in {22}{\draw[->,thick,red!70!black] (\s,0.98) -- (\s,1.35);}
\end{scope}
\foreach \x in {0,50,100}{\draw[gray!60!black] (\x,-0.06) -- (\x,0.06) node[below,font=\scriptsize] {\x\,ms};}
\end{tikzpicture}
\caption{A mesma entrada, dois destinatários diferentes. Cada disparo eleva a tensão em $w$, e depois ela escoa, como no modelo de neurônio~\eqref{eq:glutneuron}; as setas vermelhas são os disparos do destinatário. O destinatário lento ($\tau=10\ms$, $\theta=2.5w$) dispara onde há \emph{muitos} disparos e não nota o par. O rápido ($\tau=2\ms$, $\theta=1.5w$) dispara só com o par dentro da janela~\eqref{eq:window} e deixa passar o fluxo frequente, mas não coincidente.}
\label{fig:reader}
\end{figure}

As medições mostram que, no córtex ativo, a condutância da membrana aumenta várias vezes em relação ao repouso~\cite{Destexhe2003}. Logo, lá $\tau$ é mais curto, e os neurônios do córtex em regime de trabalho estão mais perto de detectores de coincidência do que de integradores.

\begin{keyblock}
\begin{proposition}[O código é definido pelo destinatário]
\label{prop:reader}
A mesma sequência de disparos carrega, para um destinatário lento, a frequência; para um rápido, os instantes; para o volume em que o neurotransmissor é liberado, um nível de concentração suavizado ao longo de segundos (capítulo~\ref{sec:serocomputer}). Qual código é usado não é determinado pelo remetente, mas pela constante de tempo do destinatário, e ela é ajustada pela inibição.
\end{proposition}
\end{keyblock}
\section{Pontos e traços: a sinapse como filtro}

Até agora, para nós, a sinapse era um fio com peso. Na verdade, sua resposta depende dos disparos anteriores, e isso dá à linguagem dos neurônios um segundo sinal.

\emph{Depressão: a sinapse transmite mudanças.} Cada liberação gasta uma fração $U$ do estoque de neurotransmissor pronto para liberação~\cite{Tsodyks1997}, e o estoque se recupera com constante de tempo $\tau_{\text{rec}}$ da ordem de centenas de milissegundos. Com frequência $r$, a amplitude da resposta a um disparo se estabiliza aproximadamente no nível $A(r)=A_0/(1+U r\tau_{\text{rec}})$, em que $A_0$ é a resposta da sinapse descansada. A corrente que o destinatário recebe por unidade de tempo é
\begin{empheq}[box=\keybox]{equation}
r\,A(r) \;=\; \frac{A_0\,r}{1+U r\,\tau_{\text{rec}}}
\;\xrightarrow[r\to\infty]{}\; \frac{A_0}{U\tau_{\text{rec}}} .
\label{eq:depression}
\end{empheq}
Com $U=0.5$ e $\tau_{\text{rec}}=0.8\,\text{s}$, a saturação começa já em $1/(U\tau_{\text{rec}})=2.5$ disparos por segundo. Se a frequência subiu de cinco para vinte, a corrente de regime cresce só um terço. Em compensação, os primeiros disparos na nova frequência chegam enquanto o estoque ainda é o antigo, e a corrente salta momentaneamente quatro vezes e depois cai em $\tau_{\text{rec}}/(1+Ur\tau_{\text{rec}})\approx90\ms$ (fig.~\ref{fig:depression}).

Essa sinapse não transmite a frequência, mas sua \emph{mudança}, essencialmente a mudança relativa $\Delta r/r$~\cite{Abbott1997depression}. Para um eletrônico, é um filtro passa-altas instalado no próprio fio.

\begin{figure}[t]
\centering
\begin{tikzpicture}[>=Stealth,x=12.5cm,y=1cm]
% input rate
\draw[->,gray!70!black] (0,2.6) -- (0.9,2.6) node[right,black] {\small $t$};
\draw[->,gray!70!black] (0,2.6) -- (0,4.3) node[above,black] {\small $r$};
\draw[very thick,blue!60!black] (0,2.9) -- (0.2,2.9) -- (0.2,3.8) -- (0.8,3.8);
\node[left,font=\scriptsize] at (0,2.9) {$5$};
\node[left,font=\scriptsize] at (0,3.8) {$20$};
\node[right,font=\small,blue!60!black] at (0.4,4.1) {frequência na entrada da sinapse};
% output current
\draw[->,gray!70!black] (0,0) -- (0.9,0) node[right,black] {\small $t$};
\draw[->,gray!70!black] (0,0) -- (0,2.45) node[left,black] {\small $rA$};
\draw[densely dashed,gray!60!black] (0,0.667) -- (0.8,0.667);
\draw[very thick,red!70!black] (0,0.5) -- (0.2,0.5) -- (0.2,2.0)
   plot[domain=0.2:0.8,samples=60] (\x,{0.667+(2.0-0.667)*exp(-(\x-0.2)/0.089)});
\node[left,font=\scriptsize] at (0,0.5) {$1.7$};
\node[left,font=\scriptsize] at (0,2.0) {$6.7$};
\node[right,font=\scriptsize] at (0.8,0.667) {$2.2$};
\node[right,font=\small,red!70!black] at (0.28,1.6) {corrente ao destinatário: salto e queda em $\approx90\ms$};
\draw[gray!60!black] (0.2,-0.08) -- (0.2,0.08); \node[below,font=\scriptsize] at (0.2,0) {$0.2\,\text{s}$};
\draw[gray!60!black] (0.8,-0.08) -- (0.8,0.08); \node[below,font=\scriptsize] at (0.8,0) {$0.8\,\text{s}$};
\end{tikzpicture}
\caption{Sinapse com depressão pela fórmula~\eqref{eq:depression}, com $U=0.5$, $\tau_{\text{rec}}=0.8\,\text{s}$; corrente em unidades de $A_0$ por segundo; a linha tracejada é a corrente de regime: ela subiu só de $1.7$ para $2.2$, e no instante do salto, até $6.7$. A sinapse informa ao destinatário que a frequência mudou, não qual ela é.}
\label{fig:depression}
\end{figure}

\emph{Facilitação: a rajada como traço.} Em outras sinapses, a probabilidade de liberação é pequena, mas após cada disparo cresce por dezenas de milissegundos. Um disparo isolado por uma sinapse dessas se perde na maioria das vezes. Já uma \emph{rajada}, vários disparos seguidos com intervalos de poucos milissegundos, passa quase com certeza. Mesmo sem facilitação, a probabilidade de que se percam todos os $k$ disparos da rajada é
\begin{empheq}[box=\keybox]{equation}
P_{\text{perda}} \;=\; (1-p)^k ,
\label{eq:burst}
\end{empheq}
com $p=0.2$, isso dá $0.8$ para um disparo isolado e $0.41$ para uma rajada de quatro; a facilitação a reduz ainda mais. Assim o neurônio ganha um segundo sinal: disparo isolado, ponto; rajada, traço. A sinapse com depressão transmite melhor o primeiro disparo após uma pausa; a facilitadora deixa passar rajadas e abafa pontos isolados. Que as rajadas são transmitidas de forma mais confiável foi medido; que o cérebro as usa de fato como um sinal à parte é uma hipótese plausível~\cite{Lisman1997}.

\section{Como falam os neurônios \nt{GABA}}

Os mais numerosos deles no córtex são os rápidos~\cite{Tremblay2016}. Seus disparos são mais curtos que os dos neurônios \nt{GLUT}, e eles podem emiti-los de forma regular e por muito tempo, a centenas por segundo, quase sem se cansar. Suas sinapses são mais confiáveis: a conexão com cada destinatário consiste em muitos pontos de liberação, e um disparo quase nunca se perde. Suas saídas se instalam perto do local onde nasce o disparo do destinatário, de modo que controlam não como o destinatário soma as entradas, mas se ele vai disparar e quando.

Eles mesmos recebem excitação de muitos neurônios \nt{GLUT} vizinhos e informam a atividade total ao redor: exatamente o que é preciso para a divisão do capítulo~\ref{sec:gaba}. Por fim, os neurônios \nt{GABA} rápidos são ligados entre si e disparam juntos com facilidade; são eles que definem o ritmo local de que se falou acima~\cite{Cardin2009}.

Na linguagem do código Morse, os neurônios \nt{GABA} não respondem pelas letras, mas pelas \emph{pausas}. Eles cortam o fluxo contínuo de disparos \nt{GLUT} em janelas dentro das quais o destinatário pode disparar, estreitam as janelas de coincidência e abafam o fluxo inteiro quando ele fica alto demais.

\begin{keyblock}
\begin{proposition}[Divisão de papéis]
\label{prop:punctuation}
Os neurônios \nt{GLUT} levam o conteúdo da mensagem: \emph{o quê} e \emph{quanto}. Os neurônios \nt{GABA} rápidos levam sua marcação: \emph{quando} se pode falar e \emph{com que volume}; mais uma classe, inibindo os inibidores, decide \emph{para quem} o portão está aberto. Partes isoladas dessa divisão foram medidas; como regra geral, é uma hipótese de trabalho.
\end{proposition}
\end{keyblock}

Há também outros neurônios \nt{GABA}, mais lentos, cujas saídas se instalam em entradas individuais do destinatário. Eles funcionam como o veto~\eqref{eq:veto} do capítulo~\ref{sec:gaba}: fecham um fluxo de disparos \nt{GLUT} sem tocar nos demais.

A divisão em neurônios \nt{GABA} rápidos e lentos é um quadro antigo, e incompleto. Hoje se distinguem no córtex pelo menos três grandes classes~\cite{Tremblay2016}, e a terceira é organizada de modo inesperado: ela não inibe neurônios \nt{GLUT}, mas outros neurônios \nt{GABA}, sobretudo os que fecham entradas~\cite{Pfeffer2013}. Inibição da inibição é uma dupla negação. Um neurônio desses \emph{abre} o portão sem enviar ao destinatário um único disparo excitatório. Ele é ligado por sinais de outras áreas do córtex e por canais de difusão, e assim funciona como a entrada de habilitação de toda uma região da rede: enquanto está ativo, os vetos são suspensos e o fluxo passa~\cite{Pi2013}. No apêndice~\ref{sec:othermediators}, esse truque é chamado de desinibição.
\section{Uma linguagem econômica}

A última coisa que se sabe com segurança sobre a linguagem dos neurônios é que ela é cara. Um disparo, junto com o trabalho sináptico que provoca, custa cerca de $10^9$ moléculas de ATP (estimativa para o córtex de roedores~\cite{Attwell2001}), e uma molécula dá cerca de $10^{-19}$ joule. Logo, um disparo custa da ordem de $E_{\text{disp}}\sim10^{-10}$ joule. O cérebro inteiro consome cerca de $20$ watts, e, se metade vai para a transmissão de sinais, $P\sim10$ watts, a frequência média de um neurônio é
\begin{empheq}[box=\keybox]{equation}
\bar r \;\approx\; \frac{P}{N\,E_{\text{disp}}}
\;\sim\; \frac{10\ \text{W}}{8.6\cdot10^{10}\cdot10^{-10}\ \text{J}}
\;\approx\; 1\ \text{disparo por segundo}.
\label{eq:energy}
\end{empheq}
Estimativas mais detalhadas para o córtex dão ainda menos~\cite{Lennie2003}. O código do cérebro tem de ser \emph{esparso}: só uma pequena fração dos neurônios pode trabalhar rápido.

De~\eqref{eq:spikeentropy} se vê que isso não é tão ruim. Com precisão de $1\ms$, o disparo de um neurônio que trabalha a $100$ por segundo carrega no máximo cinco bits, e o de um neurônio que trabalha uma vez por segundo, até onze. Se é preciso pagar pelos disparos, vale a pena falar pouco. O córtex de fato troca precisão por energia: quando falta alimento, os neurônios mantêm a frequência de disparos, mas gastam menos com correntes sinápticas, e suas respostas ficam menos precisas~\cite{Padamsey2022}.

No código Morse, as letras frequentes são curtas: a letra mais frequente do texto inglês, o E, é um único ponto. O código dos neurônios, até onde se sabe, segue a mesma regra de outra forma: o que é esperado custa poucos disparos~\cite{Barlow1961}. Com um estímulo constante, o neurônio reduz aos poucos a frequência; os sensores do capítulo~\ref{sec:glutcomputer} respondem a mudanças, não a níveis; e os neurônios \nt{DOPA} do capítulo~\ref{sec:neurontypes} informam só o erro de predição de recompensa.

\begin{keyblock}
\begin{proposition}[Economia]
\label{prop:economy}
A energia limita a frequência média de um neurônio a algo da ordem de um disparo por segundo. Por isso a linguagem dos neurônios \nt{GLUT} é esparsa e gasta disparos com notícias: com o que mudou ou não foi previsto. O limite energético foi medido; que o código do cérebro esteja ajustado para o máximo de bits por joule é uma hipótese bem fundamentada.
\end{proposition}
\end{keyblock}

\section{Resumo}

\begin{table}[t]
\centering
\footnotesize
\setlength{\tabcolsep}{4pt}
\renewcommand{\arraystretch}{1.15}
\begin{tabular}{>{\raggedright}p{2.6cm}>{\raggedright}p{2.3cm}>{\raggedright}p{3.0cm}>{\raggedright}p{2.0cm}>{\raggedright\arraybackslash}p{3.6cm}}
\toprule
Modo de escrita & O que carrega & Quem lê & Tempo por mensagem & O que se sabe \\
\midrule
número do neurônio que disparou (código de lugar) & qual valor & qualquer destinatário; detectores de coincidência & um atraso & estabelecido nos neurônios sensoriais e na determinação da direção do som \\
frequência de um neurônio & uma grandeza & integrador com $\tau$ grande; volume (cap.~\ref{sec:serocomputer}) & centenas de ms a segundos & estabelecido; lento demais para um estágio de processamento de $10$--$20\ms$ \\
frequência de um grupo & uma grandeza & neurônio com milhares de entradas & $10$--$50\ms$ & estabelecido; o ganho é limitado pelas correlações~\eqref{eq:correlated} \\
tempo do primeiro disparo, ordem & grandeza, ordem & detectores de coincidência com atrasos & um atraso & estabelecido na entrada do sistema nervoso; no córtex, em discussão \\
fase no ritmo local & grandeza, posição & neurônios com ritmo comum & ciclo de $12$--$50\ms$ ou mais lento & observado em algumas áreas; papel geral é hipótese \\
rajada (``traço'') & ``importante'': entrega confiável & sinapse facilitadora & $\sim10\ms$ & confiabilidade medida; sinal à parte é hipótese \\
mudança de frequência & notícia & sinapse com depressão & $\sim0.1\,\text{s}$ & estabelecido \\
disparos dos neurônios \nt{GABA} rápidos & quando e com que volume & todos os neurônios \nt{GLUT} vizinhos & $1$--$30\ms$ & ritmo e divisão medidos; ``pontuação'' como regra é hipótese \\
\bottomrule
\end{tabular}
\caption{Modos pelos quais os neurônios \nt{GLUT} e \nt{GABA} escrevem mensagens com disparos. A coluna da direita separa o medido do suposto.}
\label{tab:code}
\end{table}

A tabela~\ref{tab:code} reúne o que sabemos; dela se vê também o que não sabemos. Não sabemos quanto o córtex usa o tempo exato dos disparos, e não só as frequências dos grupos. Não sabemos se os neurônios têm ``palavras'': combinações estáveis de disparos de muitos neurônios, lidas como um todo. E não sabemos se a linguagem é a mesma em diferentes partes do cérebro. O código Morse se aprende numa noite, porque sua tabela é conhecida. A tabela da linguagem dos neurônios ainda está por ser montada, e quem vai montá-la, muito provavelmente, são engenheiros de comunicações.

\section{Exercícios}

\begin{exercise}[\lvlA]
\label{ex:04:1}
\emph{A capacidade em números.} $\Delta t=1\ms$. (a)~Quantos bits por joule dá um neurônio com frequência de $1$ disparo por segundo, com o custo do disparo de~\eqref{eq:energy}? (b)~Com $r=10$, quantas vezes menos que o limite~\eqref{eq:spikeentropy} extrai um destinatário que só conta os disparos num segundo?
\end{exercise}

\begin{exercise}[\lvlB]
\label{ex:04:2}
\emph{Frequência ótima.} Um neurônio gasta a potência $P_0+rE_{\text{disp}}$, em que $P_0$ é a outra metade dos $20$~W dividida por todos os neurônios, $E_{\text{disp}}=10^{-10}$~J. Com que frequência $r$ o número de bits por joule $R_{\max}(r)/(P_0+rE_{\text{disp}})$ por~\eqref{eq:spikeentropy}, com $\Delta t=1\ms$, é máximo?
\end{exercise}

\begin{exercise}[\lvlB]
\label{ex:04:3}
\emph{Que correlação é nociva.} $N=1000$ neurônios, variância $\sigma^2=1$, correlação par a par $c=0.1$. Calcule a informação de Fisher $\mathcal I=f'^{\mathsf T}\Sigma^{-1}f'$ ($\Sigma$ é a matriz de covariâncias) para inclinações iguais $f'=\mathbf 1$ e para inclinações $\pm1$ com $\mathbf 1^{\mathsf T}f'=0$. Quantas vezes elas diferem?
\end{exercise}

\begin{exercise}[\lvlB]
\label{ex:04:4}
\emph{Modelagem: detector de coincidência.} O neurônio~\eqref{eq:glutneuron} recebe $1000$ entradas de Poisson a $5$ disparos por segundo cada, $w=1$; o limiar $\theta$ é tal que a tensão livre o cruza uma vez por segundo. Por simulação, encontre quantas entradas simultâneas $m$ acendem o neurônio com probabilidade $1/2$ para $\tau=10$ e $2\ms$.
\end{exercise}

\begin{exercise}[\lvlB]
\label{ex:04:5}
\emph{Projeto: detector de mudanças relativas.} Ajuste a sinapse~\eqref{eq:depression}: a corrente de regime a $40$ disparos por segundo não mais que $10\,\%$ acima da corrente a $10$, e, após um salto para $40$, a corrente chega ao novo nível com constante de tempo de no máximo $50\ms$. Qual é o menor $\tau_{\text{rec}}$, e com que $U$?
\end{exercise}

\begin{exercise}[\lvlA]
\label{ex:04:6}
\emph{Tempo do primeiro disparo e rajadas.} (a)~Por~\eqref{eq:latency}, encontre a entrada com a qual o primeiro disparo chega exatamente após uma constante de tempo. De que ela depende? (b)~Com probabilidade de liberação $p=0.2$, quantos disparos a rajada precisa ter para que se percam todos com probabilidade abaixo de $5\,\%$?
\end{exercise}

\begin{exercise}[\lvlC]
\label{ex:04:7}
\emph{Pesquisa: quantos bits há na disposição.} O neurônio é feito de células independentes~\eqref{eq:spikeentropy}, $r=10$, $\Delta t=1\ms$. (a)~Decomponha a entropia de uma ``palavra'' de $20$ células na entropia do número de disparos e no resto. (b)~Simule a estimativa da entropia pelas frequências das palavras a partir de $N=30$ e $10^4$ registros. Quanto ela é subestimada?
\end{exercise}

\chapter{O que os neurônios \nt{SERO} calculam}
\label{sec:serocomputer}
\begin{chaptergoals}
\item como um marca-passo com receptor inibitório faz NÃO e NOU com um só neurônio \nt{SERO};
\item por que a liberação num volume deixa poucos fios independentes, dados por $\ell\sim\sqrt{D\tau}$;
\item como a concentração suaviza os disparos num nível, e como a autoinibição poderia mantê-lo estável;
\item por que o \nt{SERO} é candidato ao horizonte de planejamento, e como o horizonte define a paciência.
\end{chaptergoals}

\section{O primeiro canal de difusão}

Até aqui, todos os neurônios do livro falaram pelo barramento de endereços: \nt{GLUT} e \nt{GABA} enviavam disparos a destinatários determinados, e descobrimos o que essa rede calcula e em que linguagem ela fala. Agora passamos ao segundo barramento da seção~\ref{sec:buses}, o de difusão. Seu primeiro canal é o \nt{SERO}. Como antes, só nos permitimos o que existe num organismo vivo.

Neste capítulo aparecem três dispositivos que depois serão usados por todos os canais de difusão: um neurônio que, com um aporte de fundo constante, funciona sozinho até ser parado; um fio que, na verdade, é um volume de tecido; e uma concentração lenta no lugar de disparos isolados. \nt{DOPA}, \nt{NORA} e \nt{ACET}, nos capítulos~\ref{sec:dofa}--\ref{sec:acet}, serão montados com as mesmas peças.

Do ponto de vista do engenheiro, o neurônio \nt{SERO} tem quatro diferenças importantes.

\emph{Primeira: quem escolhe o sinal é o destinatário.} A serotonina tem receptores dos dois sinais, e a regra~\eqref{eq:dalesign} não vale aqui: o sinal não é definido pelo remetente, mas pelo receptor do destinatário (proposição~\ref{prop:receptor})~\cite{BarnesSharp1999}.

\emph{Segunda: o neurônio \nt{SERO} funciona sozinho se houver um fundo constante.} Na vigília, ele emite disparos regulares, algumas vezes por segundo, e não precisa de um empurrão para cada disparo: é um marca-passo. Mas precisa de um aporte de fundo permanente. Numa fatia de cérebro, onde esse aporte falta, a maioria dessas células fica em silêncio; o ritmo regular volta se aplicarmos noradrenalina pelos receptores $\alpha_1$, e o bloqueio desses receptores o apaga~\cite{VandermaelenAghajanian1983}. No organismo, esse fundo vem do \nt{NORA}. Para o circuito, é um oscilador que precisa de alimentação: enquanto ela existe, ele funciona sozinho até que a inibição o pare.

\emph{Terceira: ele é lento.} Quase todos os receptores de serotonina são metabotrópicos: não abrem um canal por conta própria, mas disparam uma cadeia de reações dentro da célula. A resposta é lenta: a corrente inibitória pelo receptor principal sobe e desce em cerca de um segundo~\cite{CourtneyFord2016}, dezenas de vezes mais que a resposta de uma sinapse rápida \nt{GLUT}.

\emph{Quarta: ele libera o neurotransmissor não num fio, mas num volume.} Nas regiões para onde vão as saídas dos neurônios \nt{SERO}, a serotonina muitas vezes não sai numa fenda estreita, mas no espaço extracelular, e se espalha em volta~\cite{Bunin1998}. O destinatário sente a concentração e não tem como saber de qual remetente veio a molécula.

Com esses tempos, os disparos isolados já não importam, por isso vamos descrever a rede por taxas $r_j$ e concentrações. A concentração $c$ de serotonina no volume em que os neurônios $j$ liberam o neurotransmissor cresce com a liberação e diminui porque o neurotransmissor é recaptado:
\begin{empheq}[box=\keybox]{equation}
\tau_c\,\frac{\dd c}{\dd t} \;=\; -c \;+\; \kappa\sum_j r_j ,
\qquad
r_i \;=\; f_i\bigl(b_i + s_i\,g\,c\bigr), \quad s_i=\pm1 .
\label{eq:seroneuron}
\end{empheq}
Essa é mais uma maneira de ler um trem de disparos, complementando a proposição~\ref{prop:reader}: o volume não vê nem os disparos isolados nem sua ordem, apenas a taxa média no tempo $\tau_c$. A primeira equação é o mesmo circuito RC da membrana: $\tau_c$ é o tempo de vida do neurotransmissor no volume; ele não foi medido diretamente, e o tomamos da ordem de um segundo; $\kappa$ é quanto neurotransmissor um disparo fornece. A segunda descreve o destinatário: $b_i$ é seu próprio aporte, positivo no marca-passo (inclui também a entrada de fundo constante vinda do \nt{NORA}); $g$ é a sensibilidade do receptor; o sinal $s_i$ é escolhido pelo destinatário; $f_i$ é uma característica de transferência não decrescente.

\section{NÃO com um só neurônio}

Tomemos um marca-passo com receptor inibitório, $s=-1$. Enquanto não há serotonina em volta, ele funciona com seu próprio aporte $b$. Quando a serotonina aparece, o aporte cai em $gc$ e, se $gc>b$, o neurônio se cala. Isso é um NÃO: há saída quando não há entrada. Gastamos um único neurônio (fig.~\ref{fig:serogates}). A proposição~\ref{prop:monotone} não se aplica aqui: ela exigia pesos positivos.

Se sobre o mesmo marca-passo age serotonina de dois volumes diferentes, ele fica em silêncio quando há serotonina em pelo menos um deles. Isso é um NOU, e com NOU se monta qualquer função lógica. O código de dois fios da rede \nt{GLUT} não é necessário aqui: a ausência de sinal é calculada por neurônios que funcionam até serem parados.

\begin{figure}[t]
\centering
\begin{tikzpicture}[>=Stealth,
  nrn/.style={circle,draw,very thick,fill=blue!6,minimum size=0.95cm,inner sep=0pt},
  pace/.style={nrn,double,double distance=1.2pt},
  inh/.style={-{Bar[width=7pt]},thick,red!70!black}]
\node[pace] (N1) at (0,0) {};
\draw[inh] (-1.6,0) node[left,black] {$c$} -- (N1);
\draw[->,very thick,red!70!black] (N1) -- (1.3,0) node[right,black] {$\bar c$};
\node[below=0.55cm] at (0,0) {\small NÃO};
\node[pace] (N2) at (5,0) {};
\draw[inh] (3.4,0.45) node[left,black] {$c_1$} -- (N2);
\draw[inh] (3.4,-0.45) node[left,black] {$c_2$} -- (N2);
\draw[->,very thick,red!70!black] (N2) -- (6.3,0) node[right,black] {$\overline{c_1\vee c_2}$};
\node[below=0.55cm] at (5,0) {\small NOU};
\end{tikzpicture}
\caption{Lógica com neurônios \nt{SERO}. Círculo duplo: marca-passo; com um aporte de fundo constante, funciona sozinho até ser inibido. Linha com barra: receptor inibitório, $s=-1$; $c$, $c_1$, $c_2$ são as concentrações de serotonina nos volumes em torno do destinatário. À esquerda, NÃO; à direita, NOU.}
\label{fig:serogates}
\end{figure}

\begin{keyblock}
\begin{proposition}[Completude da lógica \nt{SERO}]
\label{prop:serocomplete}
Se a rede tem marca-passos com receptores inibitórios e as liberações em volumes diferentes não se misturam, a rede~\eqref{eq:seroneuron} pode calcular qualquer função lógica, e cada bit é transmitido por um único fio.
\end{proposition}
\end{keyblock}

Logicamente, a rede \nt{SERO} é mais forte que a \nt{GLUT}, mas a condição ``as liberações em volumes diferentes não se misturam'' não se cumpre no cérebro (seção~\ref{sec:volume}).

\section{Realimentação negativa}

Os neurônios serotoninérgicos têm receptores inibitórios neles mesmos~\cite{Sprouse1987}. A serotonina que liberaram volta a eles e os freia. Para o engenheiro, o esquema é conhecido: um amplificador com realimentação negativa (fig.~\ref{fig:serofeedback}).

O que aqui é medido e o que é modelo. No próprio núcleo da rafe, onde ficam os neurônios \nt{SERO}, a serotonina inibe os vizinhos não por um volume comum, mas ponto a ponto, por sinapses: o transportador a remove depressa e não a deixa acumular fora da fenda. Uma liberação isolada produz apenas uma pausa curta na descarga, e a taxa média não muda com isso~\cite{CourtneyFord2016}. Por isso o laço abaixo, fechado por um volume comum, é um modelo: calculamos o que essa realimentação faria se operasse no nível das taxas médias.

\begin{figure}[t]
\centering
\begin{tikzpicture}[>=Stealth,
  nrn/.style={circle,draw,very thick,fill=blue!6,minimum size=0.95cm,inner sep=0pt},
  pace/.style={nrn,double,double distance=1.2pt},
  blk/.style={draw,very thick,rounded corners=2pt,fill=orange!10,minimum height=0.9cm,minimum width=2.2cm,align=center,font=\small},
  inh/.style={-{Bar[width=7pt]},thick,red!70!black}]
\node[pace] (N) at (0,0) {};
\node[blk] (V) at (3.6,0) {volume\\$\tau_c$};
\draw[->,very thick] (N) -- node[above] {\small $r$} (V);
\draw[->,very thick] (V) -- (6.0,0) node[right] {\small $c$: a todos os destinatários};
\draw[inh] (V.south) -- ++(0,-0.9) -| node[pos=0.25,below] {\small $-g\,c$} (N.south);
\draw[->,thick,blue!60!black] (-1.7,0) node[left,black] {\small $b$} -- (N);
\end{tikzpicture}
\caption{Neurônio \nt{SERO} com realimentação pelo próprio neurotransmissor: a concentração $c$ vai a todos os destinatários e inibe o próprio neurônio. No modelo, é um regulador de nível; no cérebro, esse papel do laço é uma hipótese.}
\label{fig:serofeedback}
\end{figure}

Vamos calcular o que esse laço faz. Seja a característica de transferência linear, $f(u)=au$ para $u>0$. No equilíbrio, $c=\kappa r$ e $r=a(b-gc)$. Substituindo uma na outra, obtemos
\begin{empheq}[box=\keybox]{equation}
r \;=\; \frac{a\,b}{1+G}, \qquad G \;=\; a\,g\,\kappa .
\label{eq:feedback}
\end{empheq}
A grandeza $G$ é o ganho de laço. É a mesma fórmula de qualquer amplificador com realimentação negativa, e ela dá os mesmos dois benefícios.

\emph{Robustez à dispersão.} Suponha que a excitabilidade do neurônio $a$ tenha variado de uma fração $\varepsilon$. Sem realimentação, a taxa variaria da mesma fração. Com realimentação, para $G\gg1$, a taxa $r\approx b/(g\kappa)$ quase não depende de $a$: sua variação é $1+G$ vezes menor. Com $G=9$, a dispersão é reduzida dez vezes.

\emph{Rapidez.} Substituindo $r=a(b-gc)$ na primeira equação~\eqref{eq:seroneuron}, obtemos $\tau_c\,\dd c/\dd t=-(1+G)\,c+\kappa ab$. A concentração se aproxima do seu nível com constante de tempo $\tau_c/(1+G)$, $1+G$ vezes mais depressa que sem realimentação.

Eis o primeiro cálculo que o sistema \nt{SERO} poderia fazer: manter o nível geral de serotonina no cérebro numa marca fixa, como um termostato mantém a temperatura. É uma hipótese: no experimento, a autoinibição por enquanto só aparece como pausas curtas que não alteram a taxa média.

\section{De disparos a nível}

O segundo cálculo está escondido na primeira equação~\eqref{eq:seroneuron}. Os neurônios emitem disparos isolados, mas o destinatário sente a concentração, que os suaviza ao longo do tempo $\tau_c$. Se a taxa da fonte muda, a concentração a acompanha com atraso:
\begin{equation}
c(t) \;=\; \frac{\kappa}{\tau_c}\int_{-\infty}^{t} e^{-(t-t')/\tau_c}\,\sum_j r_j(t')\,\dd t' .
\label{eq:lowpass}
\end{equation}
É a média móvel da atividade das fontes nos últimos poucos $\tau_c$: um filtro passa-baixas (fig.~\ref{fig:lowpass}). É assim que o conversor digital-analógico mais simples passa pulsos por um circuito RC e transforma sua taxa num nível. O sistema \nt{SERO} faz o mesmo com centenas de milhares de neurônios ao mesmo tempo.

Essa grandeza é também uma memória. Depois que as fontes se calam, a concentração se mantém por um tempo da ordem de $\tau_c$. Não é um bit, mas um rastro que se desfaz aos poucos.

\begin{figure}[t]
\centering
\begin{tikzpicture}[>=Stealth,x=1.45cm,y=1cm]
\draw[gray!70!black] (0,3.2) -- (8.1,3.2);
\node[left,font=\small] at (-0.05,3.4) {disparos};
\node[above,font=\scriptsize,gray!40!black] at (1,3.65) {$2$ por segundo};
\node[above,font=\scriptsize,gray!40!black] at (3.5,3.65) {$8$ por segundo};
\node[above,font=\scriptsize,gray!40!black] at (6.5,3.65) {$2$ por segundo};
\draw[->,gray!70!black] (0,0) -- (8.3,0) node[right,black] {\small $t$};
\draw[->,gray!70!black] (0,0) -- (0,2.75) node[left,black] {\small $c$};
\foreach \s in {0.136,0.529,0.760,1.223,1.714,1.748,1.755,2.663,2.701,2.734,3.414,3.493,3.719,3.800,3.928,3.948,4.074,4.327,4.420,4.589,4.728,4.736,4.914,5.025,5.205,5.220,6.224,6.544,7.178}{\draw[thick,blue!60!black] (\s,3.2) -- (\s,3.6);}
\foreach \a/\b/\v in {0.000/0.136/0.000,0.136/0.529/1.000,0.529/0.760/1.675,0.760/1.223/2.330,1.223/1.714/2.466,1.714/1.748/2.509,1.748/1.755/3.425,1.755/2.663/4.402,2.663/2.701/2.775,2.701/2.734/3.672,2.734/3.414/4.553,3.414/3.493/3.306,3.493/3.719/4.055,3.719/3.800/4.235,3.800/3.928/4.905,3.928/3.948/5.316,3.948/4.074/6.211,4.074/4.327/6.476,4.327/4.420/6.028,4.420/4.589/6.493,4.589/4.728/6.483,4.728/4.736/6.642,4.736/4.914/7.589,4.914/5.025/7.351,5.025/5.205/7.579,5.205/5.220/7.331,5.220/6.224/8.221,6.224/6.544/4.012,6.544/7.178/3.914,7.178/8.000/3.076}{\draw[very thick,orange!85!black] plot[domain=\a:\b,samples=10] (\x,{0.25*\v*exp(-(\x-\a))});}
\foreach \s/\p/\q in {0.136/0.000/1.000,0.529/0.675/1.675,0.760/1.330/2.330,1.223/1.466/2.466,1.714/1.509/2.509,1.748/2.425/3.425,1.755/3.402/4.402,2.663/1.775/2.775,2.701/2.672/3.672,2.734/3.553/4.553,3.414/2.306/3.306,3.493/3.055/4.055,3.719/3.235/4.235,3.800/3.905/4.905,3.928/4.316/5.316,3.948/5.211/6.211,4.074/5.476/6.476,4.327/5.028/6.028,4.420/5.493/6.493,4.589/5.483/6.483,4.728/5.642/6.642,4.736/6.589/7.589,4.914/6.351/7.351,5.025/6.579/7.579,5.205/6.331/7.331,5.220/7.221/8.221,6.224/3.012/4.012,6.544/2.914/3.914,7.178/2.076/3.076}{\draw[very thick,orange!85!black] (\s,0.25*\p) -- (\s,0.25*\q);}
\draw[thick,densely dashed,black] plot[domain=0:2,samples=30] (\x,{0.25*2*(1-exp(-\x))})
  plot[domain=2:5,samples=40] (\x,{0.25*(8-6.271*exp(-(\x-2)))})
  plot[domain=5:8,samples=40] (\x,{0.25*(2+5.688*exp(-(\x-5)))});
\foreach \x in {0,2,5,8}{\draw[gray!60!black] (\x,-0.05) -- (\x,0.05) node[below=3pt,font=\scriptsize] {\x\,s};}
\end{tikzpicture}
\caption{De disparos a nível. Em cima, os disparos dos neurônios \nt{SERO}: dois segundos raros, três segundos frequentes, depois raros de novo. Embaixo, a concentração de serotonina~\eqref{eq:lowpass} com $\tau_c=1\,\text{s}$. Linha tracejada: o mesmo se os disparos viessem perfeitamente regulares.}
\label{fig:lowpass}
\end{figure}

\section{O fio é um volume}
\label{sec:volume}

Voltemos à condição da proposição~\ref{prop:serocomplete}. A serotonina liberada nas proximidades por remetentes diferentes se soma numa única concentração.

\emph{O fio aqui não é uma fibra, mas um volume de tecido.} Dois sinais só continuam distintos se forem liberados em regiões mais afastadas do que a distância que o neurotransmissor consegue percorrer ao se espalhar. No tempo $\tau$, uma molécula com coeficiente de difusão $D$ percorre
\begin{empheq}[box=\keybox]{equation}
\ell \;\sim\; \sqrt{D\,\tau}\, .
\label{eq:diffusionlength}
\end{empheq}
A tortuosidade dos espaços extracelulares torna a difusão de uma molécula pequena cerca de $2.6$ vezes mais lenta~\cite{Sykova2008}, e o coeficiente de difusão da própria serotonina no tecido não foi medido; pelo tamanho da molécula, estimamos que seja de algumas unidades de $10^{-6}$~cm$^2$/s. Se o sinal dura $\tau\sim1$~s (também uma estimativa), então $\ell\sim\sqrt{3\cdot10^{-6}}$~cm $\approx20$~\textmu m. O endereço numa rede assim é um \emph{lugar}: o destinatário só sabe de quem é o sinal pelo lugar onde ele próprio está.

Os neurônios \nt{SERO} reais são construídos de modo que quase não sobram esses endereços separados. A saída de cada um se espalha por enormes áreas do cérebro: os ramos de uma única célula chegam a muitas regiões distantes entre si~\cite{GagnonParent2014}. Os locais de liberação por célula são, por estimativa, cerca de $10^5$ (tabela~\ref{tab:types}); não foram contados diretamente. Os ``fios'' de neurônios \nt{SERO} diferentes se sobrepõem e se misturam. Mesmo assim, não se fundem num só fio: os neurônios \nt{SERO} se dividem em alguns subsistemas, que enviam a saída a regiões diferentes e respondem de modo diferente ao mesmo evento~\cite{Ren2018}. Mas esses subsistemas são poucos, não centenas de milhares.

\begin{keyblock}
\begin{proposition}[Número de fios independentes]
\label{prop:volumewires}
Numa rede com transmissão por volume, o número de sinais independentes não é limitado pelo número de neurônios, mas pelo número de regiões disjuntas de tamanho $\ell\sim\sqrt{D\tau}$ nas quais eles liberam o neurotransmissor. Se a saída de cada neurônio se espalha por um volume grande, a rede inteira transmite apenas algumas variáveis comuns.
\end{proposition}
\end{keyblock}

Por isso, no cérebro não há com que montar os circuitos lógicos da proposição~\ref{prop:serocomplete}. Já o regulador~\eqref{eq:feedback} e o filtro~\eqref{eq:lowpass} não precisam de fios separados: basta-lhes uma única grandeza comum. O filtro decorre diretamente da liberação em volume medida nas regiões de destino~\cite{Bunin1998}; o regulador de nível é uma hipótese.

\section{O horizonte de planejamento}
\label{sec:horizon}

Para que pode servir uma única grandeza comum e lenta? Um bom candidato é um parâmetro que deve ser igual para toda a rede e não mudar mais depressa que em segundos ou minutos. No capítulo~\ref{sec:neurontypes}, esse parâmetro já apareceu: o coeficiente $\gamma$ da fórmula do erro~\eqref{eq:rpe}, que diz quanto uma recompensa futura vale menos que uma imediata. Em tempo contínuo, seu lugar é ocupado pelo \emph{horizonte} $\tau_\gamma$: uma recompensa $R$ que exige esperar um tempo $D$ vale agora $R\,e^{-D/\tau_\gamma}$.

O horizonte decide se vale a pena esperar. Suponha que se pode pegar já uma recompensa pequena $r_0$ ou esperar uma grande $R$. Esperar compensa enquanto
\begin{empheq}[box=\keybox]{equation}
D \;<\; \tau_\gamma\,\ln\frac{R}{r_0} .
\label{eq:patience}
\end{empheq}
Com $\tau_\gamma=2\,\text{s}$ e uma recompensa adiada três vezes maior, vale a pena esperar até $2.2\,\text{s}$; se o horizonte for duas vezes mais longo, até $4.4\,\text{s}$. A paciência é proporcional ao horizonte.

A fórmula~\eqref{eq:patience} supõe um desconto exponencial. O desconto medido com atrasos de segundos fica mais perto da hipérbole $R/(1+D/\tau_\gamma)$: é assim que, em macacos, caem tanto o valor da recompensa adiada na escolha quanto a resposta dos neurônios \nt{DOPA} ao seu sinal antecipador~\cite{KobayashiSchultz2008}. Para a hipérbole, a condição~\eqref{eq:patience} passa a ser $D<\tau_\gamma\,(R/r_0-1)$: a dependência da razão entre as recompensas deixa de ser logarítmica e passa a ser linear, e, com os mesmos números, o limite de espera se desloca quase para o dobro. A hipérbole sai da mesma exponencial se o atraso for aleatório (exercício~\ref{ex:05:7}), e a proporcionalidade entre paciência e escala de tempo se mantém.

Pela hipótese, é justamente o \nt{SERO} que define o horizonte~\cite{Doya2002}. Ele tem tudo para esse papel: um único nível comum para toda a rede, que muda em segundos. Há também experimentos diretos: se os neurônios serotoninérgicos são ativados artificialmente, o animal espera mais tempo por uma recompensa adiada antes de desistir~\cite{Miyazaki2014} e insiste por mais tempo em buscar recompensa onde ela ficou rara~\cite{Lottem2018}. Como exatamente a concentração de serotonina se converte em $\tau_\gamma$ dentro da rede não se sabe. No próximo capítulo, o horizonte entra no erro que o \nt{DOPA} difunde.

\section{Resumo}

\begin{table}[t]
\centering
\small
\begin{tabular}{>{\raggedright}p{3.6cm}>{\raggedright}p{4.3cm}>{\raggedright\arraybackslash}p{4.3cm}}
\toprule
& \nt{GLUT} & \nt{SERO} \\
\midrule
tempo de reação & milissegundos & frações de segundo a um segundo \\
sinal da ação & só $+$ & $+$ e $-$, escolhido pelo destinatário \\
sem entrada & fica em silêncio & funciona sozinho com aporte de fundo constante \\
endereçamento & ponto a ponto & volume; o endereço é o lugar \\
NÃO & só com sensores de ``desligamento'' & um neurônio \\
memória & atividade autossustentada, apaga-se por fadiga & concentração, desfaz-se no tempo $\tau_c$ \\
cálculos principais & coincidências, atrasos, lógica, sequências & suavização; regulação de nível e horizonte $\tau_\gamma$ (hipóteses) \\
\bottomrule
\end{tabular}
\caption{O que calculam as redes feitas só de neurônios \nt{GLUT} e só de neurônios \nt{SERO}.}
\label{tab:glutvssero}
\end{table}

Como elemento lógico (tabela~\ref{tab:glutvssero}), o neurônio \nt{SERO} é mais rico que o \nt{GLUT}, mas como sistema de comunicação ele é uma difusão em massa: lenta e quase sem endereços. Por isso ele não tenta ser um processador: suaviza e difunde algumas grandezas comuns, das quais falamos no capítulo~\ref{sec:neurontypes}, entre elas, pela hipótese, o horizonte de planejamento. O marca-passo, o volume e a concentração lenta vão nos aparecer mais três vezes: no \nt{DOPA}, no \nt{NORA} e no \nt{ACET}.

\section{Exercícios}

\begin{exercise}[\lvlA]
\label{ex:05:1}
\emph{O fio é um volume.} Adote para a serotonina $D=3\cdot10^{-6}$~cm$^2$/s (seção~\ref{sec:volume}). Encontre por~\eqref{eq:diffusionlength} $\ell$ para um sinal que dura $1\,\text{s}$ e o tempo para o neurotransmissor se espalhar por $5$~cm. O que garante então a difusão ampla do \nt{SERO}: a difusão molecular ou a ramificação do fio com $\sim10^5$ locais de liberação?
\end{exercise}

\begin{exercise}[\lvlA]
\label{ex:05:2}
\emph{Regulador de nível.} Mostre que em~\eqref{eq:seroneuron}, com $f(u)=au$, a sensibilidade relativa $S=(a/r)\,\partial r/\partial a$ é exatamente $1/(1+G)$, e não só para $G\gg1$. Com $G=9$, a excitabilidade $a$ subiu $20\,\%$. Em quantos por cento muda $r$ sem linearização?
\end{exercise}

\begin{exercise}[\lvlB]
\label{ex:05:3}
\emph{Receptor lento no laço.} O receptor \nt{SERO} responde com atraso: $r(t)=a\bigl(b-gc(t-T)\bigr)$, e o volume segue~\eqref{eq:seroneuron}. Mostre que, para $G>1$, o laço só é estável se $T<T^*=\tau_c\arccos(-1/G)/\sqrt{G^2-1}$. Encontre $T^*$ com $\tau_c=1\,\text{s}$ para $G=2$ e $G=9$. Que redução da dispersão é possível com um atraso de centenas de milissegundos?
\end{exercise}

\begin{exercise}[\lvlB]
\label{ex:05:4}
\emph{Ruído balístico do nível.} Cada disparo soma a $c$, por~\eqref{eq:lowpass}, uma exponencial com $\tau_c=1\,\text{s}$. Encontre o coeficiente de variação de $c$ se todos os $3\cdot10^5$ neurônios \nt{SERO} liberam num mesmo volume disparos de Poisson a $2$ por segundo. Que $\tau_c$ daria $\mathrm{CV}=10\,\%$ a um único neurônio com taxa de $4$ disparos por segundo?
\end{exercise}

\begin{exercise}[\lvlB]
\label{ex:05:5}
\emph{Simulação: realimentação contra ruído.} Simule $N=50$ marca-passos de Poisson com taxas $r_i=a_i[b-gc]_+$, $a_i$ uniformes em $[2.8,\,5.2]$, e o volume~\eqref{eq:seroneuron} com $\kappa=1/N$, $\tau_c=1\,\text{s}$. Quantas vezes a realimentação ($b=1$, $g=2.25$) reduz o CV do nível $c$ em comparação com $b=0.1$, $g=0$? Ela iguala as taxas dos neurônios?
\end{exercise}

\begin{exercise}[\lvlB]
\label{ex:05:6}
\emph{Projeto: XOR com lógica \nt{SERO}.} A porta é um marca-passo que emite $r_0$ se $b-g\sum_kc_k>0$, e $0$ caso contrário, no seu próprio volume $\tau_c\,\dd c/\dd t=-c+\kappa r$; ela calcula NOU. Monte um XOR com cinco dessas portas e encontre seu tempo de acomodação com $\tau_c=1\,\text{s}$ e $G'=g\kappa r_0/b=2$.
\end{exercise}

\begin{exercise}[\lvlB]
\label{ex:05:7}
\emph{Paciência com atraso desconhecido.} (a)~Um animal aceita esperar por uma recompensa três vezes maior se a espera não passar de $6\,\text{s}$. Que horizonte $\tau_\gamma$ isso indica? (b)~Seja o atraso $D$ distribuído exponencialmente com média $\bar D$. Mostre que esperar compensa se $\bar D<\tau_\gamma(R/r_0-1)$ e compare com um atraso fixo para $\tau_\gamma=2\,\text{s}$, $R=3r_0$.
\end{exercise}

\begin{exercise}[\lvlC]
\label{ex:05:8}
\emph{Como a concentração define o horizonte.} Um neurônio \nt{DOPA} recebe $V$ diretamente com peso $k_1$ e por um intermediário \nt{GABA}, que suaviza com $\tau_d=0.1\,\text{s}$, com peso $-k_2$; para $V$ lento, a soma é $(k_1-k_2)V+k_2\tau_d\dot V$. Escolha $k_1$, $k_2$ para~\eqref{eq:ctd} com $\tau_\gamma=2\,\text{s}$. O \nt{SERO} multiplica $k_1$ por $m$: quanto é preciso mudar $m$ para dobrar o horizonte?
\end{exercise}

\chapter{Redes neurais que sabem aprender: acrescentamos o \nt{DOPA}}
\chaptermark{Redes que aprendem}
\label{sec:dofa}

\section{As regras do jogo}

Em todos os circuitos dos capítulos anteriores, os pesos foram ajustados pelo engenheiro. No organismo não há engenheiro. Os pesos têm de se ajustar sozinhos, durante o funcionamento, e de modo que a rede faça algo útil. É isso que se chama aprendizado.

Às regras anteriores soma-se mais uma. A sinapse muda o próprio peso sozinha e só pode saber o que acontece \emph{perto dela}: a atividade do remetente $x$, a atividade do destinatário $y$ e as substâncias que chegam até ela. Ninguém pode mandar à sinapse uma correção personalizada.

Isso exclui o principal método das redes neurais artificiais, a retropropagação do erro. Nela, o erro da saída percorre a rede no sentido inverso pelos mesmos pesos, numa fase separada depois da passagem direta, e cada peso recebe sua própria correção. Para isso são necessários fios de retorno que repitam exatamente os diretos, e um relógio comum. Nenhuma das duas coisas foi encontrada no cérebro, pelo menos na forma que têm nas redes artificiais~\cite{Lillicrap2020,WhittingtonBogacz2019}.

Vamos escrever a mudança de peso como um processo lento e contínuo:
\begin{equation}
\frac{\dd w}{\dd t} \;=\; \eta\,\Phi(\text{o que a sinapse sabe}),
\label{eq:plasticity}
\end{equation}
em que $\eta$ é a taxa de aprendizado, tão pequena que durante um disparo o peso quase não muda. O capítulo inteiro trata de qual pode ser a função $\Phi$.

\section{Onde fica o peso}
\label{sec:weightstore}

O que é uma sinapse que ``muda o próprio peso''? A conexão tem dois lados: o terminal do fio do remetente (lado \emph{pré-sináptico}) e o trecho da membrana do destinatário com receptores (lado \emph{pós-sináptico}), e entre eles uma fenda estreita. Nas conexões \nt{GLUT}, a membrana do destinatário costuma ficar numa \emph{espinha}, uma pequena saliência do ramo de entrada do destinatário. É conveniente decompor o peso da conexão em três fatores~\cite{delCastillo1954}:
\begin{empheq}[box=\keybox]{equation}
w \;=\; N_s\,p\,A_1 .
\label{eq:quantal}
\end{empheq}
Aqui $N_s$ é o número de locais de liberação entre dois neurônios, $p$ é a probabilidade de um disparo liberar uma porção de neurotransmissor num local (a mesma $p$ do capítulo~\ref{sec:code}), e $A_1$ é a resposta do destinatário a uma porção, isto é, quantos receptores a capturam. O fator $p$ mora no remetente, $A_1$ no destinatário, e $N_s$ exige os dois lados: novos locais de liberação crescem, antigos desaparecem, e isso continua a vida toda.

\emph{Quem decide é o destinatário.} Numa sinapse \nt{GLUT} típica, um dos receptores de glutamato só conduz corrente se duas condições coincidirem: o glutamato chegou do remetente e a membrana do destinatário já está excitada~\cite{Nowak1984}. É a regra de Hebb embutida numa molécula: o detector de coincidência do capítulo~\ref{sec:glutcomputer}, instalado na entrada da própria sinapse. Ao ser ativado, o receptor deixa entrar cálcio no destinatário, e o cálcio desencadeia a mudança de peso. O destinatário é o único lugar onde se veem ao mesmo tempo a entrada e a saída, por isso a decisão é tomada nele.

\emph{Qualquer lado pode executar a decisão.} A potenciação de longo prazo mais estudada ocorre no destinatário: novos receptores são inseridos na membrana~\cite{Malinow2002} e a espinha cresce~\cite{Matsuzaki2004}; cresce $A_1$. Mas o destinatário também pode mudar $p$: ele libera uma substância que atravessa a fenda de volta até o remetente (o canal de retorno do apêndice~\ref{sec:othermediators}), e este passa a liberar menos neurotransmissor~\cite{WilsonNicoll2001}. Já as mudanças de peso de curto prazo, a depressão e a facilitação do capítulo~\ref{sec:code}, são mudanças de $p$ que o próprio remetente conduz segundo sua atividade recente.

Para o engenheiro, o registrador do peso está dividido entre dois chips. A regra de aprendizado é executada pelo destinatário, mas ele pode gravar o resultado tanto em si mesmo quanto no remetente, pelo canal de retorno. Por isso, a palavra ``local'' nas regras deste capítulo não significa um lado da conexão, mas a conexão inteira.

\section{Aprender sem professor}

O mais simples que uma sinapse pode fazer é se reforçar quando remetente e destinatário estão ativos juntos:
\begin{equation}
\frac{\dd w}{\dd t} \;=\; \eta\,x\,y .
\label{eq:hebb}
\end{equation}
É a regra de Hebb~\cite{Hebb1949}. Ela é local e não exige relógio. Mas tem o mesmo problema da rede \nt{GLUT} do capítulo~\ref{sec:glutcomputer}: realimentação positiva. Um peso forte aumenta $y$, um $y$ grande reforça ainda mais o peso, e os pesos crescem sem limite. O remédio é uma realimentação negativa sobre o peso: que ele também diminua em proporção a $y^2$. Para um neurônio com entradas $\mathbf x$ e saída linear $y=\mathbf w\cdot\mathbf x$, obtemos
\begin{empheq}[box=\keybox]{equation}
\frac{\dd\mathbf w}{\dd t} \;=\; \eta\,y\,\bigl(\mathbf x - y\,\mathbf w\bigr) .
\label{eq:oja}
\end{empheq}
A regra continua local: a sinapse conhece sua entrada $x_i$, a saída $y$ e o peso $w_i$.

\begin{keyblock}
\begin{proposition}[O que a regra de Hebb encontra]
\label{prop:oja}
A regra~\eqref{eq:oja} leva o vetor de pesos a ter comprimento unitário ao longo da direção em que as entradas variam mais, ou seja, ao autovetor principal da matriz de correlação das entradas $C=\langle\mathbf x\mathbf x^{\mathsf T}\rangle$~\cite{Oja1982}. O neurônio aprende a emitir a grandeza única mais expressiva na qual suas entradas podem ser comprimidas.
\end{proposition}
\end{keyblock}

Demonstração. Tomemos a média de~\eqref{eq:oja} sobre as entradas: $\dd\mathbf w/\dd t=\eta\bigl(C\mathbf w-(\mathbf w^{\mathsf T}C\mathbf w)\,\mathbf w\bigr)$. Os pontos fixos são os autovetores de $C$ de comprimento unitário. Se $\mathbf w$ está perto do autovetor $\mathbf e_k$ com autovalor $\lambda_k$, uma pequena perturbação ao longo de outro vetor $\mathbf e_j$ cresce como $e^{\eta(\lambda_j-\lambda_k)t}$. Só é estável o ponto para o qual todos os $\lambda_j<\lambda_k$, isto é, o vetor principal.

Há também uma variante da regra de Hebb que olha para o tempo dos disparos. Se o disparo do remetente chegou $s$ milissegundos \emph{antes} do disparo do destinatário, o peso cresce $A_+e^{-s/\tau_+}$; se chegou $s$ milissegundos \emph{depois}, diminui $A_-e^{-s/\tau_-}$, com $\tau_\pm$ da ordem de vinte milissegundos. Essa regra foi medida em sinapses de neurônios \nt{GLUT}~\cite{BiPoo1998}. Ela reforça as conexões que \emph{podem ter sido a causa} da resposta e enfraquece as atrasadas (fig.~\ref{fig:stdp}). Passe muitas vezes pela rede uma mesma sequência de excitações, e nela crescerão sozinhas conexões de cada elo para o seguinte: a cadeia do capítulo~\ref{sec:glutcomputer}, montada sem engenheiro.

\begin{figure}[t]
\centering
\begin{tikzpicture}[>=Stealth,x=0.07cm,y=1.3cm]
\draw[->,gray!70!black] (-66,0) -- (68,0) node[right,black] {\small $s$};
\draw[->,gray!70!black] (0,-1.2) -- (0,1.35) node[above,black] {\small mudança de peso};
\draw[very thick,blue!60!black] plot[domain=0.5:60,samples=50] (\x,{exp(-\x/20)});
\draw[very thick,red!70!black] plot[domain=-60:-0.5,samples=50] (\x,{-0.6*exp(\x/20)});
\draw[blue!60!black,thick] (0.5,0) -- (0.5,1);
\draw[red!70!black,thick] (-0.5,0) -- (-0.5,-0.6);
\foreach \x in {-40,-20,20,40}{\draw[gray!60!black] (\x,-0.04) -- (\x,0.04);}
\node[below,font=\scriptsize] at (20,-0.04) {$20\ms$};
\node[above,font=\scriptsize] at (-20,0.04) {$-20\ms$};
\node[align=left,font=\small,blue!60!black,anchor=west] at (22,0.95) {remetente antes:\\ conexão se reforça};
\node[align=right,font=\small,red!70!black,anchor=east] at (-12,-0.95) {remetente depois:\\ conexão se enfraquece};
\end{tikzpicture}
\caption{Regra de Hebb temporal. No eixo horizontal, $s=t_{\text{dest}}-t_{\text{rem}}$, o quanto o disparo do destinatário ficou atrás do disparo do remetente; $\tau_\pm=20\ms$, $A_-=0.6A_+$. A janela, com largura de algumas dezenas de milissegundos, é da mesma ordem que a janela de coincidência~\eqref{eq:window}.}
\label{fig:stdp}
\end{figure}

Mas todas as regras desta seção ensinam o que é \emph{frequente}, e não o que é \emph{útil}. Uma sinapse que só conhece $x$ e $y$ não sabe se a ação trouxe suco ou dor. Para aprender com o resultado, a sinapse precisa de um terceiro sinal.

\section{O terceiro fator}

O terceiro sinal é trazido pelo \nt{DOPA}: os neurônios dopaminérgicos difundem um único número, o erro de predição de recompensa $\delta$~\eqref{eq:rpe} (capítulo~\ref{sec:neurontypes}). Seja a mudança de peso proporcional ao produto de três fatores: a atividade do remetente, a atividade do destinatário e $\delta$.

A dificuldade é que a recompensa não chega quando a rede escolhia a ação, mas centenas de milissegundos ou segundos depois. Quando $\delta$ chega, o produto $xy$ já é zero há muito tempo, e a sinapse precisa de uma memória de que participou. A memória mais simples é o nosso conhecido circuito RC:
\begin{empheq}[box=\keybox]{equation}
\tau_e\,\frac{\dd e}{\dd t} \;=\; -e \;+\; x\,y ,
\qquad
\frac{\dd w}{\dd t} \;=\; \eta\,\delta(t)\,e(t) .
\label{eq:threefactor}
\end{empheq}
A grandeza $e$ se chama \emph{traço de elegibilidade}~\cite{Fremaux2016}. A atividade conjunta deixa na sinapse uma marca que se desfaz no tempo $\tau_e$. Se nesse tempo chega dopamina, o peso muda com o sinal de $\delta$; se não chega, a marca some sem deixar rastro (fig.~\ref{fig:trace}). Em sinapses \nt{GLUT} foi medido: a dopamina que chega um ou dois segundos depois da atividade conjunta a converte em reforço da conexão, e a que chega mais tarde, não~\cite{Yagishita2014}. Janelas de plasticidade de segundos também foram encontradas onde o terceiro sinal não é a dopamina, mas uma forte excitação do próprio destinatário~\cite{Bittner2017}.

\begin{figure}[t]
\centering
\begin{tikzpicture}[>=Stealth,x=7cm,y=1cm]
\fill[orange!12] (0.7,-0.2) rectangle (0.8,5.2);
% rows: x*y, delta, traces, weights
\foreach \y/\lab in {4.4/{$x\,y$},3.1/{$\delta$},1.4/{traço $e$},0/{peso $w$}}{
  \draw[gray!70!black] (0,\y) -- (1.42,\y);
  \node[left,font=\small] at (-0.02,\y+0.3) {\lab};}
\draw[very thick,blue!60!black] (0,4.4) -- (0,4.9) -- (0.2,4.9) -- (0.2,4.4);
\node[right,font=\scriptsize,blue!60!black] at (0.21,4.8) {remetente e destinatário ativos juntos};
\draw[very thick,orange!85!black] (0,3.1) -- (0.7,3.1) -- (0.7,3.7) -- (0.8,3.7) -- (0.8,3.1) -- (1.4,3.1);
\node[right,font=\scriptsize,orange!85!black] at (0.81,3.6) {chegou dopamina};
% traces, each in units of its own maximum
\draw[very thick,blue!60!black] plot[domain=0:0.2,samples=20] (\x,{1.4+1.1*(1-exp(-\x))/(1-exp(-0.2))})
  plot[domain=0.2:1.4,samples=60] (\x,{1.4+1.1*exp(-(\x-0.2))});
\draw[thick,densely dashed,gray!50!black] plot[domain=0:0.2,samples=30] (\x,{1.4+1.1*(1-exp(-\x/0.05))/(1-exp(-4))})
  plot[domain=0.2:1.4,samples=80] (\x,{1.4+1.1*exp(-(\x-0.2)/0.05)});
\node[right,font=\scriptsize,blue!60!black] at (1.0,2.15) {$\tau_e=1\,\text{s}$};
\node[right,font=\scriptsize,gray!40!black] at (0.3,1.65) {$\tau_e=50\ms$};
% weights
\draw[very thick,blue!60!black] (0,0.25) -- (0.7,0.25) -- (0.8,0.8) -- (1.4,0.8) node[right,font=\scriptsize] {$\tau_e=1\,\text{s}$: peso cresceu};
\draw[thick,densely dashed,gray!50!black] (0,0.2) -- (1.4,0.2) node[right,font=\scriptsize,gray!40!black] {$\tau_e=50\ms$: sem mudança};
% time axis marks
\foreach \x/\l in {0/0,0.2/0.2,0.7/0.7,1.4/1.4}{\draw[gray!60!black] (\x,-0.05) -- (\x,0.05) node[below=3pt,font=\scriptsize] {\l\,s};}
\draw[<->,thick,gray!50!black] (0.2,5.35) -- node[above,font=\scriptsize,black] {atraso da recompensa $0.5\,\text{s}$} (0.7,5.35);
\end{tikzpicture}
\caption{Traço de elegibilidade pela regra~\eqref{eq:threefactor}. A atividade conjunta dura $0.2\,\text{s}$, e a dopamina chega meio segundo após seu fim. Os traços são mostrados em frações do próprio máximo. O traço lento ($\tau_e=1\,\text{s}$) ainda está vivo nesse momento; o rápido ($\tau_e=50\ms$) já se desfez há muito.}
\label{fig:trace}
\end{figure}

\begin{keyblock}
\begin{proposition}[Aprendizado por sinal de difusão]
\label{prop:threefactor}
A regra~\eqref{eq:threefactor} muda só as sinapses que participaram do funcionamento da rede nos últimos $\tau_e$, no sentido dado pelo sinal comum $\delta$. O sinal de difusão, junto com o traço local, produz uma correção endereçada. Um atraso entre ação e resultado é admissível enquanto não passar de $\tau_e$.
\end{proposition}
\end{keyblock}

\section{Por que funciona: o ruído como busca}
\label{sec:dofagradient}

Por que a regra~\eqref{eq:threefactor} leva a uma melhora? A resposta vem do ruído do capítulo~\ref{sec:code}. Seja a saída do neurônio $y=f(u)+\xi$, com $u=\sum_jw_jx_j$, e $\xi$ uma flutuação aleatória de média zero e variância $\sigma^2$. A recompensa $R$ depende de $y$. Tomemos como traço não simplesmente $x_jy$, mas sua parte aleatória $x_j\xi$, e como terceiro fator o erro $\delta=R-\bar R$, em que $\bar R$ é a recompensa esperada. Com ruído pequeno, $R\approx R\bigl(f(u)\bigr)+R'\,\xi$, portanto
\begin{empheq}[box=\keybox]{equation}
\bigl\langle \delta\,\xi\,x_j \bigr\rangle \;=\; \sigma^2\,R'\,x_j \;=\; \frac{\sigma^2}{f'(u)}\,\frac{\partial\langle R\rangle}{\partial w_j} .
\label{eq:perturbation}
\end{empheq}
O passo médio segue o gradiente da recompensa esperada~\cite{Williams1992}. Ninguém calculou derivadas: a rede se desviou ao acaso, a dopamina informou se melhorou, e as sinapses que participaram se deslocaram no sentido vantajoso (fig.~\ref{fig:perturbation}). O ruído, que no capítulo~\ref{sec:code} atrapalhava a transmissão de mensagens, aqui se torna busca.

\begin{figure}[t]
\centering
\begin{tikzpicture}[>=Stealth,x=2cm,y=3cm]
\draw[->,gray!70!black] (0,0) -- (4.2,0) node[below left,font=\small,black] {saída $y$};
\draw[->,gray!70!black] (0,0) -- (0,1.2) node[left,font=\small,black] {recompensa $R$};
% reward hill
\draw[very thick,blue!60!black] plot[domain=0:4,samples=60] (\x,{max(0,1-((\x-3)/2.2)^2)});
\node[font=\scriptsize,blue!60!black,anchor=south] at (3,1.02) {melhor saída};
% noise around f(u)
\draw[thick,gray!60!black] plot[domain=0.6:2.4,samples=50] (\x,{0.28*exp(-((\x-1.5)/0.3)^2/2)});
\draw[densely dashed,gray!60!black] (1.5,0) -- (1.5,0.535);
\node[below,font=\small] at (1.5,0) {$f(u)$};
\node[font=\scriptsize,gray!50!black] at (1.5,-0.2) {flutuação $\xi$};
% expected reward
\draw[densely dotted,gray!60!black] (0.5,0.535) node[left,font=\scriptsize,black] {$\bar R$} -- (2.1,0.535);
% two random deviations
\fill[green!45!black] (2.1,0.833) circle (0.035);
\draw[very thick,green!45!black] (2.1,0.535) -- (2.1,0.833);
\node[font=\scriptsize,green!45!black,anchor=west,align=left] at (2.18,0.6) {$\xi>0$: melhor,\\ $\delta>0$};
\fill[red!70!black] (0.9,0.089) circle (0.035);
\draw[very thick,red!70!black] (0.9,0.535) -- (0.9,0.089);
\node[font=\scriptsize,red!70!black,anchor=east,align=right] at (0.83,0.27) {$\xi<0$: pior,\\ $\delta<0$};
% resulting step
\draw[->,very thick,orange!85!black] (1.55,0.4) -- (1.95,0.4) node[right,font=\scriptsize,black] {passo};
\end{tikzpicture}
\caption{O ruído como busca~\eqref{eq:perturbation}. A saída do neurônio flutua em torno de $f(u)$. Um desvio aleatório para a direita (verde) deu mais que o esperado, $\delta>0$; para a esquerda (vermelho), menos, $\delta<0$. Nos dois casos o produto $\delta\,\xi$ é positivo, e o peso se desloca para a direita, onde a recompensa cresce. Em média, o passo é proporcional à inclinação $R'$: no topo da colina, desvios para os dois lados são igualmente ruins, e os passos se anulam.}
\label{fig:perturbation}
\end{figure}

\begin{keyblock}
\begin{proposition}[Descida de gradiente sem fios de retorno]
\label{prop:gradient}
Se a rede tem desvios aleatórios de atividade e o sinal de difusão é igual ao erro de predição de recompensa, com média zero em cada situação, a regra~\eqref{eq:threefactor} muda, em média, os pesos segundo o gradiente da recompensa esperada. Para isso não são necessários fios de retorno nem uma fase separada de aprendizado.
\end{proposition}
\end{keyblock}

Agora fica claro por que a dopamina não informa a recompensa, mas o \emph{erro} de sua predição. Se o terceiro fator fosse a própria recompensa $R\ge0$, a média em~\eqref{eq:perturbation} não mudaria, mas surgiriam dois problemas. Primeiro, à parte útil se somaria o termo $\bar R\,\xi x_j$: zero em média, mas um ruído forte. Segundo, e pior, uma sinapse real não sabe separar no traço $x_jy$ a parte aleatória da não aleatória. Com as entradas dadas, $x_jf(u)$ é o mesmo número de uma tentativa para outra; se $\delta$ tem média zero nessa situação, essa parte do traço não muda nada, mas com $\delta\ge0$ ela reforça, vez após vez, tudo o que a rede fez, o certo e o errado. Por isso $\bar R$ deve ser a expectativa \emph{na situação dada}, e não a média de toda a vida. Calculá-la é tarefa do crítico, de que falamos abaixo.

Verificamos isso num modelo. Dois grupos de neurônios \nt{GLUT} escolhem uma de duas ações por inibição mútua, como na fig.~\ref{fig:wta}; os pesos do sinal ``começar'' para os grupos são inicialmente iguais, e o ruído decide quem vence. A ação $A$ traz suco com probabilidade $0.8$, a ação $B$ com probabilidade $0.2$, e o suco chega meio segundo depois da escolha. Com $\tau_e=1\,\text{s}$ e $\delta=R-\bar R$, a fração de escolhas de $A$ em quinhentas tentativas subiu de $0.65$--$0.8$ na primeira centena para $0.96$--$1.0$ na última, e os pesos ficaram entre $0.85$ e $1.35$. Com $\tau_e=50\ms$, menor que o atraso da recompensa, a rede não aprendeu nada: a fração de escolhas de $A$ ficou em torno da metade. Com $\delta=R$, sem subtrair a expectativa, a rede também passou a escolher $A$, mas o peso do vencedor cresceu mil vezes nas mesmas quinhentas tentativas e continuou crescendo; e com o mesmo $\delta=R$ e uma taxa de aprendizado dez vezes maior, a rede, numa de três execuções, fixou para sempre sua primeira escolha aleatória, a pior.

\section{O preço de um único fio}

O sinal $\delta$ é um só para todos, e o ruído que ele avalia é a soma das flutuações de todos os $N$ neurônios que influem no resultado. Cada sinapse vê em $\delta$ sua parte útil e o ruído alheio de $N-1$ vizinhos. Para separar o que é seu, é preciso fazer a média sobre muitas tentativas, e o tempo de aprendizado cresce proporcionalmente a $N$~\cite{Werfel2005}. A retropropagação não cobra esse preço.

\begin{keyblock}
\begin{proposition}[O preço da difusão]
\label{prop:broadcastcost}
O tempo de aprendizado por um único sinal comum cresce proporcionalmente ao número de neurônios cujo ruído influi no resultado. Esse aprendizado serve para circuitos pequenos, que escolhem entre poucas opções, mas não para ajustar bilhões de pesos.
\end{proposition}
\end{keyblock}

O cérebro, ao que parece, divide o trabalho assim: as saídas dos neurônios \nt{DOPA} vão sobretudo às regiões que escolhem ações (capítulo~\ref{sec:neurontypes}), e a enorme rede \nt{GLUT} do córtex, onde está a esmagadora maioria dos pesos, aprende principalmente por regras locais, sem professor externo. Antes, elas eram reduzidas às regras de Hebb~\eqref{eq:oja}. Hoje há cada vez mais dados de que o córtex tem um objetivo próprio, \emph{prever} sua próxima entrada, e então o erro é calculado no próprio local, como a diferença entre o previsto e o que chegou~\cite{Keller2018}: o professor está embutido na própria rede. Janelas de plasticidade de segundos, em que o terceiro sinal é uma forte excitação do próprio destinatário~\cite{Bittner2017}, mostram que um sinal de erro local nas sinapses de fato existe. Até que ponto essa é uma regra geral do córtex é uma hipótese.

\section{De onde vem o erro: o crítico em tempo contínuo}

Resta a pergunta de como os próprios neurônios \nt{DOPA} calculam $\delta$. Nos programas de aprendizado por reforço, o erro de predição é calculado em passos $t,t+1,\dots$ No organismo não há passos, por isso vamos escrevê-lo em tempo contínuo~\cite{Doya2000}. Seja $r(t)$ o fluxo de recompensa e $V(t)$ a expectativa de toda a recompensa futura, em que a mais distante vale menos:
\begin{equation}
V(t) \;=\; \int_t^{\infty} e^{-(s-t)/\tau_\gamma}\,r(s)\,\dd s .
\end{equation}
Derivando, obtemos $\dd V/\dd t=V/\tau_\gamma-r$. O resíduo dessa igualdade é justamente o erro de predição:
\begin{empheq}[box=\keybox]{equation}
\delta(t) \;=\; r(t) \;+\; \frac{\dd V}{\dd t} \;-\; \frac{V}{\tau_\gamma} .
\label{eq:ctd}
\end{empheq}
A constante $\tau_\gamma$ é o horizonte de planejamento, análogo contínuo do coeficiente $\gamma$; pela hipótese, é definida pelo \nt{SERO} (seção~\ref{sec:horizon}), e então~\eqref{eq:patience} é a regra de quanto vale a pena esperar. Vamos verificar três casos (fig.~\ref{fig:ctdcases}). Acendeu-se uma luz que promete suco: $V$ subiu de repente, $\dd V/\dd t$ é grande, pico. Chegou o suco prometido: $r>0$, mas a expectativa se esgotou no mesmo instante, $\dd V/\dd t\approx-r$, e não há pico. O suco não veio: a expectativa sumiu sem recompensa, $\dd V/\dd t<0$, queda.

\begin{figure}[t]
\centering
\begin{tikzpicture}[>=Stealth,x=1.15cm,y=1cm]
\foreach \col/\ttl/\lampon/\juice in {0/{suco inesperado}/0/1, 1/{suco prometido}/1/1, 2/{o suco não veio}/1/0}{
 \begin{scope}[xshift=\col*4.6cm]
  \node[font=\small] at (1.5,3.55) {\ttl};
  \foreach \yy in {2.4,1.2,0}{\draw[gray!60!black] (0,\yy) -- (3.1,\yy);}
  % reward r
  \ifnum\juice=1
    \fill[green!45!black] (1.95,2.4) rectangle (2.05,2.85);
    \pic[scale=0.55] at (2.0,3.1) {juice};
  \else
    \pic[scale=0.55] at (2.0,3.1) {nojuice};
  \fi
  % expectation V and error delta
  \ifnum\lampon=1
    \pic[scale=0.55] at (1.0,3.1) {lamp};
    \draw[very thick,blue!60!black] (0,1.2) -- (1,1.2) -- (1,1.45)
      plot[domain=1:2,samples=20] (\x,{1.2+0.25*exp((\x-1)/1.5)}) -- (2,{1.2+0.25*exp(1/1.5)}) -- (2,1.2) -- (3.1,1.2);
    \draw[very thick,red!70!black] (0,0) -- (0.95,0) -- (0.95,0.5) -- (1.05,0.5) -- (1.05,0) -- (3.1,0);
    \ifnum\juice=0
      \draw[very thick,red!70!black] (1.95,0) -- (1.95,-0.45) -- (2.05,-0.45) -- (2.05,0);
      \node[font=\scriptsize,red!70!black,anchor=west] at (2.1,-0.35) {queda};
    \else
      \node[font=\scriptsize,gray!50!black,anchor=north,align=center] at (2.0,-0.08) {$r$ e a queda de $V$\\ se anulam};
    \fi
  \else
    \draw[very thick,blue!60!black] (0,1.2) -- (3.1,1.2);
    \draw[very thick,red!70!black] (0,0) -- (1.95,0) -- (1.95,0.5) -- (2.05,0.5) -- (2.05,0) -- (3.1,0);
  \fi
  \ifnum\col=0
    \node[font=\small,anchor=east] at (-0.1,2.55) {$r$};
    \node[font=\small,anchor=east] at (-0.1,1.35) {$V$};
    \node[font=\small,anchor=east] at (-0.1,0.1) {$\delta$};
  \fi
 \end{scope}}
\end{tikzpicture}
\caption{Três casos para o erro~\eqref{eq:ctd}; de cima para baixo, recompensa $r$, expectativa $V$ e erro $\delta$. À esquerda não há expectativa, e o suco é inteiramente uma surpresa. No meio, a luz eleva $V$ de repente: é um salto de $\dd V/\dd t$, e o pico de $\delta$ cai na luz. Depois, $V$ cresce exatamente como o horizonte exige, e $\delta=0$; no instante do suco, a recompensa $r$ e a queda de $V$ se anulam. À direita, $V$ cai sem recompensa: queda. São o mesmo pico, a mesma transferência e a mesma queda da fig.~\ref{fig:rpe}, mas agora se vê de onde vêm.}
\label{fig:ctdcases}
\end{figure}

O que é preciso, dos circuitos que já temos, para~\eqref{eq:ctd}? Três coisas.

\emph{A própria estimativa $V$.} Ela é emitida por um grupo de neurônios \nt{GLUT}, o crítico, cujos pesos aprendem pela mesma regra~\eqref{eq:threefactor} com o mesmo $\delta$: se $\delta>0$, a expectativa estava subestimada, e as conexões ativas logo antes se reforçam.

\emph{A derivada $\dd V/\dd t$.} Ela é calculada por qualquer circuito que responda às variações, e não ao nível: excitação direta junto com inibição atrasada por um intermediário \nt{GABA}, como no detector de direção do capítulo~\ref{sec:gaba}, ou uma sinapse com depressão~\eqref{eq:depression} do capítulo~\ref{sec:code}.

\emph{O sinal num único fio.} O erro pode ser positivo ou negativo, mas a taxa de disparos só pode ser positiva. A solução é a mesma do neurônio \nt{SERO} do capítulo~\ref{sec:serocomputer}: o neurônio \nt{DOPA} é um marca-passo. Seus poucos disparos regulares por segundo significam $\delta=0$; uma rajada, $\delta>0$; uma pausa, $\delta<0$. Para o erro negativo sobra menos espaço: a taxa não pode cair abaixo de zero. Nos neurônios \nt{DOPA} reais, as quedas são de fato menores que os picos~\cite{Bayer2005}.

Que a dopamina se comporta como $\delta$ foi medido. Como exatamente o cérebro calcula $\dd V/\dd t$ é uma hipótese.

\section{Ator e crítico}

Juntemos tudo (fig.~\ref{fig:actorcritic}). Os neurônios de contexto excitam grupos de ação, que escolhem um vencedor via \nt{GABA} (proposição~\ref{prop:wta}), o \emph{ator}, e o crítico, que emite $V$~\cite{SuttonBarto2018}. O neurônio \nt{DOPA} recebe a recompensa $r$ e $V$, calcula~\eqref{eq:ctd} e difunde $\delta$ a todas as sinapses do ator e do crítico.

\begin{figure}[t]
\centering
\begin{tikzpicture}[>=Stealth,
  nrn/.style={circle,draw,very thick,fill=blue!6,minimum size=0.9cm,inner sep=0pt,font=\small},
  gab/.style={circle,draw,very thick,fill=red!10,minimum size=0.6cm,inner sep=0pt},
  dofa/.style={circle,draw,very thick,double,double distance=1.2pt,fill=orange!15,minimum size=1.35cm,inner sep=0pt,font=\scriptsize},
  inh/.style={-{Bar[width=7pt]},thick,red!70!black},
  bc/.style={->,thick,densely dashed,orange!85!black}]
\node[nrn] (S1) at (0,2.4) {$x_1$};
\node[nrn] (S2) at (0,1.2) {$x_2$};
\node[nrn] (S3) at (0,0) {$x_3$};
\node[left=0.15cm,align=right,font=\small] at (-0.5,1.2) {contexto};
\node[nrn] (A) at (4,2.6) {$A$};
\node[nrn] (B) at (4,1.0) {$B$};
\node[gab] (G) at (5.2,1.8) {};
\draw[->,thick] (A) -- (G); \draw[->,thick] (B) -- (G);
\draw[inh] (G) to[bend left=35] (A);
\draw[inh] (G) to[bend right=35] (B);
\node[above,font=\small] at (4.6,3.05) {ator: escolha da ação};
\node[nrn] (V) at (4,-1.1) {$V$};
\node[below,font=\small] at (4,-1.6) {crítico};
\foreach \s in {S1,S2,S3}{\foreach \t in {A,B,V}{\draw[->,gray!60!black] (\s) -- (\t);}}
\node[dofa] (D) at (8.0,-1.1) {\nt{DOPA}};
\draw[->,thick] (V) -- node[above,font=\scriptsize] {$V,\ \dd V/\dd t$} (D);
\draw[->,thick,green!45!black] (10.0,-1.1) node[right,black,font=\small] {recompensa $r$} -- (D);
\pic[scale=1.1] at (10.6,-0.5) {juice};
\draw[bc] (D.north) -- (8.0,4.0) -- node[above,font=\small,black] {$\delta$ a todas as sinapses do ator e do crítico} (2.2,4.0) -- (2.2,3.0);
\draw[bc] (D.south) -- (8.0,-2.4) -- (2.2,-2.4) -- (2.2,-0.6);
\end{tikzpicture}
\caption{Uma rede que aprende a escolher ações. Setas cinza: sinapses \nt{GLUT} com pesos variáveis; círculo vermelho: neurônio \nt{GABA}; círculo duplo: marca-passo \nt{DOPA}, que calcula $\delta$ por~\eqref{eq:ctd}; setas tracejadas: difusão de $\delta$.}
\label{fig:actorcritic}
\end{figure}

O sinal da ação do neurotransmissor é escolhido pelo destinatário (proposição~\ref{prop:receptor}), e na região de escolha de ações parte dos neurônios tem receptores de dopamina de uma família, e parte, de outra. Em experimentos com fatias, a dopamina junto com a atividade conjunta reforça as sinapses nos primeiros e as enfraquece nos segundos~\cite{Shen2008}; no modelo, isso significa que nos segundos o sinal de $\delta$ em~\eqref{eq:threefactor} está invertido. Os primeiros aprendem o que \emph{vale a pena fazer}; os segundos, o que \emph{não vale a pena fazer}.

\section{Resumo}

\begin{table}[t]
\centering
\footnotesize
\setlength{\tabcolsep}{4pt}
\renewcommand{\arraystretch}{1.15}
\begin{tabular}{>{\raggedright}p{2.6cm}>{\raggedright}p{2.7cm}>{\raggedright}p{3.1cm}>{\raggedright\arraybackslash}p{5.0cm}}
\toprule
Regra & O que a sinapse sabe & O que a rede aprende & O que se sabe \\
\midrule
Hebb por taxas~\eqref{eq:oja} & $x$, $y$, o próprio peso & comprimir as entradas: direções principais & reforço com atividade conjunta medido; o mecanismo de limitação dos pesos é objeto de pesquisa \\
Hebb temporal & ordem dos disparos de $x$ e $y$ dentro de $\sim20\ms$ & conexões causais, sequências & medido em sinapses \nt{GLUT} \\
três fatores~\eqref{eq:threefactor} & $x$, $y$, traço $e$, sinal comum $\delta$ & ações que trazem recompensa & efeito da dopamina nos segundos após a atividade medido; como principal mecanismo do aprendizado de escolha, hipótese bem fundamentada \\
crítico~\eqref{eq:ctd} & o mesmo & expectativa de recompensa $V$ & semelhança da dopamina com $\delta$ medida; o circuito que calcula $\dd V/\dd t$ é hipótese \\
retropropagação do erro & correção personalizada para cada peso & tudo, mas exige fios de retorno e relógio & não encontrada no cérebro nessa forma; procuram-se aproximações \\
\bottomrule
\end{tabular}
\caption{Regras de aprendizado numa rede de neurônios \nt{GLUT}, \nt{GABA} e \nt{DOPA}.}
\label{tab:learning}
\end{table}

As regras de aprendizado estão reunidas na tab.~\ref{tab:learning}. O \nt{GLUT} transporta os dados, o \nt{GABA} controla o fluxo, e o \nt{DOPA} decide quais mudanças na rede ficam.

\section{Exercícios}

\begin{exercise}[\lvlA]
\label{ex:06:1}
\emph{Quanto dura o traço.} A atividade conjunta $xy=1$ dura $T_a=0.2\,\text{s}$, a partir de $e=0$, e a dopamina chega $d=0.5\,\text{s}$ após seu fim. (a)~Quantas vezes o traço~\eqref{eq:threefactor} nesse instante é maior com $\tau_e=1\,\text{s}$ do que com $\tau_e=50\ms$? (b)~Com que $\tau_e$ ele é máximo?
\end{exercise}

\begin{exercise}[\lvlA]
\label{ex:06:2}
\emph{Deriva da regra de Hebb.} Remetente e destinatário emitem trens de Poisson independentes de $10$ disparos por segundo, e cada par de disparos muda o peso pela janela da fig.~\ref{fig:stdp} com $\tau_\pm=20\ms$, $A_-=0.6A_+$. Quantos $A_+$ o peso cresce em uma hora de atividade totalmente aleatória? Com que $\tau_-$ a deriva desaparece?
\end{exercise}

\begin{exercise}[\lvlB]
\label{ex:06:3}
\emph{Correlações, não covariâncias.} A regra da proposição~\ref{prop:oja} encontra o autovetor principal de $C=\langle\mathbf x\mathbf x^{\mathsf T}\rangle$. As taxas são positivas: $\mathbf x=\mathbf m+\mathbf n$, $\mathbf m=(\mu,0)$, covariância do ruído $\operatorname{diag}(s_1^2,s_2^2)$. Sob que condição o neurônio aprende $\mathbf e_1$? O que aprende com $\mu=1$, $s_1=0.1$, $s_2=0.5$?
\end{exercise}

\begin{exercise}[\lvlA]
\label{ex:06:4}
\emph{O horizonte no pico.} A luz acende num instante imprevisível, e um suco de tamanho $R$ chega após $L=1\,\text{s}$. Mostre que, para a expectativa aprendida por~\eqref{eq:ctd}, entre a luz e o suco $\delta\equiv0$, e encontre a área do pico na luz com $\tau_\gamma=2\,\text{s}$ e $\tau_\gamma=4\,\text{s}$.
\end{exercise}

\begin{exercise}[\lvlB]
\label{ex:06:5}
\emph{De onde vem o preço da difusão.} $N$ neurônios flutuam independentemente, $\xi_i\sim\mathcal N(0,\sigma^2)$, o erro é $\delta=a\sum_i\xi_i$, e a sinapse $j$ estima o gradiente como $\delta\,\xi_jx_j$. Encontre a relação sinal-ruído de uma tentativa. Um neurônio aprende em $100$ tentativas de $5\,\text{s}$: quanto leva o aprendizado com $N=10^4$ e $N=10^{10}$?
\end{exercise}

\begin{exercise}[\lvlB]
\label{ex:06:6}
\emph{Projeto: circuito que calcula $\delta$.} O neurônio \nt{DOPA} recebe $r$, $V$ com peso $k_1$ e uma cópia de $V$ suavizada com $\tau_d$, com peso $-k_2$. (a)~Escolha $k_1$, $k_2$ para que, com $V$ lento, a saída seja~\eqref{eq:ctd}. (b)~Com $V=V_0e^{t/\tau_\gamma}$, o erro verdadeiro é $\delta=0$. Com que $\tau_d$ o erro espúrio do circuito não passa de $5\,\%$ de $V/\tau_\gamma$, se $\tau_\gamma=2\,\text{s}$?
\end{exercise}

\begin{exercise}[\lvlB]
\label{ex:06:7}
\emph{Simulação: atraso máximo da recompensa.} Acrescente um atraso de recompensa $d$ em \texttt{run()} numa cópia de \texttt{models/\allowbreak dofa\_learning.py}. Ache a fração de escolhas de $A$ nas últimas $100$ de $500$ tentativas com $d=0.5\,\text{s}$ e $\tau_e=0.05$, $0.2$, $1$, $4\,\text{s}$, e com $\tau_e=1\,\text{s}$, $d=3\,\text{s}$. Com que fração do traço vivo até a recompensa (exercício~\ref{ex:06:1}) o aprendizado some?
\end{exercise}

\begin{exercise}[\lvlC]
\label{ex:06:8}
\emph{Dá para escapar do preço da difusão?} $N$ neurônios: $y_i=w_i+\xi_i$, $\xi_i\sim\mathcal N(0,\sigma^2)$, recompensa $R=-\sum_i(y_i-w_i^*)^2$, regra $\Delta w_i=\eta\,(R-\langle R\rangle)\,\xi_i$. Em quantas tentativas, com o melhor $\eta$, o erro $\sum_i(w_i-w_i^*)^2$ cai $e$ vezes? E se só $m$ neurônios aleatórios de $N$ flutuarem? A busca esparsa ajuda?
\end{exercise}

\chapter{Teorias alternativas modernas do aprendizado baseado em dopamina}
\chaptermark{Outras teorias do \nt{DOPA}}
\label{sec:dofaalt}
\begin{chaptergoals}
\item como neurônios \nt{DOPA} com respostas assimétricas ao erro registram a distribuição da recompensa;
\item como alguns neurônios \nt{DOPA} informam a importância ou o perigo, e não o sinal do erro;
\item como um fio leva um erro rápido que ensina e um nível lento que define a velocidade das ações;
\item por que o sinal \nt{DOPA} é um feixe, e não um número, e que teoria contesta o próprio erro.
\end{chaptergoals}

\section{O que realmente passa pelo fio}

No capítulo~\ref{sec:dofa}, o sistema \nt{DOPA} era para nós um único fio pelo qual passa um único número: o erro de predição de recompensa $\delta$~\eqref{eq:ctd}. Esse quadro explica muito, mas quanto mais precisamente se mede a dopamina, mais se encontra o que não cabe nele. Alguns neurônios respondem ao desagradável como a uma recompensa; a concentração de dopamina sobe enquanto o animal corre para o alvo, embora nada de inesperado aconteça; neurônios \nt{DOPA} diferentes se comportam de modo diferente.

Para o engenheiro, todas essas descobertas se resumem a uma pergunta: \emph{qual é o formato do sinal no barramento de difusão?} Um número ou um vetor? Um sinal no fio ou vários, separados em frequência? Ele fala do futuro, como $\delta$, ou do passado? A seguir, seis respostas modernas; para cada uma, vamos separar o medido do suposto.

\section{Não um número, mas uma distribuição}

O erro~\eqref{eq:ctd} ensina ao crítico a expectativa \emph{média} de recompensa. Mas meia gota de suco garantida e uma gota com probabilidade de um meio têm a mesma média. Para distingui-las, é preciso a distribuição inteira da recompensa.

Tomemos o problema mais simples: depois de um sinal antecipador chega uma recompensa de tamanho aleatório $r$. Suponha que há muitos neurônios \nt{DOPA} e cada um, o $i$-ésimo, cuida de sua própria estimativa $V_i$, de modo que no instante da recompensa seu erro é $\delta_i=r-V_i$. E suponha que cada um conta o erro positivo com peso $\alpha_i^+$ e o negativo com peso $\alpha_i^-$. Depois de cada recompensa, a estimativa se desloca em
\begin{equation}
\Delta V_i \;=\; \eta\,\bigl(\alpha_i^+\,[\delta_i]_+ \;-\; \alpha_i^-\,[-\delta_i]_+\bigr).
\label{eq:asymmetric}
\end{equation}
A estimativa deixa de mudar quando, em média, as correções positivas equilibram as negativas:
\begin{empheq}[box=\keybox]{equation}
\alpha_i^+\,\bigl\langle[r-V_i]_+\bigr\rangle \;=\; \alpha_i^-\,\bigl\langle[V_i-r]_+\bigr\rangle ,
\qquad
\kappa_i \;=\; \frac{\alpha_i^+}{\alpha_i^++\alpha_i^-} .
\label{eq:expectile}
\end{empheq}
Com $\kappa_i=1/2$, é a média comum. Com $\kappa_i>1/2$, o neurônio é ``otimista'': sua estimativa fica deslocada para o extremo superior da distribuição; com $\kappa_i<1/2$, é ``pessimista''.

Exemplo: o suco chega com probabilidade $1/2$, gota de tamanho $1$. Então~\eqref{eq:expectile} dá $\kappa_i\cdot\tfrac12(1-V_i)=(1-\kappa_i)\cdot\tfrac12V_i$, isto é, $V_i=\kappa_i$ (fig.~\ref{fig:expectiles}). Um conjunto de neurônios com $\kappa_i$ diferentes registra a distribuição da recompensa assim como um conjunto de quantis registra uma distribuição em estatística.

\begin{figure}[t]
\centering
\begin{tikzpicture}[>=Stealth,x=7cm,y=3cm]
\draw[->,gray!70!black] (-0.08,0) -- (1.15,0) node[right,black] {\small recompensa};
\draw[->,gray!70!black] (-0.08,0) -- (-0.08,0.75) node[left,black] {\small probabilidade};
\fill[blue!25] (-0.03,0) rectangle (0.03,0.5);
\fill[blue!25] (0.97,0) rectangle (1.03,0.5);
\node[below,font=\small] at (0,0) {$0$};
\node[below,font=\small] at (1,0) {$1$};
\node[left,font=\scriptsize] at (-0.08,0.5) {$1/2$};
\foreach \k/\c/\lab/\h in {0.2/blue!60!black/{pessimista, $\kappa=0.2$}/0.62, 0.5/gray!50!black/{média, $\kappa=0.5$}/0.42, 0.8/red!70!black/{otimista, $\kappa=0.8$}/0.22}{
  \draw[very thick,\c] (\k,0) -- (\k,\h);
  \node[above,font=\scriptsize,\c] at (\k,\h) {\lab};}
\end{tikzpicture}
\caption{A distribuição da recompensa registrada por um conjunto de neurônios \nt{DOPA}. Recompensa $0$ ou $1$ com probabilidade $1/2$ (barras azuis); o neurônio com fração $\kappa$ chega à estimativa $V=\kappa$~\eqref{eq:expectile}.}
\label{fig:expectiles}
\end{figure}

O que foi medido: neurônios \nt{DOPA} diferentes têm ``pontos de reversão'' diferentes (o tamanho de recompensa em que a resposta passa de pico a queda), e os neurônios com maior assimetria entre as respostas a erros positivos e negativos se invertem em recompensas maiores, como exige~\eqref{eq:expectile}~\cite{Dabney2020}. A partir das respostas de um grupo de neurônios é possível reconstruir a forma da distribuição. Que essa seja a função principal da dispersão, e não um efeito colateral, é uma hipótese.

\begin{keyblock}
\begin{proposition}[Distribuição a partir da assimetria]
\label{prop:expectile}
Se os neurônios \nt{DOPA} diferem na fração $\kappa_i$ com que contam os erros positivos, suas estimativas convergem para pontos diferentes~\eqref{eq:expectile} da distribuição da recompensa. O grupo de neurônios transmite não a média, mas a distribuição e, com ela, o risco.
\end{proposition}
\end{keyblock}

\section{Não só o erro, mas também a importância}

Pela fórmula do erro~\eqref{eq:ctd}, um evento desagradável (um choque, um gosto amargo) deveria provocar uma queda: saiu pior que o esperado. Parte dos neurônios \nt{DOPA} faz isso, mas outra parte responde ao desagradável com um \emph{pico}, como a uma recompensa~\cite{Matsumoto2009}. Esses neurônios não informam o sinal do erro, mas que aconteceu algo \emph{importante}: inesperado, forte, que exige atenção~\cite{BrombergMartin2010}.

Ao engenheiro convém pensar nisso como dois sinais. O primeiro é o erro com sinal $\delta$: para onde mudar os pesos. O segundo é seu módulo $|\delta|$ ou a novidade do evento: \emph{quanto} mudá-los, isto é, a taxa de aprendizado; depois de um evento inesperado é preciso aprender mais depressa. Pelo sentido, ele se aproxima do ganho noradrenérgico do capítulo~\ref{sec:neurontypes}.

Há ainda uma terceira variante: um fio separado para um eixo separado. Os neurônios \nt{DOPA} que vão para uma das regiões de escolha de ações não respondem à recompensa, mas à novidade e à intensidade do estímulo, e são necessários para que o animal aprenda a evitar o perigoso~\cite{Menegas2018}: pelo fio não passa ``melhor ou pior'', mas ``perigoso ou não''.

O que foi medido: grupos diferentes de neurônios \nt{DOPA} respondem de modo diferente ao desagradável e ao novo, e saídas diferentes ensinam coisas diferentes. Como exatamente o cérebro usa o sinal de importância (como taxa de aprendizado, como sinal de atenção ou como erro separado no eixo do perigo) está em discussão.

\section{Não só ensinar, mas também apressar}

A dopamina tem também um efeito imediato: quanto mais dela, mais disposto e mais rápido o animal age. Como conciliar isso com o aprendizado?

A resposta de engenharia é a \emph{separação em frequência}. Picos e quedas rápidos, com duração de centenas de milissegundos, levam $\delta$ e ensinam. O nível lento, que muda em minutos, leva outra grandeza: a recompensa média por unidade de tempo $\bar R$~\cite{Niv2007}. Suponha que uma ação executada no tempo $\tau_a$ exige um esforço $C/\tau_a$: quanto mais rápida, mais cara. Em compensação, cada segundo de demora custa $\bar R$, a recompensa que se poderia obter nele. O custo total $C/\tau_a+\bar R\,\tau_a$ é mínimo em
\begin{empheq}[box=\keybox]{equation}
\tau_a^{*} \;=\; \sqrt{\frac{C}{\bar R}} .
\label{eq:vigor}
\end{empheq}
Quanto mais rico o ambiente, mais caro o tempo e mais vale a pena agir depressa: se $\bar R$ dobrou, o tempo de ação cai $\sqrt2\approx1.4$ vezes. Note que $\bar R$ é a mesma recompensa esperada que subtraímos no capítulo~\ref{sec:dofa}.

A liberação de dopamina no destino pode mudar mesmo sem mudança na taxa de disparos: ela é ajustada por sinais locais nas próprias saídas~\cite{Mohebi2019}. Mas há componentes lentos também nos próprios disparos: nos experimentos da próxima seção, o crescimento suave ao se aproximar da recompensa aparece também na taxa dos neurônios \nt{DOPA}~\cite{Kim2020}. Logo, o nível lento se compõe de duas fontes, a taxa de disparos e o ajuste local da liberação, e qual delas é a principal parece depender da saída e da tarefa.

\begin{keyblock}
\begin{proposition}[Dois sinais num só fio]
\label{prop:multiplex}
O sistema \nt{DOPA} transmite pelo menos duas grandezas separadas no tempo: o erro rápido $\delta$, que ensina, e o nível lento, que define o preço do tempo~\eqref{eq:vigor} e, com isso, a velocidade das ações. Foi medido que os componentes rápido e lento se comportam de modo diferente; que o lento seja exatamente $\bar R$ é uma hipótese.
\end{proposition}
\end{keyblock}

\section{Crescimento na aproximação}

Se o animal corre por um corredor rumo a uma recompensa distante, a concentração de dopamina na região de escolha de ações sobe suavemente durante segundos, enquanto o alvo se aproxima~\cite{Hamid2016}. Pelo capítulo~\ref{sec:dofa}, isso não deveria acontecer: tudo corre conforme o plano, e $\delta$ deveria ser zero.

Primeira explicação: esse crescimento não é $\delta$, mas a própria expectativa $V$, o sinal ``o alvo está perto, vale a pena se esforçar'' da seção anterior~\cite{Hamid2016}. A segunda explicação preserva $\delta$~\cite{Kim2020}. Se a expectativa está correta, com horizonte $\tau_\gamma$ ela cresce no caminho até a recompensa $R$ como $V=Re^{-(T-t)/\tau_\gamma}$, e pela fórmula do erro~\eqref{eq:ctd} $\dd V/\dd t-V/\tau_\gamma=0$. Mas suponha que de longe a expectativa esteja subestimada e que a estimativa se refine com a aproximação. Então a expectativa cresce mais íngreme, digamos como $V=Re^{-(T-t)/\tau_v}$ com $\tau_v<\tau_\gamma$, e
\begin{empheq}[box=\keybox]{equation}
\delta(t) \;=\; \frac{\dd V}{\dd t}-\frac{V}{\tau_\gamma} \;=\; V\Bigl(\frac{1}{\tau_v}-\frac{1}{\tau_\gamma}\Bigr) \;>\; 0 .
\label{eq:ramp}
\end{empheq}
O erro é positivo e cresce junto com $V$: é o crescimento suave (fig.~\ref{fig:ramp}). Com $\tau_\gamma=4\,\text{s}$ e $\tau_v=1\,\text{s}$, obtém-se $\delta=0.75\,V$ por segundo. Essa versão foi testada num corredor virtual, transportando o animal instantaneamente para mais perto do alvo: a dopamina respondeu com um pico, como convém a uma derivada, e não com um nível suave~\cite{Kim2020}.

\begin{figure}[t]
\centering
\begin{tikzpicture}[>=Stealth,x=1.6cm,y=2.6cm]
\draw[->,gray!70!black] (-4.2,0) -- (0.5,0) node[below left,black] {\small $t$};
\draw[->,gray!70!black] (-4.2,0) -- (-4.2,1.2);
\foreach \x/\l in {-4/{$T-4$},-2/{$T-2$},0/{$T$}}{\draw[gray!60!black] (\x,-0.03) -- (\x,0.03); \node[below,font=\scriptsize] at (\x,0) {\l\,s};}
\draw[thick,densely dashed,gray!60!black] plot[domain=-4:0,samples=40] (\x,{exp(\x/4)});
\draw[very thick,blue!60!black] plot[domain=-4:0,samples=60] (\x,{exp(\x)});
\draw[very thick,red!70!black] plot[domain=-4:0,samples=60] (\x,{0.75*exp(\x)});
\node[anchor=south east,font=\small,gray!50!black] at (-1.3,0.77) {$V$ com $\tau_v=\tau_\gamma$: $\delta=0$};
\node[right,font=\small,blue!60!black] at (0.05,1.0) {$V$ com $\tau_v=1\,\text{s}$};
\node[right,font=\small,red!70!black] at (0.05,0.72) {$\delta$ por~\eqref{eq:ramp}};
\end{tikzpicture}
\caption{O crescimento da dopamina ao se aproximar da recompensa como erro de predição, $\tau_\gamma=4\,\text{s}$. Linha tracejada: expectativa que cresce exatamente como o horizonte exige ($\delta=0$); azul: expectativa que cresce mais íngreme; vermelha: o erro~\eqref{eq:ramp}.}
\label{fig:ramp}
\end{figure}

O que foi medido: o crescimento suave existe, tanto na concentração de dopamina quanto na taxa de disparos dos neurônios \nt{DOPA}~\cite{Kim2020}, e saltos instantâneos do ambiente causam saltos de dopamina. Qual das duas explicações é a correta, ou se ambas valem em saídas diferentes, está em discussão.

\section{Olhar para trás}

Todas as teorias acima olham \emph{para a frente}. Uma das teorias modernas inverte a direção~\cite{Jeong2022}. Suponha que o cérebro não aprenda a prever a recompensa pelo sinal antecipador, mas que, ao receber a recompensa, \emph{olhe para trás}: que evento a precedeu com mais frequência que os outros? A medida dessa ligação é a diferença de duas probabilidades:
\begin{empheq}[box=\keybox]{equation}
M \;=\; P(\text{sinal antes da recompensa}) \;-\; P(\text{sinal numa janela aleatória}) .
\label{eq:retro}
\end{empheq}
Se o sinal antecipador aparece muitas vezes sem recompensa, o segundo termo é grande, e a ligação é fraca. Nessa teoria, a dopamina não informa o erro de predição, mas o quanto um dado evento é uma provável \emph{causa} da recompensa.

O mecanismo do olhar para trás exige o mesmo que a regra de três fatores do capítulo~\ref{sec:dofa}: cada evento deve deixar um traço que se desfaz, para que no instante da recompensa se possa ler o que veio antes dela. A diferença não está na sinapse, mas no que o \nt{DOPA} transmite.

Em experimentos em que a teoria~\eqref{eq:retro} e o erro~\eqref{eq:ctd} preveem coisas diferentes, a liberação de dopamina seguiu, em vários casos, a primeira~\cite{Jeong2022}. Mas em outros experimentos a resposta da dopamina durante o aprendizado se deslocou gradualmente do instante da recompensa para o instante do sinal antecipador, como exige o aprendizado pelo erro~\eqref{eq:ctd}~\cite{Amo2022}. Qual teoria descreve o cérebro ainda não se sabe.

\section{Um erro não só na recompensa}

A última extensão é sobre \emph{do que} é o erro. Se o suco esperado é trocado por outro, igualmente valioso mas de outro sabor, o valor não mudou, e o erro~\eqref{eq:ctd} é zero. Mas os neurônios \nt{DOPA} respondem a essa troca~\cite{Takahashi2017}. E picos de dopamina provocados artificialmente podem ensinar ao animal uma ligação entre dois estímulos neutros, sem nenhuma recompensa~\cite{Sharpe2017}. Parece que o \nt{DOPA} informa o erro de predição não só da recompensa, mas também de \emph{o que} vai acontecer.

Formalmente, é uma generalização simples de~\eqref{eq:ctd}: no lugar de um único fluxo de recompensa $r$, toma-se um conjunto de características do que acontece $\varphi_k$ (sabor, lugar, som) e, para cada uma, sua expectativa $M_k$ e seu erro~\cite{Gardner2018}:
\begin{empheq}[box=\keybox]{equation}
\delta_k(t) \;=\; \varphi_k(t) \;+\; \frac{\dd M_k}{\dd t} \;-\; \frac{M_k}{\tau_\gamma} ,
\qquad k=1,\dots,K .
\label{eq:vectortd}
\end{empheq}
A recompensa é só uma das características, e o vetor de erros ensina não só o que é bom, mas também o que segue o quê: um modelo do mundo.

Isso concorda com a heterogeneidade dos neurônios \nt{DOPA}: numa mesma tarefa, neurônios diferentes levam grandezas diferentes, uns ligados à recompensa, outros ao movimento, outros ainda à direção da atenção~\cite{Engelhard2019}.

\begin{keyblock}
\begin{proposition}[Um feixe em vez de um fio]
\label{prop:bundle}
O sinal do sistema \nt{DOPA} não é um único número. Ele é distribuído entre os neurônios (a distribuição~\eqref{eq:expectile}, características diferentes~\eqref{eq:vectortd}, um eixo separado de perigo) e no tempo (erro rápido e nível lento~\eqref{eq:vigor}). O erro escalar~\eqref{eq:ctd} é uma primeira aproximação desse sinal, assim como o somador com limiar é uma primeira aproximação do neurônio.
\end{proposition}
\end{keyblock}

\section{Resumo}

\begin{table}[t]
\centering
\footnotesize
\setlength{\tabcolsep}{4pt}
\renewcommand{\arraystretch}{1.15}
\begin{tabular}{>{\raggedright}p{2.5cm}>{\raggedright}p{3.2cm}>{\raggedright}p{3.6cm}>{\raggedright\arraybackslash}p{4.2cm}}
\toprule
Teoria & O que passa pelo fio & Medição principal & O que se sabe \\
\midrule
erro de predição (cap.~\ref{sec:dofa}) & escalar $\delta$~\eqref{eq:ctd} & pico, transferência para o sinal antecipador, queda & medido; primeira aproximação \\
distribuição & conjunto de $\delta_i$ com assimetrias diferentes & pontos de reversão diferentes em neurônios diferentes & há concordância com o experimento; o papel da dispersão é hipótese \\
importância & $|\delta|$, novidade; eixo do perigo & picos ao desagradável em parte dos neurônios & medido; como é usado está em discussão \\
preço do tempo & nível lento $\bar R$ & a dopamina acelera as ações; componente lento na liberação nas saídas e na taxa de disparos & separação entre rápido e lento medida; $\bar R$ é hipótese \\
crescimento na aproximação & $V$ ou $\delta$ com refinamento da estimativa~\eqref{eq:ramp} & crescimento suave, saltos ao transportar no corredor & crescimento medido; explicação em discussão \\
olhar para trás & medida de ligação causal~\eqref{eq:retro} & experimentos em que as previsões das duas teorias divergem & resultados contraditórios \\
vetor de erros & $\delta_k$ por característica~\eqref{eq:vectortd} & resposta à troca de sabor; aprendizado sem recompensa & medido; a forma vetorial é hipótese \\
\bottomrule
\end{tabular}
\caption{O que o sistema \nt{DOPA} transmite: o quadro padrão do capítulo~\ref{sec:dofa} e suas extensões modernas.}
\label{tab:dofaalt}
\end{table}

As teorias da tabela~\ref{tab:dofaalt} não se excluem: quase todas mantêm o erro de predição como base e ampliam seu formato. Só o olhar para trás compete seriamente com ele, e essa disputa será decidida por experimentos. Para o engenheiro, a conclusão é simples: o preço da difusão da proposição~\ref{prop:broadcastcost} cai se o fio for, na verdade, muitos fios com destinatários diferentes.

\section{Exercícios}

\begin{exercise}[\lvlA]
\label{ex:07:1}
\emph{Pontos da distribuição para uma gota.} Um suco de tamanho $1$ chega com probabilidade $p=0.2$; caso contrário, a recompensa é $0$. Encontre por~\eqref{eq:expectile} $V_i$ para $\kappa_i=0.2$, $0.5$, $0.8$. Qual é a mediana dessa recompensa, e em que o ponto~\eqref{eq:expectile} difere de um quantil?
\end{exercise}

\begin{exercise}[\lvlB]
\label{ex:07:2}
\emph{Distribuição a partir de dois neurônios.} A recompensa é $0$ ou $x$, e a probabilidade de $x$ é $p$. Dois neurônios com a regra~\eqref{eq:asymmetric} informaram $V(1/2)=0.25$ e $V(0.8)=0.5$. Reconstrua $p$, $x$ e o desvio padrão da recompensa. O que informará um neurônio com $\kappa=0.2$?
\end{exercise}

\begin{exercise}[\lvlB]
\label{ex:07:3}
\emph{Simulação: distribuição a partir da assimetria.} Simule nove neurônios \nt{DOPA} com $\kappa_i=0.1,\dots,0.9$ e a regra~\eqref{eq:asymmetric}, $\alpha_i^+=2\kappa_i$, $\alpha_i^-=2(1-\kappa_i)$, com recompensa uniforme em $[0,1]$. Como a dispersão de $V_i$ depende de $\eta$, e que $\eta$ é preciso para distinguir essa distribuição de duas gotas equiprováveis $0.25$ e $0.75$?
\end{exercise}

\begin{exercise}[\lvlA]
\label{ex:07:4}
\emph{O preço do tempo.} (a)~Deduza~\eqref{eq:vigor} e encontre o custo total no ótimo. (b)~O cérebro estimou $\bar R$ com erro de um fator $k$. Encontre o custo excedente relativo para $k=2$ e $k=4$. Com que precisão o nível lento de dopamina deve informar $\bar R$?
\end{exercise}

\begin{exercise}[\lvlB]
\label{ex:07:5}
\emph{Pico no transporte.} A expectativa cresce por~\eqref{eq:ramp} com $\tau_v=1\,\text{s}$, $\tau_\gamma=4\,\text{s}$; o animal é transportado instantaneamente do ponto $T-2\,\text{s}$ para $T-1\,\text{s}$. Encontre a área do pico de $\delta$ em frações de $R$ e quantos segundos de rampa antes do transporte ele substitui. O que o transporte mostrará se a dopamina informar o próprio $V$?
\end{exercise}

\begin{exercise}[\lvlB]
\label{ex:07:6}
\emph{Projeto: dois sinais no fio.} Pelo fio \nt{DOPA} passam um nível que cresce até $s=1/60$ disparo por segundo a cada segundo e eventos de $A=3$ disparos. Uma sinapse com traço $\tau_e=1\,\text{s}$ subtrai da taxa uma cópia dela suavizada com $\tau_f$. Com que $\tau_f$ se perde no máximo $10\,\%$ do evento e a contribuição do nível no tempo do traço fica abaixo de $10\,\%$ da do evento?
\end{exercise}

\begin{exercise}[\lvlB]
\label{ex:07:7}
\emph{Projeto de experimento: olhar para trás contra erro.} Numa janela, o sinal antecipador ocorre com probabilidade $0.1$, e depois dele a recompensa com probabilidade $0.8$. Compare $\Delta P=P(\text{recomp.}\mid\text{sinal})-P(\text{recomp.}\mid\text{sem})$ e $M$ de~\eqref{eq:retro} se (a)~forem acrescentadas recompensas sem sinal, $0.08$ por janela; (b)~a frequência do sinal sem recompensas for dobrada.
\end{exercise}

\begin{exercise}[\lvlC]
\label{ex:07:8}
\emph{O feixe e o preço da difusão.} No modelo do exercício~\ref{ex:06:8}, $N$ neurônios são divididos em $K$ grupos com fios próprios, e em cada fio vaza uma fração $\rho$ dos erros alheios. Mostre que a constante de aprendizado é $N/K+2+\rho^2N(1-1/K)$. Quantas tentativas ela dá com $K\to\infty$, $N=10^{10}$ e $\rho=0.01$?
\end{exercise}

\chapter{Quando buscar e quando decidir: acrescentamos o \nt{NORA}}
\chaptermark{Acrescentamos o \nt{NORA}}
\label{sec:nora}

\section{O que falta}

No capítulo~\ref{sec:dofa}, a rede aprendia graças ao ruído: os desvios aleatórios de atividade eram a busca, e o \nt{DOPA} informava se eles tinham melhorado as coisas. Mas lá ficou sem definir um parâmetro: \emph{quanto} buscar. Se o ruído é forte demais, a rede se debate e não aproveita o que já sabe. Se é fraco demais, ela escolhe sempre a mesma coisa e não percebe que, em algum lugar perto, as coisas melhoraram. No aprendizado por reforço, isso se chama o dilema entre aproveitamento e busca, e todo algoritmo tem um botão para ele. No capítulo~\ref{sec:neurontypes}, supusemos que no cérebro quem gira esse botão é o \nt{NORA}.

Os neurônios noradrenérgicos no ser humano são algumas dezenas de milhares (tabela~\ref{tab:types}), quase todos reunidos num único grupo pequeno, e suas saídas se espalham por praticamente o cérebro inteiro, embora partes diferentes desse grupo, como se descobriu, atendam em parte regiões diferentes~\cite{Poe2020}. É um canal de difusão ainda mais estreito que o \nt{DOPA}. Ele funciona em dois modos~\cite{AstonJones2005}. No modo \emph{de fundo}, os neurônios emitem disparos regulares, com taxa de frações a alguns disparos por segundo, e essa taxa muda lentamente com a vigília. No modo \emph{de rajadas}, eles respondem com uma rajada curta a um evento importante para a tarefa atual, mais ou menos no instante da decisão. Um ponto de teste cômodo, mas impreciso, é o diâmetro da pupila: ele muda com a atividade desses neurônios, mas também com a atividade de outras regiões~\cite{Joshi2016} e das saídas colinérgicas no córtex~\cite{Reimer2016}. A pupila mostra o nível geral de excitação, e não só o \nt{NORA}.

\section{O ganho}

O que foi medido é isto: a noradrenalina suprime a atividade de fundo dos neurônios mais do que sua resposta ao estímulo, e a relação entre sinal e fundo aumenta~\cite{Foote1975NE}. Que a inclinação da característica de um neurônio isolado seja literalmente multiplicada por um coeficiente não decorre do experimento. Como no capítulo~\ref{sec:neurontypes}, tomemos um modelo simples que reproduz esse efeito: a noradrenalina muda a inclinação da característica de transferência do destinatário~\cite{ServanSchreiber1990}. As entradas fracas, sublimiares (o fundo), passam a dar ainda menos, e as fortes, ainda mais. Seja a resposta do neurônio a taxa
\begin{empheq}[box=\keybox]{equation}
r \;=\; F\bigl(g\,(u-\theta)\bigr), \qquad F(z)=\frac{r_{\max}}{1+e^{-z}} ,
\label{eq:gain}
\end{empheq}
em que $u$ é a corrente de entrada, $\theta$ é o limiar, e $g$ é o ganho, que no modelo é controlado pelo \nt{NORA}. Com $g$ pequeno, a taxa acompanha suavemente a entrada numa faixa ampla: o neurônio é um \emph{amplificador analógico}. Com $g$ grande, quase qualquer entrada acima do limiar dá a taxa máxima, e abaixo dele, zero: o neurônio é um \emph{comparador}, um elemento quase digital (fig.~\ref{fig:gaincurves}). Um único botão comuta o mesmo neurônio entre dois modos que, em eletrônica, são feitos por circuitos diferentes.

\begin{figure}[t]
\centering
\begin{tikzpicture}[>=Stealth,x=1.1cm,y=2.4cm]
\draw[->,gray!70!black] (-3.3,0) -- (3.4,0) node[below left,black] {\small $u$};
\draw[->,gray!70!black] (-3.3,0) -- (-3.3,1.2) node[left,black] {\small $r$};
\draw[densely dashed,gray!60!black] (0,0) -- (0,1.1);
\node[below,font=\small] at (0,0) {$\theta$};
\draw[densely dotted,gray!60!black] (-3.3,1) -- (3.2,1) node[right,black,font=\small] {$r_{\max}$};
\draw[very thick,blue!60!black] plot[domain=-3.2:3.2,samples=60] (\x,{1/(1+exp(-0.8*\x))});
\draw[very thick,orange!85!black] plot[domain=-3.2:3.2,samples=60] (\x,{1/(1+exp(-2*\x))});
\draw[very thick,red!70!black] plot[domain=-3.2:3.2,samples=120] (\x,{1/(1+exp(min(-8*\x,9)))});
\node[blue!60!black,font=\small,anchor=west] at (0.5,0.12) {$g$ pequeno: amplificador};
\node[red!70!black,font=\small,anchor=east] at (-0.3,0.85) {$g$ grande: comparador};
\end{tikzpicture}
\caption{Característica de transferência~\eqref{eq:gain} com três valores do ganho $g$.}
\label{fig:gaincurves}
\end{figure}

Um ganho grande ajuda contra o ruído? Nem sempre. Suponha que duas entradas diferem em $\Delta u$ e é preciso dizer qual é maior. Se o ruído $\sigma_{\text{antes}}$ é somado à entrada \emph{antes} do amplificador, amplificam-se o sinal e o ruído, e sua relação não muda. Mas se o ruído $\sigma_{\text{depois}}$ é somado \emph{depois} do amplificador (vindo de sinapses pouco confiáveis e de disparos aleatórios da própria rede, como no capítulo~\ref{sec:code}), a relação sinal-ruído
\begin{empheq}[box=\keybox]{equation}
d' \;=\; \frac{g\,\Delta u}{\sqrt{g^2\sigma_{\text{antes}}^2+\sigma_{\text{depois}}^2}}
\label{eq:gainsnr}
\end{empheq}
cresce com $g$ até que $g\sigma_{\text{antes}}$ se iguale a $\sigma_{\text{depois}}$~\cite{ServanSchreiber1990}.

\begin{keyblock}
\begin{proposition}[Quando o ganho ajuda]
\label{prop:gainnoise}
Aumentando o ganho, o \nt{NORA} melhora a discriminação das entradas só contra o ruído que surge \emph{depois} do amplificador, na própria rede. O ruído que chegou junto com a entrada, o ganho não remove.
\end{proposition}
\end{keyblock}

\section{O ganho na escolha entre duas opções}

Voltemos à escolha do capítulo~\ref{sec:gaba}: dois grupos de neurônios \nt{GLUT} com entradas $I_1$, $I_2$ se inibem mutuamente com força $\beta$. Agora cada grupo tem um ganho $g$: nas equações da escolha~\eqref{eq:wta}, a expressão entre parênteses é multiplicada por $g$. Enquanto os dois grupos funcionam, as taxas estacionárias saem de $r_1=g(I_1-\beta r_2)$, $r_2=g(I_2-\beta r_1)$. Somando e subtraindo essas igualdades, obtemos a soma $r_1+r_2=g(I_1+I_2)/(1+g\beta)$ e a diferença $r_1-r_2=g(I_1-I_2)/(1-g\beta)$. A razão entre elas diz com que nitidez a rede prefere uma entrada à outra:
\begin{empheq}[box=\keybox]{equation}
\frac{r_1-r_2}{r_1+r_2} \;=\; \varepsilon\,\frac{1+g\beta}{1-g\beta}, \qquad \varepsilon=\frac{I_1-I_2}{I_1+I_2}, \quad g\beta<1 .
\label{eq:wtagain}
\end{empheq}
Uma pequena diferença de entradas $\varepsilon$ é amplificada $(1+g\beta)/(1-g\beta)$ vezes, e esse fator cresce sem limite quando $g\beta$ se aproxima de um. Com $g\beta>1$, o estado em que os dois grupos funcionam é instável, como na proposição~\ref{prop:wta}: um vence, o outro se cala (fig.~\ref{fig:pitchfork}).

\begin{figure}[t]
\centering
\begin{tikzpicture}[>=Stealth,x=3.2cm,y=1.5cm]
\draw[->,gray!70!black] (0,0) -- (2.15,0) node[below left,black] {\small $g\beta$};
\draw[->,gray!70!black] (0,-1.25) -- (0,1.35) node[above,black] {\small $(r_1-r_2)/(r_1+r_2)$};
\draw[gray!60!black] (1,-0.04) -- (1,0.04); \node[below right,font=\small] at (1,0) {$1$};
\node[left,font=\scriptsize] at (0,1) {$1$}; \node[left,font=\scriptsize] at (0,-1) {$-1$};
\draw[very thick,blue!60!black] plot[domain=0:0.905,samples=60] (\x,{0.05*(1+\x)/(1-\x)}) -- (1,1) -- (2.05,1);
\draw[very thick,blue!60!black] (1.105,-1) -- (2.05,-1);
\draw[thick,densely dashed,gray!60!black] (1,0) -- (2.05,0);
\node[font=\small,align=center] at (0.45,-0.62) {pondera:\\os dois grupos\\funcionam};
\node[font=\small,align=center] at (1.55,0.45) {decidiu:\\só um funciona};
\node[font=\small,blue!60!black,anchor=south west] at (1.05,-0.97) {a entrada fraca venceu antes e se mantém};
\end{tikzpicture}
\caption{Escolha entre duas opções com ganhos diferentes; as entradas diferem em $\varepsilon=0.05$. Com $g\beta>1$, as duas decisões são estáveis (linhas contínuas; a vitória da entrada fraca, com $g\beta>(1+\varepsilon)/(1-\varepsilon)\approx1.1$), e a rede mantém a que já tomou; o estado simétrico (linha tracejada) é instável.}
\label{fig:pitchfork}
\end{figure}

\begin{keyblock}
\begin{proposition}[O ganho comuta o modo de escolha]
\label{prop:gainwta}
Numa rede de escolha com inibição mútua $\beta$, o ganho $g$ leva a rede a atravessar a fronteira $g\beta=1$. Abaixo dela, a rede ``pondera'': os dois grupos funcionam, e a diferença de entradas é amplificada segundo~\eqref{eq:wtagain}. Acima, a rede ``decidiu'': um grupo funciona, e a decisão se mantém. Mudando $g$, o \nt{NORA} pode comutar a rede entre esses modos sem tocar em nenhum peso.
\end{proposition}
\end{keyblock}

\section{O ganho como temperatura}

Agora acrescentemos à escolha um ruído que surge na própria rede, depois do amplificador. Qual grupo vence é decidido pelo sinal da grandeza $g(u_A-u_B)+\eta$, em que $u_A-u_B$ é a diferença de entradas, isto é, a diferença dos pesos aprendidos no capítulo~\ref{sec:dofa}, e $\eta$ é um ruído de escala $\sigma$. Se o ruído tem distribuição logística (para a gaussiana, a curva é quase a mesma), a probabilidade de escolher $A$ é
\begin{empheq}[box=\keybox]{equation}
P(A) \;=\; \frac{1}{1+e^{-g\,(u_A-u_B)/\sigma}} .
\label{eq:softmax}
\end{empheq}
O programador reconhece aqui a escolha por softmax com temperatura $T=\sigma/g$. O ganho é o inverso da temperatura. Com $g$ pequeno, mesmo uma opção claramente melhor é escolhida só um pouco mais vezes que a pior: a rede busca muito. Com $g$ grande, a melhor opção conhecida é escolhida quase sempre: a rede aproveita o que sabe.

Precisemos o que é $g$ em~\eqref{eq:wtagain} e~\eqref{eq:softmax}: é o ganho \emph{efetivo} da diferença de entradas da tarefa atual no instante da escolha. Os dois modos do \nt{NORA} agem sobre ele em sentidos opostos. A rajada curta no instante da decisão o eleva, e a rede consolida a escolha. A atividade de fundo alta, ao contrário, acompanha a busca: em humanos, antes de uma escolha ao acaso, a pupila basal é mais larga~\cite{JepmaNieuwenhuis2011}, e em ratos a entrada do \nt{NORA} no córtex cingulado anterior leva a escolha a um modo aleatório~\cite{Tervo2014stochastic}. Logo, ao \nt{NORA} de fundo alto corresponde um $g$ efetivo \emph{pequeno}, e à rajada, um $g$ grande.

\section{Quanto buscar}

Qual ganho é melhor? Verificamos isso num modelo. A rede escolhe uma de duas ações por~\eqref{eq:softmax}; cada uma traz recompensa com certa probabilidade, e a rede a estima pelas recompensas da ação escolhida, como no capítulo~\ref{sec:dofa}. De tempos em tempos o mundo muda: as duas probabilidades são sorteadas de novo. Pode mudar tanto a ação escolhida quanto aquela que a rede não experimenta há muito tempo; desta última a rede só pode saber pela busca. A fig.~\ref{fig:invertedU} mostra que fração da melhor recompensa possível a rede obtém com diferentes $g$ constantes.

\begin{figure}[t]
\centering
\begin{tikzpicture}[>=Stealth,x=1.2cm,y=1cm]
\draw[->,gray!70!black] (0,0) -- (7.6,0) node[below left,black] {\small $g$};
\draw[->,gray!70!black] (0,0) -- (0,4.4) node[above,black,font=\small] {fração da melhor recompensa};
\foreach \x/\l in {0/0.5,1/1,2/2,3/4,4/8,5/16,6/32,7/64}{\draw[gray!60!black] (\x,-0.06) -- (\x,0.06); \node[below,font=\scriptsize] at (\x,0) {\l};}
\foreach \v/\y in {0.8/0.8,0.9/2.4,1.0/4.0}{\draw[gray!60!black] (-0.06,\y) -- (0.06,\y); \node[left,font=\scriptsize] at (0,\y) {\v};}
\draw[very thick,blue!60!black] plot[mark=*,mark size=1.3pt] coordinates {(0,0.368) (1,0.880) (2,1.728) (3,2.784) (4,3.456) (5,3.616) (6,3.488) (7,3.264)};
\draw[very thick,red!70!black] plot[mark=*,mark size=1.3pt] coordinates {(0,0.368) (1,0.672) (2,1.184) (3,1.840) (4,2.176) (5,1.920) (6,1.456) (7,1.104)};
\node[blue!60!black,font=\small,anchor=south] at (5,3.7) {mundo calmo};
\node[red!70!black,font=\small,anchor=north] at (5.1,1.25) {mundo que muda depressa};
\draw[->,thick,blue!60!black] (5,3.25) -- (5,3.55);
\draw[->,thick,red!70!black] (4,1.8) -- (4,2.1);
\node[font=\small\itshape,gray!35!black,anchor=west] at (0.15,2.3) {busca às cegas};
\node[font=\small\itshape,gray!35!black] at (6.45,2.55) {não nota mudanças};
\end{tikzpicture}
\caption{Fração da melhor recompensa possível com ganho $g$ constante (escala logarítmica em $g$). Curva azul: o mundo muda em média uma vez a cada mil tentativas; vermelha, uma vez a cada vinte. As setas marcam o melhor $g$: num mundo que muda depressa, ele é duas vezes menor. Média de 60 execuções de 4000 tentativas.}
\label{fig:invertedU}
\end{figure}

Com $g=0.5$, a rede quase não aproveita o conhecimento e obtém cerca de três quartos do possível; com $g$ muito grande, perde as mudanças. O melhor valor: $g=16$ quando o mundo muda uma vez a cada mil tentativas (fração $0.976$) e $g=8$ quando muda uma vez a cada vinte ($0.886$).

\begin{keyblock}
\begin{proposition}[O U invertido]
\label{prop:explore}
A recompensa que a rede colhe depende do ganho como um U invertido: um $g$ pequeno demais gasta tentativas em busca, um grande demais não nota mudanças. O melhor $g$ é tanto menor quanto mais depressa muda o mundo.
\end{proposition}
\end{keyblock}

O U invertido também aparece nos experimentos: os animais resolvem a tarefa melhor com atividade de fundo moderada dos neurônios \nt{NORA} e respostas em rajada nítidas, e com atividade de fundo alta se distraem e alternam entre tarefas~\cite{AstonJones2005}. Isso concorda com a diferença de modos da seção anterior: a rajada no instante da decisão dá um $g$ efetivo grande exatamente quando é preciso escolher, e a atividade de fundo alta, sem rajadas, amplifica tudo, inclusive entradas alheias, de modo que o $g$ efetivo para a tarefa atual cai: é a busca.

\section{Quem gira o botão}

Já que o melhor ganho depende da rapidez com que o mundo muda, seria bom ajustá-lo. Mas como os neurônios \nt{NORA} saberiam que o mundo mudou? O indício natural é a \emph{surpresa}: os erros de predição ficaram maiores que o habitual. É importante distinguir dois tipos de incerteza. Se a recompensa é simplesmente aleatória, os erros são sempre grandes, e isso é esperado. Mas se os erros de repente cresceram em relação ao seu tamanho habitual, algo mudou. Segundo uma hipótese, é exatamente essa incerteza inesperada que o \nt{NORA} informa, e a esperada, o \nt{ACET}~\cite{YuDayan2005}.

O que fazer com esse sinal? Pode-se reduzir o ganho e buscar mais. Pode-se elevar a taxa de aprendizado, para esquecer mais depressa as estimativas obsoletas: experimentos mostram que depois de uma mudança brusca as pessoas aprendem mais depressa, e a pupila se dilata~\cite{Nassar2012}; lembremos que a pupila é indicador não só do \nt{NORA}, mas também do \nt{ACET}~\cite{Reimer2016}. E pode-se fazer as duas coisas de uma vez: segundo outra hipótese, a rajada do \nt{NORA} funciona como uma \emph{interrupção}, um sinal pelo qual a rede zera o estado atual e se reconfigura para a nova situação~\cite{BouretSara2005}, como a linha geral de interrupção de um processador.

Digamos com franqueza o que não encontramos. No modelo, testamos as duas formas simples de ajuste: reduzir $g$ ou elevar a taxa de aprendizado quando o erro suavizado passa da sua média de longo prazo. Num mundo que ora fica parado por muito tempo, ora muda com frequência, as melhores configurações das duas formas deram $0.926$--$0.927$ da melhor recompensa, quase o mesmo que o melhor $g$ constante ($0.925$ com $g=8$) ou uma taxa de aprendizado constante, mas grande o bastante ($0.928$), e as configurações ruins deram menos. As curvas da fig.~\ref{fig:invertedU} são achatadas perto do topo, e num mundo assim quase não há o que ajustar. O ajuste deve compensar onde regimes diferentes exigem configurações muito diferentes. Que lei de controle exatamente o \nt{NORA} implementa ainda não foi estabelecido.

\section{Resumo}

\begin{table}[t]
\centering
\footnotesize
\setlength{\tabcolsep}{4pt}
\renewcommand{\arraystretch}{1.15}
\begin{tabular}{>{\raggedright}p{3.0cm}>{\raggedright}p{4.4cm}>{\raggedright\arraybackslash}p{6.0cm}}
\toprule
O que o \nt{NORA} faz & Como & O que se sabe \\
\midrule
muda a inclinação da resposta dos neurônios & ganho $g$ em~\eqref{eq:gain} & aumento da relação sinal-ruído medido; o modelo multiplicativo é uma simplificação \\
comuta a escolha & leva a rede através de $g\beta=1$ & decorre do modelo~\eqref{eq:wtagain}; não há medições diretas do limiar \\
define quanto buscar & $g$ é o inverso da temperatura~\eqref{eq:softmax}; fundo: $g$ efetivo pequeno (busca); rajada: grande (decisão) & o U invertido e os dois modos de funcionamento foram medidos; a ligação com a busca é hipótese bem fundamentada \\
informa mudanças & incerteza inesperada & pupila e taxa de aprendizado crescem após mudanças: medido; que seja um sinal do \nt{NORA} é hipótese \\
interrompe e zera & rajada em resposta a evento inesperado & rajadas em eventos importantes medidas; a ``interrupção'' é hipótese \\
\bottomrule
\end{tabular}
\caption{O que o \nt{NORA} acrescenta à rede de \nt{GLUT}, \nt{GABA} e \nt{DOPA}.}
\label{tab:nora}
\end{table}

O \nt{DOPA} diz à rede \emph{para onde} mudar os pesos. O \nt{NORA} não muda nenhum peso, mas decide \emph{como} a rede usa o que já sabe (tabela~\ref{tab:nora}): um único botão, o ganho, transforma o neurônio de amplificador em comparador, a rede de escolha de ``ponderando'' em ``decidida'', e a escolha, de busca em aproveitamento. Como o \nt{NORA} descobre a rapidez com que o mundo muda é uma questão em aberto.

\section{Exercícios}

\begin{exercise}[\lvlA]
\label{ex:08:1}
\emph{Um botão para o cérebro inteiro.} (a)~Pela tabela~\ref{tab:types}, estime quantos neurônios do cérebro cabem a cada neurônio \nt{NORA}. (b)~Ruído antes do amplificador $\sigma_{\text{antes}}=0.5$, depois, $\sigma_{\text{depois}}=4$. Com que $g$ o $d'$ de~\eqref{eq:gainsnr} chega a $90\,\%$ do seu limite?
\end{exercise}

\begin{exercise}[\lvlB]
\label{ex:08:2}
\emph{Escolha com ganho.} $\tau\,\dd r_1/\dd t=-r_1+g[I_1-\beta r_2]_+$, $\tau\,\dd r_2/\dd t=-r_2+g[I_2-\beta r_1]_+$, $\tau=20\ms$. (a)~Em quanto tempo $r_1-r_2$ decai com $g\beta=0.5$ e $0.9$? (b)~Com $\varepsilon=0.05$, encontre o $g\beta$ abaixo do qual os dois grupos funcionam e acima do qual a vitória da entrada fraca é estável.
\end{exercise}

\begin{exercise}[\lvlA]
\label{ex:08:3}
\emph{Softmax a partir do ruído.} Cada um dos $K$ grupos recebe seu próprio ruído de Gumbel, $P(\eta_k<z)=\exp(-e^{-z/\sigma})$, e vence o maior $gu_k+\eta_k$. Encontre $P(k)$. Com que $g$ a melhor de duas opções com $u_A-u_B=0.1$, $\sigma=1$, é escolhida em $90\,\%$ dos casos?
\end{exercise}

\begin{exercise}[\lvlB]
\label{ex:08:4}
\emph{Estimativa do melhor ganho.} No modelo da fig.~\ref{fig:invertedU}, as probabilidades de recompensa após uma mudança são uniformes em $[0,1]$. A rede experimenta a pior ação uma vez a cada $e^{g\Delta}$ tentativas; igualando isso a $1/h$, obtemos $g^*\approx\ln(1/h)/\bar\Delta$. Encontre $\bar\Delta=\langle|p_A-p_B|\rangle$ e $g^*$ com $h=0.001$, $0.01$, $0.05$.
\end{exercise}

\begin{exercise}[\lvlB]
\label{ex:08:5}
\emph{Simulação: ganho e rapidez das mudanças.} Encontre $g^*(h)$ para $h$ de $0.001$ a $0.2$ com uma cópia de \texttt{models/\allowbreak nora\_explore.py} (ela grava arquivos na pasta atual). Onde vale a estimativa do exercício~\ref{ex:08:4}, e por que, com mudanças rápidas, $g^*$ deixa de diminuir?
\end{exercise}

\begin{exercise}[\lvlB]
\label{ex:08:6}
\emph{Projeto: interrupção para a rede de escolha.} A rede do exercício~\ref{ex:08:2} com $\beta=2$, $\tau=20\ms$, $g=1$ mantém $r_1=0.9$, $r_2=0$, embora as entradas agora sejam $I_1=0.9$, $I_2=1.0$. O \nt{NORA} reduz $g$ a $g_{\text{lo}}$ por um tempo $T_p$. Encontre o menor $T_p$ com $g_{\text{lo}}=0$ analiticamente e, com $g_{\text{lo}}=0.25$, por simulação.
\end{exercise}

\begin{exercise}[\lvlC]
\label{ex:08:7}
\emph{Quando o ajuste compensa.} Um oráculo sabe se a época é calma ou instável e usa em cada uma seu próprio $g$. (a)~Por simulação, encontre seu ganho sobre o melhor $g$ constante no mundo do capítulo (épocas de $1000$ tentativas, $h=0.001$ e $0.05$). (b)~Construa um mundo em que o ganho seja maior e encontre que parte dele o ajuste pela surpresa captura.
\end{exercise}

\chapter{Quando escrever e quando ler: acrescentamos o \nt{ACET}}
\chaptermark{Acrescentamos o \nt{ACET}}
\label{sec:acet}

\section{O que falta}

Um chip de memória, além dos fios de endereço e de dados, tem mais uma linha: a habilitação de escrita. Enquanto ela está desligada, a memória só é lida; quando está ligada, o que estiver nas linhas de dados é gravado. Sem essa linha a memória não funciona: se cada leitura grava algo ao mesmo tempo, o que foi lido, inclusive o que foi lido com erro, é imediatamente gravado de volta.

No cérebro, esse problema é mais agudo que no chip. No capítulo~\ref{sec:glutcomputer}, a memória era um grupo de neurônios \nt{GLUT} mantido ligado por suas próprias conexões (fig.~\ref{fig:bistable}). No capítulo~\ref{sec:dofa}, vimos que essas conexões aprendem pela regra de Hebb durante o próprio funcionamento. Logo, as mesmas sinapses guardam o antigo e recebem o novo, e não há uma fase separada de escrita, assim como não há relógio. Como a rede distingue o momento em que é preciso memorizar o que se vê do momento em que é preciso lembrar o que se sabe? No capítulo~\ref{sec:neurontypes}, supusemos que essa linha é mantida pelo \nt{ACET}. Neste capítulo, vamos verificar como essa linha deve funcionar e por que não se pode passar sem ela.

\section{Três ações de um só fio}

Os neurônios colinérgicos que atuam dentro do cérebro são poucos e estão reunidos em alguns grupos pequenos, e suas saídas se espalham por todo o córtex. É um canal de difusão, como o \nt{DOPA} e o \nt{NORA}. O nível de acetilcolina no córtex é alto na vigília atenta e baixo no sono de ondas lentas~\cite{Hasselmo1999}. Como na dopamina, o sentido do sinal é definido pelo receptor do destinatário (proposição~\ref{prop:receptor}): a acetilcolina tem tanto receptores rápidos, que abrem um canal, quanto lentos, que disparam reações dentro da célula. Para nós, importam três ações.

\emph{Primeira: enfraquecer as conexões internas.} Experimentos em fatias do córtex olfativo mostraram que a acetilcolina suprime fortemente a transmissão pelas conexões entre neurônios de uma mesma região e quase não afeta as entradas que chegam de fora~\cite{HasselmoBower1992}.

\emph{Segunda: reforçar as entradas.} A acetilcolina aumenta a excitabilidade dos neurônios e facilita a transmissão por algumas vias de entrada; o sinal dos órgãos dos sentidos fica mais alto que a atividade própria da rede~\cite{Hasselmo2006}.

\emph{Terceira: facilitar a mudança dos pesos.} Com nível alto de acetilcolina, as conexões mudam com mais facilidade~\cite{Hasselmo2006}.

Para o engenheiro, as três ações são uma única operação: passar do modo ``ler'' para o modo ``escrever''. A rede deixa de ouvir a si mesma, passa a ouvir a entrada e memoriza o que ouve.

\section{Uma rede que completa}

Tomemos $N$ neurônios \nt{GLUT} ligados entre si. Guardemos em suas conexões alguns padrões (conjuntos de $pN$ neurônios ativos ao mesmo tempo) pela regra de Hebb~\cite{Hebb1949}. Se aplicarmos a essa rede metade de um padrão, as conexões excitarão a metade que falta: a rede \emph{completa} o padrão a partir de uma parte, assim como o grupo da fig.~\ref{fig:bistable} se sustentava sozinho. É uma memória associativa: o endereço é uma parte do conteúdo. O \nt{GABA}, como no capítulo~\ref{sec:gaba}, mantém o número total de neurônios ativos em torno de $pN$, para que dois padrões não possam acender juntos.

Denotemos o nível de acetilcolina por $a$, de $0$ a $1$, e escrevamos suas três ações diretamente no modelo:
\begin{empheq}[box=\keybox]{equation}
\tau\frac{\dd r_i}{\dd t} = -r_i + F\Bigl(\tfrac{1+a}{2}\,x_i + (1-a)\sum_j W_{ij}\,r_j - g\Bigr),
\qquad
\frac{\dd W_{ij}}{\dd t} = \eta\,a\,(r_i-p)(r_j-p) .
\label{eq:acetnet}
\end{empheq}
Aqui $r_i$ é a taxa do neurônio $i$, $x_i$ é sua entrada vinda dos órgãos dos sentidos, $F$ é uma característica de transferência com limiar, e $g$ é a inibição geral do \nt{GABA}, que cresce quando há mais de $pN$ neurônios ativos. O fator $\frac{1+a}{2}$ é o reforço da entrada, $1-a$ é o enfraquecimento das conexões internas, e $\eta a$ é a taxa de aprendizado. A segunda equação é a regra de Hebb~\eqref{eq:hebb} numa forma conveniente para padrões esparsos~\cite{Tsodyks1988}: o peso cresce quando os dois neurônios estão mais ativos que o habitual e diminui quando só um está ativo. A parte negativa dos pesos, como sempre, fica a cargo de intermediários \nt{GABA}. Não há relógio: tanto os neurônios quanto os pesos mudam continuamente.

\begin{figure}[t]
\centering
\begin{tikzpicture}[>=Stealth,
  nrn/.style={circle,draw,very thick,fill=blue!6,minimum size=0.8cm,inner sep=0pt,font=\small},
  acet/.style={circle,draw,very thick,double,double distance=1.2pt,fill=orange!15,minimum size=1.35cm,inner sep=0pt,font=\scriptsize},
  bc/.style={->,thick,densely dashed,orange!85!black}]
\foreach \i/\y in {1/2.4,2/1.2,3/0}{
  \node[nrn] (N\i) at (3,\y) {$r_\i$};
  \draw[->,thick,blue!60!black] (0,\y) node[left,black] {\small $x_\i$} -- (N\i);}
\node[left,align=right,font=\small] at (-0.5,1.2) {dos órgãos\\dos sentidos};
\draw[->,thick,gray!60!black] (N1) to[bend left=30] (N2);
\draw[->,thick,gray!60!black] (N2) to[bend left=30] (N1);
\draw[->,thick,gray!60!black] (N2) to[bend left=30] (N3);
\draw[->,thick,gray!60!black] (N3) to[bend left=30] (N2);
\draw[->,thick,gray!60!black] (N1.east) .. controls (4.6,2.0) and (4.6,0.4) .. (N3.east);
\node[right,font=\small,gray!40!black,align=left] at (4.7,1.2) {conexões\\internas $W$};
\node[acet] (A) at (1.2,-1.9) {\nt{ACET}};
\draw[bc] (A) -- (1.6,-0.15);
\node[font=\small,orange!60!black,anchor=east] at (0.95,-0.9) {$+$ entrada};
\draw[bc] (A) -- (4.3,0.6);
\node[font=\small,orange!60!black,anchor=west] at (4.3,0.1) {$-$ conexões internas, $+$ taxa de aprendizado};
\node[right,font=\small,align=left] at (2.2,-1.9) {nível alto: ``escrever''\\nível baixo: ``ler''};
\end{tikzpicture}
\caption{O \nt{ACET} como linha de habilitação de escrita. Os neurônios $r_i$ recebem a entrada $x_i$ dos órgãos dos sentidos (setas azuis) e se excitam mutuamente pelas conexões internas $W$ (cinza), nas quais os padrões ficam guardados. A acetilcolina (setas tracejadas) reforça a entrada, enfraquece as conexões internas e facilita a mudança dos pesos, como em~\eqref{eq:acetnet}.}
\label{fig:acetnet}
\end{figure}

\section{Escrita e leitura atrapalham uma à outra}

Vamos testar o modelo~\eqref{eq:acetnet} num problema em que escrita e leitura realmente entram em conflito. Uma rede de $200$ neurônios já guarda cinco padrões de $20$ neurônios ativos. Agora é preciso memorizar um sexto, que coincide pela metade com um dos antigos: dez neurônios são comuns aos dois. Isso acontece o tempo todo: um quarto novo parece um conhecido, um rosto novo parece um antigo.

\emph{Escrita com $a$ baixo.} Primeiro, separemos as duas primeiras ações da acetilcolina da terceira: deixemos a taxa de aprendizado plena, e só o reforço da entrada e o enfraquecimento das conexões internas vão variar. Com $a=0$, a entrada acende os dez neurônios comuns, e eles, pelas conexões internas, acendem a metade restante do padrão \emph{antigo}. O \nt{GABA} não deixa os dois acesos ao mesmo tempo, e o padrão antigo, sustentado por suas próprias conexões, expulsa o novo, que só se apoia na entrada. A rede não vê o que está diante dela, mas o que lembra, e a regra de Hebb grava diligentemente o que foi visto.

No fim, o padrão novo não foi memorizado de modo algum: a partir de sua metade distintiva, a rede não lembra nada. E o antigo ficou corrompido: evocado por sua metade distintiva, ele arrasta em média seis neurônios alheios de dez. Com $a=0.1$, o padrão novo já é memorizado, mas se funde com o antigo, e a corrupção é até maior: cada um dos dois, evocado por sua metade, arrasta cerca de oito neurônios alheios; a partir de $a=0.2$ a escrita é limpa: menos de um neurônio a mais na lembrança (média de dez execuções).

\emph{Escrita com a taxa de aprendizado real.} Agora deixemos a acetilcolina, como em~\eqref{eq:acetnet}, definir também a taxa de aprendizado. Numa única apresentação, o padrão novo é memorizado por inteiro só com $a\ge0.9$; com $a=0.75$, em oito décimos; com $a\le0.5$, nada.

\emph{Leitura.} Aplicamos à rede com cinco padrões escritos de forma limpa metade de um padrão e depois a retiramos. Com $a\le0.2$, a rede completa o padrão inteiro e o mantém sozinha; com $a=0.3$, quase inteiro; com $a\ge0.4$, as conexões internas estão enfraquecidas demais e, depois que a dica some, a atividade se apaga (fig.~\ref{fig:acetsim}).

\begin{figure}[t]
\centering
\begin{tikzpicture}[>=Stealth,x=9cm,y=3.2cm]
\draw[->,gray!70!black] (0,0) -- (1.1,0) node[right,black] {\small $a$};
\draw[->,gray!70!black] (0,0) -- (0,1.2);
\foreach \x in {0,0.2,0.4,0.6,0.8,1.0}{\draw[gray!60!black] (\x,-0.02) -- (\x,0.02); \node[below,font=\scriptsize] at (\x,0) {$\x$};}
\foreach \y in {0,0.5,1}{\draw[gray!60!black] (-0.01,\y) -- (0.01,\y); \node[left,font=\scriptsize] at (0,\y) {$\y$};}
\node[rotate=90,font=\small] at (-0.1,0.6) {fração do padrão recuperada};
\draw[very thick,blue!60!black] plot[mark=*,mark size=1.6pt] coordinates {(0,1.00) (0.1,1.00) (0.2,1.00) (0.3,0.96) (0.4,0.04) (0.5,0.00) (0.75,0.00) (1.0,0.00)};
\draw[very thick,red!70!black] plot[mark=*,mark size=1.6pt] coordinates {(0.1,0.00) (0.2,0.00) (0.3,0.00) (0.5,0.00) (0.75,0.80) (0.9,1.00) (1.0,1.00)};
\node[font=\small,blue!60!black,anchor=west,align=left] at (0.03,0.5) {leitura:\\padrão\\pela metade};
\node[font=\small,red!70!black,anchor=east,align=right] at (1.0,0.35) {escrita:\\padrão novo\\de uma vez};
\end{tikzpicture}
\caption{Modelo~\eqref{eq:acetnet}: $200$ neurônios, padrões de $20$ ativos, média de dez execuções. Pontos azuis, leitura: que fração do padrão é recuperada a partir de sua metade depois que a dica é retirada. Pontos vermelhos, escrita: que fração de um padrão novo, parecido pela metade com um antigo, é lembrada após uma apresentação, se a acetilcolina define também a taxa de aprendizado. A leitura funciona com $a\le0.3$, a escrita com $a\ge0.9$.}
\label{fig:acetsim}
\end{figure}

\begin{keyblock}
\begin{proposition}[Escrita e leitura exigem modos diferentes]
\label{prop:writeread}
No modelo~\eqref{eq:acetnet}, em que as mesmas conexões guardam o antigo e aprendem o novo e a acetilcolina define também a taxa de aprendizado, não há nível $a$ em que escrita e leitura funcionem igualmente bem. Para escrever, é preciso enfraquecer as conexões internas, senão a rede grava não a entrada, mas o que lembrou; para ler, é preciso reforçá-las, senão ela não completa o padrão a partir de uma parte. É necessário um comutador, e um único fio de difusão pode servir como tal.
\end{proposition}
\end{keyblock}

Se a acetilcolina não mudasse a taxa de aprendizado, uma faixa estreita $0.2\le a\le0.3$ serviria para as duas tarefas. Mas então a rede aprenderia também a cada leitura, isto é, regravaria o que leu junto com seus erros. A própria ideia de suprimir as conexões internas durante o aprendizado vem de modelos construídos a partir de medições em fatias~\cite{HasselmoAndersonBower1992}. Os números acima se referem ao nosso modelo simplificado; onde exatamente ficam as fronteiras dos modos no cérebro não se sabe.

\section{Quanto acreditar nos olhos}

O comutador ``escrever/ler'' é um caso extremo de uma questão mais geral: quanto acreditar na entrada e quanto na memória? Para essa questão os engenheiros têm uma resposta exata, e ela ajuda a entender por que um único fio controla ao mesmo tempo o modo e a taxa de aprendizado.

Suponha que a rede acompanha uma grandeza $s(t)$ que vagueia lentamente: a cada segundo sua variância cresce $q$. Os órgãos dos sentidos informam $y(t)=s(t)+$ruído com intensidade $R$. A melhor estimativa $\hat s$ em tempo contínuo é dada pelo filtro de Kalman--Bucy~\cite{KalmanBucy1961}:
\begin{empheq}[box=\keybox]{equation}
\frac{\dd\hat s}{\dd t} \;=\; \frac{P}{R}\,\bigl(y-\hat s\bigr),
\qquad
\frac{\dd P}{\dd t} \;=\; q-\frac{P^2}{R} ,
\label{eq:kalman}
\end{empheq}
em que $P$ é a variância do erro da estimativa, isto é, o quanto a própria rede está insegura do que sabe. A primeira equação é o mesmo circuito RC da membrana no modelo de neurônio~\eqref{eq:glutneuron}, mas com uma constante de tempo $R/P$ que varia. Quando a rede está insegura, $P$ é grande, a constante de tempo é pequena, e a estimativa acompanha depressa a entrada: é ``escrever''. Quando está segura, $P$ é pequeno, e a estimativa quase não ouve a entrada, apegando-se à memória: é ``ler''.

Em regime estacionário, $P=\sqrt{qR}$, e a constante de tempo da memória é $\sqrt{R/q}$: se o mundo se desloca uma unidade em cem segundos, $q=0.01$, e o ruído dos órgãos dos sentidos é $R=1$, a memória deve lembrar cerca de dez segundos.

O principal aqui é que um único número, $P/R$, define duas coisas ao mesmo tempo: o peso da entrada em relação à memória e a velocidade com que a estimativa guardada muda. É exatamente o que a acetilcolina faz em~\eqref{eq:acetnet}. Pela hipótese~\cite{YuDayan2005}, a acetilcolina informa à rede uma grandeza semelhante a $P$: a incerteza \emph{esperada}, aquela que a rede conhece de antemão. Esse canal nem sempre é lento: parte dos neurônios \nt{ACET} responde à recompensa e à punição em duas dezenas de milissegundos, tanto mais forte quanto mais inesperado o reforço~\cite{Hangya2015}.

\begin{keyblock}
\begin{proposition}[Um único botão para o peso da entrada e a taxa de aprendizado]
\label{prop:kalman}
No filtro ótimo~\eqref{eq:kalman}, o peso com que a entrada se mistura à memória e a taxa de aprendizado são a mesma grandeza $P/R$. Por isso é razoável controlá-los com um único sinal. Com propriedades do mundo inalteradas, esse sinal se estabiliza num nível constante $\sqrt{q/R}$, e só é preciso ajustá-lo quando as propriedades do mundo mudam.
\end{proposition}
\end{keyblock}

A segunda frase da proposição explica o resultado do capítulo~\ref{sec:nora}, em que o ajuste da taxa de aprendizado pela surpresa não superou a melhor taxa constante: num mundo com propriedades de mudança inalteradas, a melhor taxa de aprendizado é justamente constante. O ganho aparece quando o mundo de repente fica diferente. Então a estimativa se vê subitamente longe da entrada, e $P$ não sabe disso. A ação correta é elevar $P$ bruscamente e, com isso, passar ao modo de escrita. Pela hipótese que discutimos no capítulo~\ref{sec:nora}, quem informa essa mudança \emph{inesperada} é o \nt{NORA}~\cite{YuDayan2005}. A divisão de trabalho resulta natural: o \nt{NORA} percebe que o mundo mudou, e o \nt{ACET} mantém a rede no modo de escrita até que a nova situação seja aprendida.

\section{O sono como modo de leitura}

Se o nível alto de acetilcolina é o modo de escrita, o que a rede faz com nível baixo? No sono de ondas lentas, o nível de acetilcolina cai, as conexões internas se livram da supressão, e quase não há entrada dos órgãos dos sentidos. A rede em modo de leitura, sem dica, reproduz o que está gravado nela. O registro da atividade mostra que neurônios que funcionavam juntos durante o dia voltam a disparar juntos no sono de ondas lentas~\cite{WilsonMcNaughton1994}, e até na mesma ordem~\cite{Skaggs1996}. Pela hipótese~\cite{Hasselmo1999}, essas repetições transcrevem o que foi aprendido depressa durante o dia para a memória de longo prazo (prolongar artificialmente os surtos curtos de repetição melhora a memória~\cite{FernandezRuiz2019}): de dia, escrita a partir da entrada; à noite, regravação a partir de dentro. Para o engenheiro, isso lembra gravar num buffer rápido durante o dia e descarregá-lo numa memória permanente lenta à noite.

\section{Resumo}

\begin{table}[t]
\centering
\footnotesize
\setlength{\tabcolsep}{4pt}
\renewcommand{\arraystretch}{1.15}
\begin{tabular}{>{\raggedright}p{3.2cm}>{\raggedright}p{4.2cm}>{\raggedright\arraybackslash}p{6.0cm}}
\toprule
O que o \nt{ACET} faz & Como & O que se sabe \\
\midrule
enfraquece as conexões internas & fator $1-a$ em~\eqref{eq:acetnet} & medido em fatias \\
reforça a entrada & fator $\frac{1+a}{2}$ & aumento da excitabilidade e da transmissão pelas vias de entrada medido \\
facilita o aprendizado & taxa $\eta a$ & medido \\
comuta ``escrever/ler'' & proposição~\ref{prop:writeread} & decorre dos modelos; não há medição direta dos modos \\
informa a incerteza esperada & $P$ em~\eqref{eq:kalman} & hipótese \\
libera a rede para repetições no sono & $a$ baixo no sono de ondas lentas & nível baixo no sono e repetições medidos; a regravação na memória de longo prazo é hipótese \\
\bottomrule
\end{tabular}
\caption{O que o \nt{ACET} acrescenta à rede de \nt{GLUT}, \nt{GABA}, \nt{DOPA} e \nt{NORA}.}
\label{tab:acet}
\end{table}

O \nt{DOPA} diz à rede para onde mudar os pesos; o \nt{NORA}, com que decisão usar o que ela sabe; e o \nt{ACET}, se ela deve ouvir o mundo ou a si mesma (tabela~\ref{tab:acet}). É a última linha de controle por difusão da tabela~\ref{tab:types}: a habilitação de escrita, sem a qual uma memória que guarda padrões nas mesmas conexões com que os lê se corromperia a cada leitura.

\section{Exercícios}

\begin{exercise}[\lvlA]
\label{ex:09:1}
\emph{Quanto lembrar.} O filtro~\eqref{eq:kalman} com $q=0.01$, $R=1$ lembra cerca de $10$~s. Encontre $P_\infty$ e a memória $\sqrt{R/q}$ se (a)~o ruído do sensor dobrou em amplitude; (b)~o mundo passou a vaguear cem vezes mais depressa. (c)~Quantas vezes o ruído deve crescer para a rede lembrar um minuto?
\end{exercise}

\begin{exercise}[\lvlB]
\label{ex:09:2}
\emph{Modo de escrita após uma mudança.} O mundo mudou, e a incerteza do filtro~\eqref{eq:kalman} com $q=0.01$, $R=1$ saltou para $P_0\to\infty$. (a)~Com que constante de tempo $P$ volta a $P_\infty$? (b)~Depois de quantos segundos a taxa de aprendizado passa a exceder a estacionária em no máximo $10\,\%$?
\end{exercise}

\begin{exercise}[\lvlB]
\label{ex:09:3}
\emph{Taxa de aprendizado constante.} A rede aprende como $\dd\hat s/\dd t=(y-\hat s)/T$; o mundo vagueia, $\dd s/\dd t=w$, $y=s+n$, em que $w$ e $n$ são ruídos brancos de intensidades $q$ e $R$. Encontre o melhor $T$ e a variância do erro com ele. Quantas vezes ela cresce se $T$ estiver errado por um fator $2$ e $10$?
\end{exercise}

\begin{exercise}[\lvlB]
\label{ex:09:4}
\emph{Onde fica a fronteira da escrita.} Em \texttt{models/\allowbreak acet\_memory.py}, o limiar é $\theta=0.35$, e o peso dentro de um padrão, $w=0.041$ (idealmente $0.045$). Um padrão novo compartilha $m=10$ neurônios com um antigo. Com que $a$ em~\eqref{eq:acetnet} os demais neurônios do padrão antigo não se acendem a partir desses $m$ (limiar abrupto)?
\end{exercise}

\begin{exercise}[\lvlB]
\label{ex:09:5}
\emph{Simulação: capacidade da memória.} Modifique em \texttt{models/\allowbreak acet\_memory.py} a função \texttt{read\_sweep}: com $a=1$, grave $M$ padrões ($M$ de $5$ a $100$) e depois, com $a=0$, leia os dez primeiros pela metade. Com que $M$ surgem erros, e como o $M_{\max}$ limite depende de $N$ e $p$?
\end{exercise}

\begin{exercise}[\lvlB]
\label{ex:09:6}
\emph{Projeto: escrita sem vazamento.} Com taxa de aprendizado $\eta a$, a rede aprende também durante a leitura. Proponha uma $\eta(a)\le\eta$ monótona, nula para $a\le0.3$ e não pior que $\eta a$ para $a\ge0.9$. Verifique: após $20$ leituras com $a=0.2$ e uma dica com três neurônios alheios, quantos deles entraram no padrão?
\end{exercise}

\begin{exercise}[\lvlC]
\label{ex:09:7}
\emph{Como regravar o buffer à noite.} (a)~No sono, $a=0.05$, e uma repetição dura $0.1\,\text{s}$, contra $0.2\,\text{s}$ de dia com $a=1$. Quantas repetições acumulam a mudança de pesos de uma apresentação diurna com taxa $\eta a$ e com a $\eta(a)$ do exercício~\ref{ex:09:6}? (b)~O armazenamento aprende $100$ vezes mais devagar que o buffer e ouve o padrão $20$ vezes por noite, $0.1\,\text{s}$ de cada vez. Quantas noites são necessárias?
\end{exercise}

\chapter{A máquina inteira}
\chaptermark{A máquina inteira}
\label{sec:synthesis}

\section{O que montamos}

Pegamos um tipo de neurônio de cada vez da tabela~\ref{tab:types} e perguntamos o que ele acrescenta à máquina de calcular. As regras foram sempre as mesmas: só o que existe num organismo vivo, e nenhum relógio comum. Agora vamos juntar as peças numa só máquina e ver como elas funcionam em conjunto.

O quadro do capítulo~\ref{sec:neurontypes} se confirmou e ficou mais preciso. Uma enorme rede endereçada de neurônios \nt{GLUT} transporta os dados. Os neurônios \nt{GABA} controlam o fluxo desses dados. E sobre a rede pairam quatro canais de difusão, cada um com o seu botão: \nt{DOPA} diz quais mudanças de peso manter, \nt{SERO} mantém o nível geral e, por hipótese, o horizonte de planejamento, \nt{NORA} escolhe entre explorar e decidir, \nt{ACET} escolhe entre escrita e leitura (fig.~\ref{fig:machine}).

\begin{figure}[t]
\centering
\begin{tikzpicture}[>=Stealth,
  blk/.style={draw,very thick,rounded corners=3pt,align=center,font=\small},
  reg/.style={draw,very thick,double,double distance=1.2pt,circle,minimum size=1.25cm,inner sep=0pt,font=\scriptsize,fill=orange!15},
  bc/.style={->,thick,densely dashed,orange!85!black}]
% sensors, bus, actions
\node[blk,fill=green!8,text width=1.7cm] (S) at (0.9,0) {sensores\\\scriptsize liga/desliga};
\draw[very thick,rounded corners=4pt,fill=blue!5] (2.6,-1.35) rectangle (9.0,1.35);
\node[font=\small] at (5.8,0.95) {barramento \nt{GLUT}: dados};
\node[font=\scriptsize,align=center] at (4.3,0.25) {coincidências, atrasos,\\lógica (cap.~\ref{sec:glutcomputer})};
\node[font=\scriptsize,align=center] at (7.4,0.25) {memória, código\\(cap.~\ref{sec:code}, \ref{sec:acet})};
\draw[thick,red!70!black,fill=red!8,rounded corners=2pt] (2.8,-1.15) rectangle (8.8,-0.55);
\node[font=\scriptsize,red!60!black] at (5.8,-0.85) {\nt{GABA} (cap.~\ref{sec:gaba}): veto, escolha, divisão, ritmo};
\node[blk,fill=blue!5,text width=2.0cm] (A) at (10.9,0) {escolha\\da ação};
\draw[->,very thick] (S) -- (2.6,0);
\draw[->,very thick] (9.0,0) -- (A);
\draw[->,very thick] (A) -- (12.6,0) node[right,font=\small] {músculos};
\pic[scale=1.1] at (13.2,-0.5) {spring};
\pic[scale=1.15] at (0.4,0.85) {eyepic};
\pic[scale=1.05] at (1.35,0.85) {speaker};
% registers
\node[reg] (D) at (3.4,3.2) {\nt{DOPA}};
\node[reg] (C) at (5.6,3.2) {\nt{SERO}};
\node[reg] (N) at (7.8,3.2) {\nt{NORA}};
\node[reg] (H) at (10.0,3.2) {\nt{ACET}};
\node[font=\scriptsize,align=center,above=1pt] at (D.north) {erro $\delta$};
\node[font=\scriptsize,align=center,above=1pt] at (C.north) {nível, $\tau_\gamma$};
\node[font=\scriptsize,align=center,above=1pt] at (N.north) {ganho $g$};
\node[font=\scriptsize,align=center,above=1pt] at (H.north) {escrita $a$};
\draw[bc] (D.south) -- (3.4,1.35);
\draw[bc] (C.south) -- (5.6,1.35);
\draw[bc] (N.south) -- (7.8,1.35);
\draw[bc] (H.south) -- (10.0,2.1) -- (8.8,1.35);
\draw[bc] (H.south east) to[bend left=20] (A.north east);
\draw[->,thick] (2.85,1.35) -- (2.85,2.75) -- (D.west) node[pos=0.25,left,font=\scriptsize] {$V$};
\draw[->,thick,green!45!black] (0.4,3.2) node[left,black,font=\footnotesize] {recompensa $r$} -- (2.2,3.2) -- (D.west);
\pic[scale=0.9] at (-0.45,3.7) {juice};
\draw[->,thick,densely dotted,gray!50!black] ([yshift=0.45cm]D.north) to[bend left=18] node[above,font=\scriptsize,black] {surpresa} ([yshift=0.45cm]N.north);
\end{tikzpicture}
\caption{A máquina inteira. O barramento \nt{GLUT} leva os dados dos sensores à escolha da ação; \nt{GABA} controla o fluxo dentro dele. Os círculos duplos são os canais de difusão; as setas tracejadas, a difusão dos seus sinais para toda a rede, cada um com o seu botão. \nt{DOPA} recebe a recompensa $r$ e a expectativa $V$ do crítico dentro do barramento (cap.~\ref{sec:dofa}); a seta pontilhada é a hipótese de que \nt{NORA} fica sabendo da surpresa por erros anormalmente grandes (cap.~\ref{sec:nora}).}
\label{fig:machine}
\end{figure}

\section{Uma fórmula para todos os botões}

Os botões podem ser reunidos numa só fórmula esquemática. Não é um modelo novo, e sim um resumo dos modelos dos capítulos~\ref{sec:glutcomputer}--\ref{sec:acet}: cada símbolo vem do seu capítulo.
\begin{empheq}[box=\keybox]{equation}
\begin{aligned}
\tau\,\frac{\dd r_i}{\dd t} \;&=\; -r_i + F\Bigl(g\,\Bigl[\tfrac{1+a}{2}\,x_i + (1-a)\sum_j W_{ij}\,r_j - \sum_k M_{ik}\,q_k - \theta\Bigr]\Bigr),\\
\frac{\dd W_{ij}}{\dd t} \;&=\; \eta\,a\,\delta(t)\,e_{ij}(t),
\qquad \tau_e\,\frac{\dd e_{ij}}{\dd t} = -e_{ij} + r_i\,r_j,
\qquad \delta = r + \frac{\dd V}{\dd t} - \frac{V}{\tau_\gamma} .
\end{aligned}
\label{eq:machine}
\end{empheq}
Aqui $r_i$ são as taxas de disparo dos neurônios \nt{GLUT}, $x_i$ é a entrada vinda dos sensores, $W_{ij}\ge0$ são os pesos entre eles (capítulo~\ref{sec:glutcomputer}); $q_k$ são as taxas dos neurônios \nt{GABA}, $M_{ik}\ge0$ é a força da sua inibição (capítulo~\ref{sec:gaba}); $e_{ij}$ é o traço de elegibilidade, $\delta$ é o erro de \nt{DOPA} (capítulo~\ref{sec:dofa}); $\tau_\gamma$ é o horizonte, por hipótese definido por \nt{SERO}; $g$ é o ganho de \nt{NORA} (capítulo~\ref{sec:nora}); $a$ é o nível de \nt{ACET} (capítulo~\ref{sec:acet}).

Veja o que não está em~\eqref{eq:machine}. Não há ciclo de relógio: tudo está escrito em equações diferenciais, e cada neurônio muda por conta própria. Não há memória separada: quem lembra são os mesmos $W_{ij}$ pelos quais passa o cálculo, e a mesma atividade $r_i$. Não há correções individuais para a imensa maioria das sinapses: o seu aprendizado é o produto de um traço local por um único número comum; um fio de erro próprio para cada saída só existe no circuito que ensina os movimentos (capítulo~\ref{sec:motorlearning}). E há só quatro tipos de sinais de controle, $\delta$, $\tau_\gamma$, $g$, $a$, para oitenta bilhões de neurônios. Na verdade cada um deles não é um número só, mas um pequeno feixe de linhas paralelas (proposição~\ref{prop:bundle}); só que as linhas num feixe são poucas.

\begin{keyblock}
\begin{proposition}[Arquitetura]
\label{prop:architecture}
A máquina do cérebro, até onde a entendemos hoje, pode ser descrita em três camadas. Os dados correm pelo barramento de endereços \nt{GLUT}. O fluxo de dados é dirigido, limitado e normalizado pelos neurônios \nt{GABA} endereçados. O modo de funcionamento é definido por alguns canais de difusão, cada um sendo um pequeno feixe de linhas paralelas, comuns a grandes regiões da rede. As duas primeiras camadas estão firmemente estabelecidas; a divisão de papéis entre os canais de difusão é, em grande parte, hipótese.
\end{proposition}
\end{keyblock}

\section{Todos os tipos numa tabela}

A tabela~\ref{tab:synthesis} reúne todos os tipos de neurônios do livro. Os quatro tipos que não ganharam capítulo próprio ocupam nela uma linha cada; dois deles encontraram um papel nos capítulos sobre as saídas da máquina: \nt{GLYC} nos circuitos da medula espinhal (capítulo~\ref{sec:muscles}) e \nt{ADRE} na saída química (capítulo~\ref{sec:chemoutput}). O papel dos outros dois no cálculo é simples ou desconhecido. As substâncias que não definem um tipo de neurônio estão reunidas no apêndice~\ref{sec:othermediators}.

\begin{table}[t]
\centering
\footnotesize
\setlength{\tabcolsep}{3pt}
\renewcommand{\arraystretch}{1.15}
\begin{tabular}{>{\raggedright}p{1.3cm}>{\raggedright}p{1.0cm}>{\raggedright}p{3.9cm}>{\raggedright}p{1.5cm}>{\raggedright}p{1.8cm}>{\raggedright\arraybackslash}p{3.8cm}}
\toprule
Tipo & Cap. & O que calcula & Tempo & Barramento & O que se sabe \\
\midrule
\nt{GLUT} & \ref{sec:glutcomputer}, \ref{sec:code}, \ref{sec:sensors}, \ref{sec:motorlearning}, \ref{sec:dendrites} & dados: coincidências, atrasos, lógica, memória, sequências & ms & de endereços & estabelecido \\
\nt{GABA} & \ref{sec:gaba}, \ref{sec:code}, \ref{sec:pain}, \ref{sec:motorlearning} & controle de fluxo: veto, escolha, divisão, estabilidade, ritmo; porta da dor & ms & de endereços & estabelecido; divisão da taxa só junto com o ruído de fundo \\
\nt{SERO} & \ref{sec:serocomputer}, \ref{sec:pain}, \ref{sec:sleep} & regulador de nível, suavização; horizonte $\tau_\gamma$; volume da dor; modos do sono & segundos & volume & autoinibição medida; horizonte é hipótese \\
\nt{DOPA} & \ref{sec:dofa}, \ref{sec:dofaalt} & erro de predição $\delta$; nível lento é o preço do tempo & $0.1$\,s e minutos & difusão & $\delta$ medido; extensões no cap.~\ref{sec:dofaalt} \\
\nt{NORA} & \ref{sec:nora}, \ref{sec:pain}, \ref{sec:chemoutput}, \ref{sec:sleep} & ganho $g$: explorar ou decidir; surpresa; ``acelerador'' dos órgãos & segundos & difusão & aumento da relação sinal-ruído medido; papel na exploração é hipótese \\
\nt{ACET} & \ref{sec:acet}, \ref{sec:muscles}, \ref{sec:chemoutput}, \ref{sec:sleep} & escrita ou leitura; taxa de aprendizado; saída para o músculo; ``freio'' dos órgãos & segundos & ambos & três efeitos medidos; troca de modos é hipótese \\
\nt{GLYC} & \ref{sec:muscles}, \ref{sec:pain} & peso negativo na medula e no tronco: veto do músculo antagonista & ms & de endereços & estabelecido \\
\nt{HIST} & --- & sinal global ``fique acordado'' & lento & difusão & papel no cálculo não estudado \\
\nt{ADRE} & \ref{sec:chemoutput} & ``acelerador'' do corpo todo pelo sangue; no cérebro, desconhecido & lento & difusão & estabelecido fora do cérebro; papel no cérebro incerto \\
\nt{ASPA} & --- & peso positivo fraco & ms & de endereços & papel controverso \\
\bottomrule
\end{tabular}
\caption{Todos os tipos de neurônios do livro: o que cada um calcula e quão bem isso é conhecido.}
\label{tab:synthesis}
\end{table}

\section{Como os botões trabalham juntos}

Até aqui cada capítulo girou um botão. Vamos agora acompanhar um evento através da máquina inteira (fig.~\ref{fig:timeline}). O animal aprendeu há muito tempo que a ação $A$ dá suco. Num certo momento o mundo muda: $A$ deixa de dar suco, e quem passa a dá-lo é a ação $B$, que o animal quase não experimenta.

\emph{Erro.} O suco depois de $A$ não veio, e os neurônios \nt{DOPA} respondem com uma queda, $\delta<0$ pela fórmula do erro~\eqref{eq:ctd}. As sinapses que acabaram de escolher $A$ carregam o traço e enfraquecem.

\emph{Surpresa.} Uma queda é um acaso comum. Mas as quedas se repetem, e os erros ficam maiores que o normal. Pela hipótese do capítulo~\ref{sec:nora}, é justamente esse o sinal de \nt{NORA}: o mundo mudou. O ganho $g$ cai, a escolha fica suave, e o ruído leva com mais frequência a experimentar $B$.

\emph{Escrita.} Pela hipótese do capítulo~\ref{sec:acet}, depois da mudança \nt{ACET} sobe e põe a rede em modo de escrita: as conexões internas, que guardam o hábito antigo, ficam enfraquecidas, a entrada dos sensores é reforçada, e os pesos mudam com facilidade.

\emph{Hábito novo.} Experimentar $B$ dá um suco que ninguém esperava: um pico, $\delta>0$, e as sinapses que escolheram $B$ se fortalecem. Algumas tentativas assim, e a expectativa $V$ alcança o mundo novo; os erros voltam a ser pequenos.

\emph{Retorno.} A surpresa passou, \nt{NORA} devolve o ganho alto, a escolha volta a ser rígida; \nt{ACET} baixa, e a rede, em modo de leitura, usa o que aprendeu. Durante todo esse tempo \nt{SERO}, por hipótese, define o horizonte $\tau_\gamma$: até que ponto ainda vale a pena esperar por um suco distante.

\begin{figure}[t]
\centering
\begin{tikzpicture}[>=Stealth,x=1.1cm,y=0.95cm]
\foreach \y/\lab in {4/{$\delta$ (\nt{DOPA})},3/{$g$ (\nt{NORA})},2/{$a$ (\nt{ACET})},1/{escolha de $B$}}{
  \draw[gray!60!black] (0,\y-0.35) -- (10,\y-0.35);
  \node[left,font=\scriptsize] at (0,\y) {\lab};}
\draw[densely dashed,gray!60!black] (2,0.4) -- (2,4.55) node[above,font=\scriptsize,black] {o mundo mudou};
\draw[densely dashed,gray!60!black] (5,0.4) -- (5,4.55) node[above,font=\scriptsize,black] {primeiro suco por $B$};
\pic[scale=0.85] at (2,5.25) {nojuice};
\pic[scale=0.85] at (5,5.25) {juice};
% delta: dips after the change, burst at the first juice for B, then back to zero
\draw[thick,red!70!black] (0,3.65) -- (2.3,3.65) -- (2.3,3.3) -- (2.5,3.3) -- (2.5,3.65) -- (3.2,3.65) -- (3.2,3.3) -- (3.4,3.3) -- (3.4,3.65) -- (4.1,3.65) -- (4.1,3.35) -- (4.3,3.35) -- (4.3,3.65) -- (5.0,3.65) -- (5.0,4.35) -- (5.2,4.35) -- (5.2,3.65) -- (5.8,3.65) -- (5.8,4.1) -- (6.0,4.1) -- (6.0,3.65) -- (6.6,3.65) -- (6.6,3.85) -- (6.8,3.85) -- (6.8,3.65) -- (10,3.65);
% gain: high, drops after surprise, recovers
\draw[thick,blue!60!black] (0,3.2) -- (3.0,3.2) -- (3.4,2.75) -- (5.8,2.75) -- (7.2,3.2) -- (10,3.2);
% ACh: low, rises after the change, falls after learning
\draw[thick,green!45!black] (0,1.75) -- (2.8,1.75) -- (3.3,2.25) -- (6.8,2.25) -- (7.8,1.75) -- (10,1.75);
% choice of B
\draw[thick,black] (0,0.7) -- (3.0,0.7) plot[domain=3:10,samples=40] (\x,{0.7+0.6/(1+exp(-(\x-5.3)*1.8))});
\draw[->,gray!70!black] (0,0.2) -- (10.2,0.2) node[below left,font=\scriptsize,black] {tempo};
\end{tikzpicture}
\caption{Um evento acompanhado através da máquina inteira. É um esquema qualitativo, baseado nos modelos e hipóteses dos capítulos~\ref{sec:dofa}--\ref{sec:acet}, e não o resultado de um cálculo. Depois da mudança \nt{DOPA} responde com quedas; as quedas repetidas, por hipótese, elevam o sinal de surpresa, \nt{NORA} reduz o ganho, \nt{ACET} liga a escrita; o primeiro suco por $B$ dá um pico, a escolha passa para $B$, e os botões voltam ao lugar.}
\label{fig:timeline}
\end{figure}

Note que nenhum botão sabe o que os outros fazem. Cada um responde aos seus próprios indícios (o erro, o seu tamanho fora do comum, a incerteza), e a ordem das ações se forma sozinha, como num circuito assíncrono, em que cada bloco dispara quando as suas entradas estão prontas. Esse trabalho coordenado como um todo, até onde sabemos, ainda não foi testado: os botões foram medidos um a um, e o seu funcionamento conjunto é hipótese.

\section{Em que esta máquina difere de um computador}

\begin{table}[t]
\centering
\footnotesize
\setlength{\tabcolsep}{4pt}
\renewcommand{\arraystretch}{1.15}
\begin{tabular}{>{\raggedright}p{2.2cm}>{\raggedright}p{4.2cm}>{\raggedright\arraybackslash}p{6.6cm}}
\toprule
& Computador comum & Cérebro \\
\midrule
tempo & relógio comum, cerca de um gigahertz & não há relógio comum; cada neurônio muda por conta própria, as janelas são definidas por eventos e ritmos locais (cap.~\ref{sec:glutcomputer}, \ref{sec:gaba}, \ref{sec:code}) \\
elementos & portas binárias confiáveis & elementos analógicos de limiar; a sinapse perde a maior parte dos disparos (cap.~\ref{sec:code}) \\
sinal da conexão & qualquer & definido pelo neurônio remetente (cap.~\ref{sec:neurontypes}) \\
memória & separada do processador & nas mesmas conexões que o cálculo e na atividade autossustentada (cap.~\ref{sec:glutcomputer}, \ref{sec:acet}) \\
aprendizado & retropropagação do erro & regras locais, um único número de difusão $\delta$ (cap.~\ref{sec:dofa}) e fios de erro para as saídas do circuito motor (cap.~\ref{sec:motorlearning}) \\
controle & registradores e barramento de instruções & quatro canais de difusão, cada um um feixe de algumas linhas (cap.~\ref{sec:serocomputer}, \ref{sec:dofa}--\ref{sec:acet}) \\
energia & dezenas a centenas de watts por processador & cerca de $20$ watts para o cérebro inteiro, em média cerca de um disparo por segundo por neurônio (cap.~\ref{sec:code}) \\
\bottomrule
\end{tabular}
\caption{O cérebro e um computador comum.}
\label{tab:computer}
\end{table}

A tabela~\ref{tab:computer} resume as principais diferenças. Nenhuma delas é um defeito que a natureza não teve tempo de corrigir: todas fazem parte de uma mesma solução. Sem um relógio comum, são necessários sensores de ``liga--desliga'' e detectores de coincidência. Elementos pouco confiáveis precisam da redundância de muitos neurônios. A memória que coincide com o cálculo precisa de \nt{ACET}, para que a escrita não estrague a leitura. O aprendizado sem fios de retorno precisa de \nt{DOPA} e de ruído. E o limite de energia de um disparo por segundo obriga toda a linguagem a ser econômica e a gastar disparos com novidades.

\section{O que não sabemos}

O mais honesto é terminar este resumo com a lista do que não sabemos. Cada item é um problema para um engenheiro.

\emph{O código.} Conhecemos a capacidade do canal de disparos e sabemos que o receptor escolhe que parte da mensagem ler (capítulo~\ref{sec:code}). Mas não se sabe o quanto o córtex usa o instante exato dos disparos, nem se os neurônios têm ``palavras''.

\emph{Aprendizado através de muitas camadas.} Um sinal de difusão ensina devagar, e mais devagar a cada neurônio que influi no resultado (proposição~\ref{prop:broadcastcost}). Então não se sabe como o córtex ajusta bilhões de pesos em circuitos de muitas camadas; há modelos em que o erro entre camadas é levado por rajadas de disparos~\cite{Payeur2021}. Talvez parte da resposta esteja nos cálculos dentro dos ramos do próprio neurônio, onde um neurônio se comporta como uma pequena rede de várias camadas: para reproduzir a resposta do modelo de um único neurônio cortical, uma rede artificial precisa de cinco a oito camadas~\cite{Beniaguev2021}, e os ramos dos neurônios humanos conseguem calcular até o OU exclusivo~\cite{Gidon2020}. Outro candidato é o autoaprendizado: se o córtex aprende a prever as suas próprias entradas, o erro surge no local, em cada camada, e nenhum sinal comum é necessário (capítulo~\ref{sec:dofa})~\cite{Keller2018}.

\emph{O que os canais de difusão transmitem de fato.} Para \nt{DOPA} há um quadro padrão e seis extensões dele (capítulo~\ref{sec:dofaalt}); para \nt{NORA} e \nt{ACET}, hipóteses plausíveis (capítulos~\ref{sec:nora}, \ref{sec:acet}); para \nt{SERO}, sobretudo palpites.

\emph{Coordenação no tempo.} Vimos como calcular uma função, medir o tempo a partir de um evento e cortar o fluxo em janelas com um ritmo local. Mas como uma rede enorme sem relógio comum coordena o trabalho de regiões distantes, só se entende em parte.

\emph{Energia.} Vinte watts para tudo. Entendemos por que o código deve ser esparso (proposição~\ref{prop:economy}), mas construir uma máquina que faça o mesmo com a mesma potência ninguém ainda sabe~\cite{MehonicKenyon2022}.

O cérebro é uma máquina de calcular montada sem engenheiro, mas segundo leis que qualquer engenheiro reconhece: circuitos RC, limiares, realimentações, filtros, barramentos e registradores de difusão. Lemos o seu esquema só em parte. Quem vai terminar de lê-lo são os que sabem ler esquemas.

\section{O que vem a seguir no livro}

Este capítulo reuniu o núcleo de cálculo da máquina. Os demais capítulos vão além dele. Os capítulos~\ref{sec:sensors} e~\ref{sec:recognition} tratam das entradas: como os sensores transformam luz, som e moléculas em disparos e como a máquina reconhece objetos neles. Os capítulos~\ref{sec:muscles}--\ref{sec:chemoutput} tratam das saídas e do que as atende: os músculos, a dor como sinal de alarme, o aprendizado de movimentos por fios de erro separados e a lenta saída química para o corpo. O capítulo~\ref{sec:debugging} trata de como a máquina é depurada de fora, inclusive com interfaces cérebro-computador. Os capítulos~\ref{sec:eetypes} e~\ref{sec:dendrites} classificam os neurônios mais uma vez, agora pela função elétrica e pelo número de entradas. O capítulo~\ref{sec:sleep} trata do que a máquina faz quando está desligada do mundo, e os capítulos~\ref{sec:livingmachine} e~\ref{sec:dishbrain}, de máquinas montadas com neurônios vivos num chip.

\section{Exercícios}

Os exercícios deste capítulo juntam resultados de capítulos diferentes: cada um se apoia em pelo menos dois deles.

\begin{exercise}[\lvlA]
\label{ex:10:1}
\emph{Barramento e registradores.} Os quatro canais de difusão são feixes de cerca de cinco linhas (proposição~\ref{prop:bundle}); cada linha muda uma vez por segundo e tem quatro níveis distinguíveis. Quantos anos o controle levaria para transmitir o que o barramento \nt{GLUT} ($8.6\cdot10^{10}$ neurônios com a taxa~\eqref{eq:energy}, precisão de $1\ms$ por~\eqref{eq:spikeentropy}) transmite num segundo?
\end{exercise}

\begin{exercise}[\lvlB]
\label{ex:10:2}
\emph{Dois botões no mesmo laço.} No laço~\eqref{eq:ei} os botões de~\eqref{eq:machine} atuam na excitação: $\tau_E\dot E=-E+g[(1-a)w_{EE}E-w_{EI}I+h]$, e não controlam \nt{GABA}. Com os números da fig.~\ref{fig:ei}, uma rajada de \nt{NORA} eleva $g$ até $6$. Qual o menor $a$ com que o laço é estável? Os botões de \nt{NORA} e \nt{ACET} podem ser girados independentemente?
\end{exercise}

\begin{exercise}[\lvlB]
\label{ex:10:3}
\emph{De onde vem o preço da difusão.} As saídas de $N$ neurônios são $y_k=f(u_k)+\xi_k$, com ruído gaussiano independente de variância $\sigma^2$; $R\approx R_0+\sum_kR'_k\xi_k$ com $R'_k$ iguais, e \nt{DOPA} difunde $\delta=R-R_0$. Ache a relação sinal-ruído da correção $\delta\,\xi_jx_j$ numa tentativa. Como o número de tentativas cresce com $N$?
\end{exercise}

\begin{exercise}[\lvlB]
\label{ex:10:4}
\emph{Ligar o som ao clarão.} O som chega ao barramento em $5\ms$, o clarão em $30\ms$ mais a latência do primeiro disparo~\eqref{eq:latency} ($\tau=10\ms$, $U$ de $1.2\theta$ a $4\theta$); o fundo em cada entrada é de $5$ disparos por segundo. Escolha o atraso da linha do som e a menor janela de coincidência para todos os contrastes. Quantas ligações falsas por segundo o fundo produz?
\end{exercise}

\begin{exercise}[\lvlB]
\label{ex:10:5}
\emph{Simulação: quem percebe a mudança.} O crítico vê $y=s+\text{ruído}$ ($R=1$); com probabilidade $1/200$, $s$ salta para um novo valor com variância $J^2=4$. Simule um filtro de Kalman com $q=0$ que eleva $P$ até $J^2$ quando $E\leftarrow0.9E+\delta^2/(P+R)$ passa de $20$. Quanto o seu erro é menor que com o melhor $\alpha$ constante?
\end{exercise}

\begin{exercise}[\lvlA]
\label{ex:10:6}
\emph{Quantas tentativas para trocar de hábito.} A rede aprendeu $Q_A=0.8$, $Q_B=0.2$ e escolhe por~\eqref{eq:softmax}, $\sigma=1$, $g=8$. Agora $A$ dá suco com probabilidade $0.2$; $Q_A\leftarrow Q_A+\alpha(r-Q_A)$, $Q_B$ não muda. Após quantas escolhas de $A$ a probabilidade de escolher $A$ cai para $0.6$, com $\alpha=0.1$ e $0.5$? E se \nt{NORA} baixar $g$ para $2$?
\end{exercise}

\begin{exercise}[\lvlC]
\label{ex:10:7}
\emph{Um feixe em vez de um fio.} Saídas $y_k=w_k+\xi_k$, recompensa $R=-\sum_ky_k^2$; cada linha difunde a um grupo de $G$ neurônios o erro de suas saídas, $w_k\leftarrow w_k+\eta\,\delta_l\xi_k$. Mostre que, com $\eta$ ótimo, chegar a $\sum w_k^2=\varepsilon\sum w_k^2(0)$ exige $\approx(G+2)\ln(1/\varepsilon)$ tentativas. Com $\varepsilon=0.01$, quanto tempo aprende um circuito de $10^6$ neurônios e $5$ linhas, uma tentativa a cada $3$~s?
\end{exercise}

\chapter{As entradas da máquina: como os olhos, os ouvidos e outros sensores estão ligados a ela}
\chaptermark{As entradas da máquina}
\label{sec:sensors}

\section{O sensor como transdutor}

Na fig.~\ref{fig:machine} os dados entravam na máquina pela esquerda, pelo bloco ``sensores''. Agora vamos abrir esse bloco. Para o engenheiro, cada órgão dos sentidos é um \emph{transdutor} com um estágio de entrada analógico e um conversor analógico-digital no fim; só que os dígitos na saída não são números, e sim disparos.

O esquema é o mesmo para todos os órgãos dos sentidos (fig.~\ref{fig:transducer}). Uma grandeza física (luz, pressão, vibração, uma molécula) atua sobre uma proteína receptora na membrana da célula sensorial. Essa proteína funciona como uma comporta: ela mesma, ou por uma curta cadeia de reações, abre um canal iônico. Surge uma \emph{corrente receptora} e, com ela, uma tensão analógica na membrana. Em seguida a tensão passa por um controle de ganho e chega ao codificador de disparos, o já conhecido neurônio integrador~\eqref{eq:glutneuron}, que transforma o nível em taxa e em instantes de disparo. A saída é quase sempre a mesma: os neurônios sensoriais do olho, do ouvido, do olfato e da pele liberam glutamato, de modo que as entradas se ligam diretamente ao barramento \nt{GLUT}.

\begin{figure}[t]
\centering
\begin{tikzpicture}[>=Stealth,
  blk/.style={draw,very thick,rounded corners=2pt,align=center,font=\footnotesize,minimum height=1.0cm,text width=1.9cm,fill=blue!5}]
\node[align=center,font=\small] (P) at (0.1,0) {luz, som,\\pressão,\\molécula};
\node[blk] (R) at (3.0,0) {receptor-\\comporta};
\node[blk] (G) at (6.5,0) {controle\\de ganho};
\node[blk] (E) at (10.0,0) {codificador\\de disparos};
\node[align=center,font=\small] (B) at (13.4,0) {barramento\\\nt{GLUT}};
\draw[->,very thick] (P) -- (R);
\foreach \ic/\xx in {sun/-1.1,speaker/-0.3,press/0.5,molecule/1.3}{\pic[scale=0.85] at (\xx,1.05) {\ic};}
\draw[->,very thick] (R) -- node[above,font=\scriptsize] {corrente} (G);
\draw[->,very thick] (G) -- node[above,font=\scriptsize] {nível} (E);
\draw[->,very thick,red!70!black] (E) -- node[above,font=\scriptsize,black] {disparos} (B);
\draw[decorate,decoration={brace,amplitude=5pt,mirror},gray!60!black] (1.8,-0.8) -- (7.7,-0.8) node[midway,below=5pt,font=\scriptsize,black] {parte analógica};
\draw[decorate,decoration={brace,amplitude=5pt,mirror},gray!60!black] (8.8,-0.8) -- (11.2,-0.8) node[midway,below=5pt,font=\scriptsize,black] {disparos};
\end{tikzpicture}
\caption{Qualquer órgão dos sentidos como uma cadeia de conversões. A proteína receptora abre um canal e gera corrente; o controle de ganho a adapta às circunstâncias; o codificador~\eqref{eq:glutneuron} transforma o nível em disparos, que seguem para o barramento \nt{GLUT}. No olho e no ouvido os primeiros estágios não emitem disparo algum: o sinal continua analógico até precisar ser levado longe.}
\label{fig:transducer}
\end{figure}

Um detalhe é bem conhecido de quem trabalha com eletrônica. No olho e no ouvido as primeiras células não emitem disparo algum: passam às vizinhas uma tensão contínua, como um estágio analógico numa placa. Os disparos só aparecem onde o sinal precisa ir longe. Um sinal analógico, ao longo de milímetros, se atenua e acumula ruído, enquanto um disparo, como vimos no capítulo~\ref{sec:neurontypes}, chega ao fim do fio com a mesma altura. Digitaliza-se onde começa a linha longa.

\section{No limite da física}

Os sensores do cérebro trabalham perto dos limites físicos de sensibilidade. O sensor de luz no escuro responde a \emph{um único fóton}: a absorção de uma só partícula de luz dá uma corrente de cerca de um picoampere, perceptível acima do seu próprio ruído~\cite{Baylor1979}. O sensor de som percebe um deslocamento do seu feixe sensível menor que um nanômetro~\cite{Hudspeth1989}.

Quando há pouca luz, a precisão não é limitada pelo sensor, e sim pela própria luz: os fótons chegam ao acaso, num fluxo de Poisson. Se durante o tempo de acumulação $T$ chegam ao sensor em média $N=\Phi T$ fótons, o seu número flutua em $\sqrt N$. Para que um acréscimo $\Delta N$ seja perceptível, ele deve ser algumas vezes maior que essa flutuação, e a fração distinguível de variação do brilho é
\begin{empheq}[box=\keybox]{equation}
\frac{\Delta I}{I} \;\approx\; \frac{k}{\sqrt{\Phi\,T}} .
\label{eq:photonnoise}
\end{empheq}
É a mesma fórmula da contagem de disparos~\eqref{eq:rateerror}: o ruído balístico (\emph{shot noise}) dos fótons e o dos disparos obedecem à mesma lei. No escuro, o olho só pode melhorar a precisão de um jeito: acumulando fótons por mais tempo, e o tempo de acumulação de fato cresce até alguns décimos de segundo. Por isso enxergamos mais devagar no crepúsculo.

\section{Faixa dinâmica: controle de ganho}

O problema de engenharia mais difícil dos sensores é a faixa. A iluminação, de uma noite estrelada até a neve ao sol, varia cerca de dez ordens de grandeza; a intensidade do som, do limiar de audição até o limiar da dor, doze. Já a taxa de disparos vai de zero a algumas centenas por segundo, menos de três ordens. Não dá para encaixar linearmente essa entrada nessa saída.

A solução é a mesma de qualquer receptor: \emph{controle automático de ganho}. As respostas das células sensoriais do olho são bem descritas pela fórmula
\begin{empheq}[box=\keybox]{equation}
r \;=\; r_{\max}\,\frac{I}{I+\sigma} ,
\label{eq:nakarushton}
\end{empheq}
em que $\sigma$ é o brilho em que a resposta chega à metade~\cite{NakaRushton1966}. Por si só,~\eqref{eq:nakarushton} cobre apenas duas ou três ordens. Mas $\sigma$ não é constante: com um trabalho prolongado sobre um fundo claro, ela cresce acompanhando o brilho médio. A curva de resposta se desloca ao longo do eixo logarítmico do brilho e sempre coloca o seu trecho íngreme onde a entrada está agora (fig.~\ref{fig:adaptation}). É a mesma divisão pela atividade total da inibição por derivação do capítulo~\ref{sec:gaba}~\cite{Carandini2012}, só que aqui o denominador é o brilho médio dos últimos segundos.

\begin{figure}[t]
\centering
\begin{tikzpicture}[>=Stealth,x=1.25cm,y=2.8cm]
\draw[->,gray!70!black] (-2,0) -- (6.4,0) node[below left,font=\small,black] {$\log_{10} I$};
\draw[->,gray!70!black] (-2,0) -- (-2,1.15) node[left,font=\small,black] {$r/r_{\max}$};
\foreach \u in {-2,0,2,4,6}{\draw[gray!60!black] (\u,-0.02) -- (\u,0.02); \node[below,font=\scriptsize] at (\u,0) {$\u$};}
\draw[densely dashed,gray!60!black] (-2,0.5) -- (6.2,0.5);
\foreach \s/\c/\lab in {0/blue!60!black/{fundo escuro},2/green!45!black/{médio},4/red!70!black/{claro}}{
  \pgfmathsetmacro{\lo}{max(-2,\s-3)}
  \draw[very thick,\c] (-2,0) -- (\lo,0) plot[domain=\lo:6.2,samples=80] (\x,{1/(1+10^(\s-\x))});
  \node[font=\scriptsize,\c,anchor=west] at (\s+0.15,0.42) {\lab};}
\foreach \s/\ic in {0/moon,2/cloud,4/sun}{\pic at (\s+0.6,0.64) {\ic};}
\end{tikzpicture}
\caption{Controle de ganho do sensor de luz segundo~\eqref{eq:nakarushton}. Cada curva cobre duas ou três ordens de brilho; ao passar para um fundo mais claro, $\sigma$ cresce e a curva se desloca ao longo do eixo logarítmico. Juntas, as curvas que se deslocam cobrem a faixa inteira.}
\label{fig:adaptation}
\end{figure}

O controle tem um preço, conhecido do capítulo~\ref{sec:code}: o sensor não informa o brilho, e sim o brilho \emph{em relação ao fundo}. Uma entrada constante deixa aos poucos de provocar resposta, e cada mudança provoca. As entradas da máquina falam, desde o início, de mudanças e não de níveis, e nisso está a mesma economia de disparos da proposição~\ref{prop:economy}~\cite{Barlow1961}. Daí também os sensores de ligar e de desligar do capítulo~\ref{sec:glutcomputer}~\cite{Schiller1992}: uns informam que ficou mais claro, outros que ficou mais escuro.

Os engenheiros repetiram esse truque nas \emph{câmeras de eventos}. Essa câmera não tira quadros ao ritmo de um relógio: cada pixel, por conta própria e sem ciclo comum, emite um evento ``mais claro'' ou ``mais escuro'' quando o logaritmo do brilho muda por uma fração dada~\cite{Lichtsteiner2008}. É exatamente o código de dois fios liga--desliga do capítulo~\ref{sec:glutcomputer}, realizado em silício, e ele dá à câmera uma faixa e uma rapidez inalcançáveis para um sensor de imagem comum~\cite{Gallego2022}.

\section{O olho: compressão até caber no fio}

O olho tem cerca de $10^8$ sensores de luz~\cite{Curcio1990}, e os neurônios de saída, que levam a imagem ao cérebro, são cerca de um milhão. Antes que o sinal saia do olho, ele é comprimido cerca de cem vezes. Os principais truques de compressão já nos são conhecidos. Cada neurônio de saída responde à diferença entre o brilho numa pequena mancha e o brilho em volta dela: é a subtração dos vizinhos por inibição, como no capítulo~\ref{sec:gaba}. Trechos uniformes da imagem quase não custam nada, enquanto bordas e mudanças são transmitidas. Os diferentes neurônios de saída são especializados: uns informam que clareou, outros que escureceu, outros ainda o movimento numa certa direção (fig.~\ref{fig:direction}); no camundongo foram contados mais de trinta desses canais de saída~\cite{Baden2016}.

Qual é a capacidade do cabo resultante? No porquinho-da-índia, a partir de registros dos neurônios de saída, mediram-se cerca de $875$ mil bits por segundo para $\sim10^5$ desses neurônios; recalculando para o milhão de neurônios de saída humanos, obtém-se da ordem de $10^7$ bits por segundo~\cite{Koch2006}, o mesmo que um velho cabo de rede de dez megabits. Com uma taxa média de alguns disparos por segundo por neurônio, isso dá alguns bits por disparo, como já saía no capítulo~\ref{sec:code}.

\section{O ouvido: um banco de filtros}

O ouvido resolve outro problema: o som precisa ser decomposto em frequências. Um engenheiro poria um banco de filtros passa-faixa, e o ouvido faz exatamente isso, mecanicamente. A vibração sonora percorre uma membrana elástica cujas propriedades mudam suavemente ao longo do seu comprimento, e cada trecho ressoa na sua própria frequência. Os sensores assentados num trecho respondem apenas à sua banda (fig.~\ref{fig:filterbank}). As frequências se distribuem ao longo do lugar de forma aproximadamente logarítmica: cada oitava recebe uma fração parecida dos sensores. O número do sensor que disparou é a frequência: o código de lugar do capítulo~\ref{sec:code}.

\begin{figure}[t]
\centering
\begin{tikzpicture}[>=Stealth,x=1.35cm,y=2.2cm]
\draw[->,gray!70!black] (-0.3,0) -- (7.6,0) node[above left,font=\small,black] {frequência, Hz};
\draw[->,gray!70!black] (-0.3,0) -- (-0.3,1.15) node[left,font=\small,black] {resposta};
\foreach \k/\f in {0/125,1/250,2/500,3/1000,4/2000,5/4000,6/8000}{
  \draw[gray!60!black] (\k,-0.02) -- (\k,0.02); \node[below,font=\scriptsize] at (\k,0) {\f};
  \draw[thick,blue!60!black] plot[domain={max(\k-1.2,-0.3)}:{\k+0.7},samples=60] (\x,{1/(1+((\x-\k)/(\x<\k?0.45:0.2))^4)});}
% a piano keyboard: seven white keys per octave, like the octaves on the axis
\foreach \i in {0,...,48}{\draw[gray!50!black,fill=white,line width=0.3pt] ({\i/7},-0.44) rectangle ({(\i+1)/7},-0.26);}
\foreach \o in {0,...,6}{\foreach \j in {0,1,3,4,5}{\fill[black] ({\o+(\j+1)/7-0.035},-0.35) rectangle ({\o+(\j+1)/7+0.035},-0.26);}}
\end{tikzpicture}
\caption{O ouvido como banco de filtros passa-faixa, esquema. Cada curva é a resposta dos sensores de um trecho a sons de frequências diferentes; as frequências centrais estão separadas por uma oitava, como teclas numa escala logarítmica. Os filtros reais têm uma encosta íngreme nas altas frequências e suave nas baixas. Embaixo, um teclado: nele, como no ouvido, cada oitava recebe o mesmo espaço.}
\label{fig:filterbank}
\end{figure}

Os ressonadores do ouvido não são passivos. Parte dos sensores sacode a própria membrana no ritmo do som e amplifica as vibrações fracas milhares de vezes em amplitude (em $66$--$76$\,dB), e com o aumento do volume o ganho cai~\cite{Ruggero1997,Robles2001}. Para o engenheiro, é um amplificador com realimentação positiva colocado bem perto do limite da autoexcitação: a mesma vida no limite da rede do capítulo~\ref{sec:gaba}; perto do limite, o sistema é especialmente sensível a sinais pequenos. A compressão de volume aqui é feita pelo mesmo amplificador: o que é baixo é amplificado muito, o que é alto, pouco.

O ouvido tem ainda um segundo código. Nas frequências baixas, os disparos dos neurônios vêm em fase com o som: o neurônio dispara num certo trecho de cada onda, embora não em toda onda. No porquinho-da-índia, o travamento de fase começa a enfraquecer em torno de $600$\,Hz e ainda é perceptível até cerca de $3.5$\,kHz~\cite{PalmerRussell1986}. Assim, os instantes dos disparos guardam o tempo exato do som e, a partir dele, a diferença de chegada do som aos dois ouvidos, pela qual a rede do capítulo~\ref{sec:glutcomputer} encontra a direção~\cite{Grothe2010}. Nas frequências altas os disparos não acompanham a onda, e resta apenas o código de lugar. Um mesmo ouvido usa os dois códigos do capítulo~\ref{sec:code}: o tempo onde os neurônios conseguem acompanhar e o lugar onde não conseguem.

\section{Os demais sensores}

\begin{table}[t]
\centering
\footnotesize
\setlength{\tabcolsep}{3pt}
\renewcommand{\arraystretch}{1.15}
\begin{tabular}{>{\raggedright}p{1.9cm}>{\raggedright}p{2.6cm}>{\raggedright}p{4.0cm}>{\raggedright\arraybackslash}p{4.6cm}}
\toprule
Sentido & O que mede & Como é a entrada & Truque do livro \\
\midrule
visão & luz & $\sim10^8$ sensores, compressão de $\sim100$ vezes, $\sim10^7$ bit/s & controle de ganho, subtração dos vizinhos, dois fios, direção do movimento \\
audição & pressão do ar & banco de filtros mecânico, amplificador ativo & código de lugar e código temporal, atrasos \\
tato & pressão, vibração & sensores de adaptação rápida e lenta & par ``filtro passa-altas --- filtro passa-baixas'' \\
temperatura & calor & sensores separados de calor e de frio & código de dois fios \\
olfato & moléculas & $\sim400$ tipos de receptores, um tipo por sensor & código combinatório: o cheiro é o conjunto de tipos que dispararam \\
paladar & substâncias no alimento & cinco qualidades: doce, amargo, salgado, azedo, umami & linhas rotuladas; parte das células libera ATP ou serotonina, e não glutamato \\
posição do corpo & estiramento dos músculos & sensores de comprimento e de velocidade de estiramento & realimentação para o regulador do movimento \\
\bottomrule
\end{tabular}
\caption{Como os sensores estão ligados. Todos os dispositivos da terceira coluna foram medidos; a coluna da direita os liga aos truques dos capítulos anteriores.}
\label{tab:sensors}
\end{table}

Os demais sentidos estão reunidos na tabela~\ref{tab:sensors}, e em quase toda linha se reconhece um truque dos capítulos anteriores. O tato usa um par de filtros: alguns sensores de pressão respondem ao toque e logo se calam, outros mantêm a resposta enquanto a pressão existe. A temperatura é transmitida por dois fios, separados para o calor e para o frio. A entrada mais incomum é o olfato. O ser humano tem cerca de quatrocentos tipos de receptores olfativos, e cada neurônio olfativo carrega um só tipo~\cite{Mombaerts2004}. Um cheiro não ativa um tipo, e sim um conjunto, e a mensagem é \emph{quais} tipos dispararam. Para o engenheiro, é um vetor binário de comprimento cerca de quatrocentos, com um número astronômico de valores possíveis.

\begin{keyblock}
\begin{proposition}[Entradas da máquina]
\label{prop:sensors}
Cada sentido é um transdutor que trabalha no limite físico de sensibilidade, comprime uma enorme faixa dinâmica por controle automático de ganho e transmite pelo barramento \nt{GLUT} sobretudo as \emph{mudanças}. Para isso, os sensores usam os mesmos truques da própria máquina: o código de dois fios, o código de lugar, o código temporal e a divisão pelo fundo. A construção dos sensores foi medida; que os seus códigos estejam ajustados para a máxima informação por disparo é uma hipótese bem fundamentada.
\end{proposition}
\end{keyblock}

As entradas, como todo o resto, funcionam de forma assíncrona. Cada sensor emite um disparo quando tem algo a informar, e os diferentes sentidos chegam com atrasos diferentes: o som em milissegundos, a luz em dezenas de milissegundos. Como a máquina os junta numa única imagem do mundo, sem um relógio comum, é um dos problemas de coordenação no tempo do capítulo~\ref{sec:synthesis}.

\section{Exercícios}

\begin{exercise}[\lvlA]
\label{ex:11:1}
\emph{Quanto tempo acumular fótons.} No crepúsculo chegam ao sensor $100$ fótons por segundo; um acréscimo é perceptível se for $3$ vezes maior que a flutuação~\eqref{eq:photonnoise}. (a)~Quanto tempo um sensor precisa para perceber uma mudança de brilho de $10\,\%$? (b)~Quantos sensores é preciso somar para conseguir em $0.2\,\text{s}$, e quantas vezes cai a resolução espacial?
\end{exercise}

\begin{exercise}[\lvlA]
\label{ex:11:2}
\emph{Quantos bits há num disparo.} (a)~O olho transmite $\sim10^7$ bit/s por $10^6$ neurônios de saída com taxa de $3$--$5$ disparos por segundo. Que fração da capacidade~\eqref{eq:spikeentropy} com $\Delta t=1\ms$ ele usa? (b)~Um cheiro ativa $20$ de $\approx400$ tipos de receptores. Quantos bits leva esse conjunto e quantos tipos em comum têm, em média, dois cheiros ao acaso?
\end{exercise}

\begin{exercise}[\lvlB]
\label{ex:11:3}
\emph{O que uma só curva~\eqref{eq:nakarushton} consegue.} (a)~Em que faixa de brilho a resposta sobe de $10\,\%$ a $90\,\%$ de $r_{\max}$? Quantas ordens de grandeza são? (b)~Quantas posições da curva deslocável são necessárias para cobrir dez ordens de brilho? E doze ordens de volume?
\end{exercise}

\begin{exercise}[\lvlB]
\label{ex:11:4}
\emph{O limite do código temporal no ouvido.} A diferença de chegada do som aos dois ouvidos vai até $0.22\,\text{m}/343\,\text{m/s}$. (a)~Os disparos vêm em fase com o tom: até que frequência a direção é determinada por eles sem ambiguidade? (b)~O disparo flutua $0.5\ms$, e é preciso distinguir $20$~\textmu s. Quantos neurônios de $100$ disparos por segundo é preciso promediar em $50\ms$?
\end{exercise}

\begin{exercise}[\lvlB]
\label{ex:11:5}
\emph{Uma câmera de eventos no cabo do olho.} Uma câmera de $10^6$ pixels com saída de $10^7$ bit/s envia um evento (endereço e sinal) quando $\ln I$ num pixel muda de $\ln1.2$. Com que velocidade máxima pode correr uma borda de $500$ pixels com salto de brilho $4$? Quantos quadros por segundo uma câmera comum, com $8$ bits por pixel, daria pela mesma saída?
\end{exercise}

\begin{exercise}[\lvlB]
\label{ex:11:6}
\emph{Simulação: controle de ganho.} O sensor~\eqref{eq:nakarushton} ajusta $\sigma$ com $\tau=1\,\text{s}$ em logaritmos, $\tau\,\dd\ln\sigma/\dd t=\ln I-\ln\sigma$, ou linearmente, $\tau\,\dd\sigma/\dd t=I-\sigma$; o fundo salta $1\to100\to1$. Após quantos segundos a resposta a um teste de $+10\,\%$ volta à metade da resposta adaptada, depois do salto para cima e para baixo, com cada lei?
\end{exercise}

\begin{exercise}[\lvlC]
\label{ex:11:7}
\emph{Para que mundo a curva está ajustada.} Com um ruído de saída pequeno e constante, a informação sobre o brilho é máxima quando $r(I)=r_{\max}\Phi_I(I)$, em que $\Phi_I$ é a função de distribuição dos brilhos. (a)~Para que distribuição a curva~\eqref{eq:nakarushton} é ótima, e qual a dispersão de $\log_{10}I$? (b)~Que curva é ótima com ruído de saída de Poisson?
\end{exercise}

\chapter{Reconhecimento de imagens e vídeo}
\chaptermark{Reconhecimento}
\label{sec:recognition}

\section{O problema: a mesma coisa com pixels diferentes}

O capítulo~\ref{sec:sensors} levou a imagem até a entrada da máquina: cerca de um milhão de neurônios de saída do olho, um fluxo comprimido que transmite sobretudo bordas e mudanças. Mas entre o sensor e a ação há um passo sobre o qual ainda não falamos. Um gato numa imagem não é um conjunto de pixels. Desloque-o meio quadro, gire-o, afaste-o, apague metade da luz: quase todos os pixels mudam, e a resposta ``gato'' precisa continuar a mesma. O engenheiro chama isso de \emph{invariância}: a resposta do classificador não deve mudar com deslocamento, escala, rotação e iluminação.

A dificuldade aparece bem num experimento simples. Tomemos uma ``retina'' unidimensional de $24$ pontos e duas figuras de cinco pontos cada, que podem estar em qualquer lugar. Um classificador que compara a entrada com um modelo memorizado num só lugar acerta em $49$--$52\%$ dos casos, o mesmo que uma moeda. O deslocamento quebra por completo a comparação pixel a pixel. É preciso outro esquema.

\section{Uma linha de montagem de estágios}

A velocidade com que o cérebro faz isso já conhecemos do capítulo~\ref{sec:code}: um animal numa imagem mostrada por um instante é reconhecido em cerca de $150\ms$, e o sinal passa por cerca de uma dezena de estágios, com $10$--$20\ms$ para cada um~\cite{Thorpe1996}. Em cada estágio o neurônio só tem tempo de emitir zero ou um disparo. Portanto, o reconhecimento rápido é uma única passagem pela linha de montagem, sem longas reflexões nem repetições. A questão é o que cada estágio faz.

A resposta sugerida pelas medições nos primeiros estágios da visão~\cite{HubelWiesel1962} consiste em duas operações alternadas (fig.~\ref{fig:conveyor}).

\emph{Seletividade.} Um neurônio do estágio $S$ olha para um pequeno trecho da entrada e dispara só se ali houver uma certa combinação de partes: esta borda aqui e aquela ao lado. É um E sobre as partes: o neurônio de limiar do capítulo~\ref{sec:glutcomputer}, com limiar mais alto do que qualquer parte isolada fornece.

\emph{Invariância.} Um neurônio do estágio $C$ reúne as saídas de muitos neurônios $S$ que procuram a mesma combinação em lugares diferentes e dispara se ela for encontrada em algum lugar. É um OU sobre as posições, ou, mais exatamente, a escolha do máximo.

Para o modelo unidimensional, as duas operações se escrevem assim:
\begin{empheq}[box=\keybox]{equation}
s_{f,p} \;=\; \Bigl[\sum_i t_{f,i}\,x_{p+i} - \theta\Bigr]_+ ,
\qquad
c_f \;=\; \max_p\, s_{f,p} ,
\label{eq:scstage}
\end{empheq}
em que $x$ é a entrada, $t_{f,i}$ é o modelo do atributo $f$, $p$ é a posição e $\theta$ é o limiar. A primeira expressão torna a resposta seletiva, a segunda a torna independente da posição. O máximo de muitas entradas a rede obtém com os recursos do capítulo~\ref{sec:gaba}: a inibição mútua deixa passar a entrada mais forte (proposição~\ref{prop:wta}), e a divisão pela atividade total equaliza as respostas com brilhos diferentes~\cite{Carandini2012}.

\begin{figure}[t]
\centering
\begin{tikzpicture}[>=Stealth,
  px/.style={draw,minimum size=0.36cm,inner sep=0pt},
  on/.style={px,fill=blue!55!black},
  snode/.style={circle,draw,very thick,minimum size=0.62cm,inner sep=0pt,font=\scriptsize,fill=blue!6},
  cnode/.style={circle,draw,very thick,minimum size=0.8cm,inner sep=0pt,font=\scriptsize,fill=red!10}]
% two inputs: the same shape at two positions
\foreach \row/\sh/\lab in {0/1/{quadro 1},1/7/{quadro 2}}{
  \begin{scope}[yshift=-\row*1.0cm]
  \node[left,font=\scriptsize] at (-0.25,0) {\lab};
  \foreach \i in {0,...,13}{\node[px] at (\i*0.4,0) {};}
  \foreach \d in {0,1,3,4}{\node[on] at ({(\sh+\d)*0.4},0) {};}
  \end{scope}}
% S stage: detectors of the shape at several positions
\foreach \k in {0,...,5}{
  \node[snode] (S\k) at ({0.8+0.8*\k},1.7) {$S$};}
\node[font=\scriptsize,left] at (-0.25,1.7) {estágio $S$: E sobre as partes};
\draw[->,thick,blue!60!black] (1.2,0.25) -- (S1.south);
\draw[->,thick,orange!85!black] (3.6,0.25) -- (S4.south);
\node[font=\scriptsize,blue!60!black,anchor=east] at (1.25,0.85) {quadro 1};
\node[font=\scriptsize,orange!85!black,anchor=west] at (3.9,0.85) {quadro 2};
% C stage
\node[cnode] (C) at (2.8,3.25) {$C$};
\node[font=\scriptsize,left] at (-0.25,3.25) {estágio $C$: OU sobre as posições};
\foreach \k in {0,...,5}{\draw[->,gray!60!black] (S\k.north) -- (C);}
\draw[->,very thick,red!70!black] (C.north) -- ++(0,0.6) node[above,font=\scriptsize,black] {``a figura está lá'' nos dois quadros};
% next stages
\draw[decorate,decoration={brace,amplitude=5pt},gray!60!black] (6.3,1.4) -- (6.3,-1.3);
\node[right,font=\scriptsize,align=left] at (6.5,0.05) {o mesmo\\objeto, outros\\pixels};
\draw[->,thick,densely dashed,gray!60!black] (3.6,3.25) -- (5.9,3.25) node[right,font=\scriptsize,black,align=left] {próximos pares\\$S$, $C$: maiores,\\mais complexos, mais invariantes};
\end{tikzpicture}
\caption{Um par de estágios da linha de montagem. Os neurônios $S$ procuram a mesma combinação de partes em lugares diferentes; o neurônio $C$ responde se ela for encontrada em algum lugar~\eqref{eq:scstage}. A figura no quadro 1 e no quadro 2 excita neurônios $S$ diferentes, mas o mesmo $C$. Os próximos pares de estágios repetem o mesmo com combinações cada vez maiores e mais complexas.}
\label{fig:conveyor}
\end{figure}

No mesmo modelo unidimensional, o par de estágios~\eqref{eq:scstage} reconhece a figura em qualquer lugar sem erros, e se cada ponto da entrada for invertido ao acaso com probabilidade $0.1$, em $86\%$ dos casos; com probabilidade $0.2$, em $70\%$. A moeda continua dando metade. Empilhe alguns desses pares: o primeiro procura bordas, o seguinte combinações de bordas, mais acima partes de objetos, no topo objetos inteiros. A cada par as combinações são maiores, e a resposta depende cada vez menos de onde e como o objeto está~\cite{Riesenhuber1999}.

\begin{keyblock}
\begin{proposition}[Linha de montagem de seletividade e invariância]
\label{prop:conveyor}
O reconhecimento rápido é uma única passagem por uma dezena de estágios. Em cada par de estágios, o E sobre as partes~\eqref{eq:scstage} torna a resposta seletiva, e o OU sobre as posições a torna independente do deslocamento. A alternância dessas duas operações dá uma resposta que distingue objetos e quase não depende de onde e como eles são vistos. A construção dos primeiros estágios e a velocidade da passagem foram medidas; que toda a cadeia se reduza a essas duas operações é uma simplificação.
\end{proposition}
\end{keyblock}

Se essa construção parece familiar, é porque deve parecer. Foi justamente essa alternância (procurar um atributo em todos os lugares e depois juntar lugares vizinhos) que os engenheiros levaram às redes artificiais \emph{convolucionais}~\cite{Fukushima1980}. E recentemente fizeram o caminho de volta: redes convolucionais treinadas para reconhecer objetos se revelaram o melhor modelo conhecido das respostas dos neurônios nos estágios superiores da visão~\cite{Yamins2014}. O resultado no topo se lê de forma simples: pelas taxas somadas de cerca de uma centena de neurônios do último estágio, um bloco de decisão linear reconhece o objeto um décimo de segundo depois da apresentação~\cite{Hung2005}. É o código de taxa de um grupo do capítulo~\ref{sec:code}.

\section{O professor que não existe: aprender com vídeo}

Uma rede convolucional é treinada com milhões de imagens rotuladas. Ao cérebro quase ninguém dá rótulos: a criança vê objetos por horas, e ouve a palavra ``xícara'' só de vez em quando. Então como os estágios $C$ sabem quais neurônios $S$ juntar? Afinal, para juntar ``esta figura aqui'' com ``esta figura ali'', é preciso já saber que é a mesma figura.

A pista vem do próprio vídeo. O mundo muda de forma contínua: um objeto no campo de visão continua sendo o mesmo objeto enquanto se desloca, gira e muda de iluminação, e só de vez em quando o olhar passa para outro. Tudo o que muda \emph{devagar} provavelmente pertence ao objeto, e tudo o que muda depressa, à sua posição e aparência~\cite{WiskottSejnowski2002}. Portanto basta ao neurônio $C$ a regra de Hebb com traço do capítulo~\ref{sec:dofa}: ligar o que veio em sequência, e não só o que veio ao mesmo tempo~\cite{Foldiak1991}:
\begin{empheq}[box=\keybox]{equation}
\tau_{\text{tr}}\,\frac{\dd\bar y_k}{\dd t} \;=\; -\bar y_k + y_k ,
\qquad
\frac{\dd\mathbf w_k}{\dd t} \;=\; \eta\,\bar y_k\,\bigl(\mathbf x - \mathbf w_k\bigr) .
\label{eq:traceinvariance}
\end{empheq}
Aqui $y_k$ é a atividade do neurônio $C$ número $k$, que venceu a competição pela inibição, $\bar y_k$ é o seu traço, o mesmo circuito RC da regra~\eqref{eq:threefactor}, e $\mathbf x$ são as saídas dos neurônios $S$. O neurônio que venceu na primeira vista do objeto ainda se lembra disso quando chega a segunda vista e a puxa para si também.

\begin{figure}[t]
\centering
\begin{tikzpicture}[>=Stealth,x=1.0cm,y=1.0cm]
\foreach \i/\lab in {0/1,1/2,2/3,3/4}{
  \node[draw,thick,fill=blue!8,minimum width=0.9cm,minimum height=0.55cm,font=\scriptsize] at (\i+0.5,2.2) {$A$ em \lab};}
\foreach \i/\lab in {0/4,1/3,2/2,3/1}{
  \node[draw,thick,fill=orange!12,minimum width=0.9cm,minimum height=0.55cm,font=\scriptsize] at (\i+6.5,2.2) {$B$ em \lab};}
\node[font=\scriptsize,gray!40!black] at (5.0,2.2) {pausa};
\draw[->,gray!70!black] (0,0) -- (10.4,0) node[below left,font=\scriptsize,black] {tempo};
% traces
\draw[thick,blue!60!black] (0,0.05) .. controls (0.4,1.2) and (0.8,1.3) .. (1.2,1.35) -- (4,1.4) .. controls (4.6,0.5) and (5.2,0.15) .. (6,0.08) -- (10,0.05);
\draw[thick,orange!85!black] (0,0.05) -- (6,0.05) .. controls (6.4,1.2) and (6.8,1.3) .. (7.2,1.35) -- (10,1.4);
\node[font=\scriptsize,blue!60!black,anchor=west] at (1.4,1.65) {traço $\bar y_1$: dura toda a passagem de $A$};
\node[font=\scriptsize,orange!85!black,anchor=west] at (7.2,1.65) {traço $\bar y_2$};
\draw[decorate,decoration={brace,amplitude=4pt},gray!60!black] (4.0,2.6) -- (0,2.6);
\draw[decorate,decoration={brace,amplitude=4pt,mirror},gray!60!black] (0,-0.35) -- (4,-0.35) node[midway,below=4pt,font=\scriptsize,black] {todas as vistas de $A$ se ligam ao neurônio 1};
\end{tikzpicture}
\caption{Como o vídeo ensina a invariância, esquema. O objeto $A$ passa por quatro posições, depois vem uma pausa, depois o objeto $B$. O traço~\eqref{eq:traceinvariance} do neurônio que venceu na primeira vista de $A$ dura toda a passagem, e todas as vistas de $A$ acabam ligadas a um só neurônio. A pausa apaga o traço, e $B$ fica com outro neurônio.}
\label{fig:tracevideo}
\end{figure}

Testamos isso num modelo (fig.~\ref{fig:tracevideo}). Dois neurônios $C$ competem pelas entradas de $16$ neurônios $S$: duas figuras em oito posições. A figura percorre todas as posições, $50\ms$ em cada uma, depois vem uma pausa de $0.3\,\text{s}$ e a próxima figura ao acaso. Nenhum rótulo. Se o traço é curto, $\tau_{\text{tr}}=5\ms$, os neurônios dividem a entrada pela \emph{posição}, e a resposta coincide com a figura não melhor que uma moeda: $52\%$ em média em dez execuções. Se o traço é comparável ao tempo de exibição de uma posição, a partir de $\tau_{\text{tr}}\approx20\ms$, cada neurônio reúne \emph{todas} as posições da sua figura (com $20$, $50$ e $200\ms$, em todas as dez execuções): a invariância foi aprendida só com o vídeo. A pausa entre objetos importa: sem ela, o traço passa para o objeto seguinte e os confunde.

O mesmo se vê em experimentos. Se durante um experimento o objeto é trocado discretamente enquanto o olho salta de um lugar para outro, os neurônios dos estágios superiores, em uma ou duas horas, passam a responder ao objeto trocado como se fosse o mesmo: eles ligam o que veio em sequência~\cite{LiDiCarlo2008}. A invariância no cérebro é de fato aprendida a partir da continuidade do tempo. Até que ponto essa regra explica todo o aprendizado visual é hipótese.

\section{Vídeo é movimento}

O vídeo não é só um fluxo de quadros para aprender, mas também a própria tarefa: o que se move, para onde e com que rapidez. O primeiro passo já conhecemos: o detector de direção do capítulo~\ref{sec:gaba} (fig.~\ref{fig:direction}), em que a inibição atrasada cancela o movimento num sentido. Esses detectores estão espalhados por todo o campo visual, e depois são tratados como os atributos de forma: estágios que reúnem muitos detectores da mesma direção em lugares diferentes respondem ao movimento de um objeto grande independentemente de onde ele está. É o mesmo par E e OU~\eqref{eq:scstage}, só que o atributo não é ``borda'', e sim ``borda que se deslocou no tempo $d$''.

O vídeo também é previsível: o próximo quadro é quase igual ao anterior. Pela hipótese da \emph{codificação preditiva}, os estágios superiores enviam aos inferiores uma predição, e para cima sobe principalmente o erro, aquilo que a predição não levou em conta~\cite{RaoBallard1999}. É a mesma economia de disparos da proposição~\ref{prop:economy}: pagar pelas novidades, e não pelo que já se sabe. Algumas observações concordam com essa hipótese, mas como construção geral da linha de montagem ela não está provada.

\section{O cérebro e a rede convolucional}

Até onde vai a semelhança? Para o reconhecimento rápido de um só objeto, ela é de fato profunda~\cite{Yamins2014}. Mas o cérebro tem também o que uma rede convolucional simples não tem.

\emph{Realimentação.} A linha de montagem do cérebro não é só direta: cada estágio tem conexões para trás e para os lados. Para imagens difíceis (com sobreposição, com pouca luz) a resposta dos estágios superiores se forma mais tarde, e essas respostas tardias são mais bem descritas por redes com realimentação do que por redes diretas~\cite{Kar2019}. Uma passagem resolve os casos fáceis; os difíceis exigem várias voltas.

\emph{Robustez.} Uma rede artificial pode ser enganada por uma falsificação de pixels imperceptível ao olho~\cite{Szegedy2014}: uma mudança que o ser humano não vê transforma um ``panda'' num ``gibão''~\cite{Goodfellow2015}. A visão humana é muito mais robusta a essas falsificações, mas não invulnerável: imagens que enganam muitas redes ao mesmo tempo, mostradas brevemente, também desviam a escolha humana~\cite{Elsayed2018}. Por que a robustez difere tanto é uma questão aberta; os candidatos são o ruído, a divisão pela atividade total e a realimentação.

\emph{O professor.} A rede precisa de milhões de exemplos rotulados; o cérebro, sobretudo de vídeo sem rótulos~\eqref{eq:traceinvariance} e de algumas dicas de \nt{DOPA} sobre o que importa.

\emph{Energia.} O cérebro inteiro gasta cerca de $20$ watts (proposição~\ref{prop:economy}), e o reconhecimento ocupa só parte disso; uma rede convolucional treinada num acelerador gráfico gasta centenas de watts.

\section{Ilusões: onde a máquina erra e por quê}
\label{sec:illusions}

O engenheiro que quer entender um circuito alheio aplica nele sinais incomuns e observa onde ele erra. Para a visão, esses sinais foram inventados há muito tempo: são as ilusões visuais. Uma ilusão não é uma avaria. Cada uma aponta para um certo mecanismo da máquina, que funciona corretamente em condições normais e se denuncia numa imagem escolhida de propósito.

\emph{Expectativas razoáveis.} A entrada é ruidosa, e a máquina a completa com o que costuma acontecer no mundo. Movimentos lentos são mais frequentes que rápidos; quando a entrada não é confiável (por exemplo, quando a imagem tem pouco contraste), a estimativa de velocidade se desloca para o lento, e um objeto pálido parece se mover mais devagar que um brilhante~\cite{Weiss2002}. Muitas ilusões de brilho se explicam da mesma forma: não vemos o brilho do pixel, e sim a superfície a que esse brilho mais frequentemente correspondeu no passado~\cite{Purves2011}. A foto do vestido que uns veem branco e dourado e outros azul e preto é o mesmo mecanismo: a imagem é ambígua, e as pessoas divergem na iluminação que supõem sem perceber~\cite{LaferSousa2015,Wallisch2017}. As diferenças entre pessoas foram medidas; que o cérebro aqui combine a entrada com a expectativa pelas regras da teoria das probabilidades é uma hipótese bem fundamentada.

\emph{Controle de ganho.} Olhe por um minuto para uma cachoeira e depois para as pedras paradas: as pedras vão flutuar para cima. Os canais ``para baixo'' e ``para cima'' são um par de fios, como em~\eqref{eq:dualrail}; o controle de ganho~\eqref{eq:nakarushton} abafou o canal cansado, e o equilíbrio do par se deslocou. As ilusões de inclinação e de contraste, em que o entorno muda a aparência do centro, se explicam pela divisão pela atividade dos vizinhos~\cite{Schwartz2009}, como no capítulo~\ref{sec:gaba}. Há também um exemplo instrutivo de cautela: as manchas escuras nos cruzamentos das faixas brancas de uma grade foram explicadas por muito tempo pela simples inibição dos vizinhos, mas experimentos com uma grade modificada mostraram que essa explicação não funciona~\cite{SchillerCarvey2005}.

\emph{Compensação de atrasos.} O caminho do olho até os estágios superiores leva dezenas de milissegundos, e um objeto em movimento se desloca nesse tempo. Se um ponto parado piscar ao lado dele, o objeto parece estar à frente do clarão, embora coincidissem~\cite{Nijhawan1994}. Por uma das explicações, a máquina prolonga o movimento para a frente para compensar o atraso, o mesmo que no capítulo~\ref{sec:muscles}: com realimentação atrasada, é preciso prever. Essa ilusão tem várias explicações, e a disputa não está resolvida.

\emph{Erro no código temporal.} Uma imagem parada de anéis com transições de cor bruscas parece girar. Trechos contrastados são processados mais rápido que os pálidos, e a diferença de latências~\eqref{eq:latency} é lida pelos detectores de direção como movimento~\cite{Conway2005}. A ilusão é disparada pelos pequenos saltos involuntários do olho e pelas piscadas: cada salto renova a entrada, e os detectores veem de novo ``movimento''~\cite{OteroMillan2012}. É o código temporal do capítulo~\ref{sec:code}, lido de forma errada.

\emph{Duas imagens numa.} A armação de um cubo que ora nos mostra uma face, ora outra, ou imagens diferentes nos dois olhos: a percepção salta a cada poucos segundos. Os melhores modelos são a inibição mútua de duas interpretações, uma fadiga lenta e ruído~\cite{MorenoBote2007}, ou seja, o mesmo esquema~\eqref{eq:halfcentre} que no capítulo~\ref{sec:muscles} gerava o ritmo do passo. O tempo das trocas é definido pelo ruído e pela fadiga; pelo modelo~\cite{MorenoBote2007}, o papel principal é do ruído, e a fadiga só ajuda.

E as redes artificiais? Uma rede treinada para prever o próximo quadro de um vídeo vê rotação nos mesmos anéis parados e não a vê nas imagens de controle~\cite{Watanabe2018}. Redes treinadas para limpar o ruído das imagens e aumentar a sua nitidez caem nas mesmas ilusões de cor e brilho que as pessoas~\cite{GomezVilla2020}. Portanto, parte das ilusões é consequência da própria tarefa que a máquina aprende, e não de particularidades do tecido nervoso. Quais ilusões são exclusivas do cérebro é um bom teste para qualquer modelo da visão.

\begin{keyblock}
\begin{proposition}[Ilusões como impressões dos mecanismos]
\label{prop:illusions}
Uma ilusão surge quando um mecanismo útil no mundo comum recebe uma entrada incomum: a expectativa vence uma entrada pouco confiável, o controle de ganho desloca um par de canais, a compensação de atraso leva o objeto para a frente, a diferença de latências é lida como movimento, a inibição mútua com ruído e fadiga salta. Cada ilusão é um sinal de teste que aponta para o seu mecanismo. As próprias ilusões foram medidas; as suas explicações vão das bem fundamentadas às controversas.
\end{proposition}
\end{keyblock}

\section{Resumo}

\begin{table}[t]
\centering
\footnotesize
\setlength{\tabcolsep}{4pt}
\renewcommand{\arraystretch}{1.15}
\begin{tabular}{>{\raggedright}p{3.0cm}>{\raggedright}p{4.6cm}>{\raggedright\arraybackslash}p{5.2cm}}
\toprule
O que a máquina faz & Como & O que se sabe \\
\midrule
reconhece em $150\ms$ & uma passagem por uma dezena de estágios & velocidade medida \\
torna a resposta seletiva & E sobre as partes, estágio $S$~\eqref{eq:scstage} & medido nos primeiros estágios \\
torna a resposta invariante & OU sobre as posições, estágio $C$; inibição e divisão & medido nos primeiros estágios; para os superiores, simplificação \\
dá a resposta & taxas de cerca de uma centena de neurônios do último estágio, leitura linear & medido \\
aprende sem rótulos & regra de Hebb com traço~\eqref{eq:traceinvariance} sobre vídeo contínuo & a troca do objeto muda as respostas: medido; como regra principal, hipótese \\
vê o movimento & detectores de direção, depois o mesmo par E/OU & medido \\
economiza no previsível & para cima vai o erro de predição & hipótese \\
resolve casos difíceis & realimentação, várias voltas & respostas tardias e papel da realimentação medidos \\
erra de modo previsível & ilusões: expectativas, controle de ganho, atrasos, inibição mútua com ruído e fadiga & ilusões medidas; explicações das fundamentadas às controversas \\
\bottomrule
\end{tabular}
\caption{Como a máquina reconhece imagens e vídeo.}
\label{tab:recognition}
\end{table}

O reconhecimento é montado com peças que já tínhamos (tabela~\ref{tab:recognition}): o limiar dá o E sobre as partes, a inibição mútua dá o OU sobre as posições, a divisão dá a independência do brilho, a regra de Hebb com traço substitui o professor que não existe. Os engenheiros tomaram da visão a alternância de seletividade e invariância e construíram com ela as redes convolucionais; o caminho de volta ainda não trouxe o principal: aprender com vídeo sem rótulos e quase sem energia.

\section{Exercícios}

\begin{exercise}[\lvlA]
\label{ex:12:1}
\emph{A taxa num estágio.} Um estágio da linha de montagem dura $15\ms$, os neurônios trabalham com taxa de $20$ disparos por segundo, os ruídos são correlacionados com $c=0.1$~\eqref{eq:correlated}. Qual é o erro relativo da taxa de um grupo de cem neurônios e o seu limite para um grupo arbitrariamente grande? Quantos níveis de taxa se distinguem num estágio?
\end{exercise}

\begin{exercise}[\lvlA]
\label{ex:12:2}
\emph{A energia de um reconhecimento.} Numa passagem de $150\ms$, dez estágios de $10^7$ neurônios emitem no máximo um disparo cada, ao custo de $10^{-10}$~J. (a)~Que fração da energia do cérebro inteiro nesses $150\ms$ custa o reconhecimento? (b)~Quantas vezes mais gasta um acelerador de $300$~W que reconhece a imagem em $10\ms$?
\end{exercise}

\begin{exercise}[\lvlB]
\label{ex:12:3}
\emph{Os limiares do estágio $S$.} ``Retina'' de $24$ pontos, figuras $A=11011$ e $B=10101$; o modelo~\eqref{eq:scstage} vale $+1$ nos uns da figura e $-1$ nos zeros. (a)~Com que limiares o neurônio $S$ perdoa um ponto invertido, mas fica calado com a figura alheia? (b)~Quantos neurônios $S$ são necessários para dez atributos de $5\times5$ numa imagem de $100\times100$?
\end{exercise}

\begin{exercise}[\lvlB]
\label{ex:12:4}
\emph{Quanto dura uma das duas imagens.} No esquema de saltos~\eqref{eq:halfcentre} com $\tau\ll\tau_a$, enquanto o grupo~1 vence, $r_2=0$ e $r_1=I-a_1$. (a)~Ache a duração da vitória com $I=1$, $\beta=2$, $b=2$, $\tau_a=0.4\,\text{s}$. (b)~Que $\tau_a$ dá uma troca de imagem a cada três segundos?
\end{exercise}

\begin{exercise}[\lvlB]
\label{ex:12:5}
\emph{Simulação: o preço da invariância.} Com a função \texttt{two\_stage} de \texttt{models/\allowbreak recognition\_models.py} ($4000$ tentativas), ache com que probabilidade de inversão de ponto $p$ o par de estágios $S$, $C$ com $N=24$ erra em um de cada quatro casos. Como muda a fração de acertos com $p=0.1$ de $N=8$ a $N=192$?
\end{exercise}

\begin{exercise}[\lvlB]
\label{ex:12:6}
\emph{Simulação: a janela do traço.} Com dez execuções de \texttt{trace\_learning} de \texttt{models/\allowbreak recognition\_models.py} e pausa de $0.3\,\text{s}$, ache para quais $\tau_{\text{tr}}$ entre $5\ms$ e $2\,\text{s}$ todas as execuções aprendem a invariância sem erro. De onde vem o limite superior dessa janela?
\end{exercise}

\begin{exercise}[\lvlB]
\label{ex:12:7}
\emph{Movimento em qualquer lugar.} Em pontos vizinhos de uma fileira, com passo de $0.1^\circ$, há detectores de direção do capítulo~\ref{sec:gaba}: a inibição de $A$ com atraso $d$ cancela a excitação de $B$ se elas se separarem no máximo $5\ms$; acima deles, um estágio $C$. Que atrasos são necessários para que $C$ fique calado com movimento inverso de $5$ a $20^\circ$ por segundo?
\end{exercise}

\begin{exercise}[\lvlC]
\label{ex:12:8}
\emph{Falsificação de pixels.} Quantos pontos invertidos são necessários para mudar a resposta do modelo \texttt{two\_stage} ($N=24$) numa figura limpa? E a do classificador ``mais de $3.5$ uns na entrada: é $A$'', também invariante ao deslocamento? Qual a sua fração de acertos com $p=0.1$, contra $0.86$ do par de estágios?
\end{exercise}

\chapter{As saídas da máquina: como o cérebro controla os músculos}
\chaptermark{As saídas da máquina}
\label{sec:muscles}
\begin{chaptergoals}
\item como o músculo transforma disparos em força: um conversor digital-analógico passa-baixas;
\item por que ligar unidades por tamanho dá a força em ponto flutuante, com passo proporcional à força;
\item como um par de músculos que só puxam define o torque pela diferença e a rigidez pela soma;
\item por que o atraso da realimentação limita a rapidez da correção e como inibição e fadiga dão ritmo.
\end{chaptergoals}

\section{A saída rápida}

Tudo o que a máquina faz rápido no mundo, ela faz com os músculos. Uma palavra, um olhar, um passo, um sorriso são contrações musculares. A segunda saída, a lenta, a química do corpo, é tratada no capítulo~\ref{sec:chemoutput}. Na fig.~\ref{fig:machine} a saída era a seta ``músculos''; agora vamos ver o que há por trás dessa seta.

O último elo da cadeia são os neurônios \nt{ACET}, justamente aqueles que na tabela~\ref{tab:types} têm a anotação ``saída para o músculo''. Cada fibra muscular recebe entrada de exatamente um desses neurônios, e um neurônio controla um grupo de fibras, de algumas centenas a um par de milhares~\cite{Feinstein1955mu}. Nos músculos que precisam de ajuste fino, as unidades são menores; nos grandes músculos das pernas, maiores. O neurônio com as suas fibras se chama \emph{unidade motora}: é o menor pedaço de músculo que o cérebro pode ligar separadamente.

A sinapse do neurônio no músculo é construída de um jeito bem diferente das sinapses do córtex do capítulo~\ref{sec:code}. Lá a maior parte dos disparos se perdia. Aqui cada disparo libera várias vezes mais acetilcolina do que o necessário para a fibra responder~\cite{WoodSlater2001}. É uma sinapse com \emph{margem de segurança}: um disparo do neurônio, exatamente uma contração das fibras. Dentro da máquina é possível se dar ao luxo de conexões pouco confiáveis e corrigir os erros com redundância; na saída, um erro custa um movimento, e a confiabilidade é comprada diretamente no hardware.

\section{A força: taxa e número de unidades}

Um disparo dá uma contração isolada, um abalo: a força cresce em algumas dezenas de milissegundos e decai. Uma boa aproximação para ele é~\cite{Fuglevand1993}
\begin{equation}
h(t) \;=\; F_1\,\frac{t}{T_c}\,e^{\,1-t/T_c},
\label{eq:twitch}
\end{equation}
em que $T_c$ é o tempo de subida, da ordem de $50\ms$, e $F_1$ é a força de um abalo. Se os disparos vêm com mais frequência, os abalos se somam:
\begin{empheq}[box=\keybox]{equation}
F(t) \;=\; \sum_k h\bigl(t-t_k\bigr) .
\label{eq:forcesum}
\end{empheq}
É o mesmo filtro passa-baixas que transformava disparos em concentração no capítulo~\ref{sec:serocomputer}, só que agora a saída não é concentração, e sim força. O músculo é um conversor digital-analógico de disparos em força (fig.~\ref{fig:tetanus}). No nosso modelo, com cinco disparos por segundo os abalos quase não se sobrepõem, e a força salta de $0.2F_1$ a $1.1F_1$; com quinze, ela já se mantém entre $1.8F_1$ e $2.2F_1$; com quarenta, fica quase uniforme, em torno de $5.4F_1$. O músculo real satura com disparos frequentes, de modo que a soma linear~\eqref{eq:forcesum} só vale até algumas dezenas de disparos por segundo.

\begin{figure}[t]
\centering
\begin{tikzpicture}[>=Stealth,x=20cm,y=0.55cm]
\draw[->,gray!70!black] (0,0) -- (0.66,0) node[right,font=\small,black] {$t$, s};
\draw[->,gray!70!black] (0,0) -- (0,6.4) node[left,font=\small,black] {$F/F_1$};
\foreach \t in {0,0.2,0.4,0.6}{\draw[gray!60!black] (\t,-0.1) -- (\t,0.1); \node[below,font=\scriptsize] at (\t,0) {\t};}
\foreach \f in {1,3,5}{\draw[gray!60!black] (-0.004,\f) -- (0.004,\f); \node[left,font=\scriptsize] at (0,\f) {\f};}
\draw[thick,blue!60!black] plot file {figures/muscle_twitch_5.dat};
\draw[thick,green!45!black] plot file {figures/muscle_twitch_15.dat};
\draw[thick,red!70!black] plot file {figures/muscle_twitch_40.dat};
\node[font=\scriptsize,blue!60!black,anchor=west] at (0.605,0.9) {5 disp./s};
\node[font=\scriptsize,green!45!black,anchor=west] at (0.605,2.1) {15 disp./s};
\node[font=\scriptsize,red!70!black,anchor=west] at (0.605,5.4) {40 disp./s};
\end{tikzpicture}
\caption{O músculo como conversor digital-analógico: a força~\eqref{eq:forcesum} com disparos regulares de uma unidade motora, $T_c=50\ms$. Disparos raros dão abalos separados, frequentes se fundem numa força uniforme, assim como disparos frequentes se fundem numa concentração uniforme na fig.~\ref{fig:lowpass}.}
\label{fig:tetanus}
\end{figure}

O cérebro tem dois botões de força: a taxa de disparos de cada unidade e o número de unidades ligadas. E o segundo botão é muito engenhoso. As unidades de um músculo são diferentes: as mais fracas são cem vezes mais fracas que as mais fortes. Quando é preciso cada vez mais força, as unidades são ligadas \emph{em ordem de tamanho}, das pequenas para as grandes~\cite{Henneman1957}. Ninguém calcula essa ordem: os neurônios pequenos simplesmente se excitam com mais facilidade pela mesma entrada comum, de modo que a ordem de ativação é definida pelo próprio hardware.

O que isso dá? Suponha que cada unidade seguinte na ordem seja $q$ vezes mais forte que a anterior, com $q$ um pouco maior que um. A força de todas as unidades ligadas é a soma de uma progressão geométrica, e quase toda ela cabe às últimas unidades. Então a próxima unidade ligada acrescenta
\begin{empheq}[box=\keybox]{equation}
\frac{\Delta F}{F} \;\approx\; q-1 \;=\; \text{const} ,
\label{eq:sizeprinciple}
\end{empheq}
ou seja, o passo de força é proporcional à própria força. Num esforço fraco, o cérebro o dosa em porções pequenas; num forte, em porções grandes. É um \emph{número de ponto flutuante}: o expoente é dado por quantas unidades estão ligadas, e a mantissa, pela taxa dos seus disparos. É a mesma escala logarítmica dos sensores do capítulo~\ref{sec:sensors}: a entrada e a saída da máquina medem o mundo em frações relativas.

O ponto flutuante tem um outro lado. A dispersão da força também cresce proporcionalmente a ela: um movimento forte é inevitavelmente menos preciso que um fraco. Segundo uma hipótese influente, é justamente esse ruído dependente do sinal que explica por que os nossos movimentos são suaves: um comando suave dá a menor dispersão no ponto final~\cite{HarrisWolpert1998}.

\section{Puxar, mas não empurrar: dois fios por articulação}

O músculo só sabe puxar. Empurrar ele não pode, assim como um neurônio \nt{GLUT} não pode emitir um sinal negativo. A solução é conhecida do capítulo~\ref{sec:glutcomputer}: \emph{dois fios}. Cada articulação é atendida por um par de músculos que puxam em sentidos opostos (fig.~\ref{fig:antagonists}). Se as suas forças são $F_1$ e $F_2$ e o braço de alavanca é $\ell$, o torque e a rigidez da articulação são
\begin{empheq}[box=\keybox]{equation}
M \;=\; \ell\,(F_1-F_2), \qquad K \;\approx\; \kappa_m\,(F_1+F_2),
\label{eq:antagonists}
\end{empheq}
em que $\kappa_m$ é o coeficiente que liga a força do músculo à sua rigidez: um músculo tenso é mais elástico que um relaxado. A diferença das forças dá o torque, a soma dá a rigidez~\cite{Hogan1984}. Com dois fios o cérebro controla duas grandezas independentes: para onde girar a articulação e com que firmeza segurá-la. Quando é preciso segurar uma xícara sem derramar, contraímos os dois músculos ao mesmo tempo: o torque é o mesmo, a rigidez é maior, e um empurrão ao acaso desvia menos a mão.

\begin{figure}[t]
\centering
\begin{tikzpicture}[>=Stealth,
  nrn/.style={circle,draw,very thick,minimum size=0.95cm,inner sep=0pt,font=\scriptsize},
  inh/.style={-{Bar[width=7pt]},thick,red!70!black}]
% the joint: a bar on a pivot, two springs pulling down on either side
\fill[gray!30] (4.6,-0.25) -- (5.4,-0.25) -- (5,0.15) -- cycle;
\draw[line width=3pt,gray!50!black] (2.4,0.2) -- (7.6,0.2);
\fill[black] (5,0.2) circle (0.08);
\draw[decorate,decoration={coil,segment length=5pt,amplitude=4pt},very thick,blue!60!black] (3.0,0.2) -- (3.0,-1.6);
\draw[decorate,decoration={coil,segment length=5pt,amplitude=4pt},very thick,orange!85!black] (7.0,0.2) -- (7.0,-1.6);
\draw[very thick] (2.5,-1.6) -- (3.5,-1.6); \draw[very thick] (6.5,-1.6) -- (7.5,-1.6);
\node[font=\small,blue!60!black] at (2.35,-0.75) {$F_1$};
\node[font=\small,orange!85!black] at (7.65,-0.75) {$F_2$};
\draw[->,thick] (5.9,0.75) arc (60:120:1.8) node[above,font=\scriptsize,pos=0.5] {$M=\ell(F_1-F_2)$};
% motor neurons and reflex circuit
\node[nrn,fill=green!12] (M1) at (1.6,-3.0) {\nt{ACET}};
\node[nrn,fill=green!12] (M2) at (8.4,-3.0) {\nt{ACET}};
\node[nrn,fill=red!10] (G) at (5.0,-3.0) {\nt{GLYC}};
\node[nrn,fill=blue!6] (S) at (1.6,-5.0) {\nt{GLUT}};
\draw[->,very thick] (M1) -- (3.0,-1.7);
\draw[->,very thick] (M2) -- (7.0,-1.7);
\draw[->,thick] (S) -- (M1);
\draw[->,thick] (S) -- (G);
\draw[inh] (G) -- (M2);
\draw[->,thick,densely dashed,gray!60!black] (3.15,-1.2) to[bend left=25] node[pos=0.35,right,font=\scriptsize,black] {estiramento} (S.north east);
\draw[->,thick,blue!60!black] (-0.3,-3.0) node[left,font=\scriptsize,black,align=right] {comando\\do cérebro} -- (M1);
\draw[->,thick,blue!60!black] (10.3,-3.0) node[right,font=\scriptsize,black,align=left] {comando\\do cérebro} -- (M2);
\end{tikzpicture}
\caption{Dois músculos numa articulação e o laço de realimentação mais rápido. Os neurônios \nt{ACET} recebem os comandos do cérebro e puxam a articulação em sentidos opostos; a diferença das forças a gira, a soma a torna mais rígida. O sensor de estiramento (\nt{GLUT}) do músculo esquerdo excita o seu próprio neurônio \nt{ACET} e, por meio de um intermediário \nt{GLYC}, inibe o neurônio do músculo antagonista: o veto do capítulo~\ref{sec:gaba}.}
\label{fig:antagonists}
\end{figure}

\section{Realimentação com atraso}

Os músculos não trabalham às cegas. Em cada um há sensores de estiramento, mencionados na tabela~\ref{tab:sensors}, e o laço mais curto se fecha diretamente na medula espinhal (fig.~\ref{fig:antagonists}). O músculo foi estirado de repente: o sensor de estiramento excita o seu próprio neurônio \nt{ACET}, o músculo se contrai e devolve a articulação. Ao mesmo tempo, por um neurônio intermediário inibitório, o músculo antagonista é desligado para não atrapalhar. Esse intermediário é \nt{GLYC}, justamente o tipo da tabela~\ref{tab:types} que não ganhou capítulo próprio: o seu lugar é exatamente aqui, nos laços rápidos da medula.

Cada laço tem um atraso $d$: no espinhal, cerca de $25$--$30\ms$ para o braço; no laço pelo córtex, uma vez e meia a três vezes mais; no laço pelo olho, $100$--$200\ms$. E a realimentação atrasada, como vimos na proposição~\ref{prop:ei}, faz o sistema oscilar. Para o regulador mais simples, que corrige o desvio $x$ com velocidade $G$ usando informação de $d$ segundos atrás, $\dd x/\dd t=-G\,x(t-d)$, a condição de estabilidade é~\cite{Hayes1950}:
\begin{empheq}[box=\keybox]{equation}
G\,d \;<\; \frac{\pi}{2} ,
\label{eq:delaystable}
\end{empheq}
e no limite o sistema oscila com frequência $1/(4d)$. Verificamos isso num modelo: com $d=30\ms$, o desvio se amortece com $G=40$ por segundo e, com $G=55$ por segundo, um pouco acima do limite de $52$, oscila cada vez mais.

Daí duas consequências. Primeira: quanto mais longo o laço, mais devagar ele é obrigado a corrigir. Pelo olho, com um atraso de cento e cinquenta milissegundos, não se pode corrigir a mão mais rápido que em cerca de um décimo de segundo, senão ela vai tremer. Segunda: o limite de estabilidade do laço espinhal, $1/(4\cdot30\ms)\approx8$ oscilações por segundo, é próximo da frequência do tremor fino que toda pessoa saudável tem: de oito a doze oscilações por segundo. O laço de estiramento é uma das fontes prováveis desse tremor~\cite{McAuleyMarsden2000}.

Então como o cérebro se move rápido e com precisão, se a realimentação chega atrasada? Segundo uma hipótese difundida, ele \emph{prevê}: mantém um modelo interno do corpo e calcula qual será o resultado do comando sem esperar pelos sensores~\cite{WolpertGhahramani2000}. Os sensores servem para corrigir o modelo, e não cada movimento. Um engenheiro chamaria isso de controle antecipativo com observador de estado.

\section{Geradores de ritmo dos movimentos}

Andar, respirar, mastigar são movimentos rítmicos. Eles precisam de um relógio? Não. Na medula e no tronco há pequenas redes que geram ritmo sozinhas, mesmo que nada rítmico chegue a elas~\cite{MarderBucher2001}. A rede mais simples desse tipo é montada com peças que já temos: dois grupos de neurônios \nt{GLUT} que se inibem mutuamente por intermediários \nt{GABA} ou \nt{GLYC}, como na fig.~\ref{fig:wta}, e se cansam lentamente, como a memória do capítulo~\ref{sec:glutcomputer}:
\begin{empheq}[box=\keybox]{equation}
\tau\,\frac{\dd r_i}{\dd t} = -r_i + \bigl[I - \beta\,r_j - a_i\bigr]_+ ,
\qquad
\tau_a\,\frac{\dd a_i}{\dd t} = -a_i + b\,r_i ,
\label{eq:halfcentre}
\end{empheq}
em que $a_i$ é a fadiga do grupo $i$, e $j$ é o outro grupo. A inibição mútua com $\beta>1$ escolhe um vencedor (proposição~\ref{prop:wta}). O vencedor trabalha e se cansa; quando a fadiga consome a sua entrada, o perdedor, já descansado, toma a frente e começa ele próprio a se cansar. Os grupos trabalham alternadamente, e cada um controla o seu músculo do par (fig.~\ref{fig:halfcentre}).

\begin{figure}[t]
\centering
\begin{tikzpicture}[>=Stealth,x=3.2cm,y=2.4cm]
\draw[->,gray!70!black] (0,0) -- (4.2,0) node[right,font=\small,black] {$t$, s};
\draw[->,gray!70!black] (0,0) -- (0,1.2) node[left,font=\small,black] {$r$};
\foreach \t in {0,1,2,3,4}{\draw[gray!60!black] (\t,-0.03) -- (\t,0.03); \node[below,font=\scriptsize] at (\t,0) {\t};}
\draw[thick,blue!60!black] plot file {figures/halfcentre1.dat};
\draw[thick,orange!85!black] plot file {figures/halfcentre2.dat};
\node[font=\scriptsize,blue!60!black,anchor=west] at (1.6,1.28) {grupo 1: ``para a frente''};
\node[font=\scriptsize,orange!85!black,anchor=west] at (2.7,1.28) {grupo 2: ``para trás''};
\end{tikzpicture}
\caption{Gerador de ritmo segundo o modelo~\eqref{eq:halfcentre}: $I=1$, $\beta=2$, $b=2$, $\tau=20\ms$, $\tau_a=0.4\,\text{s}$. Dois grupos com inibição mútua e fadiga trabalham alternadamente com período de $0.84\,\text{s}$, embora a entrada receba um sinal constante.}
\label{fig:halfcentre}
\end{figure}

No nosso modelo, com entrada constante, os grupos se revezam com período de $0.84\,\text{s}$, cerca do dobro do tempo de fadiga $\tau_a$. Um comando constante vindo de cima, ``ande'', vira um ritmo no local. É mais um ritmo local, como o ritmo do laço \nt{GLUT}--\nt{GABA} no capítulo~\ref{sec:gaba}: ele surge só quando a rede trabalha e marca o compasso apenas dos músculos que controla.

\begin{keyblock}
\begin{proposition}[A saída da máquina]
\label{prop:muscles}
O cérebro controla os músculos como uma hierarquia de reguladores. No topo, escolhe-se a ação (capítulo~\ref{sec:dofa}); abaixo, geradores locais transformam comandos constantes em ritmo, e os laços rápidos da medula seguram as articulações. A força é dada em ponto flutuante: o expoente, pelo número de unidades ligadas por ordem de tamanho; a mantissa, pela taxa de disparos. Cada articulação é controlada por um par de músculos que puxam, cuja diferença de forças dá o torque e cuja soma dá a rigidez. Esses dispositivos foram medidos; que o cérebro se apoie num modelo interno do corpo é uma hipótese bem fundamentada.
\end{proposition}
\end{keyblock}

\section{Resumo}

\begin{table}[t]
\centering
\footnotesize
\setlength{\tabcolsep}{4pt}
\renewcommand{\arraystretch}{1.15}
\begin{tabular}{>{\raggedright}p{3.0cm}>{\raggedright}p{4.6cm}>{\raggedright\arraybackslash}p{5.2cm}}
\toprule
O que a saída faz & Como & O que se sabe \\
\midrule
transforma disparos em força & soma de abalos~\eqref{eq:forcesum}, filtro passa-baixas & medido; a soma linear é uma simplificação \\
dosa a força & ativação das unidades por tamanho, ponto flutuante~\eqref{eq:sizeprinciple} & ordem de ativação medida \\
empurra sabendo só puxar & par de músculos: torque é a diferença, rigidez é a soma~\eqref{eq:antagonists} & medido \\
segura a articulação & laço de estiramento, veto do antagonista via \nt{GLYC} & medido \\
não treme & $Gd<\pi/2$~\eqref{eq:delaystable} & decorre da teoria de controle; contribuição ao tremor é causa provável \\
move-se rápido & predição por um modelo interno & hipótese bem fundamentada \\
anda e respira com ritmo & gerador de ritmo~\eqref{eq:halfcentre} & geradores medidos; o esquema mais simples é modelo \\
\bottomrule
\end{tabular}
\caption{Como a máquina controla os músculos.}
\label{tab:muscles}
\end{table}

A saída da máquina é construída com os mesmos truques de todo o resto (tabela~\ref{tab:muscles}): dois fios em vez de sinal, um filtro passa-baixas em vez de DAC, escala logarítmica em vez de linear, inibição mútua e fadiga em vez de um gerador de relógio. A entrada (capítulo~\ref{sec:sensors}) e a saída medem o mundo em frações relativas, e entre elas trabalha a máquina a que este livro é dedicado.

\section{Exercícios}

\begin{exercise}[\lvlA]
\label{ex:13:1}
\emph{Ponto flutuante em números.} Um músculo tem $N=100$ unidades motoras com forças $f_n=f_1q^{\,n-1}$; a mais forte é $100$ vezes mais forte que a mais fraca. (a)~Ache $q$, o passo relativo~\eqref{eq:sizeprinciple} e a faixa dinâmica em decibéis. (b)~Qual é o passo com força $10f_1$ num ``DAC linear'' de $100$ unidades iguais com a mesma força total?
\end{exercise}

\begin{exercise}[\lvlA]
\label{ex:13:2}
\emph{O preço da margem de segurança.} A cada disparo, a sinapse numa fibra muscular libera um número de Poisson de quanta com média $m$, e a fibra responde com $n_{\text{th}}=20$ quanta ou mais; a margem de segurança é $S=m/n_{\text{th}}$. Quantas falhas por dia dá uma unidade que trabalha a $20$ disparos por segundo, com $S=2$ e com $S=4$?
\end{exercise}

\begin{exercise}[\lvlB]
\label{ex:13:3}
\emph{O músculo como filtro.} O abalo~\eqref{eq:twitch} com $T_c=50\ms$ é a resposta ao impulso de um sistema linear. (a)~Ache a sua função de transferência $H(s)$ e a frequência de corte. (b)~Com que taxa de disparos regulares a amplitude das pulsações de força é $10\%$ da força média?
\end{exercise}

\begin{exercise}[\lvlB]
\label{ex:13:4}
\emph{Mais rápido que o atraso não se corrige.} Regulador $\dd x/\dd t=-G\,x(t-d)$. (a)~Mostre que o desvio se amortece mais rápido sem oscilar com $Gd=1/e$, com taxa $1/d$. (b)~Ache esse $G$ e a constante de tempo do amortecimento para o laço espinhal, $d=30\ms$, e para o laço pelo olho, $d=150\ms$.
\end{exercise}

\begin{exercise}[\lvlB]
\label{ex:13:5}
\emph{O período do gerador de ritmo.} Com $\tau\ll\tau_a$, o período $T$ do modelo~\eqref{eq:halfcentre} é dado pela equação $(\beta-1)(1-E^{2+b})/(\beta-E)=b\,(1-E^{1+b})/(1+b)$, com $E=e^{-T/(2\tau_a)}$. (a)~Ache $T$ com $\beta=b=2$, $\tau_a=0.4$~s. (b)~Mostre que $T$ não depende da entrada $I$ e que o ritmo só é possível com $b>\beta-1$.
\end{exercise}

\begin{exercise}[\lvlB]
\label{ex:13:6}
\emph{Simulação do gerador.} Importe a função \texttt{halfcentre} de \texttt{models/\allowbreak muscle\_models.py} (não execute o próprio arquivo) e meça o período, descartando os dois primeiros ciclos, com $b=0.8$, $1.05$, $1.2$, $1.5$, $2$, $4$ e com $I=0.5$, $1$, $2$. O comando ``ande mais rápido'' pode mudar o passo através de $I$?
\end{exercise}

\begin{exercise}[\lvlB]
\label{ex:13:7}
\emph{Rigidez contra reflexo.} Uma articulação está num laço de estiramento: $\dd\theta/\dd t=-a\theta(t)-G\theta(t-d)$, $a=K/B$, $d=30\ms$. Reduza-o ao exercício~\ref{ex:13:4} e ache o menor $a$ com que o desvio se amortece sem oscilar com constante de tempo de $15\ms$, e o $G$ necessário. Como mudar $G$ ao aumentar a co-contração dos músculos?
\end{exercise}

\begin{exercise}[\lvlC]
\label{ex:13:8}
\emph{O botão do andamento.} Pelos exercícios~\ref{ex:13:5} e~\ref{ex:13:6}, a entrada $I$ do gerador~\eqref{eq:halfcentre} muda só a amplitude, e não o andamento. Proponha a menor mudança do modelo, com peças do livro, em que abaixo de um certo $I_{\min}$ não há ritmo e acima dele o período diminui com $I$, pelo menos à metade quando $I$ triplica. Ache $I_{\min}$ com $\tau\ll\tau_a$ e verifique por simulação.
\end{exercise}

\chapter{A dor: um sinal de alarme}
\chaptermark{A dor}
\label{sec:pain}
\begin{chaptergoals}
\item por que a dor é um sinal de alarme de alta prioridade, e não um canal de medição;
\item como um fio rápido e um lento informam onde está o dano e quão grave ele é;
\item como a porta na medula deixa o tato vetar a dor, e como a difusão define o volume do alarme;
\item como a sensibilização eleva o ganho após uma lesão, e como a dor ensina e interrompe.
\end{chaptergoals}

\section{Alarme, e não medição}

Todos os sensores do capítulo~\ref{sec:sensors} mediam o mundo: quanta luz, que altura tem o som, que pressão. A dor funciona de outro modo. Ela não informa ``o que há lá'', e sim ``perigo, largue tudo''. Para o engenheiro, não é um canal de medição, mas um \emph{sinal de alarme}: uma linha de prioridade máxima, que precisa atravessar qualquer trabalho em curso da máquina.

Um alarme difere de um canal de medição em três propriedades, e a dor tem as três. Primeiro, é difícil abafá-la por adaptação: os sensores de dor se adaptam muito menos que os sensores de toque, e depois de uma lesão leve ficam até mais sensíveis~\cite{LaMotte1982hyper}; o enfraquecimento da resposta após um aquecimento forte foi medido em apenas um dos seus tipos~\cite{Treede1995heat}. O sensor de luz se cala sobre um fundo constante (fig.~\ref{fig:adaptation}), e um alarme que se cala depois de um minuto é inútil. Segundo, ela tem limiar alto: os sensores de dor só disparam quando o estímulo está perto de ser perigoso; para o calor, por exemplo, os limiares dos diferentes sensores vão de $41^\circ$ a $46$--$48^\circ$, conforme o tipo~\cite{Treede1995heat, Weidner1999cfib}. Terceiro, e isso é o principal para este capítulo, o volume do alarme não é decidido só pelo sensor, mas pela própria máquina. A mesma ferida dói de modo diferente em situações diferentes. Vejamos com que peças já conhecidas isso é montado.

Os sensores de dor são neurônios \nt{GLUT}, como os demais neurônios sensoriais do capítulo~\ref{sec:sensors}. Parte deles libera, junto com o glutamato, a substância~P do apêndice~\ref{sec:othermediators}, que reforça lentamente a transmissão.

\section{Dois fios: a dor rápida e a lenta}

Bata o dedo, e você sentirá a dor duas vezes: primeiro curta e aguda, depois de um segundo, surda e longa. São dois fios diferentes. O fio rápido é isolado e conduz disparos com velocidade $v$ de $5$ a $30$ metros por segundo; o lento é fino e nu, $0.5$--$2$ metros por segundo. O sinal do pé percorre cerca de um metro, e o tempo de chegada
\begin{empheq}[box=\keybox]{equation}
t \;=\; \frac{L}{v}
\label{eq:paintime}
\end{empheq}
é de dezenas de milissegundos para o fio rápido e de cerca de um segundo para o lento. Os dois fios levam duas mensagens. O rápido diz \emph{onde} e \emph{quando}: os seus disparos chegam antes e com mais precisão, como no código temporal do capítulo~\ref{sec:code}. O lento diz \emph{quão ruim é}, e a sua dor longa e surda mantém a máquina em modo de alarme enquanto a lesão não passa.

\section{A porta na medula espinhal}

O primeiro lugar aonde chega o sinal de dor é a medula espinhal, e lá ele passa por uma \emph{porta}~\cite{MelzackWall1965}; os neurônios concretos dessa porta foram encontrados relativamente há pouco tempo~\cite{Duan2014}. No neurônio \nt{GLUT} de saída, que transmite a dor adiante para o cérebro, convergem o fio da dor e o fio do toque. E ao lado fica um neurônio intermediário inibitório, \nt{GABA} ou \nt{GLYC}, que o toque excita (fig.~\ref{fig:gate}). O toque fecha a porta. Na teoria original da porta~\cite{MelzackWall1965}, a dor, ao contrário, desliga esse intermediário, e uma dor forte abre a porta; isso é uma suposição da teoria e não foi mostrado diretamente nos neurônios da porta já encontrados.

\begin{figure}[t]
\centering
\begin{tikzpicture}[>=Stealth,
  nrn/.style={circle,draw,very thick,minimum size=1.0cm,inner sep=0pt,font=\scriptsize},
  inh/.style={-{Bar[width=7pt]},thick,red!70!black},
  bc/.style={->,thick,densely dashed,orange!85!black}]
\node[nrn,fill=blue!6] (Z) at (6,0) {\nt{GLUT}};
\node[nrn,fill=red!10] (Q) at (3.6,-1.6) {\nt{GABA}};
\draw[->,very thick,red!70!black] (0,0.6) node[left,font=\small,black,align=right] {dor $\nu$} -- (5.45,0.15);
\draw[->,very thick,blue!60!black] (0,-2.6) node[left,font=\small,black,align=right] {toque $\mu$} -- (2.9,-2.0);
\draw[->,thick,blue!60!black] (1.8,-2.35) -- (5.5,-0.35) node[pos=0.8,below right=-2pt,font=\scriptsize] {$w_{z\mu}$};
\draw[inh] (1.6,0.5) -- (3.35,-1.1) node[pos=0.5,left=1pt,font=\scriptsize] {$w_{q\nu}$};
\draw[inh] (Q) -- (Z) node[pos=0.5,above left=-1pt,font=\scriptsize] {$\beta$};
\draw[->,very thick] (Z) -- (8.2,0) node[right,font=\small,align=left] {ao cérebro: $z$};
\draw[bc] (6,2.1) node[above,font=\scriptsize,black,align=center] {de cima: \nt{SERO}, \nt{NORA}, opioides} -- (Z.north) node[pos=0.45,right,font=\scriptsize,black] {ganho $g$};
\end{tikzpicture}
\caption{A porta da dor na medula espinhal segundo o modelo~\eqref{eq:gate}. O neurônio \nt{GLUT} de saída recebe a dor $\nu$ e uma excitação fraca do toque $\mu$. O intermediário inibitório (\nt{GABA} ou \nt{GLYC}) é excitado pelo toque e, pela suposição da teoria da porta, desligado pela dor. De cima, pelos canais de difusão, chega o ganho $g$.}
\label{fig:gate}
\end{figure}

Vamos escrever isso para taxas, como no capítulo~\ref{sec:gaba}. A taxa do intermediário $q$ e a taxa de saída $z$ são
\begin{empheq}[box=\keybox]{equation}
q \;=\; \bigl[w_{q\mu}\,\mu + q_0 - w_{q\nu}\,\nu\bigr]_+ ,
\qquad
z \;=\; \bigl[\,g\,(\nu + w_{z\mu}\,\mu) - \beta\,q\,\bigr]_+ ,
\label{eq:gate}
\end{empheq}
em que $\nu$ é a entrada de dor, $\mu$ é a entrada de toque, $w$ são os pesos das conexões (o primeiro índice é o destino, o segundo, a origem), $\beta$ é a força da inibição, $q_0$ é a atividade de fundo própria do intermediário, e $g$ é o ganho definido de cima (falaremos dele adiante). É o veto~\eqref{eq:veto} do capítulo~\ref{sec:gaba}: ``dor, mas não com toque'', com uma correção tomada da teoria da porta como suposição: uma dor forte desliga ela mesma o neurônio que veta.

Calculamos o modelo com $w_{q\mu}=1$, $q_0=0.2$, $w_{q\nu}=0.5$, $w_{z\mu}=0.3$, $\beta=1.5$ (fig.~\ref{fig:gatecurves}). Sem toque, a dor começa a passar a partir de $\nu\approx0.18$. Com toque, o limiar sobe para $0.86$, e com $\nu=1$ a saída cai a um quarto, de $1$ para $0.25$. Já com uma dor forte, $\nu=2$, o toque não muda mais nada: a porta está escancarada. Assim é também na vida: esfregar uma pancada ajuda, mas numa lesão grave, não. O toque sozinho, sem dor, não passa pela porta: o alarme não dispara com um simples contato.

\begin{figure}[t]
\centering
\begin{tikzpicture}[>=Stealth,x=5.2cm,y=1.35cm]
\draw[->,gray!70!black] (0,0) -- (2.12,0) node[right,font=\small,black] {dor $\nu$};
\draw[->,gray!70!black] (0,0) -- (0,4.3) node[left,font=\small,black] {$z$};
\foreach \t in {0,0.5,1,1.5,2}{\draw[gray!60!black] (\t,-0.05) -- (\t,0.05); \node[below,font=\scriptsize] at (\t,0) {\t};}
\foreach \f in {1,2,3,4}{\draw[gray!60!black] (-0.01,\f) -- (0.01,\f); \node[left,font=\scriptsize] at (0,\f) {\f};}
\draw[very thick,red!70!black] plot file {figures/pain_gate_notouch.dat};
\draw[very thick,blue!60!black] plot file {figures/pain_gate_touch.dat};
\draw[very thick,orange!85!black,densely dashed] plot file {figures/pain_gate_sens.dat};
\draw[very thick,red!70!black] (0.08,3.9) -- (0.2,3.9) node[right,font=\scriptsize,black] {sem toque};
\draw[very thick,blue!60!black] (0.08,3.45) -- (0.2,3.45) node[right,font=\scriptsize,black] {com toque};
\draw[very thick,orange!85!black,densely dashed] (0.08,3.0) -- (0.2,3.0) node[right,font=\scriptsize,black] {após sensibilização, $g=2$};
\end{tikzpicture}
\caption{A saída da porta~\eqref{eq:gate} em função da força da dor. O toque eleva o limiar e abafa a dor fraca e média, mas não a forte. A sensibilização (linha tracejada) dobra o ganho: a mesma dor é transmitida com o dobro do volume, e o limiar desce.}
\label{fig:gatecurves}
\end{figure}

\section{O botão de volume vindo de cima}

A porta não é controlada só por baixo. Do tronco encefálico descem até ela vias que mudam o ganho $g$ em~\eqref{eq:gate}: \nt{SERO}, \nt{NORA} e os peptídeos opioides do apêndice~\ref{sec:othermediators}~\cite{Fields2004}. São os mesmos canais de difusão dos capítulos~\ref{sec:serocomputer}--\ref{sec:acet}, só que aqui o seu botão é o volume do alarme. Eles sabem girá-lo nos dois sentidos: os mesmos circuitos descendentes podem tanto abafar a dor quanto amplificá-la.

Por que a máquina baixaria o alarme? Porque o alarme é uma interrupção, e uma interrupção nem sempre é oportuna. Um animal que foge de um predador não deve parar por causa de uma pata ferida: primeiro correr, depois lamber a ferida. A expectativa de alívio também abafa a dor, e parte desse efeito some se os receptores opioides forem bloqueados~\cite{Levine1978}: a expectativa calculada no cérebro desce pelo canal opioide até a porta. Esse quadro em si foi medido; como exatamente o cérebro decide qual deve ser o volume é hipótese.

\section{Sensibilização: o alarme que aprende}

Depois de uma lesão, os neurônios de saída da porta ficam mais sensíveis: a mesma entrada provoca uma saída maior, e o limiar desce~\cite{Woolf1983}. No modelo~\eqref{eq:gate} isso é um aumento do ganho $g$, e a descrição mais simples dele é um circuito RC lento carregado pela própria dor:
\begin{empheq}[box=\keybox]{equation}
\tau_s\,\frac{\dd g}{\dd t} \;=\; -(g-1) + k_s\,z ,
\label{eq:sensitization}
\end{empheq}
em que $\tau_s$ é o tempo para a sensibilidade voltar ao normal; ele é maior que a duração da própria dor, porque a sensibilização persiste depois que o estímulo lesivo cessou~\cite{Woolf1983}; e $k_s$ é a força da carga. Enquanto o tecido está lesado, a saída da porta mantém $g$ acima de um, e tudo o que acontece perto da ferida é transmitido mais alto (linha tracejada na fig.~\ref{fig:gatecurves}: com $g=2$, o limiar desce para $0.11$ e a saída dobra). É um ajuste razoável do alarme: enquanto a lesão não sarou, é melhor pecar pelo excesso de cautela.

Para o engenheiro, é uma realimentação positiva com fuga lenta: a dor aumenta o ganho, o ganho aumenta a dor. Enquanto o laço é fraco, $g$ volta ao normal quando a ferida sara. Mas se o laço é forte, o alarme pode travar no estado ligado mesmo depois que a lesão não existe mais, como o grupo com realimentação forte da fig.~\ref{fig:bistable}. O modelo~\eqref{eq:sensitization} é uma simplificação; a existência da sensibilização central foi medida.

\section{A dor como professor e como interrupção}

Na máquina, a dor tem dois trabalhos além do alarme.

\emph{Professor.} A dor é uma recompensa negativa: na fórmula do erro~\eqref{eq:ctd}, ela entra como $r<0$. Tudo o que foi dito no capítulo~\ref{sec:dofa} vale aqui também: um prenúncio da dor passa a provocar o erro com antecedência, e a rede aprende a evitar o que leva à dor. Experimentos em humanos mostram que a atividade cerebral no aprendizado com dor segue justamente o erro de diferença temporal~\cite{Seymour2004}, e parte dos neurônios \nt{DOPA} responde também a eventos desagradáveis (capítulo~\ref{sec:dofaalt}).

\emph{Interrupção.} A dor arranca a atenção da tarefa em curso: essa é a sua função, e não um efeito colateral~\cite{EcclestonCrombez1999}. No capítulo~\ref{sec:nora}, o mesmo papel de ``interrupção'' foi discutido para a rajada fásica de \nt{NORA}. E a resposta mais rápida nem chega ao cérebro: retirar a mão de algo quente se fecha diretamente na medula, com o mesmo par de músculos e o mesmo veto do antagonista da fig.~\ref{fig:antagonists}.

\begin{keyblock}
\begin{proposition}[Sinal de alarme]
\label{prop:pain}
A dor não é uma medição, e sim um sinal de alarme de alta prioridade. Ela se adapta muito menos que os canais de medição, segue por dois fios (o rápido, ``onde''; o lento, ``quão ruim''), passa por uma porta que o toque fecha (e que uma dor forte, pela suposição da teoria da porta, abre), e o seu volume é definido de cima pelos canais de difusão. Esses dispositivos foram medidos, exceto a abertura da porta pela dor; os modelos simples~\eqref{eq:gate} e~\eqref{eq:sensitization} são simplificações.
\end{proposition}
\end{keyblock}

\section{Resumo}

\begin{table}[t]
\centering
\footnotesize
\setlength{\tabcolsep}{4pt}
\renewcommand{\arraystretch}{1.15}
\begin{tabular}{>{\raggedright}p{3.0cm}>{\raggedright}p{4.9cm}>{\raggedright\arraybackslash}p{4.9cm}}
\toprule
O que a dor faz & Como & O que se sabe \\
\midrule
dispara o alarme & limiar alto, adaptação menor que a do toque & limiar medido; a comparação da adaptação é estimativa \\
informa ``onde'' e ``quão ruim'' & dois fios de velocidades diferentes~\eqref{eq:paintime} & medido \\
passa pela porta & veto via \nt{GABA}/\nt{GLYC}~\eqref{eq:gate} & porta medida; abertura pela dor é suposição da teoria; o modelo é simplificação \\
muda o volume & \nt{SERO}, \nt{NORA}, opioides descendentes & medido; a regra de escolha do volume é hipótese \\
aprende após a lesão & aumento do ganho~\eqref{eq:sensitization} & sensibilização medida; o modelo é simplificação \\
ensina a evitar & recompensa negativa em~\eqref{eq:ctd} & semelhança com o erro de diferença temporal medida \\
interrompe o trabalho & captura da atenção; reflexo de retirada & medido \\
\bottomrule
\end{tabular}
\caption{A dor como sinal de alarme.}
\label{tab:pain}
\end{table}

A dor é montada com peças que já vimos (tabela~\ref{tab:pain}): sensores \nt{GLUT}, dois fios, veto por um intermediário inibitório, um botão de volume por difusão, um circuito RC lento, o erro de predição e a interrupção. Só uma coisa nela é nova: é a única entrada da máquina cujo volume a própria máquina define, um alarme que o dono pode tanto abafar quanto reconfigurar.

\section{Exercícios}

\begin{exercise}[\lvlA]
\label{ex:14:1}
\emph{Dois fios.} O sinal vem do pé, $L=1$~m. (a)~Uma salva percorre um feixe de fios de um só tipo, cujas velocidades preenchem a faixa~\eqref{eq:paintime}. Quanto ela se espalha até chegar à medula, para o tipo rápido e o lento? (b)~As ondas de dor pelos fios de $15$ e $1$~m/s se distinguem com diferença a partir de $0.1$~s. A partir de que distância elas se fundem?
\end{exercise}

\begin{exercise}[\lvlB]
\label{ex:14:2}
\emph{Quando o toque dói.} Porta~\eqref{eq:gate} com os parâmetros do texto; o ganho $g$ cresce com a sensibilização. Até que $g$ o toque sem dor não passa pela porta com nenhuma força $\mu$? Com que $g$ o toque $\mu=1$ começa a passar?
\end{exercise}

\begin{exercise}[\lvlB]
\label{ex:14:3}
\emph{Estabilidade da sensibilização.} As entradas são constantes, e a saída da porta é linear em $g$: $z=gA-C$, $A=\nu+w_{z\mu}\mu$, $C=\beta q\ge0$. (a)~Ache o equilíbrio $g^*$ do modelo~\eqref{eq:sensitization} e a condição da sua estabilidade. (b)~Com $k_s=0.5$, ache a força da dor acima da qual o modelo dispara, sem toque e com toque $\mu=1$.
\end{exercise}

\begin{exercise}[\lvlA]
\label{ex:14:4}
\emph{O prenúncio da dor.} Um sinal no instante $0$ prediz dor no instante $T=2$~s; em~\eqref{eq:ctd}, $r(t)=-P\,\delta(t-T)$, $\tau_\gamma=2$~s. (a)~Ache a área do pico de $\delta$ no instante do sinal. (b)~Uma ação no instante $t_e=1$~s cancela a dor. Qual é o pico de $\delta$ nesse instante e o que ele ensinará à rede?
\end{exercise}

\begin{exercise}[\lvlB]
\label{ex:14:5}
\emph{O alarme que trava.} Complete a porta de \texttt{models/\allowbreak pain\_gate.py} com a dinâmica~\eqref{eq:sensitization} e a saturação $z\le4$. O toque $\mu=1$ está sempre presente, a lesão $\nu=2$ dura um tempo $T$, depois $\nu=0$. Por simulação, ache o menor $T$ (em unidades de $\tau_s$) depois do qual o alarme continua ligado, com $k_s=4$, $4.6$, $5$, $6$, $8$.
\end{exercise}

\begin{exercise}[\lvlB]
\label{ex:14:6}
\emph{Projeto da porta.} Com $\beta=1.5$, $w_{z\mu}=0.3$, $g=1$, escolha em~\eqref{eq:gate} $q_0$, $w_{q\nu}$, $w_{q\mu}$ de modo que o limiar de dor seja $0.2$ sem toque e $1.0$ com toque $\mu=1$, e que o toque deixe de abafar a dor em $\nu_\times=2.5$. Até que ganho $g$ o toque sem dor não passa pela porta?
\end{exercise}

\begin{exercise}[\lvlC]
\label{ex:14:7}
\emph{O botão de volume.} O trabalho rende um valor $V$ por unidade de tempo, trabalhar com a lesão $\nu$ custa $c\nu$, de modo que interromper compensa com $\nu>V/c$; a dor interrompe com $z>\theta=0.5$. (a)~Que ganho $g^*$ da porta~\eqref{eq:gate} sem toque dá esse limiar com $V/c=0.25$, $0.5$, $2$? (b)~A partir de que sinais do livro a máquina poderia estimar $V$ e $c$?
\end{exercise}

\chapter{Como a máquina aprende a se mover}
\chaptermark{Como a máquina aprende a se mover}
\label{sec:motorlearning}
\begin{chaptergoals}
\item por que mover-se exige um terceiro professor, um fio de erro para cada saída, e o que isso custa;
\item como a camada de expansão \nt{GLUT} e um fio de erro por saída dão a regra dos mínimos quadrados;
\item como o circuito aprendido vira um modelo interno, um filtro adaptativo que prevê o corpo;
\item por que um processo rápido e outro lento explicam o pós-efeito e a recuperação espontânea.
\end{chaptergoals}

\section{Três professores}

No capítulo~\ref{sec:muscles} vimos como a máquina move os músculos e paramos numa hipótese: movimentos rápidos e precisos exigem um modelo interno do corpo, que prevê o resultado do comando~\cite{WolpertGhahramani2000}. De onde vem esse modelo? Ele precisa ser aprendido, e aprendido de um jeito diferente de tudo o que aprendemos até agora.

O livro já teve dois professores. O primeiro é \emph{nenhum}: a regra de Hebb (capítulo~\ref{sec:dofa}) reforça o que acontece junto com frequência e comprime as entradas. O segundo é a \emph{recompensa}: \nt{DOPA} difunde para todos um único número, ``melhor ou pior do que o esperado''. O movimento precisa de um terceiro professor: o \emph{erro}. A mão errou a xícara por dois centímetros à esquerda: isso não é ``pior'', é um vetor que diz a cada saída para onde e quanto ela errou. Um engenheiro chamaria isso de aprendizado supervisionado.

A diferença entre o segundo e o terceiro professor é a diferença entre um fio para todos e um fio próprio para cada saída. Na proposição~\ref{prop:broadcastcost}, o aprendizado por um único número comum ficava mais lento a cada neurônio cujo ruído influía no resultado. Mas se cada saída $i$ recebe o seu próprio erro $e_i$, os pesos podem mudar pela regra mais simples:
\begin{empheq}[box=\keybox]{equation}
\frac{\dd w_{ij}}{\dd t} \;=\; -\eta\,e_i(t)\,x_j(t) .
\label{eq:lms}
\end{empheq}
É a regra dos mínimos quadrados, conhecida há muito na técnica dos filtros adaptativos~\cite{WidrowHoff1960}: o peso muda proporcionalmente ao produto da sua entrada pelo erro da sua saída e, em média, segue o gradiente do erro quadrático. Não é preciso nenhum ruído de busca, como no capítulo~\ref{sec:dofa}: a direção da correção chega diretamente pelo fio de erro. E a velocidade não cai com o número de saídas, porque os erros alheios não chegam à sinapse.

\begin{keyblock}
\begin{proposition}[Três professores]
\label{prop:teachers}
A regra de Hebb ensina sem professor o que é frequente; o sinal de \nt{DOPA} ensina, com um número para a rede inteira, o que é vantajoso, e mais devagar a cada neurônio; o fio de erro para cada saída ensina, pela regra~\eqref{eq:lms}, o que é preciso, e a velocidade não depende do número de saídas. O preço do terceiro modo é um fio separado para cada saída e alguém que conheça a resposta certa.
\end{proposition}
\end{keyblock}

\section{O fio de erro}

Quem pode conhecer a resposta certa para um movimento? Os próprios sensores. Se a mão foi para a esquerda do alvo, o olho e os sensores de estiramento do capítulo~\ref{sec:sensors} informam isso. É preciso um circuito que leve esse erro às saídas certas.

O cérebro tem uma região que, ao que parece, é construída exatamente para isso~\cite{Marr1969,Albus1971}. O seu circuito é incomumente regular, quase cristalino, e nele é fácil reconhecer a regra~\eqref{eq:lms} (fig.~\ref{fig:errorwire}).

\emph{Camada de expansão da entrada.} Os sinais sobre o comando e sobre o estado do corpo chegam a uma enorme camada de pequenos neurônios \nt{GLUT}. São cerca de setenta bilhões, mais do que em todo o resto do cérebro somado~\cite{Azevedo2009}. Cada um mistura algumas entradas, e o sinal se decompõe num vetor muito longo. O engenheiro conhece esse truque: aumentando a dimensão da entrada, no novo espaço quase qualquer dependência desejada pode ser obtida com uma simples soma ponderada~\cite{Cover1965}.

\emph{Neurônios de saída.} A soma ponderada é calculada por cerca de trinta milhões de grandes neurônios de saída~\cite{Andersen1992cereb}, cada um com da ordem de cem mil entradas vindas da camada de expansão. E são neurônios \nt{GABA}: o resultado do aprendizado é passado adiante não por excitação, e sim por inibição, como no veto do capítulo~\ref{sec:gaba}.

\emph{O fio de erro.} Cada neurônio de saída recebe mais uma entrada, especial: exatamente um fio \nt{GLUT}, que dispara raramente, cerca de uma vez por segundo, mas a cada vez provoca no neurônio de saída uma descarga poderosa~\cite{Ito2001}. Esse é o $e_i$: a probabilidade da descarga depende da direção do erro de movimento~\cite{Herzfeld2018}. A coincidência da descarga de erro com a atividade de uma entrada enfraquece essa entrada: exatamente o termo $-\eta\,e_i x_j$ em~\eqref{eq:lms}.

\begin{figure}[t]
\centering
\begin{tikzpicture}[>=Stealth,
  gran/.style={circle,draw,thick,fill=blue!6,minimum size=0.32cm,inner sep=0pt},
  outn/.style={circle,draw,very thick,fill=red!10,minimum size=1.1cm,inner sep=0pt,font=\scriptsize},
  inh/.style={-{Bar[width=7pt]},very thick,red!70!black}]
% inputs
\foreach \y/\lab in {1.6/{comando},0.4/{estado do corpo}}{
  \draw[->,thick,blue!60!black] (-0.6,\y) node[left,font=\scriptsize,black] {\lab} -- (0.35,\y);}
% expansion layer
\foreach \i in {0,...,11}{\node[gran] (g\i) at ({0.8+0.28*mod(\i,2)},{-0.3+0.23*\i}) {};}
\foreach \i in {0,...,11}{\pgfmathsetmacro{\yy}{mod(\i,3)>0 ? 1.6 : 0.4}\draw[gray!60!black] (0.35,\yy) -- (g\i);}
\node[font=\scriptsize,align=center] at (0.95,-1.05) {camada de expansão\\$\sim7\cdot10^{10}$ \nt{GLUT}};
% parallel wires to the output neuron
\foreach \i in {0,...,11}{\draw[->,gray!70!black] (g\i.east) -- (4.1,{1.0+0.07*(\i-5.5)});}
\node[outn] (P) at (4.6,1.0) {\nt{GABA}};
\node[font=\scriptsize,align=center,above=2pt] at (P.north) {neurônio de saída\\$\sim10^5$ entradas $x_j$, pesos $w_{ij}$};
\draw[inh] (P) -- (6.8,1.0) node[right,font=\scriptsize,align=left] {correção\\do movimento};
% error wire
\draw[->,very thick,red!70!black,densely dashed] (4.6,-1.4) node[below,font=\scriptsize,black,align=center] {um fio de erro $e_i$\\\nt{GLUT}, $\sim1$ vez/s} -- (P.south);
\end{tikzpicture}
\caption{O esquema de aprendizado supervisionado como, ao que parece, o cérebro o realiza. O comando e o estado do corpo são decompostos por uma enorme camada de neurônios \nt{GLUT} num vetor longo $x_j$; o neurônio \nt{GABA} de saída o soma com os pesos $w_{ij}$. O único fio de erro traz $e_i$; a coincidência do erro com uma entrada ativa enfraquece o peso dessa entrada, como na regra~\eqref{eq:lms}.}
\label{fig:errorwire}
\end{figure}

Uma dificuldade nos é conhecida do capítulo~\ref{sec:dofa}: o erro chega depois da entrada que o causou. O desvio é visto dezenas a centenas de milissegundos depois do comando. Portanto, a sinapse precisa lembrar que esteve ativa: é preciso um traço, como na regra~\eqref{eq:threefactor}. E o comprimento desse traço, a julgar pelos experimentos, está ajustado ao atraso da realimentação em cada tarefa: onde o erro chega após cem milissegundos, a entrada mais enfraquecida é a que esteve ativa cem milissegundos antes dele~\cite{Suvrathan2016}.

\begin{keyblock}
\begin{proposition}[O fio de erro]
\label{prop:errorwire}
No circuito mostrado na fig.~\ref{fig:errorwire}, cada neurônio de saída recebe exatamente um fio de erro, e a coincidência do erro com uma entrada enfraquece essa entrada. É a regra dos mínimos quadrados~\eqref{eq:lms}, aplicada a uma entrada expandida até uma dimensão enorme. A construção do circuito e o enfraquecimento por coincidência foram medidos; que o circuito inteiro funcione como aprendizado supervisionado é a hipótese principal.
\end{proposition}
\end{keyblock}

\section{O modelo interno como filtro adaptativo}

O que esse circuito aprende? A resposta mais natural é: a prever. Suponha que chegue à entrada o comando $u(t)$, passado por um conjunto de filtros com diferentes constantes de tempo, como na camada de expansão: $x_j(t)=(g_j*u)(t)$. A saída é a sua soma ponderada, a predição do que os sensores vão informar:
\begin{empheq}[box=\keybox]{equation}
\hat y(t) \;=\; \sum_j w_j\,(g_j*u)(t), \qquad e(t) \;=\; \hat y(t) - y(t) ,
\label{eq:adaptivefilter}
\end{empheq}
e os pesos aprendem por~\eqref{eq:lms}. É um \emph{filtro adaptativo}: um circuito que ajusta sozinho a sua função de transferência para reproduzir um sistema desconhecido~\cite{Fujita1982}. Aqui o sistema desconhecido é o próprio corpo, com a sua massa, elasticidade e atrasos. O filtro aprendido é o modelo interno do capítulo~\ref{sec:muscles}: ele diz o que os sensores vão informar, sem esperar por eles.

Esse modelo é útil duas vezes. Primeiro, a sua predição pode entrar no lugar da realimentação atrasada, e então a condição de estabilidade~\eqref{eq:delaystable} deixa de amarrar as mãos: pode-se corrigir depressa, pelo modelo. Segundo, a diferença entre o previsto e o recebido é um sinal do \emph{inesperado}. Tudo o que a máquina fez por si mesma o modelo previu e subtraiu; resta o que o mundo fez. Por isso, segundo uma das hipóteses, não conseguimos fazer cócegas em nós mesmos~\cite{Blakemore1998}.

\section{Quando o mundo muda: o rápido e o lento}

Vamos testar o circuito com um experimento fácil de montar. Uma pessoa estende a mão para um alvo, e o mundo muda de repente: por exemplo, uma força lateral age sobre a mão. Os primeiros movimentos erram, depois o erro diminui. Se a força for retirada, a mão primeiro erra para o \emph{outro} lado: o modelo aprendeu a força e continua a compensá-la. É uma prova direta de que foi aprendido um modelo, e não simplesmente ``passou a dar certo''.

Os experimentos mostram também algo mais sutil: aprendem ao mesmo tempo dois processos, um rápido, que aprende rápido e esquece rápido, e um lento, que aprende devagar e lembra por muito tempo~\cite{Smith2006}. As correções são atualizadas depois de cada movimento, por evento:
\begin{empheq}[box=\keybox]{equation}
e = p - (x_f + x_s), \qquad x_f \leftarrow A_f\,x_f + B_f\,e, \qquad x_s \leftarrow A_s\,x_s + B_s\,e ,
\label{eq:tworate}
\end{empheq}
em que $p$ é a perturbação, $x_f$ e $x_s$ são as correções aprendidas pelos dois processos, $A$ indica quanto a correção se conserva até o próximo movimento, e $B$, quanto ela aprende com o erro. No processo rápido, $B$ é grande e $A$ é bem menor que um; no lento, o contrário.

\begin{figure}[t]
\centering
\begin{tikzpicture}[>=Stealth,x=0.034cm,y=1.9cm]
\draw[->,gray!70!black] (0,0) -- (355,0) node[right,font=\small,black] {movimento};
\draw[->,gray!70!black] (0,-1.15) -- (0,1.25);
\foreach \v in {-1,1}{\draw[gray!60!black] (-3,\v) -- (3,\v); \node[left,font=\scriptsize] at (0,\v) {$\v$};}
\foreach \n in {100,200,300}{\draw[gray!60!black] (\n,-0.04) -- (\n,0.04); \node[below,font=\scriptsize] at (\n,0) {\n};}
\draw[thin,gray!70!black] plot file {figures/tworate_p.dat};
\draw[thick,orange!85!black] plot file {figures/tworate_fast.dat};
\draw[thick,blue!60!black] plot file {figures/tworate_slow.dat};
\draw[very thick,black] plot file {figures/tworate_net.dat};
\node[font=\scriptsize,gray!50!black] at (120,1.12) {perturbação $p$};
\node[font=\scriptsize,black,anchor=west] at (238,0.52) {soma: retorno};
\node[font=\scriptsize,blue!60!black,anchor=west] at (150,0.39) {lento $x_s$};
\node[font=\scriptsize,orange!85!black,anchor=west] at (60,0.08) {rápido $x_f$};
\draw[densely dashed,gray!60!black] (226,-1.15) -- (226,1.25);
\node[font=\scriptsize,align=center,anchor=south west] at (228,-1.1) {sem mais\\erros};
\end{tikzpicture}
\caption{Dois processos de aprendizado segundo o modelo~\eqref{eq:tworate}, $A_f=0.92$, $B_f=0.1$, $A_s=0.996$, $B_s=0.02$. Duzentos movimentos com perturbação $p=1$, depois seis com a perturbação inversa, $p=-1$, até a soma voltar a zero. Depois da linha tracejada os erros deixam de chegar, e a soma volta sozinha à correção antiga: o processo rápido esqueceu a perturbação inversa, o lento lembra a direta.}
\label{fig:tworate}
\end{figure}

O modelo~\eqref{eq:tworate} explica um experimento difícil de entender sem ele (fig.~\ref{fig:tworate}). No nosso cálculo, em duzentos movimentos com perturbação, a soma das correções chega a $0.84$, e a maior parte é mantida pelo processo lento, $0.63$. Depois, seis movimentos com a perturbação inversa devolvem a soma a zero: o processo rápido foi para o negativo, $-0.53$, o lento ainda está no positivo, $0.46$. Se agora os erros deixarem de ser informados, o processo rápido esquece a sua correção em dezenas de movimentos, o lento, muito mais devagar, e a soma sobe \emph{sozinha} até $0.37$ em quarenta movimentos: a habilidade antiga volta sem treino algum. Essa recuperação espontânea foi medida em humanos; um modelo com um só processo não consegue produzi-la: sem erros, ele simplesmente decai a zero.

Por que dois processos? Uma resposta de engenharia plausível vem do filtro ótimo do capítulo~\ref{sec:acet}. O corpo muda por razões diferentes, em velocidades diferentes: os músculos se cansam em minutos, crescem e enfraquecem em meses. Se as perturbações são rápidas e lentas, a melhor estimativa consiste em dois filtros com constantes de tempo diferentes, cada um pegando a sua~\cite{Kording2007}. O processo rápido é uma correção temporária, ``a mão está cansada''; o lento, ``a mão ficou diferente''. Por enquanto isso é hipótese, mas explica por que uma habilidade aprendida num dia se perde em parte até o dia seguinte e por que, na segunda vez, ela é aprendida mais depressa.

\section{Resumo}

\begin{table}[t]
\centering
\footnotesize
\setlength{\tabcolsep}{4pt}
\renewcommand{\arraystretch}{1.15}
\begin{tabular}{>{\raggedright}p{3.0cm}>{\raggedright}p{4.6cm}>{\raggedright\arraybackslash}p{5.2cm}}
\toprule
O que o circuito faz & Como & O que se sabe \\
\midrule
aprende com professor & fio de erro para cada saída, regra~\eqref{eq:lms} & construção do circuito e enfraquecimento por coincidência medidos; aprendizado supervisionado é a hipótese principal \\
expande a entrada & setenta bilhões de neurônios \nt{GLUT} & número de neurônios medido; papel da expansão é teoria \\
leva em conta o atraso do erro & traço ajustado ao atraso & medido \\
constrói o modelo interno & filtro adaptativo~\eqref{eq:adaptivefilter} & hipótese bem fundamentada \\
aprende em dois ritmos & processos rápido e lento~\eqref{eq:tworate} & pós-efeito e retorno espontâneo medidos; dois processos são modelo \\
\bottomrule
\end{tabular}
\caption{Como a máquina aprende a se mover.}
\label{tab:motorlearning}
\end{table}

Agora a máquina tem três professores (tabela~\ref{tab:motorlearning}). Hebb comprime o que é frequente; \nt{DOPA}, com um só número, ensina a escolher o vantajoso; o fio de erro ensina a precisão, e ensina depressa, porque cada saída recebe o seu próprio erro. O cérebro, ao que parece, usa os três, cada um onde ele é mais barato: o sinal de difusão para poucas decisões, fios separados para o ajuste fino do corpo, em que a resposta certa é dada pelos próprios sensores.

\section{Exercícios}

\begin{exercise}[\lvlA]
\label{ex:15:1}
\emph{A capacidade do circuito com fio de erro.} São $3\cdot10^7$ neurônios de saída com $10^5$ entradas, cada um com seu fio de erro ($1$~disparo/s). Um neurônio com $n$ entradas aprende toda rotulação de $\approx2n$ vetores~\cite{Cover1965}. (a)~Quantas associações o circuito guarda? (b)~Estime por~\eqref{eq:spikeentropy}, com $\Delta t=10\ms$, o menor tempo para encher a capacidade do neurônio.
\end{exercise}

\begin{exercise}[\lvlB]
\label{ex:15:2}
\emph{A velocidade da regra dos mínimos quadrados.} Os modos da regra~\eqref{eq:lms} decaem com taxas $\eta\lambda_k(C)$. Para cinco filtros~\eqref{eq:adaptivefilter} com ruído branco, $C_{jk}=2\sqrt{\tau_j\tau_k}/(\tau_j+\tau_k)$. Ache a razão entre a taxa mais rápida e a mais lenta se os $\tau_j$ crescem $2$ vezes ($10$--$160\ms$) e $4$ vezes.
\end{exercise}

\begin{exercise}[\lvlB]
\label{ex:15:3}
\emph{Análise dos dois processos.} Escreva o modelo~\eqref{eq:tworate} com os parâmetros da fig.~\ref{fig:tworate} como $\mathbf x_{n+1}=M\mathbf x_n+\mathbf b\,p$. (a)~Ache, pelos autovalores de $M$, as constantes de tempo em movimentos. (b)~Ache os valores estacionários de $x_f$, $x_s$ e a sua soma com $p=1$ constante.
\end{exercise}

\begin{exercise}[\lvlB]
\label{ex:15:4}
\emph{Simulação: a segunda vez é mais rápida.} Pelo modelo~\eqref{eq:tworate} com os parâmetros de \texttt{models/\allowbreak motor\_learning.py}, ache o número de movimentos até $x_f+x_s\ge0.8$ com $p=1$: (a)~numa máquina nova; (b)~depois de $200$ movimentos com $p=1$ e da perturbação inversa até zerar a soma, como na fig.~\ref{fig:tworate}. Um modelo de um só processo acelera assim?
\end{exercise}

\begin{exercise}[\lvlB]
\label{ex:15:5}
\emph{Regulador com modelo interno.} O objeto $\dd x/\dd t=ku$ é visto com atraso $d$; o regulador $u=-G\hat x$ prevê $\hat x(t)=x(t-d)+\hat k\int_{t-d}^{t}u\,\dd s$. (a)~Com $\hat k=k=1$, $d=150\ms$, ache $G$ para uma constante de tempo de $20\ms$ e compare com o limite~\eqref{eq:delaystable}. (b)~O esquema é estável com $\hat k=0.4k$?
\end{exercise}

\begin{exercise}[\lvlB]
\label{ex:15:6}
\emph{O filtro adaptativo ensina o músculo.} O objeto é o abalo~\eqref{eq:twitch} de área unitária, $T_c=50\ms$; os filtros~\eqref{eq:adaptivefilter} são circuitos RC com $\tau_j=12.5$, $25$, $50$, $100$, $200\ms$; a entrada é um novo número normal a cada $10\ms$. Em quantos segundos a regra~\eqref{eq:lms} com $\eta=50$ e $200$ reduz o erro a $0.5\%$? Por que os pesos, mesmo após $300$~s, estão longe dos melhores?
\end{exercise}

\begin{exercise}[\lvlC]
\label{ex:15:7}
\emph{Dois processos como filtro ótimo.} A perturbação é a soma dos processos $p^{f,s}_{n+1}=A_{f,s}\,p^{f,s}_n+w^{f,s}_n$ com variâncias de ruído $q_f$, $q_s$, e o erro é observado com ruído de variância $R$; o filtro de Kalman dá~\eqref{eq:tworate}. Com que $q_f/R$, $q_s/R$, a partir de $A_f=0.92$, $A_s=0.996$, resulta $B_f=0.1$, $B_s=0.02$? Como mudam $B_f$, $B_s$ se $q_s$ for quadruplicado?
\end{exercise}

\chapter{A segunda saída: como o cérebro controla a química do corpo}
\chaptermark{A saída química}
\label{sec:chemoutput}
\begin{chaptergoals}
\item como o acelerador \nt{NORA} e o freio \nt{ACET} ajustam um marca-passo, e por que o freio é mais rápido;
\item como o sangue funciona como barramento de difusão, e o \nt{ADRE} põe o corpo em modo de corrida;
\item por que uma cascata hormonal com realimentação não pode ser precisa e estável ao mesmo tempo: $K<8$;
\item como um receptor que se cansa lê a frequência de pulsos, e a luz trava a fase do ritmo diário.
\end{chaptergoals}

\section{A saída rápida e a lenta}

No capítulo~\ref{sec:muscles} vimos a saída rápida da máquina: os músculos. Mas ela tem uma segunda saída, lenta. O cérebro mantém dentro dos limites certos a temperatura do corpo, a frequência cardíaca, a quantidade de água e de açúcar, e prepara o corpo para a corrida ou para o sono. Para isso ele não move articulações: libera sinais químicos. Há dois caminhos. O primeiro são os nervos que vão direto aos órgãos internos: ao coração, aos vasos, ao intestino, às glândulas. O segundo é o sangue: alguns neurônios e glândulas não liberam a substância na fenda ao lado do vizinho, mas direto na corrente sanguínea.

O sangue é o barramento de difusão mais amplo de todos os que apareceram no livro. Os canais de difusão dos capítulos~\ref{sec:serocomputer}--\ref{sec:acet} enviavam o sinal a grandes regiões do cérebro; o sangue o envia a cada célula do corpo. O sangue dá a volta no corpo em cerca de um minuto, de modo que a mensagem chega a todos em dezenas de segundos e se mantém por minutos e horas, até a substância se degradar. Esse envio não tem endereço nenhum. Como na proposição~\ref{prop:receptor}, quem escolhe o sentido da mensagem é o destinatário: só respondem as células que têm receptor para essa substância.

\section{Acelerador e freio: dois fios para os órgãos}

O melhor exemplo é o coração. O coração tem seu próprio marca-passo, como o neurônio \nt{SERO} do capítulo~\ref{sec:serocomputer}: ele mesmo define o ritmo, e, se todos os nervos forem desligados, o coração bate cerca de cem vezes por minuto. O cérebro não gera o ritmo, só o ajusta, e por dois fios de sinais opostos (fig.~\ref{fig:gasbrake}). Por um vem o \nt{NORA}: é o \emph{acelerador}, que acelera o ritmo. Pelo outro, o \nt{ACET}: é o \emph{freio}, que o desacelera. Grosso modo,
\begin{empheq}[box=\keybox]{equation}
f \;\approx\; f_0 \;+\; a_S\,S \;-\; a_P\,P ,
\label{eq:gasbrake}
\end{empheq}
onde $f_0\approx100$ batimentos por minuto é o ritmo próprio do marca-passo, e $S$ e $P$ são a atividade do acelerador e do freio. Em repouso o freio está pisado, e o coração costuma bater de sessenta a setenta vezes por minuto; no esforço, o freio é solto e o acelerador é pisado.

\begin{figure}[t]
\centering
\begin{tikzpicture}[>=Stealth,
  blk/.style={draw,very thick,rounded corners=2pt,align=center,font=\small},
  pace/.style={circle,draw,very thick,double,double distance=1.2pt,minimum size=1.8cm,inner sep=0pt,font=\scriptsize,fill=red!8,align=center},
  inh/.style={-{Bar[width=7pt]},thick,red!70!black}]
\node[blk,fill=blue!5,text width=1.6cm] (B) at (0,0) {comando\\do cérebro};
\node[pace] (H) at (6.2,0) {marca-\\passo};
\draw[->,very thick,orange!85!black] (B.north east) to[bend left=20] node[above,font=\scriptsize,black,align=center] {\nt{NORA}: acelerador, segundos} (H.north west);
\draw[inh,very thick,blue!60!black] (B.south east) to[bend right=20] node[below,font=\scriptsize,black,align=center] {\nt{ACET}: freio, frações de segundo} (H.south west);
\draw[->,very thick] (H) -- (8.4,0) node[right,font=\small] {frequência $f$};
\node[blk,fill=orange!10,text width=1.6cm] (A) at (0,-2.6) {\nt{ADRE}\\no sangue};
\draw[->,thick,orange!85!black] (B) -- (A);
\foreach \y/\lab in {-2.0/{coração},-2.6/{vasos},-3.2/{fígado, músculos}}{
  \draw[->,thick,densely dashed,orange!85!black] (A.east) -- (4.3,\y) node[right,font=\scriptsize,black] {\lab};}
\end{tikzpicture}
\caption{Dois fios de sinais opostos para o marca-passo do coração: o acelerador \nt{NORA} é lento, o freio \nt{ACET} é rápido. Embaixo, o mesmo acelerador pelo barramento de difusão: as células \nt{ADRE} liberam adrenalina no sangue, e o sinal chega a todos os órgãos com o receptor adequado (setas tracejadas).}
\label{fig:gasbrake}
\end{figure}

É o código de dois fios do capítulo~\ref{sec:glutcomputer} e o par de músculos do capítulo~\ref{sec:muscles}, só que agora os fios não controlam uma articulação, e sim o andamento do marca-passo. E os fios têm velocidades diferentes. Os dois funcionam como filtros passa-baixas, mas o freio deixa passar as mudanças várias vezes mais depressa que o acelerador~\cite{Berger1989}: no \nt{ACET} o caminho é curto, pois a proteína ligada ao receptor abre ela mesma um canal de potássio, sem segundo mensageiro dentro da célula~\cite{Logothetis1987gbg}, enquanto o \nt{NORA} age por uma cadeia de reações dentro da célula. O engenheiro reconhece aqui a técnica dos dois atuadores de velocidades diferentes: o rápido faz o ajuste fino instantâneo, o lento faz as mudanças grandes e demoradas.

Aqui também encontra seu papel o \nt{ADRE}, sobre o qual a tabela~\ref{tab:types} dizia que seu papel no cérebro não é claro. Fora do cérebro ele é claro. Parte das células da via do acelerador não manda um fio ao órgão, mas libera adrenalina direto no sangue. É o mesmo acelerador, mas pelo barramento de difusão: em dezenas de segundos ele chega ao coração, aos vasos, ao fígado e aos músculos ao mesmo tempo e se mantém por minutos. O acelerador endereçado \nt{NORA} ajusta um órgão; o acelerador difuso \nt{ADRE} põe o corpo inteiro em modo de corrida.

\section{Cascata com realimentação}

Muitos hormônios não são liberados por uma só glândula, mas por uma cascata. Um pequeno grupo de neurônios na base do cérebro lança no sangue o primeiro sinal; ele faz uma glândula intermediária liberar o segundo; este faz a glândula final liberar o hormônio que de fato age no corpo. E o hormônio final volta aos primeiros estágios e os inibe. O resultado é um amplificador de três estágios envolvido por realimentação negativa: um regulador de nível, como o do sistema \nt{SERO} na fórmula~\eqref{eq:feedback}.

Descrevemos cada estágio por um circuito RC com constante de tempo $\tau$ e ganho $g$, e a realimentação por um coeficiente $k$:
\begin{equation}
\tau\,\frac{\dd x_1}{\dd t} = -x_1 + g\bigl[u - k\,x_3\bigr]_+ ,\qquad
\tau\,\frac{\dd x_2}{\dd t} = -x_2 + g\,x_1 ,\qquad
\tau\,\frac{\dd x_3}{\dd t} = -x_3 + g\,x_2 ,
\label{eq:cascade}
\end{equation}
onde $u$ é o comando do cérebro e $x_3$, o nível do hormônio final. O ganho do laço é $K=g^3k$. No equilíbrio, $x_3=g^3u/(1+K)$ e, como em qualquer regulador com realimentação, um laço grande torna o nível insensível à dispersão dos ganhos dos estágios. Mas cada estágio desloca a fase: na frequência $\omega=\sqrt3/\tau$, os três estágios iguais juntos dão meia volta de defasagem e uma atenuação de oito vezes. Daí o resultado clássico da teoria de controle:
\begin{empheq}[box=\keybox]{equation}
K \;<\; 8 ,
\label{eq:cascadestable}
\end{empheq}
senão o regulador começa a oscilar, com período de cerca de $2\pi\tau/\sqrt3\approx3.6\,\tau$.

\begin{keyblock}
\begin{proposition}[Precisão contra estabilidade]
\label{prop:cascade}
Numa cascata de três estágios com realimentação proporcional, o erro de nível é $1/(1+K)$, e ela só é estável para $K<8$. Por isso esse regulador não mantém o nível com precisão melhor que uns dez por cento; a tentativa de aumentar o ganho o transforma num oscilador.
\end{proposition}
\end{keyblock}

Verificamos isso no modelo~\eqref{eq:cascade} (fig.~\ref{fig:cascade}). Com $K=4$, o nível oscila um pouco depois de ligado e se estabiliza. Com $K=7$ também se estabiliza, mas devagar. Com $K=12$, as oscilações não se amortecem nem crescem sem limite: um estágio não consegue emitir sinal negativo, e a cascata entra em pulsações regulares entre $0.83$ e $1.24$ do nível médio, com período de $3.7$--$3.9\,\tau$.

\begin{figure}[t]
\centering
\begin{tikzpicture}[>=Stealth,x=0.3cm,y=1.2cm]
\draw[->,gray!70!black] (0,0) -- (41.5,0) node[right,font=\small,black] {$t/\tau$};
\draw[->,gray!70!black] (0,0) -- (0,2.8) node[left,font=\small,black] {$x_3/x_3^*$};
\foreach \t in {0,10,20,30,40}{\draw[gray!60!black] (\t,-0.05) -- (\t,0.05); \node[below,font=\scriptsize] at (\t,0) {\t};}
\foreach \f in {1,2}{\draw[gray!60!black] (-0.3,\f) -- (0.3,\f); \node[left,font=\scriptsize] at (0,\f) {\f};}
\draw[densely dashed,gray!60!black] (0,1) -- (40,1);
\draw[thick,blue!60!black] plot file {figures/cascade_K4.dat};
\draw[thick,red!70!black] plot file {figures/cascade_K12.dat};
\node[font=\scriptsize,blue!60!black,anchor=west] at (16,0.55) {$K=4$: estabiliza-se};
\node[font=\scriptsize,red!70!black,anchor=west] at (16,1.75) {$K=12$: pulsa};
\end{tikzpicture}
\caption{Cascata de três estágios com realimentação segundo~\eqref{eq:cascade}: nível do hormônio final em frações do valor de equilíbrio $x_3^*$. Com ganho do laço abaixo de oito, o nível se estabiliza; acima, a própria cascata gera pulsações regulares.}
\label{fig:cascade}
\end{figure}

Os hormônios reais de cascata de fato saem em pulsos: o nível do hormônio do estresse, por exemplo, oscila com período de cerca de uma hora. Segundo uma das hipóteses, essas pulsações nascem da própria realimentação atrasada entre os estágios, e não é preciso um gerador de ritmo separado para elas~\cite{Walker2010}. É a mesma causa de oscilação do laço \nt{GLUT}--\nt{GABA} na proposição~\ref{prop:ei} e do laço de estiramento do músculo na condição~\eqref{eq:delaystable}: uma realimentação negativa que chega atrasada.

\section{O código de pulsos dos hormônios}

Os pulsos não são só um efeito colateral: também podem ser um código. Os neurônios que controlam a reprodução liberam seu sinal no sangue em porções curtas, cerca de uma vez por hora. Se a glândula receptora recebe o mesmo sinal não em porções, mas continuamente, ela não trabalha com força total, como com os pulsos: ela se cala; se as porções voltam a vir uma vez por hora, o trabalho se restabelece~\cite{Belchetz1978}. Portanto, o receptor não lê o nível, mas a \emph{frequência} das porções.

Para o engenheiro não há mistério aqui. Um receptor que se cansa e deixa de responder sob ação prolongada da substância é um filtro passa-altas, como a sinapse que se esgota~\eqref{eq:depression} do capítulo~\ref{sec:code}. Ele suprime o sinal contínuo e deixa passar as mudanças. A esse destinatário só se pode comunicar algo por pulsos, e a informação passa do nível para a frequência e a forma das porções. Além disso, frequências diferentes de porções agem de modo diferente sobre células diferentes da mesma glândula, de modo que por uma única linha química se pode transmitir mais de uma grandeza, assim como pelo único fio \nt{DOPA} do capítulo~\ref{sec:dofaalt} passavam uma componente rápida e uma lenta.

\section{O ritmo diário: um gerador lento de modos}

O ritmo mais lento do corpo é o diário. Ele nasce de uma realimentação negativa atrasada dentro das células: genes produzem proteínas que, com um atraso de algumas horas, suprimem sua própria produção~\cite{TakahashiJS2017}. De novo uma realimentação negativa atrasada, pela quarta vez neste livro, e de novo um ritmo, desta vez com período de cerca de um dia. O gerador principal é um pequeno grupo de dezenas de milhares de neurônios na base do cérebro, cujo ritmo sincroniza as demais células do corpo.

O período próprio desse gerador no ser humano não é exatamente um dia, mas em média $24.18$ horas~\cite{Czeisler1999}. Todo dia ele é acertado pela luz que chega dos sensores do olho (capítulo~\ref{sec:sensors}). Para o eletrônico, isso é uma \emph{malha de captura de fase} (PLL): um oscilador com frequência própria $\omega_0$ se ajusta a um sinal externo de frequência $\omega$, e a diferença de fase $\varphi$ obedece à equação
\begin{empheq}[box=\keybox]{equation}
\frac{\dd\varphi}{\dd t} \;=\; (\omega-\omega_0) \;-\; K_\varphi\,\sin\varphi .
\label{eq:pll}
\end{empheq}
A sincronização se mantém se $|\omega-\omega_0|<K_\varphi$. Um gerador com período de $24.18$ horas precisa se deslocar cerca de onze minutos por dia; a luz da manhã costuma bastar para isso. Na noite polar ou numa caverna sem luz não há ajuste, e o ritmo segue seu período próprio.

O gerador diário não é o gerador de clock de um computador. Ele não conta passos de computação nem sincroniza os disparos dos neurônios. Ele troca \emph{modos}: sono e vigília, níveis hormonais, temperatura do corpo, disposição para comer. É um registrador de difusão lento do mesmo tipo que os registradores dos capítulos~\ref{sec:serocomputer}--\ref{sec:acet}, só que seu valor muda segundo um horário acertado pelo sol.

\begin{keyblock}
\begin{proposition}[A saída química]
\label{prop:chemoutput}
A saída lenta da máquina é construída com as mesmas peças que a rápida. Os órgãos recebem dois fios de sinais opostos e velocidades diferentes, o acelerador \nt{NORA} e o freio \nt{ACET}, e o organismo inteiro recebe de uma vez sinais de difusão pelo sangue, entre eles a adrenalina \nt{ADRE}. Os níveis hormonais são mantidos por cascatas com realimentação negativa, cuja precisão é limitada pela estabilidade; parte dos hormônios é codificada pela frequência das porções; o ritmo diário é um gerador lento com captura de fase pela luz. A organização das vias e os experimentos com porções e com o período foram medidos; a explicação das pulsações da cascata pela própria realimentação é uma hipótese.
\end{proposition}
\end{keyblock}

\section{Resumo}

\begin{table}[t]
\centering
\footnotesize
\setlength{\tabcolsep}{4pt}
\renewcommand{\arraystretch}{1.15}
\begin{tabular}{>{\raggedright}p{3.1cm}>{\raggedright}p{4.8cm}>{\raggedright\arraybackslash}p{4.9cm}}
\toprule
O que a saída química faz & Como & O que se sabe \\
\midrule
ajusta os órgãos & acelerador \nt{NORA} e freio \nt{ACET}~\eqref{eq:gasbrake} & medido; o freio é mais rápido que o acelerador \\
põe o corpo inteiro em modo de corrida & adrenalina \nt{ADRE} no sangue & medido \\
mantém os níveis hormonais & cascata com realimentação, $K<8$~\eqref{eq:cascadestable} & cascatas e realimentação medidas; pulsações geradas pelo próprio laço: hipótese \\
transmite pela frequência das porções & receptor cansado como filtro passa-altas & dependência das porções medida \\
troca de modo ao longo do dia & gerador com captura de fase pela luz~\eqref{eq:pll} & período e ajuste pela luz medidos \\
\bottomrule
\end{tabular}
\caption{Como a máquina controla a química do corpo.}
\label{tab:chemoutput}
\end{table}

A máquina tem duas saídas (tabela~\ref{tab:chemoutput}). A rápida são os músculos: milissegundos, endereçada, a cada articulação. A lenta é a química do corpo: segundos, minutos e dias, por dois fios para os órgãos e pelo sangue para todos de uma vez. As peças das duas saídas são as mesmas (dois fios de sinais opostos, marca-passos, a realimentação e seu atraso), e são as mesmas com que a máquina é montada por dentro.

\section{Exercícios}

\begin{exercise}[\lvlA]
\label{ex:16:1}
\emph{O sangue como filtro.} Um hormônio é eliminado do sangue com meia-vida $t_{1/2}=10$~min e liberado em porções curtas a cada $T=60$~min. Quantas vezes o máximo da concentração é maior que o mínimo, e quanto é atenuado o harmônico fundamental das porções? Com que $t_{1/2}$ o máximo supera o mínimo só duas vezes?
\end{exercise}

\begin{exercise}[\lvlA]
\label{ex:16:2}
\emph{Sincronização diária.} O período próprio do gerador~\eqref{eq:pll} é $24.18$~h, e a luz desloca a fase no máximo uma hora por dia: $K_\varphi=2\pi/24$~rad/dia. Ache a diferença de fase estacionária $\varphi^*$ em minutos e o tempo de relaxação. Em quantos dias o desvio após um salto de $\pm6$~h no sinal externo fica menor que uma hora?
\end{exercise}

\begin{exercise}[\lvlB]
\label{ex:16:3}
\emph{Precisão contra estabilidade.} A cascata linearizada~\eqref{eq:cascade} foi estendida para $n$ estágios iguais. Ache o ganho crítico $K_c(n)$ e o menor erro atingível $1/(1+K_c)$ para $n=3$ e $n=6$. Para onde tende esse erro quando $n\to\infty$?
\end{exercise}

\begin{exercise}[\lvlB]
\label{ex:16:4}
\emph{Acelerador e freio.} Em~\eqref{eq:gasbrake}, as contribuições do acelerador e do freio são saídas de circuitos RC com $\tau_S=2$~s e $\tau_P=0.4$~s, com comandos de no máximo $80$ e $60$ batimentos por minuto. O ritmo é $f_0=100$; em repouso, freio $35$, acelerador $0$, $f=65$. Em quanto tempo se chega a $f=120$? E se o freio não for solto e só o acelerador trabalhar?
\end{exercise}

\begin{exercise}[\lvlB]
\label{ex:16:5}
\emph{Simulação: cascata com atraso.} Na realimentação da função \texttt{cascade} de \texttt{models/\allowbreak chem\_models.py} (copie-a, não execute o arquivo), acrescente um atraso $d$: o primeiro estágio vê $x_3(t-d)$. Ache $K_c$ e o período no limite para $d/\tau=0$, $0.5$, $1$, $2$ e verifique por simulação com $0.9K_c$ e $1.1K_c$.
\end{exercise}

\begin{exercise}[\lvlB]
\label{ex:16:6}
\emph{Um destinatário que lê porções.} O receptor se cansa: $y=[c-a]_+$, $\tau_a\,\dd a/\dd t=c-a$, $\tau_a=30$~min. O hormônio chega em porções de $6$~min com período $T$ e a mesma dose média $\bar c$. Ache a resposta média para $T=15$, $60$, $120$~min e $T\to\infty$. Qual é ela com concentração constante $\bar c$?
\end{exercise}

\begin{exercise}[\lvlC]
\label{ex:16:7}
\emph{Gerador ou realimentação?} Proponha um experimento que distinga as pulsações geradas pela realimentação atrasada da cascata~\eqref{eq:cascade} das de um marca-passo separado. É possível distinguir as hipóteses pelo deslocamento do período se $K$ só pode ser alterado em uma vez e meia, e o período é medido com precisão de $5\%$?
\end{exercise}

\chapter{Depurando a máquina}
\chaptermark{Depurando a máquina}
\label{sec:debugging}

\section{Como se depura uma máquina sem esquema}

Um engenheiro que recebe uma placa funcionando sem o esquema a depura com quatro técnicas. Ele põe uma \emph{sonda} e olha o que passa pelo fio. Ele \emph{aplica um sinal de teste} e olha o que mudou. Ele \emph{desliga} um pedaço do circuito e olha o que parou de funcionar. E ele \emph{lê os registradores de controle} para saber em que modo a placa está operando. Este livro inteiro é uma tentativa de ler o esquema do cérebro, e neste capítulo veremos quais das quatro técnicas os engenheiros já sabem usar e o que os capítulos anteriores dizem sobre o que esperar.

O exemplo mais visível hoje dessa depuração são as interfaces cérebro-computador: dispositivos que leem os disparos dos neurônios e os transformam em comandos para máquinas externas ou, ao contrário, enviam ao cérebro sinais de fora. A maior parte do capítulo é dedicada a elas e, antes de tudo, ao que faz a empresa Neuralink.

\section{A sonda: o que o eletrodo ouve}

A sonda do cérebro é um eletrodo fino inserido no tecido. Ele ouve os disparos dos neurônios que estão num raio da ordem de cem micrômetros: em geral, de um a alguns. No eletrodo, um disparo é um pico de dezenas a centenas de microvolts com duração de cerca de um milissegundo, e pela forma dele é possível distinguir neurônios vizinhos uns dos outros. Os mais bem ouvidos são os grandes neurônios \nt{GLUT}: eles são a maioria no córtex (tabela~\ref{tab:types}), e suas correntes são mais fortes.

A dificuldade principal aparece logo. Uma boa sonda hoje tem centenas ou milhares de eletrodos, e o cérebro tem $8.6\cdot10^{10}$ neurônios. Escutamos cerca de um centésimo de milionésimo da máquina. Por que, então, dá para entender alguma coisa a partir de uma amostra tão ínfima? A resposta veio no capítulo~\ref{sec:code}. Neurônios vizinhos têm ruído correlacionado, e a média sobre um grupo esbarra num limite (proposição~\ref{prop:pooling}): cem neurônios carregam quase tanto quanto mil. Além disso, a tarefa que o córtex motor resolve tem poucas dimensões: o braço tem poucas articulações (capítulo~\ref{sec:muscles}), e a atividade das centenas de neurônios que o controlam cabe num espaço de apenas uma ou duas dezenas de variáveis comuns~\cite{Gallego2017}. Uma sonda de cem neurônios ouve quase todas essas variáveis. Se o cérebro guardasse uma mensagem independente em cada neurônio, essa escuta seria inútil.

\section{O fluxo de dados: compressão na própria sonda}

A sonda produz um fluxo enorme. Para ver a forma do disparo, cada eletrodo precisa ser digitalizado cerca de $20$ mil vezes por segundo, com precisão de uns dez bits. Mil eletrodos dão
\begin{empheq}[box=\keybox]{equation}
\begin{aligned}
R_{\text{bruto}} \;&=\; N\,f_s\,b \;\approx\; 10^3\cdot2\cdot10^4\cdot10 \;\approx\; 2\cdot10^8\ \text{bit/s},\\
R_{\text{disp}} \;&\approx\; N\,r\,\log_2\frac{e}{r\,\Delta t} \;\approx\; 8\cdot10^4\ \text{bit/s}.
\end{aligned}
\label{eq:probedata}
\end{empheq}
A segunda estimativa é a capacidade do fluxo de disparos~\eqref{eq:spikeentropy} com $r=10$ disparos por segundo e precisão $\Delta t=1\ms$: tudo o que de fato há nele ocupa duas mil e quinhentas vezes menos. Para um dispositivo implantado, essa diferença é decisiva: o fluxo bruto não pode ser transmitido por rádio através da pele com a potência que se pode gastar dentro da cabeça sem aquecer o tecido. Por isso um dispositivo implantado precisa detectar os disparos no próprio chip, junto dos eletrodos, e enviar para fora só os eventos: o número do canal e o instante do disparo. Na primeira descrição do sistema da Neuralink isso ainda não tinha sido alcançado: os chips amplificam e digitalizam os sinais, todo o fluxo bruto sai por cabo, e os disparos são detectados por um programa do lado de fora; a compressão no chip e a comunicação sem fio aparecem ali como plano~\cite{Musk2019}.

É uma técnica conhecida. O olho comprime o sinal cem vezes antes de mandá-lo pelo cabo longo (capítulo~\ref{sec:sensors}); uma câmera de eventos não transmite quadros, mas mudanças. Para caber no fio, a sonda é obrigada a fazer o mesmo que o próprio cérebro faz: falar por disparos e só sobre novidades.

\section{O que a Neuralink faz}

O dispositivo da Neuralink é uma sonda construída sob essas restrições~\cite{Musk2019}. Em vez de uma matriz rígida de agulhas, há muitos fios flexíveis, mais finos que um cabelo, com eletrodos ao longo de cada um: na primeira descrição do sistema, até $3072$ eletrodos em $96$ fios, $32$ por fio; no dispositivo implantado em pessoas, segundo a empresa, cerca de mil canais em $64$ fios. Fios flexíveis danificam menos o tecido e toleram melhor o fato de o cérebro se deslocar ligeiramente o tempo todo dentro do crânio. Os fios são inseridos por um robô, que desvia dos vasos. Na primeira descrição, como dito acima, o fluxo bruto~\eqref{eq:probedata} saía por cabo. O dispositivo para pessoas, segundo a empresa, é diferente: a carcaça fica totalmente coberta pela pele, os disparos são detectados dentro dela, os dados saem por rádio e a energia chega sem fio.

A tarefa do primeiro dispositivo é a leitura. Os fios são postos na parte do córtex que planeja os movimentos do braço, e um decodificador transforma os disparos no movimento de um cursor na tela. Uma pessoa com os braços paralisados controla o computador imaginando o movimento. A ideia em si não é nova: interfaces com matrizes rígidas de uma centena de eletrodos já permitiam há tempos controlar um cursor~\cite{Hochberg2006}, escrever texto imaginando os movimentos da mão ao escrever~\cite{Willett2021} e até converter tentativas de falar em texto na tela a cerca de sessenta palavras por minuto~\cite{Willett2023}. A novidade da Neuralink é a engenharia: o número de canais, a ausência de fios, a carcaça fechada e a inserção robotizada, isto é, a tentativa de fazer uma sonda que possa ser usada por anos fora do laboratório. Pelo que sabemos, os resultados dos testes em pessoas foram publicados sobretudo pela própria empresa, e não em artigos revisados por pares; a empresa relatou, em particular, que parte dos fios se deslocou depois da inserção e o número de canais funcionando diminuiu. Para a depuração, é um contratempo comum: uma sonda que se move em relação à placa já ouve outros fios.

A segunda linha da empresa é a escrita: um dispositivo de visão que envia sinais ao córtex visual de uma pessoa que perdeu a visão. Falamos dele abaixo, na seção sobre escrita.

\section{Leitura: o decodificador como destinatário}

O decodificador da interface é um destinatário no sentido da proposição~\ref{prop:reader}: ele mesmo decide que parte da mensagem ler. Praticamente todas as interfaces leem \emph{taxas de grupo}: contam os disparos de cada canal em janelas de algumas dezenas de milissegundos e os somam com pesos ajustados para produzir o movimento pretendido. É o código de taxa de grupo da seção~\ref{sec:rate}, e a escolha não é por acaso: com poucos disparos por janela, a taxa de um só neurônio não está definida, mas a soma sobre uma centena já funciona.

Quanta informação se consegue extrair? A escrita com a ``mão mental'' chegou a cerca de noventa caracteres por minuto~\cite{Willett2021}, ou seja, um caractere e meio por segundo: mesmo com cinco bits por caractere, são menos de oito bits por segundo. O controle do cursor, segundo a Neuralink, dá da ordem de dez bits por segundo. Já o limite superior~\eqref{eq:spikeentropy} para \emph{um} neurônio é de dezenas de bits por segundo, e para cem, de milhares. A interface extrai dos neurônios escutados uma pequena fração do que eles poderiam carregar. O gargalo não é o fio, e sim quantas variáveis independentes o grupo carrega (proposição~\ref{prop:pooling}) e quão bem o decodificador entende o código, que, como vimos no capítulo~\ref{sec:code}, não foi decifrado por completo.

Há uma segunda particularidade, conhecida do capítulo~\ref{sec:motorlearning}. O decodificador e o cérebro aprendem um com o outro. O cérebro vê para onde o cursor foi, recebe o erro e ajusta seus comandos: é o mesmo aprendizado motor pelo fio de erro. E os engenheiros de tempos em tempos reajustam os pesos do decodificador, porque os neurônios sob os fios mudam e as conexões da rede reaprendem. O resultado são dois aprendizes que se ajustam um ao outro; para que esse par não entre em descontrole, convém separar suas velocidades de aprendizado: que um aprenda depressa e o outro devagar.

\begin{keyblock}
\begin{proposition}[Por que uma sonda de mil neurônios funciona]
\label{prop:probe}
Uma interface que escuta uma fração ínfima dos neurônios consegue controlar movimento porque a atividade do grupo tem poucas dimensões e seus ruídos são correlacionados: algumas centenas de neurônios carregam quase todas as variáveis comuns. Pelo mesmo motivo, a interface extrai só alguns bits por segundo: tantos quantas são as variáveis independentes da tarefa, e não quantos cabem no fluxo de disparos. O desempenho das interfaces foi medido; quanto o decodificador vai melhorar quando o código for entendido é uma questão em aberto.
\end{proposition}
\end{keyblock}

\section{Escrita: mais difícil que a leitura}

Enviar um sinal à máquina é mais difícil que lê-lo. Um eletrodo pelo qual passa corrente não excita um neurônio, mas todos os neurônios e fios em volta, sem distinguir tipo nem sinal. Com ele não se pode escrever um código em que importa \emph{qual} neurônio disparou e \emph{quando} (capítulo~\ref{sec:code}): só se pode definir grosseiramente \emph{onde} e \emph{com que intensidade}.

Para a visão, isso basta em parte. Uma corrente por um eletrodo no córtex visual é vista pela pessoa como um ponto luminoso num certo lugar do campo visual; quando os pontos eram acesos ao mesmo tempo, até dezesseis de uma vez, uma pessoa que perdeu a visão reconhecia algumas letras e distinguia bordas de objetos~\cite{Fernandez2021}. É um código de posição, e ele pode ser escrito. Mas lembre o capítulo~\ref{sec:sensors}: o olho não envia pixels ao cérebro, e sim uma descrição comprimida (bordas, mudanças, clareamento e escurecimento) por fios separados. Escrever ``pontos de luz'' no córtex é escrever pixels numa máquina que espera na entrada um código completamente diferente. Por isso a visão artificial ainda é mais grosseira do que se esperaria pelo número de eletrodos. Há também sinais de teste que dispensam eletrodo: as ilusões visuais entram pelo olho como uma entrada incomum e tiram a máquina do seu modo normal, e cada erro aponta para um mecanismo (seção~\ref{sec:illusions}).

\section{O que a sonda não vê}

O eletrodo ouve o barramento de endereços: os disparos dos neurônios \nt{GLUT} e, pior, os dos pequenos \nt{GABA}. Mas nos capítulos~\ref{sec:serocomputer}--\ref{sec:acet} vimos que o modo de operação da máquina inteira é definido pelos registradores de difusão: as concentrações de \nt{SERO}, \nt{DOPA}, \nt{NORA}, \nt{ACET}. O eletrodo não os ouve: os neurônios de origem são pouquíssimos, e seu sinal não é um disparo junto da sonda, mas uma concentração num volume. Um depurador que vê o barramento de dados mas não vê os registradores de controle vai se surpreender sempre: a mesma entrada vai dar respostas diferentes, porque em algum lugar mudou $g$, $a$ ou $\tau_\gamma$. As concentrações podem ser medidas por sensores químicos no próprio tecido~\cite{Bunin1998}, mas esses sensores ainda não são usados em interfaces. Também fica totalmente fora do alcance da sonda a segunda saída, a lenta, da máquina: os hormônios no sangue (capítulo~\ref{sec:chemoutput}), que são medidos em amostras de sangue, não por eletrodo.

A sonda também não vê os pesos, e é neles que está guardada quase toda a memória e tudo o que foi aprendido. Só dá para adivinhá-los pela maneira como as respostas mudam. E ela não vê a dor como a pessoa a sente: no capítulo~\ref{sec:pain} vimos que a intensidade do sinal de alarme é definida pela própria máquina, pela comporta na medula espinhal e pelos reguladores de cima, de modo que pelo sensor não se pode ler o quanto dói.

\section{Resumo: depuração e questões em aberto}

\begin{table}[t]
\centering
\footnotesize
\setlength{\tabcolsep}{4pt}
\renewcommand{\arraystretch}{1.15}
\begin{tabular}{>{\raggedright}p{2.2cm}>{\raggedright}p{3.6cm}>{\raggedright}p{3.4cm}>{\raggedright\arraybackslash}p{4.0cm}}
\toprule
Técnica de depuração & O que a interface faz & O que o livro explica & O que se sabe \\
\midrule
sonda & ouve centenas a milhares de neurônios & poucas dimensões e ruído correlacionado (cap.~\ref{sec:code}) & medido \\
compressão na sonda & transmite eventos em vez do sinal bruto & capacidade do fluxo de disparos~\eqref{eq:spikeentropy} & resolvido na engenharia \\
leitura & taxas de grupo $\to$ movimento, texto, fala & o decodificador é um destinatário; código de taxa de grupo & funciona; alguns bits por segundo \\
aprendizado conjunto & cérebro e decodificador se ajustam um ao outro & fio de erro (cap.~\ref{sec:motorlearning}) & observado; as regras de ajuste estão em estudo \\
escrita & a corrente acende pontos no campo visual & o código de posição pode ser escrito, o temporal não (cap.~\ref{sec:sensors}) & pontos medidos; visão útil, ainda não \\
leitura de registradores & quase não se faz & canais de difusão (cap.~\ref{sec:serocomputer}--\ref{sec:acet}) & sensores químicos existem em experimentos \\
\bottomrule
\end{tabular}
\caption{Depurando a máquina: o que as interfaces cérebro-computador sabem fazer e o que os capítulos do livro dizem sobre elas.}
\label{tab:debugging}
\end{table}

A tabela~\ref{tab:debugging} resume o capítulo. As interfaces cérebro-computador já sabem pôr uma sonda no barramento de endereços e ler dele alguns bits por segundo, e o fato de isso funcionar se explica pelas poucas dimensões e pelo ruído correlacionado do capítulo~\ref{sec:code}. Escrever na máquina elas só sabem de forma grosseira, porque o código de entrada dela é comprimido e endereçado a fios separados. E os registradores de controle e os pesos, aquilo que torna a máquina capaz de aprender e de se ajustar, ainda são quase invisíveis para a sonda.

Cada questão em aberto do capítulo~\ref{sec:synthesis} é também uma tarefa para o depurador. Que código ler se resolverá quando as sondas ficarem mais densas e os decodificadores aprenderem a usar o tempo dos disparos, e não só as taxas. Como uma rede de várias camadas aprende pode ser verificado pondo sondas em várias camadas ao mesmo tempo e observando como suas respostas mudam durante o aprendizado. O que os canais de difusão transmitem ficará visível quando sondas químicas se somarem às elétricas. Este livro é uma tentativa de ler o esquema da máquina pelo seu comportamento e pelas leis que qualquer engenheiro conhece. Os depuradores que estão sendo construídos agora permitirão verificar esse esquema com uma sonda.

\section{Exercícios}

\begin{exercise}[\lvlA]
\label{ex:17:1}
\emph{Orçamento do rádio.} Transmitir através da pele custa $1$~nJ por bit, e o transmissor pode gastar no máximo $10$~mW. Que potência é necessária para o fluxo bruto~\eqref{eq:probedata} com $N=1000$ eletrodos e para o fluxo de disparos? Que energia por bit o fluxo bruto exigiria?
\end{exercise}

\begin{exercise}[\lvlA]
\label{ex:17:2}
\emph{Quantos bits há em cem neurônios.} O decodificador soma os disparos de $N$ neurônios em janelas de $50\ms$ com $r=20$~disparos/s, correlação de ruído entre pares $c=0.1$, e o sinal muda a taxa do grupo em $30\%$ (valor quadrático médio). Estime a taxa $\tfrac12\log_2(1+\text{SNR})$ por janela para $N=10$, $100$, $1000$ e $N\to\infty$. E com $c=0$, $N=1000$?
\end{exercise}

\begin{exercise}[\lvlB]
\label{ex:17:3}
\emph{Velocidade de uma interface-teclado.} Uma vez e meia por segundo, a interface escolhe um de $N=32$ caracteres: o pretendido com probabilidade $P$, e os erros são equiprováveis. Quantos bits por segundo ela transmite com $P=0.99$, $0.94$, $0.8$? Quantos caracteres por segundo são necessários com $P=0.94$ para $10$~bit/s?
\end{exercise}

\begin{exercise}[\lvlB]
\label{ex:17:4}
\emph{Simulação: decodificador de grupo.} $N$ neurônios com $\lambda_i=[b+m\,\mathbf v\cdot\mathbf p_i]_+$, $b=20$, $m=15$~disparos/s, janelas de $50\ms$, $\mathbf v$ com desvio padrão $0.5$ por componente; todos veem um ruído comum, $\mathbf v+\boldsymbol\xi$, $\sigma_\xi=0.2$. Simule o decodificador não enviesado de vetor de população. Com que $N$ os dois ruídos se igualam, e quantos bits por componente por janela dá $N\to\infty$?
\end{exercise}

\begin{exercise}[\lvlB]
\label{ex:17:5}
\emph{Dois aprendizes.} O cérebro emite o comando $m=b\,v^*$, o decodificador, $v=d\,m$, com $\langle v^{*2}\rangle=1$. Após cada tentativa, ambos dão um passo de gradiente em $\tfrac12(v-v^*)^2$ com taxa $\eta$. Simule a partir de $b=1.2$, $d=1$ com $\eta=0.8$, $1.1$, $1.5$, $2$. Cada um sozinho é estável para $\eta<2$; com que $\eta$ o par é estável?
\end{exercise}

\begin{exercise}[\lvlB]
\label{ex:17:6}
\emph{Projeto: o pipeline da sonda.} $1024$ canais a $20$~mil amostras por segundo, neurônios a $10$~disparos/s, rádio a $1$~nJ/bit. Que limiar, sobre ruído branco das amostras, dá no máximo $0.1$~Hz de disparos falsos por canal? Transmitimos as contagens de disparos em $20\ms$ com $2$~bits: qual é a potência, e quantas vezes ela é menor que para amostras brutas de $10$~bits?
\end{exercise}

\begin{exercise}[\lvlC]
\label{ex:17:7}
\emph{O limite de velocidade da interface.} Estime $R\le DW\log_2(1+\text{SNR})$ para $D=10$ variáveis com banda $W=5$~Hz, com $\text{SNR}\approx0.9$ (ruído parecido com o sinal) e $90$ (ruído independente de $1000$ neurônios). O medido é cerca de $10$~bit/s. Que experimento distinguiria as duas explicações da diferença: ruído parecido com o sinal ou desconhecimento do código?
\end{exercise}

\chapter{Tipos de neurônios pela função elétrica}
\chaptermark{Neurônios como componentes}
\label{sec:eetypes}
\begin{chaptergoals}
\item os neurônios como componentes: integradores, filtros, conversores, osciladores e latches;
\item como uma célula vira um ressonador passa-faixa, ou um latch que o \nt{SERO} liga;
\item como a realimentação faz um integrador sem vazamento de neurônios que vazam, mas exige ajuste fino;
\item por que são modos de uma célula, e não tipos, definidos pelos canais iônicos e pelos de difusão.
\end{chaptergoals}

No capítulo~\ref{sec:neurontypes} classificamos os neurônios pelo neurotransmissor que sua saída libera. Um engenheiro eletrônico os classificaria de outro modo: pelo que o neurônio faz com seu sinal, isto é, como que componente ele funciona. O componente principal do livro nós conhecemos desde o capítulo~\ref{sec:glutcomputer}: o neurônio integrador com vazamento~\eqref{eq:glutneuron}, um circuito RC com comparador. Mas no cérebro há outros. Na tabela~\ref{tab:eetypes} eles estão reunidos em quatro grupos, e quase todos já apareceram no livro.

\begin{center}
\footnotesize
\setlength{\tabcolsep}{3pt}
\renewcommand{\arraystretch}{1.15}
\begin{longtable}{>{\raggedright}p{3.1cm}>{\raggedright}p{3.3cm}>{\raggedright}p{4.7cm}>{\raggedright\arraybackslash}p{2.4cm}}
\caption{Neurônios como componentes: nome, análogo eletrônico, o que o componente faz e onde aparece no livro.}\label{tab:eetypes}\\
\toprule
Neurônio & Componente & O que faz & Onde no livro \\
\midrule
\endfirsthead
\toprule
Neurônio & Componente & O que faz & Onde no livro \\
\midrule
\endhead
\bottomrule
\endfoot
\multicolumn{4}{l}{\emph{Filtros e integradores}} \\
neurônio integrador com vazamento & integrador RC com comparador & soma as entradas numa janela $\tau$ e dispara no limiar & cap.~\ref{sec:glutcomputer}, \eqref{eq:glutneuron} \\
detector de coincidência & porta E com janela temporal & dispara só se as entradas chegam quase juntas; $\tau$ muito pequeno & cap.~\ref{sec:glutcomputer}, \ref{sec:code}, \eqref{eq:window} \\
neurônio fásico & filtro passa-altas, detector de borda & responde à mudança da entrada e se cala & cap.~\ref{sec:sensors} \\
neurônio adaptativo & filtro passa-altas com controle automático de ganho & a taxa cai com entrada constante & cap.~\ref{sec:sensors}, \eqref{eq:nakarushton} \\
ressonador & circuito LC, filtro passa-faixa & responde melhor à entrada de sua própria frequência & seção~\ref{sec:resonator} \\
integrador sem vazamento & integrador com amplificador operacional & converte velocidade em posição e a mantém & seção~\ref{sec:perfectint} \\
\midrule
\multicolumn{4}{l}{\emph{Conversores}} \\
neurônio tônico & conversor tensão--frequência & a taxa segue a entrada linearmente e quase não se cansa & cap.~\ref{sec:code} \\
neurônio graduado (sem disparos) & amplificador analógico, buffer & transmite uma tensão contínua a curta distância & cap.~\ref{sec:sensors} \\
neurônio retificador & retificador de meia onda & a taxa nunca é negativa, $[u]_+$ & cap.~\ref{sec:glutcomputer}, \ref{sec:gaba} \\
par ``liga--desliga'' & par de retificadores, código de dois fios & um responde a ``mais'', o outro a ``menos'' & cap.~\ref{sec:glutcomputer}, \eqref{eq:dualrail} \\
neurônio com inibição por derivação & divisor analógico, controle de ganho & divide a entrada em vez de subtrair dela & cap.~\ref{sec:gaba}, \eqref{eq:shunt} \\
\midrule
\multicolumn{4}{l}{\emph{Geradores e tempo}} \\
marca-passo & oscilador de relaxação, multivibrador & emite sozinho disparos com frequência constante & cap.~\ref{sec:serocomputer}, \ref{sec:dofa} \\
marca-passo com entrada & oscilador controlado por tensão & ritmo próprio, que a entrada acelera ou desacelera & cap.~\ref{sec:chemoutput} \\
neurônio de rajadas & gerador de rajadas com porta & emite rajadas de disparos rápidos, os ``traços'' & cap.~\ref{sec:code}, \eqref{eq:burst} \\
neurônio travado em fase & detector de fase, malha de captura de fase & dispara numa fase definida da onda de entrada & cap.~\ref{sec:sensors}, \ref{sec:chemoutput} \\
axônio com atraso & linha de atraso & atrasa o sinal por um tempo dado & cap.~\ref{sec:glutcomputer}, fig.~\ref{fig:itd} \\
\midrule
\multicolumn{4}{l}{\emph{Memória e comutação}} \\
neurônio biestável & disparador Schmitt, latch & uma entrada curta o leva a um estado duradouro & seção~\ref{sec:bistable} \\
neurônio com ramos dendríticos & multiplexador, pequena rede de várias camadas & o ramo só deixa passar a entrada se houver uma entrada de habilitação & cap.~\ref{sec:synthesis} \\
\end{longtable}
\end{center}

Uma ressalva é mais importante que a tabela inteira: não são células diferentes, e sim \emph{modos} diferentes. O mesmo neurônio pode ser um integrador com entrada lenta e um detector de coincidência quando a inibição encurta seu $\tau$ (capítulo~\ref{sec:code}); um marca-passo na vigília e um receptor silencioso no sono. Que componente resulta é decidido pelos canais iônicos da membrana e pelos canais de difusão que os ajustam.

\section{Quatro componentes que ainda não apareceram no livro}

\subsection{O ressonador}
\label{sec:resonator}

Alguns neurônios têm uma corrente lenta que se opõe a qualquer mudança de tensão: se a tensão sobe, a corrente a puxa para baixo; se cai, puxa para cima, mas com atraso. Para o eletrônico, essa corrente se comporta como uma indutância, e junto com a capacitância da membrana forma um circuito ressonante. O neurônio responde mais forte à entrada de uma certa frequência, em geral de alguns hertz a uma ou duas dezenas de hertz, e mais fraco a entradas mais lentas e mais rápidas~\cite{HutcheonYarom2000}. É um filtro passa-faixa numa só célula: de uma entrada ruidosa ele escolhe o ritmo da sua frequência, por exemplo o ritmo local do capítulo~\ref{sec:gaba}.

\subsection{O integrador sem vazamento}
\label{sec:perfectint}

Um integrador com vazamento lembra a entrada apenas pelo tempo $\tau$: dezenas de milissegundos. Mas o olho, para ficar parado, precisa lembrar sua posição por segundos: o comando ``girar'' chega como uma velocidade breve, e o olho precisa ser mantido na nova posição. A rede resolve isso com a mesma realimentação da memória do capítulo~\ref{sec:glutcomputer}. Se um grupo de neurônios com constante de tempo $\tau$ excita a si mesmo com peso $w$, então $\tau\,\dd r/\dd t=-r+w\,r+I$, e a constante de tempo do grupo é
\begin{empheq}[box=\keybox]{equation}
\tau_{\text{ef}} \;=\; \frac{\tau}{1-w} .
\label{eq:perfectint}
\end{empheq}
Com $\tau=0.1\,\text{s}$ e $w=0.995$ obtém-se $20\,\text{s}$: um integrador quase ideal feito de neurônios que, cada um sozinho, esquecem em um décimo de segundo~\cite{Seung1996}. O preço é o ajuste fino: com $w$ um pouco acima de um, o integrador satura; um pouco abaixo, esquece depressa. Por isso esse integrador precisa de ajuste contínuo por aprendizado, como o descrito no capítulo~\ref{sec:motorlearning}.

\subsection{O neurônio biestável}
\label{sec:bistable}

Nos capítulos~\ref{sec:glutcomputer} e~\ref{sec:gaba}, o flip-flop era montado com um grupo de neurônios. Mas também existe flip-flop numa única célula. Alguns neurônios têm uma corrente que, uma vez ligada, sustenta sozinha a excitação: uma entrada curta leva o neurônio a um funcionamento prolongado, o \emph{platô}, e a inibição o devolve ao repouso. É um disparador Schmitt: o neurônio tem dois estados estáveis e histerese entre eles. É assim que se comportam, por exemplo, os neurônios \nt{ACET} que controlam os músculos, e essa propriedade é ligada por um canal de difusão: o platô aparece quando a serotonina chega ao neurônio~\cite{Hounsgaard1988}. A mesma célula, com \nt{SERO}, é um latch; sem ele, um integrador comum.

\subsection{Integradores e ressonadores}

A teoria divide as células excitáveis em duas classes~\cite{Izhikevich2007}. Os \emph{integradores} se parecem com um oscilador controlado por tensão que parte do zero: logo acima do limiar, o neurônio funciona tão devagar quanto se queira, e a taxa cresce suavemente com a entrada. Esse neurônio soma tudo o que chegou e transmite bem, pela taxa, a grandeza da entrada. Os \emph{ressonadores} se parecem com um oscilador de frequência mínima: não sabem funcionar devagar, com entrada no limiar já emitem disparos a uma frequência finita e preferem entradas da sua frequência. Os primeiros são os componentes naturais do código de taxa do capítulo~\ref{sec:code}; os segundos, do código temporal.

\begin{keyblock}
\begin{proposition}[Um neurônio, muitos componentes]
\label{prop:eetypes}
Pela função elétrica, os neurônios funcionam como integradores com e sem vazamento, detectores de coincidência, filtros passa-altas e passa-faixa, conversores tensão--frequência, retificadores, divisores, osciladores, linhas de atraso e flip-flops. Que componente uma dada célula se torna é definido por seus canais iônicos e pelos canais de difusão que os ajustam. Os modos em si foram medidos; o quanto o cérebro usa essa reconfiguração para calcular é, em grande parte, uma hipótese.
\end{proposition}
\end{keyblock}

Para o engenheiro, isso lembra um arranjo lógico programável: o hardware é um só, e o componente é definido pela configuração. Só que aqui a configuração não muda na gravação do firmware, mas em operação, pelos canais de difusão lentos de que tratamos nos capítulos~\ref{sec:serocomputer}--\ref{sec:acet}.

\section{Exercícios}

\begin{exercise}[\lvlA]
\label{ex:18:1}
\emph{Tolerância do integrador sem vazamento.} Neurônios com $\tau=0.1$~s mantêm a posição do olho segundo~\eqref{eq:perfectint} por $T=1$~s com erro de no máximo $5\%$ para qualquer lado. (a)~Ache a faixa admissível de $w$. (b)~$w$ é a soma dos pesos de $K$ sinapses, cada uma com erro independente de $10\%$. Com que $K$ a dispersão da soma cabe na tolerância?
\end{exercise}

\begin{exercise}[\lvlB]
\label{ex:18:2}
\emph{O ressonador como circuito.} Descrevemos a corrente lenta da seção~\ref{sec:resonator} pela variável $W$: $C\,\dd V/\dd t=-g_LV-g_wW+I$, $\tau_w\,\dd W/\dd t=V-W$. Mostre que isso é um circuito $RC$ com um ramo paralelo $R_w=1/g_w$, $L_w=\tau_w/g_w$ e, com $C/g_L=10\ms$, $\tau_w=100\ms$, $\alpha=g_w/g_L=4$, ache a frequência de ressonância e o ganho de $|Z|$ nela sobre $|Z(0)|$.
\end{exercise}

\begin{exercise}[\lvlB]
\label{ex:18:3}
\emph{Como o integrador começa a funcionar.} (a)~Neurônio $\tau\,\dd V/\dd t=-V+\mu$, $\tau=10\ms$, com limiar $\theta$ e reinício em zero. Com que $\mu/\theta-1$ a taxa é $1$~Hz? (b)~Que taxa dá o modelo $\tau\,\dd v/\dd t=v^2+\mu$ (disparo: $v$ vai a $+\infty$, reinício em $-\infty$) com $\mu=0.01$?
\end{exercise}

\begin{exercise}[\lvlB]
\label{ex:18:4}
\emph{Simulação: integrador e ressonador.} O modelo do exercício~\ref{ex:18:2} com $g_L=1$, $\tau_m=10\ms$, $\tau_w=100\ms$, limiar $1$ e reinício $V\to0$ ($W$ não muda) recebe $\mu=1.5\sin2\pi ft$, $f=1$--$40$~Hz. Em que faixa de frequências o neurônio emite disparos com $\alpha=0$ e com $\alpha=4$? Compare com a condição $1.5\,|Z(f)|>1$.
\end{exercise}

\begin{exercise}[\lvlB]
\label{ex:18:5}
\emph{Um latch de uma só célula.} Neurônio com corrente de platô: $\tau\,\dd r/\dd t=-r+F(wr+I)$, $F(x)=\min(1,\max(0,x))$, $\tau=10\ms$, $w=1.5$, fundo $I=-0.2$. (a)~Qual a menor duração de um pulso $I=0.3$ que leva o neurônio de $r=0$ ao platô? (b)~Qual a menor amplitude $B$ de um pulso longo $I=-0.2-B$ que o devolve ao repouso?
\end{exercise}

\begin{exercise}[\lvlB]
\label{ex:18:6}
\emph{Autoajuste do integrador.} Na fixação do olhar, a entrada é $I=0$, um sensor de deslizamento dá a velocidade de deriva $e=\dd r/\dd t$, e $\dd w/\dd t=-\eta\,e\,r$. Com $w=0.95$, $\tau=0.1$~s, $\langle r^2\rangle=1$, ache o $\eta$ com que, em $60$~s de fixações, $|1-w|$ entra na tolerância do exercício~\ref{ex:18:1}. O que quebra se aprender também durante os giros do olho?
\end{exercise}

\begin{exercise}[\lvlC]
\label{ex:18:7}
\emph{Por que reconfigurar o componente?} Um canal de difusão comuta o $\alpha$ do modelo do exercício~\ref{ex:18:2} entre $0$ e $4$. Quantas vezes o ressonador supera o integrador na relação sinal-ruído para uma senoide fraca em ruído branco de corrente a $11$~Hz e a $1$~Hz? Invente uma tarefa em que compense justamente a troca de modo.
\end{exercise}

\chapter{Neurônios pelo número de dendritos}
\chaptermark{Número de dendritos}
\label{sec:dendrites}

\section{O número de entradas como parâmetro do circuito}

No capítulo~\ref{sec:neurontypes} classificamos os neurônios pelo que sua saída libera, e no capítulo~\ref{sec:eetypes}, pelo componente com que funcionam. Há um terceiro critério, que o engenheiro nota logo: \emph{quantas entradas o neurônio tem}. As entradas são coletadas pelos dendritos, os ramos em que se apoiam axônios alheios (capítulo~\ref{sec:dofa}, seção~\ref{sec:weightstore}). Alguns neurônios não têm dendrito nenhum, outros têm um ou dois, e outros têm uma árvore que coleta cem mil entradas. Para o eletrônico, isso é a capacidade de carga na entrada, o fan-in, e ela quase sempre corresponde à tarefa.

Por que o número e o tamanho dos dendritos são tão importantes fica claro por uma estimativa simples. A membrana é um capacitor com capacitância específica $c_m\approx1$~\textmu F/cm$^2$, e quanto maior a área $A$ do neurônio junto com seus dendritos, mais carga é preciso trazer para elevar a tensão em $\Delta V$. O tempo que uma corrente de entrada $I$ leva para isso, sem contar o vazamento, é
\begin{empheq}[box=\keybox]{equation}
t \;\approx\; \frac{c_m\,A\,\Delta V}{I} .
\label{eq:dendriteload}
\end{empheq}
Um neurônio pequeno, com alguns dendritos curtos, tem área da ordem de $2\cdot10^{-5}$~cm$^2$ e capacitância de cerca de $20$~pF. Uma única sinapse grande, com corrente de alguns nanoampères, o eleva em $20\mV$ em menos de um décimo de milissegundo. Num neurônio com árvore grande, a área é umas quinze vezes maior e a capacitância, cerca de $300$~pF; uma sinapse comum, com corrente de cerca de $0.1$~nA, o carregaria em $60\ms$, muito mais que a constante de tempo do vazamento, ou seja, nunca o carregaria. Esse neurônio precisa de dezenas de entradas simultâneas. Os números aqui são ordens de grandeza, mas a conclusão não depende deles: \emph{poucas entradas, rápido e preciso; muitas entradas, lento e pela soma}.

\begin{figure}[t]
\centering
\begin{tikzpicture}[>=Stealth,
  soma/.style={circle,draw,very thick,fill=blue!6,minimum size=0.55cm,inner sep=0pt},
  ax/.style={->,very thick,red!70!black},
  den/.style={very thick,blue!60!black}]
% 0 dendrites: sensor on a line
\begin{scope}[xshift=0cm]
\draw[den,decorate,decoration={snake,amplitude=1pt,segment length=4pt}] (-0.9,0) -- (0,0);
\node[font=\scriptsize] at (-0.9,0.3) {sensor};
\draw[ax] (0,0) -- (1.0,0);
\draw[thick] (0.1,0) -- (0.1,0.45);
\node[soma,minimum size=0.4cm] at (0.1,0.65) {};
\node[font=\scriptsize,align=center] at (0.05,-0.75) {$0$\\linha com sensor};
\end{scope}
% 1 dendrite
\begin{scope}[xshift=3.3cm]
\draw[den] (-0.9,0) -- (-0.28,0);
\node[soma] at (0,0) {};
\draw[ax] (0.28,0) -- (0.9,0);
\node[font=\scriptsize,align=center] at (0,-0.75) {$1$\\buffer, repetidor};
\end{scope}
% 2 dendrites: one per ear
\begin{scope}[xshift=6.2cm]
\draw[den] (-0.28,0) -- (-0.9,0); \draw[den] (0.28,0) -- (0.9,0);
\node[soma] at (0,0) {};
\draw[ax] (0,-0.28) -- (0,-0.6);
\node[font=\scriptsize] at (-0.9,0.25) {esq.}; \node[font=\scriptsize] at (0.9,0.25) {dir.};
\node[font=\scriptsize,align=center] at (0,-1.0) {$2$\\``E'' no tempo};
\end{scope}
% ~4 short dendrites: mixer
\begin{scope}[xshift=9.0cm]
\foreach \a in {45,135,225,315}{\draw[den] (\a:0.28) -- (\a:0.7); \fill[blue!60!black] (\a:0.72) circle (0.06);}
\node[soma] at (0,0) {};
\draw[ax] (0.28,0) -- (1.0,0);
\node[font=\scriptsize,align=center] at (0,-1.0) {$\sim4$\\misturador};
\end{scope}
% many: summing tree
\begin{scope}[xshift=12.0cm]
\foreach \a in {60,90,120}{\draw[den] (0,0.28) -- ++(\a:0.55) coordinate (b); \foreach \d in {-20,20}{\draw[den] (b) -- ++(\a+\d:0.35);}}
\foreach \a in {-120,-60}{\draw[den] (0,-0.28) -- ++(\a:0.45);}
\node[soma] at (0,0) {};
\draw[ax] (0.28,0) -- (1.0,0);
\node[font=\scriptsize,align=center] at (0,-1.0) {$10^4$--$10^5$\\somador};
\end{scope}
\end{tikzpicture}
\caption{Cinco classes pelo número de entradas. Em azul, os dendritos (entradas); em vermelho, o axônio (saída). Quanto menos entradas, mais rápido e preciso o neurônio transmite um sinal isolado; quanto mais, melhor ele soma e classifica. É um esquema, não um desenho da forma das células.}
\label{fig:dendriteclasses}
\end{figure}

\section{Zero dendritos: sensor na ponta da linha}

Os neurônios sensoriais do tato, da dor e do estiramento muscular (capítulos~\ref{sec:sensors}, \ref{sec:pain}, \ref{sec:muscles}) não têm nenhum dendrito no corpo celular. O axônio sai do corpo como um só galho e logo se divide em dois: um ramo vai à pele ou ao músculo, e sua ponta serve ela mesma de sensor; o outro vai à medula espinhal. O disparo corre do sensor direto à medula, sem passar pelo corpo celular. Para o engenheiro, é uma linha de comunicação com o sensor na ponta distante e o corpo celular pendurado ao lado, como uma fonte de alimentação: ele abastece a linha, mas não participa da transmissão. A vantagem é velocidade e confiabilidade: no caminho não há nenhum ponto em que o sinal tenha de ser somado para decidir de novo se será transmitido.

Os sensores de luz e de som são um caso ainda mais extremo: não são neurônios coletores, mas os próprios transdutores, cuja entrada não é uma sinapse, e sim uma proteína receptora (fig.~\ref{fig:transducer}).

\section{Um dendrito: o buffer}

Um neurônio com um dendrito e um axônio recebe uma linha de entrada e a passa adiante. Assim são os neurônios olfatórios: o único dendrito sai à superfície do nariz e carrega um único tipo de receptor~\cite{Mombaerts2004}; assim são as células de transmissão da retina, entre os sensores de luz e os neurônios de saída do olho; assim são os neurônios que ligam os sensores do ouvido ao cérebro. O engenheiro os chamaria de buffer ou repetidor: eles não calculam, entregam um sinal, às vezes mudando sua forma, por exemplo convertendo a tensão contínua do sensor em disparos, como no capítulo~\ref{sec:sensors}.

\section{Dois dendritos: ``E'' no tempo}

Os detectores de coincidência que determinam a direção do som (fig.~\ref{fig:itd}) muitas vezes têm, nos mamíferos, exatamente dois dendritos voltados para lados opostos: num chegam as entradas do ouvido esquerdo, no outro, as do direito~\cite{Grothe2010}. Por que separar as entradas? Lembremos a fórmula~\eqref{eq:shunt}: quando muitos canais excitatórios estão abertos num trecho da membrana, a tensão ali se aproxima do potencial de equilíbrio, e cada nova entrada acrescenta cada vez menos. Por isso muitas entradas de um só ouvido, reunidas num dendrito, o saturam e não conseguem sozinhas levar o neurônio ao limiar. Já uma entrada de cada dendrito se soma de forma quase linear. Dois dendritos tornam o neurônio uma porta ``E'' mais estrita: dispare só se os sinais vieram dos dois lados e juntos. Modelos mostram que essa separação de fato melhora a detecção de coincidências~\cite{AgmonSnir1998}.

\section{Alguns dendritos curtos: misturador e retransmissor}

\emph{Misturadores.} No circuito de aprendizado motor do capítulo~\ref{sec:motorlearning}, cerca de setenta bilhões de pequenos neurônios \nt{GLUT} expandem a entrada num vetor muito longo. Cada um deles tem só uns quatro dendritos curtos, e cada dendrito termina numa ``garra'' que abraça o terminal de uma entrada. Assim, cada misturador vê só umas quatro entradas entre muitas e funciona como uma pequena porta ``E'' sobre a sua quadra. De $M$ linhas de entrada podem-se escolher $\binom{M}{4}$ quadras diferentes: com $M=100$, são quase quatro milhões. Para que tantas? Um elemento de limiar com $n$ entradas consegue separar corretamente em duas classes no máximo
\begin{empheq}[box=\keybox]{equation}
P_{\max} \;\approx\; 2n
\label{eq:cover}
\end{empheq}
padrões aleatórios~\cite{Cover1965}. O neurônio de saída do mesmo circuito recebe cerca de $10^5$ entradas dos misturadores, e por~\eqref{eq:cover} o limite superior para ele é da ordem de $2\cdot10^5$ associações diferentes. É o limite de um elemento ideal: se todos os pesos são positivos, como nas sinapses excitatórias, e é preciso margem para o ruído, a estimativa para esse mesmo neurônio é de cerca de $4\cdot10^4$ associações~\cite{Brunel2004}. Os misturadores de poucas entradas existem justamente para que o grande somador tenha muitas entradas diferentes e quase independentes~\cite{Marr1969,Albus1971}.

\emph{Retransmissores.} No caminho do ouvido até os detectores de direção há neurônios com poucos dendritos curtos que recebem só algumas sinapses enormes, e alguns, essencialmente uma. Por~\eqref{eq:dendriteload}, é exatamente o necessário: capacitância pequena e corrente grande, e o neurônio dispara frações de milissegundo depois de cada disparo de entrada, quase sem perder precisão temporal. Um desses retransmissores recebe \nt{GLUT} e emite \nt{GLYC}: é um inversor com precisão de dezenas de microssegundos, que fornece inibição rápida aos detectores de direção~\cite{BorstSoria2012}.

\section{Milhares de entradas: somador e classificador}

Na outra ponta da escala estão os neurônios \nt{GLUT} do córtex, com uma dezena de milhares de entradas, e os neurônios \nt{GABA} de saída do circuito de aprendizado motor, com uma centena de milhares. Por~\eqref{eq:dendriteload}, esse neurônio é lento e não consegue responder a uma entrada isolada: ele soma. É o neurônio integrador do capítulo~\ref{sec:glutcomputer} e o somador com que se constroem os classificadores. A árvore grande acrescenta também algo novo: seus ramos funcionam como subcircuitos separados, com limiares próprios, de modo que um neurônio se comporta como uma pequena rede de várias camadas~\cite{Beniaguev2021,Gidon2020}.

\section{Dendritos que também são saída}

Há neurônios sem axônio nenhum, cujos dendritos servem de entrada e de saída: recebem sinais e eles mesmos liberam neurotransmissor. O exemplo mais estudado são os neurônios \nt{GABA} da retina que participam da detecção da direção do movimento (capítulo~\ref{sec:gaba}, fig.~\ref{fig:direction}). Nesse neurônio, cada ramo, por conta própria, responde ao movimento do corpo celular para a sua ponta e emite inibição por essa ponta~\cite{Euler2002}. Um neurônio são vários detectores de direção independentes, um por ramo. Para o engenheiro, não é um elemento, e sim um chip com vários canais num só encapsulamento.

\section{Resumo}

\begin{table}[t]
\centering
\footnotesize
\setlength{\tabcolsep}{3pt}
\renewcommand{\arraystretch}{1.15}
\begin{tabular}{>{\raggedright}p{1.9cm}>{\raggedright}p{3.4cm}>{\raggedright}p{2.6cm}>{\raggedright}p{1.8cm}>{\raggedright\arraybackslash}p{3.8cm}}
\toprule
Dendritos & Exemplo & Componente & Entradas & O que se sabe \\
\midrule
$0$ & sensores de tato, dor, estiramento & linha com sensor na ponta & --- & estabelecido \\
$1$ & neurônios olfatórios, células de transmissão da retina, neurônios do ouvido & buffer, repetidor & $1$--poucas & estabelecido \\
$2$ & detectores de direção do som & ``E'' no tempo & dois feixes & estrutura estabelecida; ganho da separação das entradas mostrado em modelos \\
$\sim4$ & misturadores do circuito de aprendizado motor & pequena porta ``E'' & $\sim4$ & estabelecido; papel de expansão da entrada: hipótese bem fundamentada \\
poucos, curtos & retransmissores no caminho do ouvido & retransmissor, inversor & $1$--algumas enormes & estabelecido \\
milhares & neurônios \nt{GLUT} do córtex, neurônios de saída do aprendizado motor & somador, classificador & $10^4$--$10^5$ & estabelecido; ramos como subcircuitos medidos em exemplos isolados \\
dendritos-saída & neurônios \nt{GABA} de direção na retina & chip multicanal & muitas & medido \\
\bottomrule
\end{tabular}
\caption{Neurônios pelo número de dendritos.}
\label{tab:dendrites}
\end{table}

\begin{keyblock}
\begin{proposition}[O número de entradas segue a tarefa]
\label{prop:dendrites}
Onde é preciso velocidade e precisão para um sinal isolado (nos sensores, retransmissores e detectores de coincidência), os neurônios têm poucos dendritos, poucas entradas e sinapses grandes: a capacitância pequena~\eqref{eq:dendriteload} carrega depressa. Onde é preciso somar e classificar, os neurônios têm árvores grandes com milhares de entradas. A estrutura das classes foi medida; a explicação computacional é bem fundamentada, mas para muitas classes continua sendo hipótese.
\end{proposition}
\end{keyblock}

A tabela~\ref{tab:dendrites} completa as duas classificações anteriores: pelo neurotransmissor (tabela~\ref{tab:types}) e pelo componente (tabela~\ref{tab:eetypes}). As três olham o mesmo elemento de lados diferentes: o que ele emite, o que ele calcula e quanto ele escuta.

\section{Exercícios}

\begin{exercise}[\lvlA]
\label{ex:19:1}
\emph{Quantas entradas um neurônio grande precisa.} Um neurônio com $C=300$~pF e $\tau_m=20\ms$ recebe sinapses que são fontes de corrente de $0.1$~nA cada, com limiar $20\mV$ acima do repouso. (a)~Quantas sinapses precisam agir juntas para chegar ao limiar; em $10\ms$; em $5\ms$? (b)~Em quanto tempo chega ao limiar um neurônio com $C=20$~pF por uma sinapse de $4$~nA?
\end{exercise}

\begin{exercise}[\lvlA]
\label{ex:19:2}
\emph{Por que quatro.} O misturador é um ``E'' sobre $k$ linhas de $M=100$; num padrão, metade das linhas está ativa; há $7\cdot10^{10}$ misturadores. (a)~Quantos respondem ao padrão com $k=2$; $4$; $40$? (b)~Quantas vezes os misturadores aumentam o número de associações~\eqref{eq:cover} do neurônio de saída com $10^5$ entradas em relação a $M=100$ linhas diretas?
\end{exercise}

\begin{exercise}[\lvlB]
\label{ex:19:3}
\emph{De onde vem o $2n$.} Um hiperplano pela origem divide $P$ pontos em posição geral no espaço $n$-dimensional em duas classes de $C(P,n)=2\sum_{k=0}^{n-1}\binom{P-1}{k}$ modos. (a)~Prove que, com $P=2n$, exatamente metade das $2^P$ partições é separável. (b)~Que fração é separável com $n=100$ e $P=180$; $220$?
\end{exercise}

\begin{exercise}[\lvlB]
\label{ex:19:4}
\emph{Dois dendritos como ``E'' estrito.} O dendrito é um trecho com vazamento $g_L$; cada entrada abre nele uma condutância $g_0=0.5\,g_L$ com potencial $E_E$; o corpo vê $\kappa(V_1+V_2)$; limiar $\theta=1.1\,\kappa E_E$. Por que nenhum número de entradas num só dendrito dispara o neurônio, e quantas entradas são necessárias em cada um?
\end{exercise}

\begin{exercise}[\lvlB]
\label{ex:19:5}
\emph{Projeto de um retransmissor.} O jitter do disparo é $\sigma_t=\sigma_V/\dot V$, onde $\dot V$ é a inclinação da tensão no limiar. Quer-se $\sigma_t\le20$~\textmu s com ruído $\sigma_V=1\mV$; despreze o vazamento. Quantas sinapses síncronas de $0.1$~nA são necessárias para isso num neurônio com $C=300$~pF, e qual é o jitter de um neurônio com $C=20$~pF e uma sinapse de $4$~nA?
\end{exercise}

\begin{exercise}[\lvlB]
\label{ex:19:6}
\emph{Simulação: carga de um dendrito longo.} Simule o cabo passivo $\tau_m\partial_tV=\lambda^2\partial_x^2V-V+i/g$ de comprimento $L=\ell/\lambda$ com extremidades fechadas ($40$ segmentos, Euler). Um pulso curto de corrente entra na ponta distante: quando e até que fração da tensão da mesma carga espalhada por todo o dendrito sobe a ponta próxima, com $L=0.2$; $1$; $2$?
\end{exercise}

\begin{exercise}[\lvlC]
\label{ex:19:7}
\emph{Pesquisa: o número ótimo de entradas.} Capacitância $C=C_0+nc_s$; $m$ entradas agem ao mesmo tempo com corrente $I_s$; o neurônio memoriza $P\approx2n$ associações~\eqref{eq:cover}. Como o tempo de resposta depende de $P$ com $m$ constante e com fração constante $m=\varphi n$? Proponha um critério de custo e ache o $n$ ótimo para um retransmissor e para um classificador.
\end{exercise}

\chapter{O que os neurônios fazem durante o sono}
\chaptermark{Sono}
\label{sec:sleep}
\begin{chaptergoals}
\item por que o sono é um conjunto de modos trocados por \nt{ACET}, \nt{NORA} e \nt{SERO}, e não um desligamento;
\item como a pressão do sono e dois limiares formam um relé com histerese ajustado ao ritmo diário;
\item como a fadiga transforma o córtex num oscilador de relaxação, e como o \nt{ACET} detém as oscilações;
\item como o dia é reproduzido em avanço rápido, e por que os pesos podem ser reduzidos durante o sono.
\end{chaptergoals}

\section{O sono não é desligamento, e sim outro modo}

Um engenheiro, ao ver um computador desligado, diz: ele não faz nada. Com o cérebro adormecido não é assim. No sono os neurônios não se calam: no sono profundo o córtex emite disparos, e no sono REM quase tantos quanto de dia. O que muda não é \emph{se} os neurônios funcionam, mas \emph{como} funcionam. O sono é um conjunto de modos da máquina, e quem os troca são os mesmos canais de difusão que atuam de dia.

Para cada modo pode-se escrever o ``estado dos registradores'' (tabela~\ref{tab:sleepmodes}). De dia, \nt{ACET}, \nt{NORA} e \nt{SERO} funcionam, os sensores estão conectados, os músculos obedecem. No sono de ondas lentas, o sono profundo, os três canais se aquietam, e a entrada dos órgãos dos sentidos fica atenuada~\cite{Hasselmo1999}: a liberação de acetilcolina no córtex cai, e os neurônios \nt{NORA} e \nt{SERO} disparam menos que de dia~\cite{Marrosu1995,AstonJonesBloom1981,McGintyHarper1976}. No sono REM o quadro é estranho: o \nt{ACET} volta a subir ao nível diurno~\cite{Marrosu1995}, enquanto \nt{NORA} e \nt{SERO} se calam quase por completo~\cite{AstonJonesBloom1981,McGintyHarper1976,JacobsAzmitia1992}. No sono REM os músculos do corpo ficam desligados: os neurônios \nt{ACET} que os controlam são fortemente inibidos por \nt{GABA} e \nt{GLYC}~\cite{BrooksPeever2012}, a mesma proibição do capítulo~\ref{sec:gaba}, só que imposta à saída inteira de uma vez.

\begin{table}[t]
\centering
\footnotesize
\setlength{\tabcolsep}{4pt}
\renewcommand{\arraystretch}{1.15}
\begin{tabular}{>{\raggedright}p{3.1cm}>{\raggedright}p{2.9cm}>{\raggedright}p{3.2cm}>{\raggedright\arraybackslash}p{3.4cm}}
\toprule
& vigília & sono de ondas lentas & sono REM \\
\midrule
\nt{ACET} (escrita, cap.~\ref{sec:acet}) & alto & baixo & alto \\
\nt{NORA} (ganho, cap.~\ref{sec:nora}) & médio, surtos & baixo & quase calado \\
\nt{SERO} (cap.~\ref{sec:serocomputer}) & médio & baixo & quase calado \\
entrada dos sensores & conectada & atenuada & atenuada \\
saída para os músculos & conectada & atenuada & desligada por inibição \\
atividade do córtex & contínua & gangorra lenta: atividade --- silêncio & contínua, como de dia \\
\bottomrule
\end{tabular}
\caption{O estado dos registradores da máquina nos três modos. Todas as linhas foram medidas; os níveis de \nt{ACET}, \nt{NORA} e \nt{SERO}, por registros diretos em cada estágio do sono~\cite{Marrosu1995,AstonJonesBloom1981,McGintyHarper1976}.}
\label{tab:sleepmodes}
\end{table}

\section{Quando dormir: um relé com histerese}

O que decide quando a máquina passa para o sono? Um esquema simples de duas partes funciona bem~\cite{Borbely1982}. A primeira parte é a \emph{pressão do sono} $S$, que se acumula enquanto estamos acordados e se descarrega enquanto dormimos. É um capacitor que carrega de dia e descarrega à noite:
\begin{empheq}[box=\keybox]{equation}
\frac{\dd S}{\dd t} \;=\; \frac{1-S}{\tau_r}\ \ \text{(vigília)},\qquad
\frac{\dd S}{\dd t} \;=\; -\frac{S}{\tau_d}\ \ \text{(sono)} .
\label{eq:sleeppressure}
\end{empheq}
A segunda parte são dois limiares: a máquina adormece quando $S$ chega ao limiar superior $H(t)$ e acorda quando $S$ desce ao inferior $L(t)$. Os dois limiares oscilam suavemente no compasso do ritmo diário do capítulo~\ref{sec:chemoutput}. O resultado é um relé com histerese, o disparador Schmitt da seção~\ref{sec:bistable}, e junto com o capacitor ele forma um oscilador de relaxação ajustado em fase pelo ritmo diário~\eqref{eq:pll}.

\begin{figure}[t]
\centering
\begin{tikzpicture}[>=Stealth,x=0.16cm,y=4.2cm]
\foreach \a/\b in {0/4.3,21.2/29.3,45.4/53.4,69.5/72}{\fill[blue!8] (\a,0) rectangle (\b,1);}
\draw[->,gray!70!black] (0,0) -- (75,0) node[right,font=\small,black] {$t$, h};
\draw[->,gray!70!black] (0,0) -- (0,1.08) node[left,font=\small,black] {$S$};
\foreach \t in {0,24,48,72}{\draw[gray!60!black] (\t,-0.02) -- (\t,0.02); \node[below,font=\scriptsize] at (\t,0) {\t};}
\draw[thick,densely dashed,red!70!black] plot file {figures/sleep_H.dat};
\draw[thick,densely dashed,green!45!black] plot file {figures/sleep_L.dat};
\draw[very thick,blue!60!black] plot file {figures/sleep_S.dat};
\node[font=\scriptsize,red!70!black,anchor=west] at (73,0.72) {$H$: adormecer};
\node[font=\scriptsize,green!45!black,anchor=west] at (73,0.24) {$L$: acordar};
\node[font=\scriptsize,blue!60!black] at (25.25,0.93) {sono};
\node[font=\scriptsize,blue!60!black] at (49.4,0.93) {sono};
\end{tikzpicture}
\caption{Pressão do sono $S$~\eqref{eq:sleeppressure} com $\tau_r=18.2$~h, $\tau_d=4.2$~h. De dia, $S$ carrega até o limiar superior $H$; à noite (sombreado), descarrega até o inferior $L$; os dois limiares oscilam com o ritmo diário. A máquina chega sozinha a um regime de cerca de $16$ horas de vigília e $8$ horas de sono.}
\label{fig:twoprocess}
\end{figure}

No nosso modelo, com os valores usuais $\tau_r=18.2$~h e $\tau_d=4.2$~h, a máquina chega sozinha a $16.1$ horas de vigília e $8.0$ horas de sono (fig.~\ref{fig:twoprocess}). O modelo explica também o que parece estranho: depois de quarenta horas sem dormir, a pressão chega a $0.90$, mas o sono de recuperação no modelo dura só $8.8$ horas, e não um dia a mais que o normal. A descarga é exponencial, e uma carga alta vai embora depressa; a ``dívida'' é paga não com a duração, e sim com a profundidade do sono.

O que desempenha fisicamente o papel de $S$? Um forte candidato é a adenosina, o ``contador de carga'' do apêndice~\ref{sec:othermediators}: sua concentração cresce durante a vigília e cai durante o sono~\cite{PorkkaHeiskanen1997,Dunwiddie2001}. Que a pressão do sono seja exatamente a adenosina e nada mais é uma hipótese.

\section{Sono de ondas lentas: a rede inteira como oscilador de relaxação}

No sono profundo, os neurônios do córtex funcionam em alternância, o grupo todo junto: menos de uma vez por segundo, eles se ligam juntos por algumas centenas de milissegundos (o \emph{estado ativo}) e juntos se calam (o \emph{estado silencioso}); a frequência dessa gangorra fica abaixo de $1$~Hz, e sob anestesia costuma ser de $0.3$--$0.4$~Hz~\cite{Steriade1993}. Essa gangorra lenta pode ser montada com peças que já temos. Tomemos um grupo de neurônios \nt{GLUT} com forte realimentação sobre si mesmo, a memória do capítulo~\ref{sec:glutcomputer}, que tem dois estados estáveis, e acrescentemos uma fadiga lenta, como no gerador de ritmo do capítulo~\ref{sec:muscles}:
\begin{empheq}[box=\keybox]{equation}
\tau\,\frac{\dd r}{\dd t} = -r + F\bigl(w\,r + I - a\bigr),
\qquad
\tau_a\,\frac{\dd a}{\dd t} = -a + b\,r .
\label{eq:updown}
\end{empheq}
O grupo se liga, funciona e se cansa; a fadiga $a$ consome o estado superior, e o grupo cai no silêncio; no silêncio a fadiga passa, o estado inferior desaparece, e o grupo se liga de novo. O resultado é o mesmo oscilador de relaxação da pressão do sono, só que umas oitenta mil vezes mais rápido.

\begin{figure}[t]
\centering
\begin{tikzpicture}[>=Stealth,x=2.9cm,y=1.0cm]
\begin{scope}[yshift=1.9cm]
\draw[->,gray!70!black] (0,0) -- (4.1,0);
\draw[->,gray!70!black] (0,0) -- (0,1.25) node[left,font=\small,black] {$r$};
\draw[thick,blue!60!black] plot file {figures/updown_sleep.dat};
\node[font=\scriptsize,anchor=west] at (4.15,0.5) {sono: $b=5$};
\end{scope}
\draw[->,gray!70!black] (0,0) -- (4.1,0) node[right,font=\small,black] {$t$, s};
\draw[->,gray!70!black] (0,0) -- (0,1.25) node[left,font=\small,black] {$r$};
\draw[thick,red!70!black] plot file {figures/updown_wake.dat};
\node[font=\scriptsize,anchor=west] at (4.15,0.5) {dia: $b=1$};
\foreach \t in {0,1,2,3,4}{\draw[gray!60!black] (\t,-0.05) -- (\t,0.05); \node[below,font=\scriptsize] at (\t,0) {\t};}
\end{tikzpicture}
\caption{O mesmo grupo~\eqref{eq:updown} em dois modos: $w=5$, $I=2$, $\theta=2.5$, $\tau=10\ms$, $\tau_a=0.3$~s. Com fadiga forte ($b=5$, como no sono de ondas lentas), o grupo oscila entre atividade e silêncio com período de $1.1$~s (valor do modelo; no córtex o período passa de um segundo, e sob anestesia é de cerca de $3$~s). Se enfraquecemos a fadiga ($b=1$), que é o que a acetilcolina faz, o mesmo grupo funciona de modo contínuo.}
\label{fig:updown}
\end{figure}

E o que leva a rede da gangorra ao funcionamento contínuo do dia? A acetilcolina suprime nos neurônios as correntes lentas que criam a fadiga~\cite{McCormick1992}. No nosso modelo, com fadiga forte, $b=5$, o grupo oscila com período de $1.1$~s (mais rápido que a gangorra medida, mas da mesma ordem) e fica ativo cerca de $35\%$ do tempo, na média sobre um número inteiro de períodos; com fadiga fraca, $b=1$, o mesmo grupo funciona de modo contínuo (fig.~\ref{fig:updown}). Um só registrador, o nível de \nt{ACET}, transforma o oscilador num amplificador. Por cima da gangorra lenta, no sono, vêm também surtos mais rápidos de um ritmo de $7$--$15$ oscilações por segundo; eles são gerados pelo laço \nt{GLUT}--\nt{GABA} entre o córtex e um nó de retransmissão intermediário~\cite{SteriadeMcCormickSejnowski1993}, parente do ritmo do capítulo~\ref{sec:gaba}.

\section{Replays: o dia em avanço rápido}

No capítulo~\ref{sec:acet} vimos que, no sono de ondas lentas, os neurônios que funcionaram juntos de dia voltam a disparar juntos e na mesma ordem. Acrescentemos um detalhe de engenharia: os replays são \emph{acelerados}. Uma sequência que de dia levou segundos é reproduzida no sono em décimos de segundo~\cite{Nadasdy1999}: no hipocampo, cerca de vinte vezes mais depressa~\cite{LeeWilson2002}, no córtex, seis a sete vezes~\cite{Euston2007}. E não é reproduzida a qualquer momento: os surtos curtos de replay numa região se ajustam à gangorra de atividade e silêncio e aos surtos rápidos de ritmo no córtex. Se essa coordenação é reforçada artificialmente, a memória melhora~\cite{Maingret2016}; isso já é uma relação causal, não uma coincidência.

Para que a máquina reproduz o dia em avanço rápido? O engenheiro conhece a dificuldade chamada \emph{esquecimento catastrófico}: uma rede que aprende coisas novas em sequência, uma depois da outra, estraga as antigas. A saída é \emph{intercalar}: aprender o novo misturado com o antigo, em porções pequenas, muitas vezes. A hipótese dos dois sistemas de aprendizado~\cite{McClelland1995} diz que a memória rápida grava o dia inteiro e, no sono, o reproduz aos poucos, intercalado e muitas vezes, para o córtex lento. É o mesmo par ``buffer rápido --- armazenamento lento'' do capítulo~\ref{sec:acet} e o mesmo par de aprendizes rápido e lento do capítulo~\ref{sec:motorlearning}. Os replays e sua coordenação foram medidos; que sua finalidade seja o aprendizado intercalado da memória lenta é uma hipótese bem fundamentada.

\section{Renormalização dos pesos}

De dia, as regras de Hebb do capítulo~\ref{sec:dofa} sobretudo reforçam conexões. Se só se reforça, os pesos crescem, e com eles crescem os gastos de energia e cai a sensibilidade: um neurônio com todas as entradas fortes deixa de distingui-las. Segundo a hipótese da \emph{homeostase sináptica}~\cite{TononiCirelli2014}, o sono serve também para devolver os pesos ao nível de trabalho: todas as conexões enfraquecem pelo mesmo fator, de modo que suas \emph{razões}, isto é, o que foi aprendido, se preservam. Para o engenheiro, é uma normalização, como na regra~\eqref{eq:oja}, só que feita não em operação, mas num horário à parte.

As medições diretas não contradizem isso: em camundongos, a área de contato das sinapses do córtex após o sono é cerca de $18\%$ menor que após a vigília, a redução é proporcional ao tamanho, e os contatos maiores e mais estáveis quase não mudam~\cite{deVivo2017}. O escalonamento dos pesos foi medido; que essa seja a tarefa principal do sono é uma hipótese.

\section{Sono REM: a máquina desconectada das entradas e saídas}

No sono REM o córtex funciona quase como de dia, mas a entrada dos órgãos dos sentidos está atenuada, e a saída para os músculos está desligada por inibição. Para o engenheiro, é um modo conhecido: o \emph{teste em bancada}. O circuito funciona a plena potência, mas suas entradas estão ligadas a um gerador interno, e as saídas, desconectadas dos atuadores, para não quebrar nada. Os sonhos, segundo uma das hipóteses, são justamente essa execução interna do modelo do mundo montado durante o dia; o que exatamente a rede faz nesse momento, e se aprende, não se sabe. Aqui só foi medido o que está na tabela~\ref{tab:sleepmodes}: que canal está ligado, qual está calado, e que a saída está bloqueada.

\section{Manutenção e sono local}

Por muito tempo se considerou que, no sono, os espaços entre as células se alargam e o cérebro é lavado mais depressa dos produtos do metabolismo~\cite{Xie2013}. Experimentos recentes, que mediram a lavagem de outro modo, obtiveram o contrário: no sono ela é mais lenta~\cite{Miao2024}. A questão hoje está em aberto, e vamos deixá-la em aberto.

Mais sólido é outro fato: o sono é uma propriedade de redes locais, e não da máquina inteira de uma vez. Se um animal é mantido acordado por muito tempo, trechos isolados do seu córtex, enquanto o animal está acordado e se movendo, caem por instantes no silêncio, como no sono de ondas lentas, e nesses momentos o animal erra mais~\cite{Vyazovskiy2011}. Cada rede se cansa sozinha e se comuta sozinha; o modo geral surge quando os canais de difusão comutam todas as redes de uma vez.

\begin{keyblock}
\begin{proposition}[O sono como conjunto de modos]
\label{prop:sleep}
No sono os neurônios não se desligam: os canais de difusão levam a máquina a outros modos. Quando dormir é decidido por um relé com histerese~\eqref{eq:sleeppressure}, ajustado pelo ritmo diário. No sono de ondas lentas, a rede se torna um oscilador de relaxação~\eqref{eq:updown} e reproduz o dia em avanço rápido; no REM, funciona desconectada das entradas e saídas. Os modos, os replays e o escalonamento dos pesos foram medidos; suas finalidades, o aprendizado intercalado e a renormalização, são hipóteses bem fundamentadas.
\end{proposition}
\end{keyblock}

\section{Resumo}

\begin{table}[t]
\centering
\footnotesize
\setlength{\tabcolsep}{4pt}
\renewcommand{\arraystretch}{1.15}
\begin{tabular}{>{\raggedright}p{3.2cm}>{\raggedright}p{4.4cm}>{\raggedright\arraybackslash}p{5.2cm}}
\toprule
O que acontece & Como & O que se sabe \\
\midrule
troca de modos & registradores \nt{ACET}, \nt{NORA}, \nt{SERO} (tab.~\ref{tab:sleepmodes}) & medido \\
quando dormir & capacitor $S$ e relé com histerese~\eqref{eq:sleeppressure} & o modelo descreve bem os experimentos; adenosina como $S$: hipótese \\
gangorra lenta & grupo com realimentação e fadiga~\eqref{eq:updown} & gangorra medida; o modelo é o circuito mais simples \\
o dia em avanço rápido & replays $6$--$20$ vezes mais rápidos, coordenados com a gangorra & medido; reforçar a coordenação melhora a memória \\
para que os replays & aprendizado intercalado da memória lenta & hipótese bem fundamentada \\
renormalização dos pesos & escalonamento das conexões no sono & redução dos contatos medida; papel principal do sono: hipótese \\
sono REM & funcionamento com entradas e saída desconectadas & modo medido; finalidade desconhecida \\
lavagem & alargamento dos espaços entre as células & resultados contraditórios \\
sono local & trechos isolados caem no silêncio & medido \\
\bottomrule
\end{tabular}
\caption{O que os neurônios fazem durante o sono.}
\label{tab:sleep}
\end{table}

A tabela~\ref{tab:sleep} resume o capítulo. O sono acabou se mostrando não uma pausa no funcionamento da máquina, mas sua manutenção em operação: troca de modos pelos mesmos canais de difusão, um oscilador feito das mesmas peças que o ritmo dos passos, o dia em avanço rápido para a memória lenta e a renormalização dos pesos. De dia a máquina grava; à noite, reescreve e põe em ordem o que foi gravado.

\section{Exercícios}

\begin{exercise}[\lvlA]
\label{ex:20:1}
\emph{Relé sem ritmo diário.} Sejam os limiares constantes: $H=0.67$, $L=0.17$ (médias dos limiares do modelo da fig.~\ref{fig:twoprocess}). Deduza de~\eqref{eq:sleeppressure} as durações da vigília e do sono e o período do oscilador com $\tau_r=18.2$~h, $\tau_d=4.2$~h. Por que o modelo com limiares oscilantes chega ainda assim a $16.1+8.0=24$~h? A que componente do capítulo~\ref{sec:chemoutput} isso corresponde?
\end{exercise}

\begin{exercise}[\lvlA]
\label{ex:20:2}
\emph{A dívida de sono.} Um sono iniciado com pressão $S_0$ dura $t_s=\tau_d\ln(S_0/L)$, $\tau_d=4.2$~h, e $L$ é o limiar inferior no despertar. (a)~Após $40$~h sem dormir, $S_0=0.90$, e o sono durou $8.8$~h. Qual era $L$? Quanto duraria o sono com o $L\approx0.092$ usual? (b)~Numa noite, a área dos contatos sinápticos diminui $18\%$. Quanto os pesos devem crescer durante o dia?
\end{exercise}

\begin{exercise}[\lvlB]
\label{ex:20:3}
\emph{Projeto dos limiares.} Escolha limiares constantes $H$ e $L$ com que o relé~\eqref{eq:sleeppressure} com $\tau_r=18.2$~h e $\tau_d=4.2$~h dê exatamente $16$~h de vigília e $8$~h de sono. Em que essa máquina é pior que a de limiares oscilantes se uma noite ela for acordada após $4$~horas?
\end{exercise}

\begin{exercise}[\lvlB]
\label{ex:20:4}
\emph{O período da gangorra lenta.} No modelo~\eqref{eq:updown}, $F(x)=1/\bigl(1+e^{-(x-\theta)/k}\bigr)$, $w=5$, $I=2$, $\theta=2.5$, $k=0.25$, $b=5$, $\tau\ll\tau_a=0.3$~s. (a)~Ache os pontos de retorno $a_{\text{s}}$ e $a_{\text{i}}$ da curva $r=F(wr+I-a)$, onde desaparecem a atividade e o silêncio. (b)~Tomando $r\approx1$ na atividade e $r\approx0$ no silêncio, ache o período e a fração de atividade.
\end{exercise}

\begin{exercise}[\lvlB]
\label{ex:20:5}
\emph{Quando a gangorra se desliga.} No modelo do exercício~\ref{ex:20:4}, o ponto fixo está na interseção da curva $a=wr+I-F^{-1}(r)$ com a reta $a=br$. As oscilações existem enquanto a interseção está no ramo do meio, entre os pontos de retorno. Para que $b$ isso vale? Como o \nt{ACET}, mudando $b$, transforma o oscilador num amplificador?
\end{exercise}

\begin{exercise}[\lvlB]
\label{ex:20:6}
\emph{Simulação: dívida ou fase.} Em \texttt{sleep\_models.py}, tome o primeiro despertar após $48$~h e mantenha a máquina acordada por $D=24$, $32$, $40$, $48$, $64$~h (\texttt{stay\_awake}, \texttt{days=8}). Quanto dura o sono de recuperação para cada $D$? O que determina sua duração?
\end{exercise}

\begin{exercise}[\lvlB]
\label{ex:20:7}
\emph{Simulação: esquecimento e intercalação.} Um neurônio linear $y=\mathbf w\cdot\mathbf x$, $n=40$, aprende pela regra $\mathbf w\leftarrow\mathbf w-0.5\,(y-y^*)\,\mathbf x$. As tarefas A e B têm $10$ padrões cada, $x_j\sim\mathcal N(0,1/n)$, $y^*=\pm1$. Qual é o erro quadrático médio em A após $200$ épocas de A e depois $200$ épocas de B? E após $200$ épocas de A e B intercaladas?
\end{exercise}

\begin{exercise}[\lvlC]
\label{ex:20:8}
\emph{Pesquisa: intercalação ou renormalização?} O neurônio do exercício~\ref{ex:20:7}, com limiar; duas rotinas noturnas: (i)~todos os pesos se multiplicam por $0.82$; (ii)~os padrões antigos se repetem intercalados com os novos. O que cada um faz com as razões dos pesos e as respostas do neurônio? Como separar as hipóteses pelo tamanho das sinapses após o sono?
\end{exercise}

\chapter{A máquina de neurônios vivos}
\chaptermark{A máquina de neurônios vivos}
\label{sec:livingmachine}
\begin{chaptergoals}
\item por que uma cultura sem canais de difusão dispara em salvas, e o que define seu intervalo;
\item como redes vivas mudam suas respostas sob realimentação, e o que ainda não está provado;
\item como um organoide serve de reservatório cuja leitura linear é treinada fora, no silício;
\item como as culturas foram ensinadas até agora, e por que o professor continua fora da placa.
\end{chaptergoals}

\section{Uma bancada com o circuito à vista}

No capítulo~\ref{sec:debugging} depuramos uma máquina sem esquema: a sonda ouve mil neurônios de oitenta bilhões, e todo o resto tem de ser adivinhado. Nessa situação, um engenheiro faria algo simples: montaria um pedaço pequeno do circuito numa protoboard, onde cada componente fica à vista. Com neurônios isso também é possível.

Neurônios do córtex são dissociados e cultivados sobre um chip com uma matriz de eletrodos, de algumas dezenas a dezenas de milhares. Em poucas semanas as células emitem prolongamentos e se ligam numa rede de milhares ou centenas de milhares de neurônios, em sua maioria \nt{GLUT}, com uma fração menor de \nt{GABA} que depende da preparação: em culturas de hipocampo, cerca de $6\%$ dos neurônios~\cite{Benson1994}. Cada eletrodo sabe tanto ouvir, como a sonda do capítulo~\ref{sec:debugging}, quanto injetar corrente, como um sinal de teste. A segunda variante da bancada são os \emph{organoides}: bolinhas tridimensionais de tecido nervoso com cerca de um milímetro, cultivadas a partir de células-tronco humanas~\cite{Lancaster2013}. Nessas bancadas as redes vivas funcionam por meses; há plataformas em que se podem fazer experimentos com organoides remotamente, pela internet, por mais de cem dias seguidos~\cite{Jordan2024}. Essa área se chama \emph{computação biológica} ou \emph{inteligência organoide}~\cite{Smirnova2023}.

A bancada difere do cérebro em dois pontos importantes. Ela é minúscula: tem cem mil vezes menos neurônios. E não tem canais de difusão: numa cultura de córtex não há neurônios \nt{DOPA}, \nt{SERO}, \nt{NORA} nem \nt{ACET}. É uma máquina com barramento de dados e controle de fluxo, mas sem registradores de controle. Vejamos do que ela é capaz.

\section{O que a rede faz sozinha}

Deixada em paz, a cultura não fica calada nem produz um ruído uniforme. De tempos em tempos quase toda a rede dispara de uma vez, numa \emph{salva}, e volta a se aquietar; os intervalos entre salvas vão de um segundo a minutos, conforme a cultura~\cite{Wagenaar2006}. Para um engenheiro o quadro é familiar: um amplificador com forte realimentação positiva e sem carga vira oscilador.

O mecanismo nós já conhecemos. As fortes conexões mútuas dos neurônios \nt{GLUT} disparam uma avalanche, como no capítulo~\ref{sec:gaba} quando se violam as condições da proposição~\ref{prop:ei}. A avalanche esgota rapidamente o estoque de neurotransmissor nas sinapses~\eqref{eq:depression}, e a rede se cala. O estoque se recupera com constante de tempo $\tau_{\text{rec}}$, e quando se recupera a fração $x^*$ suficiente para uma nova avalanche, vem nova salva. O intervalo entre salvas é
\begin{empheq}[box=\keybox]{equation}
T_{\text{salva}} \;\approx\; \tau_{\text{rec}}\,\ln\frac{1}{1-x^*} .
\label{eq:burstinterval}
\end{empheq}
Com $\tau_{\text{rec}}=0.8\,\text{s}$ do capítulo~\ref{sec:code} e $x^*=0.9$, dá cerca de dois segundos; com $x^*=0.99$, cerca de quatro. É uma estimativa do modelo com o $\tau_{\text{rec}}$ do livro, e ela cai na ponta curta dos intervalos medidos. A medição direta em microculturas mostra que o estoque de neurotransmissor se recupera depois de uma salva em dezenas de segundos~\cite{CohenSegal2011}, isto é, lá $\tau_{\text{rec}}$ é maior; com esse $\tau_{\text{rec}}$ a fórmula~\eqref{eq:burstinterval} dá dezenas de segundos e minutos, como nas culturas lentas. É um oscilador de relaxação, parente do gerador de ritmo do capítulo~\ref{sec:muscles}: lá quem o recarregava era a fadiga dos grupos, aqui é o estoque de neurotransmissor.

As salvas são o que a rede faz quando não tem nada para fazer. No cérebro vivo, um quadro parecido aparece no sono profundo (capítulo~\ref{sec:sleep}) e no cérebro em desenvolvimento, antes de chegarem as entradas de verdade. A semelhança é sugestiva, mas a identidade dos mecanismos é uma hipótese.

\section{Como ensinar a rede sem o professor dopamina}

Como não há \nt{DOPA} na cultura, o sinal de ``melhor ou pior'' tem de ser dado por nós mesmos, pelos eletrodos. Os experimentos mostram que a rede na bancada de fato muda suas respostas, e o mais interessante neles é \emph{com o que} ela é ensinada.

\emph{O professor que se cala.} A rede é estimulada por um eletrodo uma vez a cada um a três segundos, até que em outro eletrodo apareça a resposta desejada $50\pm10\ms$ depois do estímulo. Assim que ela aparece, a estimulação é desligada por cinco minutos. Depois de alguns ciclos assim, a resposta desejada surge logo ao estímulo~\cite{Shahaf2001}. A recompensa aqui é o silêncio: o estímulo sacode a rede, e parar o estímulo a deixa no estado que deu a resposta certa. É a descida de gradiente por perturbações aleatórias do capítulo~\ref{sec:dofa} (proposição~\ref{prop:gradient}), em que o papel do erro é feito pelo próprio fato de sacudir: sacode-se até ficar certo.

\emph{A rede que joga tênis.} Num experimento com uma bancada de matriz densa de eletrodos, cerca de um milhão de neurônios foram ligados a um jogo de computador de tênis com uma só raquete~\cite{Kagan2022}. A posição da bola era aplicada como corrente nos eletrodos ``sensoriais'': onde está a bola, por código de lugar; a que distância, por frequência. A atividade em duas regiões ``motoras'' movia a raquete para cima ou para baixo. A regra de aprendizado era incomum. Se a raquete rebatia a bola, a rede recebia um estímulo curto e \emph{previsível}; se errava, alguns segundos de corrente aleatória \emph{imprevisível}. Em vinte minutos de jogo as sequências de bolas rebatidas ficavam mais longas. Um aumento significativo dentro da sessão só ocorreu com a realimentação completa; com silêncio no lugar do estímulo, a cultura também jogava melhor no fim da sessão do que sem realimentação, mas com efeito menor, e não houve aprendizado estável de uma sessão para outra ao longo de vários dias. A análise detalhada desse experimento está no capítulo~\ref{sec:dishbrain}.

Os autores explicaram o resultado com a hipótese de que a rede age de modo a tornar sua entrada previsível~\cite{Friston2010}. É a mesma ideia da economia de disparos no capítulo~\ref{sec:code} e da predição do quadro no capítulo~\ref{sec:recognition}: a máquina gasta menos quando a entrada é esperada. Mas tanto o experimento quanto, sobretudo, as palavras dos autores sobre a ``inteligência'' da cultura provocaram críticas duras: o efeito é pequeno e pode ser explicado de forma mais simples~\cite{Balci2023}. A avaliação honesta é esta: a cultura na bancada muda de comportamento sob realimentação, isso está medido; que ela aprenda justamente a prever sua entrada é hipótese.

\emph{A rede que conduz um corpo.} De modo parecido, uma cultura foi ligada a um ``corpo'' virtual simples e ensinada a movê-lo até um alvo, escolhendo-se o padrão espaço-temporal dos estímulos de treino~\cite{Bakkum2008}. Em nenhum desses experimentos há dopamina: o professor é um padrão de corrente escolhido pelo engenheiro. O que muda quando o professor passa a ser um programa ou a própria dopamina está nas seções~\ref{sec:dishdopamine}--\ref{sec:cartpole}.

\begin{figure}[t]
\centering
\begin{tikzpicture}[>=Stealth,
  blk/.style={draw,very thick,rounded corners=2pt,align=center,font=\footnotesize,minimum height=0.8cm,fill=blue!5}]
% (a) closed loop
\node[font=\small] at (1.5,4.3) {(a) malha fechada};
\node[blk,text width=2.6cm] (G) at (1.5,3.3) {jogo: bola e raquete};
\draw[very thick,fill=orange!10,rounded corners=3pt] (0.5,0.2) rectangle (2.5,1.6);
\foreach \x in {0.7,1.0,...,2.35}{\foreach \y in {0.4,0.7,1.0,1.3}{\fill[gray!60] (\x,\y) circle (0.04);}}
\draw[->,thick,blue!60!black] (G.west) -| (-0.5,0.9) -- (0.5,0.9);
\node[font=\scriptsize,blue!60!black,anchor=east] at (-0.6,2.1) {\shortstack{bola\\$\to$ corrente}};
\draw[->,thick,red!70!black] (2.5,0.9) -- (3.5,0.9) |- (G.east);
\node[font=\scriptsize,red!70!black,anchor=west] at (3.6,2.1) {disparos $\to$ raquete};
\node[font=\scriptsize] at (1.5,-0.2) {neurônios na matriz de eletrodos};
\node[font=\scriptsize,align=center] at (1.5,-0.85) {acerto: estímulo previsível\\erro: corrente aleatória};
% (b) reservoir
\begin{scope}[xshift=7.9cm]
\node[font=\small] at (1.6,4.3) {(b) reservatório};
\node[blk,text width=1.1cm] (X) at (-0.4,1.0) {entrada\\$u(t)$};
\draw[very thick,fill=orange!10] (1.6,1.0) circle (0.8);
\foreach \a/\r in {20/0.35,80/0.5,150/0.3,210/0.55,280/0.4,330/0.2,110/0.15}{\fill[red!60!black] ({1.6+\r*cos(\a)},{1.0+\r*sin(\a)}) circle (0.05);}
\node[font=\scriptsize] at (1.6,-0.2) {organoide: sem professor};
\node[blk,text width=1.8cm,fill=green!8] (W) at (4.0,1.0) {leitura\\$W_{\text{saída}}$};
\draw[->,thick] (X) -- (0.8,1.0);
\draw[->,thick] (2.4,1.0) -- node[above,font=\scriptsize] {$\mathbf r(t)$} (W);
\draw[->,thick] (W) -- (4.0,2.3) node[above,font=\scriptsize] {$y(t)$};
\node[font=\scriptsize,green!40!black] at (4.0,0.3) {aprende com professor};
\end{scope}
\end{tikzpicture}
\caption{Duas maneiras de fazer uma rede viva calcular. (a) Malha fechada: a posição da bola é aplicada como corrente, os disparos das regiões motoras movem a raquete, e um estímulo previsível ou aleatório depois do acerto serve de professor. (b) Reservatório~\eqref{eq:reservoir}: o organoide não recebe sinal de ensino; ele decompõe a entrada em muitas respostas não lineares com memória, e só uma leitura linear externa aprende pelo erro. As conexões do próprio organoide também mudam nesse processo, sem professor.}
\label{fig:livingmachine}
\end{figure}

\section{O reservatório vivo}

Há outro jeito de usar uma rede viva: não ensiná-la de modo algum. Em engenharia isso se chama \emph{computação de reservatório}~\cite{Maass2002,JaegerHaas2004}. Toma-se um sistema não linear pronto e com memória, qualquer um, desde que sua resposta à entrada seja rica e dependa do passado recente. A entrada $u(t)$ o excita, obtém-se um vetor de respostas $\mathbf r(t)$, e só se treina a leitura linear
\begin{empheq}[box=\keybox]{equation}
y(t) \;=\; W_{\text{saída}}\,\mathbf r(t) ,
\label{eq:reservoir}
\end{empheq}
cujos pesos são ajustados pela regra dos mínimos quadrados~\eqref{eq:lms} do capítulo~\ref{sec:motorlearning}. O reservatório faz o que faziam os pequenos neurônios \nt{GLUT} do capítulo~\ref{sec:motorlearning}: decompõe a entrada num vetor longo, em que as classes desejadas passam a ser separáveis por um plano (capítulo~\ref{sec:dendrites}).

Um organoide sobre uma matriz de eletrodos serviu bem como esse reservatório. Sons de fala aplicados a ele como padrões de corrente foram distinguidos pelo falante, depois de treinada uma única leitura linear, com acurácia de cerca de $78\%$~\cite{Cai2023}. A acurácia de $78\%$ é o número que os próprios autores relatam e que a imprensa repete. O resultado é útil, mas é importante entender onde está a ``inteligência'' nele: o organoide fornece não linearidade e memória que se apaga, e o aprendizado pelo erro é feito fora, no silício (fig.~\ref{fig:livingmachine}). Ao mesmo tempo, o organoide não fica inalterado: segundo os autores, durante o treino mudou sua própria conectividade funcional, e eles chamam isso de aprendizado não supervisionado no próprio organoide~\cite{Cai2023}. Que parte da acurácia vem dessa mudança e que parte vem da leitura não foi separado.

\section{O que a bancada diz sobre a máquina}

\emph{Energia.} Pela estimativa~\eqref{eq:energy}, um disparo custa cerca de $10^{-10}$~J. Uma cultura de um milhão de neurônios, a um disparo por segundo, gasta na transmissão de sinais
\begin{empheq}[box=\keybox]{equation}
P \;\approx\; 10^{6}\times1\,\text{s}^{-1}\times10^{-10}\,\text{J} \;=\; 10^{-4}\,\text{W},
\label{eq:dishpower}
\end{empheq}
um décimo de miliwatt. A eletrônica que estimula, registra e processa esses disparos gasta ordens de grandeza a mais. Como no capítulo~\ref{sec:debugging}, o gargalo não é o cálculo, e sim a comunicação com ele.

\emph{As peças que faltam.} A bancada mostra o quanto sabe o barramento \nt{GLUT} com controle de fluxo \nt{GABA} sem registradores de controle: auto-organizar-se, gerar salvas, mudar as respostas sob realimentação e servir como um bom reservatório. O que ele não sabe fazer sem eles é decidir sozinho o que conta como sucesso. Na bancada esse sinal é dado pelo engenheiro; no cérebro, como vimos nos capítulos~\ref{sec:dofa}--\ref{sec:acet}, pelos canais de difusão.

\emph{Ética.} Os organoides são cultivados a partir de células humanas, e quanto mais complexos ficam, mais séria é a questão de se esse tecido pode sentir algo. O olhar de engenheiro sugere uma resposta parcial: a dor, no capítulo~\ref{sec:pain}, não é o sinal de um único sensor, e sim todo um sistema de alarme com portas, botão de volume e canais de difusão, que a bancada não tem. Mas isso é um argumento, não uma prova, e a própria área propõe conduzir a avaliação ética junto com os experimentos, desde o início~\cite{Smirnova2023}.

\begin{keyblock}
\begin{proposition}[A rede viva na bancada]
\label{prop:livingmachine}
Uma rede de neurônios \nt{GLUT} e \nt{GABA} sobre uma matriz de eletrodos gera salvas sozinha~\eqref{eq:burstinterval}, muda as respostas sob realimentação e pode servir de reservatório~\eqref{eq:reservoir} para uma leitura treinada fora dela. Tudo isso está medido. Que essa rede aprenda por alguma regra geral, por exemplo prever sua entrada, é hipótese, e as palavras fortes sobre sua ``inteligência'' não são aceitas pela maioria dos especialistas.
\end{proposition}
\end{keyblock}

\section{Dopamina por um clarão de luz}
\label{sec:dishdopamine}

O único experimento direto que conhecemos em que se deu dopamina a uma cultura como professor foi feito numa plataforma de organoides à qual os pesquisadores se conectam remotamente~\cite{Jordan2024}. Ao líquido que banha os organoides acrescentou-se dopamina dentro de uma ``gaiola'' fotossensível: a molécula presa não age sobre os receptores. A concentração da dopamina enjaulada é de $1$~mM. Quando é preciso recompensar a rede, ela é iluminada com ultravioleta: a gaiola se desfaz, e a dopamina sai para o volume.

A malha funciona assim. O mesmo estímulo é aplicado repetidas vezes; se ele provocou mais disparos que antes, o clarão é acionado e a dopamina é liberada. Ao longo de $240$ repetições, o número de disparos evocados em geral cresceu. Os próprios autores escrevem que um único experimento desses não permite concluir que o crescimento se deve à dopamina~\cite{Jordan2024}.

Para um engenheiro, é a primeira tentativa de montar numa placa de cultura a regra de três fatores~\eqref{eq:threefactor}: a atividade do remetente e do destinatário deixa um traço, e o sinal comum decide se a mudança se fixa. Mas aqui o sinal só pode ser positivo: há recompensa ou não há, e a montagem não fornece erro negativo $\delta<0$. Além disso, a dopamina se espalha e é removida devagar, como a concentração em~\eqref{eq:lowpass}, de modo que clarões frequentes podem simplesmente elevar seu nível geral. Para distinguir aprendizado disso, são necessários dois controles: o mesmo número de clarões em momentos aleatórios, sem relação com a resposta da rede, e os mesmos clarões sem a dopamina enjaulada. E é preciso um teste depois de um repouso: a mudança permaneceu quando a dopamina foi embora?

\section{A dopamina ensina o silício}
\label{sec:biohybrid}

No experimento inverso, células vivas ensinam a eletrônica~\cite{Keene2020}. Células que liberam dopamina são cultivadas num microcanal sobre um elemento neuromórfico orgânico, um transistor cuja condutância do canal faz o papel de peso sináptico. A dopamina que sai das células se oxida junto ao eletrodo do elemento, e isso muda a condutância. O dispositivo reproduz também o mecanismo de remoção do neurotransmissor da fenda, por isso o peso pode não só mudar por muito tempo, mas também voltar ao valor anterior.

Em termos de circuito, é um integrador com vazamento de entrada química: o peso acumula o fluxo de dopamina, e a ``remoção'' define a rapidez com que ele esquece; é o mesmo circuito RC de~\eqref{eq:lowpass}. O essencial aqui é o sentido da ligação. A sonda do capítulo~\ref{sec:debugging} não enxerga os registradores de difusão, porque eles são químicos. Este elemento lê justamente um registrador químico e o converte num peso elétrico. Para as interfaces cérebro-computador, é o caminho para um sensor que ouve não só o barramento de dados, mas também os registradores de controle.

\section{Organoides com seus próprios neurônios de dopamina}
\label{sec:midbrainorg}

A partir de células-tronco humanas já se sabe cultivar organoides parecidos com a região do cérebro onde ficam os neurônios \nt{DOPA}. Eles têm neurônios dopaminérgicos funcionais, que liberam dopamina~\cite{Jo2016}. Esses organoides são cultivados sobretudo para estudar doenças. Não encontramos experimentos em que a dopamina deles servisse de professor para outra rede.

Para a máquina de neurônios vivos, essa é a peça que falta. A bancada do capítulo~\ref{sec:livingmachine} é barramento de dados e controle de fluxo sem registradores de controle; aqui a fonte do sinal de difusão já existe. Falta o crítico: uma rede que calcule o erro de predição~\eqref{eq:ctd} e comande esses neurônios. Ligar um organoide de córtex a um organoide desses de modo que a dopamina seja liberada conforme o erro é um problema em aberto, e por enquanto uma hipótese, não um experimento.

\section{Um organoide equilibra uma haste sem dopamina}
\label{sec:cartpole}

O mais completo dos experimentos recentes dispensa totalmente a dopamina~\cite{Robbins2026}. Organoides de córtex de camundongo foram ligados em malha fechada à tarefa do ``pêndulo no carrinho'': a atividade do organoide move o carrinho, e a posição da haste volta por estimulação. Entre as tentativas, o organoide recebe estímulos de ensino curtos e de alta frequência. Quais exatamente, quem escolhe é um programa de aprendizado por reforço.

Os resultados foram medidos:
\begin{itemize}
\item na maioria dos organoides, os sinais de ensino escolhidos pelo programa deram resultado melhor que sinais escolhidos ao acaso ou a ausência de sinal;
\item depois de $45$ minutos de repouso, a melhora não se mantinha;
\item o bloqueio dos dois tipos de receptor de glutamato, AMPA e NMDA, anulava a melhora.
\end{itemize}

O último item diz onde mora o que foi aprendido: nas sinapses \nt{GLUT}, e por meio do receptor que, na seção~\ref{sec:weightstore}, funciona como detector de coincidência de entrada e saída. O segundo item mostra o que falta: há mudança rápida, mas não há consolidação, como no aprendiz rápido do capítulo~\ref{sec:motorlearning} sem o lento. E o papel do professor mudou: o engenheiro que escolhe o padrão de corrente foi substituído por um programa, mas o professor continua do lado de fora da placa.

Entre os experimentos anteriores de ensino de culturas em matrizes de eletrodos também não encontramos nenhum que usasse dopamina: o sinal de ensino sempre foi a estimulação elétrica.

\section{Quem ensina a cultura}

\begin{table}[t]
\centering
\footnotesize
\setlength{\tabcolsep}{4pt}
\renewcommand{\arraystretch}{1.15}
\begin{tabular}{>{\raggedright}p{2.7cm}>{\raggedright}p{2.5cm}>{\raggedright}p{3.2cm}>{\raggedright\arraybackslash}p{4.4cm}}
\toprule
Experimento & Quem ensina & Sinal de ensino & O que se sabe \\
\midrule
parar a estimulação na resposta desejada & engenheiro & fim da estimulação & a resposta se fixa --- medido \\
DishBrain (cap.~\ref{sec:dishbrain}) & engenheiro & corrente previsível ou aleatória & aumento significativo dos ralis na sessão só com realimentação completa; com silêncio, efeito menor --- medido; instável de uma sessão para outra; a causa está em debate \\
pêndulo no carrinho & programa de aprendizado por reforço & estímulos de alta frequência escolhidos pelo programa & melhor que os aleatórios, mas não se mantém após o repouso --- medido \\
dopamina por clarão & regra ``mais disparos = recompensa'' & dopamina liberada pela luz & aumento dos disparos --- observação; o papel da dopamina não está provado \\
sinapse bio-híbrida & células vivas & dopamina das células & mudança duradoura e retorno do peso do elemento --- medido \\
organoides com neurônios \nt{DOPA} & --- & --- & a fonte de dopamina existe; não foi usada como professor \\
\bottomrule
\end{tabular}
\caption{Quem ensina os neurônios vivos no chip, e com quê.}
\label{tab:dishteachers}
\end{table}

Em todos os experimentos em que a própria cultura aprende, o professor por enquanto está fora (tabela~\ref{tab:dishteachers}): o engenheiro, um programa ou uma regra definida pelo engenheiro. A dopamina entrou na placa uma única vez, e seu papel ainda precisa ser provado. Montar na placa um canal de difusão de verdade, com crítico, erro e traço, como no capítulo~\ref{sec:dofa}, é o próximo passo, que ninguém ainda deu.


\section{Resumo}

\begin{table}[t]
\centering
\footnotesize
\setlength{\tabcolsep}{4pt}
\renewcommand{\arraystretch}{1.15}
\begin{tabular}{>{\raggedright}p{2.8cm}>{\raggedright}p{4.2cm}>{\raggedright\arraybackslash}p{5.8cm}}
\toprule
Experimento & O que a rede viva faz & O que se sabe \\
\midrule
rede sem entrada & salvas com intervalos de um segundo a minutos~\eqref{eq:burstinterval} & medido; o mecanismo por esgotamento das sinapses é o modelo aceito \\
professor que se cala & a resposta desejada se fixa quando a estimulação é desligada & medido \\
jogo de tênis & as sequências de bolas rebatidas ficam mais longas; aumento significativo na sessão só com realimentação completa & a mudança de comportamento está medida; o aprendizado de uma sessão para outra é instável; o tamanho do efeito e a explicação pela predição da entrada são discutíveis \\
corpo virtual & movimento até o alvo com um padrão de estímulos escolhido & medido \\
reservatório & não linearidade e memória para a leitura~\eqref{eq:reservoir} & medido; o aprendizado pelo erro é externo, no silício; as conexões do organoide também mudam \\
energia & cerca de $10^{-4}$~W por milhão de neurônios~\eqref{eq:dishpower} & estimativa por~\eqref{eq:energy} \\
\bottomrule
\end{tabular}
\caption{O que as máquinas de neurônios vivos mostraram.}
\label{tab:livingmachine}
\end{table}

A máquina de neurônios vivos (tabela~\ref{tab:livingmachine}) é uma bancada em que se vê o que sabem fazer o barramento de dados e o controle de fluxo sem registradores de controle. Sabem muita coisa, mas não sabem decidir sozinhos o que é bom. Na bancada esse papel é do engenheiro com seu padrão de corrente; no cérebro, dos poucos fios de difusão a que é dedicado o meio do livro.

\section{Exercícios}

\begin{exercise}[\lvlA]
\label{ex:21:1}
\emph{Intervalo entre salvas.} Uma salva esgota o estoque até zero; depois $\tau_{\text{rec}}\,\dd x/\dd t=1-x$, e nova salva começa em $x=x^*$; $\tau_{\text{rec}}=0.8$~s. (a)~Ache o intervalo para $x^*=0.9$ e $0.99$. (b)~O limiar $x^*$ varia ao acaso $\pm0.005$ em torno de $0.99$. Em que faixa ficam os intervalos?
\end{exercise}

\begin{exercise}[\lvlA]
\label{ex:21:2}
\emph{Energia da bancada.} (a)~Que potência a estimativa~\eqref{eq:dishpower} dá para o cérebro inteiro ($8.6\cdot10^{10}$ neurônios)? (b)~A bancada registra $1024$ eletrodos a $20$~kHz, com $10$~bits por amostra. Quantas vezes transmitir esse fluxo a $1$~nJ por bit custa mais que a potência da própria cultura?
\end{exercise}

\begin{exercise}[\lvlB]
\label{ex:21:3}
\emph{Projeto: suprimir as salvas.} Uma estimulação esparsa por muitos eletrodos faz cada sinapse liberar neurotransmissor à taxa $r_b$; o estoque segue $\tau_{\text{rec}}\,\dd x/\dd t=1-x-Ur_b\tau_{\text{rec}}x$, $U=0.5$, $\tau_{\text{rec}}=0.8$~s. Qual a menor $r_b$ com que o estoque não chega a $x^*=0.9$; $0.99$, e quanto a sinapse enfraquece com isso?
\end{exercise}

\begin{exercise}[\lvlB]
\label{ex:21:4}
\emph{A memória do reservatório não passa do número de neurônios.} Um reservatório com $N$ saídas $\mathbf r(t)$ recebe uma entrada branca $u(t)$ de média zero e variância unitária; $m_k$ é o maior quadrado da correlação da leitura $\mathbf w_k\cdot\mathbf r(t)$ com $u(t-k)$. Prove que a capacidade de memória $\mathrm{MC}=\sum_{k\ge0}m_k\le N$.
\end{exercise}

\begin{exercise}[\lvlB]
\label{ex:21:5}
\emph{Simulação: um reservatório em silício.} Reservatório $\mathbf r(t+1)=\tanh(W\mathbf r(t)+\mathbf v\,u(t)+\mathbf c)$, $N=30$, $W$ aleatória com raio espectral $0.9$, $v_i\sim U(-0.5,0.5)$, $c_i\sim U(-0.2,0.2)$, $u\sim U(-1,1)$, leitura por regressão ridge. Ache a capacidade de memória do exercício~\ref{ex:21:4} com e sem $\tanh$. Qual reservatório lembra mais?
\end{exercise}

\begin{exercise}[\lvlB]
\label{ex:21:6}
\emph{A leitura de um reservatório vivo.} Um organoide dá $1024$ canais e $P=240$ registros de treino de duas classes. (a)~Quantos atributos $n$ se podem manter para que, pela fórmula do exercício~\ref{ex:19:3}, uma rotulação aleatória seja separável com probabilidade de até $1\%$? (b)~Quantos sigmas acima do acaso está a acurácia de $78\%$ em $100$ registros de teste?
\end{exercise}

\begin{exercise}[\lvlC]
\label{ex:21:7}
\emph{Pesquisa: um professor sem canal de difusão.} Proponha um protocolo de estimulação que implemente na bancada a regra de três fatores do capítulo~\ref{sec:dofa} por $8$--$1024$ eletrodos, e um controle com leitura congelada que distinga o aprendizado dos pesos de um deslocamento do estado da rede. Que resultado refutaria a hipótese de aprendizado?
\end{exercise}

\chapter{DishBrain}
\chaptermark{DishBrain}
\label{sec:dishbrain}

\section{Uma rede na placa joga tênis}

DishBrain é o nome que os autores deram à bancada em que uma cultura de neurônios vivos, cultivada sobre um chip com eletrodos, joga um tênis simplificado com uma só raquete~\cite{Kagan2022}. É o experimento mais discutido com uma máquina de neurônios vivos (um panorama dessas máquinas está no capítulo~\ref{sec:livingmachine}), e para o nosso livro ele é especialmente conveniente. Aqui a pequena rede está acessível por inteiro: cada entrada é definida pelo experimentador, cada saída é registrada. Todas as perguntas dos capítulos anteriores (como o sinal é codificado, quem ensina, o que a sonda vê) podem ser feitas a uma rede que não tem olhos, nem músculos, nem neurônios \nt{DOPA}. Vamos analisar o experimento como um engenheiro analisa a bancada de outra pessoa: esquema, entrada, saída, professor, resultado, controvérsias.

\section{A bancada}

Os neurônios vêm do córtex de embriões de camundongo, cerca de $800\,000$ células por chip, ou são cultivados a partir de células-tronco humanas, cerca de um milhão por chip. Durante algumas semanas eles crescem no chip, entrelaçam-se e começam a produzir sozinhos disparos e rajadas. Numa área de $8$~mm$^2$ sob eles há $26\,000$ eletrodos de platina; até $1024$ deles podem ser registrados ao mesmo tempo, e a estimulação, na prática, se faz por oito.

Em torno da cultura fecha-se uma malha (fig.~\ref{fig:dishloop}). O computador simula o jogo: a bola voa, bate nas paredes, a raquete se move para cima e para baixo. A posição da bola é convertida em corrente em oito eletrodos ``sensoriais''. Os disparos de duas regiões ``motoras'' são contados em janelas de $10\ms$ e movem a raquete. Depois de cada acerto ou erro, a cultura recebe mais um estímulo: a realimentação. As janelas de $10\ms$ são o ciclo do computador da bancada, não da cultura: a própria rede continua funcionando de forma assíncrona, como todas as redes deste livro.

\begin{figure}[t]
\centering
\begin{tikzpicture}[>=Stealth,
  blk/.style={draw,very thick,rounded corners=3pt,align=center,font=\small,minimum height=1.2cm},
  lab/.style={font=\scriptsize,align=center}]
\node[blk,fill=blue!6,text width=3.4cm] (C) at (0,0) {cultura no chip\\\scriptsize $\sim10^6$ neurônios, $26\,000$ eletrodos};
\node[blk,fill=orange!10,text width=3.0cm] (G) at (7.2,0) {computador:\\jogo de tênis};
\draw[->,very thick,blue!60!black] (G.north west) to[bend right=25] node[lab,above] {bola: \emph{lugar} = qual dos 8 eletrodos,\\\emph{frequência} = 4--40 disparos/s} (C.north east);
\draw[->,very thick,red!70!black] (C.south east) to[bend right=25] node[lab,below] {raquete: contagem de disparos\\das regiões ``sobe'' e ``desce'' em $10\ms$} (G.south west);
\draw[->,thick,densely dashed,green!45!black] (G.east) -- ++(0.9,0) |- ++(0,-2.6) -| node[lab,pos=0.25,below] {acerto: estímulo previsível; erro: imprevisível} (C.south);
\end{tikzpicture}
\caption{A malha do DishBrain. Seta azul: entrada; a posição da bola é codificada por código de lugar e código de taxa (capítulo~\ref{sec:code}). Vermelha: saída; a diferença entre as contagens das duas regiões move a raquete. Verde tracejada: realimentação depois de cada rali, o único ``professor'' da bancada.}
\label{fig:dishloop}
\end{figure}

\section{Entrada e saída}

A \emph{entrada} usa ao mesmo tempo dois códigos do capítulo~\ref{sec:code}. Qual dos oito eletrodos é estimulado diz a que altura está a bola: é um código de lugar, como nos sensores do capítulo~\ref{sec:sensors}. A frequência dos pulsos de estimulação diz a que distância está a bola: $4$ por segundo quando ela está junto à parede do fundo, e linearmente mais, até $40$ por segundo, quando está junto à parede da raquete. É um código de taxa; a amplitude dos pulsos da bola é de $75$~mV. A cultura não ``conhece'' nenhum desses códigos de antemão: o código foi definido pelo experimentador, e a rede tem de aprender a lê-lo.

A \emph{saída} é lida como nas interfaces do capítulo~\ref{sec:debugging}. Os disparos de uma região motora empurram a raquete para cima, os da outra, para baixo; na forma mais simples, a cada janela
\begin{empheq}[box=\keybox]{equation}
\Delta p \;=\; c\,\bigl(n_\uparrow - n_\downarrow\bigr),
\label{eq:dishreadout}
\end{empheq}
em que $n_\uparrow$, $n_\downarrow$ são os números de disparos das duas regiões em $10\ms$, e $c$ é um coeficiente da bancada. É o código de dois fios do capítulo~\ref{sec:glutcomputer}: o sentido do movimento não depende de o sinal elétrico ser positivo ou negativo, e sim de qual dos dois fios fala mais alto.

Vejamos o ruído dessa saída. Se uma região produz ao todo $R$ disparos por segundo, numa janela $T=10\ms$ há em média $RT$ deles, e por~\eqref{eq:rateerror} a dispersão relativa é $1/\sqrt{RT}$. Com $R=1000$ disparos por segundo, isso dá cerca de $30$ por cento em cada janela; com $R=100$, mais de cem. A raquete inevitavelmente treme, e a rede só pode aprender de um jeito: deslocando a diferença \emph{média} das contagens para o lado certo. São as mesmas leis que limitam qualquer interface do capítulo~\ref{sec:debugging}, só que agora dos dois lados da malha.

\section{Um professor sem dopamina}

O mais interessante na bancada é o professor. Numa cultura de neurônios do córtex não há neurônios \nt{DOPA}, e a bancada não envia nenhum sinal de ``melhor do que o esperado''. A realimentação é outra. Se a rede rebateu a bola, os oito eletrodos sensoriais recebem ao mesmo tempo um estímulo curto e \emph{previsível}: pulsos de $75$~mV, $100$ por segundo durante $0.1\,\text{s}$. Se errou, durante quatro segundos vem um estímulo que os autores chamam de \emph{imprevisível}: pulsos de $150$~mV, $5$ por segundo, em momentos aleatórios e em eletrodos escolhidos ao acaso; depois, quatro segundos de pausa sem estimulação, e a bola parte de novo~\cite{Kagan2022}.

Os autores partiram da hipótese de que a rede age de modo a tornar sua entrada previsível~\cite{Friston2010}. Para um engenheiro o sentido é simples: a cultura tem exatamente um jeito de se livrar de quatro segundos de ruído aleatório, que é rebater a bola. A mesma ideia já apareceu na economia de disparos do capítulo~\ref{sec:code} e na predição do quadro no capítulo~\ref{sec:recognition}.

Mas é preciso mesmo a previsibilidade? Um mecanismo mais grosseiro também é possível. Suponha que o estímulo imprevisível depois de um erro simplesmente \emph{embaralhe} as mudanças recentes na rede, e que o estímulo calmo depois de um acerto não mexa nelas. Descrevamos a configuração da rede por um único vetor $\boldsymbol\psi$, e seja $P_{\text{erro}}(\boldsymbol\psi)$ a probabilidade de a rede errar na configuração $\boldsymbol\psi$. Se os empurrões forem simétricos (de $\boldsymbol\psi$ para $\boldsymbol\psi'$ é tão provável quanto o inverso), no regime estacionário o fluxo de saídas de cada configuração se equilibra com o fluxo de chegadas, e a fração do tempo que a rede passa na configuração $\boldsymbol\psi$ é
\begin{empheq}[box=\keybox]{equation}
\rho(\boldsymbol\psi) \;\propto\; \frac{1}{P_{\text{erro}}(\boldsymbol\psi)} .
\label{eq:dishdensity}
\end{empheq}
Uma configuração que erra dez vezes menos é mantida dez vezes mais. Ninguém diz à rede para que lado mudar os pesos: as boas configurações simplesmente são destruídas com menos frequência.

Verificamos isso num modelo que reduz a cultura a dois números: a raquete chega ao ponto $a\,y+b$ quando a bola chega à altura $y$, e rebate se errar por menos de $0.1$. Depois de um erro, os dois números recebem um empurrão aleatório; depois de um acerto, ficam como estão. A fração de bolas rebatidas em $2000$ ralis subiu de $0.21$ nos primeiros duzentos para $0.38$ nos últimos quinhentos (quarenta ``culturas'', fig.~\ref{fig:dishmodel}). Se o empurrão vier depois de \emph{cada} rali, a realimentação não traz informação, e a fração fica em torno de $0.12$; sem empurrão nenhum, ela fica parada em $0.19$. O experimento na bancada concorda com isso: um aumento significativo dentro da sessão só ocorreu com a realimentação completa. Com silêncio no lugar dos estímulos, a cultura ainda assim jogava melhor no fim da sessão do que sem realimentação, mas com efeito menor~\cite{Kagan2022}; note que também nessa condição um erro é seguido de uma pausa longa, e um acerto, de uma curta, de modo que a diferença entre os desfechos não desaparece por completo.

\begin{figure}[t]
\centering
\begin{tikzpicture}[x=1cm,y=5cm]
\draw[->,gray!70!black] (0,0) -- (10.6,0);
\draw[->,gray!70!black] (0,0) -- (0,0.5) node[above,font=\small,black] {fração rebatida};
\foreach \v in {0.1,0.2,0.3,0.4}{\draw[gray!60!black] (-0.06,\v) -- (0.06,\v); \node[left,font=\scriptsize] at (0,\v) {\v};}
\foreach \x/\f/\l/\lab in {1.8/0.21/0.38/{ruído após\\o erro},5.2/0.13/0.11/{ruído após\\cada rali},8.6/0.19/0.19/{sem ruído}}{
  \fill[blue!35] (\x-0.8,0) rectangle (\x-0.05,\f);
  \fill[red!60!black] (\x+0.05,0) rectangle (\x+0.8,\l);
  \node[below,font=\scriptsize,align=center] at (\x,-0.01) {\lab};
  \node[above,font=\scriptsize] at (\x-0.425,\f) {\f};
  \node[above,font=\scriptsize] at (\x+0.425,\l) {\l};}
\fill[blue!35] (7.4,0.47) rectangle (7.7,0.495); \node[right,font=\scriptsize] at (7.75,0.482) {primeiros 200};
\fill[red!60!black] (7.4,0.41) rectangle (7.7,0.435); \node[right,font=\scriptsize] at (7.75,0.422) {últimos 500};
\end{tikzpicture}
\caption{O modelo ``embaralhar no erro'' (\texttt{dishbrain\_model.py}): fração de bolas rebatidas no início e no fim de $2000$ ralis, média de quarenta culturas-modelo. Só há aprendizado se o empurrão vier depois do erro, e não do acerto, como no experimento, em que o aumento significativo na sessão só ocorreu com a realimentação completa.}
\label{fig:dishmodel}
\end{figure}

\begin{keyblock}
\begin{proposition}[O professor embaralhador]
\label{prop:dishscramble}
Se o fracasso sacode a rede e o sucesso a deixa em paz, a rede deriva sozinha para configurações que erram menos: o tempo passado numa configuração é inversamente proporcional à probabilidade de fracasso~\eqref{eq:dishdensity}. Para esse aprendizado não são necessários nem um sinal de recompensa, nem o seu sinal algébrico, nem a previsibilidade em si; basta que a realimentação depois do sucesso e depois do fracasso seja diferente. A mudança de comportamento da cultura está medida; qual mecanismo atua na placa, predição da entrada ou embaralhamento, não está estabelecido.
\end{proposition}
\end{keyblock}

Compare com o capítulo~\ref{sec:dofa}. Lá o ruído também servia de busca, mas a dopamina tinha sinal algébrico: o erro $\delta$ dizia para que lado deslocar os pesos, e o passo seguia o gradiente~\eqref{eq:perturbation}. Aqui não há sinal algébrico, só ``manter'' ou ``sacudir''. Essa busca é mais grosseira e mais lenta, mas exige menos da rede: nem neurônios \nt{DOPA}, nem crítico, nem traço estendido por segundos.

\section{O que foi medido e o que se discute}

Cada jogo durava $20$ minutos; comparavam-se os primeiros $5$ minutos com os últimos $15$. Nas culturas com realimentação completa, as sequências de bolas rebatidas ficaram mais longas, e os erros imediatos, mais raros; nos controles sem células, sem jogo e com uma ``raquete'' de computador, nada parecido aconteceu. As culturas de células humanas em média rebatiam por mais tempo que as de camundongo. A melhora é estatisticamente sólida (em nove culturas de camundongo e onze humanas, em mais de cem jogos de cada grupo), mas pequena: a rede não virou um bom jogador, ficou um pouco melhor que o acaso. Ela é sólida dentro da sessão; aprendizado estável de uma sessão para outra ao longo de vários dias os autores não obtiveram~\cite{Kagan2022}. Num trabalho seguinte do mesmo grupo, verificou-se que durante o jogo a atividade da cultura se aproxima do regime \emph{crítico}, a fronteira entre a extinção e a avalanche, aquela mesma vida no limite do capítulo~\ref{sec:gaba}, e quanto mais perto, melhor o jogo~\cite{Habibollahi2023}.

A principal controvérsia não foi causada pelos números, e sim pelas palavras. O título do artigo afirmava que os neurônios na placa ``exibem senciência'' (sentience), e um grupo de pesquisadores objetou: a mudança de comportamento numa malha fechada ainda não é inteligência nem sensação, e as conclusões deveriam ser formuladas com mais modéstia~\cite{Balci2023}. Do ponto de vista de engenharia, há duas questões em aberto, ambas sobre o mecanismo. Primeira: a rede aprende, isto é, seus pesos mudam de forma duradoura, ou a realimentação só desloca seu estado atual? Segunda: é preciso justamente a previsibilidade, ou basta qualquer diferença entre a resposta ao sucesso e ao fracasso, como na proposição~\ref{prop:dishscramble}? Uma bancada em que cada eletrodo é acessível permite responder às duas, e esse talvez seja seu principal valor.

\section{Especificação da bancada: como reproduzi-la}
\label{sec:dishspec}

Abaixo está reunido tudo o que é preciso para montar essa bancada de novo: equipamento, malha, codificação da entrada, leitura da saída, realimentação e protocolo do experimento. Cada parâmetro traz sua fonte. ``Artigo'': o valor foi tirado do texto do trabalho~\cite{Kagan2022}. ``Escolha'': o artigo não o traz, e para reproduzir será preciso defini-lo por conta própria; propomos um valor padrão e dizemos como verificá-lo.

\subsection*{Equipamento e cultura}

\begin{center}
\footnotesize
\setlength{\tabcolsep}{4pt}
\renewcommand{\arraystretch}{1.15}
\begin{tabular}{>{\raggedright}p{4.0cm}>{\raggedright}p{6.9cm}>{\raggedright\arraybackslash}p{1.6cm}}
\toprule
Parâmetro & Valor & Fonte \\
\midrule
chip & matriz MaxOne: $26\,000$ eletrodos de platina numa área de $8$~mm$^2$ & artigo \\
registro & até $1024$ canais simultâneos, $20\,000$ amostras por segundo & artigo \\
estimulação & até $32$ eletrodos tecnicamente, no experimento $8$ independentes & artigo \\
células & córtex de embrião de camundongo; neurônios humanos de células-tronco (dois métodos de cultivo); células HEK293T (não neurônios) como controle & artigo \\
semeadura & cerca de $10^6$ células por chip & artigo \\
idade da cultura no experimento & camundongo: a partir de $\sim10$ dias; humana: $14$--$24$ dias ou $\sim80$ dias, conforme o método de cultivo & artigo \\
\bottomrule
\end{tabular}
\end{center}

\subsection*{Malha e tempo}

A malha funciona em janelas de $10\ms$ ($200$ amostras): em cada janela contam-se os disparos nos eletrodos das regiões motoras, o jogo é atualizado, e sai a estimulação seguinte. O atraso entre um disparo na cultura e o estímulo de resposta é de cerca de $5\ms$. Durante o jogo, o sinal de cada eletrodo passa por um filtro passa-altas Bessel de 2.ª ordem com frequência de corte de $100$~Hz; o módulo do sinal é suavizado por um filtro passa-baixas Bessel de 1.ª ordem de $1$~Hz, e o limiar de disparo é proporcional a esse módulo suavizado. O limiar de seis desvios-padrão do fundo o artigo indica para medições isoladas de atividade, não para o jogo. Para que a estimulação não seja contada como resposta da cultura, todos os sinais são suprimidos quando chegam ao mesmo tempo mais de $15$ pulsos grandes, acima de $75$~mV. Tudo isso é do artigo.

\subsection*{Codificação da bola}

\begin{center}
\footnotesize
\setlength{\tabcolsep}{4pt}
\renewcommand{\arraystretch}{1.15}
\begin{tabular}{>{\raggedright}p{3.4cm}>{\raggedright}p{7.5cm}>{\raggedright\arraybackslash}p{1.6cm}}
\toprule
Parâmetro & Valor & Fonte \\
\midrule
região sensorial & $8$ eletrodos dispostos de modo que eletrodos vizinhos signifiquem posições vizinhas da bola & artigo \\
código de lugar & o número do eletrodo $k=1,\dots,8$ dá a posição da bola em relação ao centro da raquete & artigo \\
qual coordenada & a posição vertical da bola; para reproduzir, em relação à raquete, $k=\lceil 8(y-p+\tfrac12)\rceil$ com $y-p\in[-\tfrac12,\tfrac12]$ & artigo; fórmula: escolha \\
código de taxa & a frequência de estimulação, de $4$ a $40$~disparos/s, codifica a distância até a raquete & artigo \\
sentido da escala & $4$~disparos/s com a bola junto à parede oposta, e linearmente até $40$~disparos/s junto à parede da raquete: $f=4+36\,(1-x/L)$, $x$ é a distância até a parede da raquete, $L$ a largura da quadra & artigo \\
forma do pulso & curto, retangular, bifásico: primeiro a fase positiva, depois a negativa & artigo \\
amplitude & $75$~mV; escolhida para não forçar o disparo de células inibidas & artigo \\
duração das fases & não indicada; usar o ajuste padrão do sistema de estimulação e não mudá-lo durante o experimento & escolha \\
\bottomrule
\end{tabular}
\end{center}

\subsection*{Leitura da raquete}

No artigo há duas regiões motoras: os disparos na primeira movem a raquete para cima, na segunda, para baixo; a contribuição das regiões é ponderada para levar em conta a diferente densidade de células e a diferente atividade~\cite{Kagan2022}. A fórmula exata não é dada. Para reproduzir, usamos a regra~\eqref{eq:dishreadout}, $\Delta p=c\,(n_\uparrow-n_\downarrow)$ a cada janela de $10\ms$, com um peso que iguala as regiões pela atividade média em repouso (escolha): $n_\uparrow$ e $n_\downarrow$ são divididos pela contagem média da respectiva região por janela, medida antes do jogo. O coeficiente $c$ é escolhido de modo que uma diferença de uma contagem média desloque a raquete $1\%$ da altura da quadra por janela (escolha); assim, uma vantagem constante de uma região leva a raquete de um lado a outro da quadra em $1$~s.

A posição das regiões no chip não é dada em coordenadas no texto do artigo; há apenas um esquema. Para reproduzir bastam três grupos de eletrodos disjuntos: os oito sensoriais no centro e um grupo de eletrodos de registro de cada lado, um ``sobe'' e outro ``desce'' (escolha).

\subsection*{O jogo}

As dimensões da quadra, o comprimento da raquete e a velocidade da bola não constam do artigo; diz-se apenas que a bola se move com velocidade constante e é refletida pelas paredes. Para reproduzir (escolha): quadra $1\times1$, raquete na parede direita com comprimento $0.2$ (acerto se $|y-p|<0.1$, como no modelo~\texttt{models/dishbrain\_model.py}), a bola atravessa a quadra em $1$~s. Um erro sem nenhuma bola rebatida no rali conta como ``ace''; um rali com mais de três acertos seguidos, como ``longo'' (artigo).

\subsection*{Realimentação}

\begin{center}
\footnotesize
\setlength{\tabcolsep}{4pt}
\renewcommand{\arraystretch}{1.15}
\begin{tabular}{>{\raggedright}p{2.6cm}>{\raggedright}p{8.3cm}>{\raggedright\arraybackslash}p{1.6cm}}
\toprule
Evento & O que é aplicado & Fonte \\
\midrule
acerto & estímulo previsível: os $8$ eletrodos sensoriais ao mesmo tempo, $75$~mV, $100$~disparos/s, $100\ms$; a estimulação normal da bola é interrompida durante esse tempo & artigo \\
erro & estímulo imprevisível: $150$~mV, $5$~disparos/s, $4$~s, em momentos aleatórios e em eletrodos sorteados entre os $8$; depois, $4$~s sem estimulação, e a bola parte numa direção aleatória & artigo \\
lei do acaso & não indicada; para reproduzir, os instantes dos pulsos como fluxo de Poisson com frequência média de $5$~disparos/s & escolha \\
\bottomrule
\end{tabular}
\end{center}

\subsection*{Protocolo e controles}

Uma sessão dura $20$~min; comparam-se os primeiros $5$ e os últimos $15$~min (artigo). Três medidas de resultado: duração média do rali, fração de aces e fração de ralis longos. Condições do experimento (artigo):
\begin{itemize}
\item \emph{estímulo}: a malha completa, como descrito acima;
\item \emph{silêncio}: no lugar do estímulo de acerto ou de erro, o mesmo tempo sem nenhuma estimulação; depois a bola parte numa direção aleatória;
\item \emph{sem realimentação}: depois de um erro a bola simplesmente é refletida e segue em frente; a estimulação da posição da bola continua;
\item \emph{repouso}: o jogo corre e a raquete se move conforme a atividade, mas não há estimulação alguma;
\item \emph{controles}: chip só com meio de cultura; células HEK293T; ``cultura no computador'', uma simulação no lugar das células.
\end{itemize}
Tamanho da amostra na série principal: $9$ culturas de camundongo ($101$ sessões), $11$ culturas humanas ($138$), $6$ chips com meio ($80$), $20$ culturas em repouso ($42$), $3$ simulações ($38$), $399$ sessões ao todo. Na série sobre realimentação: $486$ sessões em $3$ dias, $3$ sessões por dia, condições ``estímulo'', ``silêncio'' e ``sem realimentação'' ($n=27$, $17$ e $15$). O aumento da duração média do rali dos primeiros $5$ minutos para os últimos $15$ só foi obtido nas culturas vivas. Na série sobre realimentação, o aumento significativo ao longo do tempo só ocorreu na condição ``estímulo''; no fim da sessão, a condição ``silêncio'' ainda superava ``sem realimentação'', mas com efeito menor, e não houve aprendizado estável de uma sessão para outra ao longo dos três dias~\cite{Kagan2022}.

\subsection*{O ciclo da bancada passo a passo}

A cada $10\ms$ o computador da bancada faz o seguinte.
\begin{enumerate}
\item Lê os números de disparos $n_\uparrow$, $n_\downarrow$ nas duas regiões motoras na janela anterior e os normaliza pela contagem média da região em repouso.
\item Desloca a raquete de $\Delta p=c\,(n_\uparrow-n_\downarrow)$ e a limita às bordas da quadra.
\item Avança a bola um passo; reflete-a nas paredes superior, inferior e esquerda.
\item Se a bola chegou à parede direita: com $|y-p|<0.1$, acerto, a contagem do rali aumenta em um e aplica-se o estímulo de acerto ($100\ms$); senão, erro, o rali é registrado, aplica-se o estímulo de erro ($4$~s), depois $4$~s de pausa, e a bola parte de novo.
\item Senão, escolhe o eletrodo $k$ pela posição vertical da bola e a frequência $f$ pela distância até a parede da raquete, e aplica ao eletrodo $k$ pulsos bifásicos de $75$~mV com frequência $f$ até a janela seguinte.
\end{enumerate}
Isso basta para montar uma bancada compatível com a publicada e testar a questão principal do capítulo: o que ensina a cultura é a previsibilidade ou a simples diferença entre a resposta ao acerto e ao erro. Para isso, vale acrescentar às três condições do artigo uma quarta, que ele não tem: depois do erro, um estímulo \emph{previsível} com a mesma duração e energia do imprevisível. Se a cultura aprender também assim, o que conta é a diferença, e não a previsibilidade.


\section{Resumo}

\begin{table}[t]
\centering
\footnotesize
\setlength{\tabcolsep}{4pt}
\renewcommand{\arraystretch}{1.15}
\begin{tabular}{>{\raggedright}p{2.1cm}>{\raggedright}p{4.2cm}>{\raggedright}p{1.9cm}>{\raggedright\arraybackslash}p{4.6cm}}
\toprule
Bloco da bancada & Como funciona & Capítulo do livro & O que se sabe \\
\midrule
entrada & lugar (qual dos 8 eletrodos) e frequência (4--40 disparos/s) & \ref{sec:code}, \ref{sec:sensors} & definida pelo experimentador \\
saída & diferença das contagens de duas regiões em $10\ms$~\eqref{eq:dishreadout} & \ref{sec:glutcomputer}, \ref{sec:debugging} & definida pelo experimentador; ruído por~\eqref{eq:rateerror} \\
professor & estímulo previsível após o acerto, imprevisível após o erro; sem dopamina & \ref{sec:dofa} & aumento significativo na sessão só com realimentação completa; com silêncio, efeito menor; instável de uma sessão para outra \\
mecanismo & predição da entrada ou embaralhamento no erro~\eqref{eq:dishdensity} & \ref{sec:dofa}, \ref{sec:recognition} & não estabelecido; ambos explicam o resultado \\
regime da rede & aproximação do regime crítico durante o jogo & \ref{sec:gaba} & a relação com a qualidade do jogo foi medida \\
``senciência'' & afirmação dos autores & --- & não demonstrada; contestada \\
\bottomrule
\end{tabular}
\caption{O DishBrain visto por um engenheiro.}
\label{tab:dishbrain}
\end{table}

O DishBrain é a máquina de neurônios vivos na forma mais simples: a entrada é definida por um código, a saída é lida por contagem, e o professor é a diferença entre a resposta ao sucesso e ao fracasso (tabela~\ref{tab:dishbrain}). Uma rede sem olhos, músculos nem dopamina aprendeu a jogar um tiquinho, e só isso já a transforma numa bancada de testes para as ideias do livro. Seja qual for o mecanismo (predição, embaralhamento ou uma terceira coisa), a resposta será encontrada como recomenda o capítulo~\ref{sec:debugging}: com sonda, sinal de teste e desligamento de partes, numa máquina que finalmente se vê por inteiro.

\section{Exercícios}

\begin{exercise}[\lvlA]
\label{ex:22:1}
\emph{Quanto é preciso deslocar a contagem.} Duas regiões produzem juntas $R_\uparrow+R_\downarrow=2000$ disparos por segundo, fluxos de Poisson. (a)~Qual a dispersão relativa da contagem de uma região em $10\ms$ para $R=1000$ e $R=100$? (b)~Em $1$~s de voo da bola, os deslocamentos~\eqref{eq:dishreadout} se somam. Qual a menor $R_\uparrow-R_\downarrow$ com que o deslocamento médio é o dobro da dispersão?
\end{exercise}

\begin{exercise}[\lvlA]
\label{ex:22:2}
\emph{Quantos ralis são precisos para ver aprendizado.} Os ralis são independentes, a fração rebatida é binomial. Quantos ralis são necessários em cada um de dois períodos para que o aumento da fração rebatida supere dois desvios-padrão da diferença: (a)~de $0.20$ para $0.25$; (b)~de $0.21$ para $0.38$, como no modelo?
\end{exercise}

\begin{exercise}[\lvlB]
\label{ex:22:3}
\emph{Dedução de~\eqref{eq:dishdensity}.} A configuração $\boldsymbol\psi$ percorre um conjunto finito; num erro (probabilidade $P(\boldsymbol\psi)>0$) a rede passa para $\boldsymbol\psi'$ com probabilidade simétrica $K(\boldsymbol\psi,\boldsymbol\psi')$; num acerto, fica. (a)~Prove que $\rho\propto1/P$ é a distribuição estacionária. (b)~Qual a fração de erros com duas configurações de $P=0.9$ e $P=0.1$?
\end{exercise}

\begin{exercise}[\lvlB]
\label{ex:22:4}
\emph{A região absorvente do modelo.} A raquete chega a $p=\min\bigl(1,\max(0,ay+b)\bigr)$, a bola à altura $y\sim U(0,1)$, acerto se $|p-y|<h=0.1$. (a)~Prove que, com $|b|<h$ e $|a-1+b|<h$, a rede rebate todas as bolas, e ache a área dessa região. (b)~Ache numericamente a fração rebatida média com $a\sim U(-1,1)$, $b\sim U(0,1)$.
\end{exercise}

\begin{exercise}[\lvlB]
\label{ex:22:5}
\emph{Simulação: uma cerca em volta das configurações.} Numa cópia de \texttt{dishbrain\_model.py}, dê a $200$ culturas $20\,000$ ralis cada. Qual a fração rebatida nos últimos $500$? E se, depois de cada empurrão, a configuração for cortada pelo retângulo $a\in[-1,2]$, $b\in[-1,1]$? Explique a diferença com o exercício~\ref{ex:22:3}.
\end{exercise}

\begin{exercise}[\lvlB]
\label{ex:22:6}
\emph{Projeto do código de entrada.} A bola é rebatida se o erro de altura for menor que $h=0.1$; a rede lê a entrada em $T=0.25$~s, até a frequência $r$ distinguem-se cerca de $\sqrt{rT}$ níveis, e a estimulação não passa de $40$ pulsos por segundo. (a)~Um código de lugar em oito eletrodos basta? (b)~Dá para transmitir a altura por frequência num só eletrodo?
\end{exercise}

\begin{exercise}[\lvlC]
\label{ex:22:7}
\emph{Pesquisa: predição ou embaralhamento?} Generalize o modelo para $d$ números ajustáveis ($p=\sum_{i<d}a_iy^i$). Como cresce com $d$ o número de ralis até a fração rebatida $0.5$, com embaralhamento e com a regra com o sinal do erro~\eqref{eq:perturbation}? Que experimento na bancada separaria as duas hipóteses?
\end{exercise}

\appendix
\renewcommand{\chaptername}{\appendixname}
\chapter{Substâncias que não definem o tipo do neurônio}
\label{sec:othermediators}

No capítulo~\ref{sec:neurontypes} classificamos os neurônios pelo neurotransmissor principal e obtivemos dez tipos (tabela~\ref{tab:types}). Agora é hora de acrescentar as complicações. O neurônio não libera só o neurotransmissor principal, e no cérebro, além dos neurotransmissores principais, atuam outras substâncias. Elas mudam a forma como a rede transmite e armazena sinais.

Além do barramento de endereços e do de difusão (seção~\ref{sec:buses}), há um terceiro: o \emph{canal reverso}. Algumas substâncias não são liberadas pelo remetente, e sim pelo \emph{destinatário}; elas atravessam a fenda de volta, até a saída do remetente, e mudam quanto neurotransmissor ele liberará da próxima vez. Em redes de comunicação, isso lembra a confirmação de recebimento que o remetente usa para ajustar sua transmissão. É justamente esse canal que os neurônios usam quando mudam os pesos de suas conexões: por ele a saída do remetente fica sabendo da atividade do destinatário, isto é, do segundo fator das regras de aprendizado do capítulo~\ref{sec:dofa}.

A lista completa de neurotransmissores conhecidos é longa: mais de cem substâncias, a maioria cadeias curtas de proteína, os neuropeptídeos. Mas todas cabem em poucas classes, e dentro de cada classe as substâncias se comportam de modo parecido. As substâncias que nunca são o neurotransmissor principal estão reunidas na tabela~\ref{tab:othermediators}, com as mesmas três perguntas que um engenheiro faria: por qual barramento o sinal passa, se age rápido e qual é seu sinal habitual. Elas não definem o tipo do neurônio: são liberadas junto com o neurotransmissor principal, ou são liberadas por células que não são neurônios, ou seguem no sentido inverso.

\begin{table}[t]
\centering
\footnotesize
\setlength{\tabcolsep}{4pt}
\renewcommand{\arraystretch}{1.12}
\begin{tabular}{>{\raggedright}p{2.25cm}>{\raggedright}p{2.8cm}>{\raggedright}p{1.8cm}>{\raggedright}p{1.5cm}>{\centering}p{0.8cm}>{\raggedright\arraybackslash}p{4.3cm}}
\toprule
Classe & Substância & Barramento & Velocidade & Sinal & Papel no cálculo \\
\midrule
Aminoácidos & D-serina & local & lenta & E & segunda chave do receptor de glutamato: condição ``E'' \\
\midrule
Monoaminas & aminas-traço & difusão & lenta & $\pm$ & ajuste fino dos canais monoaminérgicos \\
\midrule
\multirow{2}{2.25cm}{Purinas}
 & ATP & de endereços & rápida e lenta & $+$ & cotransmissor; sinal entre neurônios e glia \\
 & adenosina & por volume & lenta & $-$ & acumula-se com a atividade: contador de carga~\cite{Dunwiddie2001} \\
\midrule
\multirow{8}{2.25cm}{Neuropeptídeos (mais de cem)}
 & encefalinas, endorfinas, dinorfina & por volume & lenta & $-$ & supressão da transmissão \\
 & substância P & por volume & lenta & $+$ & reforço lento \\
 & neuropeptídeo Y & por volume & lenta & $-$ & inibição lenta \\
 & somatostatina & por volume & lenta & $-$ & inibição lenta, acompanha o GABA \\
 & peptídeo intestinal vasoativo (VIP) & por volume & lenta & $+$ & desinibição: inibição dos neurônios inibitórios \\
 & colecistocinina & por volume & lenta & $\pm$ & acompanha o GABA em parte dos neurônios inibitórios \\
 & orexina & difusão & lenta & $+$ & estabilizador do modo de vigília \\
 & ocitocina, vasopressina & difusão & lenta & $\pm$ & ajustes globais do comportamento \\
\midrule
\multirow{2}{2.25cm}{Gases}
 & óxido nítrico NO & reverso, por volume & lenta & $\pm$ & atravessa membranas; sinal reverso na mudança de pesos~\cite{Garthwaite2008} \\
 & CO, H$_2$S & por volume & lenta & $\pm$ & papel pouco estudado \\
\midrule
Lipídios & endocanabinoides (anandamida, 2-AG) & reverso & lenta & $-$ & o destinatário reduz a liberação do remetente: realimentação sobre o peso~\cite{WilsonNicoll2001} \\
\bottomrule
\end{tabular}
\caption{Substâncias que não definem o tipo do neurônio. \emph{Barramento}, \emph{velocidade} e \emph{sinal} como na tabela~\ref{tab:types}; além disso, barramento por volume: o neurotransmissor se espalha no espaço extracelular em torno do local de liberação; reverso: vai do destinatário de volta ao remetente; local: age perto do local de liberação. ``E'' é um sinal habilitador: sozinho não gera corrente, mas sem ele outra entrada não dispara, como a segunda entrada de uma porta lógica ``E''. Os neuropeptídeos quase sempre são liberados junto com o neurotransmissor principal e sobretudo em alta frequência de disparos~\cite{vandenPol2012}.}
\label{tab:othermediators}
\end{table}

\chapter{Notação}
\label{sec:notation}

Os tipos de neurônio são designados por um código de quatro letras do neurotransmissor principal: \nt{GLUT}, glutamato; \nt{GABA}, ácido gama-aminobutírico; \nt{GLYC}, glicina; \nt{ACET}, acetilcolina; \nt{DOPA}, dopamina; \nt{SERO}, serotonina; \nt{HIST}, histamina; \nt{NORA}, noradrenalina; \nt{ADRE}, adrenalina; \nt{ASPA}, aspartato (tabela~\ref{tab:types}).

Há menos letras que grandezas, e em capítulos diferentes a mesma letra às vezes significa coisas diferentes. Na tabela abaixo, cada significado tem sua linha; na coluna da direita está a fórmula em que ele é introduzido.

\begin{center}
\footnotesize
\setlength{\tabcolsep}{4pt}
\renewcommand{\arraystretch}{1.12}
\begin{longtable}{>{\raggedright}p{2.0cm}>{\raggedright}p{9.6cm}>{\raggedright\arraybackslash}p{2.4cm}}
\toprule
Símbolo & Significado & Onde é introduzido \\
\midrule
\endfirsthead
\toprule
Símbolo & Significado & Onde é introduzido \\
\midrule
\endhead
\bottomrule
\endfoot
$a$ & inclinação da característica de transferência linear, $f(u)=au$ & \eqref{eq:feedback} \\
$a$ & nível de \nt{ACET}: $0$, leitura; $1$, escrita & \eqref{eq:acetnet} \\
$a_i$, $a$ & fadiga do grupo: no gerador de ritmo; na gangorra do sono de ondas lentas & \eqref{eq:halfcentre}, \eqref{eq:updown} \\
$a_S$, $a_P$ & sensibilidade do ritmo cardíaco ao acelerador e ao freio & \eqref{eq:gasbrake} \\
$A(r)$, $A_0$ & resposta da sinapse a um disparo na frequência $r$; numa sinapse descansada & \eqref{eq:depression} \\
$A_1$ & resposta do destinatário a uma porção de neurotransmissor & \eqref{eq:quantal} \\
$A$ & área da membrana do neurônio junto com os dendritos & \eqref{eq:dendriteload} \\
$A_f$, $A_s$ & fração da correção mantida até o próximo movimento, nos processos rápido e lento & \eqref{eq:tworate} \\
$b_i$ & aporte próprio do marca-passo & \eqref{eq:seroneuron} \\
$b$ & rapidez com que o trabalho cansa o grupo no gerador de ritmo ou na gangorra do sono & \eqref{eq:halfcentre}, \eqref{eq:updown} \\
$b$ & número de bits por amostra da digitalização & \eqref{eq:probedata} \\
$B_f$, $B_s$ & intensidade com que a correção aprende com o erro, nos processos rápido e lento & \eqref{eq:tworate} \\
$c$ & concentração do neurotransmissor no volume & \eqref{eq:seroneuron} \\
$c$ & coeficiente de correlação dos ruídos de neurônios vizinhos & \eqref{eq:correlated} \\
$c$ & coeficiente com que a diferença das contagens move a raquete & \eqref{eq:dishreadout} \\
$c_f$ & resposta do neurônio $C$: o maior dos $s_{f,p}$ entre todas as posições & \eqref{eq:scstage} \\
$c_m$ & capacitância específica da membrana & \eqref{eq:dendriteload} \\
$C$ & custo do esforço: uma ação no tempo $\tau_a$ custa $C/\tau_a$ & \eqref{eq:vigor} \\
$d'$ & discriminabilidade de duas entradas: a diferença dividida pelo ruído & \eqref{eq:gainsnr} \\
$d$, $d_j$ & atraso no caminho do sinal & \eqref{eq:glutneuron}, \eqref{eq:vstar} \\
$d$ & atraso da malha de realimentação com o músculo & \eqref{eq:delaystable} \\
$D$ & coeficiente de difusão & \eqref{eq:diffusionlength} \\
$D$ & denominador nas frequências da rede \nt{GLUT}--\nt{GABA} & \eqref{eq:isn} \\
$D$ & tempo de espera por uma recompensa adiada & \eqref{eq:patience} \\
$e$, $e_{ij}$ & traço na sinapse & \eqref{eq:threefactor} \\
$e$, $e_i$ & erro da saída ou do movimento, que chega pelo fio de erro & \eqref{eq:lms}, \eqref{eq:adaptivefilter}, \eqref{eq:tworate} \\
$E_{\text{disp}}$ & energia de um disparo, incluindo o trabalho das sinapses & \eqref{eq:energy} \\
$E$ & frequência do grupo \nt{GLUT} & \eqref{eq:ei} \\
$E_E$ & potencial de equilíbrio da sinapse excitatória & \eqref{eq:shunt} \\
$f$, $F$ & característica de transferência: frequência em função da entrada & \eqref{eq:ratenet}, \eqref{eq:gain} \\
$f$, $f_0$ & frequência dos batimentos cardíacos; ritmo próprio do marca-passo & \eqref{eq:gasbrake} \\
$f_s$ & frequência de digitalização de um eletrodo & \eqref{eq:probedata} \\
$F(t)$, $F_1$ & força do músculo; força de um abalo & \eqref{eq:forcesum} \\
$F_1$, $F_2$ & forças de dois músculos que puxam a articulação em sentidos opostos & \eqref{eq:antagonists} \\
$g$ & sensibilidade do receptor à concentração & \eqref{eq:seroneuron} \\
$g_L$, $g_E$, $g_I$ & condutâncias de vazamento, excitação e inibição & \eqref{eq:shunt} \\
$g$ & ganho definido pelo \nt{NORA} & \eqref{eq:gain}, \eqref{eq:machine} \\
$g$ & inibição total do \nt{GABA} na rede de memória & \eqref{eq:acetnet} \\
$g$ & ganho do portão da dor, definido de cima & \eqref{eq:gate} \\
$g$ & ganho de um estágio da cascata hormonal & \eqref{eq:cascade} \\
$g_j$ & filtro $j$ com sua constante de tempo na entrada do filtro adaptativo & \eqref{eq:adaptivefilter} \\
$G$ & ganho da malha de realimentação & \eqref{eq:feedback} \\
$G$ & velocidade de correção na malha com atraso & \eqref{eq:delaystable} \\
$h$, $h_I$ & entrada externa do grupo \nt{GLUT} e do grupo \nt{GABA} & \eqref{eq:ei}, \eqref{eq:isn} \\
$h(t)$ & abalo isolado do músculo & \eqref{eq:twitch} \\
$H(t)$ & limiar superior da pressão de sono: ao atingi-lo, a máquina adormece & \eqref{eq:sleeppressure} \\
$I$, $I_i$ & aporte externo ao neurônio ou ao grupo & \eqref{eq:ratenet}, \eqref{eq:wta}, \eqref{eq:updown} \\
$I$ & frequência do grupo \nt{GABA} & \eqref{eq:ei} \\
$I$ & brilho ou intensidade do estímulo & \eqref{eq:photonnoise}, \eqref{eq:nakarushton} \\
$I$ & corrente de entrada que carrega a membrana & \eqref{eq:dendriteload} \\
$k$ & número de disparos na rajada & \eqref{eq:burst} \\
$k$ & quantas vezes o acréscimo deve superar o ruído & \eqref{eq:photonnoise} \\
$k$ & coeficiente de realimentação da cascata hormonal & \eqref{eq:cascade} \\
$k_s$ & intensidade com que a dor alimenta a sensibilização & \eqref{eq:sensitization} \\
$K$ & número de atributos do resultado & \eqref{eq:vectortd} \\
$K$ & rigidez da articulação & \eqref{eq:antagonists} \\
$K$ & ganho da malha da cascata hormonal, $K=g^3k$ & \eqref{eq:cascade}, \eqref{eq:cascadestable} \\
$K_\varphi$ & intensidade do ajuste do gerador circadiano ao sinal externo & \eqref{eq:pll} \\
$L$ & comprimento da linha de atraso & \eqref{eq:itd} \\
$L$ & comprimento do fio a partir do sensor de dor & \eqref{eq:paintime} \\
$L(t)$ & limiar inferior da pressão de sono: ao atingi-lo, a máquina acorda & \eqref{eq:sleeppressure} \\
$M$ & medida da relação causal entre o sinal antecipador e a recompensa & \eqref{eq:retro} \\
$M_k$ & expectativa do atributo $k$ & \eqref{eq:vectortd} \\
$M_{ik}$ & intensidade da inibição vinda do neurônio \nt{GABA} $k$ & \eqref{eq:machine} \\
$M$ & torque da articulação & \eqref{eq:antagonists} \\
$M$ & número de linhas de entrada dos misturadores & \eqref{eq:cover} \\
$n$, $\bar n$ & número de disparos na janela; sua média & \eqref{eq:rateerror} \\
$n$ & número de entradas do elemento de limiar & \eqref{eq:cover} \\
$n_\uparrow$, $n_\downarrow$ & contagem de disparos das regiões ``sobe'' e ``desce'' na janela & \eqref{eq:dishreadout} \\
$N$ & número de neurônios & \eqref{eq:correlated}, \eqref{eq:energy} \\
$N$ & número de eletrodos da sonda & \eqref{eq:probedata} \\
$N_s$ & número de locais de liberação entre dois neurônios & \eqref{eq:quantal} \\
$p$ & probabilidade de liberação de neurotransmissor por disparo & \eqref{eq:burst} \\
$p$ & fração de neurônios ativos na rede de memória & \eqref{eq:acetnet} \\
$p$ & perturbação que se aprende a compensar & \eqref{eq:tworate} \\
$P$ & potência gasta com sinais & \eqref{eq:energy}, \eqref{eq:dishpower} \\
$P$ & incerteza da estimativa armazenada & \eqref{eq:kalman} \\
$P$ & atividade do ``freio'': fios \nt{ACET} para o coração & \eqref{eq:gasbrake} \\
$P_{\max}$ & quantos padrões aleatórios o elemento de limiar consegue separar em duas classes & \eqref{eq:cover} \\
$P_{\text{erro}}$ & probabilidade de erro & \eqref{eq:dishdensity} \\
$q$ & velocidade com que o mundo muda & \eqref{eq:kalman} \\
$q_k$ & frequência do neurônio \nt{GABA} $k$ & \eqref{eq:machine} \\
$q_0$ & atividade de fundo do intermediário inibitório & \eqref{eq:gate} \\
$q$ & frequência do intermediário inibitório no portão da dor & \eqref{eq:gate} \\
$q$ & quantas vezes a próxima unidade motora é mais forte que a anterior & \eqref{eq:sizeprinciple} \\
$r$, $r_i$ & frequência de disparos do neurônio & \eqref{eq:ratenet} \\
$r$, $r_t$ & recompensa (fluxo de recompensa) & \eqref{eq:rpe}, \eqref{eq:ctd} \\
$\mathbf r(t)$ & vetor de respostas do reservatório & \eqref{eq:reservoir} \\
$R$, $\bar R$ & recompensa obtida; sua expectativa ou média por unidade de tempo & \eqref{eq:perturbation}, \eqref{eq:vigor} \\
$R$ & variância do ruído na entrada do filtro & \eqref{eq:kalman} \\
$R$, $r_0$ & recompensas adiada e imediata & \eqref{eq:patience} \\
$R_{\max}$ & taxa máxima de transmissão, bit/s & \eqref{eq:spikeentropy} \\
$r_{\max}$ & maior frequência de disparos & \eqref{eq:gain}, \eqref{eq:nakarushton} \\
$R_{\text{bruto}}$, $R_{\text{disp}}$ & fluxo de dados da sonda: bruto e só os disparos, bit/s & \eqref{eq:probedata} \\
$s_j$ & sinal das saídas do neurônio $j$ & \eqref{eq:dalesign} \\
$s_i$ & sinal do receptor do destinatário & \eqref{eq:seroneuron} \\
$s$, $\hat s$ & estado do mundo; sua estimativa & \eqref{eq:rpe}, \eqref{eq:kalman} \\
$s_{f,p}$ & resposta do neurônio $S$ ao atributo $f$ na posição $p$ & \eqref{eq:scstage} \\
$S$ & atividade do ``acelerador'': fios \nt{NORA} para o coração & \eqref{eq:gasbrake} \\
$S$ & pressão de sono & \eqref{eq:sleeppressure} \\
$t_1$ & instante do primeiro disparo & \eqref{eq:latency} \\
$t_k$ & instante do $k$-ésimo disparo & \eqref{eq:forcesum} \\
$t_{f,i}$ & modelo do atributo $f$ & \eqref{eq:scstage} \\
$T$ & janela de contagem de disparos; tempo de acumulação de fótons & \eqref{eq:rateerror}, \eqref{eq:photonnoise} \\
$T_c$ & tempo de subida do abalo muscular & \eqref{eq:twitch} \\
$T_{\text{salva}}$ & intervalo entre salvas da cultura & \eqref{eq:burstinterval} \\
$u$ & entrada total do neurônio & \eqref{eq:gain} \\
$u$, $u(t)$ & comando do cérebro: na entrada do filtro adaptativo; para a cascata hormonal & \eqref{eq:adaptivefilter}, \eqref{eq:cascade} \\
$u(t)$ & entrada do reservatório & \eqref{eq:reservoir} \\
$U$ & nível da entrada constante & \eqref{eq:latency} \\
$U$ & fração do estoque de neurotransmissor liberada por disparo & \eqref{eq:depression} \\
$v$ & velocidade do sinal no fio; velocidade de movimento da mancha & \eqref{eq:itd}, \eqref{eq:vstar} \\
$V$ & tensão na membrana & \eqref{eq:glutneuron}, \eqref{eq:shunt} \\
$V$ & expectativa da recompensa futura & \eqref{eq:rpe}, \eqref{eq:ctd} \\
$w$, $w_{ij}$, $W_{ij}$ & peso da conexão & \eqref{eq:glutneuron}, \eqref{eq:ratenet}, \eqref{eq:acetnet}, \eqref{eq:lms}, \eqref{eq:adaptivefilter} \\
$\mathbf w_k$ & pesos das entradas do neurônio $C$ $k$ & \eqref{eq:traceinvariance} \\
$w_{q\mu}$, $w_{q\nu}$, $w_{z\mu}$ & pesos das conexões no portão da dor & \eqref{eq:gate} \\
$W_{\text{saída}}$ & pesos da leitura linear do reservatório & \eqref{eq:reservoir} \\
$x$, $y$ & entradas e saída lógicas & \eqref{eq:dualrail}, \eqref{eq:veto} \\
$x$, $y$ & atividade do remetente e do destinatário & \eqref{eq:hebb} \\
$x$ & posição do detector na linha de atraso & \eqref{eq:itd} \\
$x_i$ & entrada vinda dos órgãos dos sentidos & \eqref{eq:acetnet} \\
$x_p$ & entrada na posição $p$ & \eqref{eq:scstage} \\
$\mathbf x$ & saídas dos neurônios $S$, entrada do neurônio $C$ & \eqref{eq:traceinvariance} \\
$x_j$ & entrada $j$ do filtro adaptativo: o comando depois do filtro $g_j$ & \eqref{eq:lms}, \eqref{eq:adaptivefilter} \\
$x_f$, $x_s$ & correções dos processos rápido e lento & \eqref{eq:tworate} \\
$x_1$, $x_2$, $x_3$ & níveis nos três estágios da cascata hormonal; $x_3$ é o hormônio final & \eqref{eq:cascade} \\
$x^*$ & fração recuperada do estoque de neurotransmissor com a qual começa uma nova avalanche & \eqref{eq:burstinterval} \\
$y$ & entrada ruidosa do filtro & \eqref{eq:kalman} \\
$y$, $\hat y$ & sinal dos sensores; sua predição pelo filtro adaptativo & \eqref{eq:adaptivefilter} \\
$y_k$, $\bar y_k$ & atividade do neurônio $C$ $k$; seu traço & \eqref{eq:traceinvariance} \\
$y(t)$ & saída da leitura do reservatório & \eqref{eq:reservoir} \\
$z$ & frequência da saída do portão da dor & \eqref{eq:gate} \\
\midrule
$\alpha_i^\pm$ & peso dos erros positivos e negativos & \eqref{eq:asymmetric} \\
$\beta$ & intensidade da inibição mútua & \eqref{eq:wta} \\
$\beta$ & intensidade da inibição da saída do portão da dor & \eqref{eq:gate} \\
$\gamma$ & fator de desconto & \eqref{eq:rpe} \\
$\delta(\cdot)$ & função delta: salto instantâneo & \eqref{eq:glutneuron} \\
$\delta$, $\delta_t$ & erro de predição de recompensa & \eqref{eq:rpe}, \eqref{eq:ctd} \\
$\delta t$ & diferença entre os tempos de chegada do som aos dois ouvidos & \eqref{eq:itd} \\
$\Delta$, $\Delta_{\max}$ & intervalo entre disparos; janela de coincidência & \eqref{eq:window} \\
$\Delta t$ & precisão com que o destinatário distingue o tempo & \eqref{eq:spikeentropy} \\
$\Delta F$ & acréscimo de força da próxima unidade motora & \eqref{eq:sizeprinciple} \\
$\Delta I$ & acréscimo de brilho distinguível & \eqref{eq:photonnoise} \\
$\Delta p$ & deslocamento da raquete na janela & \eqref{eq:dishreadout} \\
$\Delta u$ & diferença entre as entradas de dois grupos & \eqref{eq:gainsnr} \\
$\Delta V$ & quanto é preciso elevar a tensão da membrana & \eqref{eq:dendriteload} \\
$\Delta x$ & distância entre sensores vizinhos & \eqref{eq:vstar} \\
$\varepsilon$ & diferença relativa entre as entradas & \eqref{eq:wtagain} \\
$\eta$ & taxa de aprendizado & \eqref{eq:plasticity}, \eqref{eq:lms}, \eqref{eq:traceinvariance} \\
$\eta$ & ruído na rede depois do amplificador & \eqref{eq:softmax} \\
$\mu$ & entrada de tato para o portão da dor & \eqref{eq:gate} \\
$\nu$ & entrada de dor & \eqref{eq:gate} \\
$\theta$ & limiar de disparo & \eqref{eq:window}, \eqref{eq:scstage} \\
$\kappa$ & quantidade de neurotransmissor por disparo & \eqref{eq:seroneuron} \\
$\kappa_i$ & razão entre os pesos dos erros positivos e negativos & \eqref{eq:expectile} \\
$\kappa_m$ & relação entre a força do músculo e sua rigidez & \eqref{eq:antagonists} \\
$\ell$ & distância que o neurotransmissor percorre ao se espalhar & \eqref{eq:diffusionlength} \\
$\ell$ & braço de alavanca com que o músculo puxa a articulação & \eqref{eq:antagonists} \\
$\xi$ & tremor aleatório da saída do neurônio & \eqref{eq:perturbation} \\
$\rho(\boldsymbol\psi)$ & fração do tempo passado na configuração $\boldsymbol\psi$ & \eqref{eq:dishdensity} \\
$\sigma$ & dispersão, escala do ruído & \eqref{eq:rateerror}, \eqref{eq:softmax} \\
$\sigma$ & brilho com o qual a resposta é metade da máxima & \eqref{eq:nakarushton} \\
$\sigma_{\text{antes}}$, $\sigma_{\text{depois}}$ & ruído antes e depois do amplificador & \eqref{eq:gainsnr} \\
$\tau$ & constante de tempo da membrana & \eqref{eq:glutneuron} \\
$\tau$ & constante de tempo de um estágio da cascata hormonal & \eqref{eq:cascade} \\
$\tau_{\text{ef}}$ & constante de tempo de um grupo que excita a si mesmo & \eqref{eq:perfectint} \\
$\tau_c$ & tempo de vida do neurotransmissor no volume & \eqref{eq:seroneuron} \\
$\tau_E$, $\tau_I$ & constantes de tempo dos grupos \nt{GLUT} e \nt{GABA} & \eqref{eq:ei} \\
$\tau_e$ & tempo de vida do traço & \eqref{eq:threefactor} \\
$\tau_{\text{tr}}$ & tempo de vida do traço de atividade do neurônio $C$ & \eqref{eq:traceinvariance} \\
$\tau_s$ & tempo de retorno da sensibilidade após uma lesão & \eqref{eq:sensitization} \\
$\tau_r$, $\tau_d$ & tempo de acúmulo da pressão de sono na vigília e de sua descarga no sono & \eqref{eq:sleeppressure} \\
$\tau_{\text{rec}}$ & tempo de recuperação do estoque de neurotransmissor & \eqref{eq:depression} \\
$\tau_\gamma$ & horizonte de planejamento & \eqref{eq:ctd} \\
$\tau_v$ & rapidez com que a expectativa é refinada & \eqref{eq:ramp} \\
$\tau_a$ & tempo de fadiga no gerador de ritmo ou na gangorra do sono & \eqref{eq:halfcentre}, \eqref{eq:updown} \\
$\tau_a^*$ & melhor tempo de execução da ação & \eqref{eq:vigor} \\
$\Phi$ & lado direito da regra de aprendizado & \eqref{eq:plasticity} \\
$\Phi$ & fluxo de fótons & \eqref{eq:photonnoise} \\
$\varphi_k$ & atributo $k$ do resultado obtido & \eqref{eq:vectortd} \\
$\varphi$ & diferença de fase entre o gerador circadiano e o sinal externo & \eqref{eq:pll} \\
$\boldsymbol\psi$ & configuração da rede no modelo DishBrain & \eqref{eq:dishdensity} \\
$\omega$, $\omega_0$ & frequência do sinal externo (luz); frequência própria do gerador circadiano & \eqref{eq:pll} \\
\end{longtable}
\end{center}

\chapter{Respostas}
\label{sec:answers}

Respostas curtas aos exercícios do fim dos capítulos. As soluções completas estão num material à parte, para professores.

\answer{ex:01:1}{(a)~$8\cdot10^{10}$ \nt{GLUT}: $10$ populações da Terra; os de difusão, $\approx1.9\cdot10^6$: uma cidade do tamanho de Viena. (b)~$\approx8\cdot10^{14}$ terminais, dos quais $\approx3\cdot10^{11}$ de difusão, fração $\approx4\cdot10^{-4}$: $\approx16$ vezes maior que a fração desses neurônios ($2.2\cdot10^{-5}$).}
\answer{ex:01:2}{(a)~$\tau_c\ln(c_0/c_*)\approx4.6\ms$. (b)~A linha fica livre se $T\ge\tau_c\ln101$: até $\approx22$ liberações por segundo, em vez de $\approx220$.}
\answer{ex:01:3}{$V(s_k)=\gamma^{K-1-k}r$; $\delta_{\text{lâmpada}}=\gamma^Kr=re^{-T/\tau_\gamma}\approx0.60\,r$ para qualquer $\Delta$, $\tau_\gamma\approx1.95\,\text{s}$.}
\answer{ex:01:4}{$n^*=93$, $47$, $25$, $12$: $n^*\approx5/\alpha$ para $\alpha$ pequeno.}
\answer{ex:01:5}{\nt{GLUT}: $1\to10\to10^4\to10^7$, $\approx10^7$ neurônios, $4$ elos, $4\ms$. \nt{DOPA}: $1\to10\to3.2\cdot10^5$, $\approx3\cdot10^5$ neurônios, $3$ elos, $30$ vezes menos; a mesma ordem que o número de neurônios \nt{DOPA}.}
\answer{ex:01:6}{$\lambda=f_EJ_E-f_IJ_I$, donde $J_I=37.5$; razão das correntes $f_IJ_I/f_EJ_E=0.94$; avalanche ($\lambda\ge1$) quando a inibição é enfraquecida em apenas $6.7\,\%$.}
\answer{ex:01:7}{$n\ge57$ apresentações para cada atraso em cada condição. A mesma queda seria produzida, por exemplo, pela imprecisão da estimativa interna do tempo, que cresce com $T$.}

\answer{ex:02:1}{Fileira de $3.2$~mm com passo de $25\um$: $\approx129$ detectores. Dispara uma faixa de largura $v\,\tau\ln2=1.7$~mm, $\approx69$ detectores, mais da metade da fileira; a direção é o centro da faixa.}
\answer{ex:02:2}{Janela $-\tau\ln\frac{w_2}{\theta-w_1}\le\delta\le\tau\ln\frac{w_1}{\theta-w_2}$, isto é, $[-0.235,\,0.178]\ms$; centro $c=\frac\tau2\ln\frac{w_1(\theta-w_1)}{w_2(\theta-w_2)}=-28$~\textmu s. A direção se desloca $28$~\textmu s para o lado do ouvido mais forte: quase três passos de resolução.}
\answer{ex:02:3}{São necessários $s_i$ com $s_iw_{ij}s_j\ge0$; eles existem se e somente se os ciclos são pares. A inibição mútua se reduz (não há oscilações estáveis); o laço \nt{GLUT}--\nt{GABA}, não (pode oscilar).}
\answer{ex:02:4}{Seis neurônios $g_k=[x+y+z\ge k]$ e $\bar g_k=[\bar x+\bar y+\bar z\ge4-k]$, $k=1,2,3$; $C=g_2$, $\bar C=\bar g_2$, $S=[g_1+\bar g_2+g_3\ge2]$, $\bar S=[\bar g_1+g_2+\bar g_3\ge2]$: oito neurônios. Com E e OU, $12$.}
\answer{ex:02:5}{$T_{\min}(0.6)=0.718\,\tau$; $I_c=0.225$; perto do limiar, $T_{\min}\propto(I-I_c)^{-1/2}$.}
\answer{ex:02:6}{(a)~$500$ elos, $5\cdot10^4$ neurônios; dispersão de $4.5\ms$ ($0.45\,\%$), erro relativo $\propto T^{-1/2}$. (b)~Um fator aleatório de atraso comum a todos os elos, com dispersão de $5\,\%$.}
\answer{ex:02:7}{Renovar a cada $\tau_F\ln2=0.69\,\text{s}$: $4.3$ disparos por segundo por neurônio contra $20$, $4.6$ vezes mais econômico ($4\cdot10^{-7}$ contra $2\cdot10^{-6}$~J com $10^{-10}$~J por disparo). O flip-flop perde o bit; o capacitor o mantém se o silenciamento começar no máximo $0.49\,\text{s}$ após a renovação (probabilidade $0.71$).}

\answer{ex:03:1}{$35\to17.5\mV$ com $g_E=g_L$; $56\to40\mV$ com $g_E=4g_L$; a subtração daria $38.5\mV$: a derivação divide $g_E$ por $3$. A janela se estreita pela metade ($\tau_{\text{ef}}$ cai $2$ vezes).}
\answer{ex:03:2}{$\operatorname{tr}=\frac{w_{EE}-1}{\tau_E}-\frac1{\tau_I}$, $\det=\frac{1-w_{EE}+w_{EI}w_{IE}}{\tau_E\tau_I}$; na fronteira $\omega=\sqrt{\det}$, $\tau_I=20\ms$, período de $73\ms$.}
\answer{ex:03:3}{$d_c=\tau_E\arccos(-1/G)/\sqrt{G^2-1}$, período $2\pi\tau_E/\sqrt{G^2-1}\in(2d_c,4d_c)$. $G=2$: $d_c=12.1\ms$, período de $36\ms$; $G=5$: $d_c=3.6\ms$, período de $12.8\ms$: exatamente a faixa dos ritmos rápidos.}
\answer{ex:03:4}{(a)~Em $I_2=\beta I_1=2$ subindo e $I_2=I_1/\beta=0.5$ descendo, um laço de largura $1.5$. (b)~Em $\tau\ln10/(\beta-1)=2.3\tau$ a cada ordem decimal de $\varepsilon$.}
\answer{ex:03:5}{Três intermediários, $d_k=5$, $10$, $15\ms$: os vetos devem cobrir $15\ms$, $5\ms$ cada um.}
\answer{ex:03:6}{$n+M+1=3+4+1=8$ neurônios. Dois: \nt{GABA} $=[x+y+z\ge2]$, saída $=[x+y+z-2\,\nt{GABA}\ge1]$.}
\answer{ex:03:7}{$\binom84=70<100\le\binom94=126$: $9$ intermediários (teorema de Sperner). Sem conexões diretas, $100$: o posto de $\beta(J-I)$ é $n$.}
\answer{ex:03:8}{$h_c=\frac1{2w_{EE}}+\frac{w_{EI}w_{IE}-w_{EE}}{4w_{EE}^2}=\frac7{18}\approx0.39$; a simulação dá a troca de sinal entre $h=0.38$ e $0.40$.}

\answer{ex:04:1}{(a)~$\approx1.1\cdot10^{11}$ bits/J. (b)~$\log_2\sqrt{10}\approx1.7$ bit contra $81$: $\approx50$ vezes.}
\answer{ex:04:2}{$r^*=(P_0/E_{\text{disp}})\ln\bigl(1/(r^*\Delta t)\bigr)$, $P_0/E_{\text{disp}}=1.16\,\text{s}^{-1}$: $r^*\approx6$ disparos por segundo, $7.4\cdot10^{10}$ bits/J.}
\answer{ex:04:3}{$\mathcal I=\frac{N}{1+(N-1)c}=9.9$ (limite $10$) contra $\mathcal I=\frac{N}{1-c}=1111$: $\approx110$ vezes; a correlação só é nociva quando se parece com o sinal.}
\answer{ex:04:4}{$\theta\approx68.7$ e $19.9$ (em unidades de $w$), $m=19$ e $m=10$: o destinatário rápido precisa de quase metade das entradas coincidentes, embora sua entrada média seja cinco vezes menor.}
\answer{ex:04:5}{$U\tau_{\text{rec}}\ge0.725\,\text{s}$ e $\tau_{\text{rec}}(1-2U)\le0.05\,\text{s}$, o que exige $U\ge0.483$; o menor $\tau_{\text{rec}}=0.725\,\text{s}$ com $U=1$ ($1.45\,\text{s}$ com $U=0.5$).}
\answer{ex:04:6}{(a)~$\theta\,e/(e-1)\approx1.58\,\theta$ para quaisquer $\tau$ e $\theta$. (b)~$14$ disparos.}
\answer{ex:04:7}{(a)~$1.62$ bit por palavra: $0.77$ do número de disparos, $0.84$ de sua disposição. (b)~$1.13$ e $1.59$ bit: com $30$ palavras, a estimativa perde quase um terço.}

\answer{ex:05:1}{$\ell\approx17\um$; $\approx100$~dias: a difusão pelo cérebro é feita pela ramificação do fio; a difusão molecular só espalha cada local de liberação por $\sim20\um$.}
\answer{ex:05:2}{$r=ab/(1+G)$, $S=1/(1+G)$ exatamente; $+1.7\,\%$ em vez de $+20\,\%$.}
\answer{ex:05:3}{$T^*=1.21\,\text{s}$ para $G=2$ e $0.19\,\text{s}$ para $G=9$: na prática $G\sim2$--$4$, redução de $3$--$5$ vezes.}
\answer{ex:05:4}{$\mathrm{CV}=1/\sqrt{2\lambda\tau_c}\approx9\cdot10^{-4}$ com taxa total $\lambda$; um único neurônio precisa de $\tau_c=12.5\,\text{s}$.}
\answer{ex:05:5}{$c\approx0.40$ nos dois casos; $\mathrm{CV}=0.050$ contra $0.158$, $\sqrt{10}$ vezes menor. As taxas dos neurônios individuais continuam proporcionais a $a_i$: só a soma é regulada.}
\answer{ex:05:6}{$N_1=\overline{A\vee B}$, $N_2=\overline{A\vee N_1}$, $N_3=\overline{B\vee N_1}$, $N_4=\overline{N_2\vee N_3}$, XOR $=N_5=\overline{N_4}$; cada comutação leva $\tau_c\ln2$, um caminho de $4$ portas leva $2.8\,\text{s}$ por XOR.}
\answer{ex:05:7}{(a)~$\tau_\gamma=6/\ln3\approx5.5\,\text{s}$. (b)~$\bar D<4\,\text{s}$ contra $D<2.2\,\text{s}$: a imprevisibilidade do atraso torna mais paciente.}
\answer{ex:05:8}{$k_2=1/\tau_d=10\,\text{s}^{-1}$, $k_1=1/\tau_d-1/\tau_\gamma=9.5\,\text{s}^{-1}$; $\tau_\gamma=1/(k_2-mk_1)$: dobra com $+2.6\,\%$ (e com $+5.3\,\%$ o horizonte fica infinito).}

\answer{ex:06:1}{$e=(1-e^{-T_a/\tau_e})e^{-d/\tau_e}$. (a)~$0.110$ contra $4.5\cdot10^{-5}$, $\approx2.5\cdot10^3$ vezes. (b)~$\tau_e^*=T_a/\ln(1+T_a/d)=0.59\,\text{s}$.}
\answer{ex:06:2}{Taxa $\nu_x\nu_y(A_+\tau_+-A_-\tau_-)=0.8A_+$ por segundo, $\approx2900A_+$ por hora; $\tau_-=\tau_+/0.6=33\ms$.}
\answer{ex:06:3}{$\mathbf e_1$ se $\mu^2+s_1^2>s_2^2$; aqui $1.01>0.25$, e o neurônio aprende $\mathbf e_1$, embora a dispersão ao longo de $\mathbf e_2$ seja $25$ vezes maior: a regra encontra o vetor principal da matriz de correlação, não da de covariância.}
\answer{ex:06:4}{$V=Re^{-(t_c+L-t)/\tau_\gamma}$ entre a luz e o suco, de modo que $\dot V=V/\tau_\gamma$; pico de área $Re^{-L/\tau_\gamma}$: $0.61R$ e $0.78R$.}
\answer{ex:06:5}{Relação sinal-ruído $1/(N+1)$, tentativas $\propto N+1$: $5\cdot10^5$ tentativas $\approx1$~mês e $5\cdot10^{11}$ tentativas $\approx8\cdot10^4$~anos.}
\answer{ex:06:6}{(a)~$k_2=1/\tau_d$, $k_1=1/\tau_d-1/\tau_\gamma$. (b)~Erro espúrio $-(V/\tau_\gamma)\,\varepsilon/(1+\varepsilon)$, $\varepsilon=\tau_d/\tau_\gamma$, donde $\tau_d\le0.105\,\text{s}$.}
\answer{ex:06:7}{$0.46$, $0.90$, $0.99$, $0.95$ e $0.70$ com $d=3\,\text{s}$. Os pontos caem numa única curva da fração do traço $F=(1-e^{-T_{\text{cue}}/\tau_e})e^{-d/\tau_e}$: aprendizado com $F\gtrsim0.1$, escolha aleatória com $F<10^{-2}$.}
\answer{ex:06:8}{$\eta^*=1/\bigl(2\sigma^2(N+2)\bigr)$, em $N+2$ tentativas; com $m$ flutuando de $N$, em $N(m+2)/m\ge N+2$: a esparsidade não ajuda.}

\answer{ex:07:1}{$V=\kappa p/\bigl(\kappa p+(1-\kappa)(1-p)\bigr)$: $0.059$, $0.20$, $0.50$; a mediana é $0$, enquanto o ponto~\eqref{eq:expectile} varia suavemente com $\kappa$.}
\answer{ex:07:2}{$p=1/3$, $x=0.75$, $\sigma_r=0.35$, $V(0.2)=0.083$.}
\answer{ex:07:3}{$V=\sqrt\kappa/(\sqrt\kappa+\sqrt{1-\kappa})$, de $0.25$ a $0.75$; dispersão $\propto\sqrt\eta$; para as gotas, $V=0.25+0.5\kappa$; as distribuições diferem em $0.05$ nos neurônios extremos, é preciso $\eta\lesssim0.01$.}
\answer{ex:07:4}{(a)~$J^*=2\sqrt{C\bar R}$. (b)~Custo excedente $\bigl(k^{1/2}+k^{-1/2}\bigr)/2-1$: $6\,\%$ para $k=2$ e $25\,\%$ para $k=4$; basta uma precisão ``dentro de um fator dois''.}
\answer{ex:07:5}{Pico de área $R(e^{-1}-e^{-2})\approx0.23R$, o equivalente a $2.3\,\text{s}$ de rampa; se a dopamina informa $V$, não há pico: o nível simplesmente sobe num degrau de $e$ vezes.}
\answer{ex:07:6}{O evento desloca o peso $\propto A\tau_f/(\tau_e+\tau_f)$, o nível dá um erro $s\tau_f$: $9\,\text{s}\le\tau_f\le17\,\text{s}$.}
\answer{ex:07:7}{De início, $\Delta P=0.80$, $M=0.90$. (a)~$\Delta P=0.74$ ($-8\,\%$), $M=0.43$ ($-52\,\%$); (b)~$\Delta P=0.40$ ($-50\,\%$), $M=0.80$ ($-11\,\%$): dupla dissociação.}
\answer{ex:07:8}{$N_{\text{ef}}\to2+\rho^2N\approx10^6$ tentativas: $10^4$ vezes menos que com um único fio, mas um vazamento de $1\,\%$ ainda custa um milhão de tentativas.}

\answer{ex:08:1}{(a)~$\approx1.7\cdot10^6$. (b)~$g=\frac{0.9}{\sqrt{0.19}}\,\sigma_{\text{depois}}/\sigma_{\text{antes}}\approx16.5$: a resposta só depende da razão entre os ruídos.}
\answer{ex:08:2}{(a)~Autovalores $-(1\pm g\beta)/\tau$; a diferença decai em $\tau/(1-g\beta)$: $40\ms$ com $g\beta=0.5$ (diferença de entradas amplificada três vezes), $200\ms$ com $0.9$ ($19$ vezes). (b)~$0.905$ e $1.105$; entre as fronteiras, só a entrada forte vence.}
\answer{ex:08:3}{$P(k)=e^{gu_k/\sigma}/\sum_le^{gu_l/\sigma}$, ou seja,~\eqref{eq:softmax}; $g=10\ln9\approx22$.}
\answer{ex:08:4}{$\bar\Delta=1/3$, $g^*\approx3\ln(1/h)$: $21$, $14$, $9$ contra $16$--$23$, $11$ e $8$ no modelo.}
\answer{ex:08:5}{$g^*$ de $16$--$23$ com $h=0.001$ a $6$--$8$ com $h=0.2$; a estimativa $3\ln(1/h)$ vale dentro de um fator $1.5$ para $h\le0.05$; com $h\gtrsim\eta/10$, a regra $Q\leftarrow Q+\eta(R-Q)$ não consegue assimilar o que a busca obteve, e $g^*$ empaca perto de $8$.}
\answer{ex:08:6}{$T_p=\tau\ln9=44\ms$ com $g_{\text{lo}}=0$ (modelo, $44\ms$) e $70\ms$ com $g_{\text{lo}}=0.25$: uma rajada de $2$--$3\tau$ com ganho perto de zero.}
\answer{ex:08:7}{(a)~Melhor constante $0.925$ ($g=8$), oráculo $0.929$ ($16/8$): quase não há o que ganhar. (b)~Por exemplo, épocas calmas com recompensas raras ($p\in[0,0.3]$) e instáveis com $h=0.1$: $0.850$ contra $0.872$; o ajuste captura cerca de um quarto com $\eta=0.2$ e quase tudo com $\eta=0.05$.}

\answer{ex:09:1}{(a)~$P_\infty=0.2$, memória de $20\,\text{s}$. (b)~$P_\infty=1$, memória de $1\,\text{s}$. (c)~A memória é proporcional à amplitude do ruído: seis vezes.}
\answer{ex:09:2}{$P(t)=P_\infty\coth(t\sqrt{q/R})$. (a)~$\frac12\sqrt{R/q}=5\,\text{s}$, metade da memória. (b)~$t=\frac12\ln21\,\sqrt{R/q}\approx15\,\text{s}$: é quanto o \nt{ACET} elevado deve durar.}
\answer{ex:09:3}{Variância $qT/2+R/(2T)$; $T^*=\sqrt{R/q}$, variância $\sqrt{qR}=P_\infty$, como no filtro~\eqref{eq:kalman}; cresce $(\kappa+1/\kappa)/2$ vezes: $1.25$ e $5.05$.}
\answer{ex:09:4}{$a>a_w=1-\theta/(mw)$: $0.15$ com $w=0.041$, $0.22$ com $w=0.045$.}
\answer{ex:09:5}{Os erros surgem com $M\approx40$--$60$; com $M=80$, $1.9$ neurônio a mais; com $M=100$, fração do padrão $0.86$; a interferência dos outros padrões tem variância $\approx(M-1)p(1-p)/N$, de modo que $M_{\max}\approx80$, $M_{\max}\propto N/p$.}
\answer{ex:09:6}{Por exemplo, $\eta(a)=\eta\min\bigl(1,\max(0,(a-0.3)/0.6)\bigr)$. Com $\eta a$, os três neurônios alheios entraram no padrão; com $\eta(a)$, nenhum, e a escrita com $a\ge0.9$ é a mesma (fração $1.00$).}
\answer{ex:09:7}{(a)~$40$ repetições; com a $\eta(a)$ do exercício~\ref{ex:09:6}, nunca. (b)~$200$ repetições, $10$ noites.}

\answer{ex:10:1}{Barramento $\approx10^{12}$ bit/s ($11.4$ bits por disparo), controle $\approx40$ bit/s; $2.5\cdot10^{10}\,\text{s}$, cerca de oitocentos anos.}
\answer{ex:10:2}{$a>1/3$ (com $a=0$, oscilações crescentes de $\approx67$~Hz). Não: a estabilidade é definida pelo produto $g(1-a)$.}
\answer{ex:10:3}{Média $\sigma^2R'x_j$, variância $(N+1)\,x_j^2\sigma^4R'^2$; a relação é $1/(N+1)$, são necessárias $n\approx z^2(N+1)\propto N$ tentativas.}
\answer{ex:10:4}{$L_v\in[32.9,\,47.9]\ms$: atraso de $35.4\ms$, janela de $\pm7.5\ms$; ligações falsas $2\Delta_{\max}r_ar_v\approx0.38$ por segundo.}
\answer{ex:10:5}{$0.098$ contra $0.180$ com o melhor $\alpha=0.2$: $45\,\%$ menor.}
\answer{ex:10:6}{$n=\ln\bigl(\ln1.5/(0.6g)\bigr)/\ln(1-\alpha)$: $24$ e $4$; com $g=2$, $11$ e $2$.}
\answer{ex:10:7}{$S'=S\bigl(1-4\eta\sigma^2+4\eta^2\sigma^4(G+2)\bigr)$, o melhor é $\eta\sigma^2=1/(2(G+2))$; $\approx10^6$ tentativas, cerca de um mês.}

\answer{ex:11:1}{(a)~$T=k^2/\bigl(\Phi(\Delta I/I)^2\bigr)=9\,\text{s}$. (b)~$45$ sensores, mancha de $\approx7\times7$, resolução $\approx7$ vezes pior.}
\answer{ex:11:2}{(a)~$2$--$3.3$ bits por disparo, $22$--$34\,\%$ do limite ($9.1$--$9.8$ bits). (b)~$\log_2\binom{400}{20}\approx111$ bits; em média $1$ tipo em comum.}
\answer{ex:11:3}{(a)~$\sigma/9\le I\le9\sigma$, $1.9$ ordem de grandeza. (b)~$6$ e $7$.}
\answer{ex:11:4}{(a)~$f<343/0.44\approx780$~Hz. (b)~$625$ disparos, $125$ neurônios.}
\answer{ex:11:5}{$21$ bits por evento, $4.8\cdot10^5$ eventos/s, $7$ eventos por pixel de borda: $136$ pixels/s. Uma câmera comum dá $1.25$ quadro/s.}
\answer{ex:11:6}{Em logaritmos, $\approx1.0\,\text{s}$ nos dois sentidos (teoria $0.96\tau$); linearmente, $0.2\,\text{s}$ para cima e $3.0\,\text{s}$ para baixo (teoria $0.18\tau$ e $3.0\tau$).}
\answer{ex:11:7}{(a)~$\ln I$ tem distribuição logística com mediana $\ln\sigma$ e escala $1$: dispersão $\pi/\sqrt3$ em logaritmos naturais, $0.79$ ordem de grandeza. (b)~$r=r_{\max}\Phi_I(I)^2$.}

\answer{ex:12:1}{$0.60$ (sem correlação seria $0.18$), limite $\sqrt{c/(rT)}\approx0.58$: cem neurônios distinguem apenas ``sim/não''.}
\answer{ex:12:2}{(a)~$10^8$ disparos, $\approx10\,\text{mJ}$: $0.3\,\%$ dos $3$~J do cérebro inteiro. (b)~$3$~J, $\approx300$ vezes mais.}
\answer{ex:12:3}{(a)~$\theta_A\in[2,3)$, $\theta_B\in[1,2)$; a figura alheia dá no máximo $2$ e $1$. (b)~$10\cdot96^2\approx9.2\cdot10^4$.}
\answer{ex:12:4}{(a)~$T_d=0.85\,\tau_a=0.34\,\text{s}$. (b)~$T_d\propto\tau_a$ e não depende de $I$: $\tau_a\approx3.5\,\text{s}$ (simulação: $3.1\,\text{s}$).}
\answer{ex:12:5}{$p\approx0.17$. A fração cai devagar: $0.90$, $0.88$, $0.86$, $0.84$, $0.82$, $0.79$ com $N=8$, $12$, $24$, $48$, $96$, $192$.}
\answer{ex:12:6}{As dez execuções não têm erro com $\tau_{\text{tr}}=20$, $50$ e $200\ms$ (com $30\ms$, uma execução em dez falhou); com $5$--$10\ms$, cerca de $0.5$; com $0.5$--$2\,\text{s}$, a qualidade cai para $0.86$. Por cima: o traço deve se apagar durante a pausa, $e^{-0.3\,\text{s}/\tau_{\text{tr}}}\ll1$.}
\answer{ex:12:7}{Um só atraso não basta: a janela de $\pm5\ms$ não cobre $\Delta x/v$ de $5$ a $20\ms$. Dois ramos de inibição com $5<d_1\le10\ms$ e $15\le d_2\le d_1+10\ms$.}
\answer{ex:12:8}{\texttt{two\_stage}: $\Delta=2$, empate com uma inversão, troca de resposta a partir de duas. Contador de uns: uma inversão; $0.57$ com $p=0.1$.}

\answer{ex:13:1}{(a)~$q=100^{1/99}\approx1.048$, passo $\approx4.8\%$; faixa $\approx2.2\cdot10^3$ ($67$~dB). (b)~Passo de $22f_1$, isto é, $220\%$ com $10f_1$.}
\answer{ex:13:2}{$P_{\text{falha}}\approx1.8\cdot10^{-4}$ e $2.8\cdot10^{-16}$: cerca de $300$ falhas por dia com $S=2$ e $5\cdot10^{-10}$ por dia com $S=4$.}
\answer{ex:13:3}{(a)~$H(s)=eF_1T_c/(1+sT_c)^2$, filtro de segunda ordem com frequência de corte $1/(2\pi T_c)\approx3.2$~Hz. (b)~$r\approx22$~disparos/s (pelo primeiro harmônico, $\approx20$).}
\answer{ex:13:4}{(a)~Com $Gd=1/e$ a raiz $s=-1/d$ é dupla, e com nenhum outro $G$ a raiz mais à direita fica mais à esquerda. (b)~$d=30\ms$: $G\approx12$~s$^{-1}$, $\tau=30\ms$; $d=150\ms$: $G\approx2.5$~s$^{-1}$, $\tau=150\ms$: a constante de tempo é igual ao atraso.}
\answer{ex:13:5}{(a)~$E\approx0.427$, $T\approx1.70\,\tau_a=0.68$~s (o modelo com $\tau=20\ms$ dá $0.84$~s). (b)~A equação não contém $I$; o ritmo existe com $b>\beta-1$.}
\answer{ex:13:6}{Com $b=0.8$ não há ritmo; com $b=1.05$, $1.2$, $1.5$, $2$, $4$, $T\approx2.73$, $1.74$, $1.19$, $0.84$, $0.45$~s; com qualquer $I$, $T=0.840$~s: a entrada muda só a amplitude, e não o andamento.}
\answer{ex:13:7}{A substituição $s=u-a$ dá a maior taxa de amortecimento $a+1/d$ com $G=e^{-1-ad}/d$; $a\ge33.3$~s$^{-1}$, $G=e^{-2}/d\approx4.5$~s$^{-1}$; quanto mais forte a co-contração (maior $a$), menor deve ser $G$.}
\answer{ex:13:8}{Problema aberto. Por exemplo, um limiar no acúmulo de fadiga $b[r_i-r_0]_+$ dá $I_{\min}=\beta b r_0/(b-\beta+1)\approx0.53$ com $b=4$, $r_0=0.2$ e período de $1.8$~s com $I=0.6$ até $0.52$~s com $I=3$.}

\answer{ex:14:1}{(a)~$0.17$~s e $1.5$~s. (b)~A partir de distâncias menores que $11$~cm.}
\answer{ex:14:2}{O toque sem dor não passa com nenhum $\mu$ enquanto $g\le\beta w_{q\mu}/w_{z\mu}=5$; o toque $\mu=1$ passa com $g>6$: uma sensibilização forte transforma um contato comum em dor.}
\answer{ex:14:3}{(a)~$g^*=(1-k_sC)/(1-k_sA)$, estável com $k_sA<1$. (b)~Sem toque, com $\nu\ge2$; com toque, já com $\nu\ge1.7$.}
\answer{ex:14:4}{$V(t)=-Pe^{-(T-t)/\tau_\gamma}$. (a)~$-Pe^{-T/\tau_\gamma}\approx-0.37P$. (b)~$+Pe^{-(T-t_e)/\tau_\gamma}\approx+0.61P$; ele cresce com $t_e$ até $+P$ e ensina a evitar a dor.}
\answer{ex:14:5}{Com $k_s=4$ não há travamento (o segundo estado estável existe com $k_s\ge4.58$); $T_{\min}\approx4.35$, $1.40$, $0.65$, $0.32\,\tau_s$ com $k_s=4.6$, $5$, $6$, $8$.}
\answer{ex:14:6}{$w_{q\nu}=4/9\approx0.444$, $q_0=2/9\approx0.222$, $w_{q\mu}=49/45\approx1.089$; o toque sem dor é seguro até $g=49/9\approx5.4$.}
\answer{ex:14:7}{(a)~Limiar $\nu_c(g)=\theta/g$ com $\nu_c\ge q_0/w_{q\nu}$, senão $(\theta+\beta q_0)/(g+\beta w_{q\nu})$; $g^*\approx2.45$, $1.0$, $0.25$. (b)~Questão aberta.}

\answer{ex:15:1}{(a)~$2\cdot10^5$ por neurônio, $6\cdot10^{12}$ no circuito. (b)~$\approx8$~bit/s por fio, pelo menos $2.5\cdot10^4$~s $\approx7$~h.}
\answer{ex:15:2}{Passo $2$: correlação entre vizinhos $0.943$, $\lambda$ de $4.18$ a $7.3\cdot10^{-4}$, razão $\approx5.7\cdot10^3$; passo $4$: correlação $0.8$, razão $\approx94$.}
\answer{ex:15:3}{(a)~$\lambda=0.988$ e $0.808$: $82$ e $4.7$ movimentos. (b)~$x_s^*=0.690$, $x_f^*=0.172$, soma $0.862$.}
\answer{ex:15:4}{(a)~$116$ movimentos. (b)~$19$. Um modelo com um só processo, com soma zero, não dá economia.}
\answer{ex:15:5}{(a)~$G=50$~s$^{-1}$, contra o limite de $10.5$~s$^{-1}$ sem modelo. (b)~Não: com $G=50$~s$^{-1}$, o atraso crítico é $d_c\approx103\ms$; com $d=150\ms$ é preciso $\hat k>0.445k$ (para quaisquer $G$ e $d$, $\hat k>k/2$).}
\answer{ex:15:6}{O erro cai abaixo de $0.5\%$ em $6$--$30$~s, com piso de $\approx(6$--$9)\cdot10^{-4}$, sendo $7.5\cdot10^{-4}$ com os melhores pesos; com $\eta=50$, os pesos ainda diferem dos melhores em até $0.2$ ao longo das direções mal condicionadas de $C$ (exercício~\ref{ex:15:2}).}
\answer{ex:15:7}{$\mathbf B=A\mathbf K$; $q_f/R\approx0.035$, $q_s/R\approx8.6\cdot10^{-4}$; com $4q_s$, $B_f\approx0.088$, $B_s\approx0.047$: o processo lento aprende mais que o dobro mais depressa, e o rápido, mais devagar.}

\answer{ex:16:1}{O sangue é um circuito RC com $\tau=t_{1/2}/\ln2=14.4$~min; máximo sobre mínimo $2^{T/t_{1/2}}=64$; o harmônico fundamental cai para $0.55$; razão $2$ com $t_{1/2}=T=60$~min.}
\answer{ex:16:2}{$\varphi^*=\arcsin\bigl((\omega-\omega_0)/K_\varphi\bigr)\approx41$~min ($\omega-\omega_0\approx0.047$~rad/dia), tempo de relaxação $1/(K_\varphi\cos\varphi^*)\approx3.9$~dias; após um deslocamento de $+6$ e $-6$~h, $7.7$ e $7.9$~dias.}
\answer{ex:16:3}{$K_c(n)=\sec^n(\pi/n)$: $8$ e $2.37$, erro de $11\%$ e $30\%$; para $n\to\infty$, $K_c\to1$, erro $\to50\%$.}
\answer{ex:16:4}{Os comandos extremos (acelerador $80$, freio $0$) dão $f=120$ em $0.76$~s; ``só o acelerador'', $2.33$~s, três vezes mais.}
\answer{ex:16:5}{$3\arctan(\omega\tau)+\omega d=\pi$, $K_c=(1+\omega^2\tau^2)^{3/2}$: $K_c=8$, $3.54$, $2.50$, $1.76$ e período $3.63$, $5.46$, $6.86$, $9.27\,\tau$; com $0.9K_c$ as oscilações se amortecem, com $1.1K_c$ se estabelecem.}
\answer{ex:16:6}{$0.60\bar c$, $0.87\bar c$, $0.90\bar c$, limite $0.91\bar c$; com concentração constante, a resposta é nula: esse destinatário só lê porções.}
\answer{ex:16:7}{Questão aberta. Com realimentação, o ritmo só existe para $K>8$, e o período é proporcional ao $\tau$ dos estágios e quase não depende de $K$ ($3.64\,\tau$ com $K=8.5$, $4.23\,\tau$ com $K=40$): mudar $K$ em uma vez e meia desloca o período em $2$--$4\%$, menos que o erro de medida. O que decide é abrir a malha (hormônio final constante na entrada do primeiro estágio mata esse ritmo) ou a resposta a um breve acréscimo de hormônio em diferentes fases do ciclo.}

\answer{ex:17:1}{$0.2$~W contra $81$~\textmu W; o fluxo bruto precisaria de no máximo $50$~pJ/bit.}
\answer{ex:17:2}{Com $c=0.1$: $5.6$, $8.7$, $9.2$~bit/s, limite $9.3$~bit/s; com $c=0$ e $N=1000$: $65$~bit/s.}
\answer{ex:17:3}{$B=\log_2N+P\log_2P+(1-P)\log_2\frac{1-P}{N-1}=4.87$, $4.37$, $3.29$~bits por caractere, isto é, $7.3$, $6.6$, $4.9$~bit/s; para $10$~bit/s são necessários $2.3$~caracteres por segundo.}
\answer{ex:17:4}{Variância do erro $2b/(Nm^2T)+\sigma_\xi^2$ por componente; as contribuições se igualam com $N\approx90$; limite $\tfrac12\log_2(1+0.25/0.04)\approx1.4$~bit por componente por janela.}
\answer{ex:17:5}{O par converge para $\eta_bd^2+\eta_db^2<2$, isto é, aproximadamente para $\eta<1$: com $0.8$ converge, com $1.1$ há oscilações estáveis de período $2$, com $1.5$ não periódicas, com $2$ diverge; as velocidades dos aprendizes devem ser separadas.}
\answer{ex:17:6}{Limiar $\approx4.4\sigma$; $102$~kbit/s, $0.10$~mW contra $205$~mW para o fluxo bruto, $2000$ vezes menos; só $5.7\cdot10^{-5}$ das amostras saturam (mais de $3$ disparos).}
\answer{ex:17:7}{Questão aberta. Cerca de $46$~bit/s com $\text{SNR}\approx0.9$ e cerca de $330$~bit/s com ruído independente. Se o obstáculo é o ruído parecido com o sinal, a informação extraída satura com o aumento do número de canais; se é o desconhecimento do código, ela cresce com um decodificador melhor (por exemplo, um que leve em conta os tempos dos disparos).}

\answer{ex:18:1}{(a)~$0.9949\le w\le1.0049$, isto é, $|1-w|\lesssim0.005$. (b)~$K\ge400$ sinapses.}
\answer{ex:18:2}{$(\omega_R\tau_w)^2=\sqrt{\alpha(\alpha+2+2\varrho)}/\varrho-1$, $\varrho=\tau_m/\tau_w$; $f_R\approx11.1$~Hz; $|Z(\omega_R)|/|Z(0)|=4.6$.}
\answer{ex:18:3}{(a)~$f=1/\bigl(\tau\ln\frac{\mu}{\mu-\theta}\bigr)$; para $1$~Hz é preciso $\mu/\theta-1\approx3.7\cdot10^{-44}$. (b)~$T=\pi\tau/\sqrt\mu$, $f\approx3.2$~Hz: $f\propto\sqrt\mu$.}
\answer{ex:18:4}{O integrador funciona com $f\lesssim17.7$~Hz (filtro passa-baixas); o ressonador, só na faixa de $5.5$--$22$~Hz (filtro passa-faixa).}
\answer{ex:18:5}{(a)~$T_{\min}=\frac{\tau}{w-1}\ln\frac{0.5}{0.3}\approx10.2\ms$. (b)~$B>w-1-0.2=0.3$.}
\answer{ex:18:6}{Constante de tempo do ajuste $\tau/(\eta\langle r^2\rangle)$, $\eta=\tau\ln10/60\,\text{s}\approx3.8\cdot10^{-3}$; com $I\neq0$ o peso se desloca, e é preciso subtrair de $e$ uma cópia do comando $I/\tau$.}
\answer{ex:18:7}{Questão aberta. O ressonador ganha só $1.35$ vezes a $11$~Hz e perde $17$ vezes a $1$~Hz; o ganho principal da reconfiguração está na seleção de banda por limiar (exercício~\ref{ex:18:4}).}

\answer{ex:19:1}{(a)~Uma sinapse dá $6.7\mV$ ($R=66.7$~M$\Omega$), então para chegar ao limiar são necessárias no mínimo $4$ (três bastam só assintoticamente); $8$ em $10\ms$; $14$ em $5\ms$. (b)~$0.1\ms\ll\tau_m$; a correção pelo vazamento é $\sim0.25\%$.}
\answer{ex:19:2}{(a)~$1.75\cdot10^{10}$ (um quarto dos misturadores responde a tudo), $4.4\cdot10^9$ e $0.06$ (com $k=40$ quase nunca responde). (b)~$10^3$ vezes: $2\cdot10^5$ contra $200$.}
\answer{ex:19:3}{(a)~$C(2n,n)=2^{2n-1}$ pela simetria dos coeficientes binomiais. (b)~$0.93$ e $0.09$; a largura da transição é $\sim\sqrt{2n}$ em $P$, a relativa, $\sim1/\sqrt n$.}
\answer{ex:19:4}{Com um só dendrito, $V_{\text{corpo}}<\kappa E_E<\theta$ para qualquer número de entradas; com $g_0=0.5g_L$ são necessárias $n\ge3$ em cada dendrito.}
\answer{ex:19:5}{$I\ge C\sigma_V/\sigma_t=15$~nA, isto é, $150$ sinapses síncronas, o que é irrealista; ao neurônio pequeno bastaria $1$~nA, e com a sinapse de $4$~nA o jitter é de $5$~\textmu s.}
\answer{ex:19:6}{Máximo na ponta próxima: $0.03\,\tau_m$ e $0.97$ ($L=0.2$); $0.32\,\tau_m$ e $0.66$ ($L=1$); $0.79\,\tau_m$ e $0.33$ ($L=2$). A estimativa~\eqref{eq:dendriteload} pela capacitância total só vale para $L\ll1$.}
\answer{ex:19:7}{Questão aberta. Com $m$ fixo, o tempo de resposta cresce linearmente com o $P$ memorizado; com fração fixa $m=\varphi n$, ele deixa de depender de $n$, e a lentidão do neurônio grande é consequência da soma de muitas entradas, não do tamanho em si.}

\answer{ex:20:1}{$t_{\text{v}}=\tau_r\ln\frac{1-L}{1-H}=16.8$~h, $t_{\text{s}}=\tau_d\ln\frac HL=5.8$~h, período $22.5$~h; o oscilador é puxado para $24$~h pela oscilação diária dos limiares: captura de fase~\eqref{eq:pll}.}
\answer{ex:20:2}{(a)~$L=0.111$; com $L=0.092$ o sono duraria $9.6$~h. (b)~$22\%$ ($1/0.82=1.22$), e não $18\%$.}
\answer{ex:20:3}{$H=\frac{p-1}{p-1/q}=0.623$, $L=H/q=0.093$, onde $p=e^{16/\tau_r}$, $q=e^{8/\tau_d}$. Depois de acordar após $4$~h, a fase se desloca para sempre (adormece $7.2$~h mais cedo); com limiares oscilantes, ela volta em dois ou três dias.}
\answer{ex:20:4}{(a)~$r_\pm=\frac12\bigl(1\pm\sqrt{1-4k/w}\bigr)=0.947,\ 0.053$; $a_{\text{s}}=3.51$, $a_{\text{i}}=0.49$. (b)~Atividade $\approx0.33$~s, silêncio $\approx0.59$~s, período $\approx0.93$~s, fração de atividade $0.36$.}
\answer{ex:20:5}{$a_{\text{s}}/r_+<b<a_{\text{i}}/r_-$: $3.71<b<9.20$. O \nt{ACET} reduz $b$ para menos de $3.71$: a interseção passa ao ramo superior, e o estado ativo ($r\approx1$) fica estável.}
\answer{ex:20:6}{$4.9$, $6.8$, $8.8$, $5.5$, $8.9$~h para $D=24$, $32$, $40$, $48$, $64$: a duração é definida pela fase diária do adormecer, e não pela dívida.}
\answer{ex:20:7}{Após B, o erro em A é em média $0.51$ (de $0.19$ a $1.08$ em $20$ conjuntos; a resposta nula dá $1$). Intercalando, os dois erros ficam abaixo de $10^{-4}$.}
\answer{ex:20:8}{Questão aberta. O escalonamento uniforme preserva as razões dos pesos, mas só preserva as respostas do neurônio de limiar se o limiar for escalonado pelo mesmo fator; os replays intercalados mudam as razões. A invariância dos contatos grandes contradiz um escalonamento puramente uniforme.}

\answer{ex:21:1}{(a)~$T=\tau_{\text{rec}}\ln\frac{1}{1-x^*}$: $1.84$~s e $3.68$~s. (b)~De $3.36$ a $4.24$~s: a sensibilidade $\dd T/\dd x^*=\tau_{\text{rec}}/(1-x^*)=80$~s; com $x^*\to1$ as salvas ficam irregulares.}
\answer{ex:21:2}{(a)~$8.6$~W, como em~\eqref{eq:energy}. (b)~$205$~Mbit/s, $0.2$~W: $2\cdot10^3$ vezes mais que a cultura ($10^{-4}$~W).}
\answer{ex:21:3}{$x_{\text{est}}=1/(1+Ur_b\tau_{\text{rec}})<x^*$ para $r_b>(1/x^*-1)/(U\tau_{\text{rec}})=0.28$~Hz com $x^*=0.9$ e $0.025$~Hz com $x^*=0.99$; a força da sinapse cai para $x_{\text{est}}$, isto é, em $10\%$ e $1\%$.}
\answer{ex:21:4}{Os $u(t-k)$ são ortonormais, $m_k=\|P_Vu(t-k)\|^2$, em que $P_V$ é o projetor sobre o espaço gerado $N$-dimensional das saídas; pela desigualdade de Bessel, $\sum_k m_k\le N$.}
\answer{ex:21:5}{$\mathrm{MC}\approx9$--$11$ com $\tanh$ e $15$--$18$ para o reservatório linear (três sementes), ambos $\le30$: a não linearidade se paga com memória.}
\answer{ex:21:6}{(a)~$n\le102$. (b)~$5.6$.}
\answer{ex:21:7}{Problema em aberto. Controle-chave: congelar a leitura e verificar se a melhora persiste depois de uma pausa sem estimulação; se não persistir, mudou o estado, não os pesos.}

\answer{ex:22:1}{(a)~$32\%$ e $100\%$. (b)~$R_\uparrow-R_\downarrow\ge2\sqrt{2000\,\text{Hz}/1\,\text{s}}\approx89$~Hz, apenas $4.5\%$ da frequência total.}
\answer{ex:22:2}{(a)~$n\ge4(p_1q_1+p_2q_2)/\Delta^2\approx560$. (b)~$\approx56$.}
\answer{ex:22:3}{(a)~Balanço detalhado: $\rho PK$ é simétrico. (b)~A fração de erros é igual à média harmônica $1/\langle1/P\rangle=0.18$, contra $0.5$ sem realimentação.}
\answer{ex:22:4}{(a)~Um paralelogramo de área $4h^2=0.04$ (levando em conta o limite de $p$, numericamente $0.045$). (b)~$\approx0.18$.}
\answer{ex:22:5}{$0.49$: parte das culturas vai para uma região em que quase sempre erra e vagueia por lá. Com o limite, $1.00$: num conjunto limitado, a densidade $\propto1/P$ se concentra na região absorvente.}
\answer{ex:22:6}{(a)~São precisos $\ge5$ níveis; $8$ eletrodos bastam. (b)~Não: seria preciso $r_{\max}\ge25/T=100$~Hz.}
\answer{ex:22:7}{Problema em aberto. O tempo da busca aleatória cresce com $d$ aproximadamente como a razão de volumes $(L/h)^d$; o da busca com o sinal do erro, linearmente em $d$. Separa as hipóteses um experimento com troca: depois do erro, um estímulo previsível mas forte; depois do acerto, um imprevisível mas fraco; são necessárias algumas dezenas de culturas em cada grupo.}


\chapter{Suplemento experimental}
\label{sec:supplement}

Para cada capítulo, reunimos aqui suas principais afirmações e aquilo em que elas se apoiam: o experimento em que foram medidas, o que exatamente foi medido e onde foi publicado. O status na coluna da direita diz o quanto a afirmação é firme: \emph{medido}, há um experimento direto; \emph{teoria}, uma consequência matemática deduzida no capítulo ou num trabalho clássico; \emph{modelo}, o resultado de um cálculo com um modelo do livro (scripts na pasta \texttt{models/}); \emph{estimativa}, um número de ordem de grandeza sem medição direta; \emph{hipótese}, uma explicação plausível que ainda não foi provada por experimento. Onde não há experimento, isso é dito.

\section{Capítulo~\ref{sec:neurontypes}. Tipos de neurônios}

Os números do capítulo se apoiam em contagens diretas de células no cérebro humano, onde elas existem; a regra ``um neurônio, um neurotransmissor'', em atlas celulares e em experimentos que também encontraram as exceções; o papel de \nt{DOPA}, em registros de disparos; e os papéis dos demais canais de difusão, em hipóteses.

{\footnotesize
\setlength{\tabcolsep}{3pt}\renewcommand{\arraystretch}{1.15}
\begin{longtable}{>{\raggedright}p{0.5cm}>{\raggedright}p{4.0cm}>{\raggedright}p{5.6cm}>{\raggedright}p{2.3cm}>{\raggedright\arraybackslash}p{1.7cm}}
\toprule
N.º & Proposição & Experimento e o que foi medido & Fonte & Status \\
\midrule\endfirsthead
\toprule
N.º & Proposição & Experimento e o que foi medido & Fonte & Status \\
\midrule\endhead
1 & O cérebro humano tem cerca de $8.6\cdot10^{10}$ neurônios, e quase todos os neurônios \nt{GLUT} são pequenos neurônios do cerebelo (tab.~\ref{tab:types}) & Fracionador isotrópico: o tecido é dissolvido numa suspensão homogênea de núcleos, e contam-se os núcleos com o marcador neuronal NeuN. Em homens adultos, $86.1\pm8.1$ bilhões de neurônios; apenas $19\,\%$ deles estão no córtex cerebral, a maior parte está no cerebelo & \cite{Azevedo2009} & medido \\
2 & Quase todo neurônio tem um neurotransmissor principal (proposição~\ref{prop:dale}), mas há exceções & Atlas transcriptômico do cérebro do camundongo: cerca de $4\cdot10^6$ células, $5322$ clusters; a cada cluster se atribui um neurotransmissor pelos genes de síntese e empacotamento. As exceções foram medidas diretamente: neurônios \nt{DOPA} em fatias liberam também GABA pelo mesmo transportador vesicular e inibem neurônios do estriado; em parte dos axônios \nt{DOPA}, glutamato e dopamina saem de terminais diferentes do mesmo fio (microscopia eletrônica e optogenética) & \cite{Strata1999,Yao2023,Tritsch2012,Zhang2015} & medido \\
3 & Toda a coluna da matriz de pesos tem o mesmo sinal~\eqref{eq:dalesign} & Dedução no capítulo a partir da proposição~\ref{prop:dale} e do fato de que os receptores de glutamato são quase todos excitatórios, e os de GABA, inibitórios. Não há verificação direta de ``todos os destinatários de um neurônio recebem peso do mesmo sinal'' como regra geral; contraexemplos: a liberação conjunta de GABA por neurônios \nt{DOPA} (linha~2) e destinatários com receptores de sinais diferentes (linha~11) & dedução no capítulo & teoria \\
4 & De três quartos a quatro quintos dos neurônios do córtex são \nt{GLUT}, de um quinto a um quarto são \nt{GABA} & Imunomarcação para GABA em colunas de $50$\,\textmu m de largura por toda a espessura do córtex em $10$ áreas corticais de macaco: em $8$ áreas os neurônios GABAérgicos são cerca de $25\,\%$ de todos os neurônios; no córtex visual primário, menos de $20\,\%$ & \cite{Hendry1987} & medido \\
5 & Um neurônio \nt{GLUT} tem cerca de $10^4$ saídas & Estereologia em autópsias: no neocórtex de homens jovens, $1.64\cdot10^{14}$ sinapses (microscopia eletrônica, dissector) e cerca de $2.3\cdot10^{10}$ neurônios. A razão dá $\sim7\cdot10^3$ sinapses por neurônio; em média, o número de entradas é igual ao de saídas & \cite{Tang2001synapses,Pakkenberg1997} & medido \\
6 & A saída de um neurônio \nt{DOPA} se ramifica em $10^5$--$10^6$ terminais; é uma estimativa pela cobertura do alvo, não uma contagem & Marcação viral da membrana de neurônios \nt{DOPA} isolados de rato e reconstrução completa do axônio: um neurônio cobre de $0.45$ a $5.7\,\%$ (em média $2.7\,\%$) do volume do estriado. Mediu-se a cobertura, não o número de terminais: este foi deduzido da cobertura em ordem de grandeza, sem contagem direta dos terminais. Para \nt{SERO} e \nt{NORA}, os números de terminais são estimativas grosseiras & \cite{Matsuda2009} & medido \\
7 & Os neurônios de difusão são poucos: \nt{DOPA} $\sim5\cdot10^5$ (o principal dos dois grupos, limite inferior), \nt{SERO} $\sim3\cdot10^5$ (soma de duas contagens parciais), \nt{HIST} $\sim6\cdot10^4$ (tab.~\ref{tab:types}, fig.~\ref{fig:typepie}) & Contagens estereológicas em humanos: $550\,000$ neurônios pigmentados na substância negra (sem a área tegmental ventral); $235\,000$ neurônios no núcleo dorsal da rafe (todos os neurônios do núcleo) e mais $125\,000$ neurônios serotoninérgicos no tegmento pontino; cerca de $64\,000$ neurônios histaminérgicos do hipotálamo. Para \nt{NORA}, \nt{GLYC}, \nt{ACET} e \nt{ADRE} não há contagem direta: na tabela são estimativas grosseiras de ordem de grandeza & \cite{Pakkenberg1991,Baker1990,Baker1991,Airaksinen1991} & medido \\
8 & O glutamato na fenda sobe até cerca de um milimolar e cai em $\sim1\ms$ & Pelo deslocamento de um bloqueador competitivo de dissociação rápida dos receptores NMDA durante a transmissão sináptica (cultura de neurônios do hipocampo): pico de $1.1$\,mM, constante de decaimento de $1.2\ms$ & \cite{Clements1992} & medido \\
9 & No córtex em funcionamento, as correntes excitatória e inibitória estão quase equilibradas, e o neurônio fica perto do limiar & Teoria: uma rede com conexões fortes chega sozinha ao regime equilibrado. Experimento: no córtex pré-frontal do furão in vivo, durante os estados ``up'', o potencial de reversão da corrente sináptica se mantém em torno de $-37$\,mV por centenas de milissegundos, embora a condutância varie em $\sim21$\,nS: excitação e inibição sobem e descem juntas & \cite{vanVreeswijk1996,Haider2006} & medido \\
10 & Os neurônios \nt{DOPA} informam o erro de predição de recompensa~\eqref{eq:rpe} (fig.~\ref{fig:rpe}) & Disparos de neurônios \nt{DOPA} de macaco: rajada ao suco inesperado; após o aprendizado, rajada ao sinal e nenhuma resposta ao suco previsto; queda da frequência quando o suco prometido não é dado. Num experimento quantitativo, a frequência é proporcional à diferença entre a recompensa atual e a média ponderada das anteriores, mas só para erros positivos. A equação~\eqref{eq:rpe} em si é o algoritmo de diferenças temporais & \cite{Schultz1997,Bayer2005,SuttonBarto2018} & medido \\
11 & O sinal e a velocidade são definidos pelo receptor: ionotrópicos, frações de milissegundo; metabotrópicos, muito mais lentos (proposição~\ref{prop:receptor}) & Pulsos rápidos de glutamato em fragmentos de membrana de neurônios do hipocampo: a corrente pelos receptores ionotrópicos sobe (de $20$ a $80\,\%$) em $0.2$--$0.6\ms$ e decai em $2$--$3\ms$. O potencial inibitório lento das células piramidais é abolido por um bloqueador seletivo do receptor metabotrópico de GABA. Duas famílias de receptores de dopamina com ação oposta; no estriado, os neurônios com receptor D1 precisam de dopamina para reforçar sinapses, e os com D2, para enfraquecê-las & \cite{Colquhoun1992,DutarNicoll1988,Beaulieu2011,Shen2008} & medido \\
12 & O canal de difusão não é um fio só, mas um feixe de poucas linhas (seção~\ref{sec:buses}) & Marcação viral e genética dos neurônios serotoninérgicos do núcleo dorsal da rafe do camundongo: os neurônios que enviam saída para a amígdala e para o córtex frontal ficam em lugares diferentes do núcleo, ramificam-se em regiões complementares e respondem de modo oposto a um estímulo aversivo. Revisão: subgrupos de neurônios do locus coeruleus (\nt{NORA}) codificam processos diferentes & \cite{Ren2018,Poe2020} & medido \\
13 & \nt{NORA} define o ganho $g$ & Hipótese do ganho adaptativo; modelo de rede em que as catecolaminas aumentam a inclinação da característica de transferência. Não há medição direta de ``$g$ cresce com a noradrenalina'' na rede inteira; mediu-se que o diâmetro da pupila acompanha os disparos dos neurônios do locus coeruleus do macaco & \cite{AstonJones2005,ServanSchreiber1990,Joshi2016} & hipótese \\
14 & \nt{ACET} alterna ``escrita/leitura'' e define a taxa de aprendizado & Hipótese. Base: em fatias de córtex olfatório de rato, agonistas da acetilcolina ($100$\,\textmu M) enfraquecem as sinapses internas, de ``lembrança'', em mais de $60\,\%$, e as de entrada em menos de $15\,\%$ & \cite{Hasselmo2006,HasselmoBower1992} & hipótese \\
15 & \nt{SERO} define o fator de desconto $\gamma$; os neurônios serotoninérgicos trabalham de forma regular, alguns disparos por segundo, e quase silenciam numa das fases do sono & Hipótese ``serotonina $=\gamma$''. O argumento mais forte: a ativação optogenética dos neurônios serotoninérgicos do núcleo dorsal da rafe do camundongo reduz o número de desistências prematuras na espera de uma recompensa adiada. A frequência de disparos e o silêncio no sono com movimentos oculares rápidos foram medidos (revisão de registros) & \cite{Doya2002,Miyazaki2014,JacobsAzmitia1992} & hipótese \\
\bottomrule
\end{longtable}}

\section{Capítulo~\ref{sec:glutcomputer}. O que os neurônios \nt{GLUT} calculam}

As afirmações lógicas do capítulo são teoremas sobre circuitos de pesos não negativos, deduzidos no próprio capítulo; a experiência confirma o modelo de neurônio e o fato de que circuitos montados segundo esses teoremas (detector de coincidência, linha de atraso, grupo autossustentado, cadeia) de fato existem no cérebro.

{\footnotesize
\setlength{\tabcolsep}{3pt}\renewcommand{\arraystretch}{1.15}
\begin{longtable}{>{\raggedright}p{0.5cm}>{\raggedright}p{4.0cm}>{\raggedright}p{5.6cm}>{\raggedright}p{2.3cm}>{\raggedright\arraybackslash}p{1.7cm}}
\toprule
N.º & Proposição & Experimento e o que foi medido & Fonte & Status \\
\midrule\endfirsthead
\toprule
N.º & Proposição & Experimento e o que foi medido & Fonte & Status \\
\midrule\endhead
1 & O neurônio é um circuito RC com limiar e reset~\eqref{eq:glutneuron} & No corpo de neurônios piramidais do córtex de rato (fatias) injetou-se uma corrente ruidosa, semelhante ao fluxo sináptico no cérebro vivo. A dependência da frequência média em relação à média e à dispersão da corrente é descrita pelo neurônio integrador com vazamento, se for acrescentada uma adaptação lenta da frequência. As faixas de $\tau$ e dos atrasos são dadas no capítulo sem referência experimental & \cite{Rauch2003} & medido \\
2 & O modelo~\eqref{eq:glutneuron} é uma simplificação: os ramos de entrada emitem eles mesmos disparos locais & Registro direto dos dendritos de neurônios piramidais do córtex visual do camundongo in vivo: o estímulo visual evoca disparos dendríticos locais que dependem dos receptores NMDA; suprimi-los alarga a sintonia do neurônio à orientação & \cite{Smith2013} & medido \\
3 & Janela de coincidência $\Delta_{\max}=\tau\ln\frac{w}{\theta-w}$~\eqref{eq:window}; com $\theta=1.5w$ ela vale $0.7\tau$ & Consequência do modelo~\eqref{eq:glutneuron}. Há medição direta de uma janela estreita de soma em neurônios com inibição rápida (capítulo~\ref{sec:gaba}, linha~3 de sua tabela) & dedução no capítulo & teoria \\
4 & Uma rede só de neurônios \nt{GLUT} calcula apenas funções monótonas das entradas (proposição~\ref{prop:monotone}) & Demonstração no capítulo: iteração a partir do silêncio com lado direito não decrescente. É uma afirmação sobre o modelo; não há experimento que a teste numa rede viva sem inibição & dedução no capítulo & teoria \\
5 & Os órgãos dos sentidos fornecem um par de fios: umas células respondem ao aumento do estímulo, outras à diminuição & Retina: já nas células bipolares, o sinal dos cones se divide em canais liga e desliga. Bloqueio farmacológico das bipolares liga no macaco: os canais seguem separados até o córtex; após o bloqueio piora a detecção do aumento de luz, mas não da diminuição & \cite{Schiller1992} & medido \\
6 & Com código de dois fios, uma rede de neurônios \nt{GLUT} calcula qualquer função lógica de modo assíncrono, sem ciclo comum, ao preço do dobro de neurônios (proposição~\ref{prop:dualrail}, \eqref{eq:dualrail}, fig.~\ref{fig:dualxor}) & O código de dois fios e a detecção de prontidão pelo próprio sinal são técnica padrão de circuitos assíncronos insensíveis a atrasos. O circuito de OU exclusivo com seis neurônios é montado no capítulo. Não há experimento mostrando que o cérebro calcula funções lógicas justamente com código de dois fios & \cite{Sparso2001} & teoria \\
7 & A maior diferença de tempo de chegada do som aos ouvidos no ser humano é $\approx0.6\ms$, e o cérebro distingue cerca de dez microssegundos & Ouvintes num experimento com duas alternativas de resposta: o limiar de discriminação da diferença de tempo (75\,\% de acertos) é de $9$\,\textmu s para ruído na banda de $150$--$1700$\,Hz, $11$\,\textmu s para um tom de $1$\,kHz, $28$\,\textmu s para um clique. O limite de $0.6\ms$ é cálculo do capítulo, $0.22\,\text{m}/343\,\text{m/s}$ & \cite{KlumppEady1956} & medido \\
8 & Nas aves, a diferença de tempos vira lugar por meio de linhas de atraso e detectores de coincidência~\eqref{eq:itd} (fig.~\ref{fig:itd}) & Suindara: os axônios dos dois ouvidos entram no núcleo dos detectores de coincidência por lados opostos; o tempo de chegada dos disparos ao longo do axônio muda de forma ordenada com a profundidade, e a variação dos atrasos através do núcleo ($160$\,\textmu s em $700$\,\textmu m) coincide com a faixa de atrasos interaurais & \cite{Carr1990} & medido \\
9 & Nos mamíferos, a fileira de atrasos não foi encontrada; a mesma tarefa é resolvida por detectores de coincidência com inibição de tempo preciso vinda de neurônios \nt{GLYC} & Registro in vivo na oliva superior medial do gerbil: as respostas não formam um mapa de direções; a aplicação de glicina e de seu bloqueador por micropipeta mostra que uma inibição \emph{glicinérgica} com tempo calculado com precisão define a faixa de trabalho dos atrasos. A inibição aqui vem de \nt{GLYC}, não de \nt{GABA} & \cite{Brand2002,Grothe2010} & medido \\
10 & Um grupo com forte realimentação é um flip-flop; a atividade autossustentada dura segundos após o estímulo (fig.~\ref{fig:bistable}) & Córtex pré-frontal do macaco numa tarefa de movimento ocular adiado para um ponto memorizado (atraso de $1$--$6$\,s): $87$ de $170$ neurônios ligados à tarefa mudam de frequência durante o atraso, e em $79\,\%$ deles isso depende da direção do ponto memorizado. Que essa atividade é mantida por uma realimentação com dois estados estáveis é teoria & \cite{Funahashi1989,Wang2001} & medido \\
11 & O bit é gravado se o aporte durar mais de $\approx0.7\tau$ (fig.~\ref{fig:recurrentgroup}b) & Cálculo da equação $\tau\,\dd r/\dd t=-r+f(wr+I)$ pelo método de Euler com $w=1$, $I=0.6$; não há script em \texttt{models/}. Não há experimento & cálculo no capítulo & modelo \\
12 & A memória pode ser guardada na facilitação das sinapses, quase sem disparos & Teoria: o cálcio residual nos terminais mantém o registro por cerca de um segundo. Base: no córtex pré-frontal medial do furão (registro de vários neurônios ao mesmo tempo) há uma sub-rede de neurônios piramidais ligados por sinapses facilitadoras com efeito prolongado. Revisão das observações de intervalos ``silenciosos'' durante a retenção na memória. Qual maneira o córtex usa e quando é questão em aberto & \cite{Mongillo2008,WangMarkram2006,Stokes2015} & hipótese \\
13 & A memória é apagada pela fadiga: o estoque de neurotransmissor se esgota & Registros pareados de neurônios piramidais do córtex: a resposta da sinapse cai com disparos frequentes, e a velocidade da queda é dada pela probabilidade de liberação. Que é exatamente isso que apaga o grupo autossustentado não foi mostrado diretamente & \cite{Tsodyks1997} & medido \\
14 & Uma cadeia de grupos é um temporizador disparado por evento; nas aves canoras, os neurônios disparam estritamente em sequência & Registro de neurônios identificados do centro vocal superior do mandarim durante o sono e o canto: cada neurônio que envia saída ao núcleo motor emite uma única rajada curta num momento preciso do canto, e os neurônios disparam em sequência & \cite{Hahnloser2002} & medido \\
\bottomrule
\end{longtable}}

\section{Capítulo~\ref{sec:gaba}. \nt{GLUT} e \nt{GABA} juntos}

As afirmações lógicas e dinâmicas do capítulo são deduzidas nele mesmo, por análise linear e pela lei de Ohm; a experiência mostra que os circuitos em que elas se apoiam (veto com atraso, inibição direta logo após a excitação, estabilização pela inibição, ritmo do laço \nt{GLUT}--\nt{GABA}) de fato funcionam no cérebro.

{\footnotesize
\setlength{\tabcolsep}{3pt}\renewcommand{\arraystretch}{1.15}
\begin{longtable}{>{\raggedright}p{0.5cm}>{\raggedright}p{4.0cm}>{\raggedright}p{5.6cm}>{\raggedright}p{2.3cm}>{\raggedright\arraybackslash}p{1.7cm}}
\toprule
N.º & Proposição & Experimento e o que foi medido & Fonte & Status \\
\midrule\endfirsthead
\toprule
N.º & Proposição & Experimento e o que foi medido & Fonte & Status \\
\midrule\endhead
1 & Os neurônios \nt{GABA} são de um quinto a um quarto dos neurônios do córtex & Imunomarcação para GABA em $10$ áreas do córtex de macaco: cerca de $25\,\%$ dos neurônios em $8$ áreas, menos de $20\,\%$ no córtex visual primário. Revisão da diversidade dos neurônios inibitórios do córtex & \cite{Hendry1987,Markram2004} & medido \\
2 & O que a inibição faz é decidido em grande parte pelo local de contato: dendrito, corpo, local de nascimento do disparo, terminal do fio de outro neurônio & Revisão: as classes de neurônios \nt{GABA} do córtex diferem pelo alvo no destinatário. Registro simultâneo do corpo e do dendrito de um neurônio piramidal: a inibição direta é muito mais forte no corpo do que nos dendritos. A inibição no terminal do fio de outro neurônio, na medula espinal, reduz a liberação de neurotransmissor de uma única linha de entrada (revisão das medições) & \cite{Markram2004,Pouille2001,Rudomin1999} & medido \\
3 & A troca de sinal por um intermediário custa um atraso, cerca de um milissegundo; a inibição que chega logo após a excitação estreita a janela de coincidência & Neurônios piramidais do hipocampo de rato em fatias: a inibição por um intermediário chega logo após a excitação direta com atraso curto, e a janela em que duas entradas excitatórias se somam até o limiar é menor que $2\ms$. Nos dendritos, onde essa inibição é mais fraca, a janela é mais larga & \cite{Pouille2001} & medido \\
4 & O veto $a\wedge\bar b$~\eqref{eq:veto} com OU e a constante um é funcionalmente completo (proposição~\ref{prop:gabacomplete}) & Demonstração no capítulo. Resultado clássico: uma rede de elementos de limiar com entradas inibitórias realiza qualquer expressão lógica. Afirmação sobre o modelo, sem experimento & \cite{McCullochPitts1943} & teoria \\
5 & O veto com atraso dá um detector de direção do movimento com velocidade ótima $v^*=\Delta x/d$~\eqref{eq:vstar} & Retina do coelho: a seletividade à direção é criada por uma inibição que, no movimento na direção ``nula'', anula a excitação. O registro separado das entradas excitatória e inibitória das células seletivas mostrou que as células amácrinas estreladas (\nt{GABA}) as inibem de forma assimétrica, só pelos prolongamentos voltados para o lado ``nulo''. A fórmula de $v^*$ é dedução do capítulo & \cite{Barlow1965,Fried2002} & medido \\
6 & A inibição por derivação divide a tensão em vez de subtrair~\eqref{eq:shunt} (fig.~\ref{fig:shunt}) & Lei de Ohm para um divisor de três condutâncias, com o potencial de equilíbrio do cloreto perto do repouso; para um neurônio sem disparos & dedução no capítulo & teoria \\
7 & Sobre a frequência de disparos a derivação age por subtração, e a divisão vem da inibição junto com ruído de fundo equilibrado & Modelo: num neurônio que dispara, a corrente pela derivação é quase constante, e o efeito sobre a frequência é subtrativo. Experimento: em neurônios piramidais do córtex somatossensorial de rato (fatias) introduziram-se, por fixação dinâmica, condutâncias de fundo excitatórias e inibitórias; o aumento simultâneo e equilibrado de ambas divide a inclinação da resposta sem deslocamento notável da frequência. Revisão: a divisão pela atividade total ocorre em muitas regiões do cérebro & \cite{HoltKoch1997,Chance2002,Carandini2012} & medido \\
8 & Com $\beta>1$ a inibição mútua escolhe um único vencedor (proposição~\ref{prop:wta}, \eqref{eq:wta}) & Os autovalores $-(1\pm\beta)/\tau$ são dedução do capítulo. As frequências $0.73$ e $0.53$ com $\beta=0.5$ são o ponto fixo $r_{1,2}=(I_{1,2}-\beta I_{2,1})/(1-\beta^2)$; não há script em \texttt{models/}. O capítulo não cita experimento direto & dedução no capítulo & teoria \\
9 & Reset da memória por sinal via \nt{GABA}; dois grupos com inibição mútua formam um latch & Decorre da fig.~\ref{fig:bistable} e da proposição~\ref{prop:wta}. Sem experimento & dedução no capítulo & teoria \\
10 & O equilíbrio do laço \nt{GLUT}--\nt{GABA} é estável se a inibição for forte e rápida o bastante (proposição~\ref{prop:ei}) & Modelo de duas populações~\eqref{eq:ei}; as condições (i) e (ii) são o determinante e o traço de sua matriz, deduzidas no capítulo & \cite{WilsonCowan1972} & teoria \\
11 & Com $w_{EE}=1.5$ e $w_{EI}w_{IE}=2$, a inibição rápida ($\tau_I=2\ms$) mantém $E^*=0.67h$, e a lenta ($30\ms$) provoca oscilações (fig.~\ref{fig:ei}) & Solução numérica de~\eqref{eq:ei}, script \texttt{figures/make\_ei.py} & cálculo & modelo \\
12 & O córtex trabalha no regime de estabilização pela inibição; a assinatura: empurre os neurônios \nt{GABA}, e sua frequência cai~\eqref{eq:isn} & A teoria previu a resposta paradoxal. Experimento: registros intracelulares no córtex visual primário do gato; quando a resposta é suprimida por um estímulo ao redor, a condutância inibitória não cresce, mas cai junto com a excitatória. A excitação optogenética de neurônios PV do camundongo nos córtices visual, somatossensorial e motor reduz a frequência dos neurônios inibitórios, inclusive sem estímulo e sob anestesia leve. O coeficiente $-1/3$ com os números da fig.~\ref{fig:ei} é dedução do capítulo & \cite{Tsodyks1997isn,Ozeki2009,Sanzeni2020} & medido \\
13 & O laço ``\nt{GLUT} excita \nt{GABA}, \nt{GABA} inibe \nt{GLUT}'' gera um ritmo rápido com período de $12$--$50\ms$ & A estimulação optogenética de neurônios \nt{GABA} rápidos no córtex somatossensorial do camundongo a $8$--$200$\,Hz amplifica seletivamente o ritmo gama ($20$--$80$\,Hz, período de $12$--$50\ms$); a estimulação de neurônios \nt{GLUT} amplifica só ritmos mais lentos & \cite{Cardin2009,Buzsaki2012} & medido \\
\bottomrule
\end{longtable}}

\section{Capítulo~\ref{sec:code}. Código Morse}

As estimativas-limite do capítulo são teoria da informação e contas; foram medidos onde surge o ruído no canal, quão rápido o cérebro trabalha e quanto custa um disparo, e a questão de que código o córtex usa de fato continua em aberto.

{\footnotesize
\setlength{\tabcolsep}{3pt}\renewcommand{\arraystretch}{1.15}
\begin{longtable}{>{\raggedright}p{0.5cm}>{\raggedright}p{4.0cm}>{\raggedright}p{5.6cm}>{\raggedright}p{2.3cm}>{\raggedright\arraybackslash}p{1.7cm}}
\toprule
N.º & Proposição & Experimento e o que foi medido & Fonte & Status \\
\midrule\endfirsthead
\toprule
N.º & Proposição & Experimento e o que foi medido & Fonte & Status \\
\midrule\endhead
1 & As frequências dos neurônios \nt{GLUT} do córtex vão de frações a dezenas de disparos por segundo, e a maior parte do tempo os neurônios trabalham perto do limite inferior & Registro com pipeta encostada à célula (a escolha do neurônio não depende de sua atividade) no córtex auditivo do rato acordado: a cada momento, menos de $5\,\%$ dos neurônios têm frequência acima de $20$ disparos por segundo; parte dos neurônios tem frequência de fundo abaixo de $0.01$ por segundo; a distribuição das frequências é lognormal & \cite{Hromadka2008} & medido \\
2 & A capacidade do canal de disparos é $R_{\max}\approx r\log_2\frac{e}{r\Delta t}$~\eqref{eq:spikeentropy}, e quase toda ela está no tempo dos disparos (proposição~\ref{prop:capacity}) & Entropia de uma cadeia binária de células de comprimento $\Delta t$. Números do capítulo: $8$ bits por disparo com $r=10$ por segundo e $\Delta t=1\ms$, cerca de $5$ bits com $\Delta t=10\ms$, menos de $2$ bits por segundo para um contador de disparos & \cite{MacKay1952} & teoria \\
3 & Os neurônios sensoriais transmitem de um a alguns bits por disparo, nos melhores casos até metade do limite & Reconstrução do estímulo a partir dos disparos de neurônios sensoriais para os quais o estímulo é conhecido; resumo das medições no livro & \cite{Rieke1997} & medido \\
4 & O gerador de disparos é preciso; o tempo exato é definido pelas variações rápidas da entrada & Neurônios do córtex de rato em fatias: com uma corrente repetida com flutuações rápidas, os disparos se repetem com precisão melhor que $1\ms$; com corrente constante, os instantes se espalham & \cite{Mainen1995} & medido \\
5 & A sinapse não é confiável: mais da metade dos disparos não chega ao destinatário por causa do acaso da liberação & Estimulação mínima das entradas de neurônios piramidais do hipocampo (fatias): mais da metade dos disparos não provoca resposta. Foram excluídas flutuações do limiar de estimulação, falhas de condução nas ramificações do axônio e a baixa temperatura do experimento; resta a liberação probabilística & \cite{AllenStevens1994} & medido \\
6 & O fluxo de disparos no córtex vivo parece aleatório: os intervalos se espalham como num fluxo de Poisson, e a variância do número de disparos fica próxima da média; parte do ``ruído'' é mensagem sobre os movimentos do animal & Registros nas áreas visuais V1 e MT do macaco acordado: coeficiente de variação dos intervalos de cerca de $1$, razão entre variância e média do número de disparos como num processo aleatório; um modelo que soma pequenas entradas independentes dá uma dispersão muito menor. No córtex visual do camundongo, uma parte considerável da variação é prevista pelo vídeo dos movimentos do próprio animal & \cite{Softky1993,Stringer2019} & medido \\
7 & O código de taxa é lento: erro $1/\sqrt{rT}$~\eqref{eq:rateerror}, enquanto o cérebro reconhece uma imagem em $\sim150\ms$ & A fórmula é estatística de Poisson, dedução do capítulo. Experimento: pessoas decidiam se havia um animal numa fotografia desconhecida mostrada por $20\ms$; os potenciais cerebrais para ``sim'' e ``não'' divergem após cerca de $150\ms$. A divisão desse tempo em uma dezena de estágios de $10$--$20\ms$ é estimativa do capítulo, não medição & \cite{Thorpe1996} & medido \\
8 & A média sobre um grupo é limitada pela correlação: o erro cai no máximo $1/\sqrt c$ vezes~\eqref{eq:correlated} (proposição~\ref{prop:pooling}) & Pares de neurônios da área MT do macaco registrados ao mesmo tempo: coeficiente de correlação dos números de disparos de $0.12$ em média, e a análise teórica mostra que mesmo essa correlação limita o ganho do grupo. Teoria: o limite é imposto só pela parte da correlação que se parece com o sinal. Modelo da posição oposta: a frequência de um grupo de $50$--$100$ neurônios é lida num único intervalo entre disparos, $10$--$50\ms$, e o tempo de disparos individuais não é usado no córtex & \cite{Zohary1994,MorenoBote2014,ShadlenNewsome1998} & medido \\
9 & O número pode ser escrito no tempo do primeiro disparo~\eqref{eq:latency}, na ordem dos disparos de um grupo ou na fase de um ritmo local & Retina da salamandra: a estrutura espacial de uma imagem mostrada brevemente está escrita nas latências relativas dos primeiros disparos dos neurônios de saída, quase independentemente do contraste. Células de lugar do hipocampo do rato: ao passar pelo ``seu'' lugar, os disparos se deslocam cada vez mais cedo no ciclo do ritmo teta ($7$--$12$\,Hz), um deslocamento de $100^\circ$ a $355^\circ$. Quanto o córtex usa o tempo dos disparos é questão em aberto (linha~8) & \cite{Gollisch2008,OKeefe1993} & medido \\
10 & Qual código é lido é decidido pela constante de tempo do destinatário~\eqref{eq:reader} (proposição~\ref{prop:reader}); no córtex ativo, os neurônios estão mais perto de detectores de coincidência & A fórmula~\eqref{eq:reader} é dedução do capítulo. Revisão de registros in vivo: no córtex ativo, a condutância da membrana é várias vezes maior que em repouso, e a constante de tempo, correspondentemente mais curta. Que o córtex troca de código exatamente assim não foi mostrado diretamente & \cite{Destexhe2003} & hipótese \\
11 & A sinapse com depressão transmite a mudança de frequência, essencialmente $\Delta r/r$~\eqref{eq:depression} (fig.~\ref{fig:depression}) & Registros pareados de neurônios piramidais do córtex: a velocidade da depressão é dada pela probabilidade de liberação; com depressão lenta, a resposta reflete a frequência; com rápida, as coincidências. Modelo baseado na depressão medida no córtex de rato: mudanças relativas iguais de frequência em entradas rápidas e lentas dão a mesma resposta, isto é, a sinapse transmite $\Delta r/r$. Os números $1.7\to2.2$, $6.7$ e $90\ms$ são cálculo do capítulo & \cite{Tsodyks1997,Abbott1997depression} & medido \\
12 & A rajada é o ``traço'': sua probabilidade de se perder, $(1-p)^k$~\eqref{eq:burst}, é pequena, e a facilitação a reduz ainda mais & Revisão: em sinapses pouco confiáveis, disparos isolados se perdem com frequência, e as rajadas são transmitidas de forma confiável graças à facilitação; a rajada provoca mudanças duradouras nas sinapses. Que o cérebro lê a rajada como um sinal à parte é hipótese & \cite{Lisman1997,AllenStevens1994} & hipótese \\
13 & Os neurônios \nt{GABA} rápidos definem o ritmo e a ``pontuação''; outra classe inibe os inibidores e abre a porta (proposição~\ref{prop:punctuation}) & Revisão das propriedades das classes de neurônios \nt{GABA} do córtex. Optogenética: a excitação rítmica de neurônios \nt{GABA} rápidos provoca o ritmo gama, e a resposta a um estímulo depende da fase do ciclo. No córtex visual do camundongo, os neurônios PV inibem uns aos outros, os SST inibem todas as outras classes, e os VIP inibem sobretudo os SST. A excitação dos neurônios VIP no camundongo acordado remove a inibição dos neurônios \nt{GLUT}; eles são ligados por um sinal de reforço. A divisão de papéis como regra geral é hipótese & \cite{Tremblay2016,Cardin2009,Pfeffer2013,Pi2013} & hipótese \\
14 & A energia limita a frequência média a algo da ordem de um disparo por segundo~\eqref{eq:energy} & Orçamento energético da substância cinzenta de roedores a partir de dados anatômicos e fisiológicos: disparos e correntes sinápticas respondem por $47\,\%$ e $34\,\%$ da energia de sinalização; no máximo $\sim15\,\%$ dos neurônios podem estar ativos ao mesmo tempo. Para o córtex humano: no máximo $\sim1\,\%$ dos neurônios podem estar fortemente ativos ao mesmo tempo. O número de $\sim10^9$ ATP por disparo e a potência de $\sim10$\,W para sinalização são estimativa do capítulo com base nesses cálculos & \cite{Attwell2001,Lennie2003} & teoria \\
15 & O código é esparso e ajustado para o máximo de bits por joule; o córtex troca precisão por energia (proposição~\ref{prop:economy}) & Hipótese do código econômico. Base: em camundongos com restrição alimentar, no córtex visual caem a condutância dos receptores AMPA e o gasto sináptico de ATP ($-29\,\%$), a frequência de disparos se mantém, e a sintonia à orientação se alarga $32\,\%$ & \cite{Barlow1961,Padamsey2022} & hipótese \\
\bottomrule
\end{longtable}}

\section{Capítulo~\ref{sec:serocomputer}. Neurônios \nt{SERO}}

As propriedades dos próprios neurônios \nt{SERO} (ritmo, autoinibição, sinal definido pelo receptor, subsistemas) foram medidas diretamente; os cálculos do capítulo (regulador, filtro, lógica) decorrem de suas equações; $\tau_c$, o coeficiente de difusão da serotonina e o número de locais de liberação são estimativas; o papel do laço como regulador de nível no cérebro é uma hipótese (no núcleo da rafe a autoinibição é pontual e dá só pausas curtas), e o papel do \nt{SERO} como horizonte é uma hipótese com experimentos comportamentais a seu favor.

{\footnotesize
\setlength{\tabcolsep}{3pt}\renewcommand{\arraystretch}{1.15}
\begin{longtable}{>{\raggedright}p{0.5cm}>{\raggedright}p{4.0cm}>{\raggedright}p{5.6cm}>{\raggedright}p{2.3cm}>{\raggedright\arraybackslash}p{1.7cm}}
\toprule
N.º & Proposição & Experimento e o que foi medido & Fonte & Status \\
\midrule\endfirsthead
\toprule
N.º & Proposição & Experimento e o que foi medido & Fonte & Status \\
\midrule\endhead
1 & O sinal da ação da serotonina é escolhido pelo receptor do destinatário (proposição~\ref{prop:receptor}) & Os receptores de serotonina se dividem em famílias por farmacologia e mecanismo: o 5-HT$_{1A}$, via proteína G, abre canais de potássio e inibe a célula; o 5-HT$_{2A}$ e o ionotrópico 5-HT$_3$ excitam. Um mesmo neurotransmissor produz respostas opostas em células diferentes. & \cite{BarnesSharp1999} & medido \\
2 & Na vigília, o neurônio \nt{SERO} emite regularmente alguns disparos por segundo & Registro de neurônios do núcleo dorsal da rafe em gatos em livre movimento: descarga rítmica lenta de $2.8\pm0.2$ disparos/s na vigília tranquila; no sono de ondas lentas a taxa cai $34$--$68\,\%$, no sono REM, $98\,\%$. Em ratos anestesiados, $1.7\pm0.5$ disparos/s, muito regular. & \cite{TrulsonJacobs1979, GagnonParent2014, JacobsAzmitia1992} & medido \\
3 & O neurônio \nt{SERO} é um marca-passo: não precisa de empurrão a cada disparo, mas precisa de um aporte de fundo constante (na fatia, noradrenalina via receptores $\alpha_1$; no organismo, do \nt{NORA}) & Na fatia de cérebro, as células descarregam de forma lenta e regular; após cada disparo há uma pós-hiperpolarização de $200$--$800\ms$. Na fatia há menos células com atividade espontânea que no organismo: o ritmo regular precisa de uma entrada tônica de noradrenalina via receptores $\alpha_1$, cujo bloqueio o apaga. & \cite{VandermaelenAghajanian1983} & medido \\
4 & Quase todos os receptores são metabotrópicos; a resposta é lenta: sobe e desce em cerca de um segundo & Todas as famílias de receptores, exceto a 5-HT$_3$, são acopladas a proteínas G. Na fatia, a liberação evocada de serotonina produz, via 5-HT$_{1A}$, uma corrente inibitória que sobe e desce em um segundo. & \cite{BarnesSharp1999, CourtneyFord2016} & medido \\
5 & Nas regiões de destino, a serotonina sai no volume, e não só na fenda; no próprio núcleo da rafe, a inibição dos vizinhos é ponto a ponto & Voltametria cíclica rápida em fatias: a liberação por disparo é igual numa região com sinapses e no núcleo da rafe, onde quase não há sinapses; não depende do bloqueio da recaptação e dos receptores; há muito menos moléculas por terminal que transportadores e receptores, e a concentração extracelular coincide com a afinidade do receptor principal. No núcleo da rafe, a inibição via 5-HT$_{1A}$ é ponto a ponto: o transportador não deixa a serotonina acumular fora da fenda. & \cite{Bunin1998, CourtneyFord2016} & medido \\
6 & A concentração em~\eqref{eq:seroneuron} diminui pela recaptação; $\tau_c$ da ordem de um segundo é uma estimativa & Voltametria no cérebro vivo: a serotonina extracelular é limitada sobretudo pela recaptação e pela degradação, não pela síntese. Não há medição direta de $\tau_c$ nos trabalhos citados; a ordem de grandeza é estimativa do capítulo. & \cite{Hashemi2012serotonin} & medido; $\tau_c$: estimativa \\
7 & NÃO e NOU num único marca-passo; completude da lógica \nt{SERO} (proposição~\ref{prop:serocomplete}, fig.~\ref{fig:serogates}) & Decorre de~\eqref{eq:seroneuron}: um marca-passo inibível é um NOU, e o NOU é funcionalmente completo. Não há experimento em que se tenha montado lógica com neurônios \nt{SERO}; a condição de volumes separados não se cumpre no cérebro. & dedução no capítulo & teoria \\
8 & A serotonina inibe o próprio neurônio \nt{SERO} por receptores nele & Em ratos anestesiados, agonistas de 5-HT$_{1A}$ por via intravenosa e iontoforética inibem a descarga dos neurônios da rafe dorsal com a mesma relação dose-resposta da própria serotonina; agonistas de 5-HT$_{1B}$ quase não agem. Na fatia, a serotonina endógena das células vizinhas produz, via 5-HT$_{1A}$, uma corrente e uma pausa curta na descarga, sem mudar a taxa média. & \cite{Sprouse1987, CourtneyFord2016} & medido \\
9 & No modelo, o laço por um volume comum é um regulador de nível: a dispersão e a constante de tempo caem $1+G$ vezes~\eqref{eq:feedback}; que o laço funcione assim no cérebro é hipótese & Fórmula clássica do amplificador com realimentação negativa; deduzida de novo no capítulo. O ganho de laço $G$ do sistema \nt{SERO} não foi medido. Medido: na fatia, a autoinibição é ponto a ponto, dá uma pausa curta e não muda a taxa média. & \cite{Black1934feedback, CourtneyFord2016} & teoria; no cérebro, hipótese \\
10 & A concentração é uma média móvel da taxa, um filtro passa-baixas~\eqref{eq:lowpass}, fig.~\ref{fig:lowpass} & Solução da equação linear~\eqref{eq:seroneuron}; a figura é um cálculo por essa fórmula com $\tau_c=1\,\text{s}$. Não há modelo separado em \texttt{models/}. & dedução no capítulo & teoria \\
11 & A tortuosidade dos espaços torna a difusão de uma molécula pequena cerca de $2.6$ vezes mais lenta & Método da fonte pontual: iontoforese de tetrametilamônio e registro de sua concentração com eletrodo íon-seletivo em fatias e no cérebro vivo. Tortuosidade $\lambda\approx1.6$, ou seja, o $D$ efetivo é $\lambda^2\approx2.6$ vezes menor que o livre; fração de volume extracelular de cerca de $0.2$. O coeficiente de difusão da própria serotonina no tecido não foi medido nesses trabalhos; no capítulo, é uma estimativa. & \cite{Sykova2008} & medido \\
12 & Comprimento de um ``fio'' independente $\ell\sim\sqrt{D\tau}\approx20$~\textmu m~\eqref{eq:diffusionlength}; o número de fios é limitado pelo número dessas regiões (proposição~\ref{prop:volumewires}) & Lei da difusão e substituição das estimativas de $D$ e $\tau$ (linhas~6 e~11); a proposição~\ref{prop:volumewires} é uma consequência. Não há experimento direto que conte os canais independentes da transmissão por volume. & dedução no capítulo & teoria \\
13 & A saída de um neurônio \nt{SERO} se espalha por enormes áreas do cérebro; locais de liberação, por estimativa, $\sim10^5$ & Reconstrução completa de neurônios isolados da rafe dorsal no rato: todos os axônios ascendentes se ramificam muito, o comprimento total do axônio de uma célula chega a $18.7$~cm, e os ramos de uma única célula alcançam o estriado, o córtex pré-frontal e a amígdala. O número de locais de liberação por célula não foi contado diretamente. & \cite{GagnonParent2014} & medido; $10^5$: estimativa \\
14 & Os neurônios \nt{SERO} se dividem em alguns subsistemas com saídas e respostas diferentes & Marcação viral-genética em camundongos: os neurônios que vão para a amígdala e para o córtex frontal quase não se sobrepõem na ramificação, recebem entradas diferentes e respondem de modo oposto a um estímulo aversivo; ativar e desativar cada grupo muda o comportamento de modos diferentes. & \cite{Ren2018} & medido \\
15 & Condição de paciência $D<\tau_\gamma\ln(R/r_0)$~\eqref{eq:patience} com desconto exponencial; com o desconto medido, hiperbólico, $D<\tau_\gamma(R/r_0-1)$ & Decorre do desconto $e^{-D/\tau_\gamma}$ ou $1/(1+D/\tau_\gamma)$; dedução no capítulo. Macacos, escolha com atrasos de segundos: o desconto fica mais perto da hipérbole que da exponencial; a resposta dos neurônios \nt{DOPA} ao sinal antecipador de uma recompensa adiada também cai com o atraso de forma hiperbólica. & dedução no capítulo; \cite{KobayashiSchultz2008} & teoria \\
16 & O \nt{SERO} define o horizonte $\tau_\gamma$ (seção~\ref{sec:horizon}) & A ativação optogenética dos neurônios \nt{SERO} da rafe dorsal em camundongos prolonga a espera por uma recompensa adiada e aumenta a persistência na busca de uma recompensa que ficou rara. Em humanos, a redução da serotonina (dieta sem triptofano) torna mais íngreme o desconto da recompensa adiada. Como a concentração se converte em $\tau_\gamma$ não se sabe. & \cite{Doya2002, Miyazaki2014, Lottem2018, Schweighofer2008} & hipótese \\
\bottomrule
\end{longtable}}

\section{Capítulo~\ref{sec:dofa}. Aprendizado com \nt{DOPA}}

Os mecanismos da sinapse (porções, detector de coincidência, cálcio, janelas de plasticidade) e o comportamento da dopamina como erro de predição foram medidos diretamente; que as regras levam ao gradiente e quanto custa a difusão são teoremas; o resultado do aprendizado de escolha é um cálculo com o modelo do livro; o circuito que calcula o erro é uma hipótese.

{\footnotesize
\setlength{\tabcolsep}{3pt}\renewcommand{\arraystretch}{1.15}
\begin{longtable}{>{\raggedright}p{0.5cm}>{\raggedright}p{4.0cm}>{\raggedright}p{5.6cm}>{\raggedright}p{2.3cm}>{\raggedright\arraybackslash}p{1.7cm}}
\toprule
N.º & Proposição & Experimento e o que foi medido & Fonte & Status \\
\midrule\endfirsthead
\toprule
N.º & Proposição & Experimento e o que foi medido & Fonte & Status \\
\midrule\endhead
1 & A retropropagação do erro, na forma que tem nas redes artificiais, não foi encontrada no cérebro & Não há experimento direto: é uma afirmação de ausência. As revisões analisam o que o algoritmo exige (conexões de retorno que repetem as diretas, fase separada) e que aproximações biológicas foram propostas. & \cite{Lillicrap2020, WhittingtonBogacz2019} & hipótese \\
2 & O peso se decompõe em número de locais de liberação, probabilidade de liberação e resposta a uma porção: $w=N_s\,p\,A_1$~\eqref{eq:quantal} & Junção neuromuscular de rã com liberação reduzida: a resposta do destinatário varia em degraus múltiplos da resposta miniatura espontânea a uma porção; o número de porções por disparo segue a lei de Poisson. & \cite{delCastillo1954} & medido \\
3 & O receptor do destinatário é um detector de coincidência; o cálcio desencadeia a mudança de peso & Patch-clamp de neurônios de camundongo: íons Mg$^{2+}$ bloqueiam o canal aberto pelo glutamato no potencial de repouso e saem com a despolarização. Fatias de hipocampo: um quelante de cálcio no destinatário bloqueia a potenciação de longo prazo, e a liberação de cálcio por luz dentro da célula por si só reforça a transmissão. & \cite{Nowak1984, Malenka1988calcium} & medido \\
4 & Qualquer lado executa a decisão: cresce $A_1$ no destinatário, muda $p$ no remetente & Glutamato liberado por luz sobre uma única espinha causa um crescimento rápido e duradouro da espinha e da corrente por seus receptores; a potenciação é acompanhada da inserção de receptores AMPA. A despolarização do destinatário suprime temporariamente a liberação no remetente; o efeito é levado por endocanabinoides que voltam pela fenda (mostrado em entradas \nt{GABA}). A depressão de curto prazo depende da probabilidade de liberação no remetente. & \cite{Matsuzaki2004, Malinow2002, WilsonNicoll2001, Tsodyks1997} & medido \\
5 & A sinapse se reforça quando remetente e destinatário estão ativos juntos~\eqref{eq:hebb} & Coelho anestesiado: após séries curtas de disparos na via de entrada do hipocampo, a resposta fica reforçada por $30$~min a $10$~h. Em fatia de hipocampo, os disparos do destinatário reforçam só as sinapses ativas no mesmo momento. & \cite{Bliss1973LTP, Kelso1986hebb} & medido \\
6 & A regra~\eqref{eq:oja} encontra o autovetor principal das correlações das entradas (proposição~\ref{prop:oja}) & Teorema; demonstração no capítulo. Não há experimento em que um neurônio tenha aprendido a componente principal de suas entradas; o mecanismo que limita o crescimento dos pesos na sinapse não foi estabelecido. & \cite{Oja1982} & teoria \\
7 & Hebb temporal: janela de cerca de $\pm20\ms$ (fig.~\ref{fig:stdp}); sequências repetidas fazem crescer cadeias & Pares de neurônios de hipocampo em cultura: um disparo do destinatário até $20\ms$ depois do disparo do remetente produz reforço duradouro, até $20\ms$ antes, enfraquecimento; ambos dependem de receptores NMDA. A razão $A_-=0.6A_+$ na figura é ilustrativa. Que assim crescem cadeias na rede não foi medido diretamente. & \cite{BiPoo1998} & medido \\
8 & Traço~\eqref{eq:threefactor}: a dopamina que chega $1$--$2\,\text{s}$ após a atividade conjunta reforça a conexão; mais tarde, não (proposição~\ref{prop:threefactor}) & Fatias de estriado, estimulação optogenética separada das entradas glutamatérgicas e dopaminérgicas: a espinha só cresce se a dopamina chega $0.3$--$2\,\text{s}$ após a entrada glutamatérgica; a janela é dada pela subida e queda rápidas do AMPc. No córtex, o pareamento deixa traços separados para reforço e enfraquecimento, que receptores de monoaminas convertem em mudanças duradouras numa janela da ordem de segundos. & \cite{Yagishita2014, He2015eligibility} & medido \\
9 & Janelas de segundos também existem sem dopamina: o terceiro sinal é uma forte excitação do destinatário & Camundongo vivo, campo CA1: um único evento de platô no destinatário reforça as entradas chegadas segundos antes e depois dele e cria um campo de lugar numa só passagem. & \cite{Bittner2017} & medido \\
10 & O passo médio da regra de três fatores segue o gradiente da recompensa~\eqref{eq:perturbation} (proposição~\ref{prop:gradient}) & Teorema do gradiente para aprendizado por desvios aleatórios; deduzido no capítulo para ruído pequeno. & \cite{Williams1992} & teoria \\
11 & O ruído é busca: a rede aprende com seus próprios desvios aleatórios & Mandarins jovens: desligar o núcleo de saída do circuito dos gânglios da base elimina, de forma abrupta e reversível, a variabilidade do canto. Mandarins adultos: uma interferência é aplicada às execuções de uma sílaba com altura de um lado da média, e as aves deslocam rapidamente a altura para o outro lado, aprendendo com a própria dispersão. & \cite{Olveczky2005, TumerBrainard2007} & medido \\
12 & A escolha entre duas opções é aprendida com $\tau_e=1\,\text{s}$ e $\delta=R-\bar R$; não com $\tau_e=50\ms$; sem a subtração, os pesos crescem sem limite, e com taxa de aprendizado alta a rede pode fixar a pior escolha & \texttt{models/dofa\_learning.py}, $\eta=0.05$, $500$ tentativas, $6$ sementes. $\tau_e=1\,\text{s}$: fração de escolhas de $A$ $0.65$--$0.79\to0.96$--$1.00$, pesos $0.85$--$1.33$. $\tau_e=50\ms$: $0.38$--$0.55$, os pesos não mudam. $\delta=R$: peso do vencedor $860$--$1190$. $\delta=R$, $\eta=0.5$, $3$ sementes: uma se fixou em $B$ ($0.04\to0$). O script imprime os quatro casos; o recálculo coincidiu com o capítulo. & \texttt{models/\allowbreak dofa\_\allowbreak learning.py} & modelo \\
13 & O tempo de aprendizado por um único sinal comum cresce proporcionalmente ao número de neurônios ruidosos (proposição~\ref{prop:broadcastcost}) & Curvas de aprendizado de uma rede linear: a perturbação de neurônios é mais lenta que o gradiente exato tantas vezes quantos neurônios de saída houver; a perturbação de pesos, ainda tantas vezes quantas entradas houver. & \cite{Werfel2005} & teoria \\
14 & As saídas do \nt{DOPA} vão sobretudo às regiões de escolha de ações & Reconstrução completa dos axônios de neurônios \nt{DOPA} nigroestriatais isolados do rato: os ramos de uma célula cobrem $0.45$--$5.7\,\%$ do volume do estriado (em média $2.7\,\%$); fora do estriado quase não há ramos. & \cite{Matsuda2009} & medido \\
15 & O córtex aprende a prever sua entrada, e o erro é calculado no local & Revisão: no córtex sensorial há respostas ao desacordo entre a entrada esperada e a recebida. Que isso seja uma regra geral de aprendizado do córtex não foi demonstrado. & \cite{Keller2018} & hipótese \\
16 & A dopamina se comporta como o erro~\eqref{eq:ctd}: pico, transferência para o sinal antecipador, queda (fig.~\ref{fig:ctdcases}) & Neurônios \nt{DOPA} de macaco: pico com suco inesperado; após o aprendizado, pico no sinal antecipador e nenhuma resposta ao suco prometido; queda quando o suco não veio. A forma contínua~\eqref{eq:ctd} é teoria. & \cite{Schultz1997, Doya2000} & medido \\
17 & Circuito que calcula $\dd V/\dd t$ e subtrai a expectativa & Camundongos, registro e optogenética na área tegmental ventral: os neurônios \nt{DOPA} subtraem a expectativa da recompensa; neurônios \nt{GABA} vizinhos os inibem quando a recompensa é esperada, e excitá-los ou inibi-los muda o erro. Como a derivada é calculada não foi mostrado. & \cite{Eshel2015arithmetic} & hipótese \\
18 & O neurônio \nt{DOPA} é um marca-passo; as quedas são menores que os picos & Em fatia, os neurônios \nt{DOPA} descarregam espontaneamente e de modo muito regular, sobre um potencial marca-passo lento. No macaco, a taxa acompanha quantitativamente o erro positivo e transmite o negativo apenas fracamente. & \cite{GraceOnn1989, Bayer2005} & medido \\
19 & Ator e crítico (fig.~\ref{fig:actorcritic}) & O algoritmo vem da teoria do aprendizado por reforço. Em humanos (fMRI), o estriado ventral se comporta como crítico com e sem escolha; o dorsal, só quando é preciso escolher ações. & \cite{SuttonBarto2018, ODoherty2004actor} & hipótese \\
20 & O sinal de $\delta$ na regra é invertido nos neurônios com a segunda família de receptores & Fatias de estriado: nos neurônios com receptores D1, o pareamento produz reforço, que requer D1; nos neurônios com D2, enfraquecimento, que requer D2. & \cite{Shen2008} & medido \\
\bottomrule
\end{longtable}}

\section{Capítulo~\ref{sec:dofaalt}. Outras teorias do \nt{DOPA}}

Cada um dos quadros ampliados se apoia em suas próprias medições diretas da dopamina (dispersão dos pontos de reversão, respostas ao desagradável e ao novo, crescimento lento, respostas à troca de sabor), mas as fórmulas que as explicam são teorias, e entre duas delas os experimentos por enquanto se contradizem.

{\footnotesize
\setlength{\tabcolsep}{3pt}\renewcommand{\arraystretch}{1.15}
\begin{longtable}{>{\raggedright}p{0.5cm}>{\raggedright}p{4.0cm}>{\raggedright}p{5.6cm}>{\raggedright}p{2.3cm}>{\raggedright\arraybackslash}p{1.7cm}}
\toprule
N.º & Proposição & Experimento e o que foi medido & Fonte & Status \\
\midrule\endfirsthead
\toprule
N.º & Proposição & Experimento e o que foi medido & Fonte & Status \\
\midrule\endhead
1 & A regra assimétrica~\eqref{eq:asymmetric} converge para o ponto~\eqref{eq:expectile}; com $\kappa=1/2$ é a média; para uma gota com probabilidade $1/2$, $V=\kappa$ (fig.~\ref{fig:expectiles}) & O ponto~\eqref{eq:expectile} é um expectil, solução de um problema de mínimos quadrados com pesos diferentes para desvios positivos e negativos; para o exemplo do capítulo, a dedução está no próprio capítulo. & \cite{NeweyPowell1987}; dedução no capítulo & teoria \\
2 & Neurônios \nt{DOPA} diferentes têm pontos de reversão diferentes, ordenados pela assimetria (proposição~\ref{prop:expectile}) & Camundongos, neurônios \nt{DOPA} da área tegmental ventral marcados optogeneticamente, recompensas de tamanhos diferentes: o ponto em que o pico dá lugar à queda difere entre neurônios; nos neurônios com respostas mais assimétricas ele é mais alto; das respostas do grupo se reconstrói a forma da distribuição das recompensas. Que essa seja a função principal da dispersão é hipótese. & \cite{Dabney2020} & medido \\
3 & Parte dos neurônios \nt{DOPA} responde ao desagradável com um pico & Macacos: alguns neurônios \nt{DOPA} são excitados pelo sinal de recompensa e inibidos pelo sinal de um jato de ar no olho; outros são excitados por ambos; os grupos ficam em partes diferentes do mesencéfalo. & \cite{Matsumoto2009, BrombergMartin2010} & medido \\
4 & O módulo do erro ou a novidade define a taxa de aprendizado & Nos trabalhos citados não há experimento direto em que o sinal de importância dos neurônios \nt{DOPA} mude a taxa de aprendizado; a revisão propõe o papel de sinal ``atenção, é importante''. & \cite{BrombergMartin2010} & hipótese \\
5 & Um fio separado para o eixo do perigo & Camundongos: os neurônios \nt{DOPA} que vão à extremidade posterior do estriado respondem a estímulos novos e fortes, e não à recompensa; sua destruição enfraquece a evitação de um objeto novo. & \cite{Menegas2018} & medido \\
6 & Tempo ótimo de ação $\tau_a^*=\sqrt{C/\bar R}$~\eqref{eq:vigor} & Mínimo da soma $C/\tau_a+\bar R\tau_a$; dedução no capítulo. No modelo original, a velocidade das ações é dada pelo preço do tempo, igual à recompensa média. & \cite{Niv2007}; dedução no capítulo & teoria \\
7 & O nível lento de dopamina leva $\bar R$ e define a velocidade das ações (proposição~\ref{prop:multiplex}) & Ratos, microdiálise e voltametria no núcleo accumbens: o nível de dopamina, minuto a minuto, acompanha a frequência das recompensas e a velocidade das ações. Em humanos, a L-DOPA reforça a dependência do tempo de reação em relação à frequência média de recompensas. Que o nível lento seja exatamente $\bar R$ não foi demonstrado. & \cite{Hamid2016, Beierholm2013vigor} & hipótese \\
8 & O nível lento vem de duas fontes: a liberação muda nas saídas sem mudança na taxa de disparos, mas o crescimento lento também aparece na própria taxa & Ratos, uma mesma tarefa: a liberação no núcleo accumbens acompanha a expectativa de recompensa em mudança, mas a taxa de disparos dos neurônios \nt{DOPA} identificados da área tegmental ventral não. Camundongos num corredor virtual (linha~10): o crescimento suave ao se aproximar da recompensa aparece também nos disparos dos neurônios \nt{DOPA} identificados. & \cite{Mohebi2019, Kim2020} & medido \\
9 & A dopamina sobe suavemente enquanto o animal se aproxima de uma recompensa distante & Ratos num labirinto, voltametria no estriado: a concentração de dopamina sobe em segundos à medida que o alvo se aproxima; a inclinação depende da distância e do tamanho da recompensa. & \cite{Howe2013ramp, Hamid2016} & medido \\
10 & O crescimento é o erro $\delta$ com refinamento da estimativa~\eqref{eq:ramp}; um salto no corredor dá um pico (fig.~\ref{fig:ramp}) & A fórmula e o número $\delta=0.75\,V$ por segundo são dedução do capítulo. Experimento: camundongos num corredor virtual, transporte instantâneo para mais perto do alvo; os disparos dos corpos celulares, o cálcio dos corpos e das saídas e a concentração no estriado respondem como erro, e não como a expectativa $V$. Qual das duas explicações vale em todas as saídas está em discussão. & \cite{Kim2020} & medido \\
11 & Olhar para trás: a dopamina informa uma medida de ligação causal~\eqref{eq:retro} & Camundongos, liberação de dopamina no núcleo accumbens em experimentos em que essa teoria e o erro~\eqref{eq:ctd} preveem coisas diferentes: em vários experimentos a liberação seguiu a primeira. & \cite{Jeong2022} & hipótese \\
12 & A resposta da dopamina durante o aprendizado se desloca gradualmente da recompensa para o sinal antecipador & Camundongos, sinal das saídas dos neurônios \nt{DOPA} no estriado ventral ao longo do aprendizado: o pico passa gradualmente do instante da recompensa para o do sinal antecipador, como exige o aprendizado por~\eqref{eq:ctd}. & \cite{Amo2022} & medido \\
13 & Os neurônios \nt{DOPA} respondem à troca do sabor da recompensa com o mesmo valor & Ratos, registro de neurônios da área tegmental ventral: trocar um suco por outro de valor igual provoca resposta, embora o erro~\eqref{eq:ctd} seja zero. & \cite{Takahashi2017} & medido \\
14 & Picos artificiais ensinam a ligação entre dois estímulos neutros & Ratos, experimento de pré-condicionamento sensorial de dois estímulos sem recompensa: a ativação optogenética dos neurônios \nt{DOPA} na transição do primeiro para o segundo cria a ligação, e a desativação impede que seja aprendida. & \cite{Sharpe2017} & medido \\
15 & Vetor de erros por característica~\eqref{eq:vectortd} & Teoria do erro de predição de características; não há verificação direta da forma vetorial. & \cite{Gardner2018} & hipótese \\
16 & Neurônios \nt{DOPA} diferentes numa mesma tarefa levam grandezas diferentes & Camundongos num corredor virtual, imagem de dois fótons de neurônios \nt{DOPA}: uns ligados à recompensa, outros à velocidade e à rotação, outros aos sinais antecipadores e à precisão da escolha. & \cite{Engelhard2019} & medido \\
17 & Um feixe em vez de um fio (proposição~\ref{prop:bundle}) & Resumo das linhas~2--16: cada decomposição foi medida separadamente; não há experimento que mostre um quadro único em que todas sejam partes de um mesmo sinal. & --- & hipótese \\
\bottomrule
\end{longtable}}

\section{Capítulo~\ref{sec:nora}. \nt{NORA}}

A organização do sistema \nt{NORA}, seus dois modos, a ligação com a pupila (não só via \nt{NORA}), o aumento da relação sinal-fundo e o U invertido no comportamento foram medidos diretamente; o ganho multiplicativo $g$ é um modelo; que o ganho comute o modo de escolha e funcione como inverso da temperatura são deduções do capítulo; a vantagem do ajuste do ganho é um cálculo com o modelo do livro; ``o \nt{NORA} é o botão da busca'' e ``o \nt{NORA} é uma interrupção'' são hipóteses.

{\footnotesize
\setlength{\tabcolsep}{3pt}\renewcommand{\arraystretch}{1.15}
\begin{longtable}{>{\raggedright}p{0.5cm}>{\raggedright}p{4.0cm}>{\raggedright}p{5.6cm}>{\raggedright}p{2.3cm}>{\raggedright\arraybackslash}p{1.7cm}}
\toprule
N.º & Proposição & Experimento e o que foi medido & Fonte & Status \\
\midrule\endfirsthead
\toprule
N.º & Proposição & Experimento e o que foi medido & Fonte & Status \\
\midrule\endhead
1 & Os neurônios \nt{NORA} no ser humano são algumas dezenas de milhares, quase todos num único grupo & Cortes seriados do tronco encefálico humano, coloração para tirosina hidroxilase: $53\,900$ células no locus coeruleus e mais $6\,260$ num subnúcleo vizinho. & \cite{Baker1989LC} & medido \\
2 & As saídas se espalham por quase todo o cérebro, mas partes do grupo atendem em parte regiões diferentes & Marcação viral-genética em camundongos: os neurônios \nt{NORA} do locus coeruleus recebem entradas de amplas regiões do cérebro e enviam saídas a amplas regiões do cérebro e da medula espinhal. Em ratos: o grupo que vai à amígdala reforça o aprendizado do medo, e um grupo separado que vai ao córtex pré-frontal, sua extinção. & \cite{Schwarz2015LC, Uematsu2017LC, Poe2020} & medido \\
3 & Modo de fundo: disparos regulares, com taxa de frações a alguns por segundo, que muda com a vigília & Ratos, registro de neurônios do locus coeruleus no ciclo sono-vigília: taxa máxima na vigília, menor no sono de ondas lentas, quase zero no sono REM; as mudanças de taxa antecedem a mudança de estado. & \cite{AstonJonesBloom1981} & medido \\
4 & Modo de rajadas: rajada curta em resposta a um evento importante para a tarefa, mais ou menos no instante da decisão & Macacos, tarefa de encontrar um alvo raro: todos os neurônios do locus coeruleus respondem com rajada só ao alvo, $\sim90\ms$ depois dele e $\sim200\ms$ antes da resposta com a mão; com mau desempenho, as rajadas são mais fracas. Numa tarefa de escolha entre duas opções, a rajada está mais presa à resposta que ao estímulo e ocorre tanto em respostas certas quanto erradas. & \cite{AstonJones1994vigilance, Clayton2004LC} & medido \\
5 & O diâmetro da pupila é um ponto de teste impreciso do \nt{NORA}: ele segue o \nt{NORA}, o \nt{ACET} e outras regiões & Macacos: os disparos do locus coeruleus, até disparos isolados de uma célula, preveem a dilatação da pupila; uma ligação parecida, porém mais tardia, existe nos colículos e no córtex cingulado. Camundongos: as flutuações da pupila acompanham a atividade das saídas noradrenérgicas e colinérgicas no córtex. & \cite{Joshi2016, Reimer2016} & medido \\
6 & A noradrenalina suprime a atividade de fundo mais que a evocada; o ganho multiplicativo~\eqref{eq:gain} é um modelo que reproduz esse efeito & Macacos acordados, córtex auditivo, iontoforese de noradrenalina: a atividade de fundo dos neurônios é mais suprimida que a resposta às vocalizações de coespecíficos, e a relação sinal-ruído cresce. A sigmoide multiplicativa~\eqref{eq:gain} foi tirada de um modelo de rede; a multiplicação literal da inclinação não foi medida. & \cite{Foote1975NE, ServanSchreiber1990} & medido; $g$: modelo \\
7 & O ganho eleva $d'$ só contra o ruído depois do amplificador~\eqref{eq:gainsnr} (proposição~\ref{prop:gainnoise}) & Dedução no capítulo; no modelo de rede, o ganho melhora a detecção do sinal. Não há experimento que separe o ruído antes e depois do amplificador e teste essa diferença. & dedução no capítulo; \cite{ServanSchreiber1990} & teoria \\
8 & A fronteira $g\beta=1$ comuta a rede de escolha de ``pondera'' para ``decidiu''~\eqref{eq:wtagain} (proposição~\ref{prop:gainwta}, fig.~\ref{fig:pitchfork}) & Análise linear de dois grupos que se inibem mutuamente; dedução no capítulo. Não há medições diretas desse limiar no cérebro. & dedução no capítulo & teoria \\
9 & O ganho é o inverso da temperatura da escolha~\eqref{eq:softmax} & Com ruído logístico depois do amplificador, a probabilidade de escolha é um softmax com $T=\sigma/g$; dedução no capítulo. & dedução no capítulo & teoria \\
10 & O \nt{NORA} define quanto buscar: atividade de fundo alta, $g$ efetivo pequeno (busca); rajada no instante da decisão, $g$ grande (decisão) & Humanos: antes de uma escolha ao acaso, a pupila basal é mais larga que antes de escolher a melhor opção conhecida; pessoas com pupila mais larga tendem a buscar mais. Ratos: a passagem para o modo de escolha aleatória é controlada pela entrada do \nt{NORA} no córtex cingulado anterior. Que a rajada eleve o $g$ efetivo não foi medido diretamente. & \cite{JepmaNieuwenhuis2011, Tervo2014stochastic, AstonJones2005} & hipótese \\
11 & A recompensa depende de $g$ como um U invertido; o melhor $g$ é menor num mundo que muda depressa (proposição~\ref{prop:explore}, fig.~\ref{fig:invertedU}) & \texttt{models/nora\_explore.py}, $60$ execuções de $4000$ tentativas. Mudança a cada $1000$ tentativas: melhor $g=16$, fração $0.976$; a cada $20$: $g=8$, fração $0.886$; com $g=0.5$, $0.77$. O recálculo coincidiu com o capítulo. & \texttt{models/\allowbreak nora\_\allowbreak explore.py} & modelo \\
12 & U invertido no comportamento: melhor com atividade de fundo moderada e rajadas nítidas & Macacos, mesma tarefa da linha~4: nos períodos de bom desempenho, a taxa de fundo é moderada e as rajadas ao alvo são nítidas; com taxa de fundo alta não há rajadas e há mais respostas falsas. As mudanças da atividade de fundo e evocada acompanham as oscilações do desempenho. & \cite{Usher1999LC, AstonJones2005} & medido \\
13 & O \nt{NORA} informa a incerteza inesperada; o \nt{ACET}, a esperada & Teoria de inferência bayesiana com dois tipos de incerteza; nos trabalhos citados não há experimento direto que separe esses sinais por neurotransmissor. & \cite{YuDayan2005} & hipótese \\
14 & Depois de uma mudança brusca, as pessoas aprendem mais depressa, e a pupila se dilata & Humanos, tarefa de predição com mudanças súbitas: variações curtas da pupila acompanham a probabilidade estimada de mudança e, junto com a pupila basal, preveem o quanto os novos dados mudam a estimativa (a taxa de aprendizado); uma variação da pupila não ligada à luz nem à tarefa muda essa taxa. Que aqui a pupila mostre justamente o \nt{NORA} é uma inferência a partir dos dados indiretos da linha~5. & \cite{Nassar2012} & medido \\
15 & A rajada funciona como interrupção: a rede zera seu estado & Não há experimento direto em que uma rajada do \nt{NORA} zere o estado da rede; só foram medidas rajadas em eventos importantes (linha~4). & \cite{BouretSara2005} & hipótese \\
16 & O ajuste de $g$ ou da taxa de aprendizado pela surpresa quase não supera as melhores configurações constantes & \texttt{models/nora\_explore.py}, mundo com épocas de $1000$ tentativas (mudança a cada $1000$ e a cada $20$). Melhor constante $g=8$: $0.925$; ajuste de $g$: até $0.927$; ajuste da taxa de aprendizado: até $0.926$; taxa constante $0.4$: $0.928$. O recálculo coincidiu com o capítulo. & \texttt{models/\allowbreak nora\_\allowbreak explore.py} & modelo \\
\bottomrule
\end{longtable}}

\section{Capítulo~\ref{sec:acet}. Acrescentamos o \nt{ACET}}

As três ações da acetilcolina foram medidas em fatias e no cérebro vivo, o conflito entre escrita e leitura foi mostrado no modelo do livro, e que o \nt{ACET} sirva de linha de habilitação de escrita e informe a incerteza esperada é uma hipótese com apoio experimental parcial.

{\footnotesize
\setlength{\tabcolsep}{3pt}\renewcommand{\arraystretch}{1.15}
\begin{longtable}{>{\raggedright}p{0.5cm}>{\raggedright}p{4.0cm}>{\raggedright}p{5.6cm}>{\raggedright}p{2.3cm}>{\raggedright\arraybackslash}p{1.7cm}}
\toprule
N.º & Proposição & Experimento e o que foi medido & Fonte & Status \\
\midrule\endfirsthead
\toprule
N.º & Proposição & Experimento e o que foi medido & Fonte & Status \\
\midrule\endhead
1 & Os neurônios \nt{ACET} do cérebro são poucos, estão reunidos em alguns grupos, e as saídas se espalham por todo o córtex: um canal de difusão & Coloração com anticorpos contra a enzima de síntese da acetilcolina e marcação retrógrada no rato: córtex, hipocampo, amígdala, bulbo olfativo e tálamo recebem acetilcolina principalmente de seis grupos compactos de células do prosencéfalo basal e do tronco & \cite{Mesulam1983} & medido \\
2 & O nível de acetilcolina no córtex é alto na vigília e baixo no sono de ondas lentas & Microdiálise no córtex e no hipocampo de gatos em livre movimento, com registro de EEG: a liberação de acetilcolina na vigília tranquila é $\sim100\%$ e na ativa $\sim175\%$ maior que no sono de ondas lentas; no sono REM, no córtex, é como na vigília & \cite{Marrosu1995,Hasselmo1999} & medido \\
3 & O sentido do sinal é definido pelo receptor: a acetilcolina tem receptores-canais rápidos e receptores lentos, que disparam reações na célula (proposição~\ref{prop:receptor}) & Revisão de medições: o efeito da acetilcolina sobre excitabilidade, transmissão e plasticidade depende do local de liberação, do subtipo de receptor (nicotínicos: canais iônicos; muscarínicos: via proteínas G) e da célula destinatária & \cite{Picciotto2012} & medido \\
4 & Primeira ação: enfraquecer as conexões internas, quase sem afetar as entradas externas (fator $1-a$ em~\eqref{eq:acetnet}) & Fatias de córtex olfativo de rato: com $100$\,\textmu M de agonistas da acetilcolina, os potenciais das conexões internas (camada Ib) caem mais de $60\%$, e os da entrada externa (camada Ia), menos de $15\%$, na mesma célula; a supressão é pré-sináptica. Fatias de hipocampo (CA1): $89\%$ contra $40\%$ para as vias interna e de entrada & \cite{HasselmoBower1992,HasselmoSchnell1994} & medido \\
5 & Segunda ação: reforçar a entrada (fator $\frac{1+a}{2}$) & Fatias de neocórtex: receptores nicotínicos reforçam as sinapses da entrada vinda do tálamo e não agem sobre as intracorticais; os muscarínicos enfraquecem ambos os tipos. A acetilcolina aumenta a excitabilidade dos neurônios (revisão) & \cite{Gil1997,Hasselmo2006} & medido \\
6 & Terceira ação: facilitar a mudança dos pesos (taxa $\eta a$) & Rato: a estimulação elétrica do núcleo basal (fonte de acetilcolina para o córtex) pareada com um tom provoca, no animal adulto, uma forte reorganização do mapa do córtex auditivo em favor do som apresentado & \cite{KilgardMerzenich1998,Hasselmo2006} & medido \\
7 & A regra de Hebb para padrões esparsos na segunda equação de~\eqref{eq:acetnet} dá uma memória associativa que completa o padrão a partir de uma parte & Cálculo da capacidade de uma rede com padrões esparsos e regra de covariância: com fração pequena $p$ de neurônios ativos, a rede guarda muito mais padrões que com $p=1/2$ & \cite{Tsodyks1988} & teoria \\
8 & Escrever com $a$ baixo corrompe a memória: a rede grava o lembrado, e não o visto; com $a=0.1$, a corrupção não é menor que com $a=0$ & \texttt{models/\allowbreak acet\_memory.py}: $200$ neurônios, cinco padrões de $20$, um padrão novo compartilha $10$ neurônios com um deles, $10$ execuções. Com $a=0$, o padrão novo não é memorizado, e o antigo arrasta $5.9$ neurônios alheios de $10$; com $a=0.1$, o novo é lembrado ($0.97$), mas os dois se fundem: o novo arrasta $7.3$ neurônios alheios, o antigo $7.9$; a partir de $a=0.2$, menos de um & dedução no capítulo & modelo \\
9 & Proposição~\ref{prop:writeread}: não há nível $a$ em que escrita e leitura sejam igualmente boas (fig.~\ref{fig:acetsim}) & Mesmo script, taxa de aprendizado $\eta a$: o padrão novo é memorizado por inteiro numa apresentação com $a\ge0.9$, em $0.8$ com $a=0.75$, e não é memorizado com $a\le0.5$. Leitura pela metade do padrão: $1.0$ com $a\le0.2$, $0.96$ com $a=0.3$, $0.04$ com $a=0.4$ & dedução no capítulo & modelo \\
10 & No cérebro, um único fio de difusão, o \nt{ACET}, comuta escrita e leitura & A ideia vem de modelos construídos a partir de medições em fatias. O experimento mais direto é em humanos: o bloqueador de receptores muscarínicos escopolamina prejudica a memorização de novos pares de palavras, sobretudo os que se sobrepõem aos antigos, mas não a lembrança de pares já aprendidos. Não há medição direta dos dois modos da rede & \cite{HasselmoAndersonBower1992,Atri2004} & hipótese \\
11 & O filtro~\eqref{eq:kalman} é ótimo; o peso da entrada e a taxa de aprendizado são a mesma grandeza $P/R$; em regime estacionário, $P=\sqrt{qR}$ e a memória é $\sqrt{R/q}$ (proposição~\ref{prop:kalman}) & Não requer experimento: a otimalidade do filtro de Kalman--Bucy é um resultado clássico da teoria de estimação; o regime estacionário é deduzido no capítulo & \cite{KalmanBucy1961} & teoria \\
12 & O \nt{ACET} informa à rede uma grandeza semelhante a $P$, a incerteza esperada & Proposto num modelo probabilístico de atenção. Não há medição direta de que o nível de acetilcolina acompanhe a incerteza da estimativa & \cite{YuDayan2005} & hipótese \\
13 & Parte dos neurônios \nt{ACET} responde à recompensa e à punição em duas dezenas de milissegundos, tanto mais forte quanto mais inesperado o reforço & Camundongo, neurônios colinérgicos do prosencéfalo basal identificados optogeneticamente: resposta à recompensa e à punição após $18\pm3$\,ms, cuja magnitude cresce com a surpresa do reforço; as respostas são semelhantes nos neurônios de dois núcleos & \cite{Hangya2015} & medido \\
14 & O \nt{NORA} percebe uma mudança inesperada do mundo; o \nt{ACET} mantém a rede no modo de escrita & A divisão de trabalho foi proposta num modelo. Indiretamente: em humanos, o diâmetro da pupila, ligado ao \nt{NORA}, acompanha as mudanças e a taxa de aprendizado. Não há teste direto do par \nt{NORA}--\nt{ACET} & \cite{YuDayan2005,Nassar2012} & hipótese \\
15 & No sono de ondas lentas, a rede em modo de leitura reproduz o que foi gravado de dia, e essas repetições o transcrevem para a memória de longo prazo & Registros de conjuntos de neurônios do hipocampo de rato: células que funcionaram juntas de dia disparam juntas com mais frequência no sono de ondas lentas seguinte, e na mesma ordem: medido. Quanto à transcrição: suprimir as ondulações (ripples) do hipocampo após o aprendizado prejudica a memória espacial, e prolongar as ondulações espontâneas a melhora; que a causa seja o nível baixo de \nt{ACET} não foi demonstrado & \cite{WilsonMcNaughton1994,Skaggs1996,Girardeau2009,FernandezRuiz2019,Hasselmo1999} & hipótese \\
\bottomrule
\end{longtable}}

\section{Capítulo~\ref{sec:synthesis}. A máquina inteira}

O capítulo de síntese não traz experimentos novos: cada linha dele se apoia no capítulo em que a afirmação foi analisada e nas mesmas fontes; o barramento \nt{GLUT} e o controle de fluxo por \nt{GABA} estão firmemente estabelecidos, enquanto os papéis dos canais de difusão e o seu funcionamento conjunto são, em grande parte, hipótese.

{\footnotesize
\setlength{\tabcolsep}{3pt}\renewcommand{\arraystretch}{1.15}
\begin{longtable}{>{\raggedright}p{0.5cm}>{\raggedright}p{4.0cm}>{\raggedright}p{5.6cm}>{\raggedright}p{2.3cm}>{\raggedright\arraybackslash}p{1.7cm}}
\toprule
N.º & Afirmação & Experimento e o que foi medido & Fonte & Status \\
\midrule\endfirsthead
\toprule
N.º & Afirmação & Experimento e o que foi medido & Fonte & Status \\
\midrule\endhead
1 & Para $8.6\cdot10^{10}$ neurônios há apenas quatro tipos de sinais de controle~\eqref{eq:machine} & Contagem de núcleos celulares num homogeneizado de cérebro (fracionador isotrópico): um homem adulto tem $86.1\pm8.1$ bilhões de neurônios & \cite{Azevedo2009} & medido \\
2 & Proposição~\ref{prop:architecture}, as duas primeiras camadas: os dados correm pelo barramento de endereços \nt{GLUT}, o sinal da conexão é definido pelo emissor, o fluxo é dirigido e normalizado pelos neurônios \nt{GABA} (capítulos~\ref{sec:neurontypes}, \ref{sec:gaba}) & Um neurotransmissor principal por neurônio (revisão de medições); a inibição direta estreita a janela em que o neurônio piramidal soma as entradas; a divisão pela atividade total foi medida em muitas regiões & \cite{Strata1999,Pouille2001,Carandini2012} & medido \\
3 & Os elementos não são confiáveis: a sinapse perde a maior parte dos disparos (tab.~\ref{tab:computer}, capítulo~\ref{sec:code}) & Fatias de hipocampo, estimulação mínima: mais da metade dos disparos não provoca resposta em CA1; falhas de limiar e de condução no axônio foram excluídas, a causa é a liberação probabilística do neurotransmissor & \cite{AllenStevens1994} & medido \\
4 & O cérebro funciona com $\sim20$ watts, em média cerca de um disparo por segundo por neurônio (capítulo~\ref{sec:code}); os engenheiros ainda não construíram uma máquina assim & Estimativa~\eqref{eq:energy} a partir dos custos medidos: cerca de $10^9$ moléculas de ATP por disparo, incluindo o trabalho das sinapses (córtex de roedores); uma estimativa detalhada para o córtex dá uma taxa ainda menor. O estado da técnica segundo uma revisão & \cite{Attwell2001,Lennie2003,MehonicKenyon2022} & teoria \\
5 & O aprendizado da maioria das sinapses é o produto de um traço local por um único número comum $\delta$ (segunda equação de~\eqref{eq:machine}, capítulo~\ref{sec:dofa}) & Os neurônios \nt{DOPA} do macaco respondem ao erro de predição de recompensa; em fatias de estriado, a dopamina aumenta a espinha só se chegar $0.3$--$2$\,s depois da entrada glutamatérgica: o traço foi medido & \cite{Schultz1997,Yagishita2014} & medido \\
6 & Um fio de erro próprio para cada saída só existe no circuito de aprendizado motor (capítulo~\ref{sec:motorlearning}) & Cerebelo: a coincidência do sinal da fibra trepadeira com uma entrada enfraquece essa entrada; no macaco, os disparos complexos das células de Purkinje carregam a direção do erro de movimento e provocam a correção & \cite{Ito2001,Herzfeld2018} & medido \\
7 & Cada canal de difusão é um feixe de algumas linhas (proposição~\ref{prop:bundle}) & Para \nt{DOPA}: na mesma tarefa, neurônios diferentes carregam recompensa, movimento e atributos do estímulo; o erro rápido e o nível lento de dopamina são distintos. Para \nt{SERO}, \nt{NORA}, \nt{ACET} essa divisão foi menos estudada & \cite{Engelhard2019,Hamid2016,Mohebi2019} & medido \\
8 & \nt{SERO}: regulador de nível e, por hipótese, o horizonte $\tau_\gamma$ (capítulo~\ref{sec:serocomputer}) & Autoinibição: agonistas dos próprios receptores 5-HT$_{1A}$ inibem os neurônios serotoninérgicos da rafe (medido). O horizonte foi proposto em teoria; o experimento mais forte: a ativação optogenética desses neurônios no camundongo prolonga a espera por uma recompensa adiada & \cite{Sprouse1987,Doya2002,Miyazaki2014} & hipótese \\
9 & \nt{NORA}: o ganho $g$ escolhe entre explorar e decidir (capítulo~\ref{sec:nora}) & Aumento da relação sinal-ruído: no macaco acordado, a noradrenalina no córtex auditivo suprime a atividade de fundo mais do que a resposta a vocalizações de coespecíficos (medido); o ganho multiplicativo é um modelo; o papel na escolha entre explorar e decidir é teoria & \cite{Foote1975NE,ServanSchreiber1990,AstonJones2005} & hipótese \\
10 & \nt{ACET}: alterna entre escrita e leitura (capítulo~\ref{sec:acet}) & Três ações (enfraquecimento das conexões internas, reforço da entrada, facilitação da plasticidade) foram medidas; a troca de modos tem apoio indireto num experimento com escopolamina em humanos & \cite{HasselmoBower1992,Gil1997,Atri2004} & hipótese \\
11 & A fórmula~\eqref{eq:machine} descreve a máquina inteira & É um resumo dos modelos dos capítulos~\ref{sec:glutcomputer}--\ref{sec:acet}, cada parte se apoia no seu capítulo; como um todo, não foi testada experimentalmente & dedução no capítulo & hipótese \\
12 & Os botões trabalham de forma coordenada sem controle comum: erro, surpresa, escrita, retorno (fig.~\ref{fig:timeline}) & Não há experimento: o esquema é qualitativo. Elos isolados: a teoria da divisão de trabalho entre \nt{NORA} e \nt{ACET} e a relação entre pupila e taxa de aprendizado em humanos & \cite{YuDayan2005,Nassar2012} & hipótese \\
13 & O aprendizado por um único sinal comum é tanto mais lento quanto mais neurônios influem no resultado (proposição~\ref{prop:broadcastcost}) & Cálculo das curvas de aprendizado de uma rede linear: o aprendizado por perturbação das saídas é mais lento que a descida de gradiente tantas vezes quantas são as saídas & \cite{Werfel2005} & teoria \\
14 & Não se sabe como o córtex ajusta circuitos de muitas camadas; candidatos: rajadas de disparos, cálculos nos ramos do neurônio, autoaprendizado por predição & Rajadas como portadoras do erro: modelo; reproduzir a resposta de um modelo detalhado de neurônio cortical só uma rede de $5$--$8$ camadas conseguiu: modelo; disparos de cálcio nos ramos de neurônios do córtex humano, que dão o OU exclusivo: medido em fatias; autoaprendizado: revisão de evidências & \cite{Payeur2021,Beniaguev2021,Gidon2020,Keller2018} & hipótese \\
\bottomrule
\end{longtable}}

\section{Capítulo~\ref{sec:sensors}. As entradas da máquina}

A construção dos sensores foi medida diretamente, com registros da corrente de células isoladas, das vibrações da cóclea e com contagens de células da retina; o limite de sensibilidade e a fórmula do ruído balístico decorrem da física, e que os códigos dos sensores estejam ajustados para a máxima informação é uma hipótese com apoio forte.

{\footnotesize
\setlength{\tabcolsep}{3pt}\renewcommand{\arraystretch}{1.15}
\begin{longtable}{>{\raggedright}p{0.5cm}>{\raggedright}p{4.0cm}>{\raggedright}p{5.6cm}>{\raggedright}p{2.3cm}>{\raggedright\arraybackslash}p{1.7cm}}
\toprule
N.º & Proposição & Experimento e o que foi medido & Fonte & Status \\
\midrule\endfirsthead
\toprule
N.º & Proposição & Experimento e o que foi medido & Fonte & Status \\
\midrule\endhead
1 & O sensor é um transdutor: o receptor gera corrente, e a saída das células sensoriais do olho e do ouvido é glutamato para o barramento \nt{GLUT}; as primeiras células não emitem disparos (fig.~\ref{fig:transducer}) & Bastonetes e cones de tartaruga liberam um aminoácido excitatório: ele é captado por canais de glutamato num fragmento de membrana excisado. Células ciliadas internas da cóclea de rato: as correntes nas terminações das fibras nervosas passam por receptores AMPA de glutamato, e a taxa de liberação cresce com a despolarização gradual da célula até $150$\,Hz & \cite{CopenhagenJahr1989,GlowatzkiFuchs2002} & medido \\
2 & O sensor de luz no escuro responde a um único fóton com uma corrente de cerca de um picoampere & Eletrodo de sucção no segmento externo de bastonete de sapo: as respostas a flashes fracos são quantizadas, evento de $\approx1$\,pA ($5\%$ da corrente de escuro) com dispersão de $0.2$\,pA, eficiência $0.5$. Humanos relatam a chegada de um único fóton acima do acaso & \cite{Baylor1979,Tinsley2016} & medido \\
3 & O sensor de som percebe um deslocamento menor que um nanômetro & Método Mössbauer, membrana basilar da cóclea de porquinho-da-índia no ponto de $16$--$19$\,kHz: no limiar de resposta do nervo auditivo, a velocidade da membrana é de cerca de $0.04$\,mm/s, isto é, um deslocamento de $\approx0.35$\,nm (cálculo nosso). Mediu-se a membrana, e não o feixe do sensor & \cite{Sellick1982,Hudspeth1989} & medido \\
4 & No escuro, a precisão é limitada pelo ruído balístico dos fótons~\eqref{eq:photonnoise}; com luz fraca, a acumulação é mais longa & A fórmula decorre da estatística de Poisson. No bastonete, o ruído com luz fraca coincide com a soma de eventos quânticos independentes; sobre um fundo claro, a resposta ao flash é menor e mais curta. O valor ``décimos de segundo'' não foi verificado separadamente aqui & \cite{Baylor1979} & teoria \\
5 & A faixa de entrada é de dez ordens de grandeza para a luz e doze para o som, e a taxa de disparos, menos de três ordens & Para o som: a resposta da membrana basilar na frequência característica cresce com inclinação de até $0.2$\,dB/dB na faixa de $40$--$80$\,dB, e a compressão aparece também acima de $100$\,dB. Os próprios números de ordens são estimativas padrão, sem fonte no capítulo & \cite{Ruggero1997} & medido \\
6 & A resposta do sensor segue~\eqref{eq:nakarushton}, e $\sigma$ acompanha o brilho médio, deslocando a curva ao longo do eixo logarítmico (fig.~\ref{fig:adaptation}) & A fórmula foi ajustada primeiro aos potenciais S da retina de peixe. Cones de tartaruga: as respostas seguem $V=I V_m/(I+\sigma)$ e cobrem $\sim3.5$ ordens de brilho; após $2$--$3$ minutos de fundo, a curva se desloca na horizontal, mantendo a largura. A ligação com a divisão pela atividade total é de uma revisão & \cite{NakaRushton1966,NormannPerlman1979,Carandini2012} & medido \\
7 & Os sensores transmitem mudanças, e os seus códigos estão ajustados para a máxima informação por disparo (proposição~\ref{prop:sensors}) & O princípio foi proposto teoricamente. A evidência mais forte: na mosca, a característica ``contraste $\to$ resposta'' dos neurônios da primeira camada coincide com a distribuição acumulada de contrastes das cenas naturais, e assim se obtém a máxima informação & \cite{Barlow1961,Laughlin1981} & hipótese \\
8 & Os sensores de ligar e de desligar formam um código de dois fios; a câmera de eventos o repete em silício & Revisão de experimentos: os canais de ligar e de desligar da retina se separam já na primeira sinapse e podem ser silenciados separadamente. Câmera: $128\times128$ pixels sem quadros, faixa de $120$\,dB, latência de $15$\,\textmu s & \cite{Schiller1992,Lichtsteiner2008,Gallego2022} & medido \\
9 & O olho tem $\sim10^8$ sensores de luz e $\sim10^6$ neurônios de saída: compressão de $\sim100$ vezes & Contagem em preparações inteiras de retina humana: em média $92$ milhões de bastonetes e $4.6$ milhões de cones; células de saída (ganglionares), de $0.7$ a $1.5$ milhão conforme o olho & \cite{Curcio1990,CurcioAllen1990} & medido \\
10 & O camundongo tem mais de trinta canais de saída do olho & Camundongo: imageamento de cálcio de dois fótons de mais de $11\,000$ células de saída e agrupamento das respostas: bem mais de $30$ tipos funcionais. Para humanos, o número não foi medido & \cite{Baden2016} & medido \\
11 & O olho do porquinho-da-índia transmite $\sim875\,000$ bit/s; o recálculo para o olho humano dá $\sim10^7$ bit/s & Porquinho-da-índia, matriz de múltiplos eletrodos, movimento em cenas naturais: $6$--$13$ bit/s por célula, $\sim875\,000$ bit/s para $\sim10^5$ células de saída. Para humanos ($\sim10^6$ células), $\sim10^7$ bit/s é um recálculo dos autores & \cite{Koch2006} & medido \\
12 & O ouvido é um banco de filtros passa-faixa com um mapa de frequências quase logarítmico (fig.~\ref{fig:filterbank}) & O ponto de máxima vibração da membrana basilar depende da frequência; a função ``frequência--lugar'' é quase exponencial e descreve com uma só forma as cócleas humanas e de animais vivos & \cite{Greenwood1990,Robles2001} & medido \\
13 & Os ressonadores são ativos: parte dos sensores amplifica as vibrações fracas milhares de vezes em amplitude ($66$--$76$\,dB), e o ganho cai com o volume & Vibrometria a laser da base da cóclea de chinchila: com tons fracos na frequência característica, ganho de $66$--$76$\,dB em relação ao estribo; a morte reduz a sensibilidade em $60$--$81$\,dB ($10^3$--$10^4$ vezes em amplitude) e elimina a compressão. A fonte da força são as células ciliadas externas & \cite{Ruggero1997,Robles2001} & medido \\
14 & O amplificador do ouvido está no limite da autoexcitação & Cálculo: perto de um ponto de bifurcação de Hopf são inevitáveis a compressão da faixa, a sintonia aguda e os tons de combinação, isto é, todas as não linearidades conhecidas da audição. Que a cóclea esteja ajustada assim não foi provado diretamente & \cite{Eguiluz2000} & hipótese \\
15 & Nas frequências baixas, os disparos vêm em fase com o som (no porquinho-da-índia, o travamento enfraquece a partir de $\sim600$\,Hz e é perceptível até $\sim3.5$\,kHz), e por eles se acha a direção; nas altas, resta o código de lugar & Nervo auditivo de porquinho-da-índia: o travamento de fase começa a cair em torno de $600$\,Hz e desaparece acima de $3.5$\,kHz; no gato e no macaco o limite fica cerca de uma oitava acima. Diferença de tempo de chegada aos dois ouvidos: revisão & \cite{PalmerRussell1986,Grothe2010} & medido \\
16 & Os demais sensores (tab.~\ref{tab:sensors}): o olfato tem $\sim400$ tipos de receptores, um por neurônio, e código combinatório; o paladar funciona em parte por ATP e serotonina; o tato tem sensores de adaptação rápida e lenta; calor e frio são separados & Camundongo: um receptor reconhece muitos cheiros, um cheiro, muitos receptores, e cheiros diferentes, conjuntos diferentes. Sem receptores de ATP, o nervo do paladar não responde; os botões gustativos liberam serotonina. Registro de fibras isoladas da pele da mão humana: quatro tipos pela velocidade de adaptação. O canal do frio é distinto dos canais do calor & \cite{Mombaerts2004,Malnic1999,Finger2005,Huang2005,JohanssonVallbo1979,McKemy2002} & medido \\
\bottomrule
\end{longtable}}

\section{Capítulo~\ref{sec:recognition}. Reconhecimento}

A velocidade do reconhecimento, a construção dos primeiros estágios, a leitura linear da resposta e o aprendizado pela continuidade do tempo foram medidos em humanos e macacos; a redução de toda a cadeia a duas operações e o aprendizado com vídeo foram testados nos modelos do livro, e as explicações das ilusões vão das bem fundamentadas às controversas.

{\footnotesize
\setlength{\tabcolsep}{3pt}\renewcommand{\arraystretch}{1.15}
\begin{longtable}{>{\raggedright}p{0.5cm}>{\raggedright}p{4.0cm}>{\raggedright}p{5.6cm}>{\raggedright}p{2.3cm}>{\raggedright\arraybackslash}p{1.7cm}}
\toprule
N.º & Afirmação & Experimento e o que foi medido & Fonte & Status \\
\midrule\endfirsthead
\toprule
N.º & Afirmação & Experimento e o que foi medido & Fonte & Status \\
\midrule\endhead
1 & A comparação pixel a pixel com um modelo, sob deslocamento, não acerta mais que uma moeda; o par de estágios~\eqref{eq:scstage} reconhece a figura em qualquer lugar & \texttt{models/\allowbreak recognition\_models.py}, ``retina'' de $24$ pontos, $4000$ tentativas: modelo, $0.49$, $0.51$, $0.52$ com fração de pontos invertidos $0$, $0.1$, $0.2$; par $S$ e $C$, $1.00$, $0.86$, $0.70$ & dedução no capítulo & modelo \\
2 & Um animal numa imagem é reconhecido em $\sim150$\,ms, numa passagem por uma dezena de estágios de $10$--$20$\,ms & Humanos, tarefa ``há um animal ou não'' em fotos mostradas por $20$\,ms: os potenciais evocados das duas respostas divergem após $\sim150$\,ms. Macaco: a latência da primeira resposta cresce de estágio em estágio (V1 primeiro, depois V2 e V4). O número de estágios e ``zero ou um disparo por estágio'' são deduções desses números & \cite{Thorpe1996,Schmolesky1998} & medido \\
3 & O estágio $S$ é um E sobre as partes, o estágio $C$, o máximo sobre as posições (proposição~\ref{prop:conveyor}) & Gato, córtex visual primário: as células simples respondem a uma borda de certa orientação num certo lugar, as complexas, à mesma orientação numa área maior. Registro intracelular de células complexas: a resposta a duas barras, em muitas células, é próxima da maior das duas respostas, em média mais próxima da operação max. Máximo por inibição mútua: proposição~\ref{prop:wta}; divisão pela atividade total: revisão & \cite{HubelWiesel1962,Lampl2004,Carandini2012} & medido \\
4 & Mais acima na cadeia, as combinações são maiores e mais complexas, e a resposta depende cada vez menos da posição & Macaco, simplificação dos estímulos até o ``atributo crítico'': as células de V4 e do córtex inferotemporal posterior respondem a combinações complexas de forma, cor e textura com campo pequeno; no córtex inferotemporal anterior o campo é maior. A alternância de $S$ e $C$ no modelo reproduz esse quadro & \cite{KobatakeTanaka1994,Riesenhuber1999} & medido \\
5 & As redes convolucionais repetem essa alternância e são as que melhor predizem as respostas dos estágios superiores; o objeto é lido linearmente das taxas de um grupo de neurônios & Rede treinada para reconhecer, sem ajuste ao cérebro: a camada superior prediz as respostas dos neurônios do córtex inferotemporal do macaco, as intermediárias, as de V4. Um classificador linear sobre a atividade de $\sim100$ células escolhidas ao acaso, em janelas a partir de $12.5$\,ms, reconhece o objeto e a categoria com posição e tamanho diferentes & \cite{Fukushima1980,Yamins2014,Hung2005} & medido \\
6 & A regra de Hebb com traço~\eqref{eq:traceinvariance} ensina a invariância com vídeo sem rótulos, se o traço não for mais curto que $\sim20$\,ms & A ideia dos atributos lentos e da regra com traço é teoria. \texttt{models/\allowbreak recognition\_models.py}, dois neurônios $C$, $16$ entradas, $10$ rodadas: pausa entre objetos de $0.3$\,s. Com $\tau_{\text{tr}}=5$\,ms, coincidência com a figura de $0.52$ em média, com $10$\,ms, $0.56$; com $20$, $50$ e $200$\,ms, $1.00$ em todas as rodadas (com $30$\,ms, uma rodada em dez erra) & \cite{Foldiak1991,WiskottSejnowski2002} & modelo \\
7 & No cérebro, a invariância é aprendida pela continuidade do tempo & Macaco: o objeto era trocado discretamente durante um salto do olho para um certo lugar do campo; a invariância de posição dos neurônios do córtex inferotemporal mudava de acordo com a troca, de forma significativa já após uma hora. Que essa seja a regra principal do aprendizado visual é hipótese & \cite{LiDiCarlo2008} & medido \\
8 & Movimento: detectores de direção, depois o mesmo par E/OU sobre áreas grandes & Detectores de direção na retina do coelho; na área MT do macaco, todos os $110$ neurônios registrados são seletivos à direção, e a resposta ao movimento é mais forte que em V1 & \cite{Barlow1965,Albright1984} & medido \\
9 & Os estágios superiores enviam para baixo uma predição, e para cima sobe o erro & O modelo explica os efeitos extraclássicos dos campos do córtex visual primário; observações isoladas concordam (revisão); como construção geral da linha de montagem, não está provado & \cite{RaoBallard1999,Keller2018} & hipótese \\
10 & Imagens difíceis exigem realimentação e várias voltas & Macaco: para imagens ``difíceis'', a decisão no córtex inferotemporal se forma $\sim30$\,ms mais tarde; essas respostas tardias são mal preditas por uma rede direta e melhor por redes com realimentação ou mais profundas & \cite{Kar2019} & medido \\
11 & Uma falsificação imperceptível de pixels engana a rede; a visão humana é muito mais robusta, mas não invulnerável & Em redes: uma pequena perturbação escolhida de propósito muda a resposta; o exemplo ``panda $\to$ gibão'' é do segundo trabalho. Em humanos, com apresentação breve, falsificações que se transferem bem entre redes desviam a escolha & \cite{Szegedy2014,Goodfellow2015,Elsayed2018} & medido \\
12 & Ilusões de expectativa: a entrada pouco confiável cede à expectativa; o pálido se move ``mais devagar'', o vestido é visto de formas diferentes (seção~\ref{sec:illusions}) & Grades de menor contraste parecem se mover mais devagar. Enquete com $1401$ pessoas: o vestido é visto em três variantes estáveis, e uma dica sobre a iluminação troca a cor; em $\sim13\,000$ observadores, o que decide é a suposição sobre a iluminação. Um modelo bayesiano com expectativa de movimentos lentos explica várias ilusões de movimento & \cite{Thompson1982,LaferSousa2015,Wallisch2017,Weiss2002,Purves2011} & hipótese \\
13 & Controle de ganho: cachoeira, inclinação e contraste; a grade de Hermann não é inibição dos vizinhos & Os detectores de direção da retina do coelho, com movimento prolongado, reduzem a taxa e, depois da parada, caem abaixo do fundo. A ilusão de inclinação é reproduzida por um modelo com divisão pela atividade dos vizinhos. Mudanças na grade que não alteram a resposta dos neurônios de saída da retina eliminam as manchas: a explicação pela inibição dos vizinhos foi rejeitada & \cite{BarlowHill1963,Schwartz2009,SchillerCarvey2005} & medido \\
14 & Compensação de atrasos: um objeto em movimento parece à frente do clarão & A ilusão foi medida. A psicofísica contradiz tanto o prolongamento do movimento quanto a diferença de latências; foi proposta uma explicação retroativa, pelos eventos $\sim80$\,ms depois do clarão. A disputa não está resolvida & \cite{Nijhawan1994,EaglemanSejnowski2000} & hipótese \\
15 & ``Serpentes'' paradas giram: a diferença de latências~\eqref{eq:latency} é lida pelos detectores de direção & A ilusão funciona também sem cor; pares de elementos vizinhos a geram isoladamente; neurônios do córtex do macaco seletivos à direção dão resposta direcional a esses pares parados. Em humanos, a rotação é disparada por microssacadas e piscadas & \cite{Conway2005,OteroMillan2012} & medido \\
16 & Imagens ambíguas: inibição mútua, fadiga e ruído; as trocas são causadas principalmente pelo ruído & Modelo de grupos de neurônios em competição: nele, as trocas são causadas pelo \emph{ruído}, e sem ruído não existem; ele explica o aumento da frequência das trocas com a força do estímulo. A fadiga tem aqui um papel menor & \cite{MorenoBote2007} & hipótese \\
17 & Parte das ilusões é consequência da tarefa que a máquina aprende & Uma rede treinada para prever o próximo quadro de vídeo ``vê'' rotação em anéis parados e não a vê nas imagens de controle; redes para limpar imagens e aumentar a nitidez caem em ilusões de cor e brilho, como os humanos & \cite{Watanabe2018,GomezVilla2020} & medido \\
\bottomrule
\end{longtable}}

\section{Capítulo~\ref{sec:muscles}. A saída para os músculos}

A construção da saída (unidades motoras, margem de segurança da sinapse, ativação por tamanho, par de músculos, laço de estiramento e geradores de ritmo) foi medida diretamente; a condição de estabilidade do laço com atraso é um teorema; o modelo interno do corpo é uma hipótese bem fundamentada; os números das figuras são cálculos pelo modelo do livro.

{\footnotesize
\setlength{\tabcolsep}{3pt}\renewcommand{\arraystretch}{1.15}
\begin{longtable}{>{\raggedright}p{0.5cm}>{\raggedright}p{4.0cm}>{\raggedright}p{5.6cm}>{\raggedright}p{2.3cm}>{\raggedright\arraybackslash}p{1.7cm}}
\toprule
N.º & Proposição & Experimento e o que foi medido & Fonte & Status \\
\midrule\endfirsthead
\toprule
N.º & Proposição & Experimento e o que foi medido & Fonte & Status \\
\midrule\endhead
1 & Unidade motora: um neurônio \nt{ACET} controla um grupo de fibras, de algumas centenas a um par de milhares; nos músculos de ajuste fino as unidades são menores & Músculos humanos, contagem de axônios motores e de fibras musculares: em média cerca de $340$ fibras por neurônio no músculo interósseo da mão, cerca de $410$ no braquiorradial, cerca de $1930$ no gastrocnêmio medial. O capítulo não dá um número exato para as menores unidades. & \cite{Feinstein1955mu} & medido \\
2 & A sinapse no músculo libera várias vezes mais neurotransmissor do que o necessário para a fibra responder & Revisão de medições em sinapses neuromusculares: a margem de segurança (razão entre o neurotransmissor liberado e o limiar) em mamíferos adultos costuma ser $3$--$5$; em humanos, a margem se apoia mais nas dobras da membrana receptora do que no tamanho da liberação. & \cite{WoodSlater2001} & medido \\
3 & O abalo~\eqref{eq:twitch} sobe num tempo $T_c$ da ordem de $50\ms$ & Unidades motoras isoladas da mão humana em esforço voluntário, força promediada pelos disparos da unidade: tempo de subida de $30$--$100\ms$, abaixo de $70\ms$ em $80\,\%$ das unidades. A função $(t/T_c)e^{1-t/T_c}$ é uma aproximação de um modelo de conjunto de unidades. & \cite{MilnerBrown1973rec, Fuglevand1993} & medido \\
4 & Soma de abalos~\eqref{eq:forcesum}: com $5$, $15$ e $40$ disparos/s a força é $0.2$--$1.1F_1$, $1.8$--$2.2F_1$ e cerca de $5.4F_1$ (fig.~\ref{fig:tetanus}) & Cálculo com \texttt{models/\allowbreak muscle\_\allowbreak models.py}, $T_c=50\ms$: $0.21$--$1.10$, $1.76$--$2.19$ e $5.28$--$5.49\,F_1$ nos últimos $0.3\,\text{s}$. A soma linear é uma simplificação: o modelo não tem a saturação do músculo real. & \texttt{models/\allowbreak muscle\_\allowbreak models.py} & modelo \\
5 & As unidades são ativadas por tamanho; as mais fracas são cem vezes mais fracas que as mais fortes & Raízes ventrais do gato: com entrada reflexa crescente, disparam primeiro os neurônios com impulsos pequenos na raiz, isto é, os pequenos. Músculo interósseo da mão humana: força de abalo das unidades de $0.1$ a $10$~g, crescendo quase linearmente com a força em que a unidade é ativada. & \cite{Henneman1957, MilnerBrown1973rec} & medido \\
6 & O passo de força é proporcional à força, $\Delta F/F\approx q-1$~\eqref{eq:sizeprinciple}: ponto flutuante & Deduzido no capítulo a partir da progressão geométrica das forças. Concorda com o experimento: em esforços grandes, o mesmo incremento de força ativa muito menos unidades novas. & dedução no capítulo; \cite{MilnerBrown1973rec} & teoria \\
7 & A dispersão da força cresce proporcionalmente à própria força & Esforço isométrico voluntário do músculo do polegar: o desvio padrão da força cresce linearmente com a força média; com o mesmo esforço provocado por estimulação elétrica, a dispersão é constante. Modelo de conjunto: o crescimento linear resulta da ativação das unidades por ordem de força. & \cite{Jones2002sdn} & medido \\
8 & A suavidade dos movimentos é consequência do ruído dependente do sinal & Modelo: a trajetória que minimiza a dispersão do ponto final com esse ruído reproduz as trajetórias medidas do olho e do braço. Não há experimento direto que distinga essa causa das outras. & \cite{HarrisWolpert1998} & hipótese \\
9 & A diferença das forças do par de músculos dá o torque, a soma dá a rigidez~\eqref{eq:antagonists} & Teoria do controle de impedância por co-contração. Braço humano: a rigidez de cada articulação, medida com pequenos empurrões, é aproximadamente proporcional ao torque nela; diferentes proporções de co-contração mudam a forma e a direção da elipse de rigidez da mão. & \cite{Hogan1984, GomiOsu1998stiff} & medido \\
10 & O sensor de estiramento excita o seu neurônio \nt{ACET} e, por um intermediário inibitório, desliga o antagonista (fig.~\ref{fig:antagonists}) & Medula espinhal do gato: um único intermediário, excitado pelas fibras dos sensores de estiramento, produz potenciais inibitórios monossinápticos nos neurônios \nt{ACET} do músculo antagonista, atraso sináptico de $0.28$--$0.42\ms$. O neurotransmissor do intermediário (glicina) não foi determinado nesse experimento. & \cite{Jankowska1972Ia} & medido \\
11 & Atraso do laço: espinhal $25$--$30\ms$, pelo córtex uma vez e meia a três vezes mais, pelo olho $100$--$200\ms$ & Braço humano, empurrão na articulação: resposta muscular curta após $20$--$45\ms$, longa após $45$--$105\ms$, voluntária após $120$--$180\ms$. Deslocamento da posição visível do dedo durante o movimento: correção após $160\ms$ em média. & \cite{Pruszynski2008rapid, SaundersKnill2003vis} & medido \\
12 & $\dd x/\dd t=-Gx(t-d)$ é estável com $Gd<\pi/2$~\eqref{eq:delaystable}; no limite, a frequência é $1/(4d)$ & Teorema sobre as raízes de uma equação com retardo. Verificação com \texttt{models/\allowbreak muscle\_\allowbreak models.py}, $d=30\ms$: em $3\,\text{s}$, amplitude $3\cdot10^{-6}$ com $G=40$, $0.13$ com $G=50$, $38$ com $G=55$ por segundo (limite $52$). & \cite{Hayes1950} & teoria \\
13 & Tremor fisiológico de $8$--$12$~Hz; o laço de estiramento é uma das suas fontes & Revisão de medições: um tremor fino na faixa de cerca de $10$~Hz existe em pessoas saudáveis; a sua origem é mista (oscilações centrais, propriedades das descargas das unidades, ressonância mecânica e ressonância do laço reflexo), e em condições diferentes predomina uma ou outra. & \cite{McAuleyMarsden2000} & hipótese \\
14 & O cérebro prevê o resultado do comando por um modelo interno do corpo & Uma pessoa segura um peso em pinça e move o braço: com carga inercial, viscosa e elástica, a força de preensão muda ao mesmo tempo que a carga, e não com o atraso do laço, ou seja, a carga é prevista. Uma revisão reúne esses experimentos; não há observação direta do próprio modelo nos neurônios. & \cite{FlanaganWing1997grip, WolpertGhahramani2000} & hipótese \\
15 & O ritmo de andar, respirar e mastigar é gerado por redes locais com entrada constante & Revisão de experimentos: o sistema nervoso isolado (medula, gânglios) produz um padrão motor rítmico sem entrada rítmica e sem realimentação dos músculos. & \cite{MarderBucher2001} & medido \\
16 & A inibição mútua com fadiga~\eqref{eq:halfcentre} dá um ritmo com período de $0.84\,\text{s}\approx2\tau_a$ (fig.~\ref{fig:halfcentre}) & Teorema: uma rede com inibição mútua e adaptação sem estado estável, com entrada constante, necessariamente oscila. Cálculo com \texttt{models/\allowbreak muscle\_\allowbreak models.py} ($I=1$, $\beta=2$, $b=2$, $\tau=20\ms$, $\tau_a=0.4\,\text{s}$): período de $0.84\,\text{s}$. Que os geradores reais sejam construídos exatamente assim não foi mostrado. & \cite{Matsuoka1985osc} & modelo \\
\bottomrule
\end{longtable}}

\section{Capítulo~\ref{sec:pain}. A dor}

Os sensores de dor, os dois fios, a porta na medula espinhal, o botão de volume descendente, a sensibilização e a captura da atenção foram medidos diretamente; as fórmulas da porta e da sensibilização são modelos simplificados do livro; a regra pela qual a máquina escolhe o volume do alarme é hipótese.

{\footnotesize
\setlength{\tabcolsep}{3pt}\renewcommand{\arraystretch}{1.15}
\begin{longtable}{>{\raggedright}p{0.5cm}>{\raggedright}p{4.0cm}>{\raggedright}p{5.6cm}>{\raggedright}p{2.3cm}>{\raggedright\arraybackslash}p{1.7cm}}
\toprule
N.º & Proposição & Experimento e o que foi medido & Fonte & Status \\
\midrule\endfirsthead
\toprule
N.º & Proposição & Experimento e o que foi medido & Fonte & Status \\
\midrule\endhead
1 & Limiar alto: os limiares de calor dos diferentes sensores de dor vão de $41^\circ$ a $46$--$48^\circ$ & Registro de fibras isoladas. Pele de macaco: limiar mediano dos sensores lentos (C) de $41^\circ$, dos sensores rápidos do segundo tipo, $46^\circ$, do primeiro tipo, acima de $53^\circ$. Pele humana: sensores C que respondem também à pressão, $40.7^\circ$; os que não respondem à pressão, $48^\circ$. & \cite{Treede1995heat, Weidner1999cfib} & medido \\
2 & Os sensores de dor se adaptam muito menos que os de toque e ficam mais sensíveis após uma lesão leve; o enfraquecimento após aquecimento forte foi medido num tipo & Queimadura leve da pele ($50^\circ$, $100\,\text{s}$): após $5$--$10$~min, o limiar de dor em humanos e o limiar dos sensores C no macaco caem $2$--$6^\circ$; os outros sensores da pele não mostram esse desvio. Nos sensores rápidos do segundo tipo, a resposta após aquecimento forte fica suprimida. A comparação da adaptação com os sensores de toque é uma estimativa geral, sem experimento próprio aqui. & \cite{LaMotte1982hyper, Treede1995heat} & medido \\
3 & Dois fios: rápido de $5$--$30$~m/s, lento de $0.5$--$2$~m/s; tempo de chegada $t=L/v$~\eqref{eq:paintime} & Velocidade de condução: sensores de dor rápidos do macaco, em média $15$ e $25$~m/s (dois tipos); sensores C humanos, $0.8$--$1.0$~m/s. Com $L=1$~m, isso dá dezenas de milissegundos e cerca de um segundo. & \cite{Treede1995heat, Weidner1999cfib} & medido \\
4 & A dor chega duas vezes: primeiro rápida e precisa, depois surda e longa & Tempo de reação ao aquecimento do antebraço humano: de $1100$ a $400\ms$ com o aumento da temperatura; o aquecimento forte da pele com pelos dá dor dupla, o da palma, não. A primeira dor só pode ser levada por fibras mais rápidas que $6$~m/s: pela latência, correspondem a ela os sensores rápidos do segundo tipo, que não existem na palma. A divisão de papéis (``onde'' e ``quão ruim'') é interpretação. & \cite{CampbellLaMotte1983first, Treede1995heat} & medido \\
5 & Porta: o neurônio de saída da medula recebe dor e toque, e o intermediário inibitório é excitado pelo toque (fig.~\ref{fig:gate}) & O esquema da porta foi proposto como teoria. Camundongos, marcação genética e remoção de grupos de neurônios: os neurônios excitatórios com somatostatina são necessários para a dor mecânica; os neurônios inibitórios com dinorfina impedem que a entrada dos sensores de toque chegue a eles; sem esses neurônios, um toque leve provoca dor. & \cite{MelzackWall1965, Duan2014} & medido \\
6 & O toque abafa a dor & Humanos, pulsos de laser dolorosos na mão: o toque simultâneo reduz a avaliação da dor e a capacidade de distinguir a energia dos pulsos; o efeito enfraquece com a distância entre o ponto do toque e o ponto da dor dentro de um mesmo segmento da pele. & \cite{Mancini2014touch} & medido \\
7 & Suposição da teoria da porta: uma dor forte desliga por si o intermediário inibitório, e a porta se escancara & No capítulo, marcado como suposição da teoria original. O trabalho em camundongos descreve como a porta impede o toque de entrar no canal da dor; o desligamento do intermediário pela dor não é mostrado nele. & \cite{MelzackWall1965} & hipótese \\
8 & Porta~\eqref{eq:gate}: limiar de $0.18$ sem toque, $0.86$ com toque, $0.11$ com $g=2$; com $\nu=1$ o toque reduz a saída de $1$ para $0.25$, com $\nu=2$ não tem efeito; o toque sozinho não passa (fig.~\ref{fig:gatecurves}) & O cálculo com \texttt{models/\allowbreak pain\_\allowbreak gate.py}, com os pesos do texto, reproduz todos esses números. & \texttt{models/\allowbreak pain\_\allowbreak gate.py} & modelo \\
9 & As vias descendentes do tronco mudam o ganho da porta nos dois sentidos & Rato, bulbo, registro durante aquecimento da cauda: alguns neurônios disparam antes da retirada (``on''), outros se calam (``off''); uma estimulação fraca dessa região suprime o reflexo. Revisão: os mesmos circuitos tanto facilitam quanto inibem a transmissão da dor, e os opioides agem por meio deles. & \cite{Fields1983rvm, Fields2004} & medido \\
10 & A expectativa de alívio abafa a dor pelo canal opioide & Pessoas com dor aguda: naquelas em que a expectativa de alívio reduziu a dor, o bloqueio dos receptores opioides devolveu a dor ao nível das demais; nas demais, o bloqueio não mudou a dor. & \cite{Levine1978} & medido \\
11 & O cérebro escolhe o volume do alarme pelo benefício da interrupção & Não há experimento em que a regra de escolha do volume tenha sido medida. & --- & hipótese \\
12 & Sensibilização: após uma lesão, a saída da porta cresce e o limiar cai, em parte por causa da medula; ela persiste depois que o estímulo lesivo cessa & Rato: após lesão da pele, o limiar do reflexo de flexão cai e a resposta cresce; a análise eletrofisiológica mostra que parte do aumento surge na própria medula. O estímulo lesivo provoca alterações duradouras da sensibilidade, que sobrevivem ao próprio estímulo. O capítulo não dá o tempo exato de decaimento. & \cite{Woolf1983} & medido \\
13 & A sensibilização~\eqref{eq:sensitization} é uma realimentação positiva; com um laço forte, o alarme pode travar depois da cura & Decorre do modelo~\eqref{eq:sensitization} (exercício no fim do capítulo). Não há experimento mostrando que a dor persistente após a cura seja um laço travado exatamente desse tipo. & dedução no capítulo & hipótese \\
14 & A dor é uma recompensa negativa, e o aprendizado com ela segue o erro de diferença temporal & Ressonância funcional, humanos, cadeias de prenúncios da dor: a atividade do estriado ventral e da ínsula anterior coincide com os sinais de aprendizado por diferença temporal. Macaco: parte dos neurônios \nt{DOPA} é excitada por um sopro de ar desagradável e pelo seu prenúncio, outra parte é inibida. & \cite{Seymour2004, Matsumoto2009} & medido \\
15 & A dor interrompe o trabalho em curso & Humanos, estímulo elétrico doloroso e testes sonoros: a resposta ao teste $100\ms$ após o início da dor fica bem mais lenta; após $1.5\,\text{s}$ já não há diferença em relação ao controle. Uma revisão reúne esses experimentos num modelo de captura da atenção. & \cite{Crombez1996disrupt, EcclestonCrombez1999} & medido \\
16 & A retirada se fecha na medula espinhal & Rato: neurônios das camadas profundas do corno dorsal reproduzem o padrão espacial do reflexo de retirada de músculos individuais; eles não podem ser excitados antidromicamente a partir da região cervical, ou seja, são intermediários espinhais. & \cite{Schouenborg1995wr} & medido \\
\bottomrule
\end{longtable}}

\section{Capítulo~\ref{sec:motorlearning}. O aprendizado motor}

A regra dos mínimos quadrados e a vantagem de um fio de erro separado são teoremas; a construção do circuito com fio de erro, o enfraquecimento por coincidência, o ajuste do traço ao atraso, o pós-efeito e o retorno espontâneo da habilidade foram medidos; que o circuito inteiro funcione como aprendizado supervisionado e construa um modelo interno é a hipótese principal; os números do modelo de dois processos são cálculos pelo modelo do livro.

{\footnotesize
\setlength{\tabcolsep}{3pt}\renewcommand{\arraystretch}{1.15}
\begin{longtable}{>{\raggedright}p{0.5cm}>{\raggedright}p{4.0cm}>{\raggedright}p{5.6cm}>{\raggedright}p{2.3cm}>{\raggedright\arraybackslash}p{1.7cm}}
\toprule
N.º & Afirmação & Experimento e o que foi medido & Fonte & Status \\
\midrule\endfirsthead
\toprule
N.º & Afirmação & Experimento e o que foi medido & Fonte & Status \\
\midrule\endhead
1 & A regra~\eqref{eq:lms} segue, em média, o gradiente do erro quadrático & Resultado da teoria dos filtros adaptativos: uma correção proporcional à entrada e ao erro da saída é uma descida estocástica pelo gradiente do erro quadrático médio. & \cite{WidrowHoff1960} & teoria \\
2 & Com um fio de erro para cada saída, a velocidade não depende do número de saídas; com um único número comum, cai com o número de neurônios ruidosos (proposição~\ref{prop:teachers}) & Curvas de aprendizado de uma rede linear: o aprendizado por perturbações aleatórias dos neurônios é mais lento que o gradiente exato tantas vezes quantos são os neurônios perturbados. & \cite{Werfel2005} & teoria \\
3 & O circuito com fio de erro funciona como aprendizado supervisionado (proposição~\ref{prop:errorwire}) & Teorias deduzidas da estrutura do circuito. Os experimentos das linhas 7--10 concordam com elas; que o circuito inteiro funcione exatamente assim não está provado. & \cite{Marr1969, Albus1971} & hipótese \\
4 & Camada de expansão: cerca de $7\cdot10^{10}$ pequenos neurônios \nt{GLUT}, mais que em todo o resto do cérebro & Contagem de núcleos numa suspensão de tecido dissolvido: no cerebelo humano há $69\cdot10^9$ neurônios dos $86\cdot10^9$ do cérebro inteiro. Estereologia: $101\cdot10^9$ pequenos neurônios no cerebelo de homens idosos. & \cite{Azevedo2009, Andersen1992cereb} & medido \\
5 & Aumentar a dimensão da entrada torna linear a dependência desejada & Teorema sobre o número de dicotomias: uma soma ponderada de $n$ entradas com limiar realiza uma rotulação arbitrária de cerca de $2n$ vetores em posição geral. Que o cerebelo use exatamente isso é parte da hipótese da linha~3. & \cite{Cover1965} & teoria \\
6 & Cerca de $3\cdot10^7$ neurônios \nt{GABA} de saída, cada um com da ordem de $10^5$ entradas & Estereologia do cerebelo humano: $30.5\cdot10^6$ neurônios de saída. Microscopia eletrônica, rato: cerca de $1.75\cdot10^5$ entradas da camada de expansão por neurônio de saída. O neurotransmissor inibitório dos neurônios de saída está na revisão. & \cite{Andersen1992cereb, NapperHarvey1988, Ito2001} & medido \\
7 & Exatamente um fio de erro por neurônio de saída, cerca de uma vez por segundo, cada vez uma descarga poderosa & Revisão de registros: cada neurônio de saída recebe uma entrada \nt{GLUT} trepadeira, que provoca um disparo complexo com frequência da ordem de $1$~Hz. & \cite{Ito2001} & medido \\
8 & A probabilidade da descarga depende da direção do erro de movimento & Macaco, região oculomotora do cerebelo, adaptação de sacadas: a direção do erro visual define a probabilidade do disparo complexo, e a magnitude, o seu momento; a descarga muda os disparos simples na tentativa seguinte e desloca o movimento ao longo do erro preferido da célula. & \cite{Herzfeld2018} & medido \\
9 & A coincidência da descarga de erro com uma entrada ativa enfraquece a entrada ($-\eta e_i x_j$) & Revisão de experimentos: a estimulação conjunta de uma entrada da camada de expansão e do fio de erro provoca um enfraquecimento de longo prazo dessa entrada; a estimulação de uma só entrada, não. & \cite{Ito2001} & medido \\
10 & O comprimento do traço é ajustado ao atraso do erro: com atraso de $100\ms$, enfraquece mais a entrada que esteve ativa $100\ms$ antes & Fatias e camundongos vivos: na região responsável pelo aprendizado oculomotor, o enfraquecimento é sintonizado estreitamente num intervalo de cerca de $120\ms$ entre entrada e erro, o tempo que o sinal de erro leva nessa tarefa; em outra região, células diferentes são sintonizadas em intervalos diferentes. & \cite{Suvrathan2016} & medido \\
11 & O circuito é um filtro adaptativo~\eqref{eq:adaptivefilter} que aprende o modelo interno do corpo & Modelo deduzido da estrutura do circuito. Não há experimento direto em que o filtro aprendido tenha sido medido como função de transferência do corpo. & \cite{Fujita1982} & hipótese \\
12 & A predição subtrai as consequências dos próprios movimentos, por isso não conseguimos fazer cócegas em nós mesmos & Humanos: o próprio toque parece menos cócegas que o mesmo toque alheio; na ressonância, o córtex somatossensorial responde mais fraco ao próprio toque, e o cerebelo fica menos ativo quando o movimento toca a pele. A explicação pela subtração da predição é hipótese. & \cite{Blakemore1998} & hipótese \\
13 & Com uma força externa, o erro diminui e, após a sua retirada, a mão erra para o outro lado & Pessoas movem a manopla de um robô num campo de forças: as trajetórias voltam a ser retas; depois da retirada do campo, o pós-efeito é quase um espelho das primeiras distorções; ele se transfere também para partes não treinadas do espaço. & \cite{ShadmehrMussaIvaldi1994} & medido \\
14 & Aprendem dois processos~\eqref{eq:tworate}; depois da perturbação inversa, a habilidade volta sozinha & Humanos, campo de forças, depois um campo inverso curto e tentativas com o erro travado: a correção volta sozinha à antiga; o modelo de dois processos reproduz isso, o de um só, não. & \cite{Smith2006} & medido \\
15 & Números da fig.~\ref{fig:tworate}: $0.84$ (lento $0.63$) em $200$ movimentos; $6$ inversos; rápido $-0.53$, lento $0.46$; retorno até $0.37$ em $40$ movimentos & Cálculo com \texttt{models/\allowbreak motor\_\allowbreak learning.py}, $A_f=0.92$, $B_f=0.1$, $A_s=0.996$, $B_s=0.02$: $0.840$ ($0.634$ e $0.205$), $6$ movimentos, $-0.531$ e $0.457$, pico de $0.371$ no $40$º movimento. & \texttt{models/\allowbreak motor\_\allowbreak learning.py} & modelo \\
16 & Dois processos são necessários porque o corpo muda em velocidades diferentes & Teoria: a estimativa ótima com perturbações de vários tempos de vida dá vários filtros com constantes de tempo diferentes e explica o retorno espontâneo e o reaprendizado acelerado. & \cite{Kording2007} & hipótese \\
\bottomrule
\end{longtable}}

\section{Capítulo~\ref{sec:chemoutput}. A saída química}

Os dois fios para o coração e suas velocidades diferentes, a realimentação negativa nas cascatas hormonais, o código pela frequência das porções, o gerador diário e seu ajuste pela luz foram medidos diretamente; o limite $K<8$ é um teorema da teoria de controle, deduzido no capítulo; as pulsações geradas pela própria realimentação da cascata são uma hipótese com um experimento forte a seu favor.

{\footnotesize
\setlength{\tabcolsep}{3pt}\renewcommand{\arraystretch}{1.15}
\begin{longtable}{>{\raggedright}p{0.5cm}>{\raggedright}p{4.0cm}>{\raggedright}p{5.6cm}>{\raggedright}p{2.3cm}>{\raggedright\arraybackslash}p{1.7cm}}
\toprule
N.º & Proposição & Experimento e o que foi medido & Fonte & Status \\
\midrule\endfirsthead
\toprule
N.º & Proposição & Experimento e o que foi medido & Fonte & Status \\
\midrule\endhead
1 & O marca-passo do coração bate sozinho cerca de $100$ vezes por minuto; em repouso o freio está pisado, e a frequência costuma ficar em torno de $60$--$70$ & Seres humanos, desligamento completo dos dois fios para o coração: a frequência própria é maior que a de repouso e diminui com a idade; em adultos, da ordem de $100$ batimentos por minuto. Logo, em repouso predomina o freio. A frequência de repouso de $60$--$70$ é um valor típico, sem citação à parte. & \cite{JoseCollison1970ihr} & medido \\
2 & O acelerador \nt{NORA} e o freio \nt{ACET} são filtros passa-baixas, e o freio é mais rápido~\eqref{eq:gasbrake} & Cão, estimulação modulada em frequência do nervo vago ou do simpático cardíaco e função de transferência até a frequência atrial: as duas vias são filtros passa-baixas; a simpática tem frequência de corte muito menor e um atraso puro de cerca de $1.7\,\text{s}$. As características dependem do tônus médio, de modo que a fórmula linear~\eqref{eq:gasbrake} é uma aproximação. & \cite{Berger1989} & medido \\
3 & O freio é rápido porque o \nt{ACET} abre o canal de potássio sem segundo mensageiro, e o \nt{NORA} age por uma cadeia de reações & Fragmentos de membrana de células atriais: o canal de potássio aberto pela acetilcolina é aberto pelas subunidades $\beta\gamma$ da proteína G ligada ao receptor, sem segundo mensageiro dentro da célula. A via do \nt{NORA} pelo AMPc não é citada aqui. & \cite{Logothetis1987gbg} & medido \\
4 & O sangue dá a volta no corpo em cerca de um minuto; o \nt{ADRE} da corrente sanguínea chega a todos os órgãos com receptor & Estimativa geral de ordem de grandeza; assim está formulada no capítulo, sem fonte citada. & --- & estimativa \\
5 & A cascata hormonal é envolvida por realimentação negativa: o hormônio final inibe os primeiros estágios & Revisão de experimentos: os hormônios do córtex adrenal suprimem a liberação de ACTH pela hipófise e do hormônio liberador pelo hipotálamo em três faixas de tempo: rápida, intermediária e lenta. & \cite{KellerWood1984fb} & medido \\
6 & Três estágios iguais são estáveis para $K<8$~\eqref{eq:cascadestable}; no limite, o período é $2\pi\tau/\sqrt3\approx3.6\tau$ & Deduzido no capítulo: na frequência $\sqrt3/\tau$, os três estágios dão defasagem de $180^\circ$ e atenuação de $8$ vezes. & dedução no capítulo & teoria \\
7 & O erro de nível é $1/(1+K)$, por isso esse regulador não tem precisão melhor que uns $10\,\%$ (proposição~\ref{prop:cascade}) & Consequência da linha~6 para realimentação proporcional. O eixo real tem realimentação em vários estágios e em várias faixas de tempo (linha~5), de modo que o limite se refere ao modelo~\eqref{eq:cascade} e não foi medido. & dedução no capítulo & teoria \\
8 & Com $K=4$ e $7$ o nível se estabiliza; com $K=12$, pulsações regulares com período de $3.7$--$3.9\,\tau$ (fig.~\ref{fig:cascade}) & Cálculo com \texttt{models/\allowbreak chem\_\allowbreak models.py}: com $K=12$, períodos sucessivos de $3.9$, $3.7$, $3.8$, $3.7\,\tau$; com $K=4$, amplitude de $0.003$ no fim do cálculo. & \texttt{models/\allowbreak chem\_\allowbreak models.py} & modelo \\
9 & O hormônio do estresse sai em pulsos cerca de uma vez por hora; os pulsos nascem da própria realimentação, sem gerador separado & Ratos, infusão constante de hormônio liberador: ACTH e corticosterona oscilam com frequência fisiológica, quase horária; uma entrada rítmica do hipotálamo não é necessária para os pulsos. Um modelo hipófise--adrenal com realimentação atrasada produz essas oscilações. & \cite{Walker2012glc, Walker2010} & hipótese \\
10 & O destinatário lê a frequência das porções, não o nível & Macacos sem liberação própria de hormônio liberador de gonadotrofinas: a infusão constante não restabelece a liberação dos hormônios da hipófise, as porções de hora em hora restabelecem; a volta à infusão constante suprime de novo a resposta. A fadiga do receptor como filtro passa-altas é uma interpretação. & \cite{Belchetz1978} & medido \\
11 & Frequências diferentes de porções agem de modo diferente sobre células diferentes da mesma glândula & Os mesmos macacos: aumentar as porções para $2$--$5$ por hora reduz os dois hormônios da hipófise; reduzi-las a uma a cada $3$~horas diminui um (LH) e aumenta o outro (FSH), de modo que a razão entre eles é definida pela frequência. & \cite{Wildt1981gnrh} & medido \\
12 & O ritmo diário nasce de uma realimentação negativa intracelular com atraso de algumas horas & Revisão: as proteínas dos genes do relógio suprimem, com atraso, a transcrição dos próprios genes; o laço de transcrição e tradução dá um período de cerca de um dia. & \cite{TakahashiJS2017} & medido \\
13 & O gerador principal é um grupo de dezenas de milhares de neurônios que sincroniza o resto do corpo & Revisão: o núcleo supraquiasmático é o gerador diário principal; seus neurônios isolados em cultura são geradores autônomos, e a rede os sincroniza; no camundongo, cerca de $10^4$ neurônios de cada lado. & \cite{Welsh2010scn} & medido \\
14 & O período próprio no ser humano é de $24.18$~horas & Pessoas sob luz fraca, com horário desvinculado do dia: o período dos ritmos de melatonina, temperatura e cortisol foi em média de $24.18$~horas em jovens e idosos, com dispersão estreita. & \cite{Czeisler1999} & medido \\
15 & A luz acerta o gerador como uma malha de captura de fase~\eqref{eq:pll}; é preciso um deslocamento de cerca de $11$~min por dia & Pessoas, um único pulso de luz intensa de $6.7$~horas em diferentes horas do dia: curva de resposta de fase do tipo~1 com amplitude de $5$~horas: atraso antes do mínimo da temperatura corporal, avanço depois. A equação~\eqref{eq:pll} é o modelo padrão; $11$~min $=0.18$~hora é aritmética. & \cite{Khalsa2003prc} & medido \\
16 & Sem luz não há ajuste, e o ritmo segue o período próprio & Pessoas totalmente cegas, sem percepção de luz, vivendo em sociedade normal: em cerca de metade delas os ritmos de melatonina e cortisol seguem um período próprio, que não coincide com o dia. & \cite{Sack1992blind} & medido \\
\bottomrule
\end{longtable}}

\section{Capítulo~\ref{sec:debugging}. Depurando a máquina}

O que o eletrodo ouve, as poucas dimensões da atividade, o desempenho das interfaces com matrizes rígidas, o aprendizado conjunto do cérebro e do decodificador e os pontos visuais produzidos por estimulação foram medidos em experimentos revisados por pares; as estimativas dos fluxos de dados são aritmética do capítulo; sobre o dispositivo da Neuralink implantado em pessoas, até onde foi possível verificar, os dados foram publicados sobretudo pela própria empresa.

{\footnotesize
\setlength{\tabcolsep}{3pt}\renewcommand{\arraystretch}{1.15}
\begin{longtable}{>{\raggedright}p{0.5cm}>{\raggedright}p{4.0cm}>{\raggedright}p{5.6cm}>{\raggedright}p{2.3cm}>{\raggedright\arraybackslash}p{1.7cm}}
\toprule
N.º & Proposição & Experimento e o que foi medido & Fonte & Status \\
\midrule\endfirsthead
\toprule
N.º & Proposição & Experimento e o que foi medido & Fonte & Status \\
\midrule\endhead
1 & O eletrodo ouve os disparos dos neurônios num raio da ordem de cem micrômetros, em geral de um a alguns; eles são distinguidos pela forma & Rato, campo CA1, registro intracelular e extracelular simultâneo: a forma do disparo extracelular reproduz a largura e a amplitude do intracelular. Estimativa: um tetrodo poderia distinguir $60$--$100$ neurônios, mas na prática se isolam cerca de seis. & \cite{Henze2000extra} & medido \\
2 & O cérebro tem $8.6\cdot10^{10}$ neurônios; a sonda ouve cerca de um centésimo de milionésimo & Contagem de núcleos numa suspensão de tecido dissolvido: $86\cdot10^9$ neurônios no cérebro humano. A fração é aritmética do capítulo. & \cite{Azevedo2009} & medido \\
3 & A atividade de centenas de neurônios do córtex motor cabe em uma ou duas dezenas de variáveis comuns & Revisão de registros no córtex motor de macacos: a atividade da população é descrita por poucas variáveis latentes, comuns a muitos neurônios. & \cite{Gallego2017} & medido \\
4 & O ruído dos vizinhos é correlacionado, e cem neurônios carregam quase tanto quanto mil (proposição~\ref{prop:pooling}) & Córtex visual de macaco: coeficiente de correlação do ruído de neurônios vizinhos em torno de $0.1$, e o ganho da média satura com uma centena de neurônios. Teoria das correlações que limitam a informação. & \cite{Zohary1994, MorenoBote2014} & medido \\
5 & Fluxo bruto de $2\cdot10^8$~bit/s contra $8\cdot10^4$~bit/s em disparos~\eqref{eq:probedata} & Aritmética do capítulo pela capacidade do fluxo de disparos~\eqref{eq:spikeentropy}. Os parâmetros de digitalização são reais: o chip da Neuralink usa $19.3$~kHz e $10$~bits por canal. & dedução no capítulo; \cite{Rieke1997, Musk2019} & teoria \\
6 & Primeiro sistema da Neuralink: os chips amplificam e digitalizam, o fluxo bruto sai por cabo, os disparos são detectados por um programa externo; detecção no chip, carcaça fechada, rádio e energia sem fio estão no dispositivo para pessoas, segundo a empresa & Artigo da empresa: o fluxo bruto completo sai por cabo USB-C, e os disparos são detectados em tempo real por um programa fora do chip; a compressão embarcada, a alimentação e a transmissão sem fio são citadas como plano para dispositivos clínicos. Para o dispositivo implantado em pessoas, só há comunicados da empresa. Que a detecção de disparos no chip é necessária é uma conclusão do capítulo a partir de~\eqref{eq:probedata}. & \cite{Musk2019}; relatórios da Neuralink, sem artigo revisado por pares & medido (artigo da empresa); em pessoas, relatório da empresa \\
7 & Até $3072$ eletrodos em $96$ fios de $32$; os fios são inseridos por um robô que desvia dos vasos & Artigo da empresa: $3072$ eletrodos em $96$ fios de $32$; o robô insere $6$ fios ($192$ eletrodos) por minuto; em ratos com matriz implantada por longo prazo, até $70\,\%$ dos eletrodos ouvem disparos. & \cite{Musk2019} & medido (artigo da empresa) \\
8 & Em pessoas, cerca de mil canais em $64$ fios; parte dos fios se deslocou; cursor da ordem de $10$~bit/s & Só comunicados da empresa. Uma busca no PubMed não encontrou artigo revisado por pares com resultados em pessoas; isso concorda com a ressalva no texto do capítulo. & relatórios da Neuralink; sem artigo revisado por pares & medido (relatório da empresa) \\
9 & Uma interface com matriz de cem eletrodos controla um cursor & Pessoa com paralisia, $96$ eletrodos no córtex motor primário: a intenção de mover o braço altera os disparos três anos após a lesão; o decodificador fornece um ``cursor neural'' (e-mail, televisão), controle de uma prótese de mão e de um braço robótico. & \cite{Hochberg2006} & medido \\
10 & Escrita com a ``mão mental'': cerca de $90$ caracteres por minuto, menos de $8$~bit/s & Pessoa com paralisia da mão, tentativas de escrever à mão, decodificador recorrente: $90$ caracteres por minuto com precisão de $94.1\,\%$ em tempo real, acima de $99\,\%$ após autocorreção. A estimativa em bits é aritmética do capítulo. & \cite{Willett2021} & medido \\
11 & Tentativas de falar são convertidas em texto a cerca de $60$ palavras por minuto & Pessoa que perdeu a fala inteligível, matrizes no córtex: $62$ palavras por minuto; $9.1\,\%$ de erros de palavra com vocabulário de $50$ palavras, $23.8\,\%$ com vocabulário de $125\,000$. & \cite{Willett2023} & medido \\
12 & A interface extrai uma pequena fração da capacidade: um neurônio dá dezenas de bits por segundo, cem dão milhares (proposição~\ref{prop:probe}) & O limite superior~\eqref{eq:spikeentropy} é teoria; a comparação com as linhas~8 e 10 é conclusão do capítulo. Que o gargalo sejam as poucas dimensões e o desconhecimento do código, e não o fio, não foi verificado por experimento direto. & \cite{Rieke1997} & teoria \\
13 & O cérebro e o decodificador aprendem um com o outro & Macacos controlam um cursor; o decodificador se ajusta em malha fechada, e os neurônios reaprendem ao mesmo tempo: a precisão aumenta, a habilidade se mantém e sofre menos com os movimentos do próprio braço. O conselho de separar as velocidades de aprendizado é uma regra de engenharia; o capítulo não traz um experimento específico. & \cite{Orsborn2014coad} & medido \\
14 & A corrente num eletrodo excita os neurônios e fios em volta sem distinguir tipo & Córtex de camundongo e gato, registro de cálcio por dois fótons: mesmo uma corrente fraca excita neurônios esparsos, espalhados a distâncias de até milímetros, por excitação direta de axônios num volume de dezenas de micrômetros de diâmetro. Ou seja, a excitação não é seletiva, mas também não é contínua. & \cite{Histed2009stim} & medido \\
15 & A estimulação do córtex visual acende pontos; até $16$ pontos ao mesmo tempo permitem reconhecer algumas letras e bordas de objetos & Pessoa cega, $96$ eletrodos no córtex visual por $6$ meses: limiar de um eletrodo de $66.8\pm36.5$~\textmu A; a estimulação simultânea de grupos (até $16$ eletrodos, o limite do estimulador) reduz o limiar e produz padrões distinguíveis; algumas letras e bordas de objetos são reconhecidas. & \cite{Fernandez2021} & medido \\
16 & A segunda linha da Neuralink é um dispositivo que envia sinais ao córtex visual & Só comunicados da empresa sobre o desenvolvimento; não foram encontrados dados revisados por pares sobre seu funcionamento. & relatórios da Neuralink & hipótese \\
17 & As concentrações dos neurotransmissores de difusão podem ser medidas por um sensor químico no tecido & Fatias de cérebro de rato, voltametria cíclica rápida com fibra de carbono: medidas a liberação e a recaptação de serotonina; a substância sai para além da sinapse. & \cite{Bunin1998} & medido \\
\bottomrule
\end{longtable}}

\section{Capítulo~\ref{sec:eetypes}. Neurônios como componentes}

Os modos de funcionamento de células individuais foram medidos diretamente, com corrente por eletrodo e registro de tensão; a matemática do integrador e da divisão em classes é teoria clássica, e o uso da reconfiguração de modos pelo cérebro para calcular continua sendo uma hipótese.

{\footnotesize
\setlength{\tabcolsep}{3pt}\renewcommand{\arraystretch}{1.15}
\begin{longtable}{>{\raggedright}p{0.5cm}>{\raggedright}p{4.0cm}>{\raggedright}p{5.6cm}>{\raggedright}p{2.3cm}>{\raggedright\arraybackslash}p{1.7cm}}
\toprule
N.º & Proposição & Experimento e o que foi medido & Fonte & Status \\
\midrule\endfirsthead
\toprule
N.º & Proposição & Experimento e o que foi medido & Fonte & Status \\
\midrule\endhead
1 & Ressonador: uma corrente lenta que se opõe à mudança de tensão torna o neurônio um filtro passa-faixa com máximo entre alguns hertz e dezenas de hertz (seção~\ref{sec:resonator}) & Em fatias, injetou-se em neurônios \nt{GLUT} corrente senoidal com frequência crescente e mediu-se a impedância $|Z(f)|$: máximo em $2$--$5$~Hz a $33^\circ$C (cerca de $7$~Hz a $38^\circ$C); a ressonância é criada por correntes lentas de potássio e catiônica e amplificada pela corrente persistente de sódio. Uma revisão mostra o mesmo recurso em muitas classes de células & \cite{HuVervaekeStorm2002,HutcheonYarom2000} & medido \\
2 & A corrente lenta se comporta como indutância e, com a capacitância da membrana, forma um circuito ressonante & A resposta sublimiar do axônio à corrente foi medida e descrita por equações de condutância linearizadas; o circuito equivalente contém uma indutância ``fenomenológica'' cujo valor depende da tensão & \cite{Mauro1970} & teoria \\
3 & Constante de tempo do grupo com realimentação $\tau_{\text{ef}}=\tau/(1-w)$~\eqref{eq:perfectint}; com $\tau=0.1$~s e $w=0.995$, $20$~s & Deduzida no capítulo a partir da equação linear do grupo; proposta como mecanismo de manutenção da posição do olho num trabalho teórico & dedução no capítulo; \cite{Seung1996} & teoria \\
4 & O integrador da posição do olho mantém a entrada pela rede, e não por uma propriedade de uma célula isolada & Registro intracelular em peixe, no núcleo que mantém a posição do olho: após um giro rápido, a taxa do neurônio muda em degrau e se mantém; uma corrente breve numa célula causa só um deslocamento passageiro, e os degraus de tensão persistem abaixo do limiar: são mantidos pelas entradas sinápticas & \cite{Aksay2001} & medido \\
5 & O integrador sem vazamento precisa de ajuste contínuo por aprendizado; desajustado, ele esquece ou satura & Peixes treinados com realimentação visual proporcional a $\pm$ a posição do olho: a fixação ficava instável ou com vazamento, e a constante de tempo dos neurônios caía mais de uma ordem de grandeza, para $<1$~s; a realimentação visual normal restabelecia aos poucos a estabilidade & \cite{Major2004} & medido \\
6 & Neurônio biestável: uma entrada curta liga um platô duradouro, a inibição o desliga; a propriedade é ligada pela serotonina (seção~\ref{sec:bistable}) & Registro intracelular de neurônios \nt{ACET} que controlam músculos, em gato: uma excitação sináptica curta leva o neurônio a funcionamento prolongado, uma inibição curta o reinicia; no gato espinhal, o estado aparece após a administração de um precursor da serotonina & \cite{Hounsgaard1988} & medido \\
7 & Duas classes: integradores (a taxa cresce a partir do zero) e ressonadores (no limiar, frequência finita de imediato) & As classes decorrem da teoria de bifurcações. Experimento: em fatias de córtex, neurônios \nt{GLUT} começam a disparar com taxa tão baixa quanto se queira (tipo~1), e os neurônios \nt{GABA} rápidos, com salto para uma frequência finita (tipo~2) & \cite{Izhikevich2007,Tateno2004} & medido \\
8 & Um neurônio é integrador ou detector de coincidência: a inibição encurta a janela de soma & Em fatias de hipocampo, a janela em que as entradas excitatórias se somam até o limiar é menor que $2$~ms; ela é encurtada pela inibição direta que chega logo após a excitação; nos dendritos, onde a inibição é mais fraca, a janela é mais larga & \cite{Pouille2001} & medido \\
9 & Linha de atraso feita de axônios (tabela~\ref{tab:eetypes}) & Na coruja-das-torres foram medidos os atrasos nos axônios que vão aos detectores de coincidência: comprimentos diferentes de axônio dão atrasos diferentes, e os detectores estão sintonizados com a diferença dos tempos de chegada do som & \cite{Carr1990} & medido \\
10 & Marca-passo na vigília e célula silenciosa no sono & Os neurônios \nt{SERO} funcionam de modo quase regular de dia, desaceleram no sono de ondas lentas e quase se calam no sono REM (mais detalhes no capítulo~\ref{sec:sleep}) & \cite{JacobsAzmitia1992,McGintyHarper1976} & medido \\
11 & O componente é definido pelos canais iônicos e pelos canais de difusão que os ajustam (proposição~\ref{prop:eetypes}) & Mostrado em redes pequenas: substâncias moduladoras alteram as propriedades intrínsecas dos neurônios e a força das sinapses, de modo que o mesmo circuito produz ritmos diferentes & \cite{Marder2012} & medido \\
12 & O cérebro usa a reconfiguração de modos para calcular (proposição~\ref{prop:eetypes}) & Não há experimento direto em que a reconfiguração do modo de uma célula seja necessária para resolver uma tarefa; há só experimentos em que os moduladores mudam a saída do circuito & \cite{Marder2012} & hipótese \\
\bottomrule
\end{longtable}}

\section{Capítulo~\ref{sec:dendrites}. Número de dendritos}

A estrutura das classes e o número de entradas foram medidos por anatomia e por registros elétricos; a estimativa~\eqref{eq:dendriteload} é física simples do capacitor, e a explicação computacional do número de entradas se apoia em teoria e modelos.

{\footnotesize
\setlength{\tabcolsep}{3pt}\renewcommand{\arraystretch}{1.15}
\begin{longtable}{>{\raggedright}p{0.5cm}>{\raggedright}p{4.0cm}>{\raggedright}p{5.6cm}>{\raggedright}p{2.3cm}>{\raggedright\arraybackslash}p{1.7cm}}
\toprule
N.º & Proposição & Experimento e o que foi medido & Fonte & Status \\
\midrule\endfirsthead
\toprule
N.º & Proposição & Experimento e o que foi medido & Fonte & Status \\
\midrule\endhead
1 & Capacitância específica da membrana $c_m\approx1$~\textmu F/cm$^2$ & Retirava-se membrana do corpo do neurônio com uma pipeta, aplicava-se um degrau de tensão e, pelo decaimento da corrente capacitiva, achava-se a capacitância total, depois dividida pela área: $0.9$~\textmu F/cm$^2$ em três classes de neurônios & \cite{Gentet2000} & medido \\
2 & Tempo de carga $t\approx c_mA\Delta V/I$~\eqref{eq:dendriteload}: o neurônio grande precisa de dezenas de entradas simultâneas, ao pequeno basta uma sinapse grande & Estimativa pela lei de carga do capacitor; os números ($20$ e $300$~pF, $0.1$ e alguns nA) são ordens de grandeza & dedução no capítulo & teoria \\
3 & Uma sinapse enorme dá corrente de alguns nanoampères e leva o neurônio pequeno ao limiar de imediato & Registro simultâneo do terminal e do destinatário numa fatia: corrente de uma sinapse de $2$--$13$~nA, cada disparo de entrada provoca resposta supralimiar, atraso de cerca de $0.4$~ms a $36^\circ$C; o destinatário emite \nt{GLYC} & \cite{Borst1995,BorstSoria2012} & medido \\
4 & Zero dendritos: nos neurônios sensoriais do tato e da dor, o disparo vai do sensor à medula espinhal sem passar pelo corpo celular & A estrutura (axônio que se divide em dois junto ao corpo) foi estabelecida anatomicamente. Que a condução não exige excitabilidade do corpo foi mostrado num modelo validado desse neurônio: sem corpo excitável, o disparo atravessa a bifurcação com a mesma confiabilidade & \cite{AmirDevor2003} & teoria \\
5 & Um dendrito: o neurônio olfatório carrega um tipo de receptor e transmite uma linha de entrada & Revisão de experimentos com genes de receptores: cada neurônio olfatório expressa um gene de receptor de uma grande família & \cite{Mombaerts2004} & medido \\
6 & Dois dendritos: nos detectores de direção do som, as entradas do ouvido esquerdo e do direito chegam a dendritos diferentes & Anatomia e registros em mamíferos, reunidos numa revisão: as entradas dos dois ouvidos são separadas em dois dendritos opostos & \cite{Grothe2010} & medido \\
7 & Separar as entradas em dois dendritos torna o ``E'' mais estrito & Mostrado em modelo: a saturação~\eqref{eq:shunt} num só dendrito impede que as entradas de um lado cheguem ao limiar, e a detecção de coincidências melhora. Não há experimento direto em que se tenha mudado a distribuição das entradas & \cite{AgmonSnir1998} & teoria \\
8 & Cerca de $7\cdot10^{10}$ pequenos neurônios \nt{GLUT} do circuito de aprendizado motor, com $\sim4$ entradas cada & Contagem de células por fracionamento isotrópico: nessa parte do cérebro humano há cerca de $69$ bilhões de neurônios. Registro de uma dessas células em rato vivo: ao toque chegam rajadas de correntes discretas de poucas entradas, e a célula dispara quando elas se somam & \cite{Azevedo2009,Chadderton2004} & medido \\
9 & Um elemento de limiar com $n$ entradas separa no máximo $P_{\max}\approx2n$ padrões aleatórios~\eqref{eq:cover}; para o neurônio de saída do circuito de aprendizado motor, o limite é $\sim2\cdot10^5$, e a estimativa realista, $\sim4\cdot10^4$ associações & O limite é um resultado clássico da geometria combinatória. A estimativa realista vem da teoria de capacidade com pesos só positivos e margem para o ruído, aplicada ao número medido de entradas desse neurônio & \cite{Cover1965,Brunel2004} & teoria \\
10 & Os misturadores de poucas entradas servem para expandir a entrada; ``quatro'' está perto do ótimo & Teoria: a dimensão da representação dada pela camada de misturadores é máxima aproximadamente com o número de entradas observado na anatomia; a hipótese original da expansão está nos trabalhos clássicos & \cite{Marr1969,Albus1971,LitwinKumar2017} & hipótese \\
11 & Árvores grandes: $10^4$--$10^5$ entradas & Contagem de sinapses em imagens de microscopia eletrônica: $\approx1.7\cdot10^5$ sinapses num neurônio \nt{GABA} de saída do circuito de aprendizado motor em rato; cerca de $30\,000$ entradas excitatórias e $1\,700$ inibitórias num neurônio \nt{GLUT} do hipocampo & \cite{NapperHarvey1988,Megias2001} & medido \\
12 & Os ramos de uma árvore grande são subcircuitos separados com limiares próprios; o neurônio é uma pequena rede de várias camadas & Estimulação pontual de dois locais em dendritos finos: entradas no mesmo ramo se somam de forma sigmoide, em ramos diferentes, linearmente. Nos dendritos de neurônios do córtex humano foram encontrados disparos graduados de cálcio que permitem a uma só célula resolver uma tarefa linearmente não separável. Modelo: para reproduzir um neurônio é preciso uma rede profunda & \cite{Polsky2004,Gidon2020,Beniaguev2021} & medido \\
13 & Dendritos-saída: cada ramo do neurônio \nt{GABA} da retina é um detector de direção separado & Registro de cálcio por dois fótons em ramos isolados: o sinal no ramo é seletivo ao movimento do corpo celular para a ponta, e a tensão do corpo não & \cite{Euler2002} & medido \\
14 & O número de entradas segue a tarefa (proposição~\ref{prop:dendrites}) & A estrutura das classes foi medida (linhas 3--13); que o número de entradas foi escolhido pela velocidade ou pela classificação só foi mostrado por teoria e modelos, para algumas classes & \cite{LitwinKumar2017,AgmonSnir1998} & hipótese \\
\bottomrule
\end{longtable}}

\section{Capítulo~\ref{sec:sleep}. Sono}

Os modos do sono, os replays e a redução das sinapses foram medidos por registros de neurônios, microdiálise e microscopia eletrônica; o horário do sono e a gangorra lenta são descritos pelos modelos do livro, e a finalidade do sono é, em grande parte, hipótese.

{\footnotesize
\setlength{\tabcolsep}{3pt}\renewcommand{\arraystretch}{1.15}
\begin{longtable}{>{\raggedright}p{0.5cm}>{\raggedright}p{4.0cm}>{\raggedright}p{5.6cm}>{\raggedright}p{2.3cm}>{\raggedright\arraybackslash}p{1.7cm}}
\toprule
N.º & Proposição & Experimento e o que foi medido & Fonte & Status \\
\midrule\endfirsthead
\toprule
N.º & Proposição & Experimento e o que foi medido & Fonte & Status \\
\midrule\endhead
1 & No sono os neurônios do córtex não se calam; o que muda é o modo de funcionamento & Registros longos de neurônios do córtex em ratos: no sono de ondas lentas a taxa continua não nula, e períodos de atividade se alternam com períodos de silêncio; após vigília longa a taxa é maior em todos os estados, após o sono, menor & \cite{Vyazovskiy2009} & medido \\
2 & O \nt{ACET} é alto de dia e no sono REM, baixo no sono de ondas lentas (tabela~\ref{tab:sleepmodes}) & Microdiálise no córtex de gatos em livre movimento: na vigília, a liberação de acetilcolina é $100$--$175\%$ maior que no sono de ondas lentas, e no sono REM fica mais ou menos no nível da vigília tranquila & \cite{Marrosu1995} & medido \\
3 & \nt{NORA} e \nt{SERO} se aquietam no sono de ondas lentas e quase se calam no REM & Registro de neurônios \nt{NORA} isolados em ratos e de neurônios \nt{SERO} em gatos em cada estágio do sono: a taxa é máxima na vigília, menor no sono de ondas lentas e quase nula no REM & \cite{AstonJonesBloom1981,McGintyHarper1976} & medido \\
4 & No sono REM os músculos são desligados por inibição via \nt{GABA} e \nt{GLYC} & Em ratos, mediu-se a atividade dos neurônios \nt{ACET} do músculo masseter e aplicaram-se bloqueadores diretamente sobre eles: a paralisia só cessa se forem bloqueados tanto os receptores \nt{GABA} dos dois tipos quanto os receptores \nt{GLYC} & \cite{BrooksPeever2012} & medido \\
5 & A pressão do sono $S$ é um capacitor com dois limiares~\eqref{eq:sleeppressure}, $\tau_r=18.2$~h, $\tau_d=4.2$~h & Modelo de dois processos; as constantes de tempo foram obtidas de como a potência das ondas lentas do EEG cresce com a vigília e cai no sono, e a forma dos limiares, do horário do despertar espontâneo & \cite{Borbely1982,Daan1984} & teoria \\
6 & A máquina chega sozinha a $16.1$~h de vigília e $8.0$~h de sono (fig.~\ref{fig:twoprocess}) & Cálculo por~\eqref{eq:sleeppressure} com limiares oscilantes, script \texttt{models/sleep\_models.py}: $16.1$~h de vigília e $7.9$--$8.0$~h de sono & --- & modelo \\
7 & Após $40$~h sem dormir, $S=0.90$, mas o sono dura só $8.8$~h: a dívida é paga com profundidade, não com duração & Modelo: \texttt{models/sleep\_models.py} dá $S=0.90$ e $8.8$~h. Experimento: após $40.5$~h sem dormir, em oito pessoas, na primeira noite de recuperação, a potência das ondas lentas do EEG foi significativamente maior que o normal & \cite{Borbely1981} & modelo \\
8 & Fisicamente, $S$ é a adenosina & Microdiálise em gatos: a adenosina extracelular numa das regiões que controlam a vigília cresce durante a vigília, cresce ainda mais com a privação de sono e cai durante o sono. Que $S$ seja só adenosina não foi mostrado & \cite{PorkkaHeiskanen1997,Dunwiddie2001} & hipótese \\
9 & No sono de ondas lentas o córtex oscila: atividade por algumas centenas de ms, depois silêncio, menos de uma vez por segundo ($<1$~Hz, sob anestesia em geral $0.3$--$0.4$~Hz) & Registro intracelular em gatos anestesiados: despolarização de $0.8$--$1.5$~s e hiperpolarização longa, ritmo $<1$~Hz, em geral $0.3$--$0.4$~Hz; a mesma alternância de atividade e silêncio aparece no sono natural (linha~1) & \cite{Steriade1993,Vyazovskiy2009} & medido \\
10 & A gangorra é um grupo com realimentação e fadiga~\eqref{eq:updown}; com $b=5$, período de $1.1$~s (valor do modelo) e atividade em cerca de $35\%$ do tempo; com $b=1$, atividade contínua (fig.~\ref{fig:updown}) & Cálculo, script \texttt{models/sleep\_models.py}: período de $1.11$~s, fração de tempo ativo $0.35$ (média sobre um número inteiro de períodos); com $b=1$, fração ativa de $1.0$. O período do modelo é mais curto que o medido (linha~9) & --- & modelo \\
11 & O \nt{ACET} desliga a gangorra e leva o córtex ao funcionamento contínuo & A estimulação de um núcleo de neurônios \nt{ACET} em gato bloqueou o ritmo lento em $79\%$ dos neurônios do córtex e os levou a funcionamento contínuo, sobretudo pelo desaparecimento das hiperpolarizações longas. A acetilcolina suprime as correntes lentas de potássio que criam a fadiga & \cite{SteriadeAmzica1993,McCormick1992} & medido \\
12 & Os surtos de ritmo de $7$--$15$~Hz são gerados pela malha \nt{GLUT}--\nt{GABA} entre o córtex e um nó de retransmissão & Em fatias do nó de retransmissão visual de furão, os surtos são criados por rajadas de resposta das células de retransmissão à inibição vinda de um núcleo \nt{GABA} vizinho, que elas mesmas excitam & \cite{vonKrosigk1993,SteriadeMcCormickSejnowski1993} & medido \\
13 & Os replays do dia são acelerados: no hipocampo cerca de $20$ vezes, no córtex $6$--$7$ vezes & No sono de ondas lentas, as sequências repetidas de neurônios do hipocampo aparecem comprimidas no tempo. Depois de correr numa pista, as sequências de neurônios de lugar se repetiram na mesma ordem em surtos curtos de $\sim100$~ms, comprimidas cerca de $20$ vezes; no córtex, compressão de $6$--$7$ vezes & \cite{Nadasdy1999,LeeWilson2002,Euston2007} & medido \\
14 & Reforçar a coordenação dos replays com o córtex melhora a memória (relação causal) & Em ratos, após o treino, uma estimulação sincronizada com os replays reforçou sua ligação com as ondas lentas e os surtos de ritmo no córtex: no dia seguinte a memória era alta, e nos controles, no nível do acaso & \cite{Maingret2016} & medido \\
15 & A finalidade dos replays é o aprendizado intercalado da memória lenta & Teoria dos dois sistemas de aprendizado e cálculos em redes artificiais: o aprendizado sequencial estraga o antigo, o intercalado não. Não há experimento direto no cérebro mostrando que os replays servem justamente a isso & \cite{McClelland1995} & hipótese \\
16 & Após o sono, as sinapses diminuem proporcionalmente ao tamanho, e as grandes quase não mudam & Microscopia eletrônica tridimensional de $6920$ sinapses do córtex de camundongo: a área de contato após o sono é $\sim18\%$ menor que após a vigília, a redução é proporcional ao tamanho, e as sinapses grandes e estáveis não mudam & \cite{deVivo2017} & medido \\
17 & A tarefa principal do sono é a renormalização dos pesos & Hipótese da homeostase sináptica; as linhas 1 e 16 concordam com ela, mas não provam que seja a função principal & \cite{TononiCirelli2014} & hipótese \\
18 & O sono REM é um teste em bancada do modelo do mundo & O modo (tabela~\ref{tab:sleepmodes}) foi medido; que a rede execute o modelo do mundo é uma suposição teórica, sem experimento & \cite{Hobson2009} & hipótese \\
19 & Lavagem do cérebro no sono: questão em aberto & Um experimento: no sono, o espaço extracelular é $60\%$ mais largo e a limpeza, mais rápida. Outro, com outro método de medida: no sono e sob anestesia, a limpeza é bem mais lenta. Os resultados se contradizem & \cite{Xie2013,Miao2024} & hipótese \\
20 & Sono local: trechos do córtex de um animal acordado caem por instantes no silêncio & Em ratos, após vigília longa, os neurônios de um trecho isolado do córtex se calam brevemente, como no sono de ondas lentas, com uma onda lenta no EEG local; nesses momentos o rato erra mais na tarefa & \cite{Vyazovskiy2011} & medido \\
\bottomrule
\end{longtable}}

\section{Capítulo~\ref{sec:livingmachine}. A máquina de neurônios vivos}

O comportamento das culturas em matrizes de eletrodos foi medido em experimentos diretos com registro e estimulação; o mecanismo das salvas se apoia no modelo de esgotamento das sinapses e em experimentos com microculturas, e as afirmações sobre a dopamina na placa, em experimentos isolados, parte dos quais os próprios autores consideram apenas uma demonstração.

{\footnotesize
\setlength{\tabcolsep}{3pt}\renewcommand{\arraystretch}{1.15}
\begin{longtable}{>{\raggedright}p{0.5cm}>{\raggedright}p{4.0cm}>{\raggedright}p{5.6cm}>{\raggedright}p{2.3cm}>{\raggedright\arraybackslash}p{1.7cm}}
\toprule
N.º & Proposição & Experimento e o que foi medido & Fonte & Status \\
\midrule\endfirsthead
\toprule
N.º & Proposição & Experimento e o que foi medido & Fonte & Status \\
\midrule\endhead
1 & Cultura no chip: sobretudo \nt{GLUT}, com uma fração menor de \nt{GABA} que depende da preparação; em culturas de hipocampo, cerca de $6\%$ & Uma rede de milhares de neurônios sobre uma matriz de eletrodos é um experimento padrão (linha~3). A fração de neurônios \nt{GABA} em culturas de hipocampo foi contada por marcação para GABA: cerca de $6\%$ dos neurônios. Para culturas de córtex não damos um número verificado & \cite{Benson1994} & medido \\
2 & Organoides: tecido nervoso tridimensional com cerca de um milímetro, a partir de células-tronco humanas & A partir de células-tronco foram cultivados organoides tridimensionais com várias regiões cerebrais, entre elas um córtex com camadas de precursores e neurônios maduros & \cite{Lancaster2013} & medido \\
3 & Sem entrada, a cultura dispara em salvas com intervalos de um segundo a minutos, conforme a cultura & $58$ culturas de córtex com densidade de $3\,000$ a $50\,000$ neurônios foram registradas por cinco semanas ($963$ registros de meia hora): depois de duas semanas, em quase todas as culturas predominam salvas da rede inteira; os intervalos entre salvas vão de $1$ a $300$~s e dependem fortemente da preparação & \cite{Wagenaar2006} & medido \\
4 & A salva termina por esgotamento do neurotransmissor, intervalo $T_{\text{salva}}\approx\tau_{\text{rec}}\ln\frac{1}{1-x^*}$~\eqref{eq:burstinterval}; $2$--$4$~s é a estimativa do modelo com $\tau_{\text{rec}}=0.8$~s, a recuperação medida é mais lenta & A fórmula é deduzida no capítulo. Modelo: uma rede com sinapses que se esgotam gera sozinha salvas da rede inteira. Experimento em microculturas de $4$--$30$ neurônios: a duração da salva depende do tempo desde a anterior, e o esgotamento das vesículas de neurotransmissor elimina as salvas; conclusão dos autores: a salva é limitada pelo estoque de neurotransmissor. Lá a recuperação do estoque leva dezenas de segundos, o que, com a mesma fórmula, dá intervalos de dezenas de segundos e minutos & dedução no capítulo; \cite{Tsodyks2000,CohenSegal2011,Tsodyks1997} & teoria \\
5 & ``O professor que se cala'': a estimulação é desligada quando aparece a resposta desejada, e a resposta se fixa & A cultura foi estimulada a $0.3$--$1$~Hz até aparecer a resposta escolhida $50\pm10$~ms depois do estímulo; então a estimulação era desligada por $5$~min; depois de alguns ciclos a resposta desejada surgia de imediato; foram traçadas curvas de aprendizado & \cite{Shahaf2001} & medido \\
6 & A cultura joga tênis; aumento significativo das sequências na sessão só com realimentação completa & Em detalhe no capítulo~\ref{sec:dishbrain} e em seu suplemento: com silêncio no lugar do estímulo, a cultura também joga melhor no fim da sessão do que sem realimentação, mas com efeito menor; o aprendizado de uma sessão para outra é instável & \cite{Kagan2022} & medido \\
7 & Que a cultura aprende a prever sua entrada é hipótese; as palavras fortes sobre ``inteligência'' foram contestadas & O princípio da energia livre é uma proposta teórica; não há experimento direto que o distinga de explicações mais simples. Trinta pesquisadores, num artigo de resposta, apontaram que a mudança de comportamento na malha não mostra nem inteligência nem sensações & \cite{Friston2010,Balci2023} & hipótese \\
8 & A cultura conduz um corpo virtual até o alvo com um padrão de estímulos escolhido & O programa escolhia sozinho os padrões de estimulação de treino pelo resultado corrente; em dezenas de minutos a rede atingia os estados pedidos, e a plasticidade durava $80$~min após $2$~h de treino. Os mesmos estímulos, aplicados sem relação com o comportamento, não tinham efeito & \cite{Bakkum2008} & medido \\
9 & Reservatório: só se treina a leitura linear $y=W_{\text{saída}}\mathbf r$~\eqref{eq:reservoir} & Teoria da computação de reservatório: qualquer sistema não linear com memória que se apaga, mais uma saída linear treinada & \cite{Maass2002,JaegerHaas2004} & teoria \\
10 & O organoide como reservatório distingue falantes com acurácia de cerca de $78\%$ (segundo os autores); a leitura aprende pelo erro, e as conexões do organoide mudam sem professor & Sons de fala de vários falantes foram aplicados ao organoide como padrões de corrente na matriz de eletrodos; a leitura treinada distinguia o falante com acurácia de cerca de $78\%$, segundo os autores. Eles também relatam que durante o treino a conectividade funcional do organoide mudou, e chamam isso de aprendizado não supervisionado no próprio organoide & \cite{Cai2023} & medido \\
11 & Um milhão de neurônios gasta cerca de $10^{-4}$~W em disparos~\eqref{eq:dishpower} & Estimativa: o custo de um disparo, $10^{-10}$~J, do orçamento energético do córtex~\eqref{eq:energy}, multiplicado pelo número de disparos; não há medição direta da potência da cultura & dedução no capítulo; \cite{Attwell2001} & teoria \\
12 & Dopamina por clarão de luz: $1$~mM de dopamina enjaulada, $240$ repetições, o número de disparos evocados cresceu & Numa plataforma remota de organoides, dopamina numa gaiola fotossensível ($1$~mM) era liberada com ultravioleta ($365$~nm, $0.8$~s) quando o estímulo evocava mais disparos; em $240$ passos (cerca de $5$~h) o número de disparos em geral cresceu. Os autores escrevem que um único experimento não permite estabelecer a relação com a dopamina; não há controles & \cite{Jordan2024} & hipótese \\
13 & A dopamina de células vivas muda o peso de um transistor orgânico de forma duradoura e reversível & Células que liberam dopamina foram cultivadas sobre um elemento neuromórfico orgânico; a oxidação da dopamina junto ao eletrodo muda a condutância do canal; foram mostrados a mudança duradoura do peso e seu retorno & \cite{Keene2020} & medido \\
14 & Existem organoides com neurônios \nt{DOPA} funcionais; sua dopamina não foi usada como professor & Em organoides de células-tronco humanas foram encontrados neurônios \nt{DOPA} maduros e eletricamente ativos, e produção de dopamina. Não encontramos na literatura experimentos em que essa dopamina ensine outra rede & \cite{Jo2016} & medido \\
15 & O organoide equilibra a haste: os estímulos de ensino do programa são melhores que os aleatórios, sem consolidação, por meio das sinapses \nt{GLUT} & Organoides de córtex de camundongo em malha fechada com a tarefa do ``pêndulo no carrinho'': na maioria dos organoides, os sinais escolhidos por aprendizado por reforço deram resultado melhor que os aleatórios ou nenhum; após $45$~min de repouso a melhora não se mantinha; o bloqueio dos receptores AMPA e NMDA a anulava & \cite{Robbins2026} & medido \\
16 & Sem registradores de controle, a rede na bancada não consegue decidir sozinha o que conta como sucesso (proposição~\ref{prop:livingmachine}) & Em todos os experimentos das linhas 5--15, o sinal de ensino é definido pelo experimentador ou por um programa; não há experimento em que a cultura tenha criado sozinha um critério de sucesso; é uma conclusão pela ausência, não um experimento & --- & hipótese \\
\bottomrule
\end{longtable}}

\section{Capítulo~\ref{sec:dishbrain}. DishBrain}

A montagem da bancada e os resultados do jogo vêm de um único trabalho experimental e de sua continuação; as fórmulas do ruído e da deriva para boas configurações são deduzidas no capítulo e verificadas com o modelo do livro, e o mecanismo de aprendizado na placa não está estabelecido.

{\footnotesize
\setlength{\tabcolsep}{3pt}\renewcommand{\arraystretch}{1.15}
\begin{longtable}{>{\raggedright}p{0.5cm}>{\raggedright}p{4.0cm}>{\raggedright}p{5.6cm}>{\raggedright}p{2.3cm}>{\raggedright\arraybackslash}p{1.7cm}}
\toprule
N.º & Proposição & Experimento e o que foi medido & Fonte & Status \\
\midrule\endfirsthead
\toprule
N.º & Proposição & Experimento e o que foi medido & Fonte & Status \\
\midrule\endhead
1 & Bancada: cerca de $800\,000$ células de córtex de camundongo ou cerca de $10^6$ células humanas por chip; $26\,000$ eletrodos em $8$~mm$^2$, registro de até $1024$, estimulação por $8$ & Métodos do trabalho: chip MaxOne, $26\,000$ eletrodos de platina em $8$~mm$^2$, registro de $1024$ canais a $20\,000$ amostras por segundo; tecnicamente é possível estimular até $32$ eletrodos, de modo independente conseguiram $8$; as culturas de camundongo foram usadas a partir de cerca de $10$~dias & \cite{Kagan2022} & medido \\
2 & Malha com ciclo de $10$~ms: os disparos das regiões motoras movem a raquete (fig.~\ref{fig:dishloop}) & Os disparos são contados em janelas de $10$~ms; o atraso entre o disparo e o estímulo de resposta é de cerca de $5$~ms, definido pelo buffer do equipamento & \cite{Kagan2022} & medido \\
3 & Entrada: código de lugar (qual dos $8$ eletrodos) e de taxa, $4$--$40$~disparos/s & No experimento principal, a frequência de estimulação cresce linearmente de $4$~Hz, com a bola junto à parede do fundo, até $40$~Hz junto à raquete; a amplitude dos pulsos da bola é de $75$~mV; os pulsos são retangulares bifásicos & \cite{Kagan2022} & medido \\
4 & Saída $\Delta p=c\,(n_\uparrow-n_\downarrow)$~\eqref{eq:dishreadout}, ruído da contagem por janela $1/\sqrt{RT}$ & A regra da diferença entre duas regiões é descrita no trabalho (com pesos das regiões pela atividade), a fórmula exata não é dada. A estimativa do ruído é consequência da contagem de Poisson, deduzida no capítulo & \cite{Kagan2022}; dedução no capítulo & teoria \\
5 & Professor: após o acerto, estímulo previsível de $75$~mV, $100$~disparos/s, $0.1$~s; após o erro, imprevisível, $150$~mV, $5$~disparos/s, $4$~s em lugares e momentos aleatórios, depois $4$~s de pausa & Métodos: após o acerto, $75$~mV, $100$~Hz, $100$~ms em todos os $8$ eletrodos de uma vez; após o erro, $150$~mV, $5$~Hz em eletrodos aleatórios e momentos aleatórios durante $4$~s, depois $4$~s sem estimulação. Não há dopamina na cultura & \cite{Kagan2022} & medido \\
6 & A rede age de modo a tornar sua entrada previsível & O princípio da energia livre é uma proposta teórica da qual os autores derivaram o protocolo; o experimento é compatível com ele, mas não o distingue de mecanismos mais simples (linhas~7--8) & \cite{Friston2010,Kagan2022} & hipótese \\
7 & Se o fracasso sacode a rede e o sucesso não, o tempo numa configuração é $\rho\propto1/P_{\text{erro}}$~\eqref{eq:dishdensity} (proposição~\ref{prop:dishscramble}) & Decorre do balanço de fluxos numa cadeia de Markov com empurrões simétricos, deduzido no capítulo. Não há teste direto numa cultura & dedução no capítulo & teoria \\
8 & O modelo ``embaralhar no erro'' aprende: fração de acertos $0.21\to0.38$; com empurrão após cada rali $\approx0.12$, sem empurrões $0.19$ (fig.~\ref{fig:dishmodel}) & Cálculo, script \texttt{models/dishbrain\_model.py}, $40$ culturas-modelo com $2000$ ralis cada: $0.21\to0.38$; $0.13\to0.11$; $0.19\to0.19$ & --- & modelo \\
9 & As sequências de bolas rebatidas se alongam só nas culturas vivas na malha completa; os controles não aprendem & $399$ sessões de $20$~min: $9$ culturas de camundongo ($101$ sessões), $11$ humanas ($138$), $6$ chips com meio ($80$), $20$ culturas em repouso ($42$), $3$ simulações ($38$). A duração média do rali nos últimos $15$~min é maior que nos primeiros $5$ só nas culturas vivas; nas células que não são neurônios (HEK293T) e nos chips com meio não há efeito & \cite{Kagan2022} & medido \\
10 & Aumento significativo na sessão só com realimentação completa; com silêncio no lugar do estímulo, o jogo no fim da sessão é melhor que sem realimentação, mas o efeito é menor & $486$ sessões em $3$ dias: ``estímulo'' ($n=27$), ``silêncio'' ($17$), ``sem realimentação'' ($15$). O aumento ao longo do tempo só é significativo na condição ``estímulo''; no fim da sessão, ``silêncio'' superava ``sem realimentação'' com efeito menor & \cite{Kagan2022} & medido \\
11 & O efeito é pequeno; se a rede aprende de forma duradoura é questão em aberto & Dentro da sessão a melhora é sólida, mas, segundo os próprios autores, um aprendizado estável entre sessões ao longo de vários dias não foi observado & \cite{Kagan2022} & medido \\
12 & Durante o jogo a rede se aproxima do regime crítico, e quanto mais perto, melhor o jogo & Análise de avalanches de atividade nas mesmas culturas: a entrada estruturada aproxima a rede do regime crítico, e a qualidade do jogo se correlaciona com essa proximidade; mas a criticidade sozinha, sem informação sobre as consequências das ações, não basta para aprender & \cite{Habibollahi2023} & medido \\
13 & A ``senciência'' da cultura não foi demonstrada & Artigo de resposta de trinta pesquisadores: a mudança de comportamento numa malha fechada não prova nem inteligência nem sensações; não há experimento que o mostre & \cite{Balci2023} & hipótese \\
14 & Quarto controle proposto: estímulo previsível após o erro, com a mesma duração e energia (seção~\ref{sec:dishspec}) & Esse experimento não foi feito; é uma proposta do livro & --- & hipótese \\
\bottomrule
\end{longtable}}


\begin{thebibliography}{99}
\bibitem{Azevedo2009} F. A. C. Azevedo et al., \emph{Equal numbers of neuronal and nonneuronal cells make the human brain an isometrically scaled-up primate brain}, J. Comp. Neurol. \textbf{513}, 532 (2009).
\bibitem{Schultz1997} W. Schultz, P. Dayan, and P. R. Montague, \emph{A neural substrate of prediction and reward}, Science \textbf{275}, 1593 (1997).
\bibitem{SuttonBarto2018} R. S. Sutton and A. G. Barto, \emph{Reinforcement Learning: An Introduction}, 2nd ed. (MIT Press, Cambridge, MA, 2018).
\bibitem{Doya2002} K. Doya, \emph{Metalearning and neuromodulation}, Neural Networks \textbf{15}, 495 (2002).
\bibitem{Strata1999} P. Strata and R. Harvey, \emph{Dale's principle}, Brain Res. Bull. \textbf{50}, 349 (1999).
\bibitem{vanVreeswijk1996} C. van Vreeswijk and H. Sompolinsky, \emph{Chaos in neuronal networks with balanced excitatory and inhibitory activity}, Science \textbf{274}, 1724 (1996).
\bibitem{AstonJones2005} G. Aston-Jones and J. D. Cohen, \emph{An integrative theory of locus coeruleus-norepinephrine function: adaptive gain and optimal performance}, Annu. Rev. Neurosci. \textbf{28}, 403 (2005).
\bibitem{Beaulieu2011} J.-M. Beaulieu and R. R. Gainetdinov, \emph{The physiology, signaling, and pharmacology of dopamine receptors}, Pharmacol. Rev. \textbf{63}, 182 (2011).
\bibitem{Sparso2001} J. Spars{\o} and S. Furber (eds.), \emph{Principles of Asynchronous Circuit Design: A Systems Perspective} (Kluwer, Boston, 2001).
\bibitem{Carr1990} C. E. Carr and M. Konishi, \emph{A circuit for detection of interaural time differences in the brain stem of the barn owl}, J. Neurosci. \textbf{10}, 3227 (1990).
\bibitem{Wang2001} X.-J. Wang, \emph{Synaptic reverberation underlying mnemonic persistent activity}, Trends Neurosci. \textbf{24}, 455 (2001).
\bibitem{Hahnloser2002} R. H. R. Hahnloser, A. A. Kozhevnikov, and M. S. Fee, \emph{An ultra-sparse code underlies the generation of neural sequences in a songbird}, Nature \textbf{419}, 65 (2002).
\bibitem{Barlow1965} H. B. Barlow and W. R. Levick, \emph{The mechanism of directionally selective units in rabbit's retina}, J. Physiol. \textbf{178}, 477 (1965).
\bibitem{Carandini2012} M. Carandini and D. J. Heeger, \emph{Normalization as a canonical neural computation}, Nat. Rev. Neurosci. \textbf{13}, 51 (2012).
\bibitem{Pouille2001} F. Pouille and M. Scanziani, \emph{Enforcement of temporal fidelity in pyramidal cells by somatic feed-forward inhibition}, Science \textbf{293}, 1159 (2001).
\bibitem{Tsodyks1997isn} M. V. Tsodyks, W. E. Skaggs, T. J. Sejnowski, and B. L. McNaughton, \emph{Paradoxical effects of external modulation of inhibitory interneurons}, J. Neurosci. \textbf{17}, 4382 (1997).
\bibitem{Buzsaki2012} G. Buzs\'aki and X.-J. Wang, \emph{Mechanisms of gamma oscillations}, Annu. Rev. Neurosci. \textbf{35}, 203 (2012).
\bibitem{Rieke1997} F. Rieke, D. Warland, R. de Ruyter van Steveninck, and W. Bialek, \emph{Spikes: Exploring the Neural Code} (MIT Press, Cambridge, MA, 1997).
\bibitem{Mainen1995} Z. F. Mainen and T. J. Sejnowski, \emph{Reliability of spike timing in neocortical neurons}, Science \textbf{268}, 1503 (1995).
\bibitem{Softky1993} W. R. Softky and C. Koch, \emph{The highly irregular firing of cortical cells is inconsistent with temporal integration of random EPSPs}, J. Neurosci. \textbf{13}, 334 (1993).
\bibitem{Thorpe1996} S. Thorpe, D. Fize, and C. Marlot, \emph{Speed of processing in the human visual system}, Nature \textbf{381}, 520 (1996).
\bibitem{Zohary1994} E. Zohary, M. N. Shadlen, and W. T. Newsome, \emph{Correlated neuronal discharge rate and its implications for psychophysical performance}, Nature \textbf{370}, 140 (1994).
\bibitem{MorenoBote2014} R. Moreno-Bote et al., \emph{Information-limiting correlations}, Nat. Neurosci. \textbf{17}, 1410 (2014).
\bibitem{Gollisch2008} T. Gollisch and M. Meister, \emph{Rapid neural coding in the retina with relative spike latencies}, Science \textbf{319}, 1108 (2008).
\bibitem{OKeefe1993} J. O'Keefe and M. L. Recce, \emph{Phase relationship between hippocampal place units and the EEG theta rhythm}, Hippocampus \textbf{3}, 317 (1993).
\bibitem{Destexhe2003} A. Destexhe, M. Rudolph, and D. Par\'e, \emph{The high-conductance state of neocortical neurons in vivo}, Nat. Rev. Neurosci. \textbf{4}, 739 (2003).
\bibitem{Tsodyks1997} M. V. Tsodyks and H. Markram, \emph{The neural code between neocortical pyramidal neurons depends on neurotransmitter release probability}, Proc. Natl. Acad. Sci. USA \textbf{94}, 719 (1997).
\bibitem{Lisman1997} J. E. Lisman, \emph{Bursts as a unit of neural information: making unreliable synapses reliable}, Trends Neurosci. \textbf{20}, 38 (1997).
\bibitem{Attwell2001} D. Attwell and S. B. Laughlin, \emph{An energy budget for signaling in the grey matter of the brain}, J. Cereb. Blood Flow Metab. \textbf{21}, 1133 (2001).
\bibitem{Lennie2003} P. Lennie, \emph{The cost of cortical computation}, Curr. Biol. \textbf{13}, 493 (2003).
\bibitem{Yagishita2014} S. Yagishita et al., \emph{A critical time window for dopamine actions on the structural plasticity of dendritic spines}, Science \textbf{345}, 1616 (2014).
\bibitem{Werfel2005} J. Werfel, X. Xie, and H. S. Seung, \emph{Learning curves for stochastic gradient descent in linear feedforward networks}, Neural Comput. \textbf{17}, 2699 (2005).
\bibitem{Doya2000} K. Doya, \emph{Reinforcement learning in continuous time and space}, Neural Comput. \textbf{12}, 219 (2000).
\bibitem{Oja1982} E. Oja, \emph{Simplified neuron model as a principal component analyzer}, J. Math. Biol. \textbf{15}, 267 (1982).
\bibitem{BiPoo1998} G.-Q. Bi and M.-M. Poo, \emph{Synaptic modifications in cultured hippocampal neurons: dependence on spike timing, synaptic strength, and postsynaptic cell type}, J. Neurosci. \textbf{18}, 10464 (1998).
\bibitem{Williams1992} R. J. Williams, \emph{Simple statistical gradient-following algorithms for connectionist reinforcement learning}, Mach. Learn. \textbf{8}, 229 (1992).
\bibitem{Bayer2005} H. M. Bayer and P. W. Glimcher, \emph{Midbrain dopamine neurons encode a quantitative reward prediction error signal}, Neuron \textbf{47}, 129 (2005).
\bibitem{Shen2008} W. Shen, M. Flajolet, P. Greengard, and D. J. Surmeier, \emph{Dichotomous dopaminergic control of striatal synaptic plasticity}, Science \textbf{321}, 848 (2008).
\bibitem{Dabney2020} W. Dabney, Z. Kurth-Nelson, N. Uchida, C. K. Starkweather, D. Hassabis, R. Munos, and M. Botvinick, \emph{A distributional code for value in dopamine-based reinforcement learning}, Nature \textbf{577}, 671 (2020).
\bibitem{Matsumoto2009} M. Matsumoto and O. Hikosaka, \emph{Two types of dopamine neuron distinctly convey positive and negative motivational signals}, Nature \textbf{459}, 837 (2009).
\bibitem{BrombergMartin2010} E. S. Bromberg-Martin, M. Matsumoto, and O. Hikosaka, \emph{Dopamine in motivational control: rewarding, aversive, and alerting}, Neuron \textbf{68}, 815 (2010).
\bibitem{Menegas2018} W. Menegas, K. Akiti, R. Amo, N. Uchida, and M. Watabe-Uchida, \emph{Dopamine neurons projecting to the posterior striatum reinforce avoidance of threatening stimuli}, Nat. Neurosci. \textbf{21}, 1421 (2018).
\bibitem{Niv2007} Y. Niv, N. D. Daw, D. Joel, and P. Dayan, \emph{Tonic dopamine: opportunity costs and the control of response vigor}, Psychopharmacology \textbf{191}, 507 (2007).
\bibitem{Mohebi2019} A. Mohebi et al., \emph{Dissociable dopamine dynamics for learning and motivation}, Nature \textbf{570}, 65 (2019).
\bibitem{Hamid2016} A. A. Hamid et al., \emph{Mesolimbic dopamine signals the value of work}, Nat. Neurosci. \textbf{19}, 117 (2016).
\bibitem{Kim2020} H. R. Kim, A. N. Malik et al., \emph{A unified framework for dopamine signals across timescales}, Cell \textbf{183}, 1600 (2020).
\bibitem{Jeong2022} H. Jeong et al., \emph{Mesolimbic dopamine release conveys causal associations}, Science \textbf{378}, eabq6740 (2022).
\bibitem{Amo2022} R. Amo et al., \emph{A gradual temporal shift of dopamine responses mirrors the progression of temporal difference error in machine learning}, Nat. Neurosci. \textbf{25}, 1082 (2022).
\bibitem{Takahashi2017} Y. K. Takahashi et al., \emph{Dopamine neurons respond to errors in the prediction of sensory features of expected rewards}, Neuron \textbf{95}, 1395 (2017).
\bibitem{Sharpe2017} M. J. Sharpe et al., \emph{Dopamine transients are sufficient and necessary for acquisition of model-based associations}, Nat. Neurosci. \textbf{20}, 735 (2017).
\bibitem{Gardner2018} M. P. H. Gardner, G. Schoenbaum, and S. J. Gershman, \emph{Rethinking dopamine as generalized prediction error}, Proc. R. Soc. B \textbf{285}, 20181645 (2018).
\bibitem{Engelhard2019} B. Engelhard et al., \emph{Specialized coding of sensory, motor and cognitive variables in VTA dopamine neurons}, Nature \textbf{570}, 509 (2019).
\bibitem{Clements1992} J. D. Clements, R. A. J. Lester, G. Tong, C. E. Jahr, and G. L. Westbrook, \emph{The time course of glutamate in the synaptic cleft}, Science \textbf{258}, 1498 (1992).
\bibitem{Hasselmo2006} M. E. Hasselmo, \emph{The role of acetylcholine in learning and memory}, Curr. Opin. Neurobiol. \textbf{16}, 710 (2006).
\bibitem{JacobsAzmitia1992} B. L. Jacobs and E. C. Azmitia, \emph{Structure and function of the brain serotonin system}, Physiol. Rev. \textbf{72}, 165 (1992).
\bibitem{Schiller1992} P. H. Schiller, \emph{The ON and OFF channels of the visual system}, Trends Neurosci. \textbf{15}, 86 (1992).
\bibitem{Grothe2010} B. Grothe, M. Pecka, and D. McAlpine, \emph{Mechanisms of sound localization in mammals}, Physiol. Rev. \textbf{90}, 983 (2010).
\bibitem{Markram2004} H. Markram et al., \emph{Interneurons of the neocortical inhibitory system}, Nat. Rev. Neurosci. \textbf{5}, 793 (2004).
\bibitem{Joshi2016} S. Joshi, Y. Li, R. M. Kalwani, and J. I. Gold, \emph{Relationships between pupil diameter and neuronal activity in the locus coeruleus, colliculi, and cingulate cortex}, Neuron \textbf{89}, 221 (2016).
\bibitem{ServanSchreiber1990} D. Servan-Schreiber, H. Printz, and J. D. Cohen, \emph{A network model of catecholamine effects: gain, signal-to-noise ratio, and behavior}, Science \textbf{249}, 892 (1990).
\bibitem{YuDayan2005} A. J. Yu and P. Dayan, \emph{Uncertainty, neuromodulation, and attention}, Neuron \textbf{46}, 681 (2005).
\bibitem{Nassar2012} M. R. Nassar et al., \emph{Rational regulation of learning dynamics by pupil-linked arousal systems}, Nat. Neurosci. \textbf{15}, 1040 (2012).
\bibitem{BouretSara2005} S. Bouret and S. J. Sara, \emph{Network reset: a simplified overarching theory of locus coeruleus noradrenaline function}, Trends Neurosci. \textbf{28}, 574 (2005).
\bibitem{Hebb1949} D. O. Hebb, \emph{The Organization of Behavior: A Neuropsychological Theory} (Wiley, New York, 1949).
\bibitem{MacKay1952} D. M. MacKay and W. S. McCulloch, \emph{The limiting information capacity of a neuronal link}, Bull. Math. Biophys. \textbf{14}, 127 (1952).
\bibitem{WilsonCowan1972} H. R. Wilson and J. D. Cowan, \emph{Excitatory and inhibitory interactions in localized populations of model neurons}, Biophys. J. \textbf{12}, 1 (1972).
\bibitem{AllenStevens1994} C. Allen and C. F. Stevens, \emph{An evaluation of causes for unreliability of synaptic transmission}, Proc. Natl. Acad. Sci. USA \textbf{91}, 10380 (1994).
\bibitem{Barlow1961} H. B. Barlow, \emph{Possible principles underlying the transformations of sensory messages}, in \emph{Sensory Communication}, ed. W. A. Rosenblith (MIT Press, Cambridge, MA, 1961), pp.~217--234.
\bibitem{Cardin2009} J. A. Cardin et al., \emph{Driving fast-spiking cells induces gamma rhythm and controls sensory responses}, Nature \textbf{459}, 663 (2009).
\bibitem{Lillicrap2020} T. P. Lillicrap, A. Santoro, L. Marris, C. J. Akerman, and G. Hinton, \emph{Backpropagation and the brain}, Nat. Rev. Neurosci. \textbf{21}, 335 (2020).
\bibitem{Fremaux2016} N. Fr\'emaux and W. Gerstner, \emph{Neuromodulated spike-timing-dependent plasticity, and theory of three-factor learning rules}, Front. Neural Circuits \textbf{9}, 85 (2016).
\bibitem{Haider2006} B. Haider, A. Duque, A. R. Hasenstaub, and D. A. McCormick, \emph{Neocortical network activity in vivo is generated through a dynamic balance of excitation and inhibition}, J. Neurosci. \textbf{26}, 4535 (2006).
\bibitem{HoltKoch1997} G. R. Holt and C. Koch, \emph{Shunting inhibition does not have a divisive effect on firing rates}, Neural Comput. \textbf{9}, 1001 (1997).
\bibitem{Chance2002} F. S. Chance, L. F. Abbott, and A. D. Reyes, \emph{Gain modulation from background synaptic input}, Neuron \textbf{35}, 773 (2002).
\bibitem{Ozeki2009} H. Ozeki, I. M. Finn, E. S. Schaffer, K. D. Miller, and D. Ferster, \emph{Inhibitory stabilization of the cortical network underlies visual surround suppression}, Neuron \textbf{62}, 578 (2009).
\bibitem{HasselmoBower1992} M. E. Hasselmo and J. M. Bower, \emph{Cholinergic suppression specific to intrinsic not afferent fiber synapses in rat piriform (olfactory) cortex}, J. Neurophysiol. \textbf{67}, 1222 (1992).
\bibitem{HasselmoAndersonBower1992} M. E. Hasselmo, B. P. Anderson, and J. M. Bower, \emph{Cholinergic modulation of cortical associative memory function}, J. Neurophysiol. \textbf{67}, 1230 (1992).
\bibitem{Hasselmo1999} M. E. Hasselmo, \emph{Neuromodulation: acetylcholine and memory consolidation}, Trends Cogn. Sci. \textbf{3}, 351 (1999).
\bibitem{Tsodyks1988} M. V. Tsodyks and M. V. Feigel'man, \emph{The enhanced storage capacity in neural networks with low activity level}, Europhys. Lett. \textbf{6}, 101 (1988).
\bibitem{KalmanBucy1961} R. E. Kalman and R. S. Bucy, \emph{New results in linear filtering and prediction theory}, J. Basic Eng. \textbf{83}, 95 (1961).
\bibitem{WilsonMcNaughton1994} M. A. Wilson and B. L. McNaughton, \emph{Reactivation of hippocampal ensemble memories during sleep}, Science \textbf{265}, 676 (1994).
\bibitem{BarnesSharp1999} N. M. Barnes and T. Sharp, \emph{A review of central 5-HT receptors and their function}, Neuropharmacology \textbf{38}, 1083 (1999).
\bibitem{Bunin1998} M. A. Bunin and R. M. Wightman, \emph{Quantitative evaluation of 5-hydroxytryptamine (serotonin) neuronal release and uptake: an investigation of extrasynaptic transmission}, J. Neurosci. \textbf{18}, 4854 (1998).
\bibitem{Sprouse1987} J. S. Sprouse and G. K. Aghajanian, \emph{Electrophysiological responses of serotoninergic dorsal raphe neurons to 5-HT$_{1A}$ and 5-HT$_{1B}$ agonists}, Synapse \textbf{1}, 3 (1987).
\bibitem{Sykova2008} E. Syková and C. Nicholson, \emph{Diffusion in brain extracellular space}, Physiol. Rev. \textbf{88}, 1277 (2008).
\bibitem{WilsonNicoll2001} R. I. Wilson and R. A. Nicoll, \emph{Endogenous cannabinoids mediate retrograde signalling at hippocampal synapses}, Nature \textbf{410}, 588 (2001).
\bibitem{Garthwaite2008} J. Garthwaite, \emph{Concepts of neural nitric oxide-mediated transmission}, Eur. J. Neurosci. \textbf{27}, 2783 (2008).
\bibitem{vandenPol2012} A. N. van den Pol, \emph{Neuropeptide transmission in brain circuits}, Neuron \textbf{76}, 98 (2012).
\bibitem{Dunwiddie2001} T. V. Dunwiddie and S. A. Masino, \emph{The role and regulation of adenosine in the central nervous system}, Annu. Rev. Neurosci. \textbf{24}, 31 (2001).
\bibitem{Skaggs1996} W. E. Skaggs and B. L. McNaughton, \emph{Replay of neuronal firing sequences in rat hippocampus during sleep following spatial experience}, Science \textbf{271}, 1870 (1996).
\bibitem{Baylor1979} D. A. Baylor, T. D. Lamb, and K.-W. Yau, \emph{Responses of retinal rods to single photons}, J. Physiol. \textbf{288}, 613 (1979).
\bibitem{Hudspeth1989} A. J. Hudspeth, \emph{How the ear's works work}, Nature \textbf{341}, 397 (1989).
\bibitem{NakaRushton1966} K. I. Naka and W. A. H. Rushton, \emph{S-potentials from colour units in the retina of fish (Cyprinidae)}, J. Physiol. \textbf{185}, 536 (1966).
\bibitem{Lichtsteiner2008} P. Lichtsteiner, C. Posch, and T. Delbruck, \emph{A $128\times128$ 120 dB 15 $\mu$s latency asynchronous temporal contrast vision sensor}, IEEE J. Solid-State Circuits \textbf{43}, 566 (2008).
\bibitem{Koch2006} K. Koch et al., \emph{How much the eye tells the brain}, Curr. Biol. \textbf{16}, 1428 (2006).
\bibitem{Robles2001} L. Robles and M. A. Ruggero, \emph{Mechanics of the mammalian cochlea}, Physiol. Rev. \textbf{81}, 1305 (2001).
\bibitem{Mombaerts2004} P. Mombaerts, \emph{Genes and ligands for odorant, vomeronasal and taste receptors}, Nat. Rev. Neurosci. \textbf{5}, 263 (2004).
\bibitem{WoodSlater2001} S. J. Wood and C. R. Slater, \emph{Safety factor at the neuromuscular junction}, Prog. Neurobiol. \textbf{64}, 393 (2001).
\bibitem{Henneman1957} E. Henneman, \emph{Relation between size of neurons and their susceptibility to discharge}, Science \textbf{126}, 1345 (1957).
\bibitem{HarrisWolpert1998} C. M. Harris and D. M. Wolpert, \emph{Signal-dependent noise determines motor planning}, Nature \textbf{394}, 780 (1998).
\bibitem{Hogan1984} N. Hogan, \emph{Adaptive control of mechanical impedance by coactivation of antagonist muscles}, IEEE Trans. Autom. Control \textbf{29}, 681 (1984).
\bibitem{McAuleyMarsden2000} J. H. McAuley and C. D. Marsden, \emph{Physiological and pathological tremors and rhythmic central motor control}, Brain \textbf{123}, 1545 (2000).
\bibitem{WolpertGhahramani2000} D. M. Wolpert and Z. Ghahramani, \emph{Computational principles of movement neuroscience}, Nat. Neurosci. \textbf{3}, 1212 (2000).
\bibitem{MarderBucher2001} E. Marder and D. Bucher, \emph{Central pattern generators and the control of rhythmic movements}, Curr. Biol. \textbf{11}, R986 (2001).
\bibitem{Miyazaki2014} K. W. Miyazaki et al., \emph{Optogenetic activation of dorsal raphe serotonin neurons enhances patience for future rewards}, Curr. Biol. \textbf{24}, 2033 (2014).
\bibitem{Fuglevand1993} A. J. Fuglevand, D. A. Winter, and A. E. Patla, \emph{Models of recruitment and rate coding organization in motor-unit pools}, J. Neurophysiol. \textbf{70}, 2470 (1993).
\bibitem{Curcio1990} C. A. Curcio, K. R. Sloan, R. E. Kalina, and A. E. Hendrickson, \emph{Human photoreceptor topography}, J. Comp. Neurol. \textbf{292}, 497 (1990).
\bibitem{Hayes1950} N. D. Hayes, \emph{Roots of the transcendental equation associated with a certain difference-differential equation}, J. London Math. Soc. \textbf{25}, 226 (1950).
\bibitem{MelzackWall1965} R. Melzack and P. D. Wall, \emph{Pain mechanisms: a new theory}, Science \textbf{150}, 971 (1965).
\bibitem{Fields2004} H. Fields, \emph{State-dependent opioid control of pain}, Nat. Rev. Neurosci. \textbf{5}, 565 (2004).
\bibitem{Levine1978} J. D. Levine, N. C. Gordon, and H. L. Fields, \emph{The mechanism of placebo analgesia}, Lancet \textbf{312}, 654 (1978).
\bibitem{Woolf1983} C. J. Woolf, \emph{Evidence for a central component of post-injury pain hypersensitivity}, Nature \textbf{306}, 686 (1983).
\bibitem{Seymour2004} B. Seymour et al., \emph{Temporal difference models describe higher-order learning in humans}, Nature \textbf{429}, 664 (2004).
\bibitem{EcclestonCrombez1999} C. Eccleston and G. Crombez, \emph{Pain demands attention: a cognitive-affective model of the interruptive function of pain}, Psychol. Bull. \textbf{125}, 356 (1999).
\bibitem{WidrowHoff1960} B. Widrow and M. E. Hoff, \emph{Adaptive switching circuits}, IRE WESCON Convention Record, part 4, 96 (1960).
\bibitem{Marr1969} D. Marr, \emph{A theory of cerebellar cortex}, J. Physiol. \textbf{202}, 437 (1969).
\bibitem{Albus1971} J. S. Albus, \emph{A theory of cerebellar function}, Math. Biosci. \textbf{10}, 25 (1971).
\bibitem{Cover1965} T. M. Cover, \emph{Geometrical and statistical properties of systems of linear inequalities with applications in pattern recognition}, IEEE Trans. Electron. Comput. \textbf{EC-14}, 326 (1965).
\bibitem{Ito2001} M. Ito, \emph{Cerebellar long-term depression: characterization, signal transduction, and functional roles}, Physiol. Rev. \textbf{81}, 1143 (2001).
\bibitem{Suvrathan2016} A. Suvrathan, H. L. Payne, and J. L. Raymond, \emph{Timing rules for synaptic plasticity matched to behavioral function}, Neuron \textbf{92}, 959 (2016).
\bibitem{Fujita1982} M. Fujita, \emph{Adaptive filter model of the cerebellum}, Biol. Cybern. \textbf{45}, 195 (1982).
\bibitem{Blakemore1998} S.-J. Blakemore, D. M. Wolpert, and C. D. Frith, \emph{Central cancellation of self-produced tickle sensation}, Nat. Neurosci. \textbf{1}, 635 (1998).
\bibitem{Smith2006} M. A. Smith, A. Ghazizadeh, and R. Shadmehr, \emph{Interacting adaptive processes with different timescales underlie short-term motor learning}, PLoS Biol. \textbf{4}, e179 (2006).
\bibitem{Kording2007} K. P. K\"ording, J. B. Tenenbaum, and R. Shadmehr, \emph{The dynamics of memory as a consequence of optimal adaptation to a changing body}, Nat. Neurosci. \textbf{10}, 779 (2007).
\bibitem{Yao2023} Z. Yao et al., \emph{A high-resolution transcriptomic and spatial atlas of cell types in the whole mouse brain}, Nature \textbf{624}, 317 (2023).
\bibitem{Sanzeni2020} A. Sanzeni, B. Akitake, H. C. Goldbach, C. E. Leedy, N. Brunel, and M. H. Histed, \emph{Inhibition stabilization is a widespread property of cortical networks}, eLife \textbf{9}, e54875 (2020).
\bibitem{Padamsey2022} Z. Padamsey, D. Katsanevaki, N. Dupuy, and N. L. Rochefort, \emph{Neocortex saves energy by reducing coding precision during food scarcity}, Neuron \textbf{110}, 280 (2022).
\bibitem{Stringer2019} C. Stringer, M. Pachitariu, N. Steinmetz, C. B. Reddy, M. Carandini, and K. D. Harris, \emph{Spontaneous behaviors drive multidimensional, brainwide activity}, Science \textbf{364}, eaav7893 (2019).
\bibitem{Bittner2017} K. C. Bittner, A. D. Milstein, C. Grienberger, S. Romani, and J. C. Magee, \emph{Behavioral time scale synaptic plasticity underlies CA1 place fields}, Science \textbf{357}, 1033 (2017).
\bibitem{WhittingtonBogacz2019} J. C. R. Whittington and R. Bogacz, \emph{Theories of error back-propagation in the brain}, Trends Cogn. Sci. \textbf{23}, 235 (2019).
\bibitem{Reimer2016} J. Reimer et al., \emph{Pupil fluctuations track rapid changes in adrenergic and cholinergic activity in cortex}, Nat. Commun. \textbf{7}, 13289 (2016).
\bibitem{Beniaguev2021} D. Beniaguev, I. Segev, and M. London, \emph{Single cortical neurons as deep artificial neural networks}, Neuron \textbf{109}, 2727 (2021).
\bibitem{Gidon2020} A. Gidon et al., \emph{Dendritic action potentials and computation in human layer 2/3 cortical neurons}, Science \textbf{367}, 83 (2020).
\bibitem{Gallego2022} G. Gallego et al., \emph{Event-based vision: a survey}, IEEE Trans. Pattern Anal. Mach. Intell. \textbf{44}, 154 (2022).
\bibitem{Berger1989} R. D. Berger, J. P. Saul, and R. J. Cohen, \emph{Transfer function analysis of autonomic regulation. I. Canine atrial rate response}, Am. J. Physiol. \textbf{256}, H142 (1989).
\bibitem{Walker2010} J. J. Walker, J. R. Terry, and S. L. Lightman, \emph{Origin of ultradian pulsatility in the hypothalamic--pituitary--adrenal axis}, Proc. R. Soc. B \textbf{277}, 1627 (2010).
\bibitem{Belchetz1978} P. E. Belchetz, T. M. Plant, Y. Nakai, E. J. Keogh, and E. Knobil, \emph{Hypophysial responses to continuous and intermittent delivery of hypothalamic gonadotropin-releasing hormone}, Science \textbf{202}, 631 (1978).
\bibitem{Czeisler1999} C. A. Czeisler et al., \emph{Stability, precision, and near-24-hour period of the human circadian pacemaker}, Science \textbf{284}, 2177 (1999).
\bibitem{TakahashiJS2017} J. S. Takahashi, \emph{Transcriptional architecture of the mammalian circadian clock}, Nat. Rev. Genet. \textbf{18}, 164 (2017).
\bibitem{Musk2019} E. Musk and Neuralink, \emph{An integrated brain-machine interface platform with thousands of channels}, J. Med. Internet Res. \textbf{21}, e16194 (2019).
\bibitem{Hochberg2006} L. R. Hochberg et al., \emph{Neuronal ensemble control of prosthetic devices by a human with tetraplegia}, Nature \textbf{442}, 164 (2006).
\bibitem{Willett2021} F. R. Willett et al., \emph{High-performance brain-to-text communication via handwriting}, Nature \textbf{593}, 249 (2021).
\bibitem{Willett2023} F. R. Willett et al., \emph{A high-performance speech neuroprosthesis}, Nature \textbf{620}, 1031 (2023).
\bibitem{Gallego2017} J. A. Gallego, M. G. Perich, L. E. Miller, and S. A. Solla, \emph{Neural manifolds for the control of movement}, Neuron \textbf{94}, 978 (2017).
\bibitem{Fernandez2021} E. Fern\'andez et al., \emph{Visual percepts evoked with an intracortical 96-channel microelectrode array inserted in human occipital cortex}, J. Clin. Invest. \textbf{131}, e151331 (2021).
\bibitem{Tritsch2012} N. X. Tritsch, J. B. Ding, and B. L. Sabatini, \emph{Dopaminergic neurons inhibit striatal output through non-canonical release of GABA}, Nature \textbf{490}, 262 (2012).
\bibitem{Zhang2015} S. Zhang et al., \emph{Dopaminergic and glutamatergic microdomains in a subset of rodent mesoaccumbens axons}, Nat. Neurosci. \textbf{18}, 386 (2015).
\bibitem{Mongillo2008} G. Mongillo, O. Barak, and M. Tsodyks, \emph{Synaptic theory of working memory}, Science \textbf{319}, 1543 (2008).
\bibitem{Stokes2015} M. G. Stokes, \emph{`Activity-silent' working memory in prefrontal cortex: a dynamic coding framework}, Trends Cogn. Sci. \textbf{19}, 394 (2015).
\bibitem{Smith2013} S. L. Smith, I. T. Smith, T. Branco, and M. H\"ausser, \emph{Dendritic spikes enhance stimulus selectivity in cortical neurons in vivo}, Nature \textbf{503}, 115 (2013).
\bibitem{Baden2016} T. Baden et al., \emph{The functional diversity of retinal ganglion cells in the mouse}, Nature \textbf{529}, 345 (2016).
\bibitem{Lottem2018} E. Lottem et al., \emph{Activation of serotonin neurons promotes active persistence in a probabilistic foraging task}, Nat. Commun. \textbf{9}, 1000 (2018).
\bibitem{Payeur2021} A. Payeur, J. Guerguiev, F. Zenke, B. A. Richards, and R. Naud, \emph{Burst-dependent synaptic plasticity can coordinate learning in hierarchical circuits}, Nat. Neurosci. \textbf{24}, 1010 (2021).
\bibitem{MehonicKenyon2022} A. Mehonic and A. J. Kenyon, \emph{Brain-inspired computing needs a master plan}, Nature \textbf{604}, 255 (2022).
\bibitem{Poe2020} G. R. Poe et al., \emph{Locus coeruleus: a new look at the blue spot}, Nat. Rev. Neurosci. \textbf{21}, 644 (2020).
\bibitem{Hangya2015} B. Hangya, S. P. Ranade, M. Lorenc, and A. Kepecs, \emph{Central cholinergic neurons are rapidly recruited by reinforcement feedback}, Cell \textbf{162}, 1155 (2015).
\bibitem{FernandezRuiz2019} A. Fern\'andez-Ruiz et al., \emph{Long-duration hippocampal sharp wave ripples improve memory}, Science \textbf{364}, 1082 (2019).
\bibitem{Herzfeld2018} D. J. Herzfeld, Y. Kojima, R. Soetedjo, and R. Shadmehr, \emph{Encoding of error and learning to correct that error by the Purkinje cells of the cerebellum}, Nat. Neurosci. \textbf{21}, 736 (2018).
\bibitem{Duan2014} B. Duan et al., \emph{Identification of spinal circuits transmitting and gating mechanical pain}, Cell \textbf{159}, 1417 (2014).
\bibitem{Tremblay2016} R. Tremblay, S. Lee, and B. Rudy, \emph{GABAergic interneurons in the neocortex: from cellular properties to circuits}, Neuron \textbf{91}, 260 (2016).
\bibitem{Ren2018} J. Ren et al., \emph{Anatomically defined and functionally distinct dorsal raphe serotonin sub-systems}, Cell \textbf{175}, 472 (2018).
\bibitem{Pfeffer2013} C. K. Pfeffer, M. Xue, M. He, Z. J. Huang, and M. Scanziani, \emph{Inhibition of inhibition in visual cortex: the logic of connections between molecularly distinct interneurons}, Nat. Neurosci. \textbf{16}, 1068 (2013).
\bibitem{Pi2013} H.-J. Pi et al., \emph{Cortical interneurons that specialize in disinhibitory control}, Nature \textbf{503}, 521 (2013).
\bibitem{Keller2018} G. B. Keller and T. D. Mrsic-Flogel, \emph{Predictive processing: a canonical cortical computation}, Neuron \textbf{100}, 424 (2018).
\bibitem{HubelWiesel1962} D. H. Hubel and T. N. Wiesel, \emph{Receptive fields, binocular interaction and functional architecture in the cat's visual cortex}, J. Physiol. \textbf{160}, 106 (1962).
\bibitem{Riesenhuber1999} M. Riesenhuber and T. Poggio, \emph{Hierarchical models of object recognition in cortex}, Nat. Neurosci. \textbf{2}, 1019 (1999).
\bibitem{Fukushima1980} K. Fukushima, \emph{Neocognitron: a self-organizing neural network model for a mechanism of pattern recognition unaffected by shift in position}, Biol. Cybern. \textbf{36}, 193 (1980).
\bibitem{Yamins2014} D. L. K. Yamins et al., \emph{Performance-optimized hierarchical models predict neural responses in higher visual cortex}, Proc. Natl. Acad. Sci. USA \textbf{111}, 8619 (2014).
\bibitem{Hung2005} C. P. Hung, G. Kreiman, T. Poggio, and J. J. DiCarlo, \emph{Fast readout of object identity from macaque inferior temporal cortex}, Science \textbf{310}, 863 (2005).
\bibitem{WiskottSejnowski2002} L. Wiskott and T. J. Sejnowski, \emph{Slow feature analysis: unsupervised learning of invariances}, Neural Comput. \textbf{14}, 715 (2002).
\bibitem{Foldiak1991} P. F\"oldi\'ak, \emph{Learning invariance from transformation sequences}, Neural Comput. \textbf{3}, 194 (1991).
\bibitem{LiDiCarlo2008} N. Li and J. J. DiCarlo, \emph{Unsupervised natural experience rapidly alters invariant object representation in visual cortex}, Science \textbf{321}, 1502 (2008).
\bibitem{RaoBallard1999} R. P. N. Rao and D. H. Ballard, \emph{Predictive coding in the visual cortex: a functional interpretation of some extra-classical receptive-field effects}, Nat. Neurosci. \textbf{2}, 79 (1999).
\bibitem{Kar2019} K. Kar, J. Kubilius, K. Schmidt, E. B. Issa, and J. J. DiCarlo, \emph{Evidence that recurrent circuits are critical to the ventral stream's execution of core object recognition behavior}, Nat. Neurosci. \textbf{22}, 974 (2019).
\bibitem{Szegedy2014} C. Szegedy et al., \emph{Intriguing properties of neural networks}, in \emph{Proc. ICLR} (2014), arXiv:1312.6199.
\bibitem{Weiss2002} Y. Weiss, E. P. Simoncelli, and E. H. Adelson, \emph{Motion illusions as optimal percepts}, Nat. Neurosci. \textbf{5}, 598 (2002).
\bibitem{Purves2011} D. Purves, W. T. Wojtach, and R. B. Lotto, \emph{Understanding vision in wholly empirical terms}, Proc. Natl. Acad. Sci. USA \textbf{108}, Suppl.~3, 15588 (2011).
\bibitem{LaferSousa2015} R. Lafer-Sousa, K. L. Hermann, and B. R. Conway, \emph{Striking individual differences in color perception uncovered by `the dress' photograph}, Curr. Biol. \textbf{25}, R545 (2015).
\bibitem{Wallisch2017} P. Wallisch, \emph{Illumination assumptions account for individual differences in the perceptual interpretation of a profoundly ambiguous stimulus in the color domain: ``The dress''}, J. Vis. \textbf{17}(4), 5 (2017).
\bibitem{Schwartz2009} O. Schwartz, T. J. Sejnowski, and P. Dayan, \emph{Perceptual organization in the tilt illusion}, J. Vis. \textbf{9}(4), 19 (2009).
\bibitem{SchillerCarvey2005} P. H. Schiller and C. E. Carvey, \emph{The Hermann grid illusion revisited}, Perception \textbf{34}, 1375 (2005).
\bibitem{Nijhawan1994} R. Nijhawan, \emph{Motion extrapolation in catching}, Nature \textbf{370}, 256 (1994).
\bibitem{Conway2005} B. R. Conway, A. Kitaoka, A. Yazdanbakhsh, C. C. Pack, and M. S. Livingstone, \emph{Neural basis for a powerful static motion illusion}, J. Neurosci. \textbf{25}, 5651 (2005).
\bibitem{OteroMillan2012} J. Otero-Millan, S. L. Macknik, and S. Martinez-Conde, \emph{Microsaccades and blinks trigger illusory rotation in the ``rotating snakes'' illusion}, J. Neurosci. \textbf{32}, 6043 (2012).
\bibitem{MorenoBote2007} R. Moreno-Bote, J. Rinzel, and N. Rubin, \emph{Noise-induced alternations in an attractor network model of perceptual bistability}, J. Neurophysiol. \textbf{98}, 1125 (2007).
\bibitem{Watanabe2018} E. Watanabe, A. Kitaoka, K. Sakamoto, M. Yasugi, and K. Tanaka, \emph{Illusory motion reproduced by deep neural networks trained for prediction}, Front. Psychol. \textbf{9}, 345 (2018).
\bibitem{GomezVilla2020} A. G\'omez-Villa, A. Mart\'in, J. Vazquez-Corral, M. Bertalm\'io, and J. Malo, \emph{Color illusions also deceive CNNs for low-level vision tasks: analysis and implications}, Vision Res. \textbf{176}, 156 (2020).
\bibitem{delCastillo1954} J. del Castillo and B. Katz, \emph{Quantal components of the end-plate potential}, J. Physiol. \textbf{124}, 560 (1954).
\bibitem{Nowak1984} L. Nowak, P. Bregestovski, P. Ascher, A. Herbet, and A. Prochiantz, \emph{Magnesium gates glutamate-activated channels in mouse central neurones}, Nature \textbf{307}, 462 (1984).
\bibitem{Malinow2002} R. Malinow and R. C. Malenka, \emph{AMPA receptor trafficking and synaptic plasticity}, Annu. Rev. Neurosci. \textbf{25}, 103 (2002).
\bibitem{Matsuzaki2004} M. Matsuzaki, N. Honkura, G. C. R. Ellis-Davies, and H. Kasai, \emph{Structural basis of long-term potentiation in single dendritic spines}, Nature \textbf{429}, 761 (2004).
\bibitem{Rudomin1999} P. Rudomin and R. F. Schmidt, \emph{Presynaptic inhibition in the vertebrate spinal cord revisited}, Exp. Brain Res. \textbf{129}, 1 (1999).
\bibitem{HutcheonYarom2000} B. Hutcheon and Y. Yarom, \emph{Resonance, oscillation and the intrinsic frequency preferences of neurons}, Trends Neurosci. \textbf{23}, 216 (2000).
\bibitem{Seung1996} H. S. Seung, \emph{How the brain keeps the eyes still}, Proc. Natl. Acad. Sci. USA \textbf{93}, 13339 (1996).
\bibitem{Hounsgaard1988} J. Hounsgaard, H. Hultborn, B. Jespersen, and O. Kiehn, \emph{Bistability of alpha-motoneurones in the decerebrate cat and in the acute spinal cat after intravenous 5-hydroxytryptophan}, J. Physiol. \textbf{405}, 345 (1988).
\bibitem{Izhikevich2007} E. M. Izhikevich, \emph{Dynamical Systems in Neuroscience: The Geometry of Excitability and Bursting} (MIT Press, Cambridge, MA, 2007).
\bibitem{AgmonSnir1998} H. Agmon-Snir, C. E. Carr, and J. Rinzel, \emph{The role of dendrites in auditory coincidence detection}, Nature \textbf{393}, 268 (1998).
\bibitem{BorstSoria2012} J. G. G. Borst and J. Soria van Hoeve, \emph{The calyx of Held synapse: from model synapse to auditory relay}, Annu. Rev. Physiol. \textbf{74}, 199 (2012).
\bibitem{Euler2002} T. Euler, P. B. Detwiler, and W. Denk, \emph{Directionally selective calcium signals in dendrites of starburst amacrine cells}, Nature \textbf{418}, 845 (2002).
\bibitem{AstonJonesBloom1981} G. Aston-Jones and F. E. Bloom, \emph{Activity of norepinephrine-containing locus coeruleus neurons in behaving rats anticipates fluctuations in the sleep-waking cycle}, J. Neurosci. \textbf{1}, 876 (1981).
\bibitem{BrooksPeever2012} P. L. Brooks and J. H. Peever, \emph{Identification of the transmitter and receptor mechanisms responsible for REM sleep paralysis}, J. Neurosci. \textbf{32}, 9785 (2012).
\bibitem{Borbely1982} A. A. Borb\'ely, \emph{A two process model of sleep regulation}, Hum. Neurobiol. \textbf{1}, 195 (1982).
\bibitem{PorkkaHeiskanen1997} T. Porkka-Heiskanen et al., \emph{Adenosine: a mediator of the sleep-inducing effects of prolonged wakefulness}, Science \textbf{276}, 1265 (1997).
\bibitem{Steriade1993} M. Steriade, A. N\'u\~nez, and F. Amzica, \emph{A novel slow ($<1$ Hz) oscillation of neocortical neurons in vivo: depolarizing and hyperpolarizing components}, J. Neurosci. \textbf{13}, 3252 (1993).
\bibitem{McCormick1992} D. A. McCormick, \emph{Neurotransmitter actions in the thalamus and cerebral cortex and their role in neuromodulation of thalamocortical activity}, Prog. Neurobiol. \textbf{39}, 337 (1992).
\bibitem{SteriadeMcCormickSejnowski1993} M. Steriade, D. A. McCormick, and T. J. Sejnowski, \emph{Thalamocortical oscillations in the sleeping and aroused brain}, Science \textbf{262}, 679 (1993).
\bibitem{Nadasdy1999} Z. N\'adasdy et al., \emph{Replay and time compression of recurring spike sequences in the hippocampus}, J. Neurosci. \textbf{19}, 9497 (1999).
\bibitem{Maingret2016} N. Maingret et al., \emph{Hippocampo-cortical coupling mediates memory consolidation during sleep}, Nat. Neurosci. \textbf{19}, 959 (2016).
\bibitem{McClelland1995} J. L. McClelland, B. L. McNaughton, and R. C. O'Reilly, \emph{Why there are complementary learning systems in the hippocampus and neocortex}, Psychol. Rev. \textbf{102}, 419 (1995).
\bibitem{TononiCirelli2014} G. Tononi and C. Cirelli, \emph{Sleep and the price of plasticity: from synaptic and cellular homeostasis to memory consolidation and integration}, Neuron \textbf{81}, 12 (2014).
\bibitem{deVivo2017} L. de Vivo et al., \emph{Ultrastructural evidence for synaptic scaling across the wake/sleep cycle}, Science \textbf{355}, 507 (2017).
\bibitem{Xie2013} L. Xie et al., \emph{Sleep drives metabolite clearance from the adult brain}, Science \textbf{342}, 373 (2013).
\bibitem{Miao2024} A. Miao et al., \emph{Brain clearance is reduced during sleep and anesthesia}, Nat. Neurosci. \textbf{27}, 1046 (2024).
\bibitem{Vyazovskiy2011} V. V. Vyazovskiy et al., \emph{Local sleep in awake rats}, Nature \textbf{472}, 443 (2011).
\bibitem{Lancaster2013} M. A. Lancaster et al., \emph{Cerebral organoids model human brain development and microcephaly}, Nature \textbf{501}, 373 (2013).
\bibitem{Jordan2024} F. D. Jordan et al., \emph{Open and remotely accessible Neuroplatform for research in wetware computing}, Front. Artif. Intell. \textbf{7}, 1376042 (2024).
\bibitem{Smirnova2023} L. Smirnova et al., \emph{Organoid intelligence (OI): the new frontier in biocomputing and intelligence-in-a-dish}, Front. Sci. \textbf{1}, 1017235 (2023).
\bibitem{Wagenaar2006} D. A. Wagenaar, J. Pine, and S. M. Potter, \emph{An extremely rich repertoire of bursting patterns during the development of cortical cultures}, BMC Neurosci. \textbf{7}, 11 (2006).
\bibitem{Shahaf2001} G. Shahaf and S. Marom, \emph{Learning in networks of cortical neurons}, J. Neurosci. \textbf{21}, 8782 (2001).
\bibitem{Kagan2022} B. J. Kagan et al., \emph{In vitro neurons learn and exhibit sentience when embodied in a simulated game-world}, Neuron \textbf{110}, 3952 (2022).
\bibitem{Friston2010} K. Friston, \emph{The free-energy principle: a unified brain theory?}, Nat. Rev. Neurosci. \textbf{11}, 127 (2010).
\bibitem{Balci2023} F. Balci et al., \emph{A response to claims of emergent intelligence and sentience in a dish}, Neuron \textbf{111}, 604 (2023).
\bibitem{Bakkum2008} D. J. Bakkum, Z. C. Chao, and S. M. Potter, \emph{Spatio-temporal electrical stimuli shape behavior of an embodied cortical network in a goal-directed learning task}, J. Neural Eng. \textbf{5}, 310 (2008).
\bibitem{Maass2002} W. Maass, T. Natschl\"ager, and H. Markram, \emph{Real-time computing without stable states: a new framework for neural computation based on perturbations}, Neural Comput. \textbf{14}, 2531 (2002).
\bibitem{JaegerHaas2004} H. Jaeger and H. Haas, \emph{Harnessing nonlinearity: predicting chaotic systems and saving energy in wireless communication}, Science \textbf{304}, 78 (2004).
\bibitem{Cai2023} H. Cai et al., \emph{Brain organoid reservoir computing for artificial intelligence}, Nat. Electron. \textbf{6}, 1032 (2023).
\bibitem{Habibollahi2023} F. Habibollahi, B. J. Kagan, A. N. Burkitt, and C. French, \emph{Critical dynamics arise during structured information presentation within embodied in vitro neuronal networks}, Nat. Commun. \textbf{14}, 5287 (2023).
\bibitem{Robbins2026} A. Robbins et al., \emph{Goal-directed learning in cortical organoids}, Cell Rep. \textbf{45}, 116984 (2026).
\bibitem{Keene2020} S. T. Keene et al., \emph{A biohybrid synapse with neurotransmitter-mediated plasticity}, Nat. Mater. \textbf{19}, 969 (2020).
\bibitem{Jo2016} J. Jo et al., \emph{Midbrain-like organoids from human pluripotent stem cells contain functional dopaminergic and neuromelanin-producing neurons}, Cell Stem Cell \textbf{19}, 248 (2016).
\bibitem{Hendry1987} S. H. Hendry, H. D. Schwark, E. G. Jones, and J. Yan, \emph{Numbers and proportions of GABA-immunoreactive neurons in different areas of monkey cerebral cortex}, J. Neurosci. \textbf{7}, 1503 (1987).
\bibitem{Tang2001synapses} Y. Tang, J. R. Nyengaard, D. M. G. De Groot, and H. J. G. Gundersen, \emph{Total regional and global number of synapses in the human brain neocortex}, Synapse \textbf{41}, 258 (2001).
\bibitem{Pakkenberg1997} B. Pakkenberg and H. J. G. Gundersen, \emph{Neocortical neuron number in humans: effect of sex and age}, J. Comp. Neurol. \textbf{384}, 312 (1997).
\bibitem{Matsuda2009} W. Matsuda, T. Furuta, K. C. Nakamura, H. Hioki, F. Fujiyama, R. Arai, and T. Kaneko, \emph{Single nigrostriatal dopaminergic neurons form widely spread and highly dense axonal arborizations in the neostriatum}, J. Neurosci. \textbf{29}, 444 (2009).
\bibitem{Pakkenberg1991} B. Pakkenberg, A. M{\o}ller, H. J. G. Gundersen, A. Mouritzen Dam, and H. Pakkenberg, \emph{The absolute number of nerve cells in substantia nigra in normal subjects and in patients with Parkinson's disease estimated with an unbiased stereological method}, J. Neurol. Neurosurg. Psychiatry \textbf{54}, 30 (1991).
\bibitem{Baker1990} K. G. Baker, G. M. Halliday, and I. T\"ork, \emph{Cytoarchitecture of the human dorsal raphe nucleus}, J. Comp. Neurol. \textbf{301}, 147 (1990).
\bibitem{Baker1991} K. G. Baker, G. M. Halliday, P. Halasz, J.-P. Hornung, L. B. Geffen, R. G. H. Cotton, and I. T\"ork, \emph{Cytoarchitecture of serotonin-synthesizing neurons in the pontine tegmentum of the human brain}, Synapse \textbf{7}, 301 (1991).
\bibitem{Airaksinen1991} M. S. Airaksinen, A. Paetau, L. Palj\"arvi, K. Reinikainen, P. Riekkinen, R. Suomalainen, and P. Panula, \emph{Histamine neurons in human hypothalamus: anatomy in normal and Alzheimer diseased brains}, Neuroscience \textbf{44}, 465 (1991).
\bibitem{Colquhoun1992} D. Colquhoun, P. Jonas, and B. Sakmann, \emph{Action of brief pulses of glutamate on AMPA/kainate receptors in patches from different neurones of rat hippocampal slices}, J. Physiol. \textbf{458}, 261 (1992).
\bibitem{DutarNicoll1988} P. Dutar and R. A. Nicoll, \emph{A physiological role for GABA$_B$ receptors in the central nervous system}, Nature \textbf{332}, 156 (1988).

\bibitem{Rauch2003} A. Rauch, G. La Camera, H.-R. L\"uscher, W. Senn, and S. Fusi, \emph{Neocortical pyramidal cells respond as integrate-and-fire neurons to in vivo-like input currents}, J. Neurophysiol. \textbf{90}, 1598 (2003).
\bibitem{KlumppEady1956} R. G. Klumpp and H. R. Eady, \emph{Some measurements of interaural time difference thresholds}, J. Acoust. Soc. Am. \textbf{28}, 859 (1956).
\bibitem{Brand2002} A. Brand, O. Behrend, T. Marquardt, D. McAlpine, and B. Grothe, \emph{Precise inhibition is essential for microsecond interaural time difference coding}, Nature \textbf{417}, 543 (2002).
\bibitem{Funahashi1989} S. Funahashi, C. J. Bruce, and P. S. Goldman-Rakic, \emph{Mnemonic coding of visual space in the monkey's dorsolateral prefrontal cortex}, J. Neurophysiol. \textbf{61}, 331 (1989).
\bibitem{WangMarkram2006} Y. Wang, H. Markram, P. H. Goodman, T. K. Berger, J. Ma, and P. S. Goldman-Rakic, \emph{Heterogeneity in the pyramidal network of the medial prefrontal cortex}, Nat. Neurosci. \textbf{9}, 534 (2006).

\bibitem{McCullochPitts1943} W. S. McCulloch and W. Pitts, \emph{A logical calculus of the ideas immanent in nervous activity}, Bull. Math. Biophys. \textbf{5}, 115 (1943).
\bibitem{Fried2002} S. I. Fried, T. A. M\"unch, and F. S. Werblin, \emph{Mechanisms and circuitry underlying directional selectivity in the retina}, Nature \textbf{420}, 411 (2002).

\bibitem{Hromadka2008} T. Hrom\'adka, M. R. DeWeese, and A. M. Zador, \emph{Sparse representation of sounds in the unanesthetized auditory cortex}, PLoS Biol. \textbf{6}, e16 (2008).
\bibitem{Abbott1997depression} L. F. Abbott, J. A. Varela, K. Sen, and S. B. Nelson, \emph{Synaptic depression and cortical gain control}, Science \textbf{275}, 221 (1997).
\bibitem{ShadlenNewsome1998} M. N. Shadlen and W. T. Newsome, \emph{The variable discharge of cortical neurons: implications for connectivity, computation, and information coding}, J. Neurosci. \textbf{18}, 3870 (1998).

\bibitem{TrulsonJacobs1979} M. E. Trulson and B. L. Jacobs, \emph{Raphe unit activity in freely moving cats: correlation with level of behavioral arousal}, Brain Res. \textbf{163}, 135 (1979).
\bibitem{GagnonParent2014} D. Gagnon and M. Parent, \emph{Distribution of VGLUT3 in highly collateralized axons from the rat dorsal raphe nucleus as revealed by single-neuron reconstructions}, PLoS One \textbf{9}, e87709 (2014).
\bibitem{VandermaelenAghajanian1983} C. P. Vandermaelen and G. K. Aghajanian, \emph{Electrophysiological and pharmacological characterization of serotonergic dorsal raphe neurons recorded extracellularly and intracellularly in rat brain slices}, Brain Res. \textbf{289}, 109 (1983).
\bibitem{CourtneyFord2016} N. A. Courtney and C. P. Ford, \emph{Mechanisms of 5-HT$_{1A}$ receptor-mediated transmission in dorsal raphe serotonin neurons}, J. Physiol. \textbf{594}, 953 (2016).
\bibitem{Hashemi2012serotonin} P. Hashemi et al., \emph{Brain dopamine and serotonin differ in regulation and its consequences}, Proc. Natl. Acad. Sci. USA \textbf{109}, 11510 (2012).
\bibitem{Black1934feedback} H. S. Black, \emph{Stabilized feedback amplifiers}, Bell Syst. Tech. J. \textbf{13}, 1 (1934).
\bibitem{KobayashiSchultz2008} S. Kobayashi and W. Schultz, \emph{Influence of reward delays on responses of dopamine neurons}, J. Neurosci. \textbf{28}, 7837 (2008).
\bibitem{Schweighofer2008} N. Schweighofer et al., \emph{Low-serotonin levels increase delayed reward discounting in humans}, J. Neurosci. \textbf{28}, 4528 (2008).

\bibitem{Malenka1988calcium} R. C. Malenka, J. A. Kauer, R. S. Zucker, and R. A. Nicoll, \emph{Postsynaptic calcium is sufficient for potentiation of hippocampal synaptic transmission}, Science \textbf{242}, 81 (1988).
\bibitem{Bliss1973LTP} T. V. P. Bliss and T. L{\o}mo, \emph{Long-lasting potentiation of synaptic transmission in the dentate area of the anaesthetized rabbit following stimulation of the perforant path}, J. Physiol. \textbf{232}, 331 (1973).
\bibitem{Kelso1986hebb} S. R. Kelso, A. H. Ganong, and T. H. Brown, \emph{Hebbian synapses in hippocampus}, Proc. Natl. Acad. Sci. USA \textbf{83}, 5326 (1986).
\bibitem{He2015eligibility} K. He et al., \emph{Distinct eligibility traces for LTP and LTD in cortical synapses}, Neuron \textbf{88}, 528 (2015).
\bibitem{Olveczky2005} B. P. \"Olveczky, A. S. Andalman, and M. S. Fee, \emph{Vocal experimentation in the juvenile songbird requires a basal ganglia circuit}, PLoS Biol. \textbf{3}, e153 (2005).
\bibitem{TumerBrainard2007} E. C. Tumer and M. S. Brainard, \emph{Performance variability enables adaptive plasticity of `crystallized' adult birdsong}, Nature \textbf{450}, 1240 (2007).
\bibitem{Eshel2015arithmetic} N. Eshel et al., \emph{Arithmetic and local circuitry underlying dopamine prediction errors}, Nature \textbf{525}, 243 (2015).
\bibitem{GraceOnn1989} A. A. Grace and S.-P. Onn, \emph{Morphology and electrophysiological properties of immunocytochemically identified rat dopamine neurons recorded in vitro}, J. Neurosci. \textbf{9}, 3463 (1989).
\bibitem{ODoherty2004actor} J. O'Doherty et al., \emph{Dissociable roles of ventral and dorsal striatum in instrumental conditioning}, Science \textbf{304}, 452 (2004).

\bibitem{NeweyPowell1987} W. K. Newey and J. L. Powell, \emph{Asymmetric least squares estimation and testing}, Econometrica \textbf{55}, 819 (1987).
\bibitem{Beierholm2013vigor} U. Beierholm et al., \emph{Dopamine modulates reward-related vigor}, Neuropsychopharmacology \textbf{38}, 1495 (2013).
\bibitem{Howe2013ramp} M. W. Howe, P. L. Tierney, S. G. Sandberg, P. E. M. Phillips, and A. M. Graybiel, \emph{Prolonged dopamine signalling in striatum signals proximity and value of distant rewards}, Nature \textbf{500}, 575 (2013).

\bibitem{Baker1989LC} K. G. Baker, I. T\"ork, J.-P. Hornung, and P. Halasz, \emph{The human locus coeruleus complex: an immunohistochemical and three dimensional reconstruction study}, Exp. Brain Res. \textbf{77}, 257 (1989).
\bibitem{Schwarz2015LC} L. A. Schwarz et al., \emph{Viral-genetic tracing of the input--output organization of a central noradrenaline circuit}, Nature \textbf{524}, 88 (2015).
\bibitem{Uematsu2017LC} A. Uematsu et al., \emph{Modular organization of the brainstem noradrenaline system coordinates opposing learning states}, Nat. Neurosci. \textbf{20}, 1602 (2017).
\bibitem{AstonJones1994vigilance} G. Aston-Jones, J. Rajkowski, P. Kubiak, and T. Alexinsky, \emph{Locus coeruleus neurons in monkey are selectively activated by attended cues in a vigilance task}, J. Neurosci. \textbf{14}, 4467 (1994).
\bibitem{Clayton2004LC} E. C. Clayton, J. Rajkowski, J. D. Cohen, and G. Aston-Jones, \emph{Phasic activation of monkey locus ceruleus neurons by simple decisions in a forced-choice task}, J. Neurosci. \textbf{24}, 9914 (2004).
\bibitem{Foote1975NE} S. L. Foote, R. Freedman, and A. P. Oliver, \emph{Effects of putative neurotransmitters on neuronal activity in monkey auditory cortex}, Brain Res. \textbf{86}, 229 (1975).
\bibitem{JepmaNieuwenhuis2011} M. Jepma and S. Nieuwenhuis, \emph{Pupil diameter predicts changes in the exploration--exploitation trade-off: evidence for the adaptive gain theory}, J. Cogn. Neurosci. \textbf{23}, 1587 (2011).
\bibitem{Tervo2014stochastic} D. G. R. Tervo et al., \emph{Behavioral variability through stochastic choice and its gating by anterior cingulate cortex}, Cell \textbf{159}, 21 (2014).
\bibitem{Usher1999LC} M. Usher, J. D. Cohen, D. Servan-Schreiber, J. Rajkowski, and G. Aston-Jones, \emph{The role of locus coeruleus in the regulation of cognitive performance}, Science \textbf{283}, 549 (1999).

\bibitem{Mesulam1983} M.-M. Mesulam, E. J. Mufson, B. H. Wainer, and A. I. Levey, \emph{Central cholinergic pathways in the rat: an overview based on an alternative nomenclature (Ch1--Ch6)}, Neuroscience \textbf{10}, 1185 (1983).
\bibitem{Marrosu1995} F. Marrosu et al., \emph{Microdialysis measurement of cortical and hippocampal acetylcholine release during sleep-wake cycle in freely moving cats}, Brain Res. \textbf{671}, 329 (1995).
\bibitem{Picciotto2012} M. R. Picciotto, M. J. Higley, and Y. S. Mineur, \emph{Acetylcholine as a neuromodulator: cholinergic signaling shapes nervous system function and behavior}, Neuron \textbf{76}, 116 (2012).
\bibitem{HasselmoSchnell1994} M. E. Hasselmo and E. Schnell, \emph{Laminar selectivity of the cholinergic suppression of synaptic transmission in rat hippocampal region CA1: computational modeling and brain slice physiology}, J. Neurosci. \textbf{14}, 3898 (1994).
\bibitem{Gil1997} Z. Gil, B. W. Connors, and Y. Amitai, \emph{Differential regulation of neocortical synapses by neuromodulators and activity}, Neuron \textbf{19}, 679 (1997).
\bibitem{KilgardMerzenich1998} M. P. Kilgard and M. M. Merzenich, \emph{Cortical map reorganization enabled by nucleus basalis activity}, Science \textbf{279}, 1714 (1998).
\bibitem{Atri2004} A. Atri et al., \emph{Blockade of central cholinergic receptors impairs new learning and increases proactive interference in a word paired-associate memory task}, Behav. Neurosci. \textbf{118}, 223 (2004).
\bibitem{Girardeau2009} G. Girardeau, K. Benchenane, S. I. Wiener, G. Buzs\'aki, and M. B. Zugaro, \emph{Selective suppression of hippocampal ripples impairs spatial memory}, Nat. Neurosci. \textbf{12}, 1222 (2009).

\input{supplement/bib_10}
\bibitem{CopenhagenJahr1989} D. R. Copenhagen and C. E. Jahr, \emph{Release of endogenous excitatory amino acids from turtle photoreceptors}, Nature \textbf{341}, 536 (1989).
\bibitem{GlowatzkiFuchs2002} E. Glowatzki and P. A. Fuchs, \emph{Transmitter release at the hair cell ribbon synapse}, Nat. Neurosci. \textbf{5}, 147 (2002).
\bibitem{Tinsley2016} J. N. Tinsley et al., \emph{Direct detection of a single photon by humans}, Nat. Commun. \textbf{7}, 12172 (2016).
\bibitem{Sellick1982} P. M. Sellick, R. Patuzzi, and B. M. Johnstone, \emph{Measurement of basilar membrane motion in the guinea pig using the M\"ossbauer technique}, J. Acoust. Soc. Am. \textbf{72}, 131 (1982).
\bibitem{Ruggero1997} M. A. Ruggero, N. C. Rich, A. Recio, S. S. Narayan, and L. Robles, \emph{Basilar-membrane responses to tones at the base of the chinchilla cochlea}, J. Acoust. Soc. Am. \textbf{101}, 2151 (1997).
\bibitem{NormannPerlman1979} R. A. Normann and I. Perlman, \emph{The effects of background illumination on the photoresponses of red and green cones}, J. Physiol. \textbf{286}, 491 (1979).
\bibitem{Laughlin1981} S. Laughlin, \emph{A simple coding procedure enhances a neuron's information capacity}, Z. Naturforsch. C \textbf{36}, 910 (1981).
\bibitem{CurcioAllen1990} C. A. Curcio and K. A. Allen, \emph{Topography of ganglion cells in human retina}, J. Comp. Neurol. \textbf{300}, 5 (1990).
\bibitem{Greenwood1990} D. D. Greenwood, \emph{A cochlear frequency-position function for several species --- 29 years later}, J. Acoust. Soc. Am. \textbf{87}, 2592 (1990).
\bibitem{Eguiluz2000} V. M. Egu\'iluz, M. Ospeck, Y. Choe, A. J. Hudspeth, and M. O. Magnasco, \emph{Essential nonlinearities in hearing}, Phys. Rev. Lett. \textbf{84}, 5232 (2000).
\bibitem{PalmerRussell1986} A. R. Palmer and I. J. Russell, \emph{Phase-locking in the cochlear nerve of the guinea-pig and its relation to the receptor potential of inner hair-cells}, Hear. Res. \textbf{24}, 1 (1986).
\bibitem{Malnic1999} B. Malnic, J. Hirono, T. Sato, and L. B. Buck, \emph{Combinatorial receptor codes for odors}, Cell \textbf{96}, 713 (1999).
\bibitem{Finger2005} T. E. Finger et al., \emph{ATP signaling is crucial for communication from taste buds to gustatory nerves}, Science \textbf{310}, 1495 (2005).
\bibitem{Huang2005} Y.-J. Huang et al., \emph{Mouse taste buds use serotonin as a neurotransmitter}, J. Neurosci. \textbf{25}, 843 (2005).
\bibitem{JohanssonVallbo1979} R. S. Johansson and A. B. Vallbo, \emph{Tactile sensibility in the human hand: relative and absolute densities of four types of mechanoreceptive units in glabrous skin}, J. Physiol. \textbf{286}, 283 (1979).
\bibitem{McKemy2002} D. D. McKemy, W. M. Neuhausser, and D. Julius, \emph{Identification of a cold receptor reveals a general role for TRP channels in thermosensation}, Nature \textbf{416}, 52 (2002).

\bibitem{Schmolesky1998} M. T. Schmolesky et al., \emph{Signal timing across the macaque visual system}, J. Neurophysiol. \textbf{79}, 3272 (1998).
\bibitem{Lampl2004} I. Lampl, D. Ferster, T. Poggio, and M. Riesenhuber, \emph{Intracellular measurements of spatial integration and the MAX operation in complex cells of the cat primary visual cortex}, J. Neurophysiol. \textbf{92}, 2704 (2004).
\bibitem{KobatakeTanaka1994} E. Kobatake and K. Tanaka, \emph{Neuronal selectivities to complex object features in the ventral visual pathway of the macaque cerebral cortex}, J. Neurophysiol. \textbf{71}, 856 (1994).
\bibitem{Albright1984} T. D. Albright, \emph{Direction and orientation selectivity of neurons in visual area MT of the macaque}, J. Neurophysiol. \textbf{52}, 1106 (1984).
\bibitem{Goodfellow2015} I. J. Goodfellow, J. Shlens, and C. Szegedy, \emph{Explaining and harnessing adversarial examples}, in \emph{Proc. ICLR} (2015), arXiv:1412.6572.
\bibitem{Elsayed2018} G. F. Elsayed et al., \emph{Adversarial examples that fool both computer vision and time-limited humans}, in \emph{Advances in Neural Information Processing Systems 31} (2018), arXiv:1802.08195.
\bibitem{Thompson1982} P. Thompson, \emph{Perceived rate of movement depends on contrast}, Vision Res. \textbf{22}, 377 (1982).
\bibitem{BarlowHill1963} H. B. Barlow and R. M. Hill, \emph{Evidence for a physiological explanation of the waterfall phenomenon and figural after-effects}, Nature \textbf{200}, 1345 (1963).
\bibitem{EaglemanSejnowski2000} D. M. Eagleman and T. J. Sejnowski, \emph{Motion integration and postdiction in visual awareness}, Science \textbf{287}, 2036 (2000).

\bibitem{Feinstein1955mu} B. Feinstein, B. Lindeg{\aa}rd, E. Nyman, and G. Wohlfart, \emph{Morphologic studies of motor units in normal human muscles}, Acta Anat. \textbf{23}, 127 (1955).
\bibitem{MilnerBrown1973rec} H. S. Milner-Brown, R. B. Stein, and R. Yemm, \emph{The orderly recruitment of human motor units during voluntary isometric contractions}, J. Physiol. \textbf{230}, 359 (1973).
\bibitem{Jones2002sdn} K. E. Jones, A. F. de C. Hamilton, and D. M. Wolpert, \emph{Sources of signal-dependent noise during isometric force production}, J. Neurophysiol. \textbf{88}, 1533 (2002).
\bibitem{GomiOsu1998stiff} H. Gomi and R. Osu, \emph{Task-dependent viscoelasticity of human multijoint arm and its spatial characteristics for interaction with environments}, J. Neurosci. \textbf{18}, 8965 (1998).
\bibitem{Jankowska1972Ia} E. Jankowska and W. J. Roberts, \emph{Synaptic actions of single interneurones mediating reciprocal Ia inhibition of motoneurones}, J. Physiol. \textbf{222}, 623 (1972).
\bibitem{Pruszynski2008rapid} J. A. Pruszynski, I. Kurtzer, and S. H. Scott, \emph{Rapid motor responses are appropriately tuned to the metrics of a visuospatial task}, J. Neurophysiol. \textbf{100}, 224 (2008).
\bibitem{SaundersKnill2003vis} J. A. Saunders and D. C. Knill, \emph{Humans use continuous visual feedback from the hand to control fast reaching movements}, Exp. Brain Res. \textbf{152}, 341 (2003).
\bibitem{FlanaganWing1997grip} J. R. Flanagan and A. M. Wing, \emph{The role of internal models in motion planning and control: evidence from grip force adjustments during movements of hand-held loads}, J. Neurosci. \textbf{17}, 1519 (1997).
\bibitem{Matsuoka1985osc} K. Matsuoka, \emph{Sustained oscillations generated by mutually inhibiting neurons with adaptation}, Biol. Cybern. \textbf{52}, 367 (1985).

\bibitem{Treede1995heat} R.-D. Treede, R. A. Meyer, S. N. Raja, and J. N. Campbell, \emph{Evidence for two different heat transduction mechanisms in nociceptive primary afferents innervating monkey skin}, J. Physiol. \textbf{483}, 747 (1995).
\bibitem{Weidner1999cfib} C. Weidner et al., \emph{Functional attributes discriminating mechano-insensitive and mechano-responsive C nociceptors in human skin}, J. Neurosci. \textbf{19}, 10184 (1999).
\bibitem{CampbellLaMotte1983first} J. N. Campbell and R. H. LaMotte, \emph{Latency to detection of first pain}, Brain Res. \textbf{266}, 203 (1983).
\bibitem{LaMotte1982hyper} R. H. LaMotte, J. G. Thalhammer, H. E. Torebj{\"o}rk, and C. J. Robinson, \emph{Peripheral neural mechanisms of cutaneous hyperalgesia following mild injury by heat}, J. Neurosci. \textbf{2}, 765 (1982).
\bibitem{Mancini2014touch} F. Mancini, T. Nash, G. D. Iannetti, and P. Haggard, \emph{Pain relief by touch: a quantitative approach}, Pain \textbf{155}, 635 (2014).
\bibitem{Fields1983rvm} H. L. Fields, J. Bry, I. Hentall, and G. Zorman, \emph{The activity of neurons in the rostral medulla of the rat during withdrawal from noxious heat}, J. Neurosci. \textbf{3}, 2545 (1983).
\bibitem{Crombez1996disrupt} G. Crombez, C. Eccleston, F. Baeyens, and P. Eelen, \emph{The disruptive nature of pain: an experimental investigation}, Behav. Res. Ther. \textbf{34}, 911 (1996).
\bibitem{Schouenborg1995wr} J. Schouenborg, H.-R. Weng, J. Kalliom{\"a}ki, and H. Holmberg, \emph{A survey of spinal dorsal horn neurones encoding the spatial organization of withdrawal reflexes in the rat}, Exp. Brain Res. \textbf{106}, 19 (1995).

\bibitem{Andersen1992cereb} B. B. Andersen, L. Korbo, and B. Pakkenberg, \emph{A quantitative study of the human cerebellum with unbiased stereological techniques}, J. Comp. Neurol. \textbf{326}, 549 (1992).
\bibitem{ShadmehrMussaIvaldi1994} R. Shadmehr and F. A. Mussa-Ivaldi, \emph{Adaptive representation of dynamics during learning of a motor task}, J. Neurosci. \textbf{14}, 3208 (1994).

\bibitem{JoseCollison1970ihr} A. D. Jose and D. Collison, \emph{The normal range and determinants of the intrinsic heart rate in man}, Cardiovasc. Res. \textbf{4}, 160 (1970).
\bibitem{Logothetis1987gbg} D. E. Logothetis, Y. Kurachi, J. Galper, E. J. Neer, and D. E. Clapham, \emph{The $\beta\gamma$ subunits of GTP-binding proteins activate the muscarinic K$^+$ channel in heart}, Nature \textbf{325}, 321 (1987).
\bibitem{KellerWood1984fb} M. E. Keller-Wood and M. F. Dallman, \emph{Corticosteroid inhibition of ACTH secretion}, Endocr. Rev. \textbf{5}, 1 (1984).
\bibitem{Walker2012glc} J. J. Walker et al., \emph{The origin of glucocorticoid hormone oscillations}, PLoS Biol. \textbf{10}, e1001341 (2012).
\bibitem{Wildt1981gnrh} L. Wildt et al., \emph{Frequency and amplitude of gonadotropin-releasing hormone stimulation and gonadotropin secretion in the rhesus monkey}, Endocrinology \textbf{109}, 376 (1981).
\bibitem{Welsh2010scn} D. K. Welsh, J. S. Takahashi, and S. A. Kay, \emph{Suprachiasmatic nucleus: cell autonomy and network properties}, Annu. Rev. Physiol. \textbf{72}, 551 (2010).
\bibitem{Khalsa2003prc} S. B. S. Khalsa, M. E. Jewett, C. Cajochen, and C. A. Czeisler, \emph{A phase response curve to single bright light pulses in human subjects}, J. Physiol. \textbf{549}, 945 (2003).
\bibitem{Sack1992blind} R. L. Sack, A. J. Lewy, M. L. Blood, L. D. Keith, and H. Nakagawa, \emph{Circadian rhythm abnormalities in totally blind people: incidence and clinical significance}, J. Clin. Endocrinol. Metab. \textbf{75}, 127 (1992).

\bibitem{Henze2000extra} D. A. Henze et al., \emph{Intracellular features predicted by extracellular recordings in the hippocampus in vivo}, J. Neurophysiol. \textbf{84}, 390 (2000).
\bibitem{Histed2009stim} M. H. Histed, V. Bonin, and R. C. Reid, \emph{Direct activation of sparse, distributed populations of cortical neurons by electrical microstimulation}, Neuron \textbf{63}, 508 (2009).
\bibitem{Orsborn2014coad} A. L. Orsborn et al., \emph{Closed-loop decoder adaptation shapes neural plasticity for skillful neuroprosthetic control}, Neuron \textbf{82}, 1380 (2014).

\bibitem{HuVervaekeStorm2002} H. Hu, K. Vervaeke, and J. F. Storm, \emph{Two forms of electrical resonance at theta frequencies, generated by M-current, h-current and persistent Na$^+$ current in rat hippocampal pyramidal cells}, J. Physiol. \textbf{545}, 783 (2002).
\bibitem{Mauro1970} A. Mauro, F. Conti, F. Dodge, and R. Schor, \emph{Subthreshold behavior and phenomenological impedance of the squid giant axon}, J. Gen. Physiol. \textbf{55}, 497 (1970).
\bibitem{Aksay2001} E. Aksay et al., \emph{In vivo intracellular recording and perturbation of persistent activity in a neural integrator}, Nat. Neurosci. \textbf{4}, 184 (2001).
\bibitem{Major2004} G. Major et al., \emph{Plasticity and tuning of the time course of analog persistent firing in a neural integrator}, Proc. Natl. Acad. Sci. USA \textbf{101}, 7745 (2004).
\bibitem{Tateno2004} T. Tateno, A. Harsch, and H. P. C. Robinson, \emph{Threshold firing frequency-current relationships of neurons in rat somatosensory cortex: type 1 and type 2 dynamics}, J. Neurophysiol. \textbf{92}, 2283 (2004).
\bibitem{Marder2012} E. Marder, \emph{Neuromodulation of neuronal circuits: back to the future}, Neuron \textbf{76}, 1 (2012).
\bibitem{McGintyHarper1976} D. J. McGinty and R. M. Harper, \emph{Dorsal raphe neurons: depression of firing during sleep in cats}, Brain Res. \textbf{101}, 569 (1976).

\bibitem{Gentet2000} L. J. Gentet, G. J. Stuart, and J. D. Clements, \emph{Direct measurement of specific membrane capacitance in neurons}, Biophys. J. \textbf{79}, 314 (2000).
\bibitem{Borst1995} J. G. G. Borst, F. Helmchen, and B. Sakmann, \emph{Pre- and postsynaptic whole-cell recordings in the medial nucleus of the trapezoid body of the rat}, J. Physiol. \textbf{489}, 825 (1995).
\bibitem{AmirDevor2003} R. Amir and M. Devor, \emph{Electrical excitability of the soma of sensory neurons is required for spike invasion of the soma, but not for through-conduction}, Biophys. J. \textbf{84}, 2181 (2003).
\bibitem{Chadderton2004} P. Chadderton, T. W. Margrie, and M. H\"ausser, \emph{Integration of quanta in cerebellar granule cells during sensory processing}, Nature \textbf{428}, 856 (2004).
\bibitem{LitwinKumar2017} A. Litwin-Kumar et al., \emph{Optimal degrees of synaptic connectivity}, Neuron \textbf{93}, 1153 (2017).
\bibitem{NapperHarvey1988} R. M. A. Napper and R. J. Harvey, \emph{Number of parallel fiber synapses on an individual Purkinje cell in the cerebellum of the rat}, J. Comp. Neurol. \textbf{274}, 168 (1988).
\bibitem{Megias2001} M. Meg\'{\i}as, Z. Emri, T. F. Freund, and A. I. Guly\'as, \emph{Total number and distribution of inhibitory and excitatory synapses on hippocampal CA1 pyramidal cells}, Neuroscience \textbf{102}, 527 (2001).
\bibitem{Polsky2004} A. Polsky, B. W. Mel, and J. Schiller, \emph{Computational subunits in thin dendrites of pyramidal cells}, Nat. Neurosci. \textbf{7}, 621 (2004).
\bibitem{Brunel2004} N. Brunel et al., \emph{Optimal information storage and the distribution of synaptic weights: perceptron versus Purkinje cell}, Neuron \textbf{43}, 745 (2004).

\bibitem{Vyazovskiy2009} V. V. Vyazovskiy et al., \emph{Cortical firing and sleep homeostasis}, Neuron \textbf{63}, 865 (2009).
\bibitem{Daan1984} S. Daan, D. G. Beersma, and A. A. Borb\'ely, \emph{Timing of human sleep: recovery process gated by a circadian pacemaker}, Am. J. Physiol. \textbf{246}, R161 (1984).
\bibitem{Borbely1981} A. A. Borb\'ely et al., \emph{Sleep deprivation: effect on sleep stages and EEG power density in man}, Electroencephalogr. Clin. Neurophysiol. \textbf{51}, 483 (1981).
\bibitem{SteriadeAmzica1993} M. Steriade, F. Amzica, and A. Nu\~nez, \emph{Cholinergic and noradrenergic modulation of the slow ($\approx0.3$ Hz) oscillation in neocortical cells}, J. Neurophysiol. \textbf{70}, 1385 (1993).
\bibitem{vonKrosigk1993} M. von Krosigk, T. Bal, and D. A. McCormick, \emph{Cellular mechanisms of a synchronized oscillation in the thalamus}, Science \textbf{261}, 361 (1993).
\bibitem{LeeWilson2002} A. K. Lee and M. A. Wilson, \emph{Memory of sequential experience in the hippocampus during slow wave sleep}, Neuron \textbf{36}, 1183 (2002).
\bibitem{Euston2007} D. R. Euston, M. Tatsuno, and B. L. McNaughton, \emph{Fast-forward playback of recent memory sequences in prefrontal cortex during sleep}, Science \textbf{318}, 1147 (2007).
\bibitem{Hobson2009} J. A. Hobson, \emph{REM sleep and dreaming: towards a theory of protoconsciousness}, Nat. Rev. Neurosci. \textbf{10}, 803 (2009).

\bibitem{CohenSegal2011} D. Cohen and M. Segal, \emph{Network bursts in hippocampal microcultures are terminated by exhaustion of vesicle pools}, J. Neurophysiol. \textbf{106}, 2314 (2011).
\bibitem{Tsodyks2000} M. Tsodyks, A. Uziel, and H. Markram, \emph{Synchrony generation in recurrent networks with frequency-dependent synapses}, J. Neurosci. \textbf{20}, RC50 (2000).
\bibitem{Benson1994} D. L. Benson, F. H. Watkins, O. Steward, and G. Banker, \emph{Characterization of GABAergic neurons in hippocampal cell cultures}, J. Neurocytol. \textbf{23}, 279 (1994).

\input{supplement/bib_22}
\end{thebibliography}


\end{document}
